\documentclass[11pt]{article}
% Shared preamble for all solution files.
\usepackage[margin=1.1in]{geometry}
\usepackage{enumitem}
\usepackage{tikz}
\usepackage[hidelinks]{hyperref}
\usepackage{amsmath,amssymb,amsthm}
\newcommand{\fa}{\mathfrak a}
\newcommand{\fb}{\mathfrak b}
\newcommand{\fc}{\mathfrak c}
\newcommand{\fp}{\mathfrak p}
\newcommand{\fq}{\mathfrak q}
\newcommand{\fm}{\mathfrak m}
\newcommand{\fn}{\mathfrak n}
\newcommand{\fr}{\mathfrak r}
\newcommand{\Nil}{\mathfrak N}
\newcommand{\Jac}{\mathfrak R}
\newcommand{\rad}{r}
\newcommand{\Spec}{\operatorname{Spec}}
\newcommand{\Max}{\operatorname{Max}}
\newcommand{\Ann}{\operatorname{Ann}}
\newcommand{\Ker}{\operatorname{Ker}}
\newcommand{\Z}{\mathbf Z}
\newcommand{\Q}{\mathbf Q}
\newcommand{\R}{\mathbf R}
\newcommand{\C}{\mathbf C}
\newcommand{\ci}{\mathrm{c}}
\newcommand{\ex}{\mathrm{e}}

\setlength{\parindent}{0pt}
\setlength{\parskip}{0.55em}
% \exercise{label}{statement}  -- followed by the solution text
\newcommand{\exercise}[2]{\par\bigskip\noindent{\large\textbf{Exercise #1}}\par\smallskip
  \begin{quote}\small\textbf{Problem.}\ #2\end{quote}\par}
\newcommand{\pf}{\noindent\textbf{Solution.}\ }

\title{Hartshorne, Chapter II: Schemes\\[2pt]\large Solutions to the exercises}
\author{}
\date{}
\begin{document}
\maketitle
\noindent\textbf{Conventions.} Rings are commutative with $1$; unlike in Chapter I, the base field is not assumed algebraically closed unless an exercise says so. Results of the text are cited by their numbers in Hartshorne (``Prop.~2.5'', ``Thm.~7.6'', \dots), and results of other chapters as ``I, \dots'' or ``III, \dots''. The numbering and statements of the exercises follow the edition scanned in this repository; each exercise is restated before its solution, and starred exercises are marked ``Starred in the book'' at the end of the statement.

Honest notes. Several printed statements need an extra hypothesis to be true, and the solutions say what is assumed: for example $\operatorname{char}k\neq2$ in Ex.~3.10(b), 4.12 and 7.7(b); linear independence of the forms in Ex.~7.3(b); connectedness of $X$ in Ex.~7.9; $r\ge1$ in Ex.~5.10 and 5.14; and the exercises 6.6(b),(c), which fail at the base point $P_0$. A few exercises are open-ended (Ex.~4.12(b)(3)), and some solutions use standard facts beyond the text (cohomology from Chapter III, Krull--Akizuki, Weierstrass-type theorems), which are named where used.

\section*{1. Sheaves}

Throughout, ``sheaf'' means sheaf of abelian groups on the indicated space unless said otherwise. For a presheaf $\mathcal F$ we write $\mathcal F^+$ for the associated sheaf and $\theta\colon\mathcal F\to\mathcal F^+$ for the canonical map (Prop.--Def.~1.2); recall $\mathcal F^+_P=\mathcal F_P$, and $\theta$ is an isomorphism if $\mathcal F$ is already a sheaf. A morphism of presheaves from $\mathcal F$ to a sheaf $\mathcal G$ factors uniquely through $\theta$. Two elementary facts about sheaves will be used repeatedly.

\medskip
\noindent\textbf{Fact A.} If $\mathcal F$ is a sheaf and $s\in\mathcal F(U)$ has $s_P=0$ for all $P\in U$, then $s=0$. (Each $s_P=0$ gives a neighborhood $W_P\subseteq U$ with $s|_{W_P}=0$; the $W_P$ cover $U$; apply (3).) Hence a sheaf all of whose stalks are $0$ is the zero sheaf.

\medskip
\noindent\textbf{Fact B.} Let $\mathcal F',\mathcal F''$ be subsheaves of a sheaf $\mathcal F$ such that $\mathcal F'_P=\mathcal F''_P$ as subgroups of $\mathcal F_P$ for every $P$. Then $\mathcal F'=\mathcal F''$. Indeed, let $s\in\mathcal F'(U)$. For $P\in U$ we have $s_P\in\mathcal F''_P$, so there is an open $W_P\ni P$, $W_P\subseteq U$, and $t\in\mathcal F''(W_P)$ with $t_P=s_P$; shrinking $W_P$ we may assume $t=s|_{W_P}$. Thus $s|_{W_P}\in\mathcal F''(W_P)$ for all $P$; these sections (elements of $\mathcal F''(W_P)\subseteq \mathcal F(W_P)$) agree on overlaps because they are restrictions of $s$, so they glue in $\mathcal F''$ to some $t\in\mathcal F''(U)$ with $t|_{W_P}=s|_{W_P}$ for all $P$; by (3) for $\mathcal F$, $t=s$. So $\mathcal F'(U)\subseteq\mathcal F''(U)$, and by symmetry equality holds.

\exercise{1.1}{Let $A$ be an abelian group, and define the \emph{constant presheaf} associated to $A$ on the topological space $X$ to be the presheaf $U\mapsto A$ for all $U\neq\emptyset$, with restriction maps the identity. Show that the constant sheaf $\mathcal A$ defined in the text is the sheaf associated to this presheaf.}
\pf
Let $\mathcal P$ be the constant presheaf ($\mathcal P(\emptyset)=0$, $\mathcal P(U)=A$ otherwise). Every open neighborhood of a point is nonempty, so all groups in the direct system defining $\mathcal P_P$ equal $A$ and all maps are the identity; hence $\mathcal P_P=A$ and the germ of $a\in\mathcal P(U)$ at $P\in U$ is $a$.

By the construction in Prop.--Def.~1.2, $\mathcal P^+(U)$ is the set of functions $s\colon U\to\bigcup_{P\in U}\mathcal P_P=A$ (with $s(P)\in\mathcal P_P=A$) such that every $P\in U$ has a neighborhood $V\subseteq U$ and $t\in\mathcal P(V)=A$ with $s(Q)=t_Q=t$ for all $Q\in V$. That is, $\mathcal P^+(U)$ is exactly the set of locally constant functions $U\to A$, i.e.\ continuous maps $U\to A$ for the discrete topology on $A$; this is $\mathcal A(U)$ by definition (Example 1.0.3), with the group structure pointwise and restriction maps the restriction of functions on both sides. (For $U=\emptyset$ both groups are $0$, being the set of the single empty function.) So $\mathcal P^+=\mathcal A$ as sheaves, and under this identification $\theta\colon \mathcal P\to\mathcal A$ sends $a\in A=\mathcal P(U)$ to the constant function $a$. Thus $\mathcal A$ is the sheaf associated to $\mathcal P$.
\hfill$\square$

\exercise{1.2}{\begin{enumerate}[label=(\alph*),leftmargin=*]
\item For any morphism of sheaves $\varphi\colon\mathcal F\to\mathcal G$, show that for each point $P$, $(\ker\varphi)_P=\ker(\varphi_P)$ and $(\operatorname{im}\varphi)_P=\operatorname{im}(\varphi_P)$.
\item Show that $\varphi$ is injective (respectively, surjective) if and only if the induced map $\varphi_P$ on the stalks is injective (respectively, surjective) for all $P$.
\item Show that a sequence $\dots\to\mathcal F^{i-1}\xrightarrow{\varphi^{i-1}}\mathcal F^i\xrightarrow{\varphi^i}\mathcal F^{i+1}\to\dots$ of sheaves and morphisms is exact if and only if for each $P\in X$ the corresponding sequence of stalks is exact as a sequence of abelian groups.
\end{enumerate}}
\pf
(a) \emph{Kernel.} $\ker\varphi$ is the subsheaf of $\mathcal F$ with $(\ker\varphi)(U)=\ker\varphi(U)$. The map $(\ker\varphi)_P\to\mathcal F_P$ is injective: if $s\in\ker\varphi(U)$ has zero germ in $\mathcal F_P$ then $s|_W=0$ for some neighborhood $W$, which is zero in $(\ker\varphi)(W)$. Its image lies in $\ker\varphi_P$ since $\varphi_P(s_P)=\varphi(s)_P=0$. Conversely let $g\in\ker\varphi_P$, $g=s_P$ with $s\in\mathcal F(U)$. Then $\varphi(s)_P=0$, so there is a neighborhood $W\subseteq U$ of $P$ with $0=\varphi(s)|_W=\varphi(s|_W)$; so $s|_W\in\ker\varphi(W)$ and $g=(s|_W)_P$ is in the image of $(\ker\varphi)_P$. Hence $(\ker\varphi)_P=\ker\varphi_P$.

\emph{Image.} Let $\mathcal I$ be the presheaf $U\mapsto\varphi(U)(\mathcal F(U))$, i.e.\ $\mathcal I(U)=\varphi(U)(\mathcal F(U))\subseteq\mathcal G(U)$. Then $\operatorname{im}\varphi=\mathcal I^+$, and the natural map $\operatorname{im}\varphi\to\mathcal G$ is induced from the inclusion $\mathcal I\to\mathcal G$. As in the first paragraph, the stalk map $\mathcal I_P\to\mathcal G_P$ is injective (if $\varphi(s)$, $s\in \mathcal F(U)$, has zero germ then $\varphi(s)|_W=0$ for some $W$, and $0$ is zero in $\mathcal I(W)$) and its image is $\{\varphi(s)_P\}=\{\varphi_P(s_P)\}=\operatorname{im}\varphi_P$. Since $\mathcal I^+_P=\mathcal I_P$ compatibly with the maps to $\mathcal G_P$, $(\operatorname{im}\varphi)_P=\operatorname{im}\varphi_P$ as subgroups of $\mathcal G_P$.

(b) $\varphi$ is injective iff $\ker\varphi=0$ iff all stalks $(\ker\varphi)_P=\ker\varphi_P$ vanish (Fact A for ``if''), iff all $\varphi_P$ are injective. $\varphi$ is surjective iff $\operatorname{im}\varphi=\mathcal G$, i.e.\ iff the inclusion $\operatorname{im}\varphi\to\mathcal G$ of sheaves is an isomorphism; by Prop.~1.1 this holds iff it is an isomorphism on every stalk, i.e.\ iff $\operatorname{im}\varphi_P=\mathcal G_P$ for all $P$ by (a), i.e.\ iff every $\varphi_P$ is surjective.

(c) The sequence is exact at $\mathcal F^i$ iff $\operatorname{im}\varphi^{i-1}=\ker\varphi^i$ as subsheaves of $\mathcal F^i$. If so, taking stalks and using (a) gives $\operatorname{im}\varphi^{i-1}_P=\ker\varphi^i_P$. Conversely, if these are equal for all $P$, the two subsheaves have the same stalks (by (a)), so are equal by Fact B. \hfill$\square$

\exercise{1.3}{\begin{enumerate}[label=(\alph*),leftmargin=*]
\item Let $\varphi\colon\mathcal F\to\mathcal G$ be a morphism of sheaves on $X$. Show that $\varphi$ is surjective if and only if the following condition holds: for every open set $U\subseteq X$, and for every $s\in\mathcal G(U)$, there is a covering $\{U_i\}$ of $U$, and there are elements $t_i\in\mathcal F(U_i)$, such that $\varphi(t_i)=s|_{U_i}$ for all $i$.
\item Give an example of a surjective morphism of sheaves $\varphi\colon\mathcal F\to\mathcal G$, and an open set $U$ such that $\varphi(U)\colon\mathcal F(U)\to\mathcal G(U)$ is not surjective.
\end{enumerate}}
\pf
(a) ($\Rightarrow$) Let $\varphi$ be surjective, $s\in\mathcal G(U)$, $P\in U$. By Ex.~1.2(b), $\varphi_P$ is surjective, so $s_P=\varphi_P(t_P)$ for some germ $t_P$ represented by $t\in\mathcal F(V)$, $V\ni P$ open, $V\subseteq U$. Then $\varphi(t)_P=s_P$, so $\varphi(t)|_{V'}=s|_{V'}$ for some smaller open $V'\ni P$, and $t':=t|_{V'}$ satisfies $\varphi(t')=s|_{V'}$. Doing this for every $P\in U$ gives the covering.

($\Leftarrow$) By Ex.~1.2(b) it suffices to show each $\varphi_P$ is surjective. Let $g\in\mathcal G_P$ be represented by $s\in\mathcal G(U)$. Take a covering $\{U_i\}$ and $t_i$ as in the condition; choose $i$ with $P\in U_i$. Then $\varphi_P((t_i)_P)=(\varphi(t_i))_P=(s|_{U_i})_P=g$.

(b) Let $X=S^1=\{z\in\C:|z|=1\}$ with its usual topology. Let $\mathcal F$ be the sheaf of continuous real-valued functions on $X$ (abelian group under addition), and $\mathcal G$ the sheaf of continuous functions to $S^1$, which is an abelian group under pointwise multiplication (both are sheaves since continuity is local). Define $\varphi(U)(f)=\ex^{2\pi i f}$; this is a homomorphism compatible with restrictions, so a morphism of sheaves.

\emph{$\varphi$ is surjective:} let $g\in\mathcal G(U)$ and $P\in U$. Put $w=g(P)$ and $V=g^{-1}(S^1\setminus\{-w\})$, an open neighborhood of $P$. On the simply cut circle $S^1\setminus\{-w\}$ there is a continuous function $\alpha$ with $\ex^{2\pi i\alpha(z)}=z$ (a continuous branch of $\frac1{2\pi i}\log$; e.g.\ fix $\beta$ with $\ex^{2\pi i\beta}=w$ and let $\alpha(z)$ be the unique number in $(\beta-\tfrac12,\beta+\tfrac12)$ with $\ex^{2\pi i\alpha(z)}=z$; the map $\alpha$ is the inverse of the homeomorphism $(\beta-\tfrac12,\beta+\tfrac12)\to S^1\setminus\{-w\}$ given by $\ex^{2\pi i(\cdot)}$). Then $f=\alpha\circ g|_V\in\mathcal F(V)$ satisfies $\varphi(f)=g|_V$. By (a), $\varphi$ is surjective.

\emph{$\varphi(X)$ is not surjective:} take $g=\mathrm{id}\in\mathcal G(X)$. Suppose $f\colon S^1\to\R$ is continuous with $\ex^{2\pi if(z)}=z$ for all $z$. Then $f$ is injective (it determines $z$). Its image is compact and connected, a closed interval $[a,b]$, with $a<b$ since $f$ is injective on an infinite set. Pick $z_0$ with $f(z_0)=c\in(a,b)$ (possible: only two points map to the endpoints). Then $S^1\setminus\{z_0\}$ is connected and $f$ maps it bijectively onto $[a,b]\setminus\{c\}$, which is disconnected, contradicting that continuous images of connected sets are connected. So $\mathrm{id}$ is not in the image of $\varphi(X)$. \hfill$\square$

\exercise{1.4}{\begin{enumerate}[label=(\alph*),leftmargin=*]
\item Let $\varphi\colon\mathcal F\to\mathcal G$ be a morphism of presheaves such that $\varphi(U)\colon\mathcal F(U)\to\mathcal G(U)$ is injective for each $U$. Show that the induced map $\varphi^+\colon\mathcal F^+\to\mathcal G^+$ of associated sheaves is injective.
\item Use part (a) to show that if $\varphi\colon\mathcal F\to\mathcal G$ is a morphism of sheaves, then $\operatorname{im}\varphi$ can be naturally identified with a subsheaf of $\mathcal G$, as mentioned in the text.
\end{enumerate}}
\pf
(a) The morphism $\varphi^+$ is the unique morphism with $\varphi^+\theta_{\mathcal F}=\theta_{\mathcal G}\varphi$. Concretely, $\varphi^+(U)$ sends $s\in\mathcal F^+(U)$ (a function $P\mapsto s(P)\in\mathcal F_P$) to $P\mapsto\varphi_P(s(P))$: this function lies in $\mathcal G^+(U)$ (if $s=t_Q$ on a neighborhood $V$ with $t\in\mathcal F(V)$, then $\varphi_Q(s(Q))=\varphi(t)_Q$ there), the assignment is compatible with restriction, and it satisfies $\varphi^+\theta_{\mathcal F}=\theta_{\mathcal G}\varphi$; by uniqueness it is $\varphi^+$.

First, $\varphi_P$ is injective for each $P$: if $t\in\mathcal F(V)$ and $\varphi_P(t_P)=\varphi(t)_P=0$ then $\varphi(t|_W)=\varphi(t)|_W=0$ for some neighborhood $W$ of $P$, hence $t|_W=0$ as $\varphi(W)$ is injective, so $t_P=0$. Now if $\varphi^+(s)=0$ then $\varphi_P(s(P))=0$, so $s(P)=0$ for every $P$, i.e.\ $s=0$. So each $\varphi^+(U)$ is injective.

(b) Let $\mathcal I$ be the presheaf image $U\mapsto\varphi(U)(\mathcal F(U))\subseteq\mathcal G(U)$, so $\operatorname{im}\varphi=\mathcal I^+$. The inclusion $\iota\colon\mathcal I\to\mathcal G$ has $\iota(U)$ injective, so by (a) $\iota^+\colon\mathcal I^+\to\mathcal G^+$ is injective. Composing with the isomorphism $\theta_{\mathcal G}^{-1}\colon\mathcal G^+\to\mathcal G$ ($\mathcal G$ is a sheaf) gives an injection $\operatorname{im}\varphi\to\mathcal G$, which is the natural map of the text (it is the map obtained from the universal property applied to $\iota$). Hence $\operatorname{im}\varphi$ is identified with a subsheaf of $\mathcal G$, whose sections over $U$ are those $s\in\mathcal G(U)$ which locally (on a cover of $U$) lie in the image of $\varphi$; its stalks are $\operatorname{im}\varphi_P$ by Ex.~1.2(a). \hfill$\square$

\exercise{1.5}{Show that a morphism of sheaves is an isomorphism if and only if it is both injective and surjective.}
\pf
If $\varphi\colon\mathcal F\to\mathcal G$ is an isomorphism, each $\varphi_P$ is an isomorphism (inverse $(\varphi^{-1})_P$), hence injective and surjective, so $\varphi$ is injective and surjective by Ex.~1.2(b). Conversely if $\varphi$ is injective and surjective, then all $\varphi_P$ are bijective by Ex.~1.2(b), hence isomorphisms, and $\varphi$ is an isomorphism by Prop.~1.1. \hfill$\square$

\exercise{1.6}{\begin{enumerate}[label=(\alph*),leftmargin=*]
\item Let $\mathcal F'$ be a subsheaf of a sheaf $\mathcal F$. Show that the natural map of $\mathcal F$ to the quotient sheaf $\mathcal F/\mathcal F'$ is surjective, and has kernel $\mathcal F'$. Thus there is an exact sequence $0\to\mathcal F'\to\mathcal F\to\mathcal F/\mathcal F'\to0$.
\item Conversely, if $0\to\mathcal F'\to\mathcal F\to\mathcal F''\to0$ is an exact sequence, show that $\mathcal F'$ is isomorphic to a subsheaf of $\mathcal F$, and that $\mathcal F''$ is isomorphic to the quotient of $\mathcal F$ by this subsheaf.
\end{enumerate}}
\pf
(a) The natural map $\pi\colon\mathcal F\to\mathcal F/\mathcal F'$ is the composite of $\mathcal F\to Q$, where $Q$ is the presheaf quotient $U\mapsto\mathcal F(U)/\mathcal F'(U)$, with $\theta\colon Q\to Q^+=\mathcal F/\mathcal F'$. On stalks, $Q_P=\mathcal F_P/\mathcal F'_P=(\mathcal F/\mathcal F')_P$ (text), and $\pi_P$ is the quotient map $\mathcal F_P\to\mathcal F_P/\mathcal F'_P$. This is surjective, so $\pi$ is surjective (Ex.~1.2(b)), and $(\ker\pi)_P=\ker\pi_P=\mathcal F'_P$ (Ex.~1.2(a)). Thus $\ker\pi$ and $\mathcal F'$ are subsheaves of $\mathcal F$ with the same stalks, hence equal by Fact B. The inclusion $\mathcal F'\to\mathcal F$ is injective with image $\ker\pi$, so the sequence is exact.

(b) Let $\alpha\colon\mathcal F'\to\mathcal F$, $\beta\colon\mathcal F\to\mathcal F''$. Since $\alpha$ is injective, each $\alpha(U)$ is injective, so $\mathcal F'_1(U):=\alpha(U)(\mathcal F'(U))$ defines a subpresheaf of $\mathcal F$ and $\alpha$ induces an isomorphism of presheaves $\mathcal F'\to\mathcal F'_1$. Hence $\mathcal F'_1$ is a sheaf, a subsheaf of $\mathcal F$, isomorphic to $\mathcal F'$; it is the image of $\alpha$ (the presheaf image is already a sheaf). By exactness $\ker\beta=\operatorname{im}\alpha=\mathcal F'_1$. Since $\beta\alpha=0$, $\beta(U)$ vanishes on $\mathcal F'_1(U)$ and induces $\bar\beta(U)\colon\mathcal F(U)/\mathcal F'_1(U)\to\mathcal F''(U)$, a morphism from the presheaf quotient to the sheaf $\mathcal F''$; it factors through a unique $\gamma\colon\mathcal F/\mathcal F'_1\to\mathcal F''$. On stalks, $\gamma_P\colon\mathcal F_P/(\mathcal F'_1)_P\to\mathcal F''_P$ is induced by $\beta_P$; it is surjective since $\beta_P$ is (Ex.~1.2(b)) and injective since $\ker\beta_P=(\ker\beta)_P=(\mathcal F_1')_P$. So $\gamma$ is an isomorphism by Prop.~1.1. \hfill$\square$

\exercise{1.7}{Let $\varphi\colon\mathcal F\to\mathcal G$ be a morphism of sheaves.
\begin{enumerate}[label=(\alph*),leftmargin=*]
\item Show that $\operatorname{im}\varphi\cong\mathcal F/\ker\varphi$.
\item Show that $\operatorname{coker}\varphi\cong\mathcal G/\operatorname{im}\varphi$.
\end{enumerate}}
\pf
(a) The presheaf quotient $Q(U)=\mathcal F(U)/\ker\varphi(U)$ maps to the presheaf image $\mathcal I(U)=\varphi(U)(\mathcal F(U))$ by $\bar\varphi(U)$, and each $\bar\varphi(U)$ is bijective (first isomorphism theorem), compatible with restrictions. So $Q\cong\mathcal I$ as presheaves, whence $\mathcal F/\ker\varphi=Q^+\cong\mathcal I^+=\operatorname{im}\varphi$.

(b) By definition $\operatorname{coker}\varphi=C^+$ with $C(U)=\mathcal G(U)/\varphi(U)(\mathcal F(U))=\mathcal G(U)/\mathcal I(U)$. Regard $\operatorname{im}\varphi$ as a subsheaf of $\mathcal G$ (Ex.~1.4(b)); then $\mathcal I\subseteq\operatorname{im}\varphi\subseteq\mathcal G$ as subpresheaves, so the identity of $\mathcal G(U)$ induces a morphism of presheaves $C\to D$, $D(U)=\mathcal G(U)/(\operatorname{im}\varphi)(U)$, and hence a morphism of sheaves $\gamma\colon\operatorname{coker}\varphi=C^+\to D^+=\mathcal G/\operatorname{im}\varphi$. On stalks (direct limits are exact, which is also the justification of $(\mathcal F/\mathcal F')_P=\mathcal F_P/\mathcal F'_P$ in the text), $\gamma_P$ is the map $\mathcal G_P/\mathcal I_P\to\mathcal G_P/(\operatorname{im}\varphi)_P$ induced by the identity of $\mathcal G_P$. By the proof of Ex.~1.2(a), $\mathcal I_P=(\operatorname{im}\varphi)_P=\operatorname{im}\varphi_P$ inside $\mathcal G_P$, so $\gamma_P$ is the identity of $\mathcal G_P/\operatorname{im}\varphi_P$. By Prop.~1.1, $\gamma$ is an isomorphism. \hfill$\square$

\exercise{1.8}{For any open subset $U\subseteq X$, show that the functor $\Gamma(U,\cdot)$ from sheaves on $X$ to abelian groups is a left exact functor, i.e., if $0\to\mathcal F'\to\mathcal F\to\mathcal F''$ is an exact sequence of sheaves, then $0\to\Gamma(U,\mathcal F')\to\Gamma(U,\mathcal F)\to\Gamma(U,\mathcal F'')$ is an exact sequence of groups. The functor $\Gamma(U,\cdot)$ need not be exact; see (Ex.~1.21) below.}
\pf
Let the maps be $\alpha\colon\mathcal F'\to\mathcal F$, $\beta\colon\mathcal F\to\mathcal F''$. Exactness at $\mathcal F'$ means $\ker\alpha=0$, i.e.\ $\ker\alpha(U)=0$: $\alpha(U)$ is injective. Exactness at $\mathcal F$ means $\operatorname{im}\alpha=\ker\beta$; since $\alpha$ is injective, the presheaf image $V\mapsto\alpha(V)(\mathcal F'(V))$ is already a sheaf (Ex.~1.6(b)), hence equals $\operatorname{im}\alpha$. Thus $\ker\beta(U)=(\operatorname{im}\alpha)(U)=\alpha(U)(\mathcal F'(U))$, which says exactly that $\ker\beta(U)=\operatorname{im}\alpha(U)$. \hfill$\square$

\exercise{1.9}{\emph{Direct Sum.} Let $\mathcal F$ and $\mathcal G$ be sheaves on $X$. Show that the presheaf $U\mapsto\mathcal F(U)\oplus\mathcal G(U)$ is a sheaf. It is called the direct sum of $\mathcal F$ and $\mathcal G$, and is denoted $\mathcal F\oplus\mathcal G$. Show that it plays the role of direct sum and of direct product in the category of sheaves of abelian groups on $X$.}
\pf
Let $\mathcal H(U)=\mathcal F(U)\oplus\mathcal G(U)$ with componentwise restriction; $\mathcal H(\emptyset)=0\oplus0=0$.

(3) If $(a,b)\in\mathcal H(U)$ restricts to $0$ on all $V_i$ of a cover, then $a|_{V_i}=0=b|_{V_i}$ for all $i$, so $a=0=b$.

(4) If $(a_i,b_i)\in\mathcal H(V_i)$ agree on overlaps, then $a_i|_{V_i\cap V_j}=a_j|_{V_i\cap V_j}$ and likewise for $b$, so they glue to $a\in\mathcal F(U)$ and $b\in\mathcal G(U)$, and $(a,b)|_{V_i}=(a_i,b_i)$.

\emph{Universal properties.} There are morphisms of sheaves $p_1\colon\mathcal H\to\mathcal F$, $p_2\colon\mathcal H\to\mathcal G$ (projections) and $\iota_1\colon\mathcal F\to\mathcal H$, $a\mapsto(a,0)$, $\iota_2\colon\mathcal G\to\mathcal H$, $b\mapsto(0,b)$, with $p_k\iota_k=\mathrm{id}$, $p_1\iota_2=0=p_2\iota_1$, and $\iota_1p_1+\iota_2p_2=\mathrm{id}$ (all clearly compatible with restriction).

\emph{Product:} given a sheaf $\mathcal E$ and morphisms $f\colon\mathcal E\to\mathcal F$, $g\colon\mathcal E\to\mathcal G$, the map $(f,g)(U)\colon e\mapsto(f(e),g(e))$ is a morphism with $p_1(f,g)=f$, $p_2(f,g)=g$; if $h$ is any morphism with $p_1h=f$, $p_2h=g$, then $h(e)=(f(e),g(e))$ on every section, so $h=(f,g)$.

\emph{Coproduct:} given $f\colon\mathcal F\to\mathcal E$, $g\colon\mathcal G\to\mathcal E$, the map $(a,b)\mapsto f(a)+g(b)$ is a morphism $h\colon\mathcal H\to\mathcal E$ with $h\iota_1=f$, $h\iota_2=g$. If $h'$ also satisfies this then $h'(a,b)=h'(\iota_1a+\iota_2b)=f(a)+g(b)$, so $h'=h$. \hfill$\square$

\exercise{1.10}{\emph{Direct Limit.} Let $\{\mathcal F_i\}$ be a direct system of sheaves and morphisms on $X$. We define the direct limit of the system $\{\mathcal F_i\}$, denoted $\varinjlim\mathcal F_i$, to be the sheaf associated to the presheaf $U\mapsto\varinjlim\mathcal F_i(U)$. Show that this is a direct limit in the category of sheaves on $X$, i.e., that it has the following universal property: given a sheaf $\mathcal G$, and a collection of morphisms $\mathcal F_i\to\mathcal G$, compatible with the maps of the direct system, then there exists a unique map $\varinjlim\mathcal F_i\to\mathcal G$ such that for each $i$, the original map $\mathcal F_i\to\mathcal G$ is obtained by composing the maps $\mathcal F_i\to\varinjlim\mathcal F_i\to\mathcal G$.}
\pf
Here the direct system is indexed by a directed set $I$, with maps $f_{ij}\colon\mathcal F_i\to\mathcal F_j$ for $i\le j$. Let $\mathcal P(U)=\varinjlim_i\mathcal F_i(U)$ with the evident restriction maps (induced by those of the $\mathcal F_i$, which commute with the $f_{ij}$); $\mathcal P$ is a presheaf, and $\mathcal F:=\mathcal P^+=\varinjlim\mathcal F_i$. Let $\lambda_i\colon\mathcal F_i(U)\to\mathcal P(U)$ be the canonical maps; they are morphisms of presheaves, and $\mu_i:=\theta\lambda_i\colon\mathcal F_i\to\mathcal F$ are morphisms of sheaves with $\mu_j f_{ij}=\mu_i$.

Let $\mathcal G$ be a sheaf with morphisms $g_i\colon\mathcal F_i\to\mathcal G$, $g_jf_{ij}=g_i$. For each $U$ the universal property of the direct limit of groups gives a unique homomorphism $g(U)\colon\mathcal P(U)\to\mathcal G(U)$ with $g(U)\lambda_i(U)=g_i(U)$. For $V\subseteq U$, the maps $\rho_{UV}g(U)$ and $g(V)\rho_{UV}$ agree after composing with each $\lambda_i(U)$ (both equal $\rho_{UV}g_i(U)=g_i(V)\rho_{UV}$), so they are equal: $g\colon\mathcal P\to\mathcal G$ is a morphism of presheaves. Since $\mathcal G$ is a sheaf, $g=\bar g\theta$ for a unique morphism $\bar g\colon\mathcal F\to\mathcal G$ (Prop.--Def.~1.2). Then $\bar g\mu_i=\bar g\theta\lambda_i=g\lambda_i=g_i$.

Uniqueness: if $h\colon\mathcal F\to\mathcal G$ satisfies $h\mu_i=g_i$ for all $i$, then $h\theta\colon\mathcal P\to\mathcal G$ satisfies $h\theta\lambda_i=g_i=g\lambda_i$ for all $i$, so $h\theta=g$ on each $\mathcal P(U)=\varinjlim\mathcal F_i(U)$ (generated by the images of the $\lambda_i(U)$), and by the uniqueness in Prop.--Def.~1.2, $h=\bar g$. \hfill$\square$

\exercise{1.11}{Let $\{\mathcal F_i\}$ be a direct system of sheaves on a noetherian topological space $X$. In this case show that the presheaf $U\mapsto\varinjlim\mathcal F_i(U)$ is already a sheaf. In particular, $\Gamma(X,\varinjlim\mathcal F_i)=\varinjlim\Gamma(X,\mathcal F_i)$.}
\pf
Let $\mathcal P(U)=\varinjlim\mathcal F_i(U)$ ($I$ directed). Recall that on a noetherian space every open subset is quasi-compact (every open cover of an open set has a finite subcover: if not, choose successively members of the cover not contained in the union of the previous ones to get a strictly increasing infinite chain of open sets, contradicting the ascending chain condition on opens, equivalent to the descending chain condition on closed sets). Every element of $\mathcal P(U)$ is the class of some $s_i\in\mathcal F_i(U)$; and the class of $s_i$ is $0$ iff $f_{ik}(s_i)=0$ for some $k\ge i$.

(3) Let $U=\bigcup U_j$ and $s\in\mathcal P(U)$ with $s|_{U_j}=0$ for all $j$. Write $s$ as the class of $s_i\in\mathcal F_i(U)$. For each $j$ there is $k_j\ge i$ with $f_{ik_j}(s_i)|_{U_j}=0$ in $\mathcal F_{k_j}(U_j)$. By quasi-compactness $U=U_{j_1}\cup\dots\cup U_{j_n}$; take $k\ge k_{j_1},\dots,k_{j_n}$. Then $f_{ik}(s_i)|_{U_{j_m}}=f_{k_{j_m}k}\bigl(f_{ik_{j_m}}(s_i)|_{U_{j_m}}\bigr)=0$ for $m=1,\dots,n$, so $f_{ik}(s_i)=0$ because $\mathcal F_k$ is a sheaf. Hence $s=0$.

(4) Let $U=\bigcup U_j$ and $s^j\in\mathcal P(U_j)$ with $s^j|_{U_j\cap U_l}=s^l|_{U_j\cap U_l}$. Take a finite subcover $U_1,\dots,U_n$ (renumbering). Represent $s^m$ by $t_m\in\mathcal F_{i_m}(U_m)$ and choose $i\ge i_1,\dots,i_n$; replacing $t_m$ by $f_{i_mi}(t_m)$ we may assume all $t_m\in\mathcal F_i(U_m)$. For each pair $(m,l)$ the equality of classes of $t_m|_{U_m\cap U_l}$ and $t_l|_{U_m\cap U_l}$ means there is $k_{ml}\ge i$ with $f_{ik_{ml}}(t_m|)=f_{ik_{ml}}(t_l|)$. Let $k\ge i$ dominate the finitely many $k_{ml}$ and put $t'_m=f_{ik}(t_m)\in\mathcal F_k(U_m)$; then $t'_m|_{U_m\cap U_l}=t'_l|_{U_m\cap U_l}$ in $\mathcal F_k$. Since $\mathcal F_k$ is a sheaf, there is $t\in\mathcal F_k(U)$ with $t|_{U_m}=t'_m$. Let $s\in\mathcal P(U)$ be its class; then $s|_{U_m}=s^m$ for $m=1,\dots,n$. For an arbitrary $j$, the sections $s|_{U_j}$ and $s^j$ of $\mathcal P(U_j)$ have the same restriction to each $U_j\cap U_m$: $s|_{U_j\cap U_m}=s^m|_{U_j\cap U_m}=s^j|_{U_j\cap U_m}$. As the $U_j\cap U_m$ ($m=1..n$) cover $U_j$, property (3) just proved gives $s|_{U_j}=s^j$.

So $\mathcal P$ is a sheaf; thus $\varinjlim\mathcal F_i=\mathcal P^+\cong\mathcal P$ via $\theta$, and $\Gamma(X,\varinjlim\mathcal F_i)=\mathcal P(X)=\varinjlim\mathcal F_i(X)=\varinjlim\Gamma(X,\mathcal F_i)$. \hfill$\square$

\exercise{1.12}{\emph{Inverse Limit.} Let $\{\mathcal F_i\}$ be an inverse system of sheaves on $X$. Show that the presheaf $U\mapsto\varprojlim\mathcal F_i(U)$ is a sheaf. It is called the inverse limit of the system $\{\mathcal F_i\}$, and is denoted $\varprojlim\mathcal F_i$. Show that it has the universal property of an inverse limit in the category of sheaves.}
\pf
Let the system be indexed by a partially ordered set $I$ with maps $f_{ji}\colon\mathcal F_j\to\mathcal F_i$ for $i\le j$ (with $f_{ii}=\mathrm{id}$, $f_{ji}f_{kj}=f_{ki}$). Put $\mathcal P(U)=\varprojlim\mathcal F_i(U)=\{(s_i)\in\prod_i\mathcal F_i(U): f_{ji}(s_j)=s_i\ \forall i\le j\}$, with componentwise restriction; this is a presheaf (and $\mathcal P(\emptyset)=0$).

(3) If $s=(s_i)\in\mathcal P(U)$ and $s|_{V_a}=0$ for a cover $\{V_a\}$, then $s_i|_{V_a}=0$ for all $a$, $i$, so $s_i=0$ for all $i$ since $\mathcal F_i$ is a sheaf.

(4) Let $s^a=(s^a_i)_i\in\mathcal P(V_a)$ agree on overlaps. For each $i$, the $s^a_i\in\mathcal F_i(V_a)$ agree on overlaps, so glue to $s_i\in\mathcal F_i(U)$ with $s_i|_{V_a}=s^a_i$. For $i\le j$, $f_{ji}(s_j)$ and $s_i$ in $\mathcal F_i(U)$ both restrict on $V_a$ to $f_{ji}(s^a_j)=s^a_i$, so are equal by (3) for $\mathcal F_i$. So $s=(s_i)\in\mathcal P(U)$ and $s|_{V_a}=s^a$.

\emph{Universal property.} The projections $p_i\colon\mathcal P\to\mathcal F_i$ are morphisms with $f_{ji}p_j=p_i$. Let $\mathcal G$ be a sheaf with morphisms $g_i\colon\mathcal G\to\mathcal F_i$ such that $f_{ji}g_j=g_i$. Then $g(U)\colon\mathcal G(U)\to\mathcal P(U)$, $t\mapsto(g_i(U)(t))_i$ lands in $\mathcal P(U)$, is compatible with restrictions, and satisfies $p_ig=g_i$. If $h\colon\mathcal G\to\mathcal P$ is any morphism with $p_ih=g_i$ for all $i$, then $h(t)=(p_ih(t))_i=(g_i(t))_i=g(t)$. So $g$ is unique. \hfill$\square$

\exercise{1.13}{\emph{Espace Etale of a Presheaf.} (This exercise is included only to establish the connection between our definition of a sheaf and another definition often found in the literature. See for example Godement [1, Ch.~II, Section 1.2].) Given a presheaf $\mathcal F$ on $X$, we define a topological space $\operatorname{Spe}(\mathcal F)$, called the espace etale of $\mathcal F$, as follows. As a set, $\operatorname{Spe}(\mathcal F)=\bigcup_{P\in X}\mathcal F_P$. We define a projection map $\pi\colon\operatorname{Spe}(\mathcal F)\to X$ by sending $s\in\mathcal F_P$ to $P$. For each open set $U\subseteq X$ and each section $s\in\mathcal F(U)$, we obtain a map $\bar s\colon U\to\operatorname{Spe}(\mathcal F)$ by sending $P\mapsto s_P$, its germ at $P$. This map has the property that $\pi\circ\bar s=\mathrm{id}_U$, in other words, it is a ``section'' of $\pi$ over $U$. We now make $\operatorname{Spe}(\mathcal F)$ into a topological space by giving it the strongest topology such that all the maps $\bar s\colon U\to\operatorname{Spe}(\mathcal F)$ for all $U$, and all $s\in\mathcal F(U)$, are continuous. Now show that the sheaf $\mathcal F^+$ associated to $\mathcal F$ can be described as follows: for any open set $U\subseteq X$, $\mathcal F^+(U)$ is the set of continuous sections of $\operatorname{Spe}(\mathcal F)$ over $U$. In particular, the original presheaf $\mathcal F$ was a sheaf if and only if for each $U$, $\mathcal F(U)$ is equal to the set of all continuous sections of $\operatorname{Spe}(\mathcal F)$ over $U$.}
\pf
Write $E=\operatorname{Spe}(\mathcal F)$. The strongest topology making all $\bar s$ continuous is: $W\subseteq E$ is open iff $\bar s^{-1}(W)$ is open in $U$ for all open $U$ and all $s\in\mathcal F(U)$. (This is a topology, and any topology making all $\bar s$ continuous is contained in it.) A \emph{section} of $\pi$ over $U$ is a map $\sigma\colon U\to E$ with $\pi\sigma=\mathrm{id}_U$, i.e.\ $\sigma(P)\in\mathcal F_P$; it is continuous when $\sigma^{-1}(W)$ is open for all open $W\subseteq E$.

By Prop.--Def.~1.2, $\mathcal F^+(U)$ is the set of sections $\sigma$ of $\pi$ over $U$ such that each $P\in U$ has a neighborhood $V\subseteq U$ and $t\in\mathcal F(V)$ with $\sigma(Q)=t_Q$ for all $Q\in V$ (i.e. $\sigma|_V=\bar t$). We show this equals the set of continuous sections.

($\subseteq$) If $\sigma\in\mathcal F^+(U)$, then $U$ is covered by open sets $V$ on which $\sigma|_V=\bar t$ with $t\in\mathcal F(V)$; $\bar t$ is continuous by definition of the topology of $E$. Continuity is local, so $\sigma$ is continuous.

($\supseteq$) First a claim: for $t\in\mathcal F(V)$, the set $\bar t(V)$ is open in $E$. Let $s\in\mathcal F(U')$. Then $\bar s^{-1}(\bar t(V))=\{Q\in U'\cap V: s_Q=t_Q\}$ (if $s_Q=t_{Q'}$ then $Q=\pi(s_Q)=Q'$). This set is open: if $s_Q=t_Q$ there is an open $W\ni Q$, $W\subseteq U'\cap V$, with $s|_W=t|_W$, so $s_{Q'}=t_{Q'}$ for all $Q'\in W$. This proves the claim.

Now let $\sigma$ be a continuous section over $U$ and $P\in U$. Choose $t\in\mathcal F(V)$ with $t_P=\sigma(P)$. The set $N=\sigma^{-1}(\bar t(V))$ is an open neighborhood of $P$ in $U$ (by the claim and continuity), and we may assume $N\subseteq V$ because $\pi\sigma(Q)=Q$ and $\sigma(Q)\in\bar t(V)$ imply $Q\in V$. For $Q\in N$, $\sigma(Q)=t_{Q'}$ for some $Q'\in V$ and applying $\pi$ gives $Q'=Q$. So $\sigma(Q)=t_Q=(t|_N)_Q$ for all $Q\in N$. Thus $\sigma\in\mathcal F^+(U)$.

Restriction maps of $\mathcal F^+$ are restrictions of sections, so $\mathcal F^+$ is the sheaf of continuous sections of $E\to X$.

Finally, $\theta\colon\mathcal F\to\mathcal F^+$ is $s\mapsto\bar s$. If $\mathcal F$ is a sheaf, $\theta$ is an isomorphism (text), so $\mathcal F(U)\to\{\text{continuous sections over }U\}$, $s\mapsto\bar s$, is bijective for each $U$. Conversely, if all these maps are bijective then $\theta$ is an isomorphism of presheaves onto the sheaf $\mathcal F^+$, so $\mathcal F$ is a sheaf. \hfill$\square$

\exercise{1.14}{\emph{Support.} Let $\mathcal F$ be a sheaf on $X$, and let $s\in\mathcal F(U)$ be a section over an open set $U$. The \emph{support} of $s$, denoted $\operatorname{Supp}s$, is defined to be $\{P\in U\mid s_P\ne0\}$, where $s_P$ denotes the germ of $s$ in the stalk $\mathcal F_P$. Show that $\operatorname{Supp}s$ is a closed subset of $U$. We define the support of $\mathcal F$, $\operatorname{Supp}\mathcal F$, to be $\{P\in X\mid\mathcal F_P\ne0\}$. It need not be a closed subset.}
\pf
The complement of $\operatorname{Supp}s$ in $U$ is $\{P\in U:s_P=0\}$. If $s_P=0$, there is an open $W$, $P\in W\subseteq U$, with $s|_W=0$; then $s_Q=0$ for all $Q\in W$. So the complement is open and $\operatorname{Supp}s$ is closed in $U$.

\emph{Example where $\operatorname{Supp}\mathcal F$ is not closed.} Let $X=\R$ (usual topology) and let $\mathcal L$ be the constant sheaf $\Z$ (locally constant functions to $\Z$). Define $\mathcal F(V)=\{s\in\mathcal L(V): s(0)=0\text{ if }0\in V\}$. This is a subpresheaf of $\mathcal L$ defined by a pointwise condition, hence a sheaf: (3) is inherited from $\mathcal L$ and a glued section $s$ of $\mathcal L$ from sections satisfying the pointwise condition satisfies it. For $P\ne0$, small neighborhoods avoid $0$ and $\mathcal F_P=\mathcal L_P=\Z\ne0$. For $P=0$ a germ is represented by a locally constant $s$ with $s(0)=0$, hence $s=0$ on a neighborhood of $0$; so $\mathcal F_0=0$. Therefore $\operatorname{Supp}\mathcal F=\R\setminus\{0\}$, which is not closed. \hfill$\square$

\exercise{1.15}{\emph{Sheaf $\mathcal Hom$.} Let $\mathcal F,\mathcal G$ be sheaves of abelian groups on $X$. For any open set $U\subseteq X$, show that the set $\operatorname{Hom}(\mathcal F|_U,\mathcal G|_U)$ of morphisms of the restricted sheaves is an abelian group. Show that the presheaf $U\mapsto\operatorname{Hom}(\mathcal F|_U,\mathcal G|_U)$ is a sheaf. This sheaf is called the \emph{sheaf of local morphisms} of $\mathcal F$ into $\mathcal G$, ``sheaf hom'' for short, and is denoted $\mathcal Hom(\mathcal F,\mathcal G)$.}
\pf
A morphism $\varphi\colon\mathcal F|_U\to\mathcal G|_U$ is a family of homomorphisms $\varphi(W)\colon\mathcal F(W)\to\mathcal G(W)$, $W\subseteq U$ open, commuting with restrictions. Define $(\varphi+\psi)(W)=\varphi(W)+\psi(W)$; this is again a family of homomorphisms (as $\mathcal G(W)$ is abelian) commuting with restriction, with zero element the zero morphism and inverse $-\varphi$; the group axioms are inherited from $\mathcal G(W)$. So $\operatorname{Hom}(\mathcal F|_U,\mathcal G|_U)$ is an abelian group.

For open $V\subseteq U$, the restriction map sends $\varphi$ to $\varphi|_V=(\varphi(W))_{W\subseteq V}$; this is a homomorphism of groups, and $(\varphi|_V)|_{V'}=\varphi|_{V'}$, so $\mathcal Hom$ is a presheaf (value $0$ on $\emptyset$, where there is a unique morphism).

(3) Let $U=\bigcup U_i$, $\varphi\in\operatorname{Hom}(\mathcal F|_U,\mathcal G|_U)$ with $\varphi|_{U_i}=0$ for all $i$. For $W\subseteq U$ open and $s\in\mathcal F(W)$, we get $\varphi(W)(s)|_{W\cap U_i}=\varphi(W\cap U_i)(s|_{W\cap U_i})=0$ for all $i$, so $\varphi(W)(s)=0$ as $\mathcal G$ is a sheaf. Thus $\varphi=0$.

(4) Let $\varphi_i\in\operatorname{Hom}(\mathcal F|_{U_i},\mathcal G|_{U_i})$ with $\varphi_i|_{U_i\cap U_j}=\varphi_j|_{U_i\cap U_j}$. For $W\subseteq U$ open and $s\in\mathcal F(W)$ the sections $t_i=\varphi_i(W\cap U_i)(s|_{W\cap U_i})\in\mathcal G(W\cap U_i)$ satisfy $t_i|_{W\cap U_i\cap U_j}=\varphi_i(W\cap U_i\cap U_j)(s|)=\varphi_j(W\cap U_i\cap U_j)(s|)=t_j|_{W\cap U_i\cap U_j}$, so glue to a unique $t\in\mathcal G(W)$; put $\varphi(W)(s)=t$. This $\varphi(W)$ is a homomorphism (uniqueness of gluing), and compatible with restriction $W'\subseteq W$ (both $\varphi(W)(s)|_{W'}$ and $\varphi(W')(s|_{W'})$ restrict to $t_i|_{W'\cap U_i}$ on $W'\cap U_i$ and hence agree by (3)). For $W\subseteq U_i$ we get $\varphi(W)=\varphi_i(W)$ (the cover of $W$ by $W\cap U_j$ and the compatibility of the $\varphi_j$ with $\varphi_i$ show $t|_{W\cap U_j}=\varphi_i(W\cap U_j)(s|)$; use (3)). So $\varphi|_{U_i}=\varphi_i$. \hfill$\square$

\exercise{1.16}{\emph{Flasque Sheaves.} A sheaf $\mathcal F$ on a topological space $X$ is \emph{flasque} if for every inclusion $V\subseteq U$ of open sets, the restriction map $\mathcal F(U)\to\mathcal F(V)$ is surjective.
\begin{enumerate}[label=(\alph*),leftmargin=*]
\item Show that a constant sheaf on an irreducible topological space is flasque. See (I, Section 1) for irreducible topological spaces.
\item If $0\to\mathcal F'\to\mathcal F\to\mathcal F''\to0$ is an exact sequence of sheaves, and if $\mathcal F'$ is flasque, then for any open set $U$, the sequence $0\to\mathcal F'(U)\to\mathcal F(U)\to\mathcal F''(U)\to0$ of abelian groups is also exact.
\item If $0\to\mathcal F'\to\mathcal F\to\mathcal F''\to0$ is an exact sequence of sheaves, and if $\mathcal F'$ and $\mathcal F$ are flasque, then $\mathcal F''$ is flasque.
\item If $f\colon X\to Y$ is a continuous map, and if $\mathcal F$ is a flasque sheaf on $X$, then $f_*\mathcal F$ is a flasque sheaf on $Y$.
\item Let $\mathcal F$ be any sheaf on $X$. We define a new sheaf $\mathcal G$, called the sheaf of \emph{discontinuous sections} of $\mathcal F$ as follows. For each open set $U\subseteq X$, $\mathcal G(U)$ is the set of maps $s\colon U\to\bigcup_{P\in U}\mathcal F_P$ such that for each $P\in U$, $s(P)\in\mathcal F_P$. Show that $\mathcal G$ is a flasque sheaf, and that there is a natural injective morphism of $\mathcal F$ to $\mathcal G$.
\end{enumerate}}
\pf
(a) Let $\mathcal A$ be the constant sheaf with group $A$ on irreducible $X$. A nonempty open subset $U$ of an irreducible space is irreducible, hence connected. $\mathcal A(U)$ is the set of locally constant functions $U\to A$; on a connected $U$ such a function is constant (the fibres are open and partition $U$). So $\mathcal A(U)=A$ for $U\ne\emptyset$ and the restriction maps for nonempty $V\subseteq U$ are the identity of $A$, hence surjective; for $V=\emptyset$ the target is $0$. So $\mathcal A$ is flasque.

(b) Let $\alpha\colon\mathcal F'\to\mathcal F$, $\beta\colon\mathcal F\to\mathcal F''$. By Ex.~1.8 the sequence is exact at $\mathcal F'(U)$ and $\mathcal F(U)$ (in particular $\ker\beta(U)=\alpha(U)(\mathcal F'(U))$), so we must show that $\beta(U)$ is surjective. Let $s''\in\mathcal F''(U)$. Consider the set $\Sigma$ of pairs $(V,t)$, $V\subseteq U$ open, $t\in\mathcal F(V)$ with $\beta(t)=s''|_V$, ordered by $(V,t)\le(V_1,t_1)$ iff $V\subseteq V_1$ and $t_1|_V=t$. By Ex.~1.3(a) there is a cover $\{U_i\}$ of $U$ and $t_i\in\mathcal F(U_i)$ with $\beta(t_i)=s''|_{U_i}$, so $\Sigma\ne\emptyset$ (also $(\emptyset,0)\in\Sigma$). If $\{(V_a,t_a)\}$ is a chain, then with $V=\bigcup V_a$ the $t_a$ agree on overlaps (the intersection of two members is the smaller one) and glue to $t\in\mathcal F(V)$; $\beta(t)|_{V_a}=\beta(t_a)=s''|_{V_a}$ for all $a$, so $\beta(t)=s''|_V$ by (3) for $\mathcal F''$; thus $(V,t)$ is an upper bound. By Zorn's lemma there is a maximal $(V,t)$. Suppose $V\ne U$, pick $P\in U\setminus V$ and by the above an open $W\ni P$, $W\subseteq U$, and $w\in\mathcal F(W)$ with $\beta(w)=s''|_W$. On $V\cap W$, $\beta(t|_{V\cap W}-w|_{V\cap W})=0$, so $t|_{V\cap W}-w|_{V\cap W}=\alpha(r)$ with $r\in\mathcal F'(V\cap W)$ by Ex.~1.8. Since $\mathcal F'$ is flasque, $r=r'|_{V\cap W}$ with $r'\in\mathcal F'(W)$. Put $w'=w+\alpha(r')\in\mathcal F(W)$. Then $\beta(w')=s''|_W$ and $w'|_{V\cap W}=t|_{V\cap W}$. So $t$ and $w'$ glue to a section over $V\cup W$ lifting $s''$ (check as before), contradicting maximality as $P\in W\setminus V$. Hence $V=U$ and $s''=\beta(t)$.

(c) Let $V\subseteq U$ be open and $s''\in\mathcal F''(V)$. By (b) applied to $V$, $s''=\beta(t)$ for some $t\in\mathcal F(V)$. Since $\mathcal F$ is flasque there is $\tilde t\in\mathcal F(U)$ with $\tilde t|_V=t$. Then $\beta(\tilde t)\in\mathcal F''(U)$ restricts to $\beta(t)=s''$. So $\mathcal F''(U)\to\mathcal F''(V)$ is surjective.

(d) For open $V\subseteq U\subseteq Y$, $f^{-1}V\subseteq f^{-1}U$ are open in $X$, and $(f_*\mathcal F)(U)\to(f_*\mathcal F)(V)$ is the restriction $\mathcal F(f^{-1}U)\to\mathcal F(f^{-1}V)$, surjective since $\mathcal F$ is flasque.

(e) $\mathcal G$ is a presheaf under restriction of maps, and in fact a sheaf: for (3), if $s|_{U_i}=0$ for a cover then $s(P)=0$ for all $P$, i.e.\ $s=0$; for (4), compatible $s_i\in\mathcal G(U_i)$ define $s(P):=s_i(P)$ for any $i$ with $P\in U_i$, which is well defined by compatibility on overlaps; clearly $s|_{U_i}=s_i$. It is flasque: for $V\subseteq U$ and $s\in\mathcal G(V)$, extend $s$ by $0$ on $U\setminus V$ (i.e.\ $\tilde s(P)=0\in\mathcal F_P$ for $P\notin V$); then $\tilde s\in\mathcal G(U)$ and $\tilde s|_V=s$.

The map $\mathcal F\to\mathcal G$, $s\mapsto\bar s=(P\mapsto s_P)$ is a morphism of sheaves (compatible with restriction). It is injective: if $\bar s=0$ then all germs $s_P$ vanish, so $s=0$ by Fact A. \hfill$\square$

\exercise{1.17}{\emph{Skyscraper Sheaves.} Let $X$ be a topological space, let $P$ be a point, and let $A$ be an abelian group. Define a sheaf $i_P(A)$ on $X$ as follows: $i_P(A)(U)=A$ if $P\in U$, $0$ otherwise. Verify that the stalk of $i_P(A)$ is $A$ at every point $Q\in\{P\}^-$, and $0$ elsewhere, where $\{P\}^-$ denotes the closure of the set consisting of the point $P$. Hence the name ``skyscraper sheaf.'' Show that this sheaf could also be described as $i_*(A)$, where $A$ denotes the constant sheaf $A$ on the closed subspace $\{P\}^-$, and $i\colon\{P\}^-\to X$ is the inclusion.}
\pf
Restriction maps of $i_P(A)$ (call it $\mathcal S$): for $V\subseteq U$, $\rho_{UV}=\mathrm{id}_A$ if $P\in V$, and $0$ otherwise. This is a presheaf: for $W\subseteq V\subseteq U$, if $P\in W$ then $P\in V,U$ and all maps are identities; if $P\notin W$ then $\rho_{VW}\rho_{UV}$ and $\rho_{UW}$ both have target $0$. Also $\mathcal S(\emptyset)=0$.

\emph{Sheaf.} (3) Let $U=\bigcup V_i$, $s\in\mathcal S(U)$ with $s|_{V_i}=0$. If $P\notin U$, $\mathcal S(U)=0$. If $P\in U$, choose $i$ with $P\in V_i$; then $s|_{V_i}=s$ (identity), so $s=0$. (4) Let $s_i\in\mathcal S(V_i)$ be compatible. If $P\notin U$, all $V_i\not\ni P$, all $s_i=0$ and $s=0$ works. If $P\in U$, pick $i_0$ with $P\in V_{i_0}$ and put $s=s_{i_0}\in A$. For $i$ with $P\in V_i$ we have $s_i=s_i|_{V_i\cap V_{i_0}}=s_{i_0}|_{V_i\cap V_{i_0}}=s$ (identities, $P\in V_i\cap V_{i_0}$), so $s|_{V_i}=s_i$; for $i$ with $P\notin V_i$ both sides lie in $\mathcal S(V_i)=0$.

\emph{Stalks.} Let $Z=\{P\}^-$. If $Q\in Z$ then every open neighborhood $U$ of $Q$ meets $\{P\}$, i.e.\ contains $P$, so $\mathcal S(U)=A$ for all $U\ni Q$ with identity restriction maps, and $\mathcal S_Q=A$. If $Q\notin Z$, then $U_0=X\setminus Z$ is an open neighborhood of $Q$ which does not contain $P$; every neighborhood $U\subseteq U_0$ of $Q$ has $\mathcal S(U)=0$, so $\mathcal S_Q=0$.

\emph{Comparison with $i_*A$.} Let $\mathcal A$ be the constant sheaf on $Z$; $(i_*\mathcal A)(U)=\mathcal A(U\cap Z)$. The space $Z$ is irreducible (the closure of a point), so as in Ex.~1.16(a), $\mathcal A(W)=A$ for every nonempty open $W\subseteq Z$ with identity restrictions, and $\mathcal A(\emptyset)=0$. Moreover $U\cap Z\ne\emptyset$ iff $P\in U$ (as above: if $Q\in U\cap Z$ then $P\in U$; conversely $P\in U\cap Z$). Hence $(i_*\mathcal A)(U)=A$ if $P\in U$ and $0$ otherwise, and for $V\subseteq U$ the restriction map is the identity if $P\in V$, and zero (into $0$) otherwise: this is precisely $\mathcal S$. So $i_P(A)=i_*\mathcal A$. \hfill$\square$

\exercise{1.18}{\emph{Adjoint Property of $f^{-1}$.} Let $f\colon X\to Y$ be a continuous map of topological spaces. Show that for any sheaf $\mathcal F$ on $X$ there is a natural map $f^{-1}f_*\mathcal F\to\mathcal F$, and for any sheaf $\mathcal G$ on $Y$ there is a natural map $\mathcal G\to f_*f^{-1}\mathcal G$. Use these maps to show that there is a natural bijection of sets, for any sheaves $\mathcal F$ on $X$ and $\mathcal G$ on $Y$, $\operatorname{Hom}_X(f^{-1}\mathcal G,\mathcal F)=\operatorname{Hom}_Y(\mathcal G,f_*\mathcal F)$. Hence we say that $f^{-1}$ is a \emph{left adjoint} of $f_*$, and that $f_*$ is a \emph{right adjoint} of $f^{-1}$.}
\pf
Recall $f^{-1}\mathcal G=\mathcal P_{\mathcal G}^+$, where $\mathcal P_{\mathcal G}(U)=\varinjlim_{V\supseteq f(U)}\mathcal G(V)$, $V$ open in $Y$. Write $\rho$ for restriction maps in all sheaves.

\emph{Natural maps.} (i) Let $\mathcal F$ be a sheaf on $X$. For $U\subseteq X$ open and $V\supseteq f(U)$ open, $U\subseteq f^{-1}V$, and we have $\rho\colon(f_*\mathcal F)(V)=\mathcal F(f^{-1}V)\to\mathcal F(U)$. These are compatible with shrinking $V$, so they induce $\varepsilon'(U)\colon\mathcal P_{f_*\mathcal F}(U)\to\mathcal F(U)$, compatible with restriction in $U$: a morphism of presheaves $\mathcal P_{f_*\mathcal F}\to\mathcal F$. Since $\mathcal F$ is a sheaf, it factors uniquely through $f^{-1}f_*\mathcal F=\mathcal P^+_{f_*\mathcal F}$, giving $\varepsilon_{\mathcal F}\colon f^{-1}f_*\mathcal F\to\mathcal F$.

(ii) Let $\mathcal G$ be a sheaf on $Y$. For $V\subseteq Y$ open, $f(f^{-1}V)\subseteq V$, so $V$ is one of the open sets in the limit defining $\mathcal P_{\mathcal G}(f^{-1}V)$, giving $\mathcal G(V)\to\mathcal P_{\mathcal G}(f^{-1}V)\xrightarrow{\theta}(f^{-1}\mathcal G)(f^{-1}V)=(f_*f^{-1}\mathcal G)(V)$. These are compatible with restrictions, and define $\eta_{\mathcal G}\colon\mathcal G\to f_*f^{-1}\mathcal G$. Both constructions are clearly functorial in $\mathcal F$, resp.\ $\mathcal G$.

\emph{The bijection.} By the universal property of Prop.--Def.~1.2, $\operatorname{Hom}_X(f^{-1}\mathcal G,\mathcal F)=\operatorname{Hom}_{\mathrm{presheaves}}(\mathcal P_{\mathcal G},\mathcal F)$ (composition with $\theta$). By the universal property of the direct limit of groups, a morphism of presheaves $\alpha\colon\mathcal P_{\mathcal G}\to\mathcal F$ is the same as a family of homomorphisms $a_{U,V}\colon\mathcal G(V)\to\mathcal F(U)$, one for each pair of opens $U\subseteq X$, $V\subseteq Y$ with $f(U)\subseteq V$, such that (i) $a_{U,V'}\rho_{VV'}=a_{U,V}$ for $V'\subseteq V$ (with $f(U)\subseteq V'$), (ii) $\rho_{UU'}a_{U,V}=a_{U',V}$ for $U'\subseteq U$.

Given such a family put $b_V=a_{f^{-1}V,V}\colon\mathcal G(V)\to\mathcal F(f^{-1}V)=(f_*\mathcal F)(V)$. For $V'\subseteq V$: $\rho_{f^{-1}V,f^{-1}V'}b_V=a_{f^{-1}V',V}$ by (ii) and this equals $a_{f^{-1}V',V'}\rho_{VV'}=b_{V'}\rho_{VV'}$ by (i); so $b\colon\mathcal G\to f_*\mathcal F$ is a morphism of sheaves. Conversely, given a morphism $b\colon\mathcal G\to f_*\mathcal F$ put $a_{U,V}=\rho_{f^{-1}V,U}\,b_V$; (ii) is clear, and (i): $a_{U,V'}\rho_{VV'}=\rho_{f^{-1}V',U}b_{V'}\rho_{VV'}=\rho_{f^{-1}V',U}\rho_{f^{-1}V,f^{-1}V'}b_V=a_{U,V}$. The two constructions are inverse: from $b$ we get $a_{f^{-1}V,V}=b_V$; from $a$ we get $\rho_{f^{-1}V,U}a_{f^{-1}V,V}=a_{U,V}$ by (ii). This gives the bijection $\operatorname{Hom}_X(f^{-1}\mathcal G,\mathcal F)\cong\operatorname{Hom}_Y(\mathcal G,f_*\mathcal F)$, natural in $\mathcal F$ and $\mathcal G$ (it is compatible with composing on either side, since composition just acts on the groups $\mathcal G(V)$, $\mathcal F(U)$). Unwinding, the bijection takes $\alpha\mapsto f_*(\alpha)\circ\eta_{\mathcal G}$, and the identity of $f_*\mathcal F$ corresponds to $\varepsilon_{\mathcal F}$, and the identity of $f^{-1}\mathcal G$ corresponds to $\eta_{\mathcal G}$; i.e.\ the natural maps constructed above are the unit and counit of the adjunction. \hfill$\square$

\exercise{1.19}{\emph{Extending a Sheaf by Zero.} Let $X$ be a topological space, let $Z$ be a closed subset, let $i\colon Z\to X$ be the inclusion, and let $U=X-Z$ be the complementary open subset, and let $j\colon U\to X$ be its inclusion.
\begin{enumerate}[label=(\alph*),leftmargin=*]
\item Let $\mathcal F$ be a sheaf on $Z$. Show that the stalk $(i_*\mathcal F)_P$ of the direct image sheaf on $X$ is $\mathcal F_P$ if $P\in Z$, $0$ if $P\notin Z$. Hence we call $i_*\mathcal F$ the sheaf obtained by \emph{extending $\mathcal F$ by zero} outside $Z$. By abuse of notation we will sometimes write $\mathcal F$ instead of $i_*\mathcal F$, and say ``consider $\mathcal F$ as a sheaf on $X$,'' when we mean $i_*\mathcal F$.
\item Now let $\mathcal F$ be a sheaf on $U$. Let $j_!(\mathcal F)$ be the sheaf on $X$ associated to the presheaf $V\mapsto\mathcal F(V)$ if $V\subseteq U$, $V\mapsto0$ if $V\not\subseteq U$. Show that the stalk $(j_!(\mathcal F))_P$ is equal to $\mathcal F_P$ if $P\in U$, $0$ if $P\notin U$, and show that $j_!\mathcal F$ is the only sheaf on $X$ which has this property, and whose restriction to $U$ is $\mathcal F$. We call $j_!\mathcal F$ the sheaf obtained by \emph{extending $\mathcal F$ by zero} outside $U$.
\item Now let $\mathcal F$ be a sheaf on $X$. Show that there is an exact sequence of sheaves on $X$, $0\to j_!(\mathcal F|_U)\to\mathcal F\to i_*(\mathcal F|_Z)\to0$.
\end{enumerate}}
\pf
(a) $(i_*\mathcal F)(W)=\mathcal F(W\cap Z)$. Let $P\in Z$. The map $\Phi\colon(i_*\mathcal F)_P=\varinjlim_{W\ni P}\mathcal F(W\cap Z)\to\mathcal F_P=\varinjlim_{W'\ni P,\ W'\subseteq Z\text{ open}}\mathcal F(W')$, $\langle W,s\rangle\mapsto\langle W\cap Z,s\rangle$, is well defined. It is surjective, since every open $W'\subseteq Z$ is of the form $W\cap Z$ with $W$ open in $X$. It is injective: if $\Phi\langle W,s\rangle=0$ then $s|_{W''}=0$ for some open $W''\ni P$ of $Z$; write $W''=W_1\cap Z$ with $W_1$ open in $X$; replacing $W_1$ by $W_1\cap W$ (which still satisfies $W_1\cap Z=W''$ since $W''\subseteq W\cap Z$) we get $W_1\subseteq W$ and $s|_{W_1\cap Z}=0$, i.e.\ $\langle W,s\rangle=0$ in $(i_*\mathcal F)_P$. If $P\notin Z$, then $X\setminus Z$ is an open neighborhood of $P$, and for all open $W\ni P$ contained in it $(i_*\mathcal F)(W)=\mathcal F(\emptyset)=0$, so the stalk is $0$.

(b) Let $\mathcal P$ be the presheaf in the statement; $\mathcal P^+=j_!\mathcal F$ and $\mathcal P^+_P=\mathcal P_P$. If $P\in U$, the open neighborhoods contained in $U$ are cofinal and $\mathcal P(V)=\mathcal F(V)$ for them, so $\mathcal P_P=\mathcal F_P$. If $P\notin U$, no neighborhood of $P$ is contained in $U$, so $\mathcal P(V)=0$ for all $V\ni P$ and $\mathcal P_P=0$.

Restriction to $U$: $\mathcal P|_U=\mathcal F$ is a sheaf on $U$, and $\theta|_U\colon\mathcal P|_U=\mathcal F\to(j_!\mathcal F)|_U$ is a morphism of sheaves on $U$ which on stalks (at $Q\in U$) is $\mathcal P_Q\to\mathcal P^+_Q$, an isomorphism; by Prop.~1.1 $(j_!\mathcal F)|_U\cong\mathcal F$ (compatibly with the identification of stalks).

Uniqueness: let $\mathcal G$ be a sheaf on $X$ with $\mathcal G_P=0$ for $P\notin U$ and an isomorphism $\psi\colon\mathcal F\to\mathcal G|_U$. Define $\kappa\colon\mathcal P\to\mathcal G$ by $\kappa(V)=\psi(V)$ for $V\subseteq U$ and $\kappa(V)=0$ otherwise. This commutes with restrictions: for $V'\subseteq V$, if $V\subseteq U$ it is the compatibility of $\psi$; if $V\not\subseteq U$ then $\mathcal P(V)=0$ and both composites $\mathcal P(V)\to\mathcal G(V')$ are $0$. By Prop.--Def.~1.2, $\kappa$ factors through $\bar\kappa\colon j_!\mathcal F=\mathcal P^+\to\mathcal G$. On stalks at $Q\in U$, $\bar\kappa_Q=\kappa_Q=\psi_Q$ is an isomorphism; at $Q\notin U$, both stalks are zero. By Prop.~1.1, $\bar\kappa$ is an isomorphism, restricting to $\psi$ on $U$ (up to the identification above). So $j_!\mathcal F$ is the unique such sheaf.

(c) \emph{The maps.} Let $\mathcal P(V)=\mathcal F(V)$ if $V\subseteq U$ and $0$ otherwise (this is the presheaf for $j_!(\mathcal F|_U)$; note $(\mathcal F|_U)(V)=\mathcal F(V)$ for $V\subseteq U$). The maps $\mathcal P(V)\to\mathcal F(V)$ (identity for $V\subseteq U$, zero otherwise) are compatible with restriction, so they induce $\alpha\colon j_!(\mathcal F|_U)\to\mathcal F$. Let $\beta=\eta_{\mathcal F}\colon\mathcal F\to i_*i^{-1}\mathcal F=i_*(\mathcal F|_Z)$ be the natural map of Ex.~1.18 (for the map $i$).

\emph{Exactness} is checked on stalks (Ex.~1.2(c)). Recall $(\mathcal F|_Z)_P=\mathcal F_P$ for $P\in Z$ (text). At $P\in U$: the stalk sequence is $0\to(j_!\mathcal F|_U)_P=\mathcal F_P\xrightarrow{\alpha_P}\mathcal F_P\to(i_*\mathcal F|_Z)_P=0$ (by (a), as $P\notin Z$); here $\alpha_P$ is the identity (computed on the presheaf $\mathcal P$, where it is the identity on germs of sections over $V\subseteq U$). This is exact. At $P\in Z$: the stalk sequence is $0\to0\to\mathcal F_P\xrightarrow{\beta_P}\mathcal F_P\to0$, since $(j_!\mathcal F|_U)_P=0$ by (b) and $(i_*\mathcal F|_Z)_P=(\mathcal F|_Z)_P=\mathcal F_P$ by (a). Here $\beta_P$ sends the germ of $s\in\mathcal F(W)$ to the germ at $P$ of the image of $s$ in $(i^{-1}\mathcal F)(W\cap Z)$, which under the identification of the stalk of $i^{-1}\mathcal F$ with $\mathcal F_P$ is $s_P$; so $\beta_P=\mathrm{id}$ and the sequence is exact. Hence $0\to j_!(\mathcal F|_U)\to\mathcal F\to i_*(\mathcal F|_Z)\to0$ is exact (in particular the composition $\beta\alpha$ has zero stalks, hence is $0$ by Fact A). \hfill$\square$

\exercise{1.20}{\emph{Subsheaf with Supports.} Let $Z$ be a closed subset of $X$, and let $\mathcal F$ be a sheaf on $X$. We define $\Gamma_Z(X,\mathcal F)$ to be the subgroup of $\Gamma(X,\mathcal F)$ consisting of all sections whose support (Ex.~1.14) is contained in $Z$.
\begin{enumerate}[label=(\alph*),leftmargin=*]
\item Show that the presheaf $V\mapsto\Gamma_{Z\cap V}(V,\mathcal F|_V)$ is a sheaf. It is called the \emph{subsheaf of $\mathcal F$ with supports in $Z$}, and is denoted $\mathcal H^0_Z(\mathcal F)$.
\item Let $U=X-Z$, and let $j\colon U\to X$ be the inclusion. Show there is an exact sequence of sheaves on $X$ \[0\to\mathcal H^0_Z(\mathcal F)\to\mathcal F\to j_*(\mathcal F|_U).\] Furthermore, if $\mathcal F$ is flasque, the map $\mathcal F\to j_*(\mathcal F|_U)$ is surjective.
\end{enumerate}}
\pf
(a) Let $\mathcal H(V)=\{s\in\mathcal F(V):\operatorname{Supp}s\subseteq Z\cap V\}=\{s\in\mathcal F(V): s_P=0\ \forall P\in V\setminus Z\}$. For $W\subseteq V$ we have $(s|_W)_P=s_P$, so restriction maps $\mathcal H(V)\to\mathcal H(W)$; $\mathcal H$ is a subpresheaf of $\mathcal F$ and a subgroup in each degree. (3) is inherited from $\mathcal F$. (4): if $s_i\in\mathcal H(V_i)$ are compatible, they glue to $s\in\mathcal F(V)$, $V=\bigcup V_i$, and for $P\in V\setminus Z$ pick $i$ with $P\in V_i$: $s_P=(s_i)_P=0$. So $s\in\mathcal H(V)$. Thus $\mathcal H=\mathcal H^0_Z(\mathcal F)$ is a sheaf, a subsheaf of $\mathcal F$.

(b) Let $\gamma\colon\mathcal F\to j_*(\mathcal F|_U)$ be given by $\gamma(V)\colon\mathcal F(V)\to\mathcal F(V\cap U)=(j_*(\mathcal F|_U))(V)$, the restriction; clearly a morphism of sheaves. We have $\ker\gamma(V)=\{s\in\mathcal F(V):s|_{V\cap U}=0\}$. By Fact A, $s|_{V\cap U}=0$ iff $s_P=0$ for all $P\in V\cap U=V\setminus Z$ (the converse implication: if $s|_{V\cap U}=0$ then its germs vanish), i.e.\ iff $s\in\mathcal H^0_Z(\mathcal F)(V)$. So $\ker\gamma=\mathcal H^0_Z(\mathcal F)$ as subsheaves of $\mathcal F$, which means that $0\to\mathcal H^0_Z(\mathcal F)\to\mathcal F\to j_*(\mathcal F|_U)$ is exact (injectivity at the left since it is a subsheaf, exactness in the middle since the image of the inclusion is $\ker\gamma$).

If $\mathcal F$ is flasque, $\gamma(V)\colon\mathcal F(V)\to\mathcal F(V\cap U)$ is surjective for every open $V$ ($V\cap U\subseteq V$), so $\gamma$ is surjective as a morphism of sheaves (by Ex.~1.3(a), with the trivial cover). \hfill$\square$

\exercise{1.21}{\emph{Some Examples of Sheaves on Varieties.} Let $X$ be a variety over an algebraically closed field $k$, as in Ch.~I. Let $\mathcal O_X$ be the sheaf of regular functions on $X$ (1.0.1).
\begin{enumerate}[label=(\alph*),leftmargin=*]
\item Let $Y$ be a closed subset of $X$. For each open set $U\subseteq X$, let $\mathcal I_Y(U)$ be the ideal in the ring $\mathcal O_X(U)$ consisting of those regular functions which vanish at all points of $Y\cap U$. Show that the presheaf $U\mapsto\mathcal I_Y(U)$ is a sheaf. It is called the \emph{sheaf of ideals} $\mathcal I_Y$ of $Y$, and it is a subsheaf of the sheaf of rings $\mathcal O_X$.
\item If $Y$ is a subvariety, then the quotient sheaf $\mathcal O_X/\mathcal I_Y$ is isomorphic to $i_*(\mathcal O_Y)$, where $i\colon Y\to X$ is the inclusion, and $\mathcal O_Y$ is the sheaf of regular functions on $Y$.
\item Now let $X=\mathbf P^1$, and let $Y$ be the union of two distinct points $P,Q\in X$. Then there is an exact sequence of sheaves on $X$, where $\mathcal F=i_*\mathcal O_P\oplus i_*\mathcal O_Q$: \[0\to\mathcal I_Y\to\mathcal O_X\to\mathcal F\to0.\] Show that the induced map on global sections $\Gamma(X,\mathcal O_X)\to\Gamma(X,\mathcal F)$ is not surjective. This shows that the global section functor $\Gamma(X,\cdot)$ is not exact (cf.\ Ex.~1.8 which shows that it is left exact).
\item Again let $X=\mathbf P^1$, and let $\mathcal O$ be the sheaf of regular functions. Let $\mathcal K$ be the constant sheaf on $X$ associated to the function field $K$ of $X$. Show that there is a natural injection $\mathcal O\to\mathcal K$. Show that the quotient sheaf $\mathcal K/\mathcal O$ is isomorphic to the direct sum $\sum_{P\in X}i_P(I_P)$, where $I_P$ is the group $K/\mathcal O_P$, and $i_P(I_P)$ denotes the skyscraper sheaf (Ex.~1.17) given by $I_P$ at the point $P$.
\item Finally show that in the case of (d) the sequence $0\to\Gamma(X,\mathcal O)\to\Gamma(X,\mathcal K)\to\Gamma(X,\mathcal K/\mathcal O)\to0$ is exact. (This is an analogue of what is called the ``first Cousin problem'' in several complex variables. See Gunning and Rossi [1, p.~248].)
\end{enumerate}}
\pf
Conventions: points of a variety are its closed points; $\mathcal O_X(U)$ is the ring of regular functions on $U$, and for $f\in\mathcal O_X(U)$, $P\in U$, $f(P)\in k$ is its value, which depends only on the germ $f_P\in\mathcal O_{X,P}$ and is compatible with restrictions.

(a) $\mathcal I_Y(U)=\{f\in\mathcal O_X(U):f(P)=0\ \forall P\in Y\cap U\}$ is an ideal (value at $P$ is a ring homomorphism), and restriction preserves the condition, so $\mathcal I_Y$ is a subpresheaf of ideals of $\mathcal O_X$. (3) is inherited from $\mathcal O_X$. (4): if $f_i\in\mathcal I_Y(U_i)$ agree on overlaps, they glue to $f\in\mathcal O_X(U)$, $U=\bigcup U_i$, and for $P\in Y\cap U$, $P\in U_i$ for some $i$ and $f(P)=f_i(P)=0$. So $f\in\mathcal I_Y(U)$.

(b) Let $Y$ be a closed subvariety (closed irreducible subset with induced structure) and $i\colon Y\to X$. For $U\subseteq X$ open, restriction of regular functions gives a ring homomorphism $r(U)\colon\mathcal O_X(U)\to\mathcal O_Y(U\cap Y)=(i_*\mathcal O_Y)(U)$ (a function regular on $U$ is locally of the form $g/h$ with $g,h$ homogeneous polynomials of the same degree, $h\ne0$ at the relevant points; its restriction to $Y\cap U$ is locally given by the same expression, hence is regular on $Y\cap U$ as a variety). Clearly $r$ is a morphism of sheaves. Its kernel is $\ker r(U)=\{f:f|_{Y\cap U}=0\}=\mathcal I_Y(U)$, so $\ker r=\mathcal I_Y$.

\emph{$r$ is surjective} (Ex.~1.3(a)): let $\varphi\in\mathcal O_Y(U\cap Y)$ and $P\in U$. If $P\notin Y$, then $W=U\setminus Y$ is an open neighborhood of $P$ with $W\cap Y=\emptyset$, so $\varphi|_W\in(i_*\mathcal O_Y)(W)=\mathcal O_Y(\emptyset)=0$ is the image of $0$. If $P\in Y\cap U$, by definition of regular function on $Y\subseteq X\subseteq\mathbf P^n$ there is an open neighborhood $W_0\ni P$ in $Y\cap U$ and homogeneous polynomials $g,h$ of the same degree with $h$ nowhere zero on $W_0$ and $\varphi=g/h$ on $W_0$. Choose $W_1$ open in $X$ with $W_1\cap Y=W_0$ (possible since $W_0$ is open in $Y$) and let $W=\{Q\in W_1\cap U:h(Q)\ne0\}$, an open neighborhood of $P$ in $X$. Then $f=g/h\in\mathcal O_X(W)$ and $r(W)(f)=f|_{W\cap Y}=\varphi|_{W\cap Y}$, as $W\cap Y\subseteq W_0$. The open sets $W$ obtained (for all $P\in U$) cover $U$, so $r$ is surjective.

Therefore $\operatorname{im}r=i_*\mathcal O_Y$ and by Ex.~1.7(a), $\mathcal O_X/\mathcal I_Y=\mathcal O_X/\ker r\cong\operatorname{im}r=i_*\mathcal O_Y$.

(c) Here $\mathcal O_P$ denotes the sheaf of regular functions on the one-point variety $\{P\}$, which is the constant sheaf $k$ on $\{P\}$; hence $i_*\mathcal O_P$ is the skyscraper sheaf $i_P(k)$ of Ex.~1.17 (note $\{P\}$ is closed). So $\mathcal F(U)=k^{\epsilon_P(U)+\epsilon_Q(U)}$ where $\epsilon_R(U)=1$ if $R\in U$ and $0$ otherwise. Define $\psi\colon\mathcal O_X\to\mathcal F$ by $\psi(U)(f)=(f(P)\text{ if }P\in U,\ f(Q)\text{ if }Q\in U)$; a morphism of sheaves (indeed (b) applied to $Y=\{P\}$ and $Y=\{Q\}$). Its kernel is $\ker\psi(U)=\{f:f(R)=0\text{ for all }R\in Y\cap U\}=\mathcal I_Y(U)$ with $Y=\{P,Q\}$, so $\ker\psi=\mathcal I_Y$, giving exactness at $\mathcal O_X$ and at $\mathcal I_Y$ (inclusion). For surjectivity use stalks (Ex.~1.2(b)): $\mathcal F_P=k$ (from $i_*\mathcal O_P$; the summand for $Q$ contributes $0$ since $P\notin\{Q\}^-=\{Q\}$), $\mathcal F_Q=k$, and $\mathcal F_R=0$ for $R\ne P,Q$ (Ex.~1.17). The map $\psi_P\colon\mathcal O_{X,P}\to k$ is evaluation, which is surjective (constants), similarly $\psi_Q$. So $\psi$ is surjective, and $0\to\mathcal I_Y\to\mathcal O_X\to\mathcal F\to0$ is exact.

On global sections: $\Gamma(X,\mathcal O_X)=k$ (the regular functions on $\mathbf P^1$ are constants, Thm.~I.3.4(a)), $\Gamma(X,\mathcal F)=k\oplus k$, and the map is $c\mapsto(c,c)$ (a constant function takes the same value at $P$ and $Q$). Its image is the diagonal, so it is not surjective. Hence $\Gamma(X,\cdot)$ is not exact.

(d) Let $X=\mathbf P^1$ and $K=k(X)=k(x)$. $X$ is irreducible, so $\mathcal K(U)=K$ for every nonempty open $U$, with identity restrictions (Ex.~1.16(a)), and $\mathcal K(\emptyset)=0$. A regular function $\varphi$ on a nonempty open $U$ is locally a quotient of homogeneous polynomials of equal degree; any two such local expressions agree on the (nonempty, by irreducibility) overlap and hence define the same element of $K$ (an element of $K$ is an equivalence class of regular functions on nonempty open sets, two being equivalent if they agree on a common nonempty open set). So there is a well-defined map $\iota(U)\colon\mathcal O(U)\to K$, a ring homomorphism compatible with restrictions. It is injective: if $\iota(U)\varphi=0$, then $\varphi$ vanishes on a nonempty open subset of $U$, which is dense in $U$, and the zero set of $\varphi$ is closed, so $\varphi=0$. Thus $\iota\colon\mathcal O\to\mathcal K$ is an injective morphism of sheaves (Ex.~1.2(b) / injectivity on sections). At the stalk $P$, $\mathcal O_P\subseteq K$ is the local ring.

\emph{The sheaf $\mathcal G$.} For $U$ open let $\mathcal G(U)=\bigoplus_{P\in U}I_P=\{(g_P)_{P\in U}\in\prod_{P\in U}I_P:\ g_P=0\text{ for all but finitely many }P\}$ with restriction ``forget the coordinates outside the smaller set.'' Then $\mathcal G$ is a presheaf. It is a sheaf: (3) is clear coordinatewise; for (4), if $g^i\in\mathcal G(U_i)$ are compatible, define $g_P=g^i_P$ for any $i$ with $P\in U_i$ (well defined by compatibility); $X=\mathbf P^1$ is noetherian so $U=\bigcup U_i$ is covered by finitely many $U_{i_1},\dots,U_{i_n}$ and the support of $g$ is contained in the union of the finite supports of $g^{i_1},\dots,g^{i_n}$, hence finite. (This $\mathcal G$ is the subsheaf of finitely supported sections of the product $\prod_Pi_P(I_P)$; it is the object denoted $\sum_Pi_P(I_P)$; this is the reading of the notation used here.)

\emph{The map.} A rational function $f\in K$ lies in $\mathcal O_P$ for all but finitely many closed points $P$ (write $f=a/b$ with $a,b$ homogeneous of the same degree, $b\ne0$; $f\in\mathcal O_P$ whenever $b(P)\ne0$, which holds outside the finite zero set of $b$). Hence for $U\ne\emptyset$ the map $\lambda(U)\colon K\to\mathcal G(U)$, $f\mapsto(f\bmod\mathcal O_P)_{P\in U}$ is well defined, a homomorphism, compatible with restriction (and $\lambda(\emptyset)=0$). Thus $\lambda\colon\mathcal K\to\mathcal G$ is a morphism of sheaves, and $\lambda\iota=0$ because $\varphi\in\mathcal O(U)$ has $\iota(\varphi)\in\mathcal O_P$ for all $P\in U$. So $\lambda$ factors through the presheaf quotient $U\mapsto\mathcal K(U)/\mathcal O(U)$, and hence through a morphism $\bar\lambda\colon\mathcal K/\mathcal O\to\mathcal G$.

\emph{$\bar\lambda$ is an isomorphism.} Check stalks (Prop.~1.1). $(\mathcal K/\mathcal O)_P=\mathcal K_P/\mathcal O_P=K/\mathcal O_P=I_P$. For $\mathcal G$, let $\pi_P\colon\mathcal G_P\to I_P$ be induced by the projection $\mathcal G(U)\to I_P$, $P\in U$. It is surjective: $a\in I_P$ is the image of the section supported at $P$ with value $a$. It is injective: if the germ of $g\in\mathcal G(U)$ has $g_P=0$ (the $P$-coordinate), then the support $S$ of $g$ is finite, so $W=U\setminus(S\setminus\{P\})$ is an open neighborhood of $P$ (finite sets of closed points are closed) and $g|_W$ has support contained in $\{P\}$ with zero $P$-coordinate, i.e.\ $g|_W=0$. So $\mathcal G_P\cong I_P$, and the composite $K/\mathcal O_P=(\mathcal K/\mathcal O)_P\xrightarrow{\bar\lambda_P}\mathcal G_P\xrightarrow{\pi_P}I_P$ is $f\bmod\mathcal O_P\mapsto f\bmod\mathcal O_P$, the identity. Hence $\bar\lambda_P$ is an isomorphism for each $P$ and $\mathcal K/\mathcal O\cong\mathcal G=\sum_Pi_P(I_P)$.

(e) By Ex.~1.6(a) (applied to the subsheaf $\iota(\mathcal O)\subseteq\mathcal K$) $0\to\mathcal O\to\mathcal K\to\mathcal K/\mathcal O\to0$ is exact, so by Ex.~1.8 the sequence $0\to\Gamma(X,\mathcal O)\to\Gamma(X,\mathcal K)\to\Gamma(X,\mathcal K/\mathcal O)$ is exact. By (d), $\Gamma(X,\mathcal K/\mathcal O)\cong\mathcal G(X)=\bigoplus_{P\in X}I_P$ and $\Gamma(X,\mathcal K)=K$, with the map $K\to\bigoplus_PK/\mathcal O_P$, $f\mapsto(f\bmod\mathcal O_P)_P$. We must show it is surjective: given finitely many distinct points $P_1,\dots,P_n$ and $g_1,\dots,g_n\in K$, find $f\in K$ with $f-g_m\in\mathcal O_{P_m}$ for each $m$, and $f\in\mathcal O_R$ for all other points $R$.

It suffices to treat each $m$ separately: find $f_m\in K$ with $f_m-g_m\in\mathcal O_{P_m}$ and $f_m\in\mathcal O_R$ for all $R\ne P_m$; then $f=\sum f_m$ works (at $P_m$, $f-g_m=(f_m-g_m)+\sum_{l\ne m}f_l\in\mathcal O_{P_m}$; at $R\notin\{P_m\}$ all $f_l\in\mathcal O_R$).

Fix $P=P_m$, $g=g_m$. Choose homogeneous coordinates on $\mathbf P^1$ (a linear change of coordinates) so that $P$ is the point at infinity $[0:1]$ and put $x$ = the coordinate function on $\mathbf A^1=\mathbf P^1\setminus\{P\}$, so $K=k(x)$, $k[x]\subseteq\mathcal O_R$ for every $R\ne P$ (a polynomial in $x$ is regular on $\mathbf A^1$), and $\mathcal O_P=\{a/b: a,b\in k[x], b\ne0,\ \deg a\le\deg b\}$. Write $g=a/b$ with $a,b\in k[x]$, $b\ne0$, and divide: $a=qb+r$ with $q,r\in k[x]$, $\deg r<\deg b$. Then $f_m:=q\in k[x]$ is regular at every $R\ne P$, and $g-f_m=r/b\in\mathcal O_P$. This proves surjectivity, so the sequence $0\to\Gamma(X,\mathcal O)\to\Gamma(X,\mathcal K)\to\Gamma(X,\mathcal K/\mathcal O)\to0$ is exact. \hfill$\square$

\exercise{1.22}{\emph{Gluing Sheaves.} Let $X$ be a topological space, let $\mathfrak U=\{U_i\}$ be an open cover of $X$, and suppose we are given for each $i$ a sheaf $\mathcal F_i$ on $U_i$, and for each $i,j$ an isomorphism $\varphi_{ij}\colon\mathcal F_i|_{U_i\cap U_j}\xrightarrow{\sim}\mathcal F_j|_{U_i\cap U_j}$, such that (1) for each $i$, $\varphi_{ii}=\mathrm{id}$, and (2) for each $i,j,k$, $\varphi_{ik}=\varphi_{jk}\circ\varphi_{ij}$ on $U_i\cap U_j\cap U_k$. Then there exists a unique sheaf $\mathcal F$ on $X$, together with isomorphisms $\psi_i\colon\mathcal F|_{U_i}\xrightarrow{\sim}\mathcal F_i$, such that for each $i,j$, $\psi_j=\varphi_{ij}\circ\psi_i$ on $U_i\cap U_j$. We say loosely that $\mathcal F$ is obtained by \emph{gluing} the sheaves $\mathcal F_i$ via the isomorphisms $\varphi_{ij}$.}
\pf
\emph{Remark on the statement.} The scan (read at high resolution) prints the compatibility as $\psi_j=\varphi_{ij}\circ\psi_i$ on $U_i\cap U_j$, which is type-correct since $\varphi_{ij}$ goes from $\mathcal F_i$ to $\mathcal F_j$; this is the condition used. Uniqueness is understood up to a unique isomorphism compatible with the $\psi_i$.

\emph{Construction.} For $V\subseteq X$ open let
\[
\mathcal F(V)=\Bigl\{(s_i)_i\in\prod_i\mathcal F_i(V\cap U_i):\ \varphi_{ij}\bigl(s_i|_{V\cap U_i\cap U_j}\bigr)=s_j|_{V\cap U_i\cap U_j}\ \text{for all }i,j\Bigr\},
\]
a subgroup of the product, with componentwise restrictions; clearly a presheaf with $\mathcal F(\emptyset)=0$.

(3) Let $V=\bigcup V_a$ and $s=(s_i)\in\mathcal F(V)$ with $s|_{V_a}=0$ for all $a$. Then for every $i$, $s_i|_{V_a\cap U_i}=0$ for all $a$, and the $V_a\cap U_i$ cover $V\cap U_i$, so $s_i=0$.

(4) Let $s^a=(s^a_i)_i\in\mathcal F(V_a)$ with $s^a|_{V_a\cap V_b}=s^b|_{V_a\cap V_b}$. Fix $i$; the sections $s^a_i\in\mathcal F_i(V_a\cap U_i)$ agree on overlaps, so glue to $s_i\in\mathcal F_i(V\cap U_i)$ with $s_i|_{V_a\cap U_i}=s^a_i$. For $i,j$, the sections $\varphi_{ij}(s_i|_{V\cap U_i\cap U_j})$ and $s_j|_{V\cap U_i\cap U_j}$ of $\mathcal F_j$ have the same restriction to each $V_a\cap U_i\cap U_j$ (namely $\varphi_{ij}(s^a_i|)=s^a_j|$), hence are equal. So $s=(s_i)\in\mathcal F(V)$ and $s|_{V_a}=s^a$. Thus $\mathcal F$ is a sheaf.

\emph{The isomorphisms $\psi_i$.} For $V\subseteq U_i$ define $\psi_i(V)\colon\mathcal F(V)\to\mathcal F_i(V)$, $(s_l)_l\mapsto s_i$ (note $V\cap U_i=V$); a morphism of sheaves $\mathcal F|_{U_i}\to\mathcal F_i$. For $V\subseteq U_i\cap U_j$, $\psi_j(V)(s)=s_j=\varphi_{ij}(s_i)=\varphi_{ij}\psi_i(V)(s)$ by the definition of $\mathcal F(V)$. So $\psi_j=\varphi_{ij}\psi_i$ on $U_i\cap U_j$. $\psi_i(V)$ is bijective: given $t\in\mathcal F_i(V)$, $V\subseteq U_i$, put $s_l=\varphi_{il}(t|_{V\cap U_l})\in\mathcal F_l(V\cap U_l)$ (here $V\cap U_l\subseteq U_i\cap U_l$); then $s_i=\varphi_{ii}(t)=t$, and by (2), $\varphi_{lm}(s_l|_{V\cap U_l\cap U_m})=\varphi_{lm}\varphi_{il}(t|)=\varphi_{im}(t|)=s_m|_{V\cap U_l\cap U_m}$, so $(s_l)\in\mathcal F(V)$ and $\psi_i(V)(s)=t$: surjective. If $(s_l)\in\mathcal F(V)$ has $s_i=0$, then $s_l=\varphi_{il}(s_i|)=0$ for all $l$: injective. A morphism of sheaves bijective on all sections is an isomorphism, so each $\psi_i$ is an isomorphism.

\emph{Uniqueness.} Let $\mathcal F'$ be another sheaf on $X$ with isomorphisms $\psi'_i\colon\mathcal F'|_{U_i}\to\mathcal F_i$ with $\psi'_j=\varphi_{ij}\psi'_i$ on $U_i\cap U_j$. For $V$ open and $s'\in\mathcal F'(V)$ put $t_i=\psi'_i(s'|_{V\cap U_i})\in\mathcal F_i(V\cap U_i)$. Then $\varphi_{ij}(t_i|_{V\cap U_i\cap U_j})=\varphi_{ij}\psi'_i(s'|_{V\cap U_i\cap U_j})=\psi'_j(s'|_{V\cap U_i\cap U_j})=t_j|_{V\cap U_i\cap U_j}$, so $\eta(V)(s'):=(t_i)_i\in\mathcal F(V)$. This defines a morphism $\eta\colon\mathcal F'\to\mathcal F$, and for $V\subseteq U_i$, $\psi_i\eta(V)(s')=t_i=\psi'_i(s')$, i.e.\ $\psi_i\eta=\psi'_i$ on $U_i$. It is the only morphism with this property, since the $i$-th component of $\eta(V)(s')$ is forced to be $\psi'_i(s'|_{V\cap U_i})$. Finally $\eta$ is an isomorphism: injective by (3) for $\mathcal F'$, since $\eta(V)(s')=0$ forces $s'|_{V\cap U_i}=0$ ($\psi'_i$ injective) for all $i$, hence $s'=0$; surjective since given $(t_i)\in\mathcal F(V)$, the sections $s'_i=\psi_i'^{-1}(t_i)\in\mathcal F'(V\cap U_i)$ satisfy $\psi'_j(s'_i|_{V\cap U_i\cap U_j})=\varphi_{ij}\psi'_i(s'_i|)=\varphi_{ij}(t_i|)=t_j|=\psi'_j(s'_j|)$, so $s'_i|_{V\cap U_i\cap U_j}=s'_j|_{V\cap U_i\cap U_j}$ ($\psi'_j$ injective), and they glue to $s'\in\mathcal F'(V)$ with $\eta(V)(s')=(t_i)$. Hence $\mathcal F$ is unique up to a unique isomorphism compatible with the $\psi_i$. \hfill$\square$

\section*{2. Schemes}

Throughout this section I use the numbering of Chapter~II of the scanned edition: Lemma~2.1, Prop.~2.2 (stalks and sections of $\mathcal O$ on $\Spec A$), Prop.~2.3 (morphisms of locally ringed spaces $\Spec B\to\Spec A$ come from ring homomorphisms), Lemma~2.4 and Prop.~2.5 (the scheme $\operatorname{Proj}S$), Prop.~2.6 (the functor $t$ from varieties to schemes). For a point $x$ of a scheme, $\mathcal O_x$ is the local ring, $\fm_x$ its maximal ideal and $k(x)=\mathcal O_x/\fm_x$ the residue field. For $f\in A$, $D(f)=\Spec A\setminus V(f)$. Recall that a morphism of schemes is a morphism of locally ringed spaces, and that a morphism of sheaves is an isomorphism iff it is an isomorphism on all stalks.

\exercise{2.1}{Let $A$ be a ring, let $X=\Spec A$, let $f\in A$ and let $D(f)\subseteq X$ be the open complement of $V((f))$. Show that the locally ringed space $(D(f),\mathcal O_X|_{D(f)})$ is isomorphic to $\Spec A_f$.}
\pf
Let $\varphi:A\to A_f$ be the localization map and $g=(g,g^\#):Y=\Spec A_f\to X$ the morphism of Prop.~2.3(b). The prime ideals of $A_f$ correspond bijectively to the primes of $A$ not containing $f$ via $\mathfrak q\mapsto\varphi^{-1}(\mathfrak q)$; thus $g$ is injective with image exactly $D(f)$. It is a homeomorphism onto $D(f)$: for an ideal $\mathfrak b\subseteq A_f$ one has $g(V(\mathfrak b))=V(\varphi^{-1}\mathfrak b)\cap D(f)$, and conversely for an ideal $\fa\subseteq A$ one has $g^{-1}(V(\fa))=V(\fa A_f)$ and $g(V(\fa A_f))=V(\fa)\cap D(f)$. So $g$ maps the closed sets of $Y$ exactly onto the closed sets of the subspace $D(f)$.

It remains to see that $g^\#:\mathcal O_X|_{D(f)}\to g_*\mathcal O_Y$ is an isomorphism of sheaves; we check this on stalks. Let $\mathfrak q\in Y$, $\fp=\varphi^{-1}(\mathfrak q)\in D(f)$. By Prop.~2.3(b) the stalk map is the localization $\varphi_{\mathfrak q}:A_\fp\to(A_f)_{\mathfrak q}$ of $\varphi$ (and $(g_*\mathcal O_Y)_{g(\mathfrak q)}=\mathcal O_{Y,\mathfrak q}$ because $g$ is a homeomorphism onto the open set $D(f)$). Since $f\notin\fp$, $f$ is a unit in $A_\fp$, so $A_\fp=(A_f)_{\fp A_f}$, and $\mathfrak q=\fp A_f$; hence $\varphi_{\mathfrak q}$ is the identity of $A_\fp$ under these identifications, in particular an isomorphism. Therefore $g$ is an isomorphism of locally ringed spaces $\Spec A_f\cong(D(f),\mathcal O_X|_{D(f)})$.\hfill$\square$

\exercise{2.2}{Let $(X,\mathcal O_X)$ be a scheme, and let $U$ be any open subset of $X$. Show that $(U,\mathcal O_X|_U)$ is a scheme. We call this the \emph{induced scheme structure} on the open set $U$, and refer to $(U,\mathcal O_X|_U)$ as an \emph{open subscheme} of $X$.}
\pf
$(U,\mathcal O_X|_U)$ is a locally ringed space (its stalks are stalks of $\mathcal O_X$). Let $x\in U$. Since $X$ is a scheme there is an open $V\ni x$ and an isomorphism of locally ringed spaces $(V,\mathcal O_X|_V)\cong\Spec A$. The set $U\cap V$ is open in $V$, so under this isomorphism it corresponds to an open subset of $\Spec A$ containing the point $\fp$ corresponding to $x$. The sets $D(h)$ form a base of the topology of $\Spec A$ (see the text after Lemma~2.1), so there is $h\in A$ with $\fp\in D(h)\subseteq U\cap V$. By Ex.~2.1, $(D(h),\mathcal O|_{D(h)})\cong\Spec A_h$ is an affine scheme. So $x$ has an open neighbourhood in $U$ which, with the restricted structure sheaf, is an affine scheme. Hence $(U,\mathcal O_X|_U)$ is a scheme.\hfill$\square$

\exercise{2.3}{\emph{Reduced schemes.} A scheme $(X,\mathcal O_X)$ is \emph{reduced} if for every open set $U\subseteq X$, the ring $\mathcal O_X(U)$ has no nilpotent elements.
\begin{enumerate}[label=(\alph*),leftmargin=*]
\item Show that $(X,\mathcal O_X)$ is reduced if and only if for every $P\in X$, the local ring $\mathcal O_{X,P}$ has no nilpotent elements.
\item Let $(X,\mathcal O_X)$ be a scheme. Let $(\mathcal O_X)_{\mathrm{red}}$ be the sheaf associated to the presheaf $U\mapsto\mathcal O_X(U)_{\mathrm{red}}$, where for any ring $A$, we denote by $A_{\mathrm{red}}$ the quotient of $A$ by its ideal of nilpotent elements. Show that $(X,(\mathcal O_X)_{\mathrm{red}})$ is a scheme. We call it the \emph{reduced scheme associated to $X$}, and denote it by $X_{\mathrm{red}}$. Show that there is a morphism of schemes $X_{\mathrm{red}}\to X$, which is a homeomorphism on the underlying topological spaces.
\item Let $f:X\to Y$ be a morphism of schemes, and assume that $X$ is reduced. Show that there is a unique morphism $g:X\to Y_{\mathrm{red}}$ such that $f$ is obtained by composing $g$ with the natural map $Y_{\mathrm{red}}\to Y$.
\end{enumerate}}
\pf
(a) Suppose $X$ is reduced and $\sigma\in\mathcal O_{X,P}$ with $\sigma^n=0$. Represent $\sigma$ by $s\in\mathcal O_X(U)$; then $s^n$ has zero germ at $P$, so $s^n|_{U'}=0$ for a smaller neighbourhood $U'\ni P$. Since $\mathcal O_X(U')$ is reduced, $s|_{U'}=0$, so $\sigma=0$. Conversely suppose all stalks are reduced and $s\in\mathcal O_X(U)$ satisfies $s^n=0$. For each $P\in U$ the germ $s_P$ is nilpotent, hence $0$; a section with all germs zero is zero (sheaf axiom). So $s=0$.

(b) \emph{Stalks.} Let $\mathcal P$ be the presheaf $U\mapsto\mathcal O_X(U)_{\mathrm{red}}$; its stalk (equal to that of its sheafification) is the filtered direct limit $\varinjlim_{U\ni P}\mathcal O_X(U)_{\mathrm{red}}$. Since a filtered direct limit is exact, and since the nilradical of $\varinjlim R_\alpha$ is $\varinjlim\mathrm{nil}(R_\alpha)$ (a nilpotent element of the limit is represented by an element that is nilpotent in some $R_\alpha$), this equals $(\mathcal O_{X,P})_{\mathrm{red}}$. Thus the stalks of $(\mathcal O_X)_{\mathrm{red}}$ are the rings $(\mathcal O_{X,P})_{\mathrm{red}}$, which are local (quotient of a local ring by an ideal inside the maximal ideal) and reduced. So $(X,(\mathcal O_X)_{\mathrm{red}})$ is a locally ringed space, and it is reduced by (a).

\emph{Natural map.} The quotient maps $\mathcal O_X(U)\to\mathcal O_X(U)_{\mathrm{red}}$ give a morphism of sheaves $\varepsilon:\mathcal O_X\to(\mathcal O_X)_{\mathrm{red}}$, which is surjective on stalks with kernel the nilradical; so on stalks it is a local homomorphism ($\varepsilon^{-1}(\fm/\mathrm{nil})=\fm$). Hence $(\mathrm{id}_X,\varepsilon):X_{\mathrm{red}}\to X$ is a morphism of locally ringed spaces, and it is a homeomorphism on spaces (the identity).

\emph{$X_{\mathrm{red}}$ is a scheme.} The construction is compatible with restriction to open sets: $((\mathcal O_X)_{\mathrm{red}})|_U=(\mathcal O_X|_U)_{\mathrm{red}}$. So it suffices to treat $X=\Spec A$. Let $\pi:\Spec A_{\mathrm{red}}\to\Spec A$ be the morphism of Prop.~2.3(b) for the quotient map $A\to A_{\mathrm{red}}$. Since $\Nil(A)$ is contained in every prime, $\pi$ is a bijection, and it is a homeomorphism ($V(\fa)\mapsto V(\fa+\Nil A)$ and the closed sets correspond; or note $\pi(V(\mathfrak b))=V(\varphi^{-1}\mathfrak b)$ as in Ex.~2.18(c)). The sheaf $\pi_*\mathcal O_{\Spec A_{\mathrm{red}}}$ has reduced rings of sections: by (a) $\Spec A_{\mathrm{red}}$ is reduced, because its stalks $(A_{\mathrm{red}})_{\fp}$ are localizations of a reduced ring. Hence $\pi^\#:\mathcal O_{\Spec A}(V)\to\mathcal O_{\Spec A_{\mathrm{red}}}(V)$ kills nilpotents and factors through $\mathcal O_{\Spec A}(V)_{\mathrm{red}}$; by the universal property of sheafification we get $\mathcal O_{\Spec A,\mathrm{red}}\to\pi_*\mathcal O_{\Spec A_{\mathrm{red}}}$. On stalks at $\fp$ this is the map $(A_\fp)_{\mathrm{red}}\to(A_{\mathrm{red}})_\fp$, an isomorphism because localization is exact and commutes with taking the nilradical ($\Nil(A_\fp)=\Nil(A)_\fp$). Hence $(\Spec A,\mathcal O_{\mathrm{red}})\cong\Spec A_{\mathrm{red}}$ is affine, and $X_{\mathrm{red}}$ is a scheme in general.

(c) \emph{Existence.} Let $i:Y_{\mathrm{red}}\to Y$ be the natural map. As $i$ is the identity on spaces, put $g=f$ as a continuous map. For $V\subseteq Y$ open, the ring $\mathcal O_X(f^{-1}V)$ is reduced, so $f^\#:\mathcal O_Y(V)\to\mathcal O_X(f^{-1}V)$ kills nilpotents and factors through $\mathcal O_Y(V)_{\mathrm{red}}$. This gives a morphism from the presheaf $V\mapsto\mathcal O_Y(V)_{\mathrm{red}}$ to the sheaf $f_*\mathcal O_X$, hence a morphism $g^\#:(\mathcal O_Y)_{\mathrm{red}}\to f_*\mathcal O_X$ with $g^\#\circ\varepsilon=f^\#$. On stalks at $x\in X$, $y=f(x)$, $g^\#_x:(\mathcal O_{Y,y})_{\mathrm{red}}\to\mathcal O_{X,x}$ is induced by the local homomorphism $f^\#_x$; since $f^\#_x(\fm_y)\subseteq\fm_x$, $g^\#_x$ maps the maximal ideal $\fm_y/\mathrm{nil}$ into $\fm_x$, so it is local. Thus $g$ is a morphism of locally ringed spaces and $i\circ g=f$ (indeed $(i\circ g)^\#=i_*(g^\#)\circ\varepsilon=f^\#$ with $\varepsilon=i^\#$).

\emph{Uniqueness.} Let $g':X\to Y_{\mathrm{red}}$ satisfy $i\circ g'=f$. Since $i$ is the identity on spaces, $g'=f=g$ as continuous maps, and $g'^\#\circ\varepsilon=f^\#=g^\#\circ\varepsilon$. As $\varepsilon$ is surjective on stalks, $g'^\#$ and $g^\#$ agree on all stalks of $(\mathcal O_Y)_{\mathrm{red}}$ (as maps into the sheaf $f_*\mathcal O_X$ whose sections are determined by germs), hence $g'^\#=g^\#$. So $g'=g$.\hfill$\square$

\exercise{2.4}{Let $A$ be a ring and $(X,\mathcal O_X)$ a scheme. Given a morphism $f:X\to\Spec A$, we have an associated map on sheaves $f^\#:\mathcal O_{\Spec A}\to f_*\mathcal O_X$. Taking global sections we obtain a homomorphism $A\to\Gamma(X,\mathcal O_X)$. Thus there is a natural map
\[\alpha:\Hom_{\mathrm{Sch}}(X,\Spec A)\to\Hom_{\mathrm{Rings}}(A,\Gamma(X,\mathcal O_X)).\]
Show that $\alpha$ is bijective (cf.\ (I, 3.5) for an analogous statement about varieties).}
\pf
Write $B=\Gamma(X,\mathcal O_X)$ and, for $s\in B$ and $x\in X$, $s_x\in\mathcal O_x$ for its germ. Recall $\Gamma(\Spec A,\mathcal O)=A$ (Prop.~2.2(c)).

\emph{Injectivity.} Let $f:X\to\Spec A$ with $\varphi=\alpha(f):A\to B$. For $x\in X$, $\fp=f(x)$, the stalk map $f^\#_x:A_\fp\to\mathcal O_x$ is local, and it composes with $A\to A_\fp$ to the germ map $A\xrightarrow{\varphi}B\to\mathcal O_x$. Locality gives $(f^\#_x)^{-1}(\fm_x)=\fp A_\fp$, so
\[f(x)=\{a\in A:\ \varphi(a)_x\in\fm_x\}.\tag{$*$}\]
Thus $f$ is determined as a map of sets by $\varphi$. Next, $f^\#$ is determined on each basic open set $D(g)$: the map $f^\#(D(g)):A_g\to\mathcal O_X(f^{-1}D(g))$ is a ring homomorphism compatible with the restrictions from $A=\Gamma(\Spec A,\mathcal O)$, so it sends $a/1$ to $\varphi(a)|_{f^{-1}D(g)}$; a ring homomorphism out of $A_g$ is determined by its values on $A$ (the images of the elements $g^{-n}$ are the inverses of the image of $g$). Since the sets $D(g)$ form a base, the morphism of sheaves $f^\#$ is determined (a morphism of sheaves is determined by its values on a base). So $\alpha$ is injective.

\emph{Surjectivity.} Let $\varphi:A\to B$ be a ring homomorphism. Define $f:X\to\Spec A$ by $(*)$: $f(x)=\varphi_x^{-1}(\fm_x)$, where $\varphi_x:A\to\mathcal O_x$ is $a\mapsto\varphi(a)_x$; this is a prime ideal, being the preimage of the prime $\fm_x$. For $g\in A$,
\[f^{-1}(D(g))=\{x:\varphi(g)_x\notin\fm_x\}=\{x:\ \varphi(g)_x\ \text{is a unit}\}=:X_{\varphi(g)}.\]
This is open: if the germ of $s\in B$ at $x$ is a unit, there is a neighbourhood $W$ and $t\in\mathcal O_X(W)$ with $st=1$ near $x$, hence on a smaller neighbourhood where $s|\cdot t=1$, and all germs there are units. So $f$ is continuous. Moreover $\varphi(g)|_{X_{\varphi(g)}}$ is a unit of $\mathcal O_X(X_{\varphi(g)})$: the local inverses exist near every point and agree on overlaps (inverses are unique), so they glue. By the universal property of localization, $A\to\mathcal O_X(X_{\varphi(g)})$, $a\mapsto\varphi(a)|$, extends uniquely to $\theta_g:A_g\to\mathcal O_X(f^{-1}D(g))$. For $D(h)\subseteq D(g)$ the diagrams with the restriction maps commute (uniqueness in the universal property). So the $\theta_g$ define a morphism of sheaves on the base $\{D(g)\}$, which extends uniquely to a morphism of sheaves $f^\#:\mathcal O_{\Spec A}\to f_*\mathcal O_X$. On the stalk at $\fp=f(x)$ it is the map $A_\fp\to\mathcal O_x$ induced from $\varphi_x:A\to\mathcal O_x$ (elements of $A\setminus\fp$ map to units); it is local, since $\varphi_x(\fp)\subseteq\fm_x$ by definition of $f(x)$. Hence $f$ is a morphism of locally ringed spaces, i.e.\ a morphism of schemes, and $\alpha(f)=\theta_1=\varphi$.\hfill$\square$

\exercise{2.5}{Describe $\Spec\Z$, and show that it is a final object for the category of schemes, i.e., each scheme $X$ admits a unique morphism to $\Spec\Z$.}
\pf
\emph{Description.} $\Z$ is a PID, so its prime ideals are $(0)$ and $(p)$, $p$ a prime number; the $(p)$ are maximal. The closed subsets $V(\fa)$ are: $V(0)=\Spec\Z$, and for $\fa=(n)$, $n\ne0$, the finite set of primes dividing $n$ (and $\emptyset$ for $n=\pm1$). So the proper closed subsets are exactly the finite sets of closed points $(p)$; the point $(0)$ is a generic point, $\overline{\{(0)\}}=\Spec\Z$, and it lies in every nonempty open set. The residue field at $(p)$ is $\mathbf F_p$, at $(0)$ it is $\Q$; the local rings are $\Z_{(p)}$ and $\Q$. Global sections: $\Gamma(\Spec\Z,\mathcal O)=\Z$.

\emph{Final object.} By Ex.~2.4, for every scheme $X$ the set $\Hom(X,\Spec\Z)$ is in bijection with $\Hom_{\mathrm{Rings}}(\Z,\Gamma(X,\mathcal O_X))$. Since rings have $1$ and homomorphisms preserve $1$, there is exactly one ring homomorphism $\Z\to R$ for every ring $R$ ($n\mapsto n\cdot1_R$). Hence there is exactly one morphism $X\to\Spec\Z$.\hfill$\square$

\exercise{2.6}{Describe the spectrum of the zero ring, and show that it is an initial object for the category of schemes. (According to our conventions, all ring homomorphisms must take $1$ to $1$. Since $0=1$ in the zero ring, we see that each ring $R$ admits a unique homomorphism to the zero ring, but that there is no homomorphism from the zero ring to $R$ unless $0=1$ in $R$.)}
\pf
The zero ring $0$ has only one ideal, namely $0=(1)$ itself, which is not a proper ideal, so it has no prime ideals: $\Spec 0=\emptyset$, with the zero sheaf (the value of a sheaf on $\emptyset$ is the zero ring). It is an affine scheme (the empty space is covered by no points, so the scheme condition is vacuous).

Let $X$ be any scheme. A morphism $\emptyset\to X$ consists of the unique (empty) continuous map and a morphism of sheaves $\mathcal O_X\to f_*\mathcal O_\emptyset$. For every open $V\subseteq X$, $(f_*\mathcal O_\emptyset)(V)=\mathcal O_\emptyset(\emptyset)=0$, and there is exactly one ring homomorphism from any ring to the zero ring; these are compatible with restrictions. So there is exactly one morphism of sheaves, and the locality condition on stalks is vacuous (there are no points). Hence $\Hom(\Spec0,X)$ has exactly one element for every scheme $X$: $\Spec 0$ is an initial object.

(Consistency check with Ex.~2.4: $\Hom(X,\Spec 0)=\Hom_{\mathrm{Rings}}(0,\Gamma(X,\mathcal O_X))$ is nonempty iff $\Gamma(X,\mathcal O_X)=0$. If $X\ne\emptyset$ this fails, since $1=0$ in $\Gamma(X,\mathcal O_X)$ would give $1=0$ in a local ring $\mathcal O_x$, which is nonzero. So only $X=\emptyset$ maps to $\Spec 0$, consistently with the above.)\hfill$\square$

\exercise{2.7}{Let $X$ be a scheme. For any point $x\in X$, let $\mathcal O_x$ be the local ring at $x$, and $\fm_x$ its maximal ideal. We define the \emph{residue field} of $x$ on $X$ to be the field $k(x)=\mathcal O_x/\fm_x$. Now let $K$ be any field. Show that to give a morphism of $\Spec K$ to $X$ is equivalent to giving a point $x\in X$ and an inclusion map $k(x)\to K$.}
\pf
I first prove a lemma that will also be used in Ex.~2.8.

\textbf{Lemma A.} Let $X$ be a scheme and $R$ a ring having exactly one prime ideal $\fn$ (e.g.\ a field, or $k[\varepsilon]$). Let $*$ be the unique point of $\Spec R$. Then $f\mapsto(f(*),f^\#_*)$ is a bijection from $\Hom(\Spec R,X)$ onto the set of pairs $(x,\theta)$ with $x\in X$ and $\theta:\mathcal O_x\to R$ a local homomorphism.

\emph{Proof.} The stalk map $f^\#_*:\mathcal O_{X,f(*)}\to\mathcal O_{\Spec R,*}=R_\fn=R$ is local by definition of a morphism of schemes, so the map is well defined.

\emph{Injective.} Let $f,f'$ give the same $(x,\theta)$. Choose an affine open $U=\Spec A\ni x$, $x\leftrightarrow\fp$. Both $f,f'$ have image $\{x\}\subseteq U$, so they factor through $U$ as morphisms of locally ringed spaces $f_U,f'_U:\Spec R\to\Spec A$. By Prop.~2.3(c) they come from ring homomorphisms $\psi,\psi':A\to R$. The stalk map of $f_U$ at $*$ is $A_\fp\to R$ induced by $\psi$, i.e.\ $\psi=\theta\circ(A\to A_\fp)$; likewise for $\psi'$. So $\psi=\psi'$ and $f=f'$.

\emph{Surjective.} Given $(x,\theta)$ and $U=\Spec A\ni x$ as above, put $\psi:A\to A_\fp=\mathcal O_x\xrightarrow{\theta}R$ and let $f$ be $\Spec R\to\Spec A=U\subseteq X$ induced by $\psi$ (Prop.~2.3(b)). Then $f(*)=\psi^{-1}(\fn)=(A\to A_\fp)^{-1}(\theta^{-1}\fn)=(A\to A_\fp)^{-1}(\fp A_\fp)=\fp$ since $\theta$ is local. The stalk map $A_\fp\to R$ of $f$ is the unique extension of $\psi$ to $A_\fp$, namely $\theta$. \hfill$\square$

Now take $R=K$. A homomorphism $\theta:\mathcal O_x\to K$ is local iff $\theta^{-1}(0)=\fm_x$, i.e.\ iff $\ker\theta=\fm_x$, i.e.\ iff $\theta$ is the composition of $\mathcal O_x\to k(x)$ with a ring homomorphism $\iota:k(x)\to K$ (which is automatically injective, $k(x)$ being a field). This identifies the pairs $(x,\theta)$ with pairs $(x,\iota:k(x)\hookrightarrow K)$. By Lemma~A, morphisms $\Spec K\to X$ correspond bijectively to a point $x\in X$ together with an inclusion $k(x)\to K$.\hfill$\square$

\exercise{2.8}{Let $X$ be a scheme. For any point $x\in X$, we define the \emph{Zariski tangent space} $T_x$ to $X$ at $x$ to be the dual of the $k(x)$-vector space $\fm_x/\fm_x^2$. Now assume that $X$ is a scheme over a field $k$, and let $k[\varepsilon]$ be the ring of dual numbers over $k$, i.e., $k[\varepsilon]=k[t]/(t^2)$. Show that to give a $k$-morphism of $\Spec k[\varepsilon]$ to $X$ is equivalent to giving a point $x\in X$, rational over $k$ (i.e., such that $k(x)=k$), and an element of $T_x$.}
\pf
The ring $k[\varepsilon]$ has a unique prime ideal $(\varepsilon)$ (it is local and $\varepsilon$ is nilpotent), with residue field $k$. So Lemma~A (Ex.~2.7) applies: morphisms $\Spec k[\varepsilon]\to X$ correspond to pairs $(x,\theta)$, $\theta:\mathcal O_x\to k[\varepsilon]$ a local homomorphism.

\emph{The $k$-structure.} Let $\pi:X\to\Spec k$ be the structure morphism; by Ex.~2.4 it corresponds to a ring map $\kappa:k\to\Gamma(X,\mathcal O_X)$, which makes each $\mathcal O_x$ a $k$-algebra. The morphism $f$ corresponding to $(x,\theta)$ is a $k$-morphism iff $\pi\circ f$ is the structure morphism of $\Spec k[\varepsilon]$; by Ex.~2.4 this says that $k\xrightarrow{\kappa}\Gamma(X,\mathcal O_X)\xrightarrow{f^\#}k[\varepsilon]$ is the inclusion. This composite factors as $k\to\mathcal O_x\xrightarrow{\theta}k[\varepsilon]$, and conversely it is determined by $\theta$ (we are only dealing with one point). So $f$ is a $k$-morphism iff $\theta$ is a homomorphism of $k$-algebras.

\emph{From $\theta$ to $(x,\delta)$.} Let $\theta:\mathcal O_x\to k[\varepsilon]$ be a local $k$-algebra homomorphism. Then $\theta(\fm_x)\subseteq(\varepsilon)$, and $\theta$ induces a $k$-algebra homomorphism of fields $\bar\theta:k(x)\to k[\varepsilon]/(\varepsilon)=k$. Both $k\to k(x)$ (a field map, hence injective) and $\bar\theta$ (injective) are $k$-algebra maps, and $\bar\theta\circ(k\to k(x))=\mathrm{id}_k$. For $y\in k(x)$ choose $c\in k$ with $\bar\theta(y)=c=\bar\theta(c\cdot1)$; injectivity of $\bar\theta$ gives $y=c\cdot 1$. So $k\to k(x)$ is onto: $k(x)=k$, i.e.\ $x$ is rational over $k$. Consequently $\mathcal O_x=k\cdot1\oplus\fm_x$ as $k$-vector spaces. For $u\in\fm_x$ write $\theta(u)=\delta(u)\varepsilon$ with $\delta(u)\in k$; $\delta:\fm_x\to k$ is $k$-linear, and for $u,v\in\fm_x$, $\delta(uv)\varepsilon=\theta(u)\theta(v)=\delta(u)\delta(v)\varepsilon^2=0$. So $\delta$ vanishes on $\fm_x^2$ and defines $\delta\in\Hom_k(\fm_x/\fm_x^2,k)=T_x$ (recall $k(x)=k$). Moreover $\theta(a+u)=a+\delta(u)\varepsilon$ for $a\in k$, $u\in\fm_x$, so $\theta$ is determined by $\delta$.

\emph{From $(x,\delta)$ to $\theta$.} Conversely let $x$ be rational over $k$ and $\delta\in T_x$, viewed as a $k$-linear map $\fm_x\to k$ vanishing on $\fm_x^2$. Define $\theta(a+u)=a+\delta(u)\varepsilon$ ($a\in k$, $u\in\fm_x$). It is $k$-linear, $\theta(1)=1$, and
\[\theta((a+u)(b+v))=ab+\bigl(b\delta(u)+a\delta(v)+\delta(uv)\bigr)\varepsilon=ab+(b\delta(u)+a\delta(v))\varepsilon=\theta(a+u)\theta(b+v),\]
since $\delta(uv)=0$. It is local: $\theta^{-1}((\varepsilon))=\fm_x$.

These two constructions are inverse to each other, so $k$-morphisms $\Spec k[\varepsilon]\to X$ correspond bijectively to pairs ($x\in X$ with $k(x)=k$, element of $T_x$).\hfill$\square$

\exercise{2.9}{If $X$ is a topological space, and $Z$ an irreducible closed subset of $X$, a \emph{generic point} for $Z$ is a point $\xi$ such that $Z=\overline{\{\xi\}}$. If $X$ is a scheme, show that every (nonempty) irreducible closed subset has a unique generic point.}
\pf
\emph{Affine case.} Let $X=\Spec A$ and $Z\subseteq X$ nonempty, closed, irreducible. Write $Z=V(\fa)$ and $I=\sqrt{\fa}$ (so $V(I)=Z$, by Lemma~2.1(c)). $I$ is a proper ideal since $Z\ne\emptyset$. It is prime: if $ab\in I$ then $Z=V(I)\subseteq V(ab)=V(a)\cup V(b)$, so $Z=(Z\cap V(a))\cup(Z\cap V(b))$ is a union of two closed subsets; by irreducibility $Z\subseteq V(a)$ or $Z\subseteq V(b)$, i.e.\ $a\in\sqrt I=I$ or $b\in I$. For a prime $\fp$ one has $\overline{\{\fp\}}=V(\fp)$ (the smallest closed set $V(\fb)$ containing $\fp$ has $\fb\subseteq\fp$, so $V(\fb)\supseteq V(\fp)$). Hence $Z=\overline{\{I\}}$: a generic point exists. If $\overline{\{\fp\}}=\overline{\{\fq\}}$ then $V(\fp)=V(\fq)$, so $\fp\supseteq\fq\supseteq\fp$ and $\fp=\fq$: it is unique.

\emph{General case.} Let $Z\subseteq X$ be nonempty, closed, irreducible, and pick $z\in Z$ and an affine open $U\ni z$. Then $Z\cap U$ is a nonempty open subset of the irreducible space $Z$, hence irreducible and dense in $Z$; and it is closed in $U$. By the affine case there is $\xi\in U$ with $Z\cap U=\overline{\{\xi\}}\cap U$ (closure in $U$ is closure in $X$ intersected with $U$). Then $\xi\in Z$, so $\overline{\{\xi\}}\subseteq Z$ ($Z$ closed); and $Z=\overline{Z\cap U}\subseteq\overline{\{\xi\}}$. So $Z=\overline{\{\xi\}}$.

\emph{Uniqueness.} If also $Z=\overline{\{\xi'\}}$, then the open set $U\cap Z$ is a nonempty open subset of $\overline{\{\xi'\}}$, so it contains $\xi'$ (an open set meeting the closure of $\{\xi'\}$ contains $\xi'$). Hence $\xi'\in U$ and $\overline{\{\xi'\}}\cap U=Z\cap U=\overline{\{\xi\}}\cap U$, i.e.\ $\xi,\xi'$ are generic points of the same irreducible closed subset of $U$. By the affine case, $\xi=\xi'$.\hfill$\square$

\exercise{2.10}{Describe $\Spec\R[x]$. How does its topological space compare to the set $\R$? To $\C$?}
\pf
$\R[x]$ is a PID, so its primes are $(0)$ and the ideals $(p)$ with $p$ irreducible monic. Over $\R$ the irreducible monic polynomials are $x-a$ ($a\in\R$) and $x^2+bx+c$ with $b^2-4c<0$ (the monic quadratics with a pair of non-real conjugate roots $z,\bar z\in\C$). Hence $\Spec\R[x]$ consists of
\begin{itemize}
\item the generic point $\eta=(0)$, with residue field $\R(x)$;
\item the \emph{real points} $(x-a)$, $a\in\R$, with residue field $\R[x]/(x-a)=\R$;
\item the \emph{complex points} $(x^2+bx+c)=(x-z)(x-\bar z)$ with $z\in\C\setminus\R$, residue field $\R[x]/(x^2+bx+c)\cong\C$, one for each conjugate pair $\{z,\bar z\}$.
\end{itemize}
All the nonzero primes are maximal, i.e.\ closed points. Thus the set of closed points is naturally in bijection with $\C/(z\sim\bar z)$, which is the closed upper half plane $\{z:\operatorname{Im}z\ge0\}$; it contains $\R$ as the real points, and it is the quotient of $\C=\Max\C[x]$ by complex conjugation (the map $\Spec\C[x]\to\Spec\R[x]$ is $2:1$ off the real line).

Topology: every nonzero ideal is $(p)$ with $p$ a product of finitely many irreducibles, so $V((p))$ is a finite set of closed points. The closed sets are therefore $\emptyset$, the finite sets of closed points, and the whole space. The open sets are the cofinite sets of closed points together with $\eta$, plus $\emptyset$. So: the closed points with the induced topology form an infinite set with the cofinite topology (just as for $\R$ or $\C$ with the Zariski topology on $\mathbf A^1$), which is not Hausdorff and is very different from the usual topology of $\R$ or $\C$. Comparing, the set of closed points is neither $\R$ nor $\C$ but the quotient $\C/\text{conj}$, which contains $\R$ and is ``half of $\C$''; and the whole space has in addition the generic point $\eta$ whose closure is everything. Compared with $\mathbf A^1_{\C}=\Spec\C[x]$, whose closed points are $\C$, the space $\Spec\R[x]$ is obtained by identifying conjugate points (and the two have the same cofinite structure).\hfill$\square$

\exercise{2.11}{Let $k=\mathbf F_p$ be the finite field with $p$ elements. Describe $\Spec k[x]$. What are the residue fields of its points? How many points are there with a given residue field?}
\pf
$k[x]$ is a PID, so $\Spec k[x]$ consists of the generic point $\eta=(0)$ and the closed points $(f)$, $f\in\mathbf F_p[x]$ monic irreducible. The topology is again: closed sets are the finite sets of closed points and the whole space. The residue fields are:
\begin{itemize}
\item at $\eta$: $\mathbf F_p(x)$, an infinite field;
\item at $(f)$, $\deg f=d$: $\mathbf F_p[x]/(f)\cong\mathbf F_{p^d}$, the field with $p^d$ elements.
\end{itemize}
\emph{Counting.} Let $N_d$ be the number of monic irreducible polynomials of degree $d$ in $\mathbf F_p[x]$. Standard fact: $x^{p^n}-x$ (which is separable, its derivative being $-1$) is the product of all monic irreducible polynomials whose degree divides $n$ (a monic irreducible $f$ of degree $d$ divides $x^{p^n}-x$ iff the class $\bar x$ of $x$ in $\mathbf F_p[x]/(f)=\mathbf F_{p^d}$ satisfies $\bar x^{p^n}=\bar x$, iff the $n$-th power of the Frobenius $F:y\mapsto y^p$ fixes the generator $\bar x$ of $\mathbf F_{p^d}$ over $\mathbf F_p$, iff $F^n=\mathrm{id}$ on $\mathbf F_{p^d}$, iff $d\mid n$, since $F$ has order exactly $d$ on $\mathbf F_{p^d}$; and $x^{p^n}-x$, being separable, is squarefree). Comparing degrees, $\sum_{d\mid n}dN_d=p^n$, and Moebius inversion gives
\[N_d=\frac1d\sum_{e\mid d}\mu(d/e)\,p^e.\]
Hence the number of points with residue field $\mathbf F_{p^d}$ is $N_d$ (e.g.\ $N_1=p$, $N_2=(p^2-p)/2$, $N_3=(p^3-p)/3$, $N_4=(p^4-p^2)/4$); these are $\ge1$ for every $d$ (since $d N_d\ge p^d-\sum_{e<d}p^e>0$). The unique point with residue field $\mathbf F_p(x)$ is $\eta$. For a field $L$ not isomorphic to any $\mathbf F_{p^d}$ or to $\mathbf F_p(x)$ there is no point with residue field $L$. Equivalently the points with residue field $\mathbf F_{p^d}$ correspond to Frobenius orbits of size exactly $d$ on $\overline{\mathbf F}_p$.\hfill$\square$

\exercise{2.12}{\emph{Gluing Lemma.} Generalize the gluing procedure described in the text (2.3.5) as follows. Let $\{X_i\}$ be a family of schemes (possibly infinite). For each $i\ne j$, suppose given an open subset $U_{ij}\subseteq X_i$, and let it have the induced scheme structure (Ex.~2.2). Suppose also given for each $i\ne j$ an isomorphism of schemes $\varphi_{ij}:U_{ij}\to U_{ji}$ such that (1) for each $i,j$, $\varphi_{ji}=\varphi_{ij}^{-1}$, and (2) for each $i,j,k$, $\varphi_{ij}(U_{ij}\cap U_{ik})=U_{ji}\cap U_{jk}$ and $\varphi_{ik}=\varphi_{jk}\circ\varphi_{ij}$ on $U_{ij}\cap U_{ik}$. Then there is a scheme $X$, together with morphisms $\psi_i:X_i\to X$ for each $i$, such that (1) $\psi_i$ is an isomorphism of $X_i$ onto an open subscheme of $X$, (2) the $\psi_i(X_i)$ cover $X$, (3) $\psi_i(U_{ij})=\psi_i(X_i)\cap\psi_j(X_j)$ and (4) $\psi_i=\psi_j\circ\varphi_{ij}$ on $U_{ij}$. We say that $X$ is obtained by \emph{glueing} the schemes $X_i$ along the isomorphisms $\varphi_{ij}$. An interesting special case is when the family $X_i$ is arbitrary, but the $U_{ij}$ are all empty. Then the scheme $X$ is called the \emph{disjoint union} of the $X_i$, and is denoted $\coprod X_i$.}
\pf
Put $U_{ii}=X_i$ and $\varphi_{ii}=\mathrm{id}$; conditions (1), (2) then hold for all $i,j,k$ (including equal indices).

\emph{The set and topology.} Let $X$ be the quotient of $\bigsqcup X_i$ by the relation $x\sim y$ ($x\in X_i$, $y\in X_j$) iff $x\in U_{ij}$ and $y=\varphi_{ij}(x)$. It is an equivalence relation: reflexivity from $\varphi_{ii}=\mathrm{id}$; symmetry from (1). Transitivity: if $y=\varphi_{ij}(x)$ and $z=\varphi_{jk}(y)$ with $y\in U_{ji}\cap U_{jk}=\varphi_{ij}(U_{ij}\cap U_{ik})$, then $x\in U_{ij}\cap U_{ik}$ (injectivity of $\varphi_{ij}$) and by (2) $z=\varphi_{jk}\varphi_{ij}(x)=\varphi_{ik}(x)$. Let $\psi_i:X_i\to X$ be the canonical maps; give $X$ the quotient topology: $W\subseteq X$ is open iff all $\psi_i^{-1}(W)$ are open. Two points of the same $X_i$ are equivalent only if equal ($\varphi_{ii}=\mathrm{id}$), so each $\psi_i$ is injective. Moreover $\psi_j^{-1}(\psi_i(X_i))=U_{ji}$, which is open in $X_j$, so $\psi_i(X_i)$ is open; and for $V\subseteq X_i$ open, $\psi_j^{-1}(\psi_i(V))=\varphi_{ij}(V\cap U_{ij})$ is open in $U_{ji}$, hence in $X_j$, so $\psi_i(V)$ is open. Thus $\psi_i$ is a homeomorphism of $X_i$ onto the open subset $X_i':=\psi_i(X_i)$ of $X$; the $X_i'$ cover $X$ (obvious) and $X'_i\cap X'_j=\psi_i(U_{ij})=\psi_j(U_{ji})$; and $\psi_i=\psi_j\circ\varphi_{ij}$ on $U_{ij}$ as maps of sets (this is the definition of $\sim$).

\emph{The sheaf.} For $W\subseteq X$ open put $W_i=\psi_i^{-1}(W)$, open in $X_i$. Since $\psi_j\varphi_{ij}=\psi_i$ on $U_{ij}$ and $\varphi_{ij}:U_{ij}\to U_{ji}$ is bijective, $\varphi_{ij}(W_i\cap U_{ij})=W_j\cap U_{ji}$, so the isomorphism of schemes $\varphi_{ij}$ gives a ring isomorphism $\varphi_{ij}^*:\mathcal O_{X_j}(W_j\cap U_{ji})\to\mathcal O_{X_i}(W_i\cap U_{ij})$. Define
\[\mathcal O_X(W)=\bigl\{(s_i)_i\in\textstyle\prod_i\mathcal O_{X_i}(W_i):\ s_i|_{W_i\cap U_{ij}}=\varphi_{ij}^*\bigl(s_j|_{W_j\cap U_{ji}}\bigr)\ \text{for all }i,j\bigr\},\]
with componentwise restriction. This is a subring of the product, and it is a sheaf: it is the kernel of a morphism from the sheaf $\prod_i(\psi_i)_*\mathcal O_{X_i}$ to a product of sheaves, namely $(s_i)\mapsto(s_i|-\varphi_{ij}^*(s_j|))_{i,j}$. Let $p_i:\mathcal O_X\to(\psi_i)_*\mathcal O_{X_i}$ be the $i$-th projection.

\emph{Claim: for $W\subseteq X_i'$ the map $p_i:\mathcal O_X(W)\to\mathcal O_{X_i}(W_i)$ is bijective.} Note $W_j\subseteq\psi_j^{-1}(X_i')=U_{ji}$ for all $j$, so $W_j\cap U_{ji}=W_j$. \emph{Injective:} the $(i,j)$ condition says $s_i|_{W_i\cap U_{ij}}=\varphi_{ij}^*(s_j)$, which determines $s_j$ from $s_i$ since $\varphi_{ij}^*$ is injective. \emph{Surjective:} given $s_i$ put $s_j=(\varphi_{ij}^*)^{-1}(s_i|_{W_i\cap U_{ij}})\in\mathcal O_{X_j}(W_j)$ ($s_i$ itself for $j=i$). The $(i,j)$ conditions hold by construction and the $(j,i)$ conditions by $\varphi_{ji}=\varphi_{ij}^{-1}$. For $(j,k)$ we must show $s_j|_{W_j\cap U_{jk}}=\varphi_{jk}^*(s_k|_{W_k\cap U_{kj}})$. Put $T=W_i\cap U_{ij}\cap U_{ik}$; by (2), $\varphi_{ij}$ maps $T$ isomorphically onto $W_j\cap U_{jk}$ (recall $W_j\cap U_{ji}=W_j$). Applying the injective map $\varphi_{ij}^*:\mathcal O_{X_j}(W_j\cap U_{jk})\to\mathcal O_{X_i}(T)$ to both sides, the left side gives $s_i|_T$ and the right side gives $(\varphi_{jk}\varphi_{ij})^*(s_k|)=\varphi_{ik}^*(s_k|_{W_k\cap U_{ki}})|_T=s_i|_T$, using $\varphi_{ik}=\varphi_{jk}\varphi_{ij}$ on $U_{ij}\cap U_{ik}$ from (2) and the definition of $s_k$. So the two sides agree. This proves the claim.

\emph{$X$ is a scheme.} By the claim, $(\psi_i,p_i)$ is an isomorphism of ringed spaces $(X_i,\mathcal O_{X_i})\to(X_i',\mathcal O_X|_{X_i'})$ (we showed $\psi_i$ is a homeomorphism onto $X_i'$). Hence the stalks of $\mathcal O_X$ are local rings, $(X,\mathcal O_X)$ is a locally ringed space, and it is covered by the open sets $X_i'$, each isomorphic to the scheme $X_i$; so each point has an affine open neighbourhood (Ex.~2.2 applied inside $X_i'$). So $X$ is a scheme, and $\psi_i=(\psi_i,p_i):X_i\to X$ is an isomorphism onto the open subscheme $X_i'$. The $X_i'$ cover $X$ and $\psi_i(U_{ij})=X_i'\cap X_j'$ (shown above). Finally (4) holds as morphisms: the sheaf map of $\psi_j\circ\varphi_{ij}|_{U_{ij}}$ sends $s=(s_k)\in\mathcal O_X(W)$ to $\varphi_{ij}^*(s_j|_{W_j\cap U_{ji}})$, while that of $\psi_i|_{U_{ij}}$ sends it to $s_i|_{W_i\cap U_{ij}}$, and these agree by the defining condition of $\mathcal O_X$. This proves properties (1)--(4).

\emph{Disjoint union.} If all $U_{ij}=\emptyset$ ($i\neq j$) the conditions (1), (2) are vacuous (all $\varphi_{ij}$ are the empty maps), the relation $\sim$ is equality, and $X=\coprod X_i$ is the topological disjoint union with $\mathcal O_X|_{X_i}=\mathcal O_{X_i}$: a scheme in which every $X_i$ is open and $X_i\cap X_j=\emptyset$.\hfill$\square$

\exercise{2.13}{A topological space is \emph{quasi-compact} if every open cover has a finite subcover.
\begin{enumerate}[label=(\alph*),leftmargin=*]
\item Show that a topological space is noetherian (Ch.~I, Section~1) if and only if every open subset is quasi-compact.
\item If $X$ is an affine scheme, show that $\operatorname{sp}(X)$ is quasi-compact, but not in general noetherian. We say that a scheme $X$ is \emph{quasi-compact} if $\operatorname{sp}(X)$ is.
\item If $A$ is a noetherian ring, show that $\operatorname{sp}(\Spec A)$ is a noetherian topological space.
\item Give an example of a non-noetherian ring $A$ such that $\operatorname{sp}(\Spec A)$ is noetherian.
\end{enumerate}}
\pf
Recall (Ch.~I, Section~1) that a space is noetherian if it satisfies the descending chain condition for closed subsets, equivalently every nonempty family of closed subsets has a minimal element.

(a) ($\Rightarrow$) Let $X$ be noetherian, $U\subseteq X$ open, and $\{V_\alpha\}$ an open cover of $U$ by open subsets of $X$. For finite subsets $F$ of the index set let $Z_F=X\setminus\bigcup_{\alpha\in F}V_\alpha$. This is a nonempty family of closed sets, so it has a minimal element $Z_F$. If $U\not\subseteq\bigcup_{\alpha\in F}V_\alpha$, pick $x\in U\setminus\bigcup_FV_\alpha$ and $\beta$ with $x\in V_\beta$; then $Z_{F\cup\{\beta\}}\subsetneq Z_F$ (it misses $x$), contradicting minimality. So $U\subseteq\bigcup_{\alpha\in F}V_\alpha$ and $U$ is quasi-compact.
($\Leftarrow$) Let $Z_1\supseteq Z_2\supseteq\cdots$ be closed. The open set $U=\bigcup_n(X\setminus Z_n)$ is the union of the increasing sequence $X\setminus Z_n$; it is quasi-compact, so $U=X\setminus Z_N$ for some $N$, i.e.\ $Z_n=Z_N$ for all $n\ge N$.

(b) Let $X=\Spec A$ and $X=\bigcup_\alpha U_\alpha$ an open cover. Each $U_\alpha$ is a union of basic sets $D(f)$, so it suffices to show that a cover $X=\bigcup_{i\in I}D(f_i)$ by basic opens has a finite subcover. Then $V(\{f_i\})=\bigcap V(f_i)=\emptyset$, so the ideal generated by the $f_i$ is $A$ (otherwise it lies in a maximal, hence prime, ideal), whence $1=\sum_{i\in F}a_if_i$ for a finite $F$, so $V(f_i:i\in F)=\emptyset$, i.e.\ $X=\bigcup_{i\in F}D(f_i)$. So $\operatorname{sp}(X)$ is quasi-compact.

\emph{Not noetherian in general.} Let $A=k[x_1,x_2,\dots]$ (infinitely many variables), $X=\Spec A$, and $U=X\setminus V(x_1,x_2,\dots)=\bigcup_iD(x_i)$. If $U$ were covered by finitely many $D(x_i)$, $i\le n$, consider $\fp=(x_1,\dots,x_n)$, prime because $A/\fp\cong k[x_{n+1},\dots]$ is a domain. Then $\fp\notin D(x_i)$ for $i\le n$, but $\fp\in U$ because $\fp\not\supseteq(x_1,x_2,\dots)$ ($x_{n+1}\notin\fp$). So $U$ is an open subset which is not quasi-compact, and by (a) $X$ is not noetherian.

(c) Let $V(\fa_1)\supseteq V(\fa_2)\supseteq\cdots$ be a descending chain of closed subsets of $\Spec A$. By Lemma~2.1(c), $V(\fa)\subseteq V(\fb)$ iff $\sqrt\fa\supseteq\sqrt\fb$, so $\sqrt{\fa_1}\subseteq\sqrt{\fa_2}\subseteq\cdots$ is an ascending chain of ideals; as $A$ is noetherian it stabilizes, and then so does $V(\fa_n)=V(\sqrt{\fa_n})$.

(d) Let $A=k[x_1,x_2,\dots]/(x_1^2,x_2^2,x_3^2,\dots)$. Every element of $A$ is a polynomial without constant term plus a constant, and every non-constant-term element is nilpotent (a finite sum of nilpotents is nilpotent). So the ideal $\fm=(x_1,x_2,\dots)$ is the nilradical, and $A/\fm=k$; hence $\fm$ is the unique prime of $A$ and $\Spec A$ is a single point, a noetherian space. But $A$ is not noetherian: $(x_1)\subsetneq(x_1,x_2)\subsetneq(x_1,x_2,x_3)\subsetneq\cdots$ is an infinite strictly ascending chain (setting $x_1=\dots=x_n=0$ gives the ring $k[x_{n+1},\dots]/(x_{j}^2)$ in which $x_{n+1}\ne0$).\hfill$\square$

\exercise{2.14}{\begin{enumerate}[label=(\alph*),leftmargin=*]
\item Let $S$ be a graded ring. Show that $\operatorname{Proj}S=\emptyset$ if and only if every element of $S_+$ is nilpotent.
\item Let $\varphi:S\to T$ be a homomorphism of graded rings (preserving degrees). Let $U=\{\fp\in\operatorname{Proj}T\mid\fp\not\supseteq\varphi(S_+)\}$. Show that $U$ is an open subset of $\operatorname{Proj}T$, and show that $\varphi$ determines a natural morphism $f:U\to\operatorname{Proj}S$.
\item The morphism $f$ can be an isomorphism even when $\varphi$ is not. For example, suppose $\varphi_d:S_d\to T_d$ is an isomorphism for all $d\ge d_0$, where $d_0$ is an integer. Then show that $U=\operatorname{Proj}T$ and the morphism $f:\operatorname{Proj}T\to\operatorname{Proj}S$ is an isomorphism.
\item Let $V$ be a projective variety with homogeneous coordinate ring $S(V)$ (I, Section~2). Show that $t(V)\cong\operatorname{Proj}S(V)$.
\end{enumerate}}
\pf
Recall (Lemma~2.4 and its proof) that a homogeneous ideal $\fp$ of $S$ is prime iff for homogeneous $a,b$, $ab\in\fp$ implies $a\in\fp$ or $b\in\fp$; $\operatorname{Proj}S$ is the set of homogeneous primes not containing $S_+$.

(a) ($\Leftarrow$) If every element of $S_+$ is nilpotent and $\fp$ is a homogeneous prime, then $S_+\subseteq\Nil(S)\subseteq\fp$, so $\fp\notin\operatorname{Proj}S$. Hence $\operatorname{Proj}S=\emptyset$.
($\Rightarrow$) Suppose some element of $S_+$ is not nilpotent. Then (since a sum of nilpotents is nilpotent) some homogeneous $g\in S_d$, $d>0$, is not nilpotent. Let $M=\{1,g,g^2,\dots\}$; $0\notin M$. By Zorn's lemma the set of homogeneous ideals disjoint from $M$ (nonempty: it contains $0$) has a maximal element $\fq$. I claim $\fq$ is prime. Let $a,b$ be homogeneous with $a,b\notin\fq$. By maximality $\fq+(a)$ and $\fq+(b)$ meet $M$; taking the homogeneous components of degrees $md$, $nd$ we get $g^m=q+ar$, $g^n=q'+bs$ with $q,q'\in\fq$ and $r,s$ homogeneous. Multiplying, $g^{m+n}\in\fq+(ab)$; so if $ab\in\fq$ then $g^{m+n}\in\fq$, a contradiction. So $ab\notin\fq$, and $\fq$ is prime. Since $g\in S_+\setminus\fq$, $\fq\not\supseteq S_+$, so $\fq\in\operatorname{Proj}S$, which is nonempty.

(b) Let $\fa\subseteq T$ be the ideal generated by $\varphi(S_+)$; it is homogeneous, and $U=\operatorname{Proj}T\setminus V(\fa)$ is open. For $\fp\in U$, $\varphi^{-1}(\fp)$ is a homogeneous prime of $S$ ($\varphi$ preserves degrees), and it does not contain $S_+$ since $\fp\not\supseteq\varphi(S_+)$. Put $f(\fp)=\varphi^{-1}(\fp)$. If $\fb\subseteq S$ is homogeneous, then $f^{-1}(V(\fb))=U\cap V(\varphi(\fb)T)$, so $f$ is continuous.

\emph{Charts.} Let $g\in S_+$ be homogeneous. Then $f^{-1}(D_+(g))=D_+(\varphi(g))\cap U=D_+(\varphi(g))$ (a prime not containing $\varphi(g)$, $g\in S_+$, certainly does not contain $\varphi(S_+)$). These sets cover $U$: if $\fp\in U$ then $\fp\not\supseteq\varphi(S_+)$, and since $\fp$ is homogeneous and $S_+$ is generated by homogeneous elements, there is a homogeneous $g\in S_+$ with $\varphi(g)\notin\fp$. The ring homomorphism $\varphi$ extends to $S_g\to T_{\varphi(g)}$ and restricts to $\varphi_{(g)}:S_{(g)}\to T_{(\varphi(g))}$. By Prop.~2.5(b), $D_+(\varphi g)\cong\Spec T_{(\varphi g)}$ and $D_+(g)\cong\Spec S_{(g)}$, so $\varphi_{(g)}$ induces (Prop.~2.3(b)) a morphism of schemes $f_g:D_+(\varphi g)\to D_+(g)\subseteq\operatorname{Proj}S$. On points it is $\fp\mapsto\varphi^{-1}(\fp)$ (the isomorphism $D_+(g)\cong\Spec S_{(g)}$ sends $\fp$ to $\fp S_g\cap S_{(g)}$ and $s/g^n\in\varphi_{(g)}^{-1}(\fp T_{\varphi g}\cap T_{(\varphi g)})$ iff $\varphi(s)\in\fp$ iff $s\in\varphi^{-1}\fp$). For two homogeneous $g,g'\in S_+$ of degrees $d,d'$, put $h=g'^d/g^{d'}\in S_{(g)}$; then $D_+(gg')=D_+(g)\cap D_+(g')$ corresponds to $D(h)\subseteq\Spec S_{(g)}$ and $S_{(gg')}=(S_{(g)})_h$ (this is part of the proof of Prop.~2.5(b)), and similarly for $T$ with $\varphi(h)=\varphi(g')^d/\varphi(g)^{d'}$. By Ex.~2.1 (naturality of $D(h)\cong\Spec(S_{(g)})_h$), the restriction of $f_g$ to $D_+(\varphi g)\cap D_+(\varphi g')=D_+(\varphi(gg'))$ is $\Spec$ of the localization $S_{(gg')}\to T_{(\varphi(gg'))}$ of $\varphi_{(g)}$ at $h$, which is the map $s/(gg')^n\mapsto\varphi(s)/\varphi(gg')^n$ induced by $\varphi$. The same description holds for $f_{g'}$ (the roles of $g,g'$ are symmetric), so $f_g$ and $f_{g'}$ agree on the overlap as morphisms into $D_+(gg')$. Morphisms of locally ringed spaces glue, so we obtain a morphism $f:U\to\operatorname{Proj}S$ with $f|_{D_+(\varphi g)}=f_g$.

(c) \emph{$U=\operatorname{Proj}T$.} Let $\fp\in\operatorname{Proj}T$; pick a homogeneous $t\in T_e\setminus\fp$ with $e>0$ (possible as $\fp\not\supseteq T_+$). Take $m$ with $me\ge d_0$. Then $t^m\notin\fp$ and $t^m\in T_{me}=\varphi(S_{me})$, so $t^m=\varphi(s)$ with $s\in S_+$; thus $\varphi(S_+)\not\subseteq\fp$.

\emph{$f$ is an isomorphism.} For a homogeneous $g\in S_+$ and $m\ge1$ we have $D_+(g)=D_+(g^m)$ and $S_{(g)}=S_{(g^m)}$, $T_{(\varphi g)}=T_{(\varphi(g)^m)}$ (as subrings of $S_g=S_{g^m}$, $T_{\varphi g}=T_{\varphi(g)^m}$); so the sets $D_+(g)$ with $g\in S_d$, $d\ge d_0$, already cover $\operatorname{Proj}S$ (each $D_+(g)$ with $\deg g=e>0$ equals $D_+(g^m)$, $me\ge d_0$). By (b), $f^{-1}(D_+(g))=D_+(\varphi g)$ and $f$ restricts to the morphism induced by $\varphi_{(g)}:S_{(g)}\to T_{(\varphi g)}$. I claim $\varphi_{(g)}$ is an isomorphism when $d=\deg g\ge d_0$.
\emph{Surjective:} an element of $T_{(\varphi g)}$ is $t/\varphi(g)^n$ with $t\in T_{nd}$; multiplying numerator and denominator by a power of $\varphi(g)$ we may assume $n\ge1$, so $nd\ge d_0$ and we can write $t=\varphi(s)$, $s\in S_{nd}$, and $t/\varphi(g)^n=\varphi_{(g)}(s/g^n)$.
\emph{Injective:} if $\varphi(s)/\varphi(g)^n=0$ in $T_{(\varphi g)}$ with $s\in S_{nd}$, there is $m\ge1$ (enlarging $m$ if necessary) with $\varphi(g)^m\varphi(s)=\varphi(g^ms)=0$; here $g^ms\in S_{(m+n)d}$ and $\varphi_{(m+n)d}$ is injective ($(m+n)d\ge d_0$), so $g^ms=0$, whence $s/g^n=0$.
So $f$ is an isomorphism over each $D_+(g)$ ($\deg g\ge d_0$), these cover $\operatorname{Proj}S$, and a morphism which is an isomorphism over the members of an open cover of the target is an isomorphism (Ex.~2.17(a)). Hence $f:\operatorname{Proj}T\to\operatorname{Proj}S$ is an isomorphism.

(d) Let $V\subseteq\mathbf P^n_k$ be a projective variety ($k$ algebraically closed, $V$ irreducible closed), $S=S(V)=k[x_0,\dots,x_n]/I(V)$, a domain since $I(V)$ is prime, and $V_i=V\cap\{x_i\ne0\}$.

\emph{Points.} By the projective Nullstellensatz (Ch.~I, Ex.~2.1, 2.3) $Y\mapsto I(Y)$ is an inclusion-reversing bijection between nonempty irreducible closed subsets $Y\subseteq V$ and homogeneous prime ideals of $S$ not containing $S_+$ (i.e.\ homogeneous primes $\fp\supseteq I(V)$ with $\fp\not\supseteq(x_0,\dots,x_n)$). So we have a bijection $\beta:t(V)\to\operatorname{Proj}S$, $Y\mapsto I(Y)$ (viewed in $S$). Moreover $Y\in t(V_i)$ iff $Y\cap V_i\ne\emptyset$ iff $x_i\notin I(Y)$ iff $\beta(Y)\in D_+(x_i)$. As $S_+$ is generated by the $x_i$, the $D_+(x_i)$ cover $\operatorname{Proj}S$, and the $t(V_i)$ cover $t(V)$.

\emph{Charts.} The ring $S_{(x_i)}=k[x_0/x_i,\dots,x_n/x_i]\subseteq S_{x_i}$ is the image of $\mu_i:k[y_0,\dots,\widehat{y_i},\dots,y_n]\to S_{x_i}$, $y_j\mapsto x_j/x_i$. Its kernel is $\{h: h^{\mathrm{hom}}\in I(V)\}$, where $h^{\mathrm{hom}}$ is the homogenization: because $S$ is a domain, $h(x_0/x_i,\dots)=0$ in $S_{x_i}$ iff $x_i^N h^{\mathrm{hom}}=0$ in $S$ for suitable $N$, iff $h^{\mathrm{hom}}\in I(V)$. If $V_i\neq\emptyset$ this kernel is $I(V_i)$: a form $F\in k[x]$ vanishes on $V$ iff it vanishes on the dense open subset $V_i$ of the irreducible space $V$ (if $F|_{V_i}=0$ then $x_iF\in I(V)$ and $x_i\notin I(V)$, a prime). If $V_i=\emptyset$, both $S_{x_i}=0$ and $D_+(x_i)=\emptyset=t(V_i)$. So $\mu_i$ induces an isomorphism of $k$-algebras $\bar\mu_i:A(V_i)\xrightarrow{\ \sim\ }S_{(x_i)}$, hence by Prop.~2.6 (where $t(V_i)\cong\Spec A(V_i)$) and Prop.~2.5(b)
\[\Theta_i:\ t(V_i)\cong\Spec A(V_i)\cong\Spec S_{(x_i)}\cong D_+(x_i)\subseteq\operatorname{Proj}S.\]
On points, $\Theta_i$ is $Y\mapsto I(Y)$: a prime $\mathfrak P$ of $A(V_i)$, being $I(Y\cap V_i)$ for $Y=\overline{Y\cap V_i}\in t(V)$, corresponds under $\bar\mu_i$ to the prime of $S_{(x_i)}$ given by the dehomogenizations $\{h:h^{\mathrm{hom}}\in I(Y)\}$ (same argument as above applied to the irreducible $Y$), i.e.\ $I(Y)S_{x_i}\cap S_{(x_i)}$, which is what Prop.~2.5(b) attaches to $I(Y)$. So all $\Theta_i$ are given by $\beta$ on points.

\emph{Gluing.} Put $W=V_i\cap V_j$ (if empty there is nothing to check). Then $t(W)=t(V_i)\cap t(V_j)$ is the principal open set $D(x_j/x_i)\subseteq\Spec A(V_i)$ and equals $D(x_i/x_j)\subseteq\Spec A(V_j)$; $W$ is an affine variety with $A(W)=A(V_i)_{x_j/x_i}=A(V_j)_{x_i/x_j}$, and by Prop.~2.6 $\Gamma(t(W),\mathcal O)=\mathcal O_V(W)=A(W)$. Under $\Theta_i$ (resp.\ $\Theta_j$) the set $t(W)$ goes isomorphically onto $D_+(x_ix_j)=D_+(x_i)\cap D_+(x_j)$, which is affine, $\cong\Spec S_{(x_ix_j)}$ with $S_{(x_ix_j)}=(S_{(x_i)})_{x_j/x_i}=(S_{(x_j)})_{x_i/x_j}$ (as in the proof of Prop.~2.5(b)). So $\Theta_i|_{t(W)}$ and $\Theta_j|_{t(W)}$ are two morphisms between the affine schemes $t(W)=\Spec A(W)$ and $\Spec S_{(x_ix_j)}$, hence by Prop.~2.3(c) (or Ex.~2.4) they are equal as soon as the ring maps $\lambda_i,\lambda_j:S_{(x_ix_j)}\to A(W)=\mathcal O_V(W)$ they induce are equal. Now $\Theta_i|_{t(W)}$ is, by Ex.~2.1 and the definition of $\Theta_i$, $\Spec$ of the localization at $x_j/x_i$ of $\bar\mu_i^{-1}$, so $\lambda_i$ sends $F/(x_i^ax_j^b)$ ($F$ a form of degree $a+b$, which equals $\bigl(F/x_i^{a+b}\bigr)/(x_j/x_i)^b$) to $\bigl(F(y)/y_i^{a+b}\bigr)\big/(x_j/x_i)^b$, i.e.\ to the regular function $P\mapsto F(P)/\bigl(x_i(P)^ax_j(P)^b\bigr)$ on $W$ (here $F/x_i^{a+b}$ corresponds under $\bar\mu_i$ to the dehomogenization of $F$, the function $P\mapsto F(P)/x_i(P)^{a+b}$). The same computation with $i,j$ interchanged gives the same function for $\lambda_j$. Hence $\lambda_i=\lambda_j$ and $\Theta_i=\Theta_j$ on $t(V_i)\cap t(V_j)$. Therefore the $\Theta_i$ glue to a morphism $\Theta:t(V)\to\operatorname{Proj}S$ with $\Theta^{-1}(D_+(x_i))=t(V_i)$ (by the bijection on points above) and $\Theta|_{t(V_i)}=\Theta_i$ an isomorphism onto $D_+(x_i)$. As the $D_+(x_i)$ cover $\operatorname{Proj}S$, $\Theta$ is an isomorphism (Ex.~2.17(a)), i.e.\ $t(V)\cong\operatorname{Proj}S$.\hfill$\square$

\exercise{2.15}{\begin{enumerate}[label=(\alph*),leftmargin=*]
\item Let $V$ be a variety over the algebraically closed field $k$. Show that a point $P\in t(V)$ is a closed point if and only if its residue field is $k$.
\item If $f:X\to Y$ is a morphism of schemes over $k$, and if $P\in X$ is a point with residue field $k$, then $f(P)\in Y$ also has residue field $k$.
\item Now show that if $V,W$ are any two varieties over $k$, then the natural map
\[\Hom_{\mathrm{Var}(k)}(V,W)\to\Hom_{\mathrm{Sch}(k)}(t(V),t(W))\]
is bijective. (Injectivity is easy. The hard part is to show it is surjective.)
\end{enumerate}
}
\pf
Recall (Prop.~2.6) the construction: $t(V)$ is the set of nonempty irreducible closed subsets of $V$; $\alpha:V\to t(V)$, $v\mapsto\overline{\{v\}}=\{v\}$, is a homeomorphism of $V$ onto the set of closed points of $t(V)$ (with the induced topology); the open subsets of $t(V)$ are the sets $t(U)$, $U\subseteq V$ open, and the structure sheaf satisfies $\mathcal O_{t(V)}(t(U))=\mathcal O_V(U)$, the ring of regular functions on $U$, compatibly with restriction. For $v\in U$ the value $\phi(v)\in k$ of $\phi\in\mathcal O_V(U)$ is the image of the section in $k(\alpha(v))=k$. If $V_0\subseteq V$ is an affine open subvariety then $t(V_0)\cong\Spec A(V_0)$, with $\alpha(v)$ corresponding to the maximal ideal $\fm_v$. The scheme $t(V)$ is a scheme over $k$ (constants are regular functions).

(a) ($\Rightarrow$) If $P=\alpha(v)$ is closed, take an affine open $V_0\ni v$; $P$ corresponds to a maximal ideal $\fm$ of the finitely generated $k$-algebra $A=A(V_0)$, and $k(P)=A_\fm/\fm A_\fm=A/\fm$. This is a field which is finitely generated as a $k$-algebra, hence a finite extension of $k$ by the Nullstellensatz (Zariski's lemma, Ch.~I Thm.~1.3A), hence equal to $k$ because $k$ is algebraically closed.
($\Leftarrow$) Suppose $k(P)=k$, meaning that the structure map $k\to k(P)$ is an isomorphism. Choose an affine open $t(V_0)=\Spec A\ni P$, $P=\fp$. Then $k\to A\to A/\fp\hookrightarrow k(P)=\operatorname{Frac}(A/\fp)$, and the image of $k$ is already all of $k(P)$; so $A/\fp=k(P)$ is a field and $\fp$ is a maximal ideal. By the description of $t(V_0)\cong\Spec A(V_0)$ above, maximal ideals correspond to the points $\alpha(v)$, $v\in V_0$. So $P\in\alpha(V)$ is a closed point of $t(V)$.

(b) A morphism of schemes over $k$ gives a local $k$-algebra homomorphism $f^\#_P:\mathcal O_{Y,f(P)}\to\mathcal O_{X,P}$, hence a $k$-algebra homomorphism of fields $\lambda:k(f(P))\to k(P)$. The composite $k\to k(f(P))\xrightarrow{\lambda}k(P)$ is the structure map of $k(P)$, an isomorphism by hypothesis. Both $k\to k(f(P))$ and $\lambda$ are field homomorphisms, hence injective. Given $y\in k(f(P))$, $\lambda(y)=\lambda(c\cdot1)$ for some $c\in k$ (surjectivity of $k\to k(P)$), so $y=c\cdot1$ by injectivity of $\lambda$. Thus $k\to k(f(P))$ is surjective, i.e.\ $k(f(P))=k$.

(c) \emph{Injectivity.} If $h:V\to W$ is a morphism of varieties, $t(h)$ acts on closed points as $\alpha_W\circ h\circ\alpha_V^{-1}$; so $t(h)=t(h')$ gives $h=h'$ as maps, hence as morphisms of varieties.

\emph{Surjectivity.} Let $g:t(V)\to t(W)$ be a morphism of schemes over $k$. By (a), (b), (a) again, $g$ maps closed points to closed points. Define $h:V\to W$ by $h=\alpha_W^{-1}\circ g\circ\alpha_V$. It is continuous, being a restriction of the continuous $g$ to the closed points, which carry the subspace topology.

(i) \emph{$g^{-1}(t(U'))=t(h^{-1}U')$ for $U'\subseteq W$ open.} The open set $g^{-1}(t(U'))$ equals $t(U'')$ for some open $U''\subseteq V$ (all open sets of $t(V)$ have this form), and $v\in U''$ iff $\alpha(v)\in g^{-1}(t(U'))$ iff $g(\alpha(v))=\alpha_W(h(v))\in t(U')$ iff $h(v)\in U'$. So $U''=h^{-1}(U')$.

(ii) \emph{$h$ is a morphism of varieties.} Let $\phi\in\mathcal O_W(U')=\mathcal O_{t(W)}(t(U'))$ and consider $\phi'=g^\#(\phi)\in\mathcal O_{t(V)}(t(h^{-1}U'))=\mathcal O_V(h^{-1}U')$, a regular function. For $v\in h^{-1}(U')$ and $P=\alpha(v)$, the value $\phi'(v)$ is the image of $\phi'_P$ in $k(P)=k$. Since $g^\#_P:\mathcal O_{g(P)}\to\mathcal O_P$ is a local $k$-algebra map inducing the $k$-linear map $k(g(P))=k\to k(P)=k$, which is the identity of $k$, we get $\phi'(v)=\phi(h(v))$. So $\phi\circ h=\phi'$ is regular on $h^{-1}(U')$, and $h$ is a morphism of varieties (Ch.~I, Section~3).

(iii) \emph{$t(h)=g$.} On closed points they agree by construction. For a non-closed point $\eta\in t(V)$, $\eta=Y$ is an irreducible closed subset, with $\overline{\{\eta\}}=t(Y)$ whose closed points are $\alpha(Y)$, and $\alpha(Y)$ is dense in $t(Y)$ (a closed set $t(Z)$ containing $\alpha(Y)$ has $Z\supseteq Y$). Let $Z=\overline{h(Y)}$, an irreducible closed subset of $W$, and $\eta'=g(\eta)=Y'\in t(W)$. Since $g$ is continuous, $g(t(Y))\subseteq\overline{\{\eta'\}}=t(Y')$, so $h(Y)\subseteq Y'$ and $Z\subseteq Y'$. Conversely, $g(\eta)\in g(\overline{\alpha(Y)})\subseteq\overline{g(\alpha(Y))}=\overline{\alpha_W(h(Y))}=t(Z)$, so $\eta'=Y'\in t(Z)$, i.e.\ $Y'\subseteq Z$. So $g(\eta)=Z=\overline{h(Y)}=t(h)(\eta)$. Hence $g=t(h)$ as continuous maps. Finally, $g^\#$ and $t(h)^\#$ both map $\mathcal O_{t(W)}(t(U'))=\mathcal O_W(U')$ into $\mathcal O_V(h^{-1}U')$, and both send $\phi$ to $\phi\circ h$ by (ii) (for $t(h)$ this is the definition); since a regular function on an open subset of a variety is determined by its values at points, $g^\#=t(h)^\#$. Hence $g=t(h)$, and the map is surjective.\hfill$\square$

\exercise{2.16}{Let $X$ be a scheme. For any $f\in\Gamma(X,\mathcal O_X)$, define $X_f$ to be the subset of points $x\in X$ such that the stalk $f_x$ of $f$ at $x$ is not contained in the maximal ideal $\fm_x$ of the local ring $\mathcal O_x$.
\begin{enumerate}[label=(\alph*),leftmargin=*]
\item If $U=\Spec B$ is an open affine subscheme of $X$, and if $\bar f\in B=\Gamma(U,\mathcal O_X|_U)$ is the restriction of $f$, show that $U\cap X_f=D(\bar f)$. Conclude that $X_f$ is an open subset of $X$.
\item Assume that $X$ is quasi-compact. Let $A=\Gamma(X,\mathcal O_X)$, and let $a\in A$ be an element whose restriction to $X_f$ is $0$. Show that for some $n>0$, $f^na=0$. [\emph{Hint}: Use an open affine cover of $X$.]
\item Now assume that $X$ has a finite cover by open affines $U_i$ such that each intersection $U_i\cap U_j$ is quasi-compact (this hypothesis is satisfied, for example, if $\operatorname{sp}(X)$ is noetherian). Let $b\in\Gamma(X_f,\mathcal O_{X_f})$. Show that for some $n>0$, $f^nb$ is the restriction of an element of $A$.
\item With the hypothesis of (c), conclude that $\Gamma(X_f,\mathcal O_{X_f})\cong A_f$.
\end{enumerate}
}
\pf
Here $\mathcal O_{X_f}=\mathcal O_X|_{X_f}$.

(a) Let $\fp\in U$. The germ $f_\fp$ is the image of $\bar f$ under $B\to B_\fp=\mathcal O_{\fp}$, and an element $b\in B$ has image in $\fp B_\fp$ iff $b\in\fp$. So $\fp\in X_f$ iff $\bar f\notin\fp$ iff $\fp\in D(\bar f)$. Thus $U\cap X_f=D(\bar f)$ is open in $U$. As $X$ is covered by affine opens, $X_f$ is open.

(b) Cover $X$ by finitely many affine opens $U_i=\Spec B_i$, $i=1,\dots,r$ ($X$ quasi-compact), and let $a_i,f_i\in B_i$ be the restrictions of $a,f$. By (a), $X_f\cap U_i=D(f_i)$ and $a$ restricts to $0$ in $\Gamma(D(f_i),\mathcal O)=(B_i)_{f_i}$ (Prop.~2.2(b)); the image of $a_i$ in $(B_i)_{f_i}$ is $0$, so $f_i^{n_i}a_i=0$ in $B_i$ for some $n_i$. With $n=\max n_i$, $f^na$ restricts to $0$ on each $U_i$, hence $f^na=0$ in $A$.

(c) Let $U_i=\Spec B_i$ ($i=1,\dots,r$) be the given cover and $f_i=f|_{U_i}$. Then $X_f\cap U_i=D(f_i)$ and $b|_{D(f_i)}\in\Gamma(D(f_i),\mathcal O)=(B_i)_{f_i}$, so $b|_{D(f_i)}=b_i'/f_i^{n_i}$ with $b_i'\in B_i$. Taking $n=\max n_i$ (and multiplying numerators by powers of $f_i$) we may write $f^nb|_{D(f_i)}=b_i|_{D(f_i)}$ with $b_i\in B_i=\Gamma(U_i,\mathcal O)$.

On $U_{ij}=U_i\cap U_j$ the sections $b_i|_{U_{ij}}$ and $b_j|_{U_{ij}}$ both restrict on $U_{ij}\cap X_f$ to $f^nb$ (since $U_{ij}\cap X_f\subseteq D(f_i)\cap D(f_j)$). The scheme $U_{ij}$ is quasi-compact by hypothesis, and $(U_{ij})_f=U_{ij}\cap X_f$. By (b) applied to the scheme $U_{ij}$, the section $f|_{U_{ij}}$, and the section $(b_i-b_j)|_{U_{ij}}$ (which vanishes on $(U_{ij})_f$), there is $m_{ij}$ with $f^{m_{ij}}(b_i-b_j)|_{U_{ij}}=0$. Let $m$ be the maximum of the finitely many $m_{ij}$. Then the sections $f^mb_i\in\mathcal O(U_i)$ agree on all overlaps, so they glue to $c\in A=\Gamma(X,\mathcal O_X)$ with $c|_{U_i}=f^mb_i$. On $D(f_i)=X_f\cap U_i$ we get $c=f^mb_i=f^{m+n}b$. So $c|_{X_f}=f^{m+n}b$ (sheaf axiom for the cover $\{D(f_i)\}$ of $X_f$).

(d) The restriction $A\to\Gamma(X_f,\mathcal O_X)$ sends $f$ to a unit (the germ of $f$ is a unit at each point of $X_f$, so $f|_{X_f}$ has a local inverse near each point and these glue), so it induces a ring homomorphism $\rho:A_f\to\Gamma(X_f,\mathcal O_X)$, $a/f^n\mapsto a|_{X_f}\cdot(f|_{X_f})^{-n}$. \emph{Surjective:} by (c), for $b\in\Gamma(X_f,\mathcal O_X)$ there are $N$ and $c\in A$ with $c|_{X_f}=f^Nb$, so $b=\rho(c/f^N)$. \emph{Injective:} if $\rho(a/f^n)=0$ then $a|_{X_f}=0$, so $f^ma=0$ by (b) ($X$ is quasi-compact, being a finite union of affine schemes), so $a/f^n=0$ in $A_f$. Hence $\Gamma(X_f,\mathcal O_{X_f})\cong A_f$.\hfill$\square$

\exercise{2.17}{\emph{A Criterion for Affineness.}
\begin{enumerate}[label=(\alph*),leftmargin=*]
\item Let $f:X\to Y$ be a morphism of schemes, and suppose that $Y$ can be covered by open subsets $U_i$, such that for each $i$, the induced map $f^{-1}(U_i)\to U_i$ is an isomorphism. Then $f$ is an isomorphism.
\item A scheme $X$ is affine if and only if there is a finite set of elements $f_1,\dots,f_r\in A=\Gamma(X,\mathcal O_X)$, such that the open subsets $X_{f_i}$ are affine, and $f_1,\dots,f_r$ generate the unit ideal in $A$. [\emph{Hint}: Use (Ex.~2.4) and (Ex.~2.16d) above.]
\end{enumerate}}
\pf
(a) Let $g_i:U_i\to f^{-1}(U_i)\subseteq X$ be the inverse of the isomorphism $f_i=f|_{f^{-1}U_i}:f^{-1}(U_i)\to U_i$ (as a morphism to $X$). On $U_i\cap U_j$, both $g_i|_{U_{ij}}$ and $g_j|_{U_{ij}}$ are morphisms $U_{ij}\to f^{-1}(U_{ij})$ that are inverse to $f|_{f^{-1}(U_{ij})}:f^{-1}(U_{ij})\to U_{ij}$ (the latter is an isomorphism, being the restriction of the isomorphism $f_i$ to the open subscheme $f^{-1}(U_{ij})$). Since an inverse is unique, $g_i=g_j$ on $U_{ij}$, and morphisms of locally ringed spaces glue: there is $g:Y\to X$ with $g|_{U_i}=g_i$. Then $f\circ g=\mathrm{id}_Y$ (check on each $U_i$), and $g\circ f=\mathrm{id}_X$ (check on each $f^{-1}(U_i)$, which cover $X$: there $g f=g_i f_i=\mathrm{id}$). So $f$ is an isomorphism.

(b) ($\Rightarrow$) If $X=\Spec A$ take $r=1$, $f_1=1$: $X_1=X$ is affine and $(1)=A$.

($\Leftarrow$) Let $f_1,\dots,f_r\in A$ generate the unit ideal with each $X_{f_i}$ affine. Let $\pi:X\to\Spec A$ be the morphism corresponding (Ex.~2.4) to $\mathrm{id}_A:A\to\Gamma(X,\mathcal O_X)$. By formula $(*)$ in the proof of Ex.~2.4, $\pi(x)=\{a\in A:a_x\in\fm_x\}$, so $\pi^{-1}(D(f_i))=\{x:(f_i)_x\notin\fm_x\}=X_{f_i}$. The $D(f_i)$ cover $\Spec A$ since the $f_i$ generate $(1)$. By (a) it suffices to show that $\pi|_{X_{f_i}}:X_{f_i}\to D(f_i)=\Spec A_{f_i}$ is an isomorphism for each $i$.

First, the hypotheses of Ex.~2.16(c) hold for $X$: the $X_{f_i}$ are finitely many open affines covering $X$ (they cover since $\pi^{-1}(D(f_i))$ cover $X$), and $X_{f_i}\cap X_{f_j}$ is the set $\{x\in X_{f_i}:(f_j)_x\notin\fm_x\}=D(f_j|_{X_{f_i}})$ by Ex.~2.16(a) applied to the affine scheme $X_{f_i}$; it is an affine scheme $\Spec(B_i)_{f_j}$, in particular quasi-compact (Ex.~2.13(b)). So by Ex.~2.16(d), the map $\rho_i:A_{f_i}\to\Gamma(X_{f_i},\mathcal O_X)=:B_i$ is an isomorphism.

Now $\pi|_{X_{f_i}}:\Spec B_i=X_{f_i}\to D(f_i)\cong\Spec A_{f_i}$ is a morphism of affine schemes; by Prop.~2.3(c) and Ex.~2.4 it is induced by its map on global sections $\Gamma(\pi^\#):A_{f_i}=\mathcal O(D(f_i))\to\mathcal O_X(X_{f_i})=B_i$. That map is a ring homomorphism whose composite with $A\to A_{f_i}$ is the restriction $A\to B_i$ (since $\Gamma(\pi^\#)=\mathrm{id}_A$ on global sections and $\pi^\#$ commutes with restrictions); hence it equals $\rho_i$, an isomorphism. So $\pi|_{X_{f_i}}$ is an isomorphism of affine schemes. By (a), $\pi:X\to\Spec A$ is an isomorphism, i.e.\ $X$ is affine.\hfill$\square$

\exercise{2.18}{In this exercise, we compare some properties of a ring homomorphism to the induced morphism of the spectra of the rings.
\begin{enumerate}[label=(\alph*),leftmargin=*]
\item Let $A$ be a ring, $X=\Spec A$, and $f\in A$. Show that $f$ is nilpotent if and only if $D(f)$ is empty.
\item Let $\varphi:A\to B$ be a homomorphism of rings, and let $f:Y=\Spec B\to X=\Spec A$ be the induced morphism of affine schemes. Show that $\varphi$ is injective if and only if the morphism of sheaves $f^\#:\mathcal O_X\to f_*\mathcal O_Y$ is injective. Show furthermore that in that case $f$ is dominant, i.e., $f(Y)$ is dense in $X$.
\item With the same notation, show that if $\varphi$ is surjective, then $f$ is a homeomorphism of $Y$ onto a closed subset of $X$, and $f^\#:\mathcal O_X\to f_*\mathcal O_Y$ is surjective.
\item Prove the converse to (c), namely, if $f:Y\to X$ is a homeomorphism onto a closed subset, and if $f^\#:\mathcal O_X\to f_*\mathcal O_Y$ is surjective, then $\varphi$ is surjective. [\emph{Hint}: Consider $X'=\Spec A/\Ker\varphi$ and use (b) and (c).]
\end{enumerate}}
\pf
(a) $D(f)=\emptyset$ iff $f$ lies in every prime ideal iff $f\in\Nil(A)$, the intersection of all primes, iff $f$ is nilpotent (Atiyah--Macdonald, Prop.~1.8).

(b) ($\Leftarrow$) Taking global sections, $f^\#$ gives $\varphi:A=\Gamma(X,\mathcal O_X)\to\Gamma(X,f_*\mathcal O_Y)=B$ (Prop.~2.2(c)); an injective morphism of sheaves is injective on the sections over every open set (the kernel presheaf is the kernel sheaf, which is $0$).
($\Rightarrow$) Assume $\varphi$ injective. For $g\in A$, $f^{-1}(D(g))=D(\varphi(g))$ and $f^\#(D(g)):A_g\to B_{\varphi(g)}$ sends $a/g^n$ to $\varphi(a)/\varphi(g)^n$. If this is $0$, then $\varphi(g)^m\varphi(a)=\varphi(g^ma)=0$ for some $m$, so $g^ma=0$ and $a/g^n=0$. Thus $f^\#$ is injective on sections over every basic open set $D(g)$, hence on all stalks (stalks are direct limits over the basic neighbourhoods, and direct limits of injections are injective), hence $f^\#$ is injective.
\emph{Dominance.} Let $\fb=\bigcap_{\mathfrak q\in Y}\varphi^{-1}(\mathfrak q)$, the intersection of the primes in $f(Y)$. Then $\overline{f(Y)}=V(\fb)$ (the closure of a set of primes $\Sigma$ is $V(\bigcap\Sigma)$), and $\fb=\varphi^{-1}(\Nil B)=\{a:\varphi(a)\text{ nilpotent}\}=\sqrt{\Ker\varphi}$, since $\varphi(a)^n=0\iff a^n\in\Ker\varphi$. Hence $\overline{f(Y)}=V(\Ker\varphi)$. If $\varphi$ is injective this is $V(0)=X$, so $f$ is dominant.

(c) Let $\varphi$ be surjective, $\fa=\Ker\varphi$, so $B\cong A/\fa$ and $f(\mathfrak q)=\varphi^{-1}(\mathfrak q)$ gives a bijection of $Y$ onto $V(\fa)$ (the primes of $A$ containing $\fa$ correspond to primes of $A/\fa$). For an ideal $\fb\subseteq B$, $f(V(\fb))=V(\varphi^{-1}\fb)$ (as $\varphi$ is onto, primes of $A$ containing $\varphi^{-1}(\fb)$ are exactly the preimages of primes containing $\fb$), a closed subset of $X$ contained in $V(\fa)$. So $f$ is a closed continuous injection, i.e.\ a homeomorphism onto the closed subset $V(\fa)$.

\emph{Surjectivity of $f^\#$.} Check on stalks. For $\fp\in V(\fa)$, $\fp=f(\mathfrak q)$, the stalk of $f_*\mathcal O_Y$ at $\fp$ is $\mathcal O_{Y,\mathfrak q}=B_{\mathfrak q}$, because $f$ is a homeomorphism onto the subspace $V(\fa)$ (the sets $f^{-1}(W)$, $W\ni\fp$ open in $X$, are cofinal among the neighbourhoods of $\mathfrak q$). The stalk map is $A_\fp\to B_{\mathfrak q}$, $a/s\mapsto\varphi(a)/\varphi(s)$ which is surjective: an element of $B_{\mathfrak q}$ is $b/c$ with $c\notin\mathfrak q$; write $b=\varphi(a)$, $c=\varphi(s)$ with $s\notin\fp$. For $\fp\notin V(\fa)=f(Y)$, since $f(Y)$ is closed there is a neighbourhood $W$ of $\fp$ with $f^{-1}(W)=\emptyset$, so the stalk of $f_*\mathcal O_Y$ at $\fp$ is $0$, and surjectivity is trivial. Hence $f^\#$ is surjective.

(d) Let $f:Y\to X$ be a homeomorphism onto a closed subset with $f^\#$ surjective. Let $\fa=\Ker\varphi$, $A'=A/\fa$, $X'=\Spec A'$, and factor $\varphi=\varphi'\circ q$ with $q:A\to A'$ and $\varphi':A'\to B$ injective. Correspondingly $f=i\circ f'$, with $i:X'\to X$ induced by $q$ and $f':Y\to X'$ induced by $\varphi'$. By (c), $i$ is a homeomorphism onto the closed set $V(\fa)$ and $i^\#:\mathcal O_X\to i_*\mathcal O_{X'}$ is surjective. By (b), $f'^\#:\mathcal O_{X'}\to f'_*\mathcal O_Y$ is injective and $f'$ is dominant.

\emph{$f'$ is a homeomorphism.} $f(Y)=i(f'(Y))$ is closed in $X$ by hypothesis, and $f'(Y)$ is dense in $X'$ so $f(Y)=i(f'(Y))$ is dense in $i(X')=V(\fa)$ ($i$ is a homeomorphism onto $V(\fa)$). Since $f(Y)$ is closed, $f(Y)=V(\fa)=i(X')$, so $f'$ is surjective. Then $f'=i^{-1}\circ f$ (where $i^{-1}:V(\fa)\to X'$) is a composite of homeomorphisms $Y\to V(\fa)\to X'$, hence a homeomorphism.

\emph{$f'^\#$ is surjective.} We have $f^\#=i_*(f'^\#)\circ i^\#$. Since $f^\#$ is surjective on stalks, so is $i_*(f'^\#):i_*\mathcal O_{X'}\to i_*f'_*\mathcal O_Y$. At a point $i(x')$ its stalk map is $f'^\#_{x'}$, since $i$ is a homeomorphism onto a closed set (as in (c)); hence $f'^\#$ is surjective on all stalks, hence surjective.

So $f'^\#$ is an isomorphism of sheaves and $f'$ is a homeomorphism: $f'$ is an isomorphism of schemes (its stalk maps are isomorphisms of local rings, being local homomorphisms that are bijective). Taking global sections, $\varphi':A'\to B$ is an isomorphism (by Prop.~2.2(c) and Prop.~2.3(c)). Thus $\varphi=\varphi'\circ q$ is surjective.\hfill$\square$

\exercise{2.19}{Let $A$ be a ring. Show that the following conditions are equivalent:
\begin{enumerate}[label=(\roman*),leftmargin=*]
\item $\Spec A$ is disconnected;
\item there exist nonzero elements $e_1,e_2\in A$ such that $e_1e_2=0$, $e_1^2=e_1$, $e_2^2=e_2$, $e_1+e_2=1$ (these elements are called \emph{orthogonal idempotents});
\item $A$ is isomorphic to a direct product $A_1\times A_2$ of two nonzero rings.
\end{enumerate}}
\pf
(i)$\Rightarrow$(ii). Write $\Spec A=Y_1\sqcup Y_2$ with $Y_1,Y_2$ nonempty, disjoint and closed, $Y_i=V(\fa_i)$. Then $V(\fa_1+\fa_2)=Y_1\cap Y_2=\emptyset$, so $\fa_1+\fa_2=A$ (an ideal not equal to $A$ lies in a prime), and we can write $1=x_1+x_2$ with $x_i\in\fa_i$. Also $V(\fa_1\fa_2)=Y_1\cup Y_2=\Spec A$ (Lemma~2.1(a)), so $\fa_1\fa_2$ lies in every prime, i.e.\ in $\Nil(A)$; in particular $(x_1x_2)^n=0$ for some $n\ge1$. Expanding $1=(x_1+x_2)^{2n-1}$, every monomial $x_1^ax_2^b$ ($a+b=2n-1$) has $a\ge n$ or $b\ge n$; so $1=u_1+u_2$ with $u_1\in(x_1^n)$, $u_2\in(x_2^n)$. Say $u_1=ux_1^n$, $u_2=vx_2^n$. Put $e_1=u_1$, $e_2=u_2$. Then $e_1+e_2=1$ and $e_1e_2=uv(x_1x_2)^n=0$, so $e_1=e_1(e_1+e_2)=e_1^2$ and $e_2=e_2^2$. Finally $e_1\ne0$: otherwise $1=e_2\in(x_2)\subseteq\fa_2$, so $Y_2=V(\fa_2)=\emptyset$, a contradiction; similarly $e_2\neq0$.

(ii)$\Rightarrow$(iii). Let $A_i=Ae_i$, a ring with identity $e_i$ (closed under addition and multiplication, with $e_ia=a$ for $a\in A_i$), nonzero since $e_i\ne0$. Define $\Phi:A\to A_1\times A_2$, $a\mapsto(ae_1,ae_2)$. It is a ring homomorphism with $\Phi(1)=(e_1,e_2)$ the identity. It is injective: $ae_1=ae_2=0$ implies $a=a(e_1+e_2)=0$. It is surjective: for $x,y\in A$, $\Phi(xe_1+ye_2)=(xe_1^2+ye_2e_1,\ xe_1e_2+ye_2^2)=(xe_1,ye_2)$. So $A\cong A_1\times A_2$.

(iii)$\Rightarrow$(ii). Under $A\cong A_1\times A_2$ take $e_1=(1,0)$, $e_2=(0,1)$: they are nonzero as $A_1,A_2\ne0$, and satisfy the required relations.

(ii)$\Rightarrow$(i). $V(e_1)\cap V(e_2)=V(e_1,e_2)=V(A)=\emptyset$ since $e_1+e_2=1$, and $V(e_1)\cup V(e_2)=V(e_1e_2)=V(0)=\Spec A$. Thus $\Spec A$ is the disjoint union of the two closed sets $V(e_1)$, $V(e_2)$. Neither is empty: if $V(e_1)=\emptyset$ then $e_1$ lies in no prime ideal, so it is a unit; but then $e_2=e_1^{-1}e_1e_2=0$, a contradiction; similarly for $e_2$. Each is open (complement of the other), so $\Spec A$ is disconnected.\hfill$\square$

\section*{3. First properties of schemes}

Numbering follows Chapter~II of the scanned edition: Prop.~3.1 (integral $=$ reduced $+$ irreducible), Prop.~3.2 (noetherian is local), Thm.~3.3 (fibred products). For a point $x$ of a scheme, $\mathcal O_x$ is its local ring and $k(x)$ its residue field. For $g$ in a ring $B$ and $V=\Spec B$, $D_V(g)=D(g)$ is the open set where $g\notin\fp$; $D(g)=\Spec B_g$ (Ex.~2.1). If $h:T\to\Spec C$ is a morphism and $c\in C$, then $h^{-1}(D(c))$ is the open set of $T$ where the section $h^\#(c)$ has a unit germ; if $T=\Spec A$ it equals $D(c\cdot 1_A)=\Spec A_c$, where $A_c=A\otimes_CC_c$. A quasi-compact scheme is a finite union of open affine subschemes. Two elementary algebraic remarks are used constantly:
\begin{itemize}
\item if $A$ is a finitely generated $C_g$-algebra then it is a finitely generated $C$-algebra (add $1/g$ to the generators);
\item if $A$ is a finitely generated $C$-algebra then $A_g$ (for $g\in A$ or $g\in C$) is a finitely generated $C$-algebra (add $1/g$) and $A_c=A\otimes_C C_c$ is a finitely generated $C_c$-algebra.
\end{itemize}

\textbf{Lemma A (basic open sets).}
\begin{itemize}
\item \textbf{(A1)} Let $V=\Spec B$ be affine, let $W=\Spec C\subseteq V$ be an open affine subscheme, let $g\in B$ with $D_V(g)\subseteq W$ and let $\bar g\in C$ be the restriction of $g$ to $W$. Then $D_V(g)=D_W(\bar g)$; in particular $B_g\cong C_{\bar g}$.
\item \textbf{(A2)} Let $Y$ be a scheme covered by open affines $V_i=\Spec B_i$, let $V=\Spec B\subseteq Y$ be open affine and $y\in V$. Then there are an index $i$, $h\in B_i$ and $g\in B$ with $y\in D_V(g)\subseteq D_{V_i}(h)\subseteq V\cap V_i$.
\end{itemize}
\emph{Proof.} (A1) For $z\in W$, $z\in D_V(g)$ iff the germ $g_z$ is a unit of $\mathcal O_z$ iff $\bar g_z$ is a unit, because $g_z=\bar g_z$. So $D_V(g)\cap W=D_W(\bar g)$, and $D_V(g)\subseteq W$. The rings of sections are $B_g$ and $C_{\bar g}$ since both sets are the same open subscheme. (A2) Pick $i$ with $y\in V_i$. The open set $V\cap V_i$ is a neighbourhood of $y$ in $V_i$, so it contains some $D_{V_i}(h)\ni y$. This is an open subset of $V$, so it contains some $D_V(g)\ni y$.\hfill$\square$


\exercise{3.1}{Show that a morphism $f:X\to Y$ is locally of finite type if and only if for \emph{every} open affine subset $V=\Spec B$ of $Y$, $f^{-1}(V)$ can be covered by open affine subsets $U_j=\Spec A_j$, where each $A_j$ is a finitely generated $B$-algebra.}
\pf
($\Leftarrow$) Cover $Y$ by open affines $V$; the hypothesis for each of them is exactly the definition.

($\Rightarrow$) By definition there is a cover $Y=\bigcup V_i$, $V_i=\Spec B_i$, and covers $f^{-1}(V_i)=\bigcup_j\Spec A_{ij}$ with $A_{ij}$ finitely generated $B_i$-algebras. Let $V=\Spec B\subseteq Y$ be any open affine and $y\in V$. By (A2) choose $i$, $h\in B_i$, $g\in B$ with $y\in D_V(g)\subseteq W:=D_{V_i}(h)\subseteq V\cap V_i$. Put $C=(B_i)_h$, so $W=\Spec C$.

\emph{Step 1.} We have $f^{-1}(W)=\bigcup_j (f^{-1}(W)\cap\Spec A_{ij})$ and $f^{-1}(W)\cap\Spec A_{ij}=D(h\cdot 1)=\Spec (A_{ij})_h$. Each $(A_{ij})_h=A_{ij}\otimes_{B_i}C$ is a finitely generated $C$-algebra. So $f^{-1}(W)$ is covered by the $\Spec A'_j$, $A'_j:=(A_{ij})_h$, finitely generated over $C=\mathcal O(W)$.

\emph{Step 2.} Let $\bar g\in C$ be the restriction of $g$ to $W$. By (A1), $D_V(g)=D_W(\bar g)\cong\Spec C_{\bar g}$ and $C_{\bar g}\cong B_g$. Applying the remark above to $f^{-1}(W)\to W=\Spec C$, $f^{-1}(D_V(g))=\bigcup_j D(\bar g)\subseteq\Spec A_j'$ is covered by the $\Spec (A'_j)_{\bar g}$, and $(A'_j)_{\bar g}$ is a finitely generated $C_{\bar g}=B_g$-algebra, hence a finitely generated $B$-algebra (add $1/g$).

Thus every $y\in V$ has an open neighbourhood $D_V(g)$ whose preimage is covered by spectra of finitely generated $B$-algebras. Since the $D_V(g)$ cover $V$, $f^{-1}(V)$ is covered by such spectra.\hfill$\square$

\exercise{3.2}{A morphism $f:X\to Y$ of schemes is \emph{quasi-compact} if there is a cover of $Y$ by open affines $V_i$ such that $f^{-1}(V_i)$ is quasi-compact for each $i$. Show that $f$ is quasi-compact if and only if for \emph{every} open affine subset $V\subseteq Y$, $f^{-1}(V)$ is quasi-compact.}
\pf
($\Leftarrow$) is clear. ($\Rightarrow$) Let $\{V_i=\Spec B_i\}$ be a cover with $f^{-1}(V_i)$ quasi-compact.

\emph{Claim Q.} If $T$ is a quasi-compact scheme, $h:T\to\Spec C$ a morphism and $c\in C$, then $h^{-1}(D(c))$ is quasi-compact. Indeed $T=U_1\cup\dots\cup U_n$ with $U_l=\Spec A_l$ affine, and $h^{-1}(D(c))\cap U_l=\Spec (A_l)_c$ is quasi-compact; a finite union of quasi-compact sets is quasi-compact.

Now let $V=\Spec B\subseteq Y$ be open affine and $y\in V$. Choose $i$, $h$, $g$ as in (A2): $y\in D_V(g)\subseteq W:=D_{V_i}(h)\subseteq V\cap V_i$, $W=\Spec C$. By Claim~Q applied to $f^{-1}(V_i)\to V_i$ and $h\in B_i$, $f^{-1}(W)$ is quasi-compact. By (A1), $D_V(g)=D_W(\bar g)$, so Claim~Q applied to $f^{-1}(W)\to W$ and $\bar g$ shows that $f^{-1}(D_V(g))$ is quasi-compact. The quasi-compact space $V$ is covered by finitely many sets $D_V(g_y)$, so $f^{-1}(V)$ is a finite union of quasi-compact sets, hence quasi-compact.\hfill$\square$

\exercise{3.3}{
\begin{enumerate}[label=(\alph*),leftmargin=*]
\item Show that a morphism $f:X\to Y$ is of finite type if and only if it is locally of finite type and quasi-compact.
\item Conclude from this that $f$ is of finite type if and only if for \emph{every} open affine subset $V=\Spec B$ of $Y$, $f^{-1}(V)$ can be covered by a finite number of open affines $U_i=\Spec A_i$, where each $A_i$ is a finitely generated $B$-algebra.
\item Show also that if $f$ is of finite type, then for \emph{every} open affine subset $V=\Spec B\subseteq Y$, and for \emph{every} open affine subset $U=\Spec A\subseteq f^{-1}(V)$, $A$ is a finitely generated $B$-algebra.
\end{enumerate}}
\pf
(a) ($\Rightarrow$) If $Y=\bigcup V_i$ and $f^{-1}(V_i)$ is a finite union of affines $\Spec A_{ij}$ with $A_{ij}$ finitely generated $B_i$-algebras, then $f$ is locally of finite type, and each $f^{-1}(V_i)$ is quasi-compact (finite union of affines), so $f$ is quasi-compact. ($\Leftarrow$) Suppose $f$ is locally of finite type with $Y=\bigcup V_i$, $f^{-1}(V_i)=\bigcup_j\Spec A_{ij}$, $A_{ij}$ finitely generated $B_i$-algebras, and quasi-compact. By Ex.~3.2, $f^{-1}(V_i)$ is quasi-compact, so finitely many of the $\Spec A_{ij}$ already cover it. Hence $f$ is of finite type.

(b) ($\Rightarrow$) By (a) $f$ is locally of finite type and quasi-compact. Let $V=\Spec B$ be open affine. By Ex.~3.1, $f^{-1}(V)$ is covered by open affines $\Spec A_j$, $A_j$ finitely generated over $B$, and by Ex.~3.2 $f^{-1}(V)$ is quasi-compact, so finitely many $\Spec A_j$ suffice. ($\Leftarrow$) Cover $Y$ by open affines; the hypothesis gives the finite covers required in the definition of finite type.

(c) Let $V=\Spec B\subseteq Y$ and $U=\Spec A\subseteq f^{-1}(V)$ be open affines. By (b), $f^{-1}(V)=U_1\cup\dots\cup U_n$ with $U_j=\Spec A_j$, $A_j$ finitely generated $B$-algebras.

\emph{Step 1: $U$ is covered by sets $D_U(g)$ with $A_g$ a finitely generated $B$-algebra.} Let $u\in U$. Apply (A2) to the scheme $f^{-1}(V)$ with the cover $\{U_j\}$ and the open affine $U$: there are $j$, $h\in A_j$, $g\in A$ with $u\in D_U(g)\subseteq D_{U_j}(h)\subseteq U\cap U_j$. Put $C=(A_j)_h$; it is a finitely generated $A_j$-algebra, hence a finitely generated $B$-algebra, and $D_{U_j}(h)=\Spec C$. By (A1), $A_g\cong C_{\bar g}$, which is a finitely generated $C$-algebra, hence a finitely generated $B$-algebra.

\emph{Step 2.} Since $U$ is quasi-compact, there are $g_1,\dots,g_m\in A$ with $D(g_1)\cup\dots\cup D(g_m)=U$, i.e.\ $(g_1,\dots,g_m)=A$, such that every $A_{g_l}$ is a finitely generated $B$-algebra. We show that this forces $A$ to be a finitely generated $B$-algebra. Choose $c_l\in A$ with $\sum c_lg_l=1$. For each $l$ choose generators of the $B$-algebra $A_{g_l}$; adding $1/g_l$ we may assume they are $z_{lk}=y_{lk}/g_l^{n_{lk}}$ with $y_{lk}\in A$ ($k=1,\dots,r_l$). Let $A'\subseteq A$ be the $B$-subalgebra generated by the $g_l$, the $c_l$ and the $y_{lk}$; it is finitely generated. Let $x\in A$. For each $l$, $x/1\in A_{g_l}$ is a polynomial $P_l(z_{l1},\dots,z_{lr_l})$ with coefficients in $B$, of degree $\le d$ say. For $N\ge d\max n_{lk}$ every monomial of $g_l^NP_l(z_l)$ is of the form $b\,g_l^{N-\sum\alpha_kn_{lk}}y_l^\alpha\in A'$. So $g_l^Nx/1=q_l/1$ in $A_{g_l}$ with $q_l\in A'$, hence $g_l^{m}(g_l^Nx-q_l)=0$ in $A$ for some $m$. Taking $M=N+m$ (uniform in $l$), $g_l^Mx\in A'$ for all $l$. Now for $T\ge m'(M-1)+1$ ($m'$ the number of $g_l$'s) expand $1=(\sum c_lg_l)^T$: every monomial has some $g_l$ with exponent $\ge M$, so $1=\sum_le_lg_l^M$ with $e_l$ a polynomial in the $c_i,g_i$, so $e_l\in A'$. Then $x=\sum e_l(g_l^Mx)\in A'$. So $A=A'$ is a finitely generated $B$-algebra.\hfill$\square$

\exercise{3.4}{Show that a morphism $f:X\to Y$ is finite if and only if for \emph{every} open affine subset $V=\Spec B$ of $Y$, $f^{-1}(V)$ is affine, equal to $\Spec A$, where $A$ is a finite $B$-module.}
\pf
($\Leftarrow$) Cover $Y$ by open affines. ($\Rightarrow$) Let $Y=\bigcup V_i$, $V_i=\Spec B_i$, with $f^{-1}(V_i)=\Spec A_i$ and $A_i$ a finite $B_i$-module. Fix an open affine $V=\Spec B\subseteq Y$ and write $X_V=f^{-1}(V)$.

\emph{Step 1: local statement.} Let $y\in V$ and choose $i$, $h\in B_i$, $g\in B$ as in (A2): $y\in D_V(g)\subseteq W:=D_{V_i}(h)\subseteq V\cap V_i$, with $W=\Spec C$, $C=(B_i)_h$. Then $f^{-1}(W)=D(h)\subseteq\Spec A_i$, i.e.\ $f^{-1}(W)=\Spec A_W$ with $A_W=(A_i)_h=A_i\otimes_{B_i}C$ a finite $C$-module. By (A1), $D_V(g)=D_W(\bar g)\cong\Spec B_g$ with $C_{\bar g}\cong B_g$, so $f^{-1}(D_V(g))=\Spec (A_W)_{\bar g}$, and $(A_W)_{\bar g}=A_W\otimes_CC_{\bar g}$ is a finite $B_g$-module.

So $V$ is covered by sets $D(g)$ with $f^{-1}(D(g))$ affine, equal to $\Spec$ of a finite $B_g$-module algebra. By quasi-compactness of $V$ we get $g_1,\dots,g_r\in B$ generating the unit ideal with this property.

\emph{Step 2: $X_V$ is affine.} Let $A=\Gamma(X_V,\mathcal O_{X_V})$ and $s_l=f^\#(g_l)\in A$. Then $s_1,\dots,s_r$ generate the unit ideal of $A$ (apply $f^\#$ to $\sum b_lg_l=1$), and the open set $(X_V)_{s_l}=f^{-1}(D(g_l))$ is affine. By the affineness criterion (Ex.~2.17(b)), $X_V$ is affine, $X_V=\Spec A$. Then $f^{-1}(D(g_l))=D(s_l)=\Spec A_{s_l}$, i.e.\ $A_{s_l}=A\otimes_BB_{g_l}$ is a finite $B_{g_l}$-module.

\emph{Step 3: $A$ is a finite $B$-module.} Let $M=A$ as a $B$-module. Each $M_{g_l}=A_{s_l}$ is finite over $B_{g_l}$. Choose finitely many generators of $M_{g_l}$ over $B_{g_l}$; multiplying by powers of $g_l$ we may take them to be images of elements of $M$. Let $N\subseteq M$ be the submodule generated by all these (finitely many) elements. Then $N_{g_l}=M_{g_l}$ for all $l$, i.e.\ $(M/N)_{g_l}=0$ by exactness of localization. So every element of $M/N$ is killed by a power of each $g_l$, hence by the unit ideal $(g_1^n,\dots,g_r^n)=B$ (for large $n$); so $M=N$ is finite over $B$.\hfill$\square$

\exercise{3.5}{A morphism $f:X\to Y$ is \emph{quasi-finite} if for every point $y\in Y$, $f^{-1}(y)$ is a finite set.
\begin{enumerate}[label=(\alph*),leftmargin=*]
\item Show that a finite morphism is quasi-finite.
\item Show that a finite morphism is \emph{closed}, i.e.\ the image of any closed subset is closed.
\item Show by example that a surjective, finite type, quasi-finite morphism need not be finite.
\end{enumerate}}
\pf
Let $f:X\to Y$ be finite, with $Y=\bigcup V_i$, $V_i=\Spec B_i$, $f^{-1}(V_i)=\Spec A_i$, $A_i$ a finite $B_i$-module.

(a) Let $y\in Y$, choose $i$ with $y\in V_i$, let $\fp\subset B_i$ be the prime of $y$ and $S=B_i\smallsetminus\fp$. A point of $f^{-1}(y)$ is a prime $\fq\subset A_i$ whose contraction to $B_i$ is $\fp$. Such primes correspond bijectively to the primes of $R:=S^{-1}A_i/\fp S^{-1}A_i=A_i\otimes_{B_i}k(y)$ (primes of $S^{-1}A_i$ are the primes of $A_i$ not meeting the image of $S$, i.e.\ with contraction $\subseteq\fp$, and those containing $\fp S^{-1}A_i$ have contraction $\supseteq\fp$; cf.\ Ex.~3.10(a)). Now $R$ is a finite-dimensional $k(y)$-algebra, say of dimension $n$. Its primes are maximal (it is artinian), and if $\fm_1,\dots,\fm_r$ are distinct maximal ideals, the Chinese remainder theorem gives a surjection $R\to\prod R/\fm_i$ onto a $k(y)$-vector space of dimension $\ge r$, so $r\le n$. Hence $f^{-1}(y)$ is finite.

(b) Let $Z\subseteq X$ be closed. A subset of $Y$ is closed iff its intersection with every $V_i$ is closed in $V_i$, and $f(Z)\cap V_i=f(Z\cap\Spec A_i)$. Write $Z\cap\Spec A_i=V(\fa)$ for an ideal $\fa\subseteq A_i$, and let $\fb=\fa\cap B_i$. The map $B_i/\fb\to A_i/\fa$ is injective and integral (since $A_i$ is a finite $B_i$-module), so by the lying-over theorem $\Spec(A_i/\fa)\to\Spec(B_i/\fb)$ is surjective. Therefore $f(V(\fa))=V(\fb)$ is closed in $V_i$. So $f(Z)$ is closed.

(c) Let $k$ be a field, $Y=\Spec k[t]=\mathbf A^1_k$, and $X=\Spec k[t,t^{-1}]\sqcup\Spec k$, with $f$ the open immersion $\Spec k[t,t^{-1}]=D(t)\hookrightarrow Y$ on the first summand and the morphism $\Spec k\to Y$ onto the closed point $t=0$ (corresponding to $k[t]\to k$, $t\mapsto0$) on the second. Then $f$ is surjective. It is quasi-finite: the fibre over a point $\fp\ne(t)$ is the single point of $D(t)$, and the fibre over $(t)$ is the single point of $\Spec k$. It is of finite type: $f^{-1}(Y)=X$ is covered by two affines $\Spec k[t,t^{-1}]$ and $\Spec k$ whose rings are finitely generated $k[t]$-algebras (this is the definition of finite type, using the cover $\{Y\}$ of $Y$). But $f$ is not finite: the first summand is a closed subset of $X$ (its complement is open), and its image $D(t)$ is not closed in $Y$, contradicting (b).\hfill$\square$

\exercise{3.6}{Let $X$ be an integral scheme. Show that the local ring $\mathcal O_\xi$ of the generic point $\xi$ of $X$ is a field. It is called the \emph{function field} of $X$, and is denoted $K(X)$. Show also that if $U=\Spec A$ is any open affine subset of $X$, then $K(X)$ is isomorphic to the quotient field of $A$.}
\pf
By Prop.~3.1, $X$ is irreducible, so it has a unique generic point $\xi$ (Ex.~2.9). Let $U=\Spec A$ be a nonempty open affine subset. Since $X=\overline{\{\xi\}}$, a nonempty open set contains $\xi$ (otherwise $\xi$ would lie in the closed set $X\smallsetminus U$, which would then contain $\overline{\{\xi\}}=X$). So $\xi\in U$, and $\xi$ is a point of $\Spec A$ whose closure in $U$ is all of $U$. If $\xi$ corresponds to the prime $\fp\subset A$, its closure is $V(\fp)$, so $\fp$ is contained in every prime of $A$, hence in the nilradical of $A$. Since $X$ is integral, $A=\mathcal O_X(U)$ is a domain, so $\fp=(0)$. Therefore $\mathcal O_\xi=\mathcal O_{U,\xi}=A_{(0)}=\operatorname{Frac}(A)$, a field, and $K(X)\cong\operatorname{Frac}A$ for every nonempty open affine $U=\Spec A$.\hfill$\square$

\exercise{3.7}{A morphism $f:X\to Y$, with $Y$ irreducible, is \emph{generically finite} if $f^{-1}(\eta)$ is a finite set, where $\eta$ is the generic point of $Y$. A morphism $f:X\to Y$ is \emph{dominant} if $f(X)$ is dense in $Y$. Now let $f:X\to Y$ be a dominant, generically finite morphism of finite type of integral schemes. Show that there is an open dense subset $U\subseteq Y$ such that the induced morphism $f^{-1}(U)\to U$ is finite. [\emph{Hint}: First show that the function field of $X$ is a finite field extension of the function field of $Y$.]}
\pf
Let $\xi,\eta$ be the generic points of $X,Y$, $L=K(X)$, $K=K(Y)$.

\emph{Step 0: $f(\xi)=\eta$.} Since $f$ is continuous, $f(X)=f(\overline{\{\xi\}})\subseteq\overline{\{f(\xi)\}}$, so $\overline{\{f(\xi)\}}\supseteq\overline{f(X)}=Y$, hence $f(\xi)=\eta$ by uniqueness of generic points.

\emph{Step 1: setup.} Let $V=\Spec B\subseteq Y$ be a nonempty open affine subset ($B$ a domain, $K=\operatorname{Frac}B$, $\eta\in V$). By Ex.~3.3(b), $f^{-1}(V)=U_1\cup\dots\cup U_N$ with $U_i=\Spec A_i$ and $A_i$ finitely generated $B$-algebras; discard empty ones, so every $U_i$ is nonempty and $A_i$ is a domain with $\operatorname{Frac}A_i=L$ (Ex.~3.6); $N\ge1$ since $\xi\in f^{-1}(V)$. The generic point $\xi\in U_i$ is the prime $(0)$ of $A_i$ and it maps to $\eta=(0)$ of $B$, so $B\to A_i$ is injective; regard $B\subseteq A_i\subseteq L$.

\emph{Step 2: $L$ is a finite extension of $K$.} With $S=B\smallsetminus\{0\}$, $f^{-1}(\eta)\cap U_1$ is the set of primes $\fq$ of $A_1$ with $\fq\cap B=0$, i.e.\ the set of primes of $D:=S^{-1}A_1$; by hypothesis it is finite. $D\subseteq L$ is a finitely generated $K$-algebra and a domain. By Noether normalization $D$ is finite over a polynomial ring $K[t_1,\dots,t_d]\subseteq D$. By lying over, $\Spec D\to\Spec K[t_1,\dots,t_d]$ is surjective. If $d\ge1$ the polynomial ring has infinitely many primes (the principal primes $(p(t_1))$ for the infinitely many irreducible $p\in K[t_1]$), contradicting finiteness. So $d=0$, $D$ is a domain finite-dimensional over $K$, hence a field. Since $A_1\subseteq D\subseteq L=\operatorname{Frac}A_1$, $D=L$, so $[L:K]<\infty$.

\emph{Step 3: shrinking to make the $A_i$ integral.} Write $A_i=B[x_{i1},\dots,x_{in_i}]\subseteq L$. Every $x_{ij}$ is algebraic over $K$; clearing denominators in a polynomial it satisfies, we get $0\ne a_{ij}\in B$ with $x_{ij}$ a root of a polynomial over $B$ with leading coefficient $a_{ij}$. Let $a=\prod a_{ij}\in B\smallsetminus\{0\}$. Then each $x_{ij}$ is integral over $B_a$, so $A_i':=(A_i)_a=B_a[x_{ij}]$ is a finite $B_a$-module and a domain. Moreover $U_i':=\Spec A_i'=U_i\cap f^{-1}(D(a))$, and $f^{-1}(D(a))=\bigcup_iU_i'$.

\emph{Step 4: a lemma.} \emph{If $R\subseteq S'$ are domains, $S'$ integral over $R$, and $0\ne r\in S'$, then $rS'\cap R\ne0$.} Take a monic equation $r^n+b_1r^{n-1}+\dots+b_n=0$ ($b_i\in R$) of minimal degree $n$. Then $b_n\ne0$ (otherwise cancel one factor $r$ in the domain $S'$), and $b_n=-r(r^{n-1}+\dots+b_{n-1})\in rS'\cap R$.

\emph{Step 5: the pieces $U_i'$ coincide over a smaller set.} Let $2\le i\le N$. The set $U_1'\cap U_i'$ is open in $U_1'=\Spec A_1'$ and contains $\xi$, so it contains a basic open set $D_{U_1'}(r)$ with $0\ne r\in A_1'$. By Step~4 (with $R=B_a$, $S'=A_1'$) there is $0\ne b'\in rA_1'\cap B_a$; writing $b'=b/a^m$ with $b\in B\smallsetminus\{0\}$ we have $b\in rA_1'$ and $D_{U_1'}(b)\subseteq D_{U_1'}(r)\subseteq U_1'\cap U_i'$. Note $D_{U_1'}(b)=U_1'\cap f^{-1}(D(b))$. So $U_1'\cap f^{-1}D(b)\subseteq U_i'$. Interchanging the roles of $1$ and $i$ we find $0\ne b''\in B$ with $U_i'\cap f^{-1}D(b'')\subseteq U_1'$. Put $c_i=bb''$. Then
\[U_1'\cap f^{-1}D(c_i)\subseteq U_i'\cap U_1'\cap f^{-1}D(c_i),\qquad U_i'\cap f^{-1}D(c_i)\subseteq U_1'\cap U_i'\cap f^{-1}D(c_i),\]
hence $U_1'\cap f^{-1}D(c_i)=U_i'\cap f^{-1}D(c_i)$.

\emph{Step 6.} Let $c=a\,c_2\cdots c_N\in B\smallsetminus\{0\}$ and $U=D(c)\subseteq V$. It is a nonempty open subset of the irreducible space $Y$, hence dense. Since $D(c)\subseteq D(a)\cap D(c_i)$, Step~5 gives $U_i'\cap f^{-1}(U)=U_1'\cap f^{-1}(U)$ for every $i$, so
\[f^{-1}(U)=\bigcup_iU_i'\cap f^{-1}(U)=U_1'\cap f^{-1}(D(c))=\Spec (A_1')_c.\]
Now $(A_1')_c$ is a finite module over $(B_a)_c=B_c$ (a localization of the finite $B_a$-module $A_1'$), and $U=\Spec B_c$. So $f^{-1}(U)\to U$ is a finite morphism (take the cover of $U$ by itself).\hfill$\square$

\exercise{3.8}{\emph{Normalization.} A scheme is \emph{normal} if all of its local rings are integrally closed domains. Let $X$ be an integral scheme. For each open affine subset $U=\Spec A$ of $X$, let $\tilde A$ be the integral closure of $A$ in its quotient field, and let $\tilde U=\Spec\tilde A$. Show that one can glue the schemes $\tilde U$ to obtain a normal integral scheme $\tilde X$, called the \emph{normalization} of $X$. Show also that there is a morphism $\tilde X\to X$, having the following universal property: for every normal integral scheme $Z$, and for every dominant morphism $f:Z\to X$, $f$ factors uniquely through $\tilde X$. If $X$ is of finite type over a field $k$, then the morphism $\tilde X\to X$ is a finite morphism. This generalizes (I, Ex.~3.17).}
\pf
Let $K=K(X)$ and $\xi$ the generic point. For a nonempty open affine $U=\Spec A$, $A$ is a domain with $\operatorname{Frac}A=K$ (Ex.~3.6); let $\tilde A_U\subseteq K$ be its integral closure, $\tilde U=\Spec\tilde A_U$, $\pi_U:\tilde U\to U$ induced by $A\subseteq\tilde A_U$. This map is surjective (lying over) and $\tilde U$ is a normal integral scheme: $\tilde A_U$ is a domain, integrally closed in its quotient field $K$, and so are its localizations. We use that integral closure commutes with localization: for $f\in A$, $\tilde{(A_f)}=(\tilde A)_f$ (Atiyah--Macdonald, 5.12).

\emph{Step 1: compatibility with open affine subsets.} For a nonempty open affine $W=\Spec A_W\subseteq X$ write $\tilde W=\Spec\tilde A_W$ and $\pi_W:\tilde W\to W$. Every $A_W=\mathcal O_X(W)$ is a subring of $K=\mathcal O_{X,\xi}$ (the map to the stalk at $\xi$ is injective), and for $W'\subseteq W$ the restriction map $A_W\to A_{W'}$ is the inclusion of subrings of $K$, since both commute with the maps to $\mathcal O_{X,\xi}$. Hence $\tilde A_W\subseteq\tilde A_{W'}$ (an element integral over $A_W$ is integral over $A_{W'}$), and we put $\rho_{W'W}=\Spec(\tilde A_W\subseteq\tilde A_{W'}):\tilde W'\to\tilde W$; it lies over the inclusion $W'\to W$, and clearly $\rho_{W''W'}$, $\rho_{W'W}$ compose to $\rho_{W''W}$ and $\rho_{WW}=\mathrm{id}$. \emph{Claim: $\rho_{W'W}$ is an open immersion with image $\pi_W^{-1}(W')$.} Let $y\in W'$ and choose $f\in A_W$ with $y\in D_W(f)\subseteq W'$. By (A1), $D_W(f)=D_{W'}(\bar f)$ and $(A_W)_f=(A_{W'})_{\bar f}$, and by (Atiyah--Macdonald, 5.12) both $(\tilde A_W)_f$ and $(\tilde A_{W'})_{\bar f}$ are the integral closure of this ring in $K$, so they are the same subring of $K$. The restriction of $\rho_{W'W}$ to the open set $\pi_{W'}^{-1}(D_{W'}(\bar f))=\Spec(\tilde A_{W'})_{\bar f}$ is the spectrum of $\tilde A_W\to(\tilde A_W)_f=(\tilde A_{W'})_{\bar f}$, i.e.\ an isomorphism onto the open set $\Spec(\tilde A_W)_f=\pi_W^{-1}(D_W(f))$. These sets $\pi_{W'}^{-1}(D_{W'}(\bar f))$ cover $\tilde W'$, so $\rho_{W'W}$ is a local isomorphism onto open subsets covering $\pi_W^{-1}(W')$; it is injective, because two points with the same image lie over the same $y\in W'$ and so in the same piece $\pi_{W'}^{-1}(D_{W'}(\bar f))$ with $y\in D_{W'}(\bar f)$, on which $\rho_{W'W}$ is injective. This proves the claim.

\emph{Step 2: gluing.} Index by the family of \emph{all} nonempty open affine subsets $U\subseteq X$, and put $\tilde U_{UU'}:=\pi_U^{-1}(U\cap U')\subseteq\tilde U$. For a nonempty open affine $W\subseteq U\cap U'$, Step~1 gives $\pi_U^{-1}(W)=\rho_{WU}(\tilde W)$, and we define $\varphi_W:=\rho_{WU'}\circ\rho_{WU}^{-1}:\pi_U^{-1}(W)\to\pi_{U'}^{-1}(W)$, an isomorphism over $W$. If $W''\subseteq W$ then $\varphi_{W''}=\varphi_W$ on $\pi_U^{-1}(W'')$, because $\rho_{W''U}=\rho_{WU}\rho_{W''W}$ and $\rho_{W''U'}=\rho_{WU'}\rho_{W''W}$. If $W_1,W_2\subseteq U\cap U'$ are open affine, then $W_1\cap W_2$ is covered by open affines $W_3\subseteq W_1\cap W_2$ and $\varphi_{W_1}$, $\varphi_{W_2}$ both agree with $\varphi_{W_3}$ on $\pi_U^{-1}(W_3)$; so they agree on $\pi_U^{-1}(W_1\cap W_2)$. Since the $W\subseteq U\cap U'$ cover $U\cap U'$, the $\varphi_W$ glue to an isomorphism $\varphi_{UU'}:\tilde U_{UU'}\to\tilde U_{U'U}$ with $\pi_{U'}\varphi_{UU'}=\pi_U$, $\varphi_{U'U}=\varphi_{UU'}^{-1}$ (same formula) and $\varphi_{UU}=\mathrm{id}$. Cocycle condition: $\tilde U_{UU'}\cap\tilde U_{UU''}=\pi_U^{-1}(U\cap U'\cap U'')$ is covered by the sets $\pi_U^{-1}(W)$ with $W\subseteq U\cap U'\cap U''$ open affine, and $\varphi_{UU'}$ maps this set onto $\pi_{U'}^{-1}(U\cap U'\cap U'')=\tilde U_{U'U}\cap\tilde U_{U'U''}$ (as $\pi_{U'}\varphi_{UU'}=\pi_U$); on $\pi_U^{-1}(W)$ we have $\varphi_{U'U''}\varphi_{UU'}=\rho_{WU''}\rho_{WU'}^{-1}\rho_{WU'}\rho_{WU}^{-1}=\varphi_{UU''}$. By Ex.~2.12 we obtain a scheme $\tilde X$, containing every $\tilde U$ as an open subscheme, and a morphism $\pi:\tilde X\to X$ restricting to $\pi_U$ on $\tilde U$. Moreover $\pi^{-1}(U)=\tilde U$ for every nonempty open affine $U$: a point of $\pi^{-1}(U)$ lies in some $\tilde U'$ over a point of $U\cap U'$, i.e.\ in $\tilde U_{U'U}$, which is identified with $\tilde U_{UU'}\subseteq\tilde U$.

$\tilde X$ is normal because its local rings are local rings of the $\tilde U_i$. It is reduced because its local rings are domains. It is irreducible: it is covered by the irreducible open sets $\tilde U_i$, and any two meet because $\pi_i$ is surjective and $U_i\cap U_j\ne\emptyset$ ($X$ irreducible). So $\tilde X$ is integral by Prop.~3.1.

\emph{Step 3: uniqueness lemma.} Let $T$ be an integral scheme with generic point $\zeta$, let $f:T\to X$ be a morphism with $f(\zeta)=\xi$ (e.g.\ $f$ dominant, Step~0 of Ex.~3.7), and let $g,g':T\to\tilde X$ with $\pi g=\pi g'=f$. Then $g=g'$. Indeed, for any nonempty open affine $U=\Spec A\subseteq X$ and nonempty open affine $V=\Spec B\subseteq f^{-1}(U)$, $A\to B$ is injective (it is the restriction of $K(X)\to K(T)$, or: the prime $(0)$ of $B$ lies over $(0)$ of $A$). Both $g|_V$ and $g'|_V$ map into $\pi^{-1}(U)=\Spec\tilde A_U$, so they correspond to ring homomorphisms $\tilde A_U\to B$ extending $A\to B$. Such an extension is unique: if $s=a/a'\in\tilde A_U$ ($a,a'\in A$, $a'\ne0$) then $\varphi(s)\varphi(a')=\varphi(a)$ and $\varphi(a')\ne0$ in the domain $B$. So $g|_V=g'|_V$. Such $V$ cover $T$, hence $g=g'$.

\emph{Step 4: universal property.} Let $Z$ be a normal integral scheme and $f:Z\to X$ dominant. Let $U=\Spec A\subseteq X$ and $V=\Spec B\subseteq f^{-1}(U)$ be nonempty open affines; $A\subseteq B$ as in Step~3, and $B$ is a domain with $\operatorname{Frac}B=K(Z)\supseteq K=\operatorname{Frac}A$. The domain $B$ is integrally closed: if $s\in\operatorname{Frac}B$ is integral over $B$, it is integral over every localization $B_\fq$ (integrally closed since $Z$ is normal), so $s\in\bigcap_{\fq}B_\fq=B$. An element of $\tilde A_U\subseteq K\subseteq K(Z)$ is integral over $A$, hence over $B$, so lies in $B$. Thus $A\to B$ extends to $\tilde A_U\to B$, i.e.\ we get $g_V:V\to\tilde U=\pi^{-1}(U)\subseteq\tilde X$ with $\pi g_V=f|_V$. For two such $V,V'$, the morphisms $g_V$ and $g_{V'}$ agree on $V\cap V'$ by Step~3 applied to the integral scheme $V\cap V'$ and the morphism $f|_{V\cap V'}$. So the $g_V$ glue to $g:Z\to\tilde X$ with $\pi g=f$, and $g$ is unique by Step~3.

\emph{Step 5: finiteness.} Let $X$ be of finite type over $k$. By Ex.~3.3(c), for every open affine $U=\Spec A\subseteq X$, $A$ is a finitely generated $k$-algebra and a domain, so its integral closure $\tilde A$ in $\operatorname{Frac}A$ is a finite $A$-module (finiteness of integral closure for finitely generated domains over a field, as in (I, Ex.~3.17); in characteristic $p$ one needs the general form of this theorem, e.g.\ Eisenbud, \emph{Commutative Algebra}, Cor.~13.13). Hence $\pi^{-1}(U)=\Spec\tilde A_U$ with $\tilde A_U$ a finite $A$-module, for all such $U$; so $\pi$ is finite.\hfill$\square$

\exercise{3.9}{\emph{The Topological Space of a Product.} Recall that in the category of varieties, the Zariski topology on the product of two varieties is not equal to the product topology (I, Ex.~1.4). Now we see that in the category of schemes, the underlying point set of a product of schemes is not even the product set.
\begin{enumerate}[label=(\alph*),leftmargin=*]
\item Let $k$ be a field, and let $\mathbf A^1_k=\Spec k[x]$ be the affine line over $k$. Show that $\mathbf A^1_k\times_{\Spec k}\mathbf A^1_k\cong\mathbf A^2_k$, and show that the underlying point set of the product is not the product of the underlying sets of the factors (even if $k$ is algebraically closed).
\item Let $k$ be a field, let $s$ and $t$ be indeterminates over $k$. Then $\Spec k(s)$, $\Spec k(t)$, and $\Spec k$ are all one-point spaces. Describe the product scheme $\Spec k(s)\times_{\Spec k}\Spec k(t)$.
\end{enumerate}}
\pf
(a) By Thm.~3.3, Step~1, $\mathbf A^1_k\times_k\mathbf A^1_k=\Spec(k[x]\otimes_kk[y])=\Spec k[x,y]=\mathbf A^2_k$, with projections $p_1,p_2$ induced by $k[x]\to k[x,y]$ and $k[y]\to k[x,y]$. Consider the two points $\fp_0=(0)$ and $\fp_1=(x-y)$ of $\mathbf A^2_k$; they are distinct primes. We have $\fp_0\cap k[x]=0$, and also $\fp_1\cap k[x]=0$: a nonzero polynomial $q(x)$ divisible by $x-y$ in $k[x,y]$ would vanish after substituting $y=x$, which gives $q(x)=0$; similarly $\fp_1\cap k[y]=0$. So $p_1(\fp_0)=p_1(\fp_1)$ and $p_2(\fp_0)=p_2(\fp_1)$, i.e.\ the canonical map
\[\mathbf A^2_k\to\mathbf A^1_k\times\mathbf A^1_k\ \text{(product of sets)},\quad z\mapsto(p_1(z),p_2(z))\]
is not injective, a fortiori not bijective: both $\fp_0$ and $\fp_1$ (the generic points of the plane and of the diagonal) map to the pair (generic point, generic point). This holds for every field $k$, in particular for $k$ algebraically closed.

(b) By Thm.~3.3, Step~1, the product is $\Spec(k(s)\otimes_kk(t))$. Since $k(s)=T_1^{-1}k[s]$ and $k(t)=T_2^{-1}k[t]$ with $T_1=k[s]\smallsetminus0$, $T_2=k[t]\smallsetminus0$, we have
\[k(s)\otimes_kk(t)=T^{-1}k[s,t],\qquad T=\{f(s)g(t):f\ne0,g\ne0\}\]
(a multiplicative set). The primes of this ring are the primes $\fp$ of $k[s,t]$ with $\fp\cap T=\emptyset$, i.e.\ (as $\fp$ is prime) with $\fp\cap k[s]=0$ and $\fp\cap k[t]=0$. In the UFD $k[s,t]$ the primes are $(0)$, the principal primes $(h)$ with $h$ irreducible, and the maximal ideals. A maximal ideal $\fm$ meets $k[s]$ nontrivially: $k[s,t]/\fm$ is a field finitely generated as a $k$-algebra, hence finite over $k$ (Zariski's lemma), so the image of $s$ is algebraic over $k$. So no maximal ideal occurs. A principal prime $(h)$ meets $k[s]$ iff $h$ divides a nonzero polynomial $f(s)$, iff (up to a unit) $h\in k[s]$ (all divisors of $f(s)$ lie in $k[s]$); similarly for $k[t]$. Hence:

\emph{$\Spec k(s)\times_{\Spec k}\Spec k(t)$ is the spectrum of the one-dimensional noetherian domain $T^{-1}k[s,t]$. Its points are the generic point $(0)$ and the points $(h)$, $h\in k[s,t]$ irreducible and involving both $s$ and $t$ (i.e.\ $h\notin k[s]\cup k[t]$ up to units); the residue field at $(h)$ is $\operatorname{Frac}(k[s,t]/(h))$.} For example the point $(s-t)$ lies in it. In particular the product is an infinite scheme, though each factor is a single point.\hfill$\square$

\exercise{3.10}{\emph{Fibres of a Morphism.}
\begin{enumerate}[label=(\alph*),leftmargin=*]
\item If $f:X\to Y$ is a morphism, and $y\in Y$ a point, show that $\operatorname{sp}(X_y)$ is homeomorphic to $f^{-1}(y)$ with the induced topology.
\item Let $X=\Spec k[s,t]/(s-t^2)$, let $Y=\Spec k[s]$, and let $f:X\to Y$ be the morphism defined by sending $s\mapsto s$. If $y\in Y$ is the point $a\in k$ with $a\ne0$, show that the fibre $X_y$ consists of two points, with residue field $k$. If $y\in Y$ corresponds to $0\in k$, show that the fibre $X_y$ is a nonreduced one-point scheme. If $\eta$ is the generic point of $Y$, show that $X_\eta$ is a one-point scheme, whose residue field is an extension of degree $2$ of the residue field of $\eta$. (Assume $k$ algebraically closed.)
\end{enumerate}}
\pf
(a) Recall $X_y=X\times_Y\Spec k(y)$ and let $p:X_y\to X$ be the projection. Choose an open affine $V=\Spec B\ni y$, with $y$ corresponding to the prime $\fp\subset B$. The morphism $\Spec k(y)\to Y$ factors through $V$, so any scheme $T$ with maps to $X$ and to $\Spec k(y)$ over $Y$ maps into $f^{-1}(V)$; hence $X_y=f^{-1}(V)\times_V\Spec k(y)$, and $p$ lands in $f^{-1}(V)$. Cover $f^{-1}(V)$ by open affines $U_j=\Spec A_j$. By Thm.~3.3 (Steps 1 and 4), $p^{-1}(U_j)=U_j\times_V\Spec k(y)=\Spec(A_j\otimes_Bk(y))$. With $S=B\smallsetminus\fp$ and $\varphi:B\to A_j$ we have $k(y)=S^{-1}B/\fp S^{-1}B$, so
\[R_j:=A_j\otimes_Bk(y)=S^{-1}A_j/\fp S^{-1}A_j.\]
Primes of $S^{-1}A_j$ correspond to primes $\fq$ of $A_j$ with $\varphi^{-1}(\fq)\cap S=\emptyset$, i.e.\ $\varphi^{-1}(\fq)\subseteq\fp$, and those containing $\fp S^{-1}A_j$ are the ones with $\varphi^{-1}(\fq)\supseteq\fp$. So $\Spec R_j\to\Spec A_j$ is a bijection onto $\{\fq:\varphi^{-1}(\fq)=\fp\}=f^{-1}(y)\cap U_j$. It is a homeomorphism onto this subset with the induced topology: it is the composite of the closed embedding $\Spec R_j\hookrightarrow\Spec S^{-1}A_j$ and the map $\Spec S^{-1}A_j\to\Spec A_j$, which is a homeomorphism onto its image with the subspace topology (the closed sets $V(S^{-1}\fa)$ are the traces of the closed sets $V(\fa)$).

Since the sets $p^{-1}(U_j)$ are the preimages under $p$ of the open sets $U_j$, $p$ is injective on all of $X_y$ (two points with the same image lie in the same $p^{-1}(U_j)$), has image $f^{-1}(y)$, and is a homeomorphism from each open set $p^{-1}(U_j)$ onto the open subset $f^{-1}(y)\cap U_j$ of $f^{-1}(y)$. These open subsets cover $f^{-1}(y)$, so $p$ is a homeomorphism of $X_y$ onto $f^{-1}(y)$.

(b) The ring map $k[s]\to k[s,t]/(s-t^2)$ identifies the target with $k[t]$, $s\mapsto t^2$; so $X=\Spec k[t]$ and $f$ is the map $t\mapsto t^2$.

\emph{Caveat.} For the statement ``two points'' we need $\operatorname{char}k\ne2$, which the printed exercise does not state (it only assumes $k$ algebraically closed); I assume it. (In characteristic $2$ the fibre over $a\ne0$ is the one-point nonreduced scheme $\Spec k[t]/(t-\alpha)^2$, $\alpha^2=a$.)

Let $y=a\ne0$, i.e.\ the maximal ideal $(s-a)$, with $k(y)=k[s]/(s-a)=k$. Then
\[X_y=\Spec\big(k[t]\otimes_{k[s]}k[s]/(s-a)\big)=\Spec k[t]/(t^2-a)=\Spec k[t]/((t-\alpha)(t+\alpha)),\]
where $\alpha^2=a$, $\alpha\ne0$, so $\alpha\ne-\alpha$. By the Chinese remainder theorem, $k[t]/(t^2-a)\cong k\times k$, so $X_y$ consists of two points (reduced), each with residue field $k$.

For $y$ corresponding to $0$, $X_y=\Spec k[t]/(t^2)$, a single point (the prime $(t)$) with $\mathcal O=k[t]/(t^2)$ nonreduced ($t\neq0$, $t^2=0$).

For the generic point $\eta=(0)$ of $Y$, $k(\eta)=k(s)$ and
\[X_\eta=\Spec\big(k[t]\otimes_{k[s]}k(s)\big)=\Spec\, k(s)[t]/(t^2-s).\]
The polynomial $t^2-s$ is irreducible over $k(s)$: otherwise $s=(u/v)^2$ with $u,v\in k[s]$, $v\ne0$, i.e.\ $sv^2=u^2$, and the left side has odd degree while the right side has even degree. Hence $k(s)[t]/(t^2-s)$ is a field, of degree $2$ over $k(s)$, and $X_\eta$ is the one-point scheme $\Spec$ of this field.\hfill$\square$

\exercise{3.11}{\emph{Closed Subschemes.}
\begin{enumerate}[label=(\alph*),leftmargin=*]
\item Closed immersions are stable under base extension: if $f:Y\to X$ is a closed immersion, and if $X'\to X$ is any morphism, then $f':Y\times_XX'\to X'$ is also a closed immersion.
\item If $Y$ is a closed subscheme of an affine scheme $X=\Spec A$, then $Y$ is also affine, and in fact $Y$ is the closed subscheme determined by a suitable ideal $\fa\subseteq A$ as the image of the closed immersion $\Spec A/\fa\to\Spec A$. [\emph{Hints}: First show that $Y$ can be covered by a finite number of open affine subsets of the form $D(f_i)\cap Y$, with $f_i\in A$. By adding some more $f_i$ with $D(f_i)\cap Y=\emptyset$, if necessary, show that we may assume that the $D(f_i)$ cover $X$. Next show that $f_1,\dots,f_r$ generate the unit ideal of $A$. Then use (Ex.~2.17b) to show that $Y$ is affine, and (Ex.~2.18d) to show that $Y$ comes from an ideal $\fa\subseteq A$.] \emph{Note}: We will give another proof of this result using sheaves of ideals later (5.10).
\item Let $Y$ be a closed subset of a scheme $X$, and give $Y$ the reduced induced subscheme structure. If $Y'$ is any other closed subscheme of $X$ with the same underlying topological space, show that the closed immersion $Y\to X$ factors through $Y'$. We express this property by saying that the reduced induced structure is the smallest subscheme structure on a closed subset.
\item Let $f:Z\to X$ be a morphism. Then there is a unique closed subscheme $Y$ of $X$ with the following property: the morphism $f$ factors through $Y$, and if $Y'$ is any other closed subscheme of $X$ through which $f$ factors, then $Y\to X$ factors through $Y'$ also. We call $Y$ the \emph{scheme-theoretic image} of $f$. If $Z$ is a reduced scheme, then $Y$ is just the reduced induced structure on the closure of the image $f(Z)$.
\end{enumerate}
(Part (b) is starred in the book.)}
\pf
I prove (b) first, since (a) uses it (the proof of (b) does not use (a)). Throughout, a closed subscheme is an equivalence class of closed immersions $i:Y\to X$, and I freely identify $Y$ with the closed subset $i(Y)$.

(b) Let $i:Y\to X=\Spec A$ be a closed immersion; put $\varphi=\alpha(i):A\to C:=\Gamma(Y,\mathcal O_Y)$ (Ex.~2.4).

\emph{Step 1: $Y$ is covered by finitely many $Y_{f_i}$ that are affine and $D(f_i)$ cover $X$.} Let $y\in Y$ and let $\Spec C'\subseteq Y$ be an open affine neighbourhood of $y$. Since $i$ is a homeomorphism onto the closed set $i(Y)$, there is an open $O\subseteq X$ with $i(\Spec C')=i(Y)\cap O$; choose $f\in A$ with $i(y)\in D(f)\subseteq O$. Then $i^{-1}(D(f))\subseteq i^{-1}(O)=\Spec C'$, and $i^{-1}(D(f))$ is the open subset $D_{\Spec C'}(\bar f)$ where $\bar f=i^\#(f)|_{\Spec C'}$ (the point $z$ lies in it iff $f\notin\fp_{i(z)}$ iff the image of $f$ in $\mathcal O_{Y,z}$ is a unit, by locality of $i^\#_z$). So $i^{-1}(D(f))=\Spec C'_{\bar f}$ is affine. As $Y$ is homeomorphic to a closed subset of the quasi-compact space $X$, it is quasi-compact, so finitely many of these cover $Y$: $Y=\bigcup_{l\le n}i^{-1}D(f_l)$. The closed set $X\smallsetminus\bigcup_{l\le n}D(f_l)$ is quasi-compact and disjoint from $i(Y)$, so it is covered by finitely many $D(g_m)$ with $D(g_m)\cap i(Y)=\emptyset$ (such sets form a base of the open set $X\smallsetminus i(Y)$). Thus, renaming, we have $f_1,\dots,f_r\in A$ with $\bigcup D(f_j)=X$ (i.e.\ $(f_1,\dots,f_r)=A$) and each $i^{-1}(D(f_j))$ affine (possibly empty).

\emph{Step 2: $Y$ is affine.} The elements $s_j=\varphi(f_j)\in C$ generate the unit ideal of $C$, and $Y_{s_j}=i^{-1}(D(f_j))$ is affine. By Ex.~2.17(b), $Y$ is affine, $Y=\Spec C$, and $i=\Spec\varphi$ (Ex.~2.4).

\emph{Step 3: $\varphi$ is surjective.} Let $\fp\in\Spec A$. The stalk of $i_*\mathcal O_Y$ at $\fp$ is $\varinjlim_{g\notin\fp}\mathcal O_Y(i^{-1}D(g))=\varinjlim_{g\notin\fp}C_{\varphi(g)}=C_\fp:=C\otimes_AA_\fp$. Since $i^\#:\mathcal O_X\to i_*\mathcal O_Y$ is a surjective map of sheaves, it is surjective on stalks, i.e.\ $A_\fp\to C_\fp$ is surjective for every prime $\fp$. Thus the $A$-module $M=C/\varphi(A)$ satisfies $M_\fp=0$ for all $\fp$, so $M=0$. Hence $\varphi$ is surjective, $C\cong A/\fa$ with $\fa=\Ker\varphi$, and $i$ is the closed immersion $\Spec A/\fa\to\Spec A$.

(a) Let $f:Y\to X$ be a closed immersion and $g:X'\to X$ any morphism, $f':Y'=Y\times_XX'\to X'$. Being a closed immersion is local on the target $X'$ (the image of $f'$ is closed iff its trace on each member of an open cover is closed, $f'$ is injective and a homeomorphism onto its image iff it is so over each member, surjectivity of $\mathcal O_{X'}\to f'_*\mathcal O_{Y'}$ can be checked on stalks). Cover $X'$ by open affines $W=\Spec A'$ each mapped by $g$ into an open affine $U=\Spec A\subseteq X$ (cover $X$ by affines and then each $g^{-1}(U)$ by affines). By (b), applied to the closed immersion $f^{-1}(U)\to U$, we have $f^{-1}(U)=\Spec A/\fa$ for an ideal $\fa$. As $W\to X$ factors through $U$, $f'^{-1}(W)=Y\times_XW=f^{-1}(U)\times_UW=\Spec(A/\fa\otimes_AA')=\Spec A'/\fa A'$, and $f'|:\Spec A'/\fa A'\to\Spec A'$ is induced by the surjection $A'\to A'/\fa A'$, a closed immersion (Example~3.2.3). Hence $f'$ is a closed immersion.

(c) Let $i':Y'\to X$ be a closed subscheme with $i'(Y')=Y$ as sets. Let $\mathcal J'=\Ker(\mathcal O_X\to i'_*\mathcal O_{Y'})$ and $\mathcal J=\Ker(\mathcal O_X\to i_*\mathcal O_Y)$ be the ideal sheaves; since $i,i'$ are homeomorphisms onto $Y$ and the maps $\mathcal O_X\to i_*\mathcal O$ are surjective, we may identify $\mathcal O_{Y'}=\mathcal O_X/\mathcal J'|_Y$ and $\mathcal O_Y=\mathcal O_X/\mathcal J|_Y$ as sheaves on the space $Y$. Let $U=\Spec A$ be an open affine of $X$. By (b), $\mathcal J'(U)=\fa'$ where $Y'\cap U=\Spec A/\fa'$, so $V(\fa')=Y\cap U$. By definition of the reduced induced structure (Example~3.2.6), $\mathcal J(U)=\fb:=\bigcap_{\fp\in Y\cap U}\fp$. Then $\fa'\subseteq\sqrt{\fa'}=\bigcap_{\fp\supseteq\fa'}\fp=\fb$. Thus $\mathcal J'(U)\subseteq\mathcal J(U)$ on all open affines, compatibly with restrictions; as the open affines form a base, $\mathcal J'\subseteq\mathcal J$. So the quotient maps give a surjective morphism of sheaves of rings $g^\#:\mathcal O_X/\mathcal J'\to\mathcal O_X/\mathcal J$ on $Y$. On stalks it is a surjection of local rings, hence a local homomorphism. So $(\mathrm{id}_Y,g^\#)$ is a morphism of locally ringed spaces, i.e.\ of schemes, $g:Y\to Y'$, and $i'\circ g=i$ (it is the identity on spaces and the composite $\mathcal O_X\to\mathcal O_X/\mathcal J'\to\mathcal O_X/\mathcal J$ is the quotient map).

(d) \emph{This part uses results of Sect.~5 (closed subschemes $\leftrightarrow$ quasi-coherent ideal sheaves, Cor.~5.10, and Prop.~5.7), as the note in the exercise indicates.} Recall that closed subschemes of $X$ correspond bijectively to quasi-coherent ideal sheaves $\mathcal I\subseteq\mathcal O_X$, via $\mathcal I\mapsto Y_{\mathcal I}=(V(\mathcal I),\mathcal O_X/\mathcal I)$ with $V(\mathcal I)=\Supp(\mathcal O_X/\mathcal I)$. If $\mathcal I'\subseteq\mathcal I$, the quotient map $\mathcal O_X/\mathcal I'\to\mathcal O_X/\mathcal I$ gives a factorization of $Y_{\mathcal I}\to X$ through $Y_{\mathcal I'}$ (as in (c)); conversely, a factorization $g:Y_{\mathcal I}\to Y_{\mathcal I'}$ gives $\mathcal O_X\to\mathcal O_X/\mathcal I'\to g_*\mathcal O_{Y_{\mathcal I}}$ equal to the quotient map by $\mathcal I$, which kills $\mathcal I'$; so $\mathcal I'\subseteq\mathcal I$.

\emph{Claim: $f:Z\to X$ factors through $Y_{\mathcal I}$ iff $\mathcal I\subseteq\mathcal K:=\Ker(f^\#:\mathcal O_X\to f_*\mathcal O_Z)$.} ($\Rightarrow$) If $f=i\circ h$ then $f^\#$ factors through $i^\#:\mathcal O_X\to i_*\mathcal O_{Y}$ whose kernel is $\mathcal I$. ($\Leftarrow$) Suppose $\mathcal I\subseteq\mathcal K$. For $z\in Z$ put $x=f(z)$. The local homomorphism $f^\#_z:\mathcal O_{X,x}\to\mathcal O_{Z,z}$ kills $\mathcal I_x$. If $\mathcal I_x=\mathcal O_{X,x}$ it would kill $1$, impossible; so $\mathcal I_x\ne\mathcal O_{X,x}$, i.e.\ $x\in\Supp(\mathcal O_X/\mathcal I)=V(\mathcal I)$. So $f(Z)\subseteq V(\mathcal I)$, defining a continuous $h:Z\to V(\mathcal I)=:Y$ with $i\circ h=f$ topologically. The map $f^\#:\mathcal O_X\to f_*\mathcal O_Z=i_*h_*\mathcal O_Z$ kills $\mathcal I$, so it factors through $\mathcal O_X/\mathcal I=i_*\mathcal O_Y$ and by adjunction gives $h^\#:\mathcal O_Y\to h_*\mathcal O_Z$ whose stalk maps are induced by the local maps $f^\#_z$, hence local. So $f=i\circ h$.

\emph{Construction.} Let $\mathcal I_f\subseteq\mathcal O_X$ be the sum of all quasi-coherent ideal sheaves contained in $\mathcal K$. It is the image of the direct sum of these sheaves, hence quasi-coherent (direct sums and images of quasi-coherent sheaves are quasi-coherent, Prop.~5.7), and obviously the largest quasi-coherent ideal contained in $\mathcal K$ (if $f$ is quasi-compact then $\mathcal K$ itself is quasi-coherent, but this is not needed). Let $Y=Y_{\mathcal I_f}$. By the claim $f$ factors through $Y$. If $f$ factors through a closed subscheme $Y'=Y_{\mathcal I'}$ then $\mathcal I'\subseteq\mathcal K$ is quasi-coherent, so $\mathcal I'\subseteq\mathcal I_f$ and $Y\to X$ factors through $Y'$. Uniqueness: if $Y_1$ has the same universal property then $Y\to X$ factors through $Y_1$ and $Y_1\to X$ through $Y$, so their ideals are comparable in both directions and equal.

\emph{The reduced case.} Let $Z$ be reduced, $W=\overline{f(Z)}$ with reduced induced structure, with ideal sheaf $\mathcal I_W$ (on an affine open $U=\Spec A$ it is $\bigcap_{\fp\in W\cap U}\fp$; it is quasi-coherent). If $s\in\mathcal I_W(U)$, then $s$ vanishes at every point of $W\cap U$, so $f^\#(s)\in\mathcal O_Z(f^{-1}U)$ vanishes at every point of $f^{-1}U$ (the value of $f^\#s$ at $z$ is the image of the value $s(f(z))=0$ in $k(f(z))\to k(z)$); on an affine open $\Spec R\subseteq f^{-1}U$ with $R$ reduced such an element lies in every prime, hence in the nilradical $=0$. So $f^\#(s)=0$ and $\mathcal I_W\subseteq\mathcal K$, hence $\mathcal I_W\subseteq\mathcal I_f$. In particular $V(\mathcal I_f)\subseteq W$. Conversely $V(\mathcal I_f)$ is closed and contains $f(Z)$ (the claim), hence contains $W$. So $V(\mathcal I_f)=W$ as sets. Finally any section $s$ of $\mathcal I_f$ has germ in $\mathfrak m_x$ at all $x\in W$ (as $(\mathcal I_f)_x\ne\mathcal O_x$ there), i.e.\ vanishes on $W$, so $\mathcal I_f\subseteq\mathcal I_W$. Hence $\mathcal I_f=\mathcal I_W$ and $Y=W$ with the reduced induced structure.\hfill$\square$

\exercise{3.12}{\emph{Closed Subschemes of $\operatorname{Proj}S$.}
\begin{enumerate}[label=(\alph*),leftmargin=*]
\item Let $\varphi:S\to T$ be a surjective homomorphism of graded rings, preserving degrees. Show that the open set $U$ of (Ex.~2.14) is equal to $\operatorname{Proj}T$, and the morphism $f:\operatorname{Proj}T\to\operatorname{Proj}S$ is a closed immersion.
\item If $I\subseteq S$ is a homogeneous ideal, take $T=S/I$ and let $Y$ be the closed subscheme of $X=\operatorname{Proj}S$ defined as the image of the closed immersion $\operatorname{Proj}S/I\to X$. Show that different homogeneous ideals can give rise to the same closed subscheme. For example, let $d_0$ be an integer, and let $I'=\bigoplus_{d\ge d_0}I_d$. Show that $I$ and $I'$ determine the same closed subscheme.
\end{enumerate}
We will see later (5.16) that every closed subscheme of $X$ comes from a homogeneous ideal $I\subseteq S$ (at least in the case where $S$ is a polynomial ring over $S_0$).}
\pf
Recall (Ex.~2.14(b)) $U=\{\fp\in\operatorname{Proj}T:\fp\not\supseteq\varphi(S_+)\}$ and $f:U\to\operatorname{Proj}S$, $\fp\mapsto\varphi^{-1}(\fp)$; for a homogeneous $s\in S_+$ we have $f^{-1}(D_+(s))=D_+(\varphi(s))$, and $f$ restricted to it is the morphism $\Spec T_{(\varphi(s))}\to\Spec S_{(s)}$ induced by $S_{(s)}\to T_{(\varphi(s))}$.

(a) As $\varphi$ is surjective and degree preserving, $\varphi(S_d)=T_d$ for all $d$ (apply $\varphi$ to the homogeneous components of a preimage), so $\varphi(S_+)=T_+$. A point of $\operatorname{Proj}T$ is a homogeneous prime not containing $T_+=\varphi(S_+)$; so $U=\operatorname{Proj}T$. The sets $D_+(s)$, $s\in S_+$ homogeneous, cover $\operatorname{Proj}S$, and $f^{-1}(D_+(s))=D_+(\varphi(s))\to D_+(s)$ is induced by $S_{(s)}\to T_{(\varphi(s))}$. This ring map is surjective: an element of $T_{(\varphi(s))}$ is $t/\varphi(s)^n$ with $t\in T_{n\deg s}$, and $t=\varphi(s')$ with $s'\in S_{n\deg s}$, so it is the image of $s'/s^n$. By Example~3.2.3 the map over each $D_+(s)$ is a closed immersion, and being a closed immersion is local on the target, so $f$ is a closed immersion.

(b) Let $d_0\in\mathbf Z$ and $I'=\bigoplus_{d\ge d_0}I_d$; it is a homogeneous ideal contained in $I$ (if $x\in I'_d$, $y\in S_e$, then $xy\in I_{d+e}$ with $d+e\ge d_0$ when $e\ge0$). Let $\psi:S/I'\to S/I$ be the natural surjection of graded rings; by (a) it gives a closed immersion $\rho:\operatorname{Proj}S/I\to\operatorname{Proj}S/I'$, and the functoriality of the construction of Ex.~2.14 shows that the composite of $\rho$ with the closed immersion $\operatorname{Proj}S/I'\to\operatorname{Proj}S$ is the closed immersion $\operatorname{Proj}S/I\to\operatorname{Proj}S$ defined by $S\to S/I$. So it suffices to show that $\rho$ is an isomorphism; then the two closed immersions differ by an isomorphism, i.e.\ $I$ and $I'$ define the same closed subscheme. As in (a), over $D_+(\bar s)\subseteq\operatorname{Proj}S/I'$ ($s\in S_+$ homogeneous of degree $e>0$, $\bar s$ its class) the morphism $\rho$ is induced by $\psi_{(s)}:(S/I')_{(\bar s)}\to(S/I)_{(\bar s)}$, which is surjective. It is injective: let $x/s^n$ ($x\in S_{ne}$) map to $0$ in $(S/I)_{(\bar s)}$. Then $s^mx\in I$ for some $m\ge0$, and enlarging $m$ we may assume $(n+m)e\ge d_0$. Then $s^mx\in I_{(n+m)e}=I'_{(n+m)e}$, so $\bar s^m\bar x=0$ in $S/I'$ and $x/s^n=0$ in $(S/I')_{(\bar s)}$. So $\psi_{(s)}$ is an isomorphism for all such $s$; these $D_+(\bar s)$ cover $\operatorname{Proj}S/I'$, so $\rho$ is an isomorphism.

(The point of the example: $I\ne I'$ in general, e.g.\ $S=k[x_0,x_1]$, $I=(x_0)$, $d_0=2$: $I'=(x_0^2,x_0x_1)\ne I$, since $x_0\in I\smallsetminus I'$; yet both define the same closed subscheme, the point $x_0=0$ of $\mathbf P^1_k$.)\hfill$\square$

\exercise{3.13}{\emph{Properties of Morphisms of Finite Type.}
\begin{enumerate}[label=(\alph*),leftmargin=*]
\item A closed immersion is a morphism of finite type.
\item A quasi-compact open immersion (Ex.~3.2) is of finite type.
\item A composition of two morphisms of finite type is of finite type.
\item Morphisms of finite type are stable under base extension.
\item If $X$ and $Y$ are schemes of finite type over $S$, then $X\times_SY$ is of finite type over $S$.
\item If $X\xrightarrow{f}Y\xrightarrow{g}Z$ are two morphisms, and if $f$ is quasi-compact, and $g\circ f$ is of finite type, then $f$ is of finite type.
\item If $f:X\to Y$ is a morphism of finite type, and if $Y$ is noetherian, then $X$ is noetherian.
\end{enumerate}}
\pf
(a) Let $f:Y\to X$ be a closed immersion and $U=\Spec A\subseteq X$ an open affine. By Ex.~3.11(b), $f^{-1}(U)\to U$ is $\Spec A/\fa\to\Spec A$, so $f^{-1}(U)$ is a single affine open with ring $A/\fa$, a finitely generated $A$-algebra. Covering $X$ by such $U$ shows $f$ is of finite type.

(b) Let $j:U\to X$ be a quasi-compact open immersion and $V=\Spec A\subseteq X$ open affine. By Ex.~3.2, $j^{-1}(V)=U\cap V$ is a quasi-compact open subset of $\Spec A$, so it is a finite union of basic open sets $D(g_l)=\Spec A_{g_l}$, and $A_{g_l}$ is a finitely generated $A$-algebra. Hence $j$ is of finite type.

(c) Let $f:X\to Y$ and $g:Y\to Z$ be of finite type, $W=\Spec C\subseteq Z$ open affine. By Ex.~3.3(b), $g^{-1}(W)=V_1\cup\dots\cup V_n$ with $V_i=\Spec B_i$, $B_i$ finitely generated $C$-algebras, and each $f^{-1}(V_i)=\bigcup_{j\le m_i}\Spec A_{ij}$ with $A_{ij}$ finitely generated $B_i$-algebras. Then $(gf)^{-1}(W)=\bigcup_{i,j}\Spec A_{ij}$ is a finite union and each $A_{ij}$ is a finitely generated $C$-algebra. By Ex.~3.3(b) (applied to the cover of $Z$ by all open affines), $gf$ is of finite type.

(d) Let $f:X\to Y$ be of finite type and $Y'\to Y$ any morphism; $X'=X\times_YY'$, $f':X'\to Y'$. Cover $Y'$ by open affines $W=\Spec C'$, each mapped into an open affine $V=\Spec B\subseteq Y$. By Ex.~3.3(b), $f^{-1}(V)=\bigcup_{i\le n}\Spec A_i$ with $A_i$ finitely generated $B$-algebras. As in Ex.~3.10(a), $f'^{-1}(W)=X\times_YW=f^{-1}(V)\times_VW$, which by Thm.~3.3 (Step~4) is covered by the $\Spec A_i\times_V W=\Spec(A_i\otimes_BC')$, $i\le n$; each $A_i\otimes_BC'$ is a finitely generated $C'$-algebra. So $f'$ is of finite type.

(e) $X\times_SY\to Y$ is the base extension of $X\to S$ by $Y\to S$, hence of finite type by (d); the composition with $Y\to S$ is of finite type by (c).

(f) We show $f$ is locally of finite type; then $f$ is of finite type by Ex.~3.3(a). Let $W=\Spec C\subseteq Z$ and $V=\Spec B\subseteq g^{-1}(W)$ be open affines; such $V$ cover $Y$. Let $U=\Spec A$ be an open affine subset of $f^{-1}(V)$; such $U$ cover $f^{-1}(V)$. Since $U\subseteq(gf)^{-1}(W)$ and $gf$ is of finite type, $A$ is a finitely generated $C$-algebra by Ex.~3.3(c). The $C$-algebra structure of $A$ factors as $C\to B\to A$, so $A$ is a finitely generated $B$-algebra (the same generators). Hence $f^{-1}(V)$ is covered by spectra of finitely generated $B$-algebras, and $f$ is locally of finite type.

(g) Cover $Y$ by finitely many $V_i=\Spec B_i$, $B_i$ noetherian. By Ex.~3.3(b), each $f^{-1}(V_i)$ is a finite union of $\Spec A_{ij}$ with $A_{ij}$ a finitely generated $B_i$-algebra, hence noetherian by the Hilbert basis theorem. So $X$ is a finite union of spectra of noetherian rings, i.e.\ noetherian.\hfill$\square$

\exercise{3.14}{If $X$ is a scheme of finite type over a field $k$, show that the closed points of $X$ are dense. Show by example that this is not true for arbitrary schemes.}
\pf
Let $U\subseteq X$ be a nonempty open subset; we show that $U$ contains a closed point of $X$. This is equivalent to the density of the set of closed points. Choose a nonempty open affine $V=\Spec A\subseteq U$; $A$ is a finitely generated $k$-algebra by Ex.~3.3(c) (applied to $X\to\Spec k$), and $A\ne0$. Let $\fm$ be a maximal ideal of $A$ and $x\in V$ the corresponding point. The field $k(x)=A/\fm$ is finitely generated as a $k$-algebra, hence finite over $k$ (Zariski's lemma / Nullstellensatz).

\emph{$x$ is closed in every open affine $W=\Spec A'$ containing it.} Let $\fq\subset A'$ be the prime of $x$. Then $A'/\fq$ is a domain which embeds in $\operatorname{Frac}(A'/\fq)=k(x)$, so it is a domain finite-dimensional over $k$, hence a field. So $\fq$ is maximal, and $x$ is closed in $W$.

\emph{$x$ is closed in $X$.} Let $z\in\overline{\{x\}}$ and choose an open affine $W\ni z$. Since $W$ is an open neighbourhood of $z\in\overline{\{x\}}$, it meets $\{x\}$, so $x\in W$. Then $z\in\overline{\{x\}}\cap W=\overline{\{x\}}^W=\{x\}$ by the previous paragraph. So $\overline{\{x\}}=\{x\}$.

Thus every nonempty open set contains a closed point of $X$.

\emph{Counterexample.} Let $R$ be a discrete valuation ring and $X=\Spec R$. It has two points, the closed point $\fm$ and the generic point $(0)$; the set of closed points is $\{\fm\}$, whose closure is $\{\fm\}\ne X$. (Of course $X$ is not of finite type over a field.)\hfill$\square$

\exercise{3.15}{Let $X$ be a scheme of finite type over a field $k$ (not necessarily algebraically closed).
\begin{enumerate}[label=(\alph*),leftmargin=*]
\item Show that the following three conditions are equivalent (in which case we say that $X$ is \emph{geometrically irreducible}): (i) $X\times_k\bar k$ is irreducible, where $\bar k$ denotes the algebraic closure of $k$. (By abuse of notation, we write $X\times_k\bar k$ to denote $X\times_{\Spec k}\Spec\bar k$.) (ii) $X\times_kk_s$ is irreducible, where $k_s$ denotes the separable closure of $k$. (iii) $X\times_kK$ is irreducible for every extension field $K$ of $k$.
\item Show that the following three conditions are equivalent (in which case we say that $X$ is \emph{geometrically reduced}): (i) $X\times_k\bar k$ is reduced. (ii) $X\times_kk_p$ is reduced, where $k_p$ denotes the perfect closure of $k$. (iii) $X\times_kK$ is reduced for all extension fields $K$ of $k$.
\item We say that $X$ is \emph{geometrically integral} if $X\times_k\bar k$ is integral. Give examples of integral schemes which are neither geometrically irreducible nor geometrically reduced.
\end{enumerate}}
\pf
Throughout, for a field extension $k\subseteq K$ and a $k$-scheme $X$ write $X_K=X\times_kK$ and $p_K:X_K\to X$ for the projection; over an affine open $\Spec A\subseteq X$, $p_K^{-1}(\Spec A)=\Spec(A\otimes_kK)$ (Thm.~3.3, Steps~1 and~4). I use the following facts.

\begin{itemize}
\item \textbf{(F1)} \emph{$p_K$ is surjective.} The fibre of $p_K$ over a point $x$ with residue field $k(x)$ is $\Spec(k(x)\otimes_kK)\ne\emptyset$ (a tensor product of two nonzero vector spaces over a field is nonzero; cf.\ Ex.~3.10). Consequently, if $X_K$ is irreducible then so is $X$ (a continuous image of an irreducible space is irreducible).
\item \textbf{(F2)} \emph{(Tensor products of domains.)} If $F$ is an algebraically closed field and $D,E$ are $F$-algebras that are integral domains, then $D\otimes_FE$ is an integral domain (Zariski--Samuel, or Lang, \emph{Algebra}; for finitely generated domains this is (I, Ex.~3.15), and the general case follows by taking the direct limit over finitely generated subalgebras).
\item \textbf{(F3)} \emph{(Base change and reducedness.)} If $A$ is a reduced algebra over a field $F$ and $L/F$ is a separable field extension (e.g.\ any extension if $F$ is perfect, or any algebraic extension of a perfect field), then $A\otimes_FL$ is reduced (Bourbaki, \emph{Alg.}~V, Sect.~15; Matsumura, \emph{Commutative Algebra}, Sect.~26).
\item \textbf{(F4)} \emph{Purely inseparable base change.} If $F\subseteq F'$ is algebraic and purely inseparable (every element of $F'$ has some power $\lambda^{q}\in F$, $q$ a power of $p=\operatorname{char}F$; in characteristic $0$ this means $F=F'$) then for every $F$-scheme $T$, $p:T_{F'}\to T$ is a homeomorphism. Indeed, on an affine piece $A\to B=A\otimes_FF'$ is injective (flatness) and integral, so $\Spec B\to\Spec A$ is surjective and closed (lying over). It is injective: for $b=\sum a_i\otimes\lambda_i$, taking $q$ large, $b^q=\sum a_i^q\lambda_i^q\otimes1\in A$ (we are in characteristic $p$), so a prime $\mathfrak Q$ of $B$ is determined by its contraction $\mathfrak P=\mathfrak Q\cap A$, namely $\mathfrak Q=\{b:b^q\in\mathfrak P\text{ for some }q\}$. A continuous closed bijection is a homeomorphism; a global statement follows because the property is local on $T$ (the preimages $p^{-1}(\Spec A)$ are saturated, and closedness is local on the target).
\end{itemize}

(a) (iii)$\Rightarrow$(i),(ii) is trivial. (i)$\Rightarrow$(ii): $X_{\bar k}\to X_{k_s}$ is surjective by (F1) (applied to $X_{k_s}$ over $k_s$), so $X_{k_s}$ is irreducible. (ii)$\Rightarrow$(i): $X_{\bar k}=(X_{k_s})_{\bar k}$ and $\bar k/k_s$ is purely inseparable (for $\alpha\in\bar k$, some $\alpha^{p^e}$ is separable over $k$, hence lies in $k_s$). By (F4), $X_{\bar k}$ is homeomorphic to $X_{k_s}$, so it is irreducible.

(i)$\Rightarrow$(iii): Let $K/k$ be an extension. The ring $K\otimes_k\bar k\ne0$; let $\fm$ be a maximal ideal and $L=(K\otimes_k\bar k)/\fm$. Then $L$ is a field containing (copies of) $K$ and $\bar k$ over $k$, so $X_L=(X_K)_L=(X_{\bar k})_L$. Put $F=\bar k$ and $Z=X_{\bar k}$, an irreducible scheme (by (i)) with a finite cover by open affines $\Spec D_i$. Consider $\Spec(D_i\otimes_FL)$. The nilradical $N_i$ of $D_i$ gives $N_i\otimes_FL\subseteq D_i\otimes_FL$ (flatness), an ideal consisting of nilpotent elements (a sum of nilpotents $n\otimes\lambda$), and $(D_i\otimes L)/(N_i\otimes L)=(D_i/N_i)\otimes_FL$. Hence $\Spec((D_i/N_i)\otimes_FL)$ and $\Spec(D_i\otimes_FL)$ have the same underlying space. Now $Z$ is irreducible, so $Z_{\mathrm{red}}$ is integral (Prop.~3.1) and its affine pieces $E_i=D_i/N_i$ are domains (nonempty $\Spec D_i$); by (F2), $E_i\otimes_FL$ is a domain, so $W_i=p_L^{-1}(\Spec D_i)\subseteq Z_L$ is irreducible. For $i,j$, $W_i\cap W_j=p_L^{-1}(\Spec D_i\cap\Spec D_j)$ is nonempty since $Z$ is irreducible (so the two nonempty opens meet) and $p_L$ is surjective (F1). A space covered by irreducible open subsets that pairwise meet is irreducible. So $X_L$ is irreducible, and then $X_K$ is irreducible by (F1) applied to $(X_K)_L\to X_K$.

(b) (iii)$\Rightarrow$(i),(ii) is trivial. (i)$\Rightarrow$(ii): over an affine open $\Spec A$, $A\otimes_kk_p\subseteq A\otimes_k\bar k$ (flatness over the field $k_p$) and a subring of a reduced ring is reduced. (ii)$\Rightarrow$(i): $A\otimes_k\bar k=(A\otimes_kk_p)\otimes_{k_p}\bar k$, where $A\otimes_kk_p$ is reduced and $\bar k/k_p$ is algebraic over the perfect field $k_p$, hence separable; by (F3) $A\otimes_k\bar k$ is reduced. (i)$\Rightarrow$(iii): let $K/k$ be an extension and $L\supseteq K,\bar k$ as in (a). By (i), $A\otimes_k\bar k$ is reduced; as $\bar k$ is perfect, (F3) shows $(A\otimes_k\bar k)\otimes_{\bar k}L=A\otimes_kL$ is reduced. Finally $A\otimes_kK\to(A\otimes_kK)\otimes_KL=A\otimes_kL$ is injective ($L$ is free over $K$), so $A\otimes_kK$ is reduced. As this holds on every affine open, $X_K$ is reduced.

(c) The following integral schemes are not geometrically integral for different reasons (the first two show that each of the two properties can fail separately; the third fails both at the same time).

\emph{(1) Not geometrically irreducible.} $k=\R$, $X=\Spec\C=\Spec\R[x]/(x^2+1)$. It is a single reduced point, hence integral. But $\C\otimes_\R\C\cong\C\times\C$ ($a\otimes b\mapsto(ab,a\bar b)$), so $X_\C$ consists of two points: $X$ is not geometrically irreducible (though it is geometrically reduced).

\emph{(2) Not geometrically reduced.} $k=\mathbf F_p(t)$, $X=\Spec k[x]/(x^p-t)$. The polynomial $x^p-t$ is irreducible, so $X=\Spec K$ with $K=k(t^{1/p})$ a field, an integral scheme. But $K\otimes_kk'=k'[x]/(x^p-t)=k'[x]/(x-t^{1/p})^p$ for $k'=k(t^{1/p})$ (or $\bar k$), which is not reduced. (This $X$ is geometrically irreducible.)

\emph{(3) Neither.} Let $k=\mathbf F_2(t)$ and
\[K=k[x,y]/(x^2+x+t,\ y^2+t),\qquad X=\Spec K.\]
First, $x^2+x+t$ is irreducible over $k$: a root $u/v$ ($u,v\in\mathbf F_2[t]$ coprime) would satisfy $u^2+uv+tv^2=0$, so $v\mid u^2$, hence $v\in\mathbf F_2^\times$, i.e.\ $v=1$ and $u^2+u+t=0$ with $u\in\mathbf F_2[t]$; if $\deg u\ge1$ the term $u^2$ has degree $2\deg u>\max(\deg u,1)$ and cannot cancel, and if $u$ is constant then $u^2+u=0$ gives $t=0$, a contradiction. So $k_1=k[x]/(x^2+x+t)$ is a field with basis $1,x$ over $k$. Next $t$ is not a square in $k_1$: $(c+dx)^2=c^2+d^2x^2=(c^2+d^2t)+d^2x$ (char $2$, $x^2=x+t$), and this equals $t$ only if $d=0$ and $c^2=t$, which is impossible in $k=\mathbf F_2(t)$ (degree parity). So $y^2+t$ is irreducible over $k_1$, and $K=k_1[y]/(y^2+t)$ is a field, so $X$ is an integral scheme. Over $\bar k$:
\[K\otimes_k\bar k=\bar k[x,y]/(x^2+x+t,\,y^2+t)=\bar k[x,y]/\big((x+a)(x+a+1),(y+b)^2\big),\]
where $a^2+a+t=0$ (so $a\ne a+1$) and $b^2=t$. By the Chinese remainder theorem this is $\bar k[y]/(y+b)^2\times\bar k[y]/(y+b)^2$: two points, each nonreduced. So $X_{\bar k}$ is neither irreducible nor reduced: $X$ is an integral scheme that is neither geometrically irreducible nor geometrically reduced.\hfill$\square$

\exercise{3.16}{\emph{Noetherian Induction.} Let $X$ be a noetherian topological space, and let $\mathcal P$ be a property of closed subsets of $X$. Assume that for any closed subset $Y$ of $X$, if $\mathcal P$ holds for every proper closed subset of $Y$, then $\mathcal P$ holds for $Y$. (In particular, $\mathcal P$ must hold for the empty set.) Then $\mathcal P$ holds for $X$.}
\pf
Suppose $\mathcal P$ fails for $X$. Let $\mathcal S$ be the set of closed subsets of $X$ for which $\mathcal P$ fails; $X\in\mathcal S$, so $\mathcal S\ne\emptyset$. Since $X$ is noetherian, closed subsets satisfy the descending chain condition, so $\mathcal S$ has a minimal element $Y$ (otherwise we could build an infinite strictly descending chain in $\mathcal S$). By minimality every proper closed subset of $Y$ has $\mathcal P$, so by hypothesis $\mathcal P$ holds for $Y$, contradicting $Y\in\mathcal S$. Hence $\mathcal P$ holds for $X$.\hfill$\square$

\exercise{3.17}{\emph{Zariski Spaces.} A topological space $X$ is a \emph{Zariski space} if it is noetherian and every (nonempty) closed irreducible subset has a unique generic point (Ex.~2.9). For example, let $R$ be a discrete valuation ring, and let $T=\operatorname{sp}(\Spec R)$. Then $T$ consists of two points $t_0=$ the maximal ideal, $t_1=$ the zero ideal. The open subsets are $\emptyset,\{t_1\}$, and $T$. This is an irreducible Zariski space with generic point $t_1$.
\begin{enumerate}[label=(\alph*),leftmargin=*]
\item Show that if $X$ is a noetherian scheme, then $\operatorname{sp}(X)$ is a Zariski space.
\item Show that any minimal nonempty closed subset of a Zariski space consists of one point. We call these closed points.
\item Show that a Zariski space $X$ satisfies the axiom $T_0$: given any two distinct points of $X$, there is an open set containing one but not the other.
\item If $X$ is an irreducible Zariski space, then its generic point is contained in every nonempty open subset of $X$.
\item If $x_0,x_1$ are points of a topological space $X$, and if $x_0\in\overline{\{x_1\}}$, then we say that $x_1$ \emph{specializes} to $x_0$, written $x_1\rightsquigarrow x_0$. We also say $x_0$ is a \emph{specialization} of $x_1$, or that $x_1$ is a \emph{generization} of $x_0$. Now let $X$ be a Zariski space. Show that the minimal points, for the partial ordering determined by $x_1>x_0$ if $x_1\rightsquigarrow x_0$, are the closed points, and the maximal points are the generic points of the irreducible components of $X$. Show also that a closed subset contains every specialization of any of its points. (We say closed subsets are \emph{stable under specialization}.) Similarly, open subsets are \emph{stable under generization}.
\item Let $t$ be the functor on topological spaces introduced in the proof of (2.6). If $X$ is a noetherian topological space, show that $t(X)$ is a Zariski space. Furthermore $X$ itself is a Zariski space if and only if the map $\alpha:X\to t(X)$ is a homeomorphism.
\end{enumerate}}
\pf
(a) $\operatorname{sp}(X)$ is a noetherian topological space for a noetherian scheme $X$ (Caution~3.1.1, from (Ex.~2.13)), and every irreducible closed subset of a scheme has a unique generic point (Ex.~2.9).

(b) Let $Z$ be a minimal nonempty closed subset. If $Z=Z_1\cup Z_2$ with $Z_i$ closed and proper in $Z$, then by minimality $Z_1=Z_2=\emptyset$ and $Z=\emptyset$, a contradiction; so $Z$ is irreducible and has a generic point $z$: $Z=\overline{\{z\}}$. If $z'\in Z$, then $\overline{\{z'\}}\subseteq Z$ is a nonempty closed subset, so $\overline{\{z'\}}=Z$, and $z'$ is also a generic point of $Z$. By uniqueness $z'=z$. So $Z=\{z\}$. (Conversely if $\{z\}$ is closed it is clearly a minimal nonempty closed set, so ``closed points'' are the points $z$ with $\{z\}$ closed.)

(c) Let $x\ne y$ and suppose every open set containing $x$ contains $y$ and vice versa. Then $y\in\overline{\{x\}}$ (every open neighbourhood of $y$ meets $\{x\}$) and $x\in\overline{\{y\}}$, so $\overline{\{x\}}=\overline{\{y\}}$: this irreducible closed set has two generic points $x,y$, contradicting uniqueness.

(d) Let $\xi$ be the generic point and $U\ne\emptyset$ open. If $\xi\notin U$ then $\xi\in X\smallsetminus U$, which is closed, so $X=\overline{\{\xi\}}\subseteq X\smallsetminus U$ and $U=\emptyset$, a contradiction.

(e) Define $x_1\ge x_0$ iff $x_0\in\overline{\{x_1\}}$. This is reflexive and transitive ($x_0\in\overline{\{x_1\}}$, $x_1\in\overline{\{x_2\}}$ give $\overline{\{x_0\}}\subseteq\overline{\{x_1\}}\subseteq\overline{\{x_2\}}$), and antisymmetric: $x_1\ge x_0\ge x_1$ gives $\overline{\{x_0\}}=\overline{\{x_1\}}$, so $x_0=x_1$ by uniqueness of generic points. Write $x_1>x_0$ if $x_1\ge x_0$, $x_1\ne x_0$.

\emph{Minimal points.} $z$ is minimal iff there is no $x_0\ne z$ with $x_0\in\overline{\{z\}}$, i.e.\ iff $\overline{\{z\}}=\{z\}$, i.e.\ iff $z$ is a closed point (by (b)).

\emph{Maximal points.} Let $Z$ be an irreducible component of $X$ (closed and irreducible, so $Z=\overline{\{z\}}$) and suppose $z\in\overline{\{x_1\}}$. The closed irreducible set $\overline{\{x_1\}}$ contains $Z$, so equals $Z$ by maximality of components; so $x_1$ is a generic point of $Z$, $x_1=z$. Thus $z$ is maximal. Conversely let $z$ be maximal. $\overline{\{z\}}$ is irreducible, hence contained in an irreducible component $Z$ of $X$ (components exist and every irreducible subset lies in one, as $X$ is noetherian), with generic point $w$. Then $z\in\overline{\{w\}}$, i.e.\ $w\ge z$, so $w=z$ by maximality; thus $z$ is the generic point of the component $Z$.

\emph{Stability.} If $Z$ is closed, $x\in Z$ and $x_0\in\overline{\{x\}}$ then $x_0\in Z$ since $\overline{\{x\}}\subseteq Z$. If $U$ is open, $x_0\in U$ and $x_0\in\overline{\{x\}}$ (i.e.\ $x$ is a generization of $x_0$), then $x\in U$: otherwise $x\in X\smallsetminus U$ (closed), so $x_0\in\overline{\{x\}}\subseteq X\smallsetminus U$, a contradiction.

(f) Recall (proof of (2.6)): $t(X)$ is the set of nonempty irreducible closed subsets of $X$; for $Y\subseteq X$ closed, $t(Y)=\{Z\in t(X):Z\subseteq Y\}$; one has $t(Y_1\cup Y_2)=t(Y_1)\cup t(Y_2)$ (an irreducible closed subset of $Y_1\cup Y_2$ lies in one of them) and $t(\bigcap Y_i)=\bigcap t(Y_i)$, and the closed sets of $t(X)$ are the $t(Y)$. The map $Y\mapsto t(Y)$ from closed subsets of $X$ to closed subsets of $t(X)$ is a bijection preserving inclusions in both directions: $Y=\bigcup_{y\in Y}\overline{\{y\}}$ is the union of the members of $t(Y)$, so $Y$ can be recovered from $t(Y)$. It also preserves finite unions. Hence:
\begin{itemize}
\item $t(X)$ is noetherian (descending chains of closed sets correspond).
\item The irreducible closed subsets of $t(X)$ are exactly the $t(Y)$ with $Y$ irreducible closed in $X$ (irreducibility is expressed by inclusions and finite unions of closed sets, which the bijection preserves), i.e.\ the sets $t(Y)$ with $Y\in t(X)$.
\item For $Y\in t(X)$ (a point of $t(X)$), $\overline{\{Y\}}$ is the smallest closed set $t(Y')$ containing $Y$, i.e.\ with $Y\subseteq Y'$; this is $t(Y)$. So $t(Y)$ has the point $Y$ as a generic point, and it is unique: if $Y'$ is also a generic point, $t(Y')=\overline{\{Y'\}}=t(Y)$ forces $Y'=Y$.
\end{itemize}
So $t(X)$ is a Zariski space.

The map $\alpha:X\to t(X)$, $x\mapsto\overline{\{x\}}$ is continuous: $\alpha^{-1}(t(Y))=\{x:\overline{\{x\}}\subseteq Y\}=Y$. It is injective iff distinct points have distinct closures iff $X$ is $T_0$; it is surjective iff every nonempty irreducible closed subset has a generic point. If $X$ is a Zariski space, then $\alpha$ is injective by (c) and surjective, hence bijective; and then $\alpha(Y)=t(Y)$ for closed $Y$ (as $\alpha^{-1}(t(Y))=Y$ and $\alpha$ is onto), so $\alpha$ maps closed sets onto closed sets and is a homeomorphism. Conversely, if $X$ is noetherian and $\alpha$ is a homeomorphism, then $X\cong t(X)$ is a Zariski space (or directly: surjectivity gives existence of generic points and injectivity gives uniqueness).

Remark on the example: for the DVR $T$, $\overline{\{t_1\}}=T$ and $\overline{\{t_0\}}=\{t_0\}$, so $t_1>t_0$, consistent with (d), (e).\hfill$\square$

\exercise{3.18}{\emph{Constructible Sets.} Let $X$ be a Zariski space (Ex.~3.17). A subset of $X$ is \emph{constructible} if it belongs to the smallest family $\mathcal F$ of subsets such that (1) every open subset is in $\mathcal F$, (2) a finite intersection of elements of $\mathcal F$ is in $\mathcal F$, and (3) the complement of an element of $\mathcal F$ is in $\mathcal F$.
\begin{enumerate}[label=(\alph*),leftmargin=*]
\item A subset of $X$ is \emph{locally closed} if it is the intersection of an open subset with a closed subset. Show that a subset of $X$ is constructible if and only if it can be written as a finite disjoint union of locally closed subsets.
\item Show that a constructible subset of an irreducible Zariski space $X$ is dense if and only if it contains the generic point. Furthermore, in that case it contains a nonempty open subset.
\item A subset $S$ of $X$ is closed if and only if it is constructible and stable under specialization. Similarly, a subset $T$ of $X$ is open if and only if it is constructible and stable under generization.
\item If $f:X\to Y$ is a continuous map of Zariski spaces, then the inverse image of any constructible subset of $Y$ is a constructible subset of $X$.
\end{enumerate}}
\pf
Note that $\mathcal F$ is closed under finite unions (A$\cup$B$=$ complement of the intersection of the complements) and contains $X$ and $\emptyset$; if $B\subseteq X$ is any subset, a family of subsets of $B$ with the analogous properties relative to $B$ is considered below.

(a) Let $\mathcal G$ be the family of finite unions of locally closed sets. Closed and open sets are locally closed, so $\mathcal G$ contains all open sets. The intersection of two locally closed sets is locally closed, so $\mathcal G$ is closed under finite intersections (distributivity). The complement of a locally closed set $O\cap C$ ($O$ open, $C$ closed) is $(X\smallsetminus O)\cup(X\smallsetminus C)$, a union of a closed set and an open set; so the complement of a finite union of locally closed sets is a finite intersection of finite unions of locally closed sets, which lies in $\mathcal G$. Thus $\mathcal G$ satisfies (1)--(3), so $\mathcal F\subseteq\mathcal G$. Conversely a locally closed set $O\cap C=O\cap(X\smallsetminus(X\smallsetminus C))$ lies in $\mathcal F$, and $\mathcal F$ is closed under finite unions; so $\mathcal G\subseteq\mathcal F$. Hence $\mathcal F=\mathcal G$.

It remains to see that every finite union of locally closed sets is a finite \emph{disjoint} union of locally closed sets. For locally closed $D=O\cap C$ and $P$ locally closed, $P\smallsetminus D=(P\cap(X\smallsetminus O))\sqcup(P\cap O\cap(X\smallsetminus C))$ is a disjoint union of two locally closed sets. Hence if $P$ is a finite disjoint union of locally closed sets, so is $P\smallsetminus D$ for any locally closed $D$ (do it for each piece). Now let $A=D_1\sqcup\dots\sqcup D_m$ be a disjoint union of locally closed sets and $L$ locally closed. Then $A\cup L=A\sqcup(L\smallsetminus A)$ and $L\smallsetminus A=(\cdots((L\smallsetminus D_1)\smallsetminus D_2)\cdots)\smallsetminus D_m$ is a finite disjoint union of locally closed sets by the above. Induction on the number of pieces gives the claim. (The converse direction, that finite disjoint unions of locally closed sets are constructible, is the inclusion $\mathcal G\subseteq\mathcal F$.)

(b) Let $\xi$ be the generic point and $S=L_1\sqcup\dots\sqcup L_n$, $L_i=O_i\cap C_i$ locally closed. Suppose $\xi\in S$, say $\xi\in L_i$. Then $\xi\in C_i$ closed, so $C_i=X$, and $L_i=O_i$ is open containing $\xi$ (so nonempty). Thus $S\supseteq O_i\ne\emptyset$ and $S\ni\xi$ is dense since $\overline{\{\xi\}}=X$. Conversely suppose $S$ is dense: $X=\overline S=\bigcup\overline{L_i}$. As $X$ is irreducible, $X=\overline{L_i}$ for some $i$ (and $L_i\ne\emptyset$). Since $C_i$ is closed and contains $L_i$, $C_i\supseteq\overline{L_i}=X$, so $L_i=O_i$ is a nonempty open set, which contains $\xi$ by Ex.~3.17(d). So $\xi\in S$, and $S$ contains the nonempty open set $O_i$.

(c) ($\Rightarrow$) A closed set is the complement of an open one, hence constructible, and it is stable under specialization (Ex.~3.17(e)); an open set is constructible and stable under generization.

($\Leftarrow$, closed) Let $S$ be constructible and stable under specialization, $Z=\overline S$. We show $Z\subseteq S$. We first note: \emph{for a closed subset $W\subseteq X$ (with the induced topology, again a Zariski space), $S\cap W$ is constructible in $W$}: the family of $B\subseteq X$ with $B\cap W$ constructible in $W$ contains the open sets, and is closed under finite intersections and complements, since $W\smallsetminus(B\cap W)=(X\smallsetminus B)\cap W$; so it contains $\mathcal F$. Let $W$ be an irreducible component of $Z$ with generic point $w$ and $O=Z\smallsetminus(\text{the other components})$, a nonempty open subset of $Z$ contained in $W$ and containing $w$ (if $w$ lay in another component $W'$, then $W=\overline{\{w\}}\subseteq W'$). Since $S$ is dense in $Z$, $S\cap O$ is dense in $O$ (every nonempty open subset of $O$ is open in $Z$) hence $\overline{S\cap O}\supseteq\overline O\supseteq\overline{\{w\}}=W$. So $S\cap W$ is a dense constructible subset of the irreducible Zariski space $W$, and by (b) $w\in S\cap W$. As $Z$ is the union of its (finitely many) irreducible components $W=\overline{\{w\}}$ and $S$ is stable under specialization, $W\subseteq S$ for every component, so $Z\subseteq S$ and $S$ is closed.

($\Leftarrow$, open) Let $T$ be constructible and stable under generization. Then $S=X\smallsetminus T$ is constructible, and stable under specialization: if $x\in S$ and $x_0\in\overline{\{x\}}\cap T$, then $x$ is a generization of $x_0\in T$, so $x\in T$, a contradiction. By the closed case $S$ is closed, so $T$ is open.

(d) The family $\{B\subseteq Y: f^{-1}(B)\in\mathcal F_X\}$ contains the open sets of $Y$ (as $f$ is continuous) and is closed under finite intersections and complements (since preimages commute with these operations). By minimality of $\mathcal F_Y$ it contains $\mathcal F_Y$.\hfill$\square$

\exercise{3.19}{The real importance of the notion of constructible subsets derives from the following theorem of Chevalley---see Cartan and Chevalley [1, exposé 7] and see also Matsumura [2, Ch.~2, Sect.~6]: let $f:X\to Y$ be a morphism of finite type of noetherian schemes. Then the image of any constructible subset of $X$ is a constructible subset of $Y$. In particular, $f(X)$, which need not be either open or closed, is a constructible subset of $Y$. Prove this theorem in the following steps.
\begin{enumerate}[label=(\alph*),leftmargin=*]
\item Reduce to showing that $f(X)$ itself is constructible, in the case where $X$ and $Y$ are affine, integral noetherian schemes, and $f$ is a dominant morphism.
\item In that case, show that $f(X)$ contains a nonempty open subset of $Y$ by using the following result from commutative algebra: let $A\subseteq B$ be an inclusion of noetherian integral domains, such that $B$ is a finitely generated $A$-algebra. Then given a nonzero element $b\in B$, there is a nonzero element $a\in A$ with the following property: if $\varphi:A\to K$ is any homomorphism of $A$ to an algebraically closed field $K$, with $\varphi(a)\ne0$, then $\varphi$ extends to a homomorphism $\varphi'$ of $B$ into $K$, such that $\varphi'(b)\ne0$. [\emph{Hint}: Prove this algebraic result by induction on the number of generators of $B$ over $A$. For the case of one generator, prove the result directly. In the application, take $b=1$.]
\item Now use noetherian induction on $Y$ to complete the proof.
\item Give some examples of morphisms $f:X\to Y$ of varieties over an algebraically closed field $k$, to show that $f(X)$ need not be either open or closed. (Part (b) is starred in the book.)
\end{enumerate}}
\pf
In what follows ``constructible'' in a noetherian scheme $Y$ refers to the Zariski space $\operatorname{sp}(Y)$ (Ex.~3.17(a)). Basic observations:
\begin{itemize}
\item \textbf{(C1)} Finite unions, finite intersections, complements and differences of constructible sets are constructible.
\item \textbf{(C2)} If $B\subseteq Y$ is open or closed (a Zariski space with the induced topology), every subset of $B$ that is constructible in $B$ is constructible in $Y$. Indeed the family $\{C\subseteq B:C\in\mathcal F_Y\}$ contains the opens of $B$ (they are $B\cap O$, locally closed in $Y$, so in $\mathcal F_Y$) and is stable under finite intersection and complements relative to $B$ ($B\smallsetminus C=B\cap(Y\smallsetminus C)$ and $B\in\mathcal F_Y$); so it contains $\mathcal F_B$. Consequently, if $Y=Y_1\cup\dots\cup Y_n$ with $Y_i$ open and $S\cap Y_i$ is constructible in $Y_i$ for each $i$, then $S=\bigcup(S\cap Y_i)$ is constructible in $Y$.
\item \textbf{(C3)} If $W\subseteq Z\subseteq Y$ with $Z$ closed in $Y$ and $W$ open in $Z$, then $W$ is locally closed in $Y$, hence constructible.
\end{itemize}

\emph{(a) Reduction.} \emph{Claim: Chevalley's theorem follows from the following statement $(\star)$: if $f:X\to Y$ is a dominant morphism of finite type between integral affine noetherian schemes, then $f(X)$ is constructible in $Y$.} I formulate the reduction so that it can be reused for closed subsets of $Y$ in (c). Fix an affine noetherian scheme $Y$ and a closed subset $Z\subseteq Y$, and consider

$(\star)_Z$: for every integral closed subscheme $Z'$ of $Y$ with $Z'\subseteq Z$ (as sets) and every dominant finite-type morphism $f:X\to Z'$ with $X$ integral affine noetherian, $f(X)$ is constructible in $Y$ (it suffices to know this in $Z'$, by (C2)).

$(T)_Z$: for every finite-type morphism $h:X\to Y$ with $X$ noetherian and $h(X)\subseteq Z$, $h(X)$ is constructible in $Y$.

\emph{Claim: $(\star)_Z\Rightarrow(T)_Z$.} Let $h:X\to Y$ be as in $(T)_Z$. Cover $X$ by finitely many open affines $X_j$; the open immersion $X_j\to X$ is quasi-compact ($X$ is noetherian), hence of finite type (Ex.~3.13(b)), so $X_j\to Y$ is of finite type (Ex.~3.13(c)), and $h(X)=\bigcup h(X_j)$. Each $X_j$ is a noetherian affine scheme with finitely many irreducible components $X_{jl}$, each given the reduced induced structure: an integral affine noetherian scheme, with $X_{jl}\to X_j$ a closed immersion, so $X_{jl}\to Y$ is of finite type (Ex.~3.13(a),(c)), and $h(X)=\bigcup h(X_{jl})$. Let $Z'=\overline{h(X_{jl})}\subseteq Z$ ($Z$ is closed), irreducible, with its reduced induced structure: an integral closed subscheme of the affine scheme $Y$, so affine. Since $X_{jl}$ is reduced, $X_{jl}\to Y$ factors through $Z'$ (Ex.~3.11(d), or directly: $B\to A_{jl}$ kills the prime $\Ker$ of $Z'$ since $A_{jl}$ is a domain and the image is dense in $Z'$); call the factored map $f$. It is dominant, and of finite type by Ex.~3.13(f) ($X_{jl}$ is noetherian, so $f$ is quasi-compact, and the composite with the closed immersion $Z'\to Y$ is of finite type). By $(\star)_Z$, $f(X_{jl})=h(X_{jl})$ is constructible; so $h(X)$ is a finite union of constructible sets.

\emph{Claim: $(T)_Y$ for all affine noetherian $Y$ implies Chevalley's theorem.} First, to show that $f(C)$ is constructible for constructible $C\subseteq X$ it is enough to show that $f(X)$ is constructible for \emph{every} finite type morphism of noetherian schemes: by Ex.~3.18(a), $C=\bigcup L_i$ with $L_i=O_i\cap C_i$ ($O_i$ open, $C_i$ closed), and $L_i$ is the underlying set of the scheme $W_i=O_i\cap C_i$, an open subscheme of the closed subscheme $C_i$ (reduced structure), which is noetherian; $W_i\to C_i\to X$ is a quasi-compact open immersion followed by a closed immersion, hence of finite type, and so is $W_i\to X\to Y$ (Ex.~3.13). Then $f(L_i)$ is the image of $W_i\to Y$. Next, for a noetherian scheme $Y$ cover it by finitely many open affines $Y_i$. The set $X_i=f^{-1}(Y_i)$ is open in $X$, hence a noetherian scheme, and the open immersion $j_i:X_i\to X$ is quasi-compact, hence of finite type. So $f\circ j_i$ is of finite type, and it equals $(Y_i\hookrightarrow Y)\circ f_i$ with $f_i:X_i\to Y_i$ quasi-compact; by Ex.~3.13(f), $f_i$ is of finite type. Since $f(X)\cap Y_i=f_i(X_i)$, (C2) shows that $f(X)$ is constructible once each $f_i(X_i)$ is constructible in the affine noetherian scheme $Y_i$, which is the case $(T)_{Y_i}$.

\emph{(b) The algebraic lemma.} Let $R\subseteq S$ be domains, $S$ a finitely generated $R$-algebra, $0\ne s\in S$. \emph{Claim (L): there is $0\ne r\in R$ such that every homomorphism $\varphi:R\to K$ into an algebraically closed field with $\varphi(r)\ne0$ extends to $\varphi':S\to K$ with $\varphi'(s)\ne0$.} (Noetherianity is not needed.)

\emph{One generator, $S=R[x]$.} Let $F=\operatorname{Frac}R$.

\emph{Case 1: $x$ transcendental over $F$.} Then $S=R[x]$ is a polynomial ring; write $s=a_0x^m+\dots+a_m$ with $a_0\ne0$ and take $r=a_0$. If $\varphi(a_0)\ne0$, $\varphi(s)=\varphi(a_0)x^m+\dots\in K[x]$ is a nonzero polynomial; $K$ is infinite (algebraically closed), so there is $c\in K$ with $\varphi(s)(c)\ne0$. Let $\varphi'$ extend $\varphi$ by $x\mapsto c$; then $\varphi'(s)=\varphi(s)(c)\ne0$.

\emph{Case 2: $x$ algebraic over $F$.} Then $\operatorname{Frac}S=F(x)=F[x]$ is a finite extension of $F$, so $x$ and $s^{-1}\in\operatorname{Frac}S$ are algebraic over $F$: there are relations $a_0x^n+a_1x^{n-1}+\dots+a_n=0$ and $b_0s^{-m}+b_1s^{-m+1}+\dots+b_m=0$ with $a_i,b_i\in R$, $a_0b_0\ne0$. Put $r=a_0b_0$. Let $\varphi(r)\ne0$. Then $\varphi$ extends to $R_r\to K$. In $S'=S_{rs}\subseteq\operatorname{Frac}S$, which is generated over $R_r$ by $x$ and $s^{-1}$, both elements are integral over $R_r$ (divide by the units $a_0,b_0$). So $R_r\subseteq S'$ is an integral extension of domains. Let $\fq=\Ker(R_r\to K)$, a prime. By lying over there is a prime $\mathfrak Q$ of $S'$ with $\mathfrak Q\cap R_r=\fq$; then $R_r/\fq\subseteq S'/\mathfrak Q$ is an integral extension of domains, and $R_r/\fq\hookrightarrow K$. Since $K$ is algebraically closed, the embedding $\operatorname{Frac}(R_r/\fq)\to K$ extends to the algebraic extension $\operatorname{Frac}(S'/\mathfrak Q)$, giving $\varphi'':S'\to K$ extending $\varphi|_{R_r}$. Restrict to $S$: $\varphi'=\varphi''|_S$ extends $\varphi$, and $\varphi'(s)\ne0$ because $s$ is a unit of $S'$.

\emph{Induction.} Let $S=R[x_1,\dots,x_n]$, $n\ge2$, and $R_1=R[x_1]\subseteq S$ (a domain), so $S=R_1[x_2,\dots,x_n]$ has $n-1$ generators over $R_1$. By induction there is $0\ne r_1\in R_1$ such that every $\psi:R_1\to K$ with $\psi(r_1)\ne0$ extends to $\psi':S\to K$ with $\psi'(s)\ne0$. By the one-generator case for $R\subseteq R_1$ and the element $r_1$, there is $0\ne r\in R$ such that every $\varphi:R\to K$ with $\varphi(r)\ne0$ extends to $\varphi_1:R_1\to K$ with $\varphi_1(r_1)\ne0$. Composing the two extensions proves (L) for $S$ (the case $n=0$ is trivial).

\emph{Application.} Let $f:X=\Spec S\to Y=\Spec R$ be dominant, with $X,Y$ integral affine noetherian and $f$ of finite type. The kernel of $R\to S$ is a prime $\fp$ with $V(\fp)=\overline{f(X)}=Y$, so $\fp=0$ and $R\subseteq S$; $S$ is a finitely generated $R$-algebra (Ex.~3.3(c)). Apply (L) with $s=1$ to get $0\ne r\in R$. We claim $D(r)\subseteq f(X)$. Let $\fq\in D(r)$ and $K=\overline{k(\fq)}$, $\varphi:R\to R/\fq\to k(\fq)\subseteq K$. Then $\varphi(r)\ne0$ (as $r\notin\fq$), so $\varphi$ extends to $\varphi':S\to K$. The kernel $\mathfrak Q$ of $\varphi'$ is a prime of $S$ with $\mathfrak Q\cap R=\Ker\varphi=\fq$, i.e.\ $f(\mathfrak Q)=\fq$. Hence $f(X)\supseteq D(r)$, a nonempty open subset of $Y$ since $r\ne0$.

\emph{(c) Noetherian induction.} Let $Y$ be an affine noetherian scheme. We prove $(\star)_Z$ for all closed $Z\subseteq Y$ by noetherian induction on $Z$ (Ex.~3.16). Assume $(\star)_{Z''}$ holds for every proper closed subset $Z''\subsetneq Z$; by (a), $(T)_{Z''}$ holds for those. Let $Z'\subseteq Z$ be an integral closed subscheme (affine, noetherian) and $f:X\to Z'$ dominant of finite type with $X$ integral affine noetherian. By (b) there is a nonempty open $W\subseteq Z'$ (open in $Z'$) with $W\subseteq f(X)$. $W$ is constructible in $Y$ by (C3). The set $Z'\smallsetminus W$ is a closed subset of $Z'$, hence of $Y$, and it is a proper subset of $Z'\subseteq Z$; so $(T)_{Z'\smallsetminus W}$ holds. Let $X''=f^{-1}(Z'\smallsetminus W)$ with its reduced induced structure: a closed subscheme of $X$, hence a noetherian scheme, and $X''\to X\to Z'\to Y$ is of finite type (Ex.~3.13(a),(c)) with image in $Z'\smallsetminus W$. By $(T)_{Z'\smallsetminus W}$, $f(X'')=f(X)\cap(Z'\smallsetminus W)$ is constructible in $Y$. Therefore $f(X)=W\cup(f(X)\cap(Z'\smallsetminus W))$ is constructible. This proves $(\star)_Z$. By Ex.~3.16, $(\star)_Y$ holds, hence $(T)_Y$ by (a) (the induction needs no base case since for $Z=\emptyset$ the statement is vacuous). By the last claim of (a), Chevalley's theorem holds for all finite type morphisms of noetherian schemes.

\emph{(d) Examples.} Let $k$ be algebraically closed and $f:\mathbf A^2_k\to\mathbf A^2_k$, $(x,y)\mapsto(x,xy)$ (a morphism of varieties). The image on closed points is $\{(a,b):a\ne0\}\cup\{(0,0)\}$: for $a\ne0$, $(a,b)=f(a,b/a)$, and $f(0,y)=(0,0)$. It is constructible ($D(x)\cup\{(0,0)\}$) but not closed: it contains the dense open set $D(x)$ but not $(0,1)$. It is not open: let $O\ni(0,0)$ be open, $O=\mathbf A^2\smallsetminus V(J)$. Since $(0,0)\notin V(J)$ there is $g\in J$ with $g(0,0)\ne0$; then $g(0,y)$ is a nonzero polynomial in $y$, so $V(J)\cap\{x=0\}$ is finite and $O$ contains infinitely many points $(0,b)$, $b\ne0$, which are not in the image. (In scheme language the image is $D(x)\cup\{\text{the origin}\}$.) Another standard example: the projection of the closed set $V(xy-1)\subset\mathbf A^2$ to the $x$-axis has image $\mathbf A^1\smallsetminus\{0\}$, which is open but not closed.\hfill$\square$

\exercise{3.20}{\emph{Dimension.} Let $X$ be an integral scheme of finite type over a field $k$ (not necessarily algebraically closed). Use appropriate results from (I, Sect.~1) to prove the following.
\begin{enumerate}[label=(\alph*),leftmargin=*]
\item For any closed point $P\in X$, $\dim X=\dim\mathcal O_{P}$, where for rings, we always mean the Krull dimension.
\item Let $K(X)$ be the function field of $X$ (Ex.~3.6). Then $\dim X=\operatorname{tr.d.}K(X)/k$.
\item If $Y$ is a closed subset of $X$, then $\operatorname{codim}(Y,X)=\inf\{\dim\mathcal O_{P,X}\mid P\in Y\}$.
\item If $Y$ is a closed subset of $X$, then $\dim Y+\operatorname{codim}(Y,X)=\dim X$.
\item If $U$ is a nonempty open subset of $X$, then $\dim U=\dim X$.
\item If $k\subseteq k'$ is a field extension, then every irreducible component of $X'=X\times_kk'$ has dimension $=\dim X$.
\end{enumerate}}
\pf
\emph{Conventions.} In (c), $P$ ranges over all (not only the closed) points of $Y$; over closed points alone the formula would be false (by (a) all closed points of $X$ have $\dim\mathcal O_P=\dim X$). The facts from (I, Sect.~1) used: for a finitely generated $k$-algebra $B$ which is a domain, (1.8A) $\dim B=\operatorname{tr.d.}\operatorname{Frac}(B)/k$, and $\operatorname{ht}\fp+\dim B/\fp=\dim B$ for every prime $\fp$; and $\dim\Spec B$ (as a topological space) equals the Krull dimension of $B$.

\emph{Topological preliminaries.}
\begin{itemize}
\item \textbf{(D1)} Let $U$ be an open subset of a topological space $T$. Then $Z\mapsto Z\cap U$ is a bijection from the irreducible closed subsets of $T$ meeting $U$ onto the irreducible closed subsets of $U$, with inverse $Z'\mapsto\overline{Z'}$, and both preserve strict inclusions (a nonempty open subset of an irreducible space is dense and irreducible). Hence $\dim U\le\dim T$.
\item \textbf{(D2)} If $T=\bigcup T_i$ is an open cover, then $\dim T=\sup\dim T_i$. Indeed ``$\ge$'' is (D1). For ``$\le$'': let $Z_0\subsetneq\dots\subsetneq Z_n$ be irreducible closed subsets of $T$ and choose $i$ with $Z_0\cap T_i\ne\emptyset$; then $Z_j\cap T_i$ ($j=0,\dots,n$) is a strictly increasing chain of irreducible closed subsets of $T_i$ by (D1), so $n\le\dim T_i$.
\item \textbf{(D3)} For any point $P$ of any scheme $X$, $\dim\mathcal O_{X,P}=\operatorname{codim}(\overline{\{P\}},X)$. Indeed let $U=\Spec A$ be an open affine containing $P$, $P\leftrightarrow\fp$. Every irreducible closed subset of $X$ containing $P$ meets $U$, so by (D1) the irreducible closed subsets of $X$ containing $\overline{\{P\}}$ correspond, preserving inclusions, to those of $U$ containing $\overline{\{P\}}\cap U=V(\fp)$, i.e.\ to the primes $\fq\subseteq\fp$. So $\operatorname{codim}(\overline{\{P\}},X)=\operatorname{ht}\fp=\dim A_\fp=\dim\mathcal O_{X,P}$.
\end{itemize}

\emph{Setup.} $X$ is covered by finitely many open affines $U_i=\Spec A_i$; by Ex.~3.3(c) every nonempty open affine $\Spec A\subseteq X$ has $A$ a finitely generated $k$-algebra, and a domain because $X$ is integral, with $\operatorname{Frac}A=K(X)$ (Ex.~3.6).

(b) For every nonempty open affine $U=\Spec A$, $\dim U=\dim A=\operatorname{tr.d.}K(X)/k=:d$ by (1.8A). By (D2) applied to a cover of $X$ by such $U$, $\dim X=d$.

(e) A nonempty open $U\subseteq X$ is covered by nonempty open affines $\Spec A$ of $X$, each of dimension $d$ by the above; by (D2), $\dim U=d=\dim X$.

(a) Let $P$ be a closed point, $U=\Spec A$ an open affine containing it, $P\leftrightarrow\fp$. As $P$ is closed in $U$, $\fp$ is maximal, so $\dim A/\fp=0$ and (1.8A) gives $\operatorname{ht}\fp=\dim A=\dim X$ by (b). Hence $\dim\mathcal O_P=\operatorname{ht}\fp=\dim X$.

(c) Recall $\operatorname{codim}(Y,X)=\inf\operatorname{codim}(Z,X)$ over irreducible closed $Z\subseteq Y$. Every irreducible closed $Z\subseteq Y$ is $\overline{\{P\}}$ for its generic point $P\in Y$, and conversely $\overline{\{P\}}\subseteq Y$ for each $P\in Y$. So by (D3), $\operatorname{codim}(Y,X)=\inf_{P\in Y}\operatorname{codim}(\overline{\{P\}},X)=\inf_{P\in Y}\dim\mathcal O_{P,X}$.

(d) \emph{Irreducible case.} Let $Z\subseteq X$ be closed irreducible, with the reduced induced structure: an integral closed subscheme, of finite type over $k$ (Ex.~3.13(a),(c)). Let $V=\Spec A\subseteq X$ be an open affine meeting $Z$, so $Z\cap V=V(\fp)=\Spec A/\fp$ for a prime $\fp$. By (e) for $Z$, $\dim Z=\dim(Z\cap V)=\dim A/\fp$. By (D1), the irreducible closed subsets of $X$ containing $Z$ correspond to those of $V$ containing $Z\cap V$, i.e.\ to primes contained in $\fp$; so $\operatorname{codim}(Z,X)=\operatorname{ht}\fp$. By (1.8A) and (b), $\dim Z+\operatorname{codim}(Z,X)=\dim A/\fp+\operatorname{ht}\fp=\dim A=\dim X$.

\emph{General case.} Let $Y=Y_1\cup\dots\cup Y_r$ be the decomposition into irreducible components. $\dim Y=\max\dim Y_j$. Every irreducible closed $Z\subseteq Y$ lies in some $Y_j$, and then $\operatorname{codim}(Z,X)\ge\operatorname{codim}(Y_j,X)$ (a chain starting with $Y_j$ can be extended by $Z\subseteq Y_j$ at the beginning if $Z\ne Y_j$). So $\operatorname{codim}(Y,X)=\min_j\operatorname{codim}(Y_j,X)$, and by the irreducible case
\[\dim Y+\operatorname{codim}(Y,X)=\max_j\dim Y_j+\min_j(\dim X-\dim Y_j)=\dim X.\]

(f) By Ex.~3.13(d), $X'$ is of finite type over $k'$. Let $X_i=\Spec A_i$ be finitely many open affines covering $X$; $X_i'=X_i\times_kk'=\Spec S_i$ with $S_i=A_i\otimes_kk'$ cover $X'$. Let $d=\dim X$.

\emph{Claim: if $P$ is a minimal prime of $S=S_i$, then $\dim S/P=d$.} By Noether normalization there is $P_0=k[t_1,\dots,t_d]\subseteq A=A_i$ with $A$ a finite $P_0$-module ($d=\operatorname{tr.d.}\operatorname{Frac}A=\dim A$). Put $R=P_0\otimes_kk'=k'[t_1,\dots,t_d]$; the inclusion $P_0\subseteq A$ stays injective after $\otimes_kk'$ (flatness over a field), so $R\subseteq S$, and $S$ is a finite $R$-module. Let $r\in R$, $r\ne0$. We show $r$ is a non-zero-divisor on $S$. Let $F_0=k(t_1,\dots,t_d)$ and $M=\operatorname{Frac}A$; $M$ is a finite-dimensional $F_0$-vector space ($A\otimes_{P_0}F_0$ is a domain finite over $F_0$, hence a field, equal to $M$). Then $S=A\otimes_kk'\subseteq M\otimes_kk'=M\otimes_{F_0}R'$ with $R'=F_0\otimes_kk'$; here $R'$ is the localization of $R$ at the multiplicative set $P_0\smallsetminus0$, a subring of $k'(t_1,\dots,t_d)$, so a domain. As $M\otimes_{F_0}R'\cong R'^{\,n}$ is free over $R'$ ($n=\dim_{F_0}M$), $r$ acts injectively on it, hence on $S$. Now a minimal prime $P$ of $S$ consists of zero-divisors (in $S_P$, which has a unique prime, every element of $PS_P$ is nilpotent; so for $x\in P$ there are $u\notin P$ and $n\ge1$ minimal with $ux^n=0$, and then $x\cdot ux^{n-1}=0$ with $ux^{n-1}\ne0$). Therefore $P\cap R=0$. Then $R\hookrightarrow S/P$ is an integral extension of domains, so $\operatorname{Frac}(S/P)$ is algebraic over $k'(t_1,\dots,t_d)$ and $\operatorname{tr.d.}_{k'}\operatorname{Frac}(S/P)=d$. Since $S/P$ is a finitely generated $k'$-domain, $\dim S/P=d$ by (1.8A).

Now let $Z'$ be an irreducible component of $X'$. The sets $Z'\cap X_i'$ that are nonempty are irreducible components of $X'_i=\Spec S_i$ (the irreducible components of an open subset are the nonempty traces of those of the space), i.e.\ $Z'\cap X_i'=V(P)$ for a minimal prime $P$ of $S_i$, of dimension $\dim S_i/P=d$ by the Claim. These sets form an open cover of $Z'$, so by (D2) $\dim Z'=d=\dim X$.\hfill$\square$

\exercise{3.21}{Let $R$ be a discrete valuation ring containing its residue field $k$. Let $X=\Spec R[t]$ be the affine line over $\Spec R$. Show that statements (a), (d), (e) of (Ex.~3.20) are false for $X$.}
\pf
Let $\pi$ be a uniformizer, $K=R[1/\pi]=\operatorname{Frac}R$, and recall $R/(\pi)\cong k$.

\emph{$\dim X=2$.} There is a chain of primes $(0)\subsetneq(\pi)\subsetneq(\pi,t)$ in $R[t]$: $R[t]/(\pi)=k[t]$ is a domain, and $R[t]/(\pi,t)=k$ is a field. So $\dim R[t]\ge2$. Since $R$ is noetherian of dimension $1$, $\dim R[t]=\dim R+1=2$ (a standard fact on polynomial rings over noetherian rings, e.g.\ Matsumura, \emph{Commutative Algebra}, Thm.~15.4, or Atiyah--Macdonald, Ex.~11.7).

\emph{(a) fails.} Consider $\fm=(\pi t-1)\subset R[t]$. Since $R_\pi=R[x]/(\pi x-1)$ (the standard presentation of a localization), $R[t]/(\pi t-1)\cong R_\pi=K$ via $t\mapsto1/\pi$. Hence $R[t]/\fm\cong K$ is a field and $\fm$ is a maximal ideal, a closed point $P$ of $X$. By Krull's principal ideal theorem, $\operatorname{ht}\fm\le1$, and $\fm\ne0$ is prime in a domain, so $\operatorname{ht}\fm=1$. Thus $\dim\mathcal O_P=1\ne2=\dim X$.

\emph{(d) fails.} Let $Y=\{P\}$ be this closed point. Then $\dim Y=0$ and $\operatorname{codim}(Y,X)=\operatorname{codim}(\overline{\{P\}},X)=\operatorname{ht}\fm=1$, by (D3) of Ex.~3.20 (valid on any scheme). So $\dim Y+\operatorname{codim}(Y,X)=1\ne2=\dim X$.

\emph{(e) fails.} Let $U=D(\pi)=\Spec R[t]_\pi=\Spec K[t]$, a nonempty open subset of $X$. Then $\dim U=\dim K[t]=1\ne2=\dim X$.\hfill$\square$

\exercise{3.22}{\emph{Dimension of the Fibres of a Morphism.} Let $f:X\to Y$ be a dominant morphism of integral schemes of finite type over a field $k$.
\begin{enumerate}[label=(\alph*),leftmargin=*]
\item Let $Y'$ be a closed irreducible subset of $Y$, whose generic point $\eta'$ is contained in $f(X)$. Let $Z$ be any irreducible component of $f^{-1}(Y')$, such that $\eta'\in f(Z)$, and show that $\operatorname{codim}(Z,X)\le\operatorname{codim}(Y',Y)$.
\item Let $e=\dim X-\dim Y$ be the \emph{relative dimension} of $X$ over $Y$. For any point $y\in f(X)$, show that every irreducible component of the fibre $X_y$ has dimension $\ge e$. [\emph{Hint}: Let $Y'=\overline{\{y\}}$, and use (a) and (Ex.~3.20b).]
\item Show that there is a dense open subset $U\subseteq X$, such that $\dim U_y=e$, for any $y\in f(U)$. [\emph{Hint}: First reduce to the case where $X$ and $Y$ are affine, say $X=\Spec A$ and $Y=\Spec B$. Then $A$ is a finitely generated $B$-algebra. Take $t_1,\dots,t_e\in A$ which form a transcendence basis of $K(X)$ over $K(Y)$, and let $X_1=\Spec B[t_1,\dots,t_e]$. Then $X_1$ is isomorphic to affine $e$-space over $Y$, and the morphism $X\to X_1$ is generically finite. Now use (Ex.~3.7) above.]
\item Going back to our original morphism $f:X\to Y$, for any integer $h$, let $E_h$ be the set of points $x\in X$ such that, letting $y=f(x)$, there is an irreducible component $Z$ of the fibre $X_y$, containing $x$, and having $\dim Z\ge h$. Show that (1) $E_e=X$ (use (b) above); (2) if $h>e$, then $E_h$ is not dense in $X$ (use (c) above); and (3) $E_h$ is closed, for all $h$ (use induction on $\dim X$).
\item Prove the following theorem of Chevalley---see Cartan and Chevalley [1, exposé 8]. For each integer $h$, let $C_h$ be the set of points $y\in Y$ such that $\dim X_y=h$. Then the subsets $C_h$ are constructible, and $C_e$ contains an open dense subset of $Y$.
\end{enumerate}
(Starred in the book.)}
\pf
\emph{Conventions.} By Ex.~3.10(a) the irreducible components of the fibre $X_y$ correspond homeomorphically to the irreducible components of the subspace $f^{-1}(y)\subseteq X$; I work with $f^{-1}(y)$ and, for dimensions, treat it as a topological space. $X_y$ is a scheme of finite type over $k(y)$ (Ex.~3.13(d)). Note that $X,Y$ are integral of finite type over $k$ with generic points $\xi,\eta$, $f(\xi)=\eta$ (Step~0 of Ex.~3.7), and $f$ is of finite type: $X\to\Spec k$ is of finite type, and $f$ is quasi-compact ($X$ is noetherian), so Ex.~3.13(f) applies. For a point $x$, $\dim\mathcal O_x=\operatorname{codim}(\overline{\{x\}},\cdot)$ by (D3) of Ex.~3.20.

(a) Let $\zeta$ be the generic point of $Z$. Since $f(Z)\subseteq\overline{\{f(\zeta)\}}$ and $\eta'\in f(Z)$, we get $\eta'\in\overline{\{f(\zeta)\}}$. Also $f(\zeta)\in Y'$ ($Z\subseteq f^{-1}(Y')$), so $\overline{\{f(\zeta)\}}\subseteq Y'=\overline{\{\eta'\}}$. Hence $\overline{\{f(\zeta)\}}=Y'$ and $f(\zeta)=\eta'$ by uniqueness of the generic point. So $f$ induces a local homomorphism of noetherian local rings $f^\#:\mathcal O_{Y,\eta'}\to\mathcal O_{X,\zeta}$.

Let $d=\dim\mathcal O_{Y,\eta'}=\operatorname{codim}(Y',Y)$. Choose $x_1,\dots,x_d\in\fm_{\eta'}$ with $\fm_{\eta'}=\sqrt{(x_1,\dots,x_d)}$ (a system of parameters). Let $I\subseteq\mathcal O_{X,\zeta}$ be the ideal generated by $f^\#(\fm_{\eta'})$. We show $\sqrt I=\fm_\zeta$. Primes of $\mathcal O_{X,\zeta}$ correspond to points $x$ of which $\zeta$ is a specialization ($\zeta\in\overline{\{x\}}$), and the contraction under $f^\#$ of the prime corresponding to $x$ is the prime of $\mathcal O_{Y,\eta'}$ corresponding to $f(x)$. If a prime $\fq_x\supseteq I$ then its contraction contains $\fm_{\eta'}$, so equals it, i.e.\ $f(x)=\eta'\in Y'$; then $x\in f^{-1}(Y')$ and $\overline{\{x\}}$ is an irreducible subset of $f^{-1}(Y')$ containing $\zeta$, hence containing $Z=\overline{\{\zeta\}}$; since $Z$ is an irreducible component, $\overline{\{x\}}=Z$ and $x=\zeta$. So the only prime containing $I$ is $\fm_\zeta$, i.e.\ $\sqrt I=\fm_\zeta=\sqrt{(f^\#x_1,\dots,f^\#x_d)}$. By Krull's height theorem, $\operatorname{ht}\fm_\zeta\le d$. Hence $\operatorname{codim}(Z,X)=\dim\mathcal O_{X,\zeta}\le d=\operatorname{codim}(Y',Y)$.

(b) Let $y\in f(X)$, $Y'=\overline{\{y\}}$, and let $W'$ be an irreducible component of $f^{-1}(y)$ with generic point $w$ (it exists and is unique because $W'$ corresponds, by Ex.~3.10(a), to an irreducible closed subset of the scheme $X_y$; thus $W'=\overline{\{w\}}\cap f^{-1}(y)$, as $f^{-1}(y)$ has the induced topology), $Z=\overline{\{w\}}\subseteq X$ (closure in $X$). Let $Z_0$ be an irreducible component of $f^{-1}(Y')$ containing $Z$ and $\zeta_0$ its generic point. As in (a), $y=f(w)\in f(Z_0)\subseteq\overline{\{f(\zeta_0)\}}\subseteq Y'$ gives $f(\zeta_0)=y$. So $\zeta_0\in f^{-1}(y)$ and $w\in\overline{\{\zeta_0\}}$; the closure of $\{\zeta_0\}$ in $f^{-1}(y)$ is an irreducible subset of $f^{-1}(y)$ containing $w$, hence containing $W'$, hence equal to $W'$ by maximality; so $\zeta_0$ is a generic point of $W'$, i.e.\ $\zeta_0=w$, and $Z=Z_0$ is an irreducible component of $f^{-1}(Y')$ with $y\in f(Z)$. By (a), $\operatorname{codim}(Z,X)\le\operatorname{codim}(Y',Y)$. Using Ex.~3.20(d) for $Z\subseteq X$ and $Y'\subseteq Y$,
\[\dim Z=\dim X-\operatorname{codim}(Z,X)\ge\dim X-\operatorname{codim}(Y',Y)=\dim X-\dim Y+\dim Y'=e+\dim Y'.\]
Now $Z$ with its reduced structure is an integral scheme of finite type over $k$ with function field $k(w)$ (the residue field of $X$ at $w$), so $\dim Z=\operatorname{tr.d.}_kk(w)$ by Ex.~3.20(b); likewise $\dim Y'=\operatorname{tr.d.}_kk(y)$. The component $W'$ of $f^{-1}(y)$, i.e.\ of $X_y$ (Ex.~3.10(a)), is an integral scheme of finite type over $k(y)$ (reduced structure), with function field the residue field at $w$ of $X_y$, which is $k(w)$. Therefore, by Ex.~3.20(b) over the field $k(y)$ and additivity of transcendence degree,
\[\dim W'=\operatorname{tr.d.}_{k(y)}k(w)=\operatorname{tr.d.}_kk(w)-\operatorname{tr.d.}_kk(y)=\dim Z-\dim Y'\ge e.\]

(c) \emph{Reduction to the affine case.} Let $V=\Spec B\subseteq Y$ be a nonempty open affine and $X_0=\Spec A\subseteq f^{-1}(V)$ a nonempty open affine (nonempty since $\xi\in f^{-1}(V)$). If we find a dense open $U\subseteq X_0$ with $\dim U_y=e$ for all $y\in f(U)$, then $U$ is dense open in $X$ and the same holds for $f$ (the fibre of $f|_U$ over $y$ is the same scheme; also $\dim X_0-\dim V=e$ by Ex.~3.20(e)). So assume $X=\Spec A$, $Y=\Spec B$, $B\subseteq A$ domains (injective as $f$ is dominant, as in Ex.~3.7), finitely generated $k$-algebras, $A$ a finitely generated $B$-algebra (Ex.~3.3(c), $f$ being of finite type).

By Ex.~3.20(b), $\operatorname{tr.d.}_{K(Y)}K(X)=\dim X-\dim Y=e$. $K(X)$ is generated as a field over $K(Y)$ by the (finitely many) generators of $A$ over $B$, so we can choose $t_1,\dots,t_e\in A$ forming a transcendence basis of $K(X)/K(Y)$. Then $B[t_1,\dots,t_e]\subseteq A$ is a polynomial ring over $B$, $X_1:=\Spec B[t_1,\dots,t_e]\to Y$ is affine $e$-space over $Y$, and $g:X\to X_1$ induced by the inclusion is dominant, of finite type ($A$ is a finitely generated $B[t]$-algebra), between integral schemes. It is generically finite: with $K_1=K(Y)(t_1,\dots,t_e)=K(X_1)$ the fibre over the generic point is $\Spec(S^{-1}A)$, $S=B[t]\smallsetminus0$; $S^{-1}A\subseteq K(X)$ is a domain generated over $K_1$ by the generators of $A$, which are algebraic over $K_1$ ($K(X)/K_1$ is algebraic), so it is finite-dimensional over $K_1$, hence a field, with exactly one prime. By Ex.~3.7 there is a dense open $U_1\subseteq X_1$ with $g^{-1}(U_1)\to U_1$ finite. Put $U=g^{-1}(U_1)$, a nonempty (it contains $\xi$) hence dense open subset of $X$. The morphism $g|_U:U\to U_1$ is finite, hence closed (Ex.~3.5(b)), and its image contains the generic point of $X_1$ ($g(\xi)$), so it is dense; therefore $g|_U$ is surjective.

Let $y\in f(U)$ and $p:X_1\to Y$ the projection. The fibre $U_y=U\times_Y\Spec k(y)=U\times_{U_1}(U_1)_y$, and $(U_1)_y$ is an open subscheme of $(X_1)_y=\Spec(B[t]\otimes_Bk(y))=\mathbf A^e_{k(y)}$. It is nonempty (as $U_y\neq\emptyset$ and $U_y\to(U_1)_y$). The morphism $U_y\to(U_1)_y$ is the base change of $g|_U$, hence finite (Spec of finite algebras) and surjective (the fibre over $z$ is $\Spec(k(z)\otimes_{k(z')}(\text{fibre over }z'))\ne\emptyset$ where $z'\in U_1$ is the image of $z$). Now:

\emph{Lemma E.} A finite surjective morphism $h:T\to T'$ of schemes has $\dim T=\dim T'$. Indeed take an open affine cover $W=\Spec C$ of $T'$ with $h^{-1}(W)=\Spec D$, $D$ a finite $C$-algebra; $\Spec D\to\Spec C$ is surjective, so $\Ker(C\to D)$ lies in every prime, and $C/\Ker\subseteq D$ is an integral extension, so $\dim D=\dim C/\Ker=\dim C$ (going up and incomparability, Atiyah--Macdonald Cor.~5.9, 5.11). Then (D2) of Ex.~3.20 for the covers $\{W\}$ of $T'$ and $\{h^{-1}W\}$ of $T$ gives $\dim T=\sup\dim D=\sup\dim C=\dim T'$.

So $\dim U_y=\dim(U_1)_y=e$, the latter by Ex.~3.20(e) ($(U_1)_y$ is a nonempty open subset of the integral $k(y)$-scheme $\mathbf A^e_{k(y)}$, of dimension $e$ by (b)).

(d) For an arbitrary finite type morphism $\phi:X\to Y$ of schemes of finite type over $k$ (not nec.\ integral) define $E_h(\phi)$ as in the statement; this depends only on the topological fibres $\phi^{-1}(y)$, their irreducible components and their dimensions as topological spaces.

(1) Let $x\in X$, $y=f(x)$. The point $x$ lies in some irreducible component $Z$ of $f^{-1}(y)$, and $\dim Z\ge e$ by (b). So $x\in E_e$. As $E_h\supseteq E_e$ for $h\le e$, we get $E_h=X$ for $h\le e$.

(2) Let $h>e$ and $U$ be as in (c). Let $x\in U$, $y=f(x)$ and $Z$ an irreducible component of $f^{-1}(y)$ containing $x$. Then $Z\cap U$ is a nonempty open subset of $Z$, and it is a closed irreducible subset of $U_y=U\cap f^{-1}(y)$; here $Z$ (reduced structure) is an integral scheme of finite type over $k(y)$, so by Ex.~3.20(e), $\dim Z=\dim(Z\cap U)\le\dim U_y=e<h$. Hence $x\notin E_h$. So $E_h\cap U=\emptyset$ and $E_h$ is not dense.

(3) We prove: for every morphism $\phi:X\to Y$ of schemes of finite type over $k$ and every $h$, $E_h(\phi)$ is closed in $X$, by induction on $\dim X$ ($X=\emptyset$ is trivial).

\emph{Reduction to $X$, $Y$ integral and $\phi$ dominant.} Let $X_1,\dots,X_m$ be the irreducible components of $X$, with reduced structures (integral, of finite type over $k$) and $\phi_j=\phi|_{X_j}$. Then $E_h(\phi)=\bigcup_jE_h(\phi_j)$. Indeed, if $x\in E_h(\phi)$ with $y=\phi(x)$ and $Z$ a component of $\phi^{-1}(y)$ containing $x$ with $\dim Z\ge h$, then $Z\subseteq X_j$ for some $j$ ($Z$ is irreducible), and $Z$ is then a component of $\phi_j^{-1}(y)=\phi^{-1}(y)\cap X_j$ (maximality persists in the smaller set), so $x\in E_h(\phi_j)$. Conversely a component $Z_j$ of $\phi_j^{-1}(y)$ lies in a component $Z$ of $\phi^{-1}(y)$ with $\dim Z\ge\dim Z_j$, so $E_h(\phi_j)\subseteq E_h(\phi)$. Also $E_h(\phi_j)$ is closed in $X_j$ iff in $X$. Next, replace $Y$ by $Y_j=\overline{\phi(X_j)}$ with the reduced structure, an integral scheme of finite type over $k$ (this doesn't change the topological fibres), and $\phi_j:X_j\to Y_j$ is dominant. So we may assume $X$, $Y$ integral and $\phi=f$ dominant (this does not increase $\dim X$; the inductive step below then lowers it).

\emph{Inductive step.} Let $f:X\to Y$ be dominant of integral finite type $k$-schemes, $h>e$ (for $h\le e$, $E_h=X$ by (1)). Let $U$ be as in (c) and $X'=X\smallsetminus U$, a proper closed subset, so $\dim X'<\dim X$ by Ex.~3.20(d) (each component $X'_i$ has codimension $\ge1$) if $X'\neq\emptyset$. By (2) $E_h\cap U=\emptyset$, so $E_h\subseteq X'$. Moreover, if $x\in E_h$ and $Z$ is a component of $f^{-1}(y)$ through $x$ with $\dim Z\ge h$, every point of $Z$ lies in $E_h$ (the same $Z$ works), so $Z\cap U=\emptyset$, i.e.\ $Z\subseteq X'$. Then $Z$ is also a component of $(f|_{X'})^{-1}(y)=f^{-1}(y)\cap X'$ and $x\in E_h(f|_{X'})$; conversely $E_h(f|_{X'})\subseteq E_h(f)$ as above. So $E_h(f)=E_h(f|_{X'})$, where $f|_{X'}:X'\to Y$ is a morphism of finite type $k$-schemes (reduced structure on $X'$) with $\dim X'<\dim X$. By the induction hypothesis $E_h(f|_{X'})$ is closed in $X'$, hence in $X$. (If $X'=\emptyset$, $E_h=\emptyset$.) This completes the induction.

(e) For each integer $h$, let $D_h=\{y\in Y:\dim X_y\ge h\}$. A point $y$ is in $D_h$ iff $y\in f(X)$ and some component of $f^{-1}(y)$ has dimension $\ge h$, iff $y\in f(E_h)$ (if $Z$ is such a component, every point of $Z$ lies in $E_h$; conversely if $x\in E_h$, $\dim f^{-1}(f(x))\ge h$). So $D_h=f(E_h)$. By (d)(3), $E_h$ is closed; with the reduced induced structure it is a noetherian scheme and $E_h\to X\to Y$ is of finite type (Ex.~3.13(a),(c)). By Chevalley's theorem (Ex.~3.19), $D_h$ is constructible. Hence for $h\ge0$, $C_h=D_h\smallsetminus D_{h+1}$ is constructible. (If $h<0$ is allowed, set $\dim\emptyset=-\infty$, so $C_h=\emptyset$.)

It remains to see that $C_e$ contains a nonempty open set. We show $\eta\in C_e$: the fibre $X_\eta$ has $\xi$ in it, and by the argument of (b) (with $y=\eta$, $Y'=Y$) every irreducible component $W'$ of $f^{-1}(\eta)$ has generic point $w$ with $Z=\overline{\{w\}}$ an irreducible component of $f^{-1}(Y)=X$, so $Z=X$, $w=\xi$, and $\dim W'=\dim Z-\dim Y=\dim X-\dim Y=e$ (formula of (b)). So $f^{-1}(\eta)$ is irreducible of dimension exactly $e$, i.e.\ $\dim X_\eta=e$, $\eta\in C_e$. Since $C_e$ is constructible and contains the generic point of the irreducible Zariski space $Y$, by Ex.~3.18(b) it contains a nonempty open subset of $Y$, which is dense.\hfill$\square$

\exercise{3.23}{If $V,W$ are two varieties over an algebraically closed field $k$, and if $V\times W$ is their product, as defined in (I, Ex.~3.15, 3.16), and if $t$ is the functor of (2.6), then $t(V\times W)=t(V)\times_{\Spec k}t(W)$.}
\pf
Recall (I, Ex.~3.15, 3.16): if $V_i\subseteq V$, $W_j\subseteq W$ are open affine subvarieties, then $V_i\times W_j$ is an affine variety with coordinate ring $A(V_i)\otimes_kA(W_j)$, the sets $V_i\times W_j=q_1^{-1}(V_i)\cap q_2^{-1}(W_j)$ are open in $V\times W$ (where $q_1,q_2$ are the projections) and cover it, and the isomorphism $A(V_i)\otimes_kA(W_j)\to A(V_i\times W_j)$ is $a\otimes b\mapsto q_1^*a\cdot q_2^*b$. Let $T=t(V)\times_{\Spec k}t(W)$ with projections $p_1,p_2$ (it exists by Thm.~3.3).

The projections $q_1,q_2$ give morphisms $t(q_1),t(q_2)$ from $t(V\times W)$ to $t(V),t(W)$ (compatible with the structure morphisms to $\Spec k$, since $t$ is a functor to schemes over $k$, Prop.~2.6), hence a morphism $\theta:t(V\times W)\to T$ with $p_i\theta=t(q_i)$. We show $\theta$ is an isomorphism, locally on $T$.

By Prop.~2.6, for an affine variety $V_i$ the scheme $t(V_i)$ is $\Spec A(V_i)$, and for an open subvariety $t$ gives the corresponding open subscheme. Thus $t(V)$ is covered by the open affine subschemes $t(V_i)=\Spec A(V_i)$ and $t(W)$ by $t(W_j)=\Spec A(W_j)$. The open sets $T_{ij}=p_1^{-1}(t(V_i))\cap p_2^{-1}(t(W_j))$ cover $T$, and by Thm.~3.3 (Steps~4 and~1) $T_{ij}=t(V_i)\times_{\Spec k}t(W_j)=\Spec(A(V_i)\otimes_kA(W_j))$. Also $\theta^{-1}(T_{ij})=t(q_1)^{-1}(t(V_i))\cap t(q_2)^{-1}(t(W_j))=t(q_1^{-1}V_i\cap q_2^{-1}W_j)=t(V_i\times W_j)=\Spec A(V_i\times W_j)$, using that $t$ takes open subvarieties to open subschemes and preserves preimages of open sets (the homeomorphism of $V\times W$ with the closed points). The restriction $\theta:\Spec A(V_i\times W_j)\to\Spec(A(V_i)\otimes_kA(W_j))$ is, by the universal property of the product (Thm.~3.3, Step~1), the morphism corresponding to the ring homomorphism $A(V_i)\otimes_kA(W_j)\to A(V_i\times W_j)$ whose restrictions to the two factors are $q_1^*$ and $q_2^*$, i.e.\ $a\otimes b\mapsto q_1^*a\cdot q_2^*b$. This is an isomorphism by (I, Ex.~3.15). Hence $\theta$ is an isomorphism over each $T_{ij}$; since the $T_{ij}$ cover $T$, $\theta$ is an isomorphism, i.e.\ $t(V\times W)\cong t(V)\times_{\Spec k}t(W)$ (compatibly with the projections).\hfill$\square$

\section*{4. Separated and proper morphisms}

% ---------------------------------------------------------------
\exercise{4.1}{Show that a finite morphism is proper.}
\pf
Let $f:X\to Y$ be finite: $Y$ is covered by open affines $V_i=\Spec B_i$ with $f^{-1}(V_i)=\Spec A_i$, $A_i$ a finitely generated $B_i$-module. We check the three conditions in the definition of ``proper''.

\emph{Finite type.} Each $A_i$ is a finitely generated $B_i$-algebra, so $f$ is of finite type.

\emph{Separated.} Put $X_i=f^{-1}(V_i)$. The open sets $X_i\times_{V_i}X_i$ cover $X\times_YX$, and $\Delta^{-1}(X_i\times_{V_i}X_i)=X_i$; the restriction of $\Delta$ is the diagonal of $X_i\to V_i$, a morphism of affine schemes, hence a closed immersion by Prop.~4.1. Being a closed immersion is local on the target, so $\Delta$ is a closed immersion.

\emph{Universally closed.} (i) A finite morphism is closed. By locality on the target (if $W\subseteq X$ is closed then $f(W)\cap V_i=f(W\cap X_i)$) we may assume $X=\Spec A\to Y=\Spec B$ with $A$ a finite $B$-algebra. A closed subset of $X$ is $V(I)$; let $J=\Ker(B\to A/I)$. Then $B/J\hookrightarrow A/I$ is an injective integral extension, so by lying over $\Spec A/I\to\Spec B/J$ is surjective; thus $f(V(I))=V(J)$ is closed.

(ii) Finite morphisms are stable under base extension. Let $Y'\to Y$ be any morphism and $X'=X\times_YY'$. Cover $Y'$ by open affines $V'=\Spec B'$ each mapping into some $V_i$. Then the preimage of $V'$ in $X'$ is $\Spec(A_i\otimes_{B_i}B')$, and $A_i\otimes_{B_i}B'$ is a finite $B'$-module. So $X'\to Y'$ is finite, hence closed by (i).

Thus $f$ is separated, of finite type, and universally closed, i.e.\ proper. \hfill$\square$

% ---------------------------------------------------------------
\exercise{4.2}{Let $S$ be a scheme, let $X$ be a reduced scheme over $S$, and let $Y$ be a separated scheme over $S$. Let $f$ and $g$ be two $S$-morphisms of $X$ to $Y$ which agree on an open dense subset of $X$. Show that $f=g$. Give examples to show that this result fails if either (a) $X$ is nonreduced, or (b) $Y$ is nonseparated. [\emph{Hint}: Consider the map $h:X\to Y\times_SY$ obtained from $f$ and $g$.]}
\pf
\textbf{Lemma.} Let $X$ be reduced, $W$ a scheme, $Z\subseteq W$ a closed subscheme and $h:X\to W$ a morphism with $h(X)\subseteq Z$ as sets. Then $h$ factors through $Z$.

\emph{Proof.} The question is local on $W$ and on $X$, so let $W=\Spec C$, $Z=\Spec C/I$, $X=\Spec A$ with $A$ reduced, $h$ given by $\varphi:C\to A$. For every prime $\mathfrak q$ of $A$ we have $\varphi^{-1}(\mathfrak q)\supseteq I$ (as $h(X)\subseteq Z$), hence $\varphi(I)\subseteq\bigcap\mathfrak q=\Nil(A)=0$. So $\varphi$ factors through $C/I$. (The local factorizations are unique, since $Z\to W$ is a monomorphism, so they glue.) $\square$

\emph{Proof of the statement.} Let $U\subseteq X$ be open dense with $f|_U=g|_U$. Let $h=(f,g)_S:X\to Y\times_SY$. Since $Y$ is separated over $S$, $\Delta:Y\to Y\times_SY$ is a closed immersion; let $Z=\Delta(Y)$ with the corresponding closed subscheme structure, so $\Delta:Y\xrightarrow{\sim}Z$. On $U$ we have $h|_U=(f|_U,f|_U)=\Delta\circ f|_U$, so $h(U)\subseteq Z$. Now $h^{-1}(Z)$ is closed in $X$ and contains the dense set $U$, hence equals $X$. By the Lemma ($X$ is reduced), $h=\Delta\circ k$ for a morphism $k:X\to Y$. Then $f=p_1\circ h=p_1\circ\Delta\circ k=k$ and $g=p_2\circ h=k$, so $f=g$. \hfill$\square$

\emph{Example (a): $X$ nonreduced.} Let $S=\Spec k$, $Y=\mathbf A^1_k=\Spec k[t]$ (affine, hence separated), and $X=\Spec A$ with $A=k[x,\varepsilon]/(x\varepsilon,\varepsilon^2)$. Since $\varepsilon$ is nilpotent, $|X|=|\Spec k[x]|$ is irreducible, so $U=D(x)$ is a nonempty open, hence dense, subset. In $A_x$ we have $\varepsilon=x^{-1}(x\varepsilon)=0$, so $A_x=k[x]_x$. Define $f,g:X\to Y$ by $f^\#(t)=0$ and $g^\#(t)=\varepsilon$. They agree on $U$ (both give $t\mapsto 0$ in $A_x$) but $\varepsilon\neq0$ in $A$, so $f\neq g$.

\emph{Example (b): $Y$ nonseparated.} Let $S=\Spec k$, $X=\mathbf A^1_k$ (reduced), and $Y$ the affine line with the origin doubled (Example~2.3.6), the union of two copies $Y_1,Y_2$ of $\mathbf A^1_k$ glued along $\mathbf A^1_k\setminus\{0\}$. Let $f,g:X\to Y$ be the inclusions of $Y_1$ and $Y_2$. They agree on the dense open set $X\setminus\{0\}$ but $f(0)\neq g(0)$ (the two origins), so $f\ne g$. \hfill$\square$

% ---------------------------------------------------------------
\exercise{4.3}{Let $X$ be a separated scheme over an affine scheme $S$. Let $U$ and $V$ be open affine subsets of $X$. Then $U\cap V$ is also affine. Give an example to show that this fails if $X$ is not separated.}
\pf
Write $S=\Spec A$. Since $S$ is affine, $U\times_SV=\Spec(\mathcal O(U)\otimes_A\mathcal O(V))$ is affine. Let $\Delta:X\to X\times_SX$ be the diagonal; $U\times_SV=p_1^{-1}(U)\cap p_2^{-1}(V)$ is open in $X\times_SX$, and
\[
\Delta^{-1}(U\times_SV)=\Delta^{-1}p_1^{-1}(U)\cap\Delta^{-1}p_2^{-1}(V)=U\cap V,
\]
because $p_1\Delta=p_2\Delta=\mathrm{id}_X$. Since $\Delta$ is a closed immersion, its restriction $U\cap V\to U\times_SV$ is a closed immersion into an affine scheme. A closed subscheme of an affine scheme is affine (it is $\Spec$ of a quotient ring). Hence $U\cap V$ is affine. \hfill$\square$

\emph{Example.} Let $k$ be a field, $S=\Spec k$, and let $X$ be the affine plane with the origin doubled: glue two copies $U=\Spec k[x,y]$ and $V=\Spec k[x,y]$ along $W=\mathbf A^2\setminus\{0\}$ (a nonseparated scheme, as in Example~4.0.1). Both $U,V$ are open affine in $X$ and $U\cap V=W$. We claim $W$ is not affine. Cover $W=D(x)\cup D(y)$. A global section of $\mathcal O_W$ is a pair of elements of $k[x,y]_x$ and $k[x,y]_y$ agreeing in $k[x,y]_{xy}$, i.e.\ an element of $k[x,y]_x\cap k[x,y]_y=k[x,y]$. So $\Gamma(W,\mathcal O_W)=k[x,y]$. If $W$ were affine it would be isomorphic to $\Spec\Gamma(W,\mathcal O_W)=\mathbf A^2$ via the canonical map, which is the open immersion $W\hookrightarrow\mathbf A^2$; but that is not surjective (it misses the origin). So $W=U\cap V$ is not affine. \hfill$\square$

% ---------------------------------------------------------------
\exercise{4.4}{Let $f:X\to Y$ be a morphism of separated schemes of finite type over a noetherian scheme $S$. Let $Z$ be a closed subscheme of $X$ which is proper over $S$. Show that $f(Z)$ is closed in $Y$, and that $f(Z)$ with its image subscheme structure (Ex.~3.11d) is proper over $S$. We refer to this result by saying that ``the image of a proper scheme is proper.'' [\emph{Hint}: Factor $f$ into the graph morphism $\Gamma_f:X\to X\times_SY$ followed by the second projection $p_2$, and show that $\Gamma_f$ is a closed immersion.]}
\pf
\textbf{Step 1: the graph is a closed immersion.} Let $f:X\to Y$ be any $S$-morphism with $Y$ separated over $S$, and $\Gamma_f=(\mathrm{id},f)_S:X\to X\times_SY$. Consider the square with $\Delta_Y:Y\to Y\times_SY$ and $f\times\mathrm{id}_Y:X\times_SY\to Y\times_SY$. For a scheme $T$, a $T$-point of $(X\times_SY)\times_{Y\times_SY}Y$ is a triple $(a_1,a_2,c)$ with $a_1:T\to X$, $a_2,c:T\to Y$ and $(fa_1,a_2)=(c,c)$; that is, $a_2=c=fa_1$. Thus it is the same as a $T$-point $a_1$ of $X$, and the map to $X\times_SY$ is $(a_1,fa_1)=\Gamma_f\circ a_1$. Hence $X\cong(X\times_SY)\times_{Y\times_SY}Y$ with $\Gamma_f$ the base change of $\Delta_Y$. Since $Y$ is separated over $S$, $\Delta_Y$ is a closed immersion, so is $\Gamma_f$ (stability of closed immersions under base change).

\textbf{Step 2.} Let $i:Z\to X$ be the closed immersion and consider $f_Z=f\circ i:Z\to Y$. By Step~1 (applied to $f_Z$), $\Gamma_{f_Z}:Z\to Z\times_SY$ is a closed immersion, and $f_Z=p_2\circ\Gamma_{f_Z}$. Now $Z\to S$ is proper, hence universally closed, so $p_2:Z\times_SY\to Y$ (the base change of $Z\to S$ by $Y\to S$) is a closed map. Thus $f(Z)=p_2(\Gamma_{f_Z}(Z))$ is the image of a closed set under a closed immersion followed by a closed map, hence closed in $Y$.

\textbf{Step 3: properness.} Let $W$ be $f(Z)$ with the image subscheme structure of Ex.~3.11d: the smallest closed subscheme of $Y$ through which $f_Z$ factors. Since $Z$ is noetherian, $f_Z$ is quasi-compact, so the underlying set of $W$ is the closure of $f(Z)$, i.e.\ $f(Z)$ by Step~2, and $f_Z=j\circ\varphi$ with $j:W\to Y$ the closed immersion and $\varphi:Z\to W$ surjective.

\emph{$W\to S$ is separated:} $j$ is a closed immersion, hence separated (for a closed immersion, which is a monomorphism, $\Delta$ is an isomorphism), and $Y\to S$ is separated; a composition of separated morphisms is separated (Cor.~4.6b).

\emph{Finite type:} $j$ is a closed immersion, hence of finite type; $Y\to S$ is of finite type; compose.

\emph{Universally closed:} Let $S'\to S$ be any morphism, $W'=W\times_SS'$, $Z'=Z\times_SS'$, $\varphi':Z'\to W'$. We claim $\varphi'$ is surjective: $Z'=Z\times_WW'$, and for $w'\in W'$ with image $w\in W$ choose $z\in Z$ over $w$ (surjectivity of $\varphi$); the scheme $\Spec(k(z)\otimes_{k(w)}k(w'))$ is nonempty (a tensor product of two nonzero vector spaces over a field is nonzero) and maps to $Z\times_WW'$, giving a point of $Z'$ over $w'$. Now let $C\subseteq W'$ be closed. Then $\varphi'^{-1}(C)$ is closed in $Z'$ and, by surjectivity, its image in $S'$ equals the image of $C$ in $S'$. But $Z'\to S'$ is closed ($Z$ is proper over $S$). So $C\to S'$ has closed image. Hence $W\to S$ is universally closed.

So $W=f(Z)$ is proper over $S$. \hfill$\square$

% ---------------------------------------------------------------
\exercise{4.5}{Let $X$ be an integral scheme of finite type over a field $k$, having function field $K$. We say that a valuation of $K/k$ (see I, Section~6) has \emph{center} $x$ on $X$ if its valuation ring $R$ dominates the local ring $\mathcal O_{x,X}$.
\begin{enumerate}[label=(\alph*),leftmargin=*]
\item If $X$ is separated over $k$, then the center of any valuation of $K/k$ on $X$ (if it exists) is unique.
\item If $X$ is proper over $k$, then every valuation of $K/k$ has a unique center on $X$.
\item (Starred in the book.) Prove the converses of (a) and (b). [\emph{Hint}: While parts (a) and (b) follow quite easily from (4.3) and (4.7), their converses will require some comparison of valuations in different fields.]
\item If $X$ is proper over $k$, and if $k$ is algebraically closed, show that $\Gamma(X,\mathcal O_X)=k$. This result generalizes (I, 3.4a). [\emph{Hint}: Let $a\in\Gamma(X,\mathcal O_X)$, with $a\notin k$. Show that there is a valuation ring $R$ of $K/k$ with $a^{-1}\in\mathfrak m_R$. Then use (b) to get a contradiction.]
\end{enumerate}
}
\pf
\textbf{Conventions.} A valuation of $K/k$ is a valuation ring $R$ of $K$ with $k\subseteq R$ (the trivial one $R=K$ is allowed). Since $X$ is integral, for every $x\in X$ the local ring $\mathcal O_x=\mathcal O_{x,X}$ is a subring of $K$; let $\eta$ be the generic point, $\varepsilon:U_K:=\Spec K\to X$ the canonical morphism with image $\eta$. $X$ is noetherian, so (4.3), (4.7) apply with $Y=\Spec k$. For a valuation ring $R$ put $T=\Spec R$, $U=\Spec K$, $i:U\to T$.

\textbf{Centers as lifts.} \emph{Claim:} the centers of $R$ on $X$ correspond bijectively to the $k$-morphisms $h:T\to X$ with $h\circ i=\varepsilon$ (the center being $h(t_0)$, $t_0$ the closed point). Indeed, apply Lemma~4.4 with $x_1=\eta$: here $Z=\{\eta\}^-=X$ (integral, so reduced), $k(\eta)=K$, and every $x\in X$ is a specialization of $\eta$. By Lemma~4.4, morphisms $h:T\to X$ with $h\circ i=\varepsilon$ correspond to points $x_0$ such that $R$ dominates $\mathcal O_{x_0,X}$ (and the inclusion $k(\eta)=K\subseteq K$ is the identity). Such an $h$ is automatically a $k$-morphism: $T\to X\to\Spec k$ restricted to $U$ is $U\to\Spec k$, and $R\subseteq K$.

\textbf{(a)} Let $x,x'$ be centers of $R$ and $h,h'$ the corresponding morphisms $T\to X$. Then $h\circ i=h'\circ i=\varepsilon$, and $h,h'$ are $k$-morphisms. As $X$ is separated over $k$, Thm~4.3 gives $h=h'$, so $x=h(t_0)=h'(t_0)=x'$.

\textbf{(b)} Let $R$ be a valuation of $K/k$. The data $\varepsilon:U\to X$ and $T\to\Spec k$ form a commutative diagram as in (4.7); since $X$ is proper over $k$, there is $h:T\to X$ with $h\circ i=\varepsilon$ (compatible with $T\to\Spec k$). By the Claim its image $h(t_0)$ is a center. Uniqueness follows from (a), since proper implies separated.

\textbf{(c)} We use two standard facts: (V1) (I,6.1A) every local subring of a field $F$ is dominated by a valuation ring of $F$, and a valuation ring is maximal for domination among local subrings of $F$; (V2) if $F\subseteq L$ are fields and $R''$ is a valuation ring of $L$, then $R''\cap F$ is a valuation ring of $F$ with maximal ideal $\mathfrak m_{R''}\cap F$ (for $0\ne x\in F$, $x$ or $x^{-1}$ lies in $R''$).

\emph{Lemma A.} Let $\xi\in X$, $Z=\{\xi\}^-$ (reduced induced structure; integral with function field $F=k(\xi)$), and let $R$ be a valuation ring of $F$ with $k\subseteq R$. Then there is a valuation ring $S$ of $K$ dominating $\mathcal O_\xi$ and a valuation ring $\tilde R\subseteq S$ of $K/k$ such that for every $x\in Z$:
\[
\tilde R\ \text{dominates}\ \mathcal O_{x,X}\iff R\ \text{dominates}\ \mathcal O_{x,Z}.
\]

\emph{Proof.} By (V1) choose a valuation ring $S$ of $K$ dominating the local subring $\mathcal O_\xi$. Then $F=\mathcal O_\xi/\mathfrak m_\xi\hookrightarrow S/\mathfrak m_S=:\kappa$, so $F\subseteq\kappa$. By (V1) choose a valuation ring $R''$ of $\kappa$ dominating the local ring $R\subseteq F\subseteq\kappa$. By (V2), $R''\cap F$ is a valuation ring of $F$ dominating $R$, so $R''\cap F=R$ by maximality. Let $\tilde R=\{s\in S:\bar s\in R''\}$ ($\bar s$ the image in $\kappa$). It is a subring of $S$ with fraction field $K$, and it is a valuation ring: let $x\in K^*$. If $x\notin S$ then $x^{-1}\in\mathfrak m_S$ so $x^{-1}\in\tilde R$. If $x\in S$ and $\bar x\in R''$ then $x\in\tilde R$; if $x\in S$ and $\bar x\notin R''$, then $\bar x\neq0$, $\bar x^{-1}\in R''$, $x\in S^*$ and $x^{-1}\in\tilde R$. Its maximal ideal is $\{s\in S:\bar s\in\mathfrak m_{R''}\}$. Since $k\subseteq R\subseteq R''$ and $k\subseteq\mathcal O_\xi\subseteq S$ map compatibly, $k\subseteq\tilde R$.

Now let $x\in Z$. Then $\xi$ is a generization of $x$, so $\mathcal O_x\subseteq\mathcal O_\xi\subseteq S$, and with $\mathfrak p_\xi=\mathfrak m_\xi\cap\mathcal O_x$ we have $\mathcal O_{x,Z}=\mathcal O_x/\mathfrak p_\xi$, which is the image of $\nu:\mathcal O_x\to\mathcal O_\xi\to F$; so $\mathcal O_{x,Z}\subseteq F$, $\mathfrak m_{x,Z}=\nu(\mathfrak m_x)$ and $\mathfrak m_x=\nu^{-1}(\mathfrak m_{x,Z})$. For $a\in\mathcal O_x$ the image of $a$ in $\kappa$ is $\nu(a)$ (because $\mathcal O_\xi\to S\to\kappa$ factors through $F$). Hence, using $R''\cap F=R$,
\[
a\in\tilde R\iff\nu(a)\in R,\qquad a\in\mathfrak m_{\tilde R}\iff\nu(a)\in\mathfrak m_{R''}\cap F=\mathfrak m_R.
\]
Therefore $\mathcal O_x\subseteq\tilde R\iff\mathcal O_{x,Z}\subseteq R$ and, in that case, $\mathfrak m_{\tilde R}\cap\mathcal O_x=\nu^{-1}(\mathfrak m_R)$ while $\mathfrak m_x=\nu^{-1}(\mathfrak m_{x,Z})$; as $\nu$ maps onto $\mathcal O_{x,Z}$, these are equal iff $\mathfrak m_R\cap\mathcal O_{x,Z}=\mathfrak m_{x,Z}$. This is the claim. $\square$

\emph{Converse of (a).} Assume every valuation of $K/k$ has at most one center on $X$; we verify the criterion of (4.3) (for $Y=\Spec k$). Let $K'$ be a field, $R'$ a valuation ring of $K'$, $T=\Spec R'$, $U=\Spec K'$, with a $k$-morphism $T\to\Spec k$ (so $k\subseteq R'$), a morphism $U\to X$, and two morphisms $h,h':T\to X$ extending it. Let $\xi$ be the image of $U$, so $h(t_1)=h'(t_1)=\xi$, and $F=k(\xi)\subseteq K'$ (the same inclusion for $h$ and $h'$, as $h|_U=h'|_U$), $Z=\{\xi\}^-$. Put $x_0=h(t_0)$, $x_0'=h'(t_0)\in Z$. By Lemma~4.4, $R'$ dominates $\mathcal O_{x_0,Z}$ and $\mathcal O_{x_0',Z}$. Let $R=R'\cap F$, a valuation ring of $F$ containing $k$ (V2); it dominates $\mathcal O_{x_0,Z}$ and $\mathcal O_{x_0',Z}$ ($\mathfrak m_R\cap\mathcal O=\mathfrak m_{R'}\cap\mathcal O=\mathfrak m_{\mathcal O}$). By Lemma~A the valuation $\tilde R$ of $K/k$ dominates both $\mathcal O_{x_0,X}$ and $\mathcal O_{x_0',X}$, so by hypothesis $x_0=x_0'$. Then $h$ and $h'$ have the same image points $t_0\mapsto x_0$, $t_1\mapsto\xi$ and the same residue field inclusion, so $h=h'$ by Lemma~4.4. By (4.3), $X$ is separated over $k$.

\emph{Converse of (b).} Assume every valuation of $K/k$ has a unique center. By the converse of (a), $X$ is separated; it is of finite type and noetherian. We verify the existence part of (4.7). Let $R'$, $K'$, $T$, $U$ be as above with a $k$-morphism $T\to\Spec k$ and $U\to X$ with image $\xi$, $F=k(\xi)\subseteq K'$, $Z=\{\xi\}^-$, $R=R'\cap F$. Let $S,\tilde R$ be as in Lemma~A. As $\tilde R$ is a valuation of $K/k$, it has a center $x_0\in X$. \emph{Claim: $x_0\in Z$.} Since $\mathcal O_{x_0}\subseteq\tilde R\subseteq S$, the prime $\mathfrak q=\mathfrak m_S\cap\mathcal O_{x_0}$ corresponds to a point $y$ (a generization of $x_0$) with $\mathcal O_y=(\mathcal O_{x_0})_{\mathfrak q}\subseteq S$ and $\mathfrak m_S\cap\mathcal O_y=\mathfrak m_y$; so $S$ (a valuation of $K/k$) has center $y$. But $S$ also has center $\xi$. By uniqueness $y=\xi$, so $x_0\in\{\xi\}^-=Z$. By Lemma~A, $R$ dominates $\mathcal O_{x_0,Z}$; then $R'$ does too ($R'$ dominates $R$). By Lemma~4.4 there is a morphism $h:T\to X$ with $h(t_0)=x_0$, $h(t_1)=\xi$ inducing $k(\xi)\subseteq K'$, i.e.\ with $h|_U$ the given morphism, and $h$ is a $k$-morphism because $R'\subseteq K'$. By (4.7), $X$ is proper over $k$. \hfill(c)

\textbf{(d)} Let $k$ be algebraically closed, $X$ proper over $k$ (so integral, finite type, separated), and $a\in\Gamma(X,\mathcal O_X)$; as $X$ is integral, $a\in K$ and $a\in\mathcal O_x$ for all $x\in X$. Suppose $a\notin k$. Since $k$ is algebraically closed, $a$ is transcendental over $k$ (an algebraic $a$ would lie in $k$). Put $b=a^{-1}$; $k[b]$ is a polynomial ring and $\Lambda=k[b]_{(b)}$ is a local subring of $K$ with maximal ideal $(b)$. By (V1) there is a valuation ring $R$ of $K$ dominating $\Lambda$; it contains $k$, so is a valuation of $K/k$, and $a^{-1}=b\in\mathfrak m_R$. By (b), $R$ has a center $x\in X$: $\mathcal O_x\subseteq R$. But $a\in\mathcal O_x\subseteq R$, so $1=a\cdot a^{-1}\in\mathfrak m_R$, a contradiction. Hence $\Gamma(X,\mathcal O_X)\subseteq k$, and obviously $k\subseteq\Gamma(X,\mathcal O_X)$. \hfill$\square$

% ---------------------------------------------------------------
\exercise{4.6}{Let $f:X\to Y$ be a proper morphism of affine varieties over $k$. Then $f$ is a finite morphism. [\emph{Hint}: Use (4.11A).]}
\pf
We regard the affine varieties as integral affine schemes of finite type over $k$ (Ch.~II, Section~2): $X=\Spec A$, $Y=\Spec B$ with $A,B$ integral domains finitely generated over $k$, and $f$ is given by a $k$-algebra homomorphism $\varphi:B\to A$; ``proper'' refers to the scheme morphism. Let $K=\operatorname{Frac}(A)$ and $B_1=\varphi(B)\subseteq A\subseteq K$.

\emph{Claim: $A$ is integral over $B_1$.} By Thm~4.11A applied to the subring $B_1$ of the field $K$, the integral closure of $B_1$ in $K$ is the intersection of all valuation rings $R$ of $K$ with $B_1\subseteq R$. So it suffices to show $A\subseteq R$ for each such $R$. Let $T=\Spec R$, $U=\Spec K$, $i:U\to T$. We have morphisms $U\to X$ (corresponding to $A\subseteq K$) and $T\to Y$ (corresponding to $B\to B_1\subseteq R$), and the diagram commutes since both composites $U\to Y$ correspond to $B\to B_1\subseteq K$. As $f$ is proper, Thm~4.7 gives $h:T\to X$, i.e.\ a ring homomorphism $\psi:A\to R$ whose composite with $R\subseteq K$ is the inclusion $A\subseteq K$ (this is the commutativity of the upper triangle). Thus $\psi$ is injective with $A\subseteq R$ as subrings of $K$.

Since $A$ is finitely generated as a $k$-algebra, it is finitely generated as a $B_1$-algebra; being also integral over $B_1$, it is a finite $B_1$-module (generated by monomials of bounded degree in integral generators). So $A$ is a finite $B$-module via $\varphi$, i.e.\ $f$ is finite. \hfill$\square$

% ---------------------------------------------------------------
\exercise{4.7}{\emph{Schemes over $\R$.} For any scheme $X_0$ over $\R$, let $X=X_0\times_{\R}\C$. Let $\alpha:\C\to\C$ be complex conjugation, and let $\sigma:X\to X$ be the automorphism obtained by keeping $X_0$ fixed and applying $\alpha$ to $\C$. Then $X$ is a scheme over $\C$, and $\sigma$ is a \emph{semi-linear} automorphism, in the sense that we have a commutative diagram
\[
\begin{array}{ccc}
X&\xrightarrow{\ \sigma\ }&X\\
\downarrow&&\downarrow\\
\Spec\C&\xrightarrow{\ \alpha\ }&\Spec\C.
\end{array}
\]
Since $\sigma^2=\mathrm{id}$, we call $\sigma$ an \emph{involution}.
\begin{enumerate}[label=(\alph*),leftmargin=*]
\item Now let $X$ be a separated scheme of finite type over $\C$, let $\sigma$ be a semilinear involution on $X$, and assume that for any two points $x_1,x_2\in X$, there is an open affine subset containing both. (This last condition is satisfied for example if $X$ is quasi-projective.) Show that there is a unique separated scheme $X_0$ of finite type over $\R$, such that $X_0\times_{\R}\C\cong X$, and such that this isomorphism identifies the given involution of $X$ with the one on $X_0\times_{\R}\C$ described above.
\item For the following statements (b)--(e), $X_0$ will denote a separated scheme of finite type over $\R$, and $X,\sigma$ will denote the corresponding scheme with involution over $\C$. Show that $X_0$ is affine if and only if $X$ is.
\item If $X_0,Y_0$ are two such schemes over $\R$, then to give a morphism $f_0:X_0\to Y_0$ is equivalent to giving a morphism $f:X\to Y$ which commutes with the involutions, i.e., $f\circ\sigma_X=\sigma_Y\circ f$.
\item If $X\cong\mathbf A^1_{\C}$, then $X_0\cong\mathbf A^1_{\R}$.
\item If $X\cong\mathbf P^1_{\C}$, then either $X_0\cong\mathbf P^1_{\R}$, or $X_0$ is isomorphic to the conic in $\mathbf P^2_{\R}$ given by the homogeneous equation $x_0^2+x_1^2+x_2^2=0$.
\end{enumerate}
}
\pf
\textbf{Notation.} $\chi:\C\to\C$ is conjugation, $\alpha=\Spec\chi$. A \emph{semilinear involution} of a $\C$-algebra $A$ (structure map $c:\C\to A$, written $z\mapsto z$) is a ring automorphism $s$ with $s^2=\mathrm{id}$ and $s(za)=\bar zs(a)$; for a semilinear involution $\sigma$ of a $\C$-scheme $X$ we mean $\pi\sigma=\alpha\pi$ ($\pi:X\to\Spec\C$), $\sigma^2=\mathrm{id}$. Then $\sigma$ is an automorphism of the scheme $X$ (not over $\C$), and for $W\subseteq X$ affine open, $\sigma(W)$ is again affine open. The check that the displayed diagram commutes for $X_0\times_{\R}\C$: $\sigma=(p_1,\alpha p_2)$, so $p_2\sigma=\alpha p_2$.

\textbf{Lemma 1 (algebra).} Let $A$ be a $\C$-algebra with a semilinear involution $s$ and $A_0=A^s$ (an $\R$-algebra).
\begin{enumerate}[label=(\roman*),leftmargin=*]
\item $A=A_0\oplus iA_0$, and $A_0\otimes_{\R}\C\to A$, $a\otimes z\mapsto az$, is an isomorphism carrying $\mathrm{id}\otimes\chi$ to $s$. Conversely, for any $\R$-algebra $A_0$ and $A=A_0\otimes_{\R}\C$, $s=\mathrm{id}\otimes\chi$ one has $A^s=A_0$.
\item $A$ is a finite $A_0$-module, so $\rho:\Spec A\to\Spec A_0$ is finite, surjective and closed, and its fibres are exactly the orbits of $s$.
\item For $f\in A_0$, $(A_f)^s=(A_0)_f$.
\item If $A$ is a finitely generated $\C$-algebra, then $A_0$ is a finitely generated $\R$-algebra.
\end{enumerate}
\emph{Proof.} (i) For $a\in A$: $a=\tfrac12(a+sa)+i\cdot\tfrac1{2i}(a-sa)$, and $s\bigl(\tfrac1{2i}(a-sa)\bigr)=\tfrac1{-2i}(sa-a)=\tfrac1{2i}(a-sa)$, so $A=A_0+iA_0$. If $a=ib$ with $a,b\in A_0$ then $a=sa=-ib=-a$, so $a=0$; so the sum is direct. Hence $A_0\otimes_{\R}\C=A_0\otimes1\oplus A_0\otimes i\to A_0\oplus iA_0=A$ is bijective; compatibility with $s$ is clear. The converse is the same decomposition ($s=1$ on $A_0\otimes1$, $-1$ on $A_0\otimes i$). (ii) $A=A_0+A_0i$ is finite over $A_0$, so $\rho$ is finite, hence closed (Ex.~4.1, step (i)); it is surjective by lying over ($A_0\subseteq A$ integral); fibres are orbits of the finite group $\{1,s\}$ acting on $A$ with $A^{s}=A_0$ (Atiyah--Macdonald, Ex.~5.13). (iii) If $a/f^n\in A_f$ is $s$-invariant, then $f^m(sa-a)=0$ for some $m$, so $f^ma\in A_0$ and $a/f^n=f^ma/f^{n+m}\in(A_0)_f$. (iv) Let $a_1,\dots,a_m$ generate $A$ over $\C$, $a_j=u_j+iv_j$ ($u_j,v_j\in A_0$), and $A_0'\subseteq A_0$ the $\R$-subalgebra generated by the $u_j,v_j$. The image of $A_0'\otimes\C\to A_0\otimes\C=A$ contains $\C$ and all $a_j$, so is $A$. As $M\otimes_{\R}\C\cong M\oplus M$ for any $\R$-module $M$, we get $(A_0/A_0')\otimes\C=0$, so $A_0=A_0'$. $\square$

\textbf{Lemma 2.} Let $W=\Spec A$ with a semilinear involution $\sigma$ (ring map $s$), $q:|W|\to|W|/\sigma$ the quotient map, and let $(q_*\mathcal O_W)^\sigma$ be the sheaf of $\sigma$-invariant sections. Then the ringed space $(|W|/\sigma,(q_*\mathcal O_W)^\sigma)$ is isomorphic to $\Spec A_0$, with $q$ corresponding to $\rho$.

\emph{Proof.} By Lemma~1(ii), $\rho$ is continuous, surjective, closed with fibres the $\sigma$-orbits, so it induces a bijection $\bar\rho:|W|/\sigma\to\Spec A_0$, which is a homeomorphism (a closed subset of $|W|/\sigma$ is $q(C)$, $C$ closed and $\sigma$-stable, and $\bar\rho(q(C))=\rho(C)$ is closed). On the basic open sets $D(f)\subseteq\Spec A_0$, $f\in A_0$, we have $\rho^{-1}D(f)=D(f)\subseteq W$ and $\mathcal O_W(D(f))^\sigma=(A_f)^s=(A_0)_f=\mathcal O_{\Spec A_0}(D(f))$ by Lemma~1(iii), compatibly with restrictions. Two sheaves with a compatible isomorphism on a basis are isomorphic. $\square$

\textbf{Lemma 3.} Let $X_0$ be any scheme over $\R$, $X=X_0\times_{\R}\C$, $\sigma=\mathrm{id}\times\alpha$, $q=p_1:X\to X_0$. Then $q\sigma=q$; $q$ is surjective, closed, and its fibres are the $\sigma$-orbits; and $q^\#:\mathcal O_{X_0}\to q_*\mathcal O_X$ is injective with image $(q_*\mathcal O_X)^\sigma$.

\emph{Proof.} Let $V=\Spec B_0\subseteq X_0$ be open affine. Then $q^{-1}(V)=\Spec(B_0\otimes_{\R}\C)$ with $\sigma$ corresponding to $s=\mathrm{id}\otimes\chi$, $A^s=B_0$ (Lemma 1(i)), and $q|$ is the map $\rho$ of Lemma 1(ii). This gives the statements on $q$. For the sheaf statement, on a basic open $D(f)\subseteq V$ ($f\in B_0$) we have $\mathcal O_{X_0}(D(f))=(B_0)_f\to\mathcal O_X(q^{-1}D(f))=A_f$ injective with image $(A_f)^s$ by Lemma 1(iii); such sets form a basis of $X_0$. $\square$

\textbf{(a) Existence.} Let $(X,\sigma)$ be as in (a). Call an open set $\sigma$-\emph{stable} if $\sigma(W)=W$.

\emph{Step 1: $X$ is covered by $\sigma$-stable open affines.} Given $x\in X$, there is an open affine $W\ni x,\sigma(x)$ by hypothesis. Then $\sigma(W)$ is open affine and contains $\sigma(x),x$. Since $X$ is separated over the affine scheme $\Spec\C$, $W'=W\cap\sigma(W)$ is affine (Ex.~4.3); it is $\sigma$-stable ($\sigma^2=\mathrm{id}$) and contains $x$.

\emph{Step 2: the ringed space.} Let $q:|X|\to|X|/\sigma$ be the quotient, and $\mathcal O_0=(q_*\mathcal O_X)^\sigma$, where $\sigma$ acts on $q_*\mathcal O_X$ through $\sigma^\#$ (note $q^{-1}(V)$ is $\sigma$-stable for every open $V$). Put $X_0=(|X|/\sigma,\mathcal O_0)$. If $W=\Spec A$ is a $\sigma$-stable open affine, then $q(W)$ is open ($q^{-1}q(W)=W$), the semilinear involution is $s=\sigma^\#_W$ on $A$ (semilinearity: $\pi|_W$ corresponds to $c:\C\to A$, and $\pi\sigma=\alpha\pi$ says $s(c(z))=c(\bar z)$), and by Lemma~2, $(q(W),\mathcal O_0|)\cong\Spec A^s$. Hence $X_0$ is a scheme, covered by affine schemes $W_0=\Spec A^s$. Because $X$ is quasi-compact (finite type over $\C$), it is covered by finitely many $\sigma$-stable affines $W_1,\dots,W_n$, and by Lemma~1(iv) each $A_j^{s}$ is a finitely generated $\R$-algebra. So $X_0$ is of finite type over $\R$ once we give it an $\R$-structure: the elements of $\R\subseteq\C\to\Gamma(X,\mathcal O_X)$ are $\sigma$-invariant, giving $\R\to\Gamma(X_0,\mathcal O_0)$ and $\psi:X_0\to\Spec\R$.

\emph{Step 3: $q$ is a morphism of schemes $X\to X_0$,} namely $(q,q^\#)$ with $q^\#:\mathcal O_0\hookrightarrow q_*\mathcal O_X$; on $W$ it is $\rho_W:\Spec A\to\Spec A_0$ (Lemma~2). It satisfies $q\sigma=q$ and $\psi q=(\Spec\C\to\Spec\R)\circ\pi$.

\emph{Step 4: $X\cong X_0\times_{\R}\C$.} Let $\theta=(q,\pi):X\to X_0\times_{\R}\C$. The open sets $W_0\times_{\R}\C=\Spec(A_0\otimes_{\R}\C)$ ($W_0=q(W)$) cover $X_0\times_{\R}\C$, $\theta^{-1}(W_0\times\C)=q^{-1}(W_0)=W$, and $\theta|_W$ is $\Spec$ of $A_0\otimes_{\R}\C\to A$, $a\otimes z\mapsto az$, an isomorphism by Lemma~1(i). So $\theta$ is an isomorphism of $\C$-schemes, and $\theta\sigma=(q\sigma,\pi\sigma)=(q,\alpha\pi)=(\mathrm{id}\times\alpha)\theta$: it identifies $\sigma$ with $\mathrm{id}\times\alpha$.

\emph{Separatedness of $X_0$} (needed for (b)--(e)). It suffices to show $\Delta_{X_0}^{-1}(W_{i,0}\times_{\R}W_{j,0})=W_{i,0}\cap W_{j,0}\to W_{i,0}\times_{\R}W_{j,0}$ is a closed immersion. Now $W_{i,0}\cap W_{j,0}=q(W_i\cap W_j)$, and $W_i\cap W_j=\Spec A_{ij}$ is affine, $\sigma$-stable, so $W_{i,0}\cap W_{j,0}=\Spec A_{ij}^s$ by Lemma 2. We need $A_i^s\otimes_{\R}A_j^s\to A_{ij}^s$ surjective. After applying $-\otimes_{\R}\C$ it becomes $A_i\otimes_{\C}A_j\to A_{ij}$ (by Lemma 1(i)), which is surjective because $X$ is separated; and $M\otimes_{\R}\C=0$ forces $M=0$.

\textbf{(a) Uniqueness.} Let $X_0'$ be another scheme over $\R$ with a $\C$-isomorphism $\kappa:X_0'\times_{\R}\C\to X$, $\kappa(\mathrm{id}\times\alpha)=\sigma\kappa$. Put $q'=p_1\kappa^{-1}:X\to X_0'$. Then $q'\sigma=q'$ and by Lemma 3 (transported by $\kappa$) $q'$ is a surjective closed map with fibres the $\sigma$-orbits, and $\mathcal O_{X_0'}\cong(q'_*\mathcal O_X)^\sigma$. Since $q$ has the same properties (Step 3), there is a unique homeomorphism $\phi:|X_0|\to|X_0'|$ with $\phi q=q'$, and an isomorphism of sheaves $\phi^\#:\mathcal O_{X_0'}\to\phi_*\mathcal O_{X_0}$, both being the invariants of $q_*\mathcal O_X$. An isomorphism of ringed spaces between schemes is an isomorphism of schemes; it is compatible with the maps to $\Spec\R$ and with the identifications with $X$. So $X_0'\cong X_0$ (compatibly with the isomorphisms to $X$).

\textbf{(b)} If $X_0=\Spec A_0$ is affine, then $X=\Spec(A_0\otimes_{\R}\C)$ is affine. Conversely if $X=\Spec A$ is affine, then $X$ itself is a $\sigma$-stable open affine, so by Lemma 3 and Lemma 2, $(X_0,\mathcal O_{X_0})\cong(|X|/\sigma,(q_*\mathcal O_X)^\sigma)\cong\Spec A^s$ is affine.

\textbf{(c)} Given $f_0:X_0\to Y_0$ put $f=f_0\times\mathrm{id}_{\C}:X\to Y$. It is a $\C$-morphism and $f\sigma_X=\sigma_Yf$ since $(f_0\times\mathrm{id})(\mathrm{id}\times\alpha)=(\mathrm{id}\times\alpha)(f_0\times\mathrm{id})$.

\emph{Injectivity.} If $f_0\times\mathrm{id}=f_0'\times\mathrm{id}$, then $f_0q_X=q_Yf=f_0'q_X$. As $q_X$ is surjective, $f_0=f_0'$ on points; on sheaves, $f_0^\#:\mathcal O_{Y_0}\to f_{0*}\mathcal O_{X_0}$ composed with the injection $f_{0*}\mathcal O_{X_0}\hookrightarrow f_{0*}q_{X*}\mathcal O_X$ equals the corresponding map for $f_0'$ (both are $(q_Yf)^\#$), so $f_0^\#=f_0'^\#$.

\emph{Surjectivity.} Let $f:X\to Y$ be a $\C$-morphism commuting with the involutions. Then $q_Yf\sigma_X=q_Y\sigma_Yf=q_Yf$ is constant on $\sigma_X$-orbits, so it induces a continuous map $f_0:|X_0|=|X|/\sigma\to|Y_0|$ with $f_0q_X=q_Yf$. For $V_0\subseteq Y_0$ open and $V=q_Y^{-1}V_0$, the map $f^\#:\mathcal O_Y(V)\to\mathcal O_X(f^{-1}V)$ commutes with the involutions (as $f\sigma_X=\sigma_Yf$), hence maps invariants to invariants:
\[
\mathcal O_{Y_0}(V_0)=\mathcal O_Y(V)^\sigma\to\mathcal O_X(f^{-1}V)^\sigma=\mathcal O_{X_0}(f_0^{-1}V_0)
\]
(using $f^{-1}V=q_X^{-1}f_0^{-1}V_0$). This defines $f_0^\#$ and a morphism of ringed spaces. It is local: for $x\in X$, $x_0=q_X(x)$, the stalk maps satisfy $(O_{Y_0,f_0x_0}\to O_{X_0,x_0}\to O_{X,x})=(O_{Y_0,f_0x_0}\to O_{Y,fx}\to O_{X,x})$, which sends the maximal ideal into $\mathfrak m_x$; since $O_{X_0,x_0}\to O_{X,x}$ is a local homomorphism, the preimage of $\mathfrak m_x$ is $\mathfrak m_{x_0}$. So $f_0$ is a morphism of schemes, over $\R$, and $(f_0\times\mathrm{id})=f$ (both are determined by their compositions $q_Yf=f_0q_X$ and $\pi_Yf=\pi_X$ to $Y_0$ and $\Spec\C$). \hfill(c)

(The correspondence is compatible with composition, so isomorphisms correspond to isomorphisms commuting with the involutions. In particular \emph{if $(X,\sigma)\cong(X',\sigma')$ then $X_0\cong X_0'$}.)

\textbf{(d)} Let $X\cong\mathbf A^1_{\C}=\Spec\C[t]$; transporting $\sigma$, we may assume $X=\Spec\C[t]$ and $\sigma$ is given by a semilinear involution $s$ of $\C[t]$ (so $s|_{\C}=\chi$). Let $d=\deg s(t)$; $s(t)$ is nonconstant (else $s(\C[t])\subseteq\C$, not surjective), and $\deg s(p)=d\deg p$ for all $p$ (coefficients are conjugated, the degree multiplies by $d$). Since $s$ is surjective some $s(p)$ has degree 1, so $d=1$: $s(t)=at+b$, $a\neq0$. The relation $s^2(t)=t$ gives $|a|^2=1$ and $\bar ab+\bar b=0$. Choose $c$ with $c^2=a$ (so $|c|=1$, $\bar ca=c$) and put $w=\bar cb$; then $w+\bar w=\bar cb+c\bar b=\bar cb-c\bar ab=\bar cb-\bar cb=0$ (using $\bar b=-\bar ab$ and $c\bar a=\bar c$). So $d_0:=w/2$ satisfies $d_0-\bar d_0=w$. Let $u=ct+d_0$. Then
\[
s(u)=\bar c(at+b)+\bar d_0=ct+(\bar cb+\bar d_0)=ct+d_0=u.
\]
Since $c\neq0$, $\C[t]=\C[u]$ with $u$ transcendental, and $s(\sum a_ku^k)=\sum\bar a_ku^k$. Hence $(\C[t],s)\cong(\R[u]\otimes_{\R}\C,\mathrm{id}\otimes\chi)$, i.e.\ $(X,\sigma)\cong(\mathbf A^1_{\R}\times_{\R}\C,\mathrm{id}\times\alpha)$, and $X_0\cong\mathbf A^1_{\R}$ by (c).

\textbf{(e)} Let $B=\C[x_0,x_1]$, $P=\mathbf P^1_{\C}=\operatorname{Proj}B$. By (c) we may replace $(X,\sigma)$ by $(P,\sigma)$ for an arbitrary $\C$-isomorphism $X\cong P$. A graded ring automorphism $S$ of $B$ (possibly semilinear: $S(zb)=\bar zS(b)$) induces an automorphism $S^*=\operatorname{Proj}S$ of the scheme $P$, with $(SS')^*=S'^*S^*$, and $\pi S^*=\alpha\pi$ if $S$ is semilinear. We use the standard fact $\operatorname{Aut}_{\C}\mathbf P^1_{\C}=PGL_2(\C)$ (II.7.1.1): every $\C$-automorphism is $N^*$ for a linear graded automorphism $N$ (substitution $x_i\mapsto\sum N_{ij}x_j$, $N\in GL_2(\C)$), and $N^*=\mathrm{id}$ iff $N|_{B_1}$ is scalar. Let $\kappa$ be the coefficient conjugation of $B$; $\kappa^*=\mathrm{id}\times\alpha$ on $\mathbf P^1_{\R}\times_{\R}\C=\operatorname{Proj}(\R[x_0,x_1]\otimes_{\R}\C)$.

\emph{Step 1.} If $\sigma$ is a semilinear involution of $P$, then $g=\sigma\kappa^*$ satisfies $\pi g=\alpha\pi\kappa^*=\alpha\alpha\pi=\pi$, so $g=N^*$ is $\C$-linear, hence $\sigma=g\kappa^*=N^*\kappa^*=(\kappa N)^*=:S^*$ with $S=\kappa\circ N$ semilinear graded. Let $L=B_1=\C x_0\oplus\C x_1$ and $\varphi=S|_L$, a semilinear bijection of $L$. Then $(S^2)^*=\sigma^2=\mathrm{id}$ and $S^2$ is linear, so $\varphi^2=\lambda\,\mathrm{id}_L$, $\lambda\in\C^*$. Applying $\varphi$ to $\varphi^2=\lambda$ gives $\varphi^3\ell=\bar\lambda\varphi\ell$ and $\varphi^3\ell=\lambda\varphi\ell$, so $\lambda\in\R$. For $c\in\C^*$, $S_c(f)=c^dS(f)$ for $f\in B_d$ is again a semilinear graded automorphism with $S_c^*=S^*$ and $(S_c|_L)^2=|c|^2\lambda$. So we may assume $\lambda=\pm1$.

\emph{Case $\lambda=1$.} $\varphi$ is a semilinear involution of the $\C$-space $L$, so by Lemma 1(i) (for the module $L$) $L=L^\varphi\oplus iL^\varphi$; take an $\R$-basis $y_0,y_1$ of $L^\varphi$; it is a $\C$-basis of $L$, $B=\C[y_0,y_1]$ and $S(\sum a_Iy^I)=\sum\bar a_Iy^I$. The graded $\C$-isomorphism $B\to B$, $x_i\mapsto y_i$, intertwines $\kappa$ and $S$. Hence $(P,\sigma)\cong(\mathbf P^1_{\R}\times\C,\kappa^*)$ and $X_0\cong\mathbf P^1_{\R}$ by (c).

\emph{Case $\lambda=-1$.} Pick $0\neq v\in L$. Then $v,\varphi v$ are $\C$-independent (if $\varphi v=cv$ then $\varphi^2v=|c|^2v\neq-v$). Put $y_0=v$, $y_1=\varphi v$: $\varphi y_0=y_1$, $\varphi y_1=-y_0$. So $(B,S)$ is isomorphic to $(\C[y_0,y_1],S)$ with $S$ the semilinear automorphism determined by $y_0\mapsto y_1$, $y_1\mapsto-y_0$.

\emph{The model.} Let $Q\subseteq\mathbf P^2_{\R}$ be $\operatorname{Proj}\R[u,v,w]/(q)$, $q=u^2+v^2+w^2$. Then $Q_{\C}=\operatorname{Proj}\C[u,v,w]/(q)$, with $\sigma_Q=\kappa_Q^*$ ($\kappa_Q$ = conjugation of coefficients). Let $B'=\C[s,t]$ and $\varepsilon:\C[u,v,w]/(q)\to B'^{(2)}$ (even-degree part), $u\mapsto\tfrac12(s^2-t^2)$, $v\mapsto\tfrac1{2i}(s^2+t^2)$, $w\mapsto st$. It is well defined: $\tfrac14(s^2-t^2)^2-\tfrac14(s^2+t^2)^2+s^2t^2=0$. It is surjective ($s^2=\varepsilon(u+iv)$, $-t^2=\varepsilon(u-iv)$, $st=\varepsilon(w)$) and injective by dimension count: in degree $d$, $\dim\C[u,v,w]_d/(q)_d=\binom{d+2}2-\binom d2=2d+1=\dim B'_{2d}$. So $\varepsilon$ is a graded isomorphism ($\deg d\mapsto2d$) and $Q_{\C}\cong\operatorname{Proj}B'^{(2)}\cong\operatorname{Proj}B'=\mathbf P^1_{\C}$ (Ex.~5.13, naturally). Let $S'$ be the semilinear graded automorphism of $B'$ with $S'(s)=it$, $S'(t)=-is$. Then $S'(\tfrac12(s^2-t^2))=\tfrac12((it)^2-(-is)^2)=\tfrac12(s^2-t^2)$, $S'(\tfrac1{2i}(s^2+t^2))=\tfrac1{-2i}((it)^2+(-is)^2)=\tfrac1{2i}(s^2+t^2)$, $S'(st)=(it)(-is)=st$. So $S'\varepsilon=\varepsilon\kappa_Q$ (both semilinear ring maps agreeing on generators), and therefore $(Q_{\C},\sigma_Q)\cong(\mathbf P^1_{\C},S'^*)$.

\emph{Comparison.} Put $y_0'=s$, $y_1'=it$ (a $\C$-basis of $B'_1$): $S'(y_0')=y_1'$ and $S'(y_1')=\bar i S'(t)=(-i)(-is)=-y_0'$. So the graded $\C$-isomorphism $\C[y_0,y_1]\to B'$, $y_0\mapsto y_0'$, $y_1\mapsto y_1'$ intertwines $S$ and $S'$. Thus $(P,\sigma)\cong(Q_{\C},\sigma_Q)$ and $X_0\cong Q$ by (c): the conic $x_0^2+x_1^2+x_2^2=0$. \hfill$\square$

(Remark: the two cases are distinct, since $Q(\R)=\emptyset$ while $\mathbf P^1_{\R}$ has real points.)

% ---------------------------------------------------------------
\exercise{4.8}{Let $\mathcal P$ be a property of morphisms of schemes such that:
\begin{enumerate}[label=(\alph*),leftmargin=*]
\item a closed immersion has $\mathcal P$;
\item a composition of two morphisms having $\mathcal P$ has $\mathcal P$;
\item $\mathcal P$ is stable under base extension. Then show that:
\item a product of morphisms having $\mathcal P$ has $\mathcal P$;
\item if $f:X\to Y$ and $g:Y\to Z$ are two morphisms, and if $g\circ f$ has $\mathcal P$ and $g$ is separated, then $f$ has $\mathcal P$;
\item if $f:X\to Y$ has $\mathcal P$, then $f_{\mathrm{red}}:X_{\mathrm{red}}\to Y_{\mathrm{red}}$ has $\mathcal P$.
\end{enumerate}
[\emph{Hint}: For (e), consider the graph morphism $\Gamma_f:X\to X\times_ZY$ and note that it is obtained by base extension from the diagonal morphism $\Delta:Y\to Y\times_ZY$.]}
\pf
\textbf{(d)} Let $f:X\to Y$ and $f':X'\to Y'$ be morphisms of $S$-schemes with $\mathcal P$. Then $f\times_Sf'$ is the composition
\[
X\times_SX'\xrightarrow{\ \mathrm{id}_X\times f'\ }X\times_SY'\xrightarrow{\ f\times\mathrm{id}_{Y'}\ }Y\times_SY'.
\]
The first map is the base extension of $f'$ by $p_2:X\times_SY'\to Y'$ (indeed $X\times_SX'=(X\times_SY')\times_{Y'}X'$), the second is the base extension of $f$ by $p_1:Y\times_SY'\to Y$. Both have $\mathcal P$ by (c), so the composition has $\mathcal P$ by (b).

\textbf{(e)} As in Ex.~4.4 (Step 1, with $S=Z$), $\Gamma_f:X\to X\times_ZY$ is the base extension of $\Delta:Y\to Y\times_ZY$ by $f\times\mathrm{id}_Y:X\times_ZY\to Y\times_ZY$. Since $g$ is separated, $\Delta$ is a closed immersion, so has $\mathcal P$ by (a), hence $\Gamma_f$ has $\mathcal P$ by (c). Next, $p_2:X\times_ZY\to Y$ is the base extension of $g\circ f:X\to Z$ by $g:Y\to Z$; since $g\circ f$ has $\mathcal P$, so does $p_2$ by (c). Finally $f=p_2\circ\Gamma_f$ has $\mathcal P$ by (b).

\textbf{(f)} Let $i:X_{\mathrm{red}}\to X$ and $j:Y_{\mathrm{red}}\to Y$ be the closed immersions. Then $f\circ i=j\circ f_{\mathrm{red}}$. Here $f\circ i$ has $\mathcal P$ by (a), (b). The closed immersion $j$ is separated (it is a monomorphism, so its diagonal is an isomorphism). By (e) applied to $f_{\mathrm{red}}$ and $g=j$, $f_{\mathrm{red}}$ has $\mathcal P$. \hfill$\square$

% ---------------------------------------------------------------
\exercise{4.9}{Show that a composition of projective morphisms is projective. [\emph{Hint}: Use the Segre embedding defined in (I, Ex.~2.14) and show that it gives a closed immersion $\mathbf P^r\times\mathbf P^s\to\mathbf P^{rs+r+s}$.] Conclude that projective morphisms have properties (a)--(f) of (Ex.~4.8) above.}
\pf
Recall $\mathbf P^n_Y=\mathbf P^n_{\Spec\Z}\times_{\Spec\Z}Y$; $f:X\to Y$ is projective if it is a closed immersion $X\to\mathbf P^n_Y$ followed by the projection.

\textbf{Step 1: Segre.} Let $N=rs+r+s$. We construct a closed immersion $\mathbf P^r_{\Z}\times_{\Z}\mathbf P^s_{\Z}\to\mathbf P^N_{\Z}$, where $\mathbf P^r=\operatorname{Proj}\Z[x_0,\dots,x_r]$, $\mathbf P^s=\operatorname{Proj}\Z[y_0,\dots,y_s]$, $\mathbf P^N=\operatorname{Proj}\Z[z_{kl}]$. For $0\le i\le r$, $0\le j\le s$, let $V_{ij}=D_+(x_i)\times D_+(y_j)=\Spec\Z[x_k/x_i,\,y_l/y_j]$ and $W_{ij}=D_+(z_{ij})=\Spec\Z[z_{kl}/z_{ij}]$, and define $\psi_{ij}:V_{ij}\to W_{ij}$ by the ring map
\[
z_{kl}/z_{ij}\longmapsto (x_k/x_i)(y_l/y_j).
\]
It is surjective on rings (it hits $x_k/x_i$ via $z_{kj}/z_{ij}$ and $y_l/y_j$ via $z_{il}/z_{ij}$, which generate). On $V_{ij}$ one has $\psi_{ij}^{-1}D(z_{i'j'}/z_{ij})=D\bigl((x_{i'}/x_i)(y_{j'}/y_j)\bigr)=V_{ij}\cap V_{i'j'}$, and on $V_{ij}\cap V_{i'j'}$ both $\psi_{ij}$ and $\psi_{i'j'}$ send $z_{kl}/z_{i'j'}=(z_{kl}/z_{ij})/(z_{i'j'}/z_{ij})$ to $x_ky_l/(x_{i'}y_{j'})$. So the $\psi_{ij}$ glue to a morphism $\psi:\mathbf P^r\times\mathbf P^s\to\mathbf P^N$ with $\psi^{-1}(W_{ij})=V_{ij}$. Since the $W_{ij}$ cover $\mathbf P^N$ and each $\psi_{ij}$ is a closed immersion (surjective ring map), $\psi$ is a closed immersion (this is local on the target). For an arbitrary base scheme $S$, base change along $S\to\Spec\Z$ gives a closed immersion $\mathbf P^r_S\times_S\mathbf P^s_S\to\mathbf P^N_S$.

\textbf{Step 2: composition.} Let $f:X\to Y$ and $g:Y\to Z$ be projective, with closed immersions $X\hookrightarrow\mathbf P^r_Y$ and $Y\hookrightarrow\mathbf P^s_Z$. Then
\[
X\hookrightarrow\mathbf P^r_Y=\mathbf P^r_{\Z}\times Y\hookrightarrow\mathbf P^r_{\Z}\times\mathbf P^s_Z=\mathbf P^r_Z\times_Z\mathbf P^s_Z\hookrightarrow\mathbf P^N_Z,
\]
where the second map is the base extension of $Y\hookrightarrow\mathbf P^s_Z$ and the third is Step~1. A composition of closed immersions is a closed immersion, and the composite $X\to\mathbf P^N_Z\to Z$ equals $g\circ f$. So $g\circ f$ is projective.

\textbf{Step 3: properties (a)--(c).} (a) A closed immersion $X\to Y=\mathbf P^0_Y$ is projective. (b) is Step~2. (c) If $X\hookrightarrow\mathbf P^n_Y$ and $Y'\to Y$ is any morphism, then $X\times_YY'\hookrightarrow\mathbf P^n_Y\times_YY'=\mathbf P^n_{Y'}$ is a closed immersion (base change), compatible with the projections to $Y'$. So projective morphisms satisfy (a), (b), (c) of Ex.~4.8, and therefore, by Ex.~4.8, also (d), (e), (f) there. \hfill$\square$

% ---------------------------------------------------------------
\exercise{4.10}{(Starred in the book.) \emph{Chow's Lemma.} This result says that proper morphisms are fairly close to projective morphisms. Let $X$ be proper over a noetherian scheme $S$. Then there is a scheme $X'$ and a morphism $g:X'\to X$ such that $X'$ is projective over $S$, and there is an open dense subset $U\subseteq X$ such that $g$ induces an isomorphism of $g^{-1}(U)$ to $U$. Prove this result in the following steps.
\begin{enumerate}[label=(\alph*),leftmargin=*]
\item Reduce to the case $X$ irreducible.
\item Show that $X$ can be covered by a finite number of open subsets $U_i$, $i=1,\dots,n$, each of which is quasi-projective over $S$. Let $U_i\to P_i$ be an open immersion of $U_i$ into a scheme $P_i$ which is projective over $S$.
\item Let $U=\bigcap U_i$, and consider the map
\[
f:U\to X\times_SP_1\times_S\cdots\times_SP_n
\]
deduced from the given maps $U\to X$ and $U\to P_i$. Let $X'$ be the closed image subscheme structure (Ex.~3.11d) $\overline{f(U)}$. Let $g:X'\to X$ be the projection onto the first factor, and let $h:X'\to P=P_1\times_S\cdots\times_SP_n$ be the projection onto the product of the remaining factors. Show that $h$ is a closed immersion, hence $X'$ is projective over $S$.
\item Show that $g^{-1}(U)\to U$ is an isomorphism, thus completing the proof.
\end{enumerate}
}
\pf
\textbf{Conventions.} $X$ is of finite type over the noetherian $S$, hence noetherian, and separated. ``Projective'' and ``quasi-projective'' are as in the text: quasi-projective means an open immersion followed by a projective morphism. By Ex.~4.9 and Ex.~4.8, projective morphisms are stable under composition and base change, closed immersions are projective, and proper schemes are stable under products etc. (Cor.~4.8). All morphisms $U\to Y$ of noetherian schemes are quasi-compact, so for a morphism $f:U\to Y$ the sheaf $\ker(\mathcal O_Y\to f_*\mathcal O_U)$ is quasi-coherent and defines the image subscheme structure of Ex.~3.11d: the smallest closed subscheme through which $f$ factors. This construction is compatible with restriction to open subsets of $Y$.

\textbf{(a) Reduction to $X$ irreducible.} Let $X_1,\dots,X_m$ be the irreducible components of $|X|$ and $W_i=X\setminus\bigcup_{j\neq i}X_j$, a nonempty open subset of $X$, irreducible. Let $\bar X_i\subseteq X$ be the closed subscheme which is the scheme-theoretic closure of the open subscheme $W_i$ (image structure of $W_i\hookrightarrow X$). Then $|\bar X_i|=X_i$, $\bar X_i|_{W_i}=X|_{W_i}$ (the ideal sheaf restricted to $W_i$ vanishes), $\bar X_i$ is irreducible, and proper over $S$ (closed immersion followed by $X\to S$, Cor.~4.8). Suppose the theorem holds for irreducible schemes, giving $g_i:X_i'\to\bar X_i$ projective over $S$ and $U_i\subseteq\bar X_i$ open dense with $g_i^{-1}(U_i)\cong U_i$. Replacing $U_i$ by $U_i\cap W_i$ (still open, dense in $\bar X_i$, since $W_i\neq\emptyset$; and the isomorphism persists over a smaller open set), we may assume $U_i\subseteq W_i$. Then $U_i$ is also an open subscheme of $X$ (as $\bar X_i|_{W_i}=X|_{W_i}$), the $U_i$ are pairwise disjoint (since $W_i\cap W_j=\emptyset$), and $U=\bigcup U_i\cong\bigsqcup U_i$ is open and dense in $X$ ($\overline{U_i}=X_i$). Let $X'=\bigsqcup X_i'$ and $g:X'\to X$ the morphism with $g|_{X_i'}=(\bar X_i\hookrightarrow X)\circ g_i$. Then $g^{-1}(U)\cap X_i'=g_i^{-1}(U\cap X_i)=g_i^{-1}(U_i)\cong U_i$ (points of $U_j$, $j\ne i$, are not in $X_i$), so $g^{-1}(U)\to U$ is an isomorphism.

It remains to see $X'$ is projective over $S$. If $X_i'\hookrightarrow\mathbf P^{n_i}_S$ are closed immersions, let $N+1=\sum(n_i+1)$ and embed $\mathbf P^{n_i}_S$ in $\mathbf P^N_S$ as the linear subspace $\Lambda_i$ where all coordinates outside the $i$-th block of $n_i+1$ coordinates vanish. The $\Lambda_i$ are closed, pairwise disjoint, so $\bigsqcup\Lambda_i\to\mathbf P^N_S$ is a homeomorphism onto a closed set and is surjective on stalks of structure sheaves (near a point of $\Lambda_i$, the other $\Lambda_j$ are absent), i.e.\ a closed immersion. Hence $\bigsqcup X_i'\hookrightarrow\mathbf P^N_S$ is a closed immersion and $X'$ is projective over $S$.

So from now on $X$ is irreducible (and nonempty).

\textbf{(b) Finite cover by quasi-projective opens.} Cover $S$ by finitely many open affines $S_k=\Spec A_k$ and $X$ by finitely many open affines $V=\Spec C$ lying over some $S_k$ (possible as $X$ is quasi-compact, $f^{-1}(S_k)$ noetherian). Here $C$ is a finitely generated $A_k$-algebra, $C=A_k[t_1,\dots,t_n]/J$. So $V\hookrightarrow\mathbf A^n_{S_k}=D_+(x_0)\subseteq\mathbf P^n_{S_k}\subseteq\mathbf P^n_S$ is an immersion (closed in an open subscheme of $\mathbf P^n_S$; here $\mathbf P^n_{S_k}$ is open in $\mathbf P^n_S$). Let $P_V\subseteq\mathbf P^n_S$ be the scheme-theoretic closure of $V$ (image structure of $V\to\mathbf P^n_S$). Since $V$ is closed in an open subscheme $G\subseteq\mathbf P^n_S$ (namely $G=\mathbf A^n_{S_k}$), $P_V\cap G=V$ and so $V\to P_V$ is an open immersion; $P_V\to\mathbf P^n_S\to S$ is a closed immersion followed by the projection, i.e.\ $P_V$ is projective over $S$. Hence each $V$ is quasi-projective over $S$. Take these finitely many $V$ as $U_1,\dots,U_n$ (nonempty ones) with open immersions $\iota_i:U_i\to P_i$, $P_i$ projective over $S$.

\textbf{(c) $h$ is a closed immersion.} Let $U=\bigcap U_i$, nonempty open, so dense in the irreducible $X$. Put $P=P_1\times_S\cdots\times_SP_n$, $Y=X\times_SP$. Each $P_i$ is projective, hence proper, hence separated, over $S$; so $P$ is projective over $S$ by Ex.~4.9, $P\to S$ is separated, and $Y\to S$ and $Y\to X$ are proper (Cor.~4.8). Let $j=(\iota_i|_U)_i:U\to P$ and $f=(\text{incl},j)_S:U\to Y$. Then $f$ is the composition of the graph $\Gamma_j:U\to U\times_SP$, a closed immersion since $P$ is separated (Ex.~4.4, Step~1), with the open immersion $U\times_SP=W_0\hookrightarrow Y$ ($W_0=p_X^{-1}(U)$). Let $X'\subseteq Y$ be the image subscheme of $f$, $g=p_X|_{X'}$ and $h=p_P|_{X'}$.

\emph{Facts.} (F1) $X'\cap W_0=f(U)$ as subschemes: restricting the image subscheme structure to the open set $W_0$ gives that of $f:U\to W_0$, which is a closed immersion, so its image subscheme is $f(U)$ itself. Hence $f:U\to X'$ is an open immersion. (F2) $X'$ is irreducible, being the closure of $f(U)$ with $U$ irreducible; its generic point $\eta'$ lies in the open set $f(U)$. (F3) $X'$ is proper over $S$ (closed in $Y$), so $h$ is proper by Cor.~4.8(e) (here $P\to S$ is separated), in particular $h$ is closed.

\emph{Local structure over $Q_i$.} Let $a_i:X'\to P_i$ be the $i$-th component of $h$, and $Q_i=p_{P_i}^{-1}(U_i)\subseteq P$, open (we identify $U_i$ with its image in $P_i$). We claim
\[
(\ast)\qquad h^{-1}(Q_i)=g^{-1}(U_i)\quad\text{and}\quad h|:g^{-1}(U_i)\to Q_i\ \text{is a closed immersion.}
\]
\emph{First the closed immersion.} Let $Y_i=p_X^{-1}(U_i)=U_i\times_SP$ (open in $Y$) and write $P=P_i\times_SP''$, $P''=\prod_{l\neq i}P_l$ (empty product $=S$), so $Y_i=(U_i\times_SP'')\times_SP_i$. The two morphisms $Y_i\to P_i$, namely the projection $\mathrm{pr}$ to $P_i$ and $\iota_i\circ p_{U_i}$, agree on $f(U)$ (indeed $\mathrm{pr}\,f=\iota_i\circ\text{incl}$). Since $P_i$ is separated over $S$, the equalizer $E_i\subseteq Y_i$ (the pullback of $\Delta_{P_i}$) is a closed subscheme through which $f$ factors. Moreover $E_i$ is the graph of the morphism $\iota_ip_{U_i}:U_i\times_SP''\to P_i$, so projecting gives an isomorphism $E_i\cong U_i\times_SP''$, and under it $h|_{E_i}$ is $\iota_i\times\mathrm{id}:U_i\times_SP''\to P_i\times_SP''=P$, an open immersion with image $Q_i$; so $E_i\cong Q_i$. As $g^{-1}(U_i)=X'\cap Y_i$ is the image subscheme of $f:U\to Y_i$, the minimality of the image structure shows that $g^{-1}(U_i)$ is a closed subscheme of $E_i\cong Q_i$, and $h|$ is the resulting closed immersion into $Q_i$.

\emph{Now the equality of sets.} Clearly $g^{-1}(U_i)\subseteq h^{-1}(Q_i)$ (by what we just saw). Conversely let $x'\in X'$ with $a_i(x')\in U_i$. Let $K'=k(\eta')$ and $Z'=X'_{\mathrm{red}}$ (integral, function field $K'$, $\eta'$ its generic point, and $x'\in Z'$ is a specialization of $\eta'$). Let $\Lambda$ be the local ring of $x'$ on $Z'$, a local subring of $K'$, and $R$ a valuation ring of $K'$ dominating $\Lambda$ (I,6.1A). By Lemma~4.4 we get $\tau:T=\Spec R\to X'$ with $\tau(t_0)=x'$, $\tau(t_1)=\eta'$, inducing the identity on $K'$, so $\tau|_{\Spec K'}$ is the generic point of the open subscheme $f(U)\cong U$. The morphism $a_i\tau:T\to P_i$ sends the closed point $t_0$ into the open set $U_i$; every point of $T$ specializes to $t_0$ and open sets are stable under generization, so $a_i\tau$ factors through $U_i$: $a_i\tau=\iota_ic$ with $c:T\to U_i$. Now $T$ has two morphisms to $X$ extending $\Spec K'\to X$, namely $g\tau$ and $(U_i\subseteq X)\circ c$: on $\Spec K'$, $g\tau$ is the generic point $\Spec K'\to U\subseteq X$, and $a_i\tau|_{\Spec K'}=\iota_i\circ(\Spec K'\to U)$ (since $a_if=\iota_i\circ\mathrm{incl}$), so $c|_{\Spec K'}=(\Spec K'\to U\subseteq U_i)$. Both lifts lie over the same morphism $T\to S$ (all maps involved are $S$-morphisms). As $X$ is separated over the noetherian $S$, Thm~4.3 gives $g\tau=(U_i\subseteq X)\circ c$. Hence $g(x')=g\tau(t_0)\in U_i$, i.e.\ $x'\in g^{-1}(U_i)$. This proves $(\ast)$.

\emph{Conclusion.} The sets $g^{-1}(U_i)$ cover $X'$ (the $U_i$ cover $X$), so $h(X')\subseteq Q:=\bigcup Q_i$, open in $P$, and by $(\ast)$ $h:X'\to Q$ is a closed immersion locally on $Q$ (indeed $h^{-1}(Q_i)\to Q_i$ is a closed immersion for each $i$), hence a closed immersion $X'\to Q$. As $h$ is proper, $h(X')$ is closed in $P$; so $h:X'\to P$ is a homeomorphism onto a closed subset $Z_0$, and $\mathcal O_P\to h_*\mathcal O_{X'}$ is surjective on stalks at points of $Z_0\subseteq Q$ (as it is for $Q$) and the target stalks vanish off $Z_0$. So $h$ is a closed immersion. Since $P$ is projective over $S$ and $h$ is a closed immersion, $X'\to P\to S$ is projective by Ex.~4.9.

\textbf{(d) $g^{-1}(U)\to U$ is an isomorphism.} By (F1), $g^{-1}(U)=X'\cap p_X^{-1}(U)=X'\cap W_0=f(U)$ as open subschemes of $X'$. The composition $U\xrightarrow{f}f(U)=g^{-1}(U)\xrightarrow{g}U$ is the identity (as $p_Xf=\mathrm{incl}$), and $f:U\to f(U)$ is an isomorphism (closed immersion onto its image subscheme). Therefore $g|:g^{-1}(U)\to U$ is an isomorphism. This completes the proof, with $X'$ projective over $S$ and $U$ open dense in $X$. \hfill$\square$

% ---------------------------------------------------------------
\exercise{4.11}{If you are willing to do some harder commutative algebra, and stick to noetherian schemes, then we can express the valuative criteria of separatedness and properness using only \emph{discrete} valuation rings.
\begin{enumerate}[label=(\alph*),leftmargin=*]
\item If $\mathcal O,\mathfrak m$ is a noetherian local domain with quotient field $K$, and if $L$ is a finitely generated field extension of $K$, then there exists a discrete valuation ring $R$ of $L$ dominating $\mathcal O$. Prove this in the following steps. By taking a polynomial ring over $\mathcal O$, reduce to the case where $L$ is a \emph{finite} extension field of $K$. Then show that for a suitable choice of generators $x_1,\dots,x_n$ of $\mathfrak m$, the ideal $\mathfrak a=(x_1)$ in $\mathcal O'=\mathcal O[x_2/x_1,\dots,x_n/x_1]$ is not equal to the unit ideal. Then let $\mathfrak p$ be a minimal prime ideal of $\mathfrak a$, and let $\mathcal O'_{\mathfrak p}$ be the localization of $\mathcal O'$ at $\mathfrak p$. This is a noetherian local domain of dimension~$1$ dominating $\mathcal O$. Let $\tilde{\mathcal O}'_{\mathfrak p}$ be the integral closure of $\mathcal O'_{\mathfrak p}$ in $L$. Use the theorem of Krull--Akizuki (see Nagata [7, p.~115]) to show that $\tilde{\mathcal O}'_{\mathfrak p}$ is a noetherian domain of dimension~$1$. Finally, take $R$ to be a localization of $\tilde{\mathcal O}'_{\mathfrak p}$ at one of its maximal ideals.
\item Let $f:X\to Y$ be a morphism of finite type of noetherian schemes. Show that $f$ is separated (respectively, proper) if and only if the criterion of (4.3) (respectively, (4.7)) holds for all \emph{discrete} valuation rings.
\end{enumerate}
}
\pf
\textbf{(a)} \emph{Caveat.} If $\mathcal O$ is a field and $L$ is algebraic over $K=\mathcal O$, then the only valuation ring of $L$ containing $\mathcal O$ is $L$ itself (an element algebraic over a field is integral over it, hence lies in every valuation ring containing it, as does its inverse), which is not a discrete valuation ring. So we assume (as is intended) that $\mathfrak m\neq0$ or $L$ is not algebraic over $K$.

\emph{Reduction to $L/K$ finite.} Choose a transcendence basis $t_1,\dots,t_r$ of $L/K$ ($L$ is finitely generated); then $L$ is finite over $K(t)=K(t_1,\dots,t_r)$. Let $\mathcal O_1=\mathcal O[t_1,\dots,t_r]_{\mathfrak n}$ with $\mathfrak n=(\mathfrak m,t_1,\dots,t_r)$; $\mathfrak n$ is a maximal ideal, since $\mathcal O[t]/\mathfrak n=\mathcal O/\mathfrak m$. So $\mathcal O_1$ is a noetherian local domain with quotient field $K(t)$, $\mathfrak n\cap\mathcal O=\mathfrak m$ so $\mathcal O_1$ dominates $\mathcal O$, and its maximal ideal is nonzero. A discrete valuation ring of $L$ dominating $\mathcal O_1$ also dominates $\mathcal O$. So we may replace $(\mathcal O,K)$ by $(\mathcal O_1,K(t))$: assume $L/K$ finite and $\mathfrak m\ne0$.

\emph{A good chart.} Let $x_1,\dots,x_n$ be nonzero generators of $\mathfrak m$ (a minimal generating set), and $\mathcal O_i=\mathcal O[x_j/x_i:j\neq i]\subseteq K$; note $\mathfrak m\mathcal O_i=x_i\mathcal O_i$. \emph{Claim: $x_i\mathcal O_i\neq\mathcal O_i$ for some $i$.} Suppose $x_i\mathcal O_i=\mathcal O_i$ for all $i$. Then $1=x_i\sum_{|\beta|\le N}c_\beta\prod_j(x_j/x_i)^{\beta_j}$ for some common $N$ and $c_\beta\in\mathcal O$; multiplying by $x_i^N$ gives $x_i^N=\sum c_\beta x_i^{N+1-|\beta|}\prod x_j^{\beta_j}\in\mathfrak m^{N+1}$. Let $M=n(N-1)+1$. Every monomial of degree $M$ in $x_1,\dots,x_n$ contains some $x_i$ with exponent $\ge N$, so lies in $x_i^N\mathfrak m^{M-N}\subseteq\mathfrak m^{M+1}$. So $\mathfrak m^M=\mathfrak m^{M+1}$, and Nakayama's lemma gives $\mathfrak m^M=0$, so $\mathfrak m=0$ in the domain $\mathcal O$, a contradiction. Renumber so that $i=1$, and put $\mathcal O'=\mathcal O_1$, $\mathfrak a=x_1\mathcal O'\neq\mathcal O'$.

\emph{A one-dimensional local ring.} $\mathcal O'\subseteq K$ is a finitely generated $\mathcal O$-algebra, a noetherian domain with quotient field $K$. Let $\mathfrak p$ be a minimal prime over the nonzero principal ideal $\mathfrak a$. By Krull's principal ideal theorem $\operatorname{ht}\mathfrak p\le1$, and $\mathfrak p\neq0$, so $\mathcal O'_{\mathfrak p}$ is a noetherian local domain of dimension~1 with quotient field $K$. It dominates $\mathcal O$: all $x_j=x_1\cdot(x_j/x_1)\in\mathfrak a\subseteq\mathfrak p$, so $\mathfrak m\subseteq\mathfrak p\cap\mathcal O$, and $\mathfrak p\cap\mathcal O$ is a proper prime, so equals $\mathfrak m$.

\emph{Normalization.} Let $B$ be the integral closure of $\mathcal O'_{\mathfrak p}$ in the finite extension $L$ of $K$. By Krull--Akizuki, $B$ (a ring between $\mathcal O'_{\mathfrak p}$ and $L$) is a noetherian domain of dimension $\le1$; as it is integral over $\mathcal O'_{\mathfrak p}$ it has dimension exactly~$1$. Take a maximal ideal $\mathfrak P$ of $B$; it lies over the maximal ideal of $\mathcal O'_{\mathfrak p}$ (integral extension), so $\mathfrak P\ne0$. Then $R=B_{\mathfrak P}$ is a noetherian local integrally closed domain of dimension~1 with quotient field $L$, hence a discrete valuation ring (I,6.2A). Since $\mathfrak P\cap\mathcal O'_{\mathfrak p}=\mathfrak p\mathcal O'_{\mathfrak p}$, $R$ dominates $\mathcal O'_{\mathfrak p}$, hence $\mathcal O$. \hfill(a)

\textbf{(b)} One direction is trivial: a discrete valuation ring is a valuation ring, so if $f$ is separated (proper), the criterion holds in particular for discrete valuation rings (note that the criterion in (4.7) includes uniqueness, so by (4.3) it implies the criterion of (4.3) for DVRs).

\emph{Separated.} Assume the criterion of (4.3) holds for all discrete valuation rings. $X$ and $X\times_YX$ are noetherian (the latter is of finite type over $Y$), so $\Delta$ is quasi-compact, and by Lemma~4.5 and Cor.~4.2 it suffices to show that $\Delta(X)$ is stable under specialization. Let $\xi_1\in\Delta(X)$ and $\xi_1\rightsquigarrow\xi_0$, $\xi_0\neq\xi_1$ (otherwise nothing to prove). Put $K=k(\xi_1)$ and let $\mathcal O$ be the local ring of $\xi_0$ on $\{\xi_1\}^-$ (reduced structure): a noetherian local domain with quotient field $K$ and $\mathfrak m_{\mathcal O}\neq0$. By (a) with $L=K$ there is a discrete valuation ring $R$ of $K$ dominating $\mathcal O$. Exactly as in the proof of (4.3), Lemma~4.4 gives $h'':T=\Spec R\to X\times_YX$ with $h''(t_0)=\xi_0$, $h''(t_1)=\xi_1$ inducing $K$; $h_1=p_1h''$, $h_2=p_2h''$ are two lifts in the diagram $U=\Spec K\to X$, $T\to Y$ (the restrictions to $U$ agree because $\xi_1\in\Delta(X)$, so $U\to X\times_YX$ factors through $\Delta$). By hypothesis $h_1=h_2$, i.e.\ $h''=\Delta\circ h_1$, so $\xi_0=h''(t_0)\in\Delta(X)$.

\emph{Proper: reduction to noetherian bases.} Assume the criterion of (4.7) holds for all DVRs. Then $f$ is separated by the previous paragraph, and finite type. We show $f$ is universally closed.

\emph{Step 1: closed after base change to noetherian $Y'$.} Let $Y'\to Y$ be a morphism with $Y'$ noetherian, $X'=X\times_YY'$ (noetherian, as $X'\to Y'$ is of finite type), $f':X'\to Y'$, $Z\subseteq X'$ a closed subset with the reduced induced structure. $f'|_Z:Z\to Y'$ is quasi-compact, so by Lemma~4.5 it suffices to show $f'(Z)$ is stable under specialization. Let $z_1\in Z$, $y_1=f'(z_1)\rightsquigarrow y_0$, $y_0\ne y_1$. Let $\mathcal O$ be the local ring of $y_0$ on $\{y_1\}^-\subseteq Y'$ (reduced): a noetherian local domain with quotient field $k(y_1)$, $\mathfrak m_{\mathcal O}\neq0$. Put $K=k(z_1)$; as $z_1$ is a point of a scheme of finite type over $k(y_1)$ (the fibre), $K$ is a finitely generated field extension of $k(y_1)$. By (a) there is a DVR $R$ of $K$ dominating $\mathcal O$. By Lemma~4.4 we get $U=\Spec K\to Z$ (the point $z_1$) and $T=\Spec R\to Y'$ forming a commutative square with $i:U\to T$ and $U\to Y'$. Compose with $Z\to X'\to X$ and $Y'\to Y$: these form a diagram as in (4.7) for the DVR $R$, so there is $h:T\to X$ making it commute. Together with $T\to Y'$ it gives $T\to X'=X\times_YY'$ (restricting to $U\to Z\subseteq X'$ on $U$); the generic point $t_1$ goes to $z_1\in Z$, so $T=\overline{\{t_1\}}$ maps into the closed set $Z$, and as $T$ is reduced, $T\to X'$ factors through $Z$ (Lemma in Ex.~4.2). The image $z_0$ of $t_0$ satisfies $f'(z_0)=y_0$, so $y_0\in f'(Z)$.

\emph{Step 2: arbitrary base change.} We show: if $f'$ is closed for every base change $Y'\to Y$ with $Y'$ noetherian, then $f$ is universally closed. Let $g:Y'\to Y$ be arbitrary and $Z\subseteq X'=X\times_YY'$ closed. Closedness of $f'(Z)$ is local on $Y'$ and we may work over an open affine $\Spec B\subseteq Y$ ($B$ noetherian), so assume $Y=\Spec B$, $Y'=\Spec C$. Let $\{C_\lambda\}$ be the directed family of finitely generated $B$-subalgebras of $C$; each $\Spec C_\lambda$ is noetherian, so $f_\lambda:X_\lambda=X\times_BC_\lambda\to\Spec C_\lambda$ is closed (Step~1). Let $\pi_\lambda:X'\to X_\lambda$, $\rho_\lambda:\Spec C\to\Spec C_\lambda$, and $\bar Z_\lambda=\overline{\pi_\lambda(Z)}$.

(i) $Z=\bigcap_\lambda\pi_\lambda^{-1}(\bar Z_\lambda)$. Take $x\notin Z$. If $X=\bigcup\Spec A_k$ (finitely many affines, $A_k$ finitely generated $B$-algebras), then $X'=\bigcup\Spec(A_k\otimes_BC)$ and $A_k\otimes_BC=\bigcup_\lambda\mathrm{Im}(A_k\otimes_BC_\lambda)$, so every basic open $D(a)$ in $\Spec(A_k\otimes C)$ is $\pi_\lambda^{-1}$ of a basic open of $X_\lambda$ for some $\lambda$ (and for all larger ones). Hence there are $\lambda$ and $V\subseteq X_\lambda$ open with $x\in\pi_\lambda^{-1}(V)\subseteq X'\setminus Z$. Then $V\cap\pi_\lambda(Z)=\emptyset$, so $V\cap\bar Z_\lambda=\emptyset$ and $x\notin\pi_\lambda^{-1}(\bar Z_\lambda)$. (The other inclusion is trivial.)

(ii) Let $y\in\bigcap_\lambda\rho_\lambda^{-1}(f_\lambda(\bar Z_\lambda))$; we show $y\in f'(Z)$. For each $\lambda$ let $F_\lambda=\pi_\lambda^{-1}(\bar Z_\lambda)\cap f'^{-1}(y)$, a closed subset of the fibre $X'_y$. It is nonempty: there is $\zeta\in\bar Z_\lambda$ with $f_\lambda(\zeta)=\rho_\lambda(y)$, and since $X'=X_\lambda\times_{\Spec C_\lambda}\Spec C$, the scheme $\Spec(k(\zeta)\otimes_{k(\rho_\lambda y)}k(y))\ne\emptyset$ gives a point of $X'$ over $\zeta$ and $y$. If $\mu\ge\lambda$ then $\pi_\mu^{-1}(\bar Z_\mu)\subseteq\pi_\lambda^{-1}(\bar Z_\lambda)$ (as $\pi_{\mu\lambda}(\bar Z_\mu)\subseteq\bar Z_\lambda$), so the $F_\lambda$ form a directed family of nonempty closed subsets of the quasi-compact space $X'_y$ (finite type over a field), and $\bigcap F_\lambda\neq\emptyset$. By (i) it lies in $Z\cap f'^{-1}(y)$, so $y\in f'(Z)$.

(iii) Since $f'(Z)\subseteq\rho_\lambda^{-1}f_\lambda(\pi_\lambda Z)\subseteq\rho_\lambda^{-1}f_\lambda(\bar Z_\lambda)$, (ii) gives $f'(Z)=\bigcap_\lambda\rho_\lambda^{-1}(f_\lambda(\bar Z_\lambda))$, an intersection of closed sets ($f_\lambda$ is closed). Hence $f'(Z)$ is closed. So $f$ is universally closed, hence proper. \hfill$\square$

% ---------------------------------------------------------------
\exercise{4.12}{\emph{Examples of Valuation Rings.} Let $k$ be an algebraically closed field.
\begin{enumerate}[label=(\alph*),leftmargin=*]
\item If $K$ is a function field of dimension~1 over $k$ (I, Section~6), then every valuation ring of $K/k$ (except for $K$ itself) is discrete. Thus the set of all of them is just the abstract nonsingular curve $C_K$ of (I, Section~6).
\item If $K/k$ is a function field of dimension two, there are several different kinds of valuations. Suppose that $X$ is a complete nonsingular surface with function field $K$.

(1) If $Y$ is an irreducible curve on $X$, with generic point $x_1$, then the local ring $R=\mathcal O_{x_1,X}$ is a discrete valuation ring of $K/k$ with center at the (nonclosed) point $x_1$ on $X$.

(2) If $f:X'\to X$ is a birational morphism, and if $Y'$ is an irreducible curve in $X'$ whose image in $X$ is a single closed point $x_0$, then the local ring $R$ of the generic point of $Y'$ on $X'$ is a discrete valuation ring of $K/k$ with center at the closed point $x_0$ on $X$.

(3) Let $x_0\in X$ be a closed point. Let $f:X_1\to X$ be the blowing-up of $x_0$ (I, Section~4) and let $E_1=f^{-1}(x_0)$ be the exceptional curve. Choose a closed point $x_1\in E_1$, let $f_2:X_2\to X_1$ be the blowing-up of $x_1$, and let $E_2=f_2^{-1}(x_1)$ be the exceptional curve. Repeat. In this manner we obtain a sequence of varieties $X_i$, with closed points $x_i$ chosen on them, and for each $i$, the local ring $\mathcal O_{x_{i+1},X_{i+1}}$ dominates $\mathcal O_{x_i,X_i}$. Let $R_0=\bigcup_{i=0}^{\infty}\mathcal O_{x_i,X_i}$. Then $R_0$ is a local ring, so it is dominated by some valuation ring $R$ of $K/k$ (I, 6.1A). Show that $R$ is a valuation ring of $K/k$, and that it has center $x_0$ on $X$. When is $R$ a discrete valuation ring?
\end{enumerate}
\emph{Note.} We will see later (V, Ex.~5.6) that in fact the $R_0$ of (3) is already a valuation ring itself, so $R_0=R$. Furthermore, every valuation ring of $K/k$ (except $K$ itself) is one of the three kinds just described.
}
\pf
\textbf{(a)} Let $R\ne K$ be a valuation ring of $K/k$. Then $\mathfrak m_R\neq0$; choose $0\neq a\in\mathfrak m_R$. Since $a\in\mathfrak m_R$, $a$ is not a unit, so $a\notin k^*$; as $k$ is algebraically closed, $a$ is transcendental over $k$. As $K$ has dimension~1 and is finitely generated over $k$, $K$ is a finite extension of $k(a)$. Let $B_0$ be the integral closure of $k[a]$ in $K$: it is a finite $k[a]$-module (I,3.9A), hence a noetherian integrally closed domain of dimension~1 (a Dedekind domain). Since $R\supseteq k[a]$ and $R$ is integrally closed (a valuation ring), $B_0\subseteq R$. Let $\mathfrak p=\mathfrak m_R\cap B_0$, a prime ideal containing $a\ne0$, hence maximal. Then $(B_0)_{\mathfrak p}$ is a noetherian integrally closed local domain of dimension~1, so a discrete valuation ring (I,6.2A), and $(B_0)_{\mathfrak p}\subseteq R\subsetneq K$ (units of $B_0\setminus\mathfrak p$ are units in $R$). A ring $R$ with $S=(B_0)_{\mathfrak p}\subseteq R\subseteq K$ equals $S$ or $K$: if $x\in R\setminus S$, write $x=u\pi^{-m}$ with $m>0$, $u$ a unit of $S$, $\pi$ a uniformizer; then $\pi^{-1}\in R$, so $R\supseteq S[\pi^{-1}]=K$. Since $R\neq K$, $R=(B_0)_{\mathfrak p}$ is a discrete valuation ring.

Conversely every discrete valuation ring of $K/k$ is a valuation ring of $K/k$. So the valuation rings $\ne K$ of $K/k$ are exactly the DVRs of $K/k$, i.e.\ the points of the abstract curve $C_K$ of (I,Section~6). \hfill(a)

\textbf{(b)(1)} Let $\mathcal O_{x_1}=\mathcal O_{x_1,X}$; it is a subring of $K$ with quotient field $K$ containing $k$. Let $P$ be a closed point of $Y$. As $X$ is nonsingular, $\mathcal O_{P,X}$ is a regular local ring (I,5.2A/II.8), and $\mathcal O_{x_1,X}$ is its localization at the prime corresponding to $Y$; a localization of a regular local ring is regular (Serre). Its dimension equals $\operatorname{codim}(Y,X)=1$ (II, Ex.~3.20). So $\mathcal O_{x_1}$ is a one-dimensional regular local ring, i.e.\ a discrete valuation ring (I,6.2A). It contains $k$, has fraction field $K$, so it is a discrete valuation ring of $K/k$, and it dominates itself: its center on $X$ is $x_1$.

\textbf{(2)} Here (as the exercise intends) $X'$ is also a nonsingular surface (otherwise the local ring at the generic point of $Y'$ need not be regular). Since $f$ is birational, $f$ induces an isomorphism $K(X')\cong K(X)=K$ compatible with $f^\#$, and we regard $\mathcal O_{x',X'}\subseteq K$ for points $x'\in X'$. Let $\eta'$ be the generic point of $Y'$. By (1) applied to $X'$, $R=\mathcal O_{\eta',X'}$ is a discrete valuation ring of $K/k$. Since $f(Y')=\{x_0\}$, $f(\eta')=x_0$, and $f^\#:\mathcal O_{x_0,X}\to\mathcal O_{\eta',X'}=R$ is a local homomorphism which, under the identification of function fields, is the inclusion $\mathcal O_{x_0,X}\subseteq R$. A local homomorphism of local rings (the preimage of $\mathfrak m_R$ is $\mathfrak m_{x_0}$) is exactly the statement that $R$ dominates $\mathcal O_{x_0,X}$. So $R$ has center $x_0$ on $X$.

\textbf{(3)} All $X_i$ are integral, and the blowing-up maps $X_{i+1}\to X_i$ are birational, so all function fields are identified with $K$ and all the local rings $\mathcal O_i:=\mathcal O_{x_i,X_i}$ are subrings of $K$. Since $x_{i+1}\mapsto x_i$, the map $\mathcal O_i\to\mathcal O_{i+1}$ is the inclusion and it is local: $\mathfrak m_{i+1}\cap\mathcal O_i=\mathfrak m_i$. So $\mathcal O_0\subseteq\mathcal O_1\subseteq\cdots$ and $\mathfrak m_0\subseteq\mathfrak m_1\subseteq\cdots$.

\emph{$R_0$ is local} with maximal ideal $M=\bigcup\mathfrak m_i$: an element of $R_0$ not in $M$ lies in some $\mathcal O_i\setminus\mathfrak m_i$, hence is a unit there. Elements of $M$ are non-units of $R_0$: if $a\in\mathfrak m_i$ and $a^{-1}\in R_0$, then $a^{-1}\in\mathcal O_l$ for some $l\ge i$ and $1=aa^{-1}\in\mathfrak m_l$, absurd. Also $M\cap\mathcal O_i=\mathfrak m_i$ for every $i$.

Let $R$ be a valuation ring of $K$ dominating $R_0$ (I,6.1A). Then $R$ is a valuation ring of $K$ with $k\subseteq\mathcal O_0\subseteq R_0\subseteq R$, so it is a valuation ring of $K/k$, and $\mathfrak m_R\cap\mathcal O_0=\mathfrak m_R\cap R_0\cap\mathcal O_0=M\cap\mathcal O_0=\mathfrak m_0$. Hence $R\supseteq\mathcal O_{x_0,X}$ and $\mathfrak m_R\cap\mathcal O_{x_0,X}=\mathfrak m_{x_0}$: $R$ has center $x_0$ on $X$.

\emph{When is $R$ discrete?} This depends on the sequence of points $x_i$ (the answer is governed by Zariski's theory of infinitely near points; I only prove the following two facts, which show that both cases occur, and do not attempt a complete classification). Take $X=\mathbf P^2_k$, $x_0=(0,0)$ in the affine chart with coordinates $x,y$. Blowing up $x_0$, the chart of $X_1$ with coordinates $(x,y_1)$, $y=xy_1$ contains $E_1=\{x=0\}$. Inductively, at each stage the chart of $X_n$ has coordinates $(x,y_n)$ with $E_n=\{x=0\}$, and the blow-up of a point $x_n=(0,a_n)$ has chart $(x,y_{n+1})$ with $y_n-a_n=xy_{n+1}$.

\emph{(i) A non-discrete case.} Choose $x_n=(x=0,y_n=0)$ for all $n$ (all $a_n=0$), so $y_n=y/x^n$ and $x_n$ lies on the strict transform of the curve $y=0$. Then $x,\ y_n\in\mathfrak m_n\subseteq M\subseteq\mathfrak m_R$ for all $n$. If $R$ had a discrete valuation $v:K^*\to\mathbf Z$ with $R=\{v\ge0\}$, then $v(x)\ge1$ and $v(y_n)\ge1$ so $v(y)=v(y_n)+nv(x)\ge n+1$ for every $n$, which is impossible. So here $R$ is \emph{not} a discrete valuation ring of rank one (one valuation ring dominating $R_0$ is the composite of $\operatorname{ord}_{\{y=0\}}$ with the valuation of $k(x)$ at $0$, with value group $\mathbf Z^2$ ordered lexicographically; by V, Ex.~5.6 it is $R_0$ itself).

\emph{(ii) A discrete case.} Let $\operatorname{char}k=0$ and let $\varphi=e^x-1=\sum_{m\ge1}a_mx^m\in k[[x]]$, which is transcendental over $k(x)$ (any $\varphi\in xk[[x]]$ transcendental over $k(x)$ would do; here is a proof for $e^x$: let $\sum_{i=0}^dp_i(x)E^i=0$, $E=e^x$, $p_i\in k[x]$, $p_d\neq0$, be a relation with $d\ge1$ minimal. Applying $d/dx$ ($E'=E$) and subtracting $\tfrac{p_d'+dp_d}{p_d}$ times the original gives a relation of degree $<d$, so all its coefficients vanish; for $i=0$ this reads $p_dp_0'-p_d'p_0=dp_dp_0$, and comparing degrees (char $0$) forces $p_0=0$; dividing by $E$ gives a relation of degree $d-1$ with $p_d\neq0$, contradicting minimality). Choose $x_n=(x=0,y_n=a_n)$, where $y_n=(y-p_{n-1}(x))/x^n$, $p_{n-1}=\sum_{m<n}a_mx^m$ (the points of the branch $y=\varphi(x)$). Since $x,\varphi$ are algebraically independent, $y\mapsto\varphi$ extends to an embedding $K=k(x,y)\hookrightarrow k((x))$, giving the discrete valuation $v=\operatorname{ord}_x$ and the DVR $R_v=K\cap k[[x]]$. For each $n$, $\mathcal O_n=k[x,y_n]_{(x,y_n-a_n)}\subseteq R_v$ and $v(y_n-a_n)>0$, $v(x)>0$, so $R_v$ dominates every $\mathcal O_n$, hence $R_0\subseteq R_v$ with domination. In fact $R_0=R_v$: let $f=P/Q\in R_v$ with $P,Q\in k[x,y]$ and put $c=\operatorname{ord}_xQ(x,\varphi)$, so $\operatorname{ord}_xP(x,\varphi)\ge c$. For $n>c$ write $y=p_{n-1}+x^ny_n$: then $Q=Q(x,p_{n-1})+x^nG$, $P=P(x,p_{n-1})+x^nH$ with $G,H\in k[x,y_n]$, and $p_{n-1}\equiv\varphi\pmod{x^n}$ gives $Q(x,p_{n-1})=x^cu(x)$ with $u(0)\neq0$ and $P(x,p_{n-1})=x^cP_1(x)$. Hence $f=(P_1+x^{n-c}H)/(u+x^{n-c}G)\in\mathcal O_n$ (the denominator is a unit of $\mathcal O_n$ since $u(0)\neq0$). So $R_0=R_v$ is a valuation ring, and since $R$ dominates the valuation ring $R_0$ (maximal for domination), $R=R_0=R_v$ is a discrete valuation ring, with center $x_0$.

In general (Zariski) $R$ can also be of rank one with non-discrete value group, or of rank two, according to the behaviour of the infinitely near points $x_i$; this classification belongs to the later chapters (V, Ex.~5.6) and is not proved here. \hfill$\square$

\section*{5. Sheaves of modules}

% ---------------------------------------------------------------- 5.1
\exercise{5.1}{Let $(X,\mathcal O_X)$ be a ringed space, and let $\mathcal E$ be a locally free $\mathcal O_X$-module of finite rank. We define the \emph{dual} of $\mathcal E$, denoted $\mathcal E^\vee$, to be the sheaf $\mathcal{H}om_{\mathcal O_X}(\mathcal E,\mathcal O_X)$.
\begin{enumerate}[label=(\alph*),leftmargin=*]
\item Show that $\mathcal E^{\vee\vee}\cong\mathcal E$.
\item For any $\mathcal O_X$-module $\mathcal F$, $\mathcal{H}om_{\mathcal O_X}(\mathcal E,\mathcal F)\cong\mathcal E^\vee\otimes_{\mathcal O_X}\mathcal F$.
\item For any $\mathcal O_X$-modules $\mathcal F,\mathcal G$, $\operatorname{Hom}_{\mathcal O_X}(\mathcal E\otimes\mathcal F,\mathcal G)\cong\operatorname{Hom}_{\mathcal O_X}(\mathcal F,\mathcal{H}om_{\mathcal O_X}(\mathcal E,\mathcal G))$.
\item (\emph{Projection formula}) If $f:(X,\mathcal O_X)\to(Y,\mathcal O_Y)$ is a morphism of ringed spaces, if $\mathcal F$ is an $\mathcal O_X$-module, and if $\mathcal E$ is a locally free $\mathcal O_Y$-module of finite rank, then there is a natural isomorphism $f_*(\mathcal F\otimes_{\mathcal O_X}f^*\mathcal E)\cong f_*\mathcal F\otimes_{\mathcal O_Y}\mathcal E$.
\end{enumerate}}
\pf
Throughout, $\mathcal{H}om(\mathcal A,\mathcal B)|_U=\mathcal{H}om(\mathcal A|_U,\mathcal B|_U)$, and if $\mathcal A=\mathcal O_U^n$ with basis $e_1,\dots,e_n$ then $\varphi\mapsto(\varphi(e_1),\dots,\varphi(e_n))$ is an isomorphism $\mathcal{H}om(\mathcal O_U^n,\mathcal B)\cong\mathcal B^n$ (a morphism out of a free module is determined by the images of a basis, which can be chosen arbitrarily). In particular $\mathcal E^\vee|_U\cong\mathcal O_U^n$ with the dual basis $e_1^*,\dots,e_n^*$, $e_i^*(e_j)=\delta_{ij}$.

\textbf{(a)} For an open set $U$ define $\theta_U:\mathcal E(U)\to\mathcal E^{\vee\vee}(U)=\operatorname{Hom}(\mathcal E^\vee|_U,\mathcal O_U)$ by $\theta_U(s)(\xi)=\xi_V(s|_V)\in\mathcal O(V)$ for $\xi\in\mathcal E^\vee(V)$, $V\subseteq U$. This is $\mathcal O(U)$-linear, compatible with restrictions, hence a morphism $\theta:\mathcal E\to\mathcal E^{\vee\vee}$ of $\mathcal O_X$-modules. To show it is an isomorphism it suffices to do so on the open sets $U$ over which $\mathcal E|_U\cong\mathcal O_U^n$ (a morphism of sheaves that is an isomorphism locally is an isomorphism). The construction is compatible with restriction to $U$ and natural with respect to isomorphisms of $\mathcal E|_U$, so we may take $\mathcal E|_U=\mathcal O_U^n$ with basis $e_i$. Then $\mathcal E^\vee|_U$ has basis $e_i^*$, and $\mathcal E^{\vee\vee}|_U$ has the dual basis $e_i^{**}$; by definition $\theta(e_i)(e_j^*)=e_j^*(e_i)=\delta_{ij}$, i.e.\ $\theta(e_i)=e_i^{**}$. So $\theta$ maps a basis to a basis and is an isomorphism.

\textbf{(b)} For $\xi\in\mathcal E^\vee(U)$ and $t\in\mathcal F(U)$ let $\mu(\xi\otimes t)\in\mathcal{H}om(\mathcal E,\mathcal F)(U)$ be the morphism $s\mapsto\xi(s)\,t|_V$ ($s\in\mathcal E(V)$, $V\subseteq U$). This is $\mathcal O(U)$-bilinear in $(\xi,t)$ and compatible with restrictions, so it defines a morphism from the presheaf $U\mapsto\mathcal E^\vee(U)\otimes_{\mathcal O(U)}\mathcal F(U)$ to the sheaf $\mathcal{H}om(\mathcal E,\mathcal F)$, hence (universal property of sheafification) a morphism $\mu:\mathcal E^\vee\otimes\mathcal F\to\mathcal{H}om(\mathcal E,\mathcal F)$. It is an isomorphism locally: if $\mathcal E|_U=\mathcal O_U^n$, then $\mathcal E^\vee\otimes\mathcal F|_U$ is the sheaf $\mathcal F^n$ (the presheaf $U'\mapsto \mathcal O(U')^n\otimes\mathcal F(U')=\mathcal F(U')^n$ is already a sheaf), and $\mu(\sum e_i^*\otimes t_i)(e_j)=t_j$, so $\mu$ corresponds to the identification $\mathcal F^n\cong\mathcal{H}om(\mathcal O_U^n,\mathcal F)$ above.

\textbf{(c)} (This holds for every $\mathcal O_X$-module $\mathcal E$.) By the definition of $\mathcal E\otimes\mathcal F$ as a sheafification, a morphism $\psi:\mathcal E\otimes\mathcal F\to\mathcal G$ is the same as a compatible family of $\mathcal O(U)$-linear maps $\psi_U:\mathcal E(U)\otimes_{\mathcal O(U)}\mathcal F(U)\to\mathcal G(U)$, i.e.\ of $\mathcal O(U)$-bilinear maps $B_U:\mathcal E(U)\times\mathcal F(U)\to\mathcal G(U)$ commuting with restrictions. Given such $B$, define $\Phi(B):\mathcal F\to\mathcal{H}om(\mathcal E,\mathcal G)$ by $t\in\mathcal F(U)\mapsto\big(s\in\mathcal E(V)\mapsto B_V(s,t|_V)\big)_{V\subseteq U}$; this is a morphism of $\mathcal O_X$-modules by bilinearity and compatibility with restrictions. Conversely, a morphism $\chi:\mathcal F\to\mathcal{H}om(\mathcal E,\mathcal G)$ gives $B_U(s,t)=\chi_U(t)_U(s)$. These two constructions are inverse to each other and additive, which gives the isomorphism of groups. (Applying this on every open set $U$ to the restrictions gives the isomorphism of sheaves $\mathcal{H}om(\mathcal E\otimes\mathcal F,\mathcal G)\cong\mathcal{H}om(\mathcal F,\mathcal{H}om(\mathcal E,\mathcal G))$.)

\textbf{(d)} Recall $f^*\mathcal E=f^{-1}\mathcal E\otimes_{f^{-1}\mathcal O_Y}\mathcal O_X$. For $e\in\mathcal E(V)$, $V\subseteq Y$ open, let $f^*e\in(f^*\mathcal E)(f^{-1}V)$ be the image of $e$ under $\mathcal E(V)\to(f^{-1}\mathcal E)(f^{-1}V)\to(f^*\mathcal E)(f^{-1}V)$, $e\mapsto 1\otimes e$. Define
\[
B_V:\ (f_*\mathcal F)(V)\times\mathcal E(V)=\mathcal F(f^{-1}V)\times\mathcal E(V)\to(f_*(\mathcal F\otimes f^*\mathcal E))(V),\qquad (t,e)\mapsto t\otimes f^*e .
\]
It is biadditive, and $\mathcal O_Y(V)$-bilinear: for $a\in\mathcal O_Y(V)$ we have $f^*(ae)=f^\#(a)\cdot f^*e$ and the $\mathcal O_Y$-module structure of $f_*\mathcal F$ is $a\cdot t=f^\#(a)t$, so $(at)\otimes f^*e=t\otimes f^\#(a)f^*e=t\otimes f^*(ae)$. The $B_V$ commute with restrictions, so by the argument of (c) they give a morphism $\lambda:f_*\mathcal F\otimes_{\mathcal O_Y}\mathcal E\to f_*(\mathcal F\otimes f^*\mathcal E)$. It is natural in $\mathcal F$ and in $\mathcal E$ (a morphism $\mathcal E\to\mathcal E'$ induces compatible $f^*$).

To see that $\lambda$ is an isomorphism we may work locally on $Y$: for $V\subseteq Y$ open both sides restrict to the corresponding sheaves for $f|_{f^{-1}V}$, and $\lambda$ is compatible with this. So assume $\mathcal E\cong\mathcal O_Y^n$. By naturality in $\mathcal E$ it suffices to treat $\mathcal E=\mathcal O_Y^n$. Then $f^*\mathcal E=\mathcal O_X^n$, $\mathcal F\otimes f^*\mathcal E=\mathcal F^n$, $f_*(\mathcal F^n)=(f_*\mathcal F)^n=f_*\mathcal F\otimes\mathcal O_Y^n$, and $\lambda(t\otimes(a_1,\dots,a_n))=(f^\#(a_k)t)_k=(a_kt)_k$, i.e.\ $\lambda$ is the identity under these identifications. \hfill$\square$

% ---------------------------------------------------------------- 5.2
\exercise{5.2}{Let $R$ be a discrete valuation ring with quotient field $K$, and let $X=\Spec R$.
\begin{enumerate}[label=(\alph*),leftmargin=*]
\item To give an $\mathcal O_X$-module is equivalent to giving an $R$-module $M$, a $K$-vector space $L$, and a homomorphism $\rho:M\otimes_RK\to L$.
\item That $\mathcal O_X$-module is quasi-coherent if and only if $\rho$ is an isomorphism.
\end{enumerate}}
\pf
Let $\pi$ be a uniformizer. $X=\{\eta,s\}$ with $\eta=(0)$ the generic point and $s=(\pi)$ the closed point. The open sets are $\emptyset$, $\{\eta\}=D(\pi)$, and $X$; $\mathcal O_X(X)=R$ and $\mathcal O_X(\{\eta\})=R_\pi=K$.

\textbf{(a)} Let $\mathcal F$ be an $\mathcal O_X$-module. Put $M=\mathcal F(X)$ (an $R$-module), $L=\mathcal F(\{\eta\})$ (a $K$-module) and let $r:M\to L$ be the restriction map. Compatibility of restriction with the module structures means $r(am)=\bar a\,r(m)$, where $\bar a$ is the image of $a\in R$ in $K$; so $r$ is $R$-linear for the $R$-module structure of $L$ coming from $R\subseteq K$, and by the universal property of $\otimes_RK$ it factors uniquely as $r=\rho\circ(m\mapsto m\otimes1)$ with $\rho:M\otimes_RK\to L$ $K$-linear, $\rho(m\otimes c)=c\,r(m)$.

Conversely, given $(M,L,\rho)$, put $\mathcal F(\emptyset)=0$, $\mathcal F(X)=M$, $\mathcal F(\{\eta\})=L$ with restriction $m\mapsto\rho(m\otimes1)$. This is a sheaf: any open covering of $X$ contains $X$ itself (the only open set containing $s$ is $X$), any covering of $\{\eta\}$ contains $\{\eta\}$, and $\mathcal F(\emptyset)=0$; so the sheaf axioms are trivial. It is an $\mathcal O_X$-module, since $r$ is semilinear as required. A morphism $\mathcal F\to\mathcal F'$ is a pair $(\varphi_X:M\to M',\varphi_\eta:L\to L')$ with $\varphi_X$ $R$-linear, $\varphi_\eta$ $K$-linear and $\varphi_\eta r=r'\varphi_X$; equivalently $(\varphi_X,\varphi_\eta)$ is a morphism of triples, i.e.\ $\varphi_\eta\rho=\rho'(\varphi_X\otimes1)$. The two constructions are mutually inverse, so we have an equivalence of categories between $\mathcal O_X$-modules and triples $(M,L,\rho)$.

\textbf{(b)} If $\rho$ is an isomorphism, then $\mathcal F\cong\widetilde M$: indeed $\widetilde M(X)=M$ and $\widetilde M(\{\eta\})=\widetilde M(D(\pi))=M_\pi=M\otimes_RK$ (5.1(c)), with restriction the canonical map $M\to M\otimes K$, i.e.\ $\widetilde M$ corresponds to the triple $(M,M\otimes K,\mathrm{id})$. So $\mathcal F$ is quasi-coherent. Conversely, if $\mathcal F$ is quasi-coherent, then since $X$ is affine $\mathcal F\cong\widetilde N$ for some $R$-module $N$ (the definition together with 5.4, or 5.5), and $N=\Gamma(X,\mathcal F)=M$ by 5.1(d). An isomorphism $\mathcal F\cong\widetilde M$ is an isomorphism of triples $(M,L,\rho)\cong(M,M\otimes K,\mathrm{id})$ and hence $\rho$ is an isomorphism. \hfill$\square$

% ---------------------------------------------------------------- 5.3
\exercise{5.3}{Let $X=\Spec A$ be an affine scheme. Show that the functors $\widetilde{\phantom{M}}$ and $\Gamma$ are adjoint, in the following sense: for any $A$-module $M$, and for any sheaf of $\mathcal O_X$-modules $\mathcal F$, there is a natural isomorphism
\[
\operatorname{Hom}_A(M,\Gamma(X,\mathcal F))\cong\operatorname{Hom}_{\mathcal O_X}(\widetilde M,\mathcal F).
\]}
\pf
Given $\varphi:\widetilde M\to\mathcal F$, let $\Phi(\varphi)=\Gamma(X,\varphi):M=\Gamma(X,\widetilde M)\to\Gamma(X,\mathcal F)$ (5.1(d)); it is $A$-linear.

Conversely let $\psi:M\to\Gamma(X,\mathcal F)$ be $A$-linear. For $f\in A$ define $\psi_f:M_f=\widetilde M(D(f))\to\mathcal F(D(f))$ by $\psi_f(m/f^n)=f^{-n}\cdot\psi(m)|_{D(f)}$, where $f$ is a unit in $\mathcal O_X(D(f))$. This is well defined: if $m/f^n=m'/f^k$ in $M_f$, then $f^N(f^km-f^nm')=0$ for some $N$, so $f^N(f^k\psi(m)-f^n\psi(m'))|_{D(f)}=0$, and multiplying by the unit $f^{-N-n-k}$ gives $f^{-n}\psi(m)|=f^{-k}\psi(m')|$. It is $A_f$-linear, and for $D(g)\subseteq D(f)$ it is compatible with restriction (both composites send $m/f^n$ to the restriction of $f^{-n}\psi(m)|_{D(f)}$ to $D(g)$). Now a section $\sigma\in\widetilde M(U)$ is, after covering $U$ by basic open sets $D(f_i)$, given by $\sigma_i\in M_{f_i}$ that agree on $D(f_if_j)$; the elements $\psi_{f_i}(\sigma_i)\in\mathcal F(D(f_i))$ agree on the overlaps $D(f_if_j)$ by the compatibility just shown, so they glue to a unique element of $\mathcal F(U)$ (independent of the covering). This defines $\widetilde\psi:\widetilde M\to\mathcal F$, a morphism of $\mathcal O_X$-modules, with $\Gamma(\widetilde\psi)=\psi$ (take $f=1$).

Then $\Phi(\widetilde\psi)=\psi$. Conversely, given $\varphi$ with $\psi=\Phi(\varphi)$: on $D(f)$, $m/f^n=f^{-n}\cdot(m|_{D(f)})$ in $\widetilde M(D(f))$ and $\varphi$ is $\mathcal O_X$-linear and commutes with restriction, so $\varphi(m/f^n)=f^{-n}\varphi(m)|_{D(f)}=\psi_f(m/f^n)$. Hence $\varphi$ and $\widetilde\psi$ agree on all $D(f)$, hence everywhere (sections are locally of that form), so $\widetilde{\Phi(\varphi)}=\varphi$. Thus $\Phi$ is bijective. It is clearly additive, and natural in $M$ and in $\mathcal F$ (the formulas above are natural). \hfill$\square$

% ---------------------------------------------------------------- 5.4
\exercise{5.4}{Show that a sheaf of $\mathcal O_X$-modules $\mathcal F$ on a scheme $X$ is quasi-coherent if and only if every point of $X$ has a neighborhood $U$, such that $\mathcal F|_U$ is isomorphic to a cokernel of a morphism of free sheaves on $U$. If $X$ is noetherian, then $\mathcal F$ is coherent if and only if it is locally a cokernel of a morphism of free sheaves of finite rank. (These properties were originally the definition of quasi-coherent and coherent sheaves.)}
\pf
Recall: $\mathcal O_X$ is quasi-coherent (Ex.\ 5.2.1 of the text), $\widetilde{\phantom{M}}$ is exact and commutes with direct sums (5.2(a),(c)), and $\widetilde A=\mathcal O_{\Spec A}$; hence for a free module $A^{(I)}$ we get $\widetilde{A^{(I)}}=\mathcal O^{(I)}$ on $\Spec A$. Kernels, cokernels of morphisms of quasi-coherent sheaves are quasi-coherent (5.7).

($\Rightarrow$) Let $\mathcal F$ be quasi-coherent. Every point has an open affine neighborhood $U=\Spec A$ with $\mathcal F|_U\cong\widetilde M$. Choose a surjection $A^{(I)}\to M$, let $N$ be its kernel and choose a surjection $A^{(J)}\to N$; this gives an exact sequence $A^{(J)}\to A^{(I)}\to M\to0$. Applying $\widetilde{\ }$ gives an exact sequence $\mathcal O_U^{(J)}\to\mathcal O_U^{(I)}\to\mathcal F|_U\to0$. If $X$ is noetherian and $\mathcal F$ is coherent, we may take $U=\Spec A$ with $A$ noetherian and $M$ finitely generated (5.4 of the text), so $I$ can be taken finite; then $N\subseteq A^{|I|}$ is finitely generated since $A$ is noetherian, so $J$ can be taken finite too.

($\Leftarrow$) Suppose $x\in U$ and $\mathcal O_U^{(J)}\xrightarrow{\ \varphi\ }\mathcal O_U^{(I)}\to\mathcal F|_U\to0$ is exact. Free sheaves are quasi-coherent on the scheme $U$, so $\mathcal F|_U=\operatorname{coker}\varphi$ is quasi-coherent by 5.7. Being quasi-coherent means covered by open affines on which the sheaf is some $\widetilde M$; covering each such $U$ by affines $\Spec A\subseteq U$ on which $\mathcal F\cong\widetilde M$, we cover $X$ by such affines, so $\mathcal F$ is quasi-coherent. If $X$ is noetherian and $I,J$ finite, take an open affine $V=\Spec A\subseteq U$, $A$ noetherian. Then $\mathcal F|_V=\widetilde M$ with $M=\Gamma(V,\mathcal F)$ the cokernel of $A^{|J|}\to A^{|I|}$ (by exactness of $\widetilde{\ }$ and 5.5, $\Gamma$ is the inverse equivalence), which is finitely generated. Hence $\mathcal F$ is coherent. \hfill$\square$

% ---------------------------------------------------------------- 5.5
\exercise{5.5}{Let $f:X\to Y$ be a morphism of schemes.
\begin{enumerate}[label=(\alph*),leftmargin=*]
\item Show by example that if $\mathcal F$ is coherent on $X$, then $f_*\mathcal F$ need not be coherent on $Y$, even if $X$ and $Y$ are varieties over a field $k$.
\item Show that a closed immersion is a finite morphism (Sect.~3).
\item If $f$ is a finite morphism of noetherian schemes, and if $\mathcal F$ is coherent on $X$, then $f_*\mathcal F$ is coherent on $Y$.
\end{enumerate}}
\pf
\textbf{(a)} Let $k$ be a field (algebraically closed if one wishes varieties in the sense of Ch.~I), $X=\mathbf A^1_k=\Spec k[t]$, $Y=\Spec k$, and $f$ the structure morphism; both are varieties over $k$. Then $\mathcal F=\mathcal O_X=\widetilde{k[t]}$ is coherent, but $f_*\mathcal F=\widetilde{(k[t])}$ with $k[t]$ viewed as a $k$-module (5.2(d)), i.e.\ the module $k[t]$, which has infinite dimension over $k$. Since coherent on $\Spec k$ means finitely generated, $f_*\mathcal F$ is quasi-coherent but not coherent. (One can also take $f:\mathbf A^2_k\to\mathbf A^1_k$ the projection: $f_*\mathcal O_X=\widetilde{k[x,y]}$ and $k[x,y]$ is not a finitely generated $k[x]$-module.)

\textbf{(b)} Let $i:Z\to X$ be a closed immersion. Cover $X$ by affine opens $V=\Spec B$. Then $i^{-1}(V)\to V$ is a closed immersion, so $i^{-1}(V)$ is a closed subscheme of $V$, and by Cor.~5.10 it is isomorphic over $V$ to $\Spec B/\fa$ for an ideal $\fa$ of $B$. Thus $i^{-1}(V)$ is affine and $B/\fa$ is a finitely generated $B$-module (generated by $1$). So $i$ is finite.

\textbf{(c)} The question is local on $Y$ ($f_*\mathcal F|_V=(f|_{f^{-1}V})_*(\mathcal F|_{f^{-1}V})$). By definition of finiteness we can cover $Y$ by open affines $V=\Spec A$ such that $f^{-1}V=\Spec B$ with $B$ a finitely generated $A$-module. Since $Y$ is noetherian, $A$ is noetherian, and $B$ (a finitely generated module, hence finitely generated algebra, over $A$) is noetherian. As $\mathcal F|_{f^{-1}V}$ is coherent on $\Spec B$, $\mathcal F|_{f^{-1}V}=\widetilde N$ with $N$ a finitely generated $B$-module. By 5.2(d), $f_*\widetilde N=({}_AN)^{\sim}$, and ${}_AN$ is a finitely generated $A$-module since $N$ is finitely generated over $B$ and $B$ over $A$. So $f_*\mathcal F|_V$ is coherent. \hfill$\square$

% ---------------------------------------------------------------- 5.6
\exercise{5.6}{\emph{Support.} Recall the notions of support of a section of a sheaf, support of a sheaf, and subsheaf with supports from (Ex.\ 1.14) and (Ex.\ 1.20).
\begin{enumerate}[label=(\alph*),leftmargin=*]
\item Let $A$ be a ring, let $M$ be an $A$-module, let $X=\Spec A$, and let $\mathcal F=\widetilde M$. For any $m\in M=\Gamma(X,\mathcal F)$, show that $\operatorname{Supp}m=V(\Ann m)$, where $\Ann m=\{a\in A\mid am=0\}$.
\item Now suppose that $A$ is noetherian, and $M$ finitely generated. Show that $\operatorname{Supp}\mathcal F=V(\Ann M)$.
\item The support of a coherent sheaf on a noetherian scheme is closed.
\item For any ideal $\fa\subseteq A$, we define a submodule $\Gamma_\fa(M)$ of $M$ by $\Gamma_\fa(M)=\{m\in M\mid\fa^nm=0\text{ for some }n>0\}$. Assume that $A$ is noetherian, and $M$ any $A$-module. Show that $\Gamma_\fa(M)^{\sim}\cong\mathcal H^0_Z(\mathcal F)$, where $Z=V(\fa)$ and $\mathcal F=\widetilde M$. [\emph{Hint}: Use (Ex.\ 1.20) and (5.8) to show a priori that $\mathcal H^0_Z(\mathcal F)$ is quasi-coherent. Then show that $\Gamma_\fa(M)\cong\Gamma_Z(\mathcal F)$.]
\item Let $X$ be a noetherian scheme, and let $Z$ be a closed subset. If $\mathcal F$ is a quasi-coherent (respectively, coherent) $\mathcal O_X$-module, then $\mathcal H^0_Z(\mathcal F)$ is also quasi-coherent (respectively, coherent).
\end{enumerate}}
\pf
Recall $\operatorname{Supp}s=\{x:s_x\ne0\}$ for a section $s$, $\operatorname{Supp}\mathcal F=\{x:\mathcal F_x\neq0\}$, and $\mathcal H^0_Z(\mathcal F)(U)=\{s\in\mathcal F(U):\operatorname{Supp}s\subseteq Z\}$ (a subsheaf of $\mathcal F$).

\textbf{(a)} The germ of $m$ at $\fp$ is $m/1\in M_\fp=(\widetilde M)_\fp$. It is zero iff there is $a\notin\fp$ with $am=0$, i.e.\ iff $\Ann m\not\subseteq\fp$. Hence $\operatorname{Supp}m=\{\fp:\Ann m\subseteq\fp\}=V(\Ann m)$.

\textbf{(b)} Let $m_1,\dots,m_r$ generate $M$. Then $M_\fp\neq0$ iff some $m_i/1\neq0$ in $M_\fp$, so $\operatorname{Supp}\mathcal F=\bigcup\operatorname{Supp}m_i=\bigcup V(\Ann m_i)=V(\bigcap\Ann m_i)$ (finite union), and $\bigcap_i\Ann m_i=\Ann M$ since $M$ is generated by the $m_i$.

\textbf{(c)} Let $\mathcal F$ be coherent on the noetherian scheme $X$. Cover $X$ by finitely many open affines $U_i=\Spec A_i$ with $\mathcal F|_{U_i}=\widetilde M_i$, $M_i$ finitely generated. By (b), $\operatorname{Supp}\mathcal F\cap U_i=\operatorname{Supp}(\mathcal F|_{U_i})=V(\Ann M_i)$ is closed in $U_i$ (stalks are local). A subset whose intersection with each $U_i$ is closed in $U_i$ is closed in $X$ (its complement is open in each $U_i$).

\textbf{(d)} \emph{Step 1: $\mathcal H^0_Z(\mathcal F)$ is the kernel of $\mathcal F\to j_*(\mathcal F|_W)$,} where $W=X\setminus Z$, $j:W\to X$ is the inclusion and the map is the restriction $\mathcal F(V)\to\mathcal F(V\cap W)$. Indeed for $s\in\mathcal F(V)$, $\operatorname{Supp}s\subseteq Z$ iff all germs of $s$ at points of $V\cap W$ vanish iff $s|_{V\cap W}=0$.

Now let $X$ be any noetherian scheme (in particular $X=\Spec A$ here) and $\mathcal F$ quasi-coherent. The scheme $W$ is noetherian and $\mathcal F|_W$ is quasi-coherent, so $j_*(\mathcal F|_W)$ is quasi-coherent by 5.8(c) (hypothesis: the source $W$ is noetherian). Hence $\mathcal H^0_Z(\mathcal F)$, the kernel of a morphism of quasi-coherent sheaves, is quasi-coherent by 5.7. (This proves the quasi-coherent part of (e) too.)

\emph{Step 2.} Now $X=\Spec A$, $\mathcal F=\widetilde M$, $Z=V(\fa)$. Since $\mathcal H^0_Z(\mathcal F)$ is quasi-coherent on the affine scheme $X$, it equals $\widetilde N$ with $N=\Gamma(X,\mathcal H^0_Z(\mathcal F))=\{m\in M:\operatorname{Supp}m\subseteq Z\}$ (5.5), as a submodule of $M$ (and $\widetilde N\to\widetilde M$ is the inclusion). By (a), $\operatorname{Supp}m\subseteq V(\fa)$ iff $V(\Ann m)\subseteq V(\fa)$ iff $\fa\subseteq\sqrt{\Ann m}$. Since $A$ is noetherian, $\fa=(a_1,\dots,a_s)$ is finitely generated; $\fa\subseteq\sqrt{\Ann m}$ means each $a_i^{n_i}\in\Ann m$, which holds iff $\fa^nm=0$ for some $n$ (if $a_i^{n_i}m=0$ then $\fa^{n}m=0$ for $n\geq\sum(n_i-1)+1$, since $\fa^n$ is generated by monomials of degree $n$ in the $a_i$, each containing some $a_i^{n_i}$; the converse is clear). Thus $N=\Gamma_\fa(M)$ and $\mathcal H^0_Z(\widetilde M)=\Gamma_\fa(M)^{\sim}$.

\textbf{(e)} Quasi-coherence was shown in Step 1. If $\mathcal F$ is coherent, let $U=\Spec A$ be an open affine of $X$ with $A$ noetherian, $\mathcal F|_U=\widetilde M$, $M$ finitely generated. Then $\mathcal H^0_Z(\mathcal F)|_U=\mathcal H^0_{Z\cap U}(\mathcal F|_U)=\Gamma_\fa(M)^{\sim}$ by (d), where $Z\cap U=V(\fa)$. Now $\Gamma_\fa(M)$ is a submodule of the finitely generated module $M$ over the noetherian ring $A$, hence finitely generated. So $\mathcal H^0_Z(\mathcal F)$ is coherent. \hfill$\square$

% ---------------------------------------------------------------- 5.7
\exercise{5.7}{Let $X$ be a noetherian scheme, and let $\mathcal F$ be a coherent sheaf.
\begin{enumerate}[label=(\alph*),leftmargin=*]
\item If the stalk $\mathcal F_x$ is a free $\mathcal O_x$-module for some point $x\in X$, then there is a neighborhood $U$ of $x$ such that $\mathcal F|_U$ is free.
\item $\mathcal F$ is locally free if and only if its stalks $\mathcal F_x$ are free $\mathcal O_x$-modules for all $x\in X$.
\item $\mathcal F$ is invertible (i.e., locally free of rank 1) if and only if there is a coherent sheaf $\mathcal G$ such that $\mathcal F\otimes\mathcal G\cong\mathcal O_X$. (This justifies the terminology \emph{invertible}: it means that $\mathcal F$ is an invertible element of the monoid of coherent sheaves under the operation $\otimes$.)
\end{enumerate}}
\pf
\textbf{(a)} Let $\mathcal F_x\cong\mathcal O_x^n$ ($n$ is finite as $\mathcal F_x$ is finitely generated). Take an affine open neighborhood $\Spec A$ of $x$, $A$ noetherian, with $\mathcal F=\widetilde M$, $M$ finitely generated, and $x=\fp$. Then $M_\fp\cong A_\fp^n$. Pick a basis of $M_\fp$; each basis element has the form $m_i/u_i$ with $m_i\in M$, $u_i\notin\fp$, and $u_i$ is a unit in $A_\fp$, so the images of $m_1,\dots,m_n$ also form a basis. Let $\varphi:A^n\to M$, $e_i\mapsto m_i$. Then $\varphi_\fp$ is an isomorphism. The modules $K=\ker\varphi$ and $C=\operatorname{coker}\varphi$ are finitely generated ($A$ noetherian) and $K_\fp=C_\fp=0$ (localization is exact). A finitely generated module with zero localization at $\fp$ is killed by some element not in $\fp$ (take the product of annihilating elements for the generators). So there is $s\notin\fp$ with $K_s=C_s=0$, i.e.\ $\varphi_s:A_s^n\to M_s$ is an isomorphism, and $\mathcal F|_{D(s)}=\widetilde{M_s}\cong\mathcal O_{D(s)}^n$.

\textbf{(b)} If $\mathcal F|_U\cong\mathcal O_U^{(I)}$ then $\mathcal F_x\cong\mathcal O_x^{(I)}$ for $x\in U$, so stalks are free. The converse is (a).

\textbf{(c)} ($\Rightarrow$) Let $\mathcal F$ be invertible and $\mathcal G=\mathcal F^\vee=\mathcal{H}om(\mathcal F,\mathcal O_X)$, which is locally $\cong\mathcal O$, hence coherent. The evaluation morphism $\mathcal F\otimes\mathcal F^\vee\to\mathcal O_X$, $s\otimes\varphi\mapsto\varphi(s)$ is an isomorphism: on $U$ with $\mathcal F|_U=\mathcal O_Ue$, $\mathcal F^\vee|_U=\mathcal O_Ue^*$ and $e\otimes e^*\mapsto1$.

($\Leftarrow$) Suppose $\mathcal F\otimes\mathcal G\cong\mathcal O_X$ with $\mathcal F,\mathcal G$ coherent. Fix $x$, $R=\mathcal O_x$ with maximal ideal $\fm$, $k=R/\fm$. Then $F=\mathcal F_x$, $G=\mathcal G_x$ are finitely generated and $F\otimes_RG\cong R$ (stalk of a tensor product). Tensoring with $k$: $(F/\fm F)\otimes_k(G/\fm G)\cong k$, so both have dimension $1$. By Nakayama $F$ and $G$ are cyclic: $F\cong R/I$, $G\cong R/J$, so $R\cong F\otimes G\cong R/(I+J)$, which forces $I+J=0$ (annihilators), hence $I=0$ and $F\cong R$. By (b), $\mathcal F$ is locally free, of rank $1$. \hfill$\square$

% ---------------------------------------------------------------- 5.8
\exercise{5.8}{Again let $X$ be a noetherian scheme and $\mathcal F$ a coherent sheaf on $X$. We will consider the function
\[
\varphi(x)=\dim_{k(x)}\mathcal F_x\otimes_{\mathcal O_x}k(x),
\]
where $k(x)=\mathcal O_x/\fm_x$ is the residue field at the point $x$. Use Nakayama's lemma to prove the following results.
\begin{enumerate}[label=(\alph*),leftmargin=*]
\item The function $\varphi$ is \emph{upper semi-continuous}, i.e., for any $n\in\mathbf Z$, the set $\{x\in X\mid\varphi(x)\ge n\}$ is closed.
\item If $\mathcal F$ is locally free, and $X$ is connected, then $\varphi$ is a constant function.
\item Conversely, if $X$ is reduced, and $\varphi$ is constant, then $\mathcal F$ is locally free.
\end{enumerate}}
\pf
\textbf{(a)} We show: every $x$ has a neighborhood on which $\varphi\le\varphi(x)$. Then $\{\varphi<n\}$ is open for every $n$ (if $\varphi(x)<n$ the neighborhood lies in $\{\varphi<n\}$), i.e.\ $\{\varphi\ge n\}$ is closed.

Let $\varphi(x)=r$. Take an affine neighborhood $\Spec A$ of $x$, $A$ noetherian, $\mathcal F=\widetilde M$, $M$ finitely generated, $x=\fp$. The $k(\fp)$-vector space $M_\fp/\fp M_\fp$ has dimension $r$; choose $m_1,\dots,m_r\in M$ whose images span it. By Nakayama ($M_\fp$ is a finitely generated $A_\fp$-module) the $m_i$ generate $M_\fp$. Then $N=M/\sum Am_i$ is finitely generated with $N_\fp=0$, so $N_s=0$ for some $s\notin\fp$. Thus on $D(s)$ the module $M_s$ is generated by the $m_i$, so for every $\fq\in D(s)$, $M_\fq$ is generated by $r$ elements, whence $\varphi(\fq)=\dim M_\fq/\fq M_\fq\le r$.

\textbf{(b)} If $\mathcal F|_U\cong\mathcal O_U^{(I)}$, then $\mathcal F_y\otimes k(y)=k(y)^{(I)}$ for $y\in U$; as $\mathcal F_y$ is finitely generated $I$ is finite and $\varphi=|I|$ on $U$. So $\varphi$ is locally constant, hence constant on the connected space $X$.

\textbf{(c)} Let $\varphi\equiv r$. Fix $x$ and, as in (a), an affine neighborhood $\Spec A$ ($A$ noetherian and reduced since $X$ is reduced), $\mathcal F=\widetilde M$, $x=\fp$, and $m_1,\dots,m_r\in M$ generating $M_\fp$ and, after shrinking to some $D(s)$ and replacing $A$ by $A_s$, $M$ by $M_s$, generating $M$ (see (a)). So we have a surjection $\pi:A^r\to M$ with kernel $K\subseteq A^r$. Let $\fq$ be any prime of $A$. Tensoring $A_\fq^r\to M_\fq\to0$ with $k(\fq)$ gives a surjection $k(\fq)^r\to M_\fq\otimes k(\fq)$ between vector spaces of the same dimension $r$ (because $\varphi(\fq)=r$), which is therefore an isomorphism. From the exact sequence $K_\fq\to A_\fq^r\to M_\fq\to0$ and right exactness of $\otimes k(\fq)$, the image of $K_\fq\otimes k(\fq)\to k(\fq)^r$ is the kernel of this isomorphism, i.e.\ zero; this says $K_\fq\subseteq\fq A_\fq^r$.

Let $\kappa=(a_1,\dots,a_r)\in K\subseteq A^r$. For every prime $\fq$ the image of $\kappa$ in $A_\fq^r$ lies in $\fq A_\fq^r$, so $a_i/1\in\fq A_\fq$, i.e.\ $a_i\in\fq$. Hence $a_i\in\bigcap\fq=\Nil(A)=0$ ($A$ reduced). So $K=0$, $M\cong A^r$ is free, and $\mathcal F$ is free in a neighborhood of $x$. \hfill$\square$

% ---------------------------------------------------------------- 5.9
\exercise{5.9}{Let $S$ be a graded ring, generated by $S_1$ as an $S_0$-algebra, let $M$ be a graded $S$-module, and let $X=\operatorname{Proj}S$.
\begin{enumerate}[label=(\alph*),leftmargin=*]
\item Show that there is a natural homomorphism $\alpha:M\to\Gamma_*(\widetilde M)$.
\item Assume now that $S_0=A$ is a finitely generated $k$-algebra for some field $k$, that $S_1$ is a finitely generated $A$-module, and that $M$ is a finitely generated $S$-module. Show that the map $\alpha$ is an isomorphism in all large enough degrees, i.e., there is a $d_0\in\mathbf Z$ such that for all $d\ge d_0$, $\alpha_d:M_d\to\Gamma(X,\widetilde M(d))$ is an isomorphism. [\emph{Hint}: Use the methods of the proof of (5.19).]
\item With the same hypotheses, we define an equivalence relation $\approx$ on graded $S$-modules by saying $M\approx M'$ if there is an integer $d$ such that $M_{\ge d}\cong M'_{\ge d}$. Here $M_{\ge d}=\bigoplus_{n\ge d}M_n$. We will say that a graded $S$-module $M$ is \emph{quasi-finitely generated} if it is equivalent to a finitely generated module. Now show that the functors $\widetilde{\ }$ and $\Gamma_*$ induce an equivalence of categories between the category of quasi-finitely generated graded $S$-modules modulo the equivalence relation $\approx$, and the category of coherent $\mathcal O_X$-modules.
\end{enumerate}}
\pf
\textbf{(a)} By 5.12(b), $\widetilde M(n)=\widetilde M\otimes\mathcal O_X(n)\cong\widetilde{M(n)}$. For $m\in M_n=M(n)_0$ let $\alpha(m)$ be the global section of $\widetilde{M(n)}$ whose value at $\fp$ is $m/1\in M(n)_{(\fp)}$; it is a section (locally of the form $m/f$ with $f=1$). This defines $\alpha_n:M_n\to\Gamma(X,\widetilde M(n))$, hence $\alpha:M\to\Gamma_*(\widetilde M)=\bigoplus_{n\in\mathbf Z}\Gamma(X,\widetilde M(n))$, additive. It is $S$-linear: $s\in S_d$ acts on $\Gamma_*$ through the section $s/1\in\Gamma(\mathcal O_X(d))$ and the maps $\mathcal O_X(d)\otimes\widetilde M(n)\to\widetilde M(n+d)$, which on stalks at $\fp$ send $(s/1)\otimes(m/1)$ to $sm/1$; so $\alpha(sm)=s\alpha(m)$. It is natural in $M$: a graded homomorphism $\varphi:M\to N$ induces $\Gamma_*(\widetilde\varphi)\circ\alpha_M=\alpha_N\circ\varphi$. Also, if $M\to M'$ is an inclusion of a graded submodule $M_{\ge d}$ etc.\ the corresponding $\alpha$'s are compatible by naturality.

\medskip
\emph{A lemma (the method of 5.19).} Let $T$ be a graded ring which is an integral domain, with $T_0$ a finitely generated algebra over a field $k$, $T$ generated over $T_0$ by $T_1$, $T_1$ a finitely generated $T_0$-module, $T_1\neq0$. Choose nonzero generators $x_0,\dots,x_m$ of the $T_0$-module $T_1$, put $X_T=\operatorname{Proj}T$ (covered by the $D_+(x_i)$ since the $x_i$ generate $T_+$), $K=\operatorname{Frac}T$. Then $T$ is a finitely generated $k$-algebra, hence noetherian. Let $L\subseteq K$ be the localization of $T$ at the nonzero homogeneous elements; it is a graded ring and each $T_{x_i}$ is a graded subring of $L$. For $n\ge0$ let $T'_n=\{s\in L_n: s\in T_{x_i}\text{ for all }i\}$ and $T'=\bigoplus_{n\ge0}T'_n$, a graded subring of $L\subseteq K$ containing $T$.
\begin{enumerate}[label=(\roman*),leftmargin=*]
\item $\Gamma(X_T,\widetilde{T}(n))=T'_n$ for $n\ge0$ (multiplication matching the product in $\Gamma_*$), and $\alpha_n:T_n\to T'_n$ is the inclusion.
\item For $s\in T'_d$ ($d\ge0$) there is $N$ such that $y\,s\in T$ for all $y\in T_{\ge N}$.
\item $T'$ is integral over $T$ and is a finitely generated $T$-module.
\item $T'_n=T_n$ for all $n\gg0$.
\end{enumerate}
\emph{Proof.} (i) By 5.12(b) and 5.11(b), the sections of $\widetilde T(n)=\widetilde{T(n)}$ over $D_+(x_i)$ are $T(n)_{(x_i)}=\{a/x_i^k:a\in T_{n+k}\}=(T_{x_i})_n$, and over $D_+(x_ix_j)$ they are $(T_{x_ix_j})_n$. All rings $T_{x_i}\to T_{x_ix_j}$ are subrings of $K$ ($T$ is a domain and $x_i\ne0$), so by the sheaf axiom a global section is a tuple $(t_i)$ with $t_i\in (T_{x_i})_n$ and $t_i=t_j$ in $K$, i.e.\ an element of $\bigcap_i(T_{x_i})_n=T'_n$. The map $\alpha$ sends $t\in T_n$ to the tuple $(t,\dots,t)$.

(ii) Write $s=a_i/x_i^{k_i}$ with $a_i\in T$; so $x_i^{n_0}s\in T$ for $n_0=\max k_i$ and all $i$. Let $N=(m+1)(n_0-1)+1$. Any monomial of degree $e\ge N$ in $x_0,\dots,x_m$ is divisible by some $x_i^{n_0}$ (pigeonhole), and $T_e$ is spanned over $T_0$ by such monomials. Hence $ys\in T$ for $y\in T_e$, $e\ge N$.

(iii) Let $s\in T'_d$ ($d\ge0$) and $N$ as in (ii). For $y\in T_{\ge N}$, $ys\in T$ and $\deg(ys)\ge N$, so by induction $ys^q\in T$ for all $q\ge1$. With $y=x_0^N\neq0$ we get $T[s]\subseteq y^{-1}T$, a finitely generated $T$-module ($T$ is noetherian), so $s$ is integral over $T$. As $T'$ is generated by such homogeneous elements it is integral over $T$, hence contained in the integral closure $\overline T$ of $T$ in $K$, which is a finite $T$-module (I, 3.9A, since $T$ is a finitely generated $k$-algebra). So $T'$, a submodule, is finitely generated.

(iv) Let $u_1,\dots,u_t$ be homogeneous generators of $T'$ over $T$, $D=\max\deg u_j$, and by (ii) $N'$ with $T_{\ge N'}u_j\subseteq T$ for all $j$. For $n\ge D+N'$: $T'_n=\sum_jT_{n-\deg u_j}u_j\subseteq T$, so $T'_n=T_n$. $\square$

\medskip
\textbf{(b)} $S$ is a finitely generated $A$-algebra, hence a finitely generated $k$-algebra and noetherian. Let $P(M)$ be the assertion ``$\alpha_d$ is bijective for all $d\gg0$''.

\emph{Domain case.} Let $\fp\subseteq S$ be a homogeneous prime and $T=S/\fp$. Then $T$ satisfies the hypotheses of the lemma, except possibly $T_1\neq0$. If $T_1=0$, then $T=T_0$, $\operatorname{Proj}T=\emptyset$, $\widetilde T=0$ and $T_d=0$ for $d\ge1$, so $P(T)$ holds. If $T_1\ne0$, then by the lemma (i), (iv) $\alpha_d:T_d\to\Gamma(\widetilde T(d))$ is bijective for $d\gg0$. (Here $\operatorname{Proj}T$ is the closed subscheme $j:\operatorname{Proj}T\hookrightarrow X$ defined by $\fp$; for $f\in S_1$ we have $(S/\fp)_{(f)}=T_{(\bar f)}$ and $j^{-1}D_+(f)=D_+(\bar f)$, so the $\mathcal O_X$-module $\widetilde{S/\fp}$ is $j_*$ of the sheaf $\widetilde T$ on $\operatorname{Proj}T$, and likewise for all twists; hence the global sections of $(\widetilde{S/\fp})(d)$ over $X$ are those of $\widetilde T(d)$ over $\operatorname{Proj}T$, and under this identification $\alpha_d$ for $S/\fp$ is $\alpha_d$ for $T$, since both send $m$ to the section with value $m/1$.) For the twist $T(l)$, $\alpha_d$ for $T(l)$ is $\alpha_{d+l}$ for $T$, so $P(T(l))$ holds too.

\emph{Extensions.} Let $0\to M'\to M\to M''\to0$ be an exact sequence of graded modules with $P(M')$ and $P(M'')$. Twisting the exact sequence $0\to\widetilde{M'}\to\widetilde M\to\widetilde{M''}\to0$ by the invertible sheaf $\mathcal O(d)$ (exactness of $\widetilde{\ }$ follows from exactness of localization on stalks $M_{(\fp)}$) and applying the left exact functor $\Gamma$ gives a commutative diagram with exact rows
\[
\begin{array}{ccccccccc}
0&\to&M'_d&\to&M_d&\to&M''_d&\to&0\\
&&\downarrow\alpha'&&\downarrow\alpha&&\downarrow\alpha''&&\\
0&\to&\Gamma(\widetilde{M'}(d))&\to&\Gamma(\widetilde M(d))&\to&\Gamma(\widetilde{M''}(d))&&
\end{array}
\]
For $d\gg0$, $\alpha'$ and $\alpha''$ are isomorphisms. Then $\alpha$ is injective (if $\alpha(m)=0$ then $\alpha''(\bar m)=0$, so $\bar m=0$, $m\in M'_d$ and $\alpha'(m)=0$, so $m=0$). It is surjective: given $t\in\Gamma(\widetilde M(d))$, its image in $\Gamma(\widetilde{M''}(d))$ equals $\alpha''(\bar m)$ for some $m\in M_d$, so $t-\alpha(m)$ maps to $0$, hence comes from $\Gamma(\widetilde{M'}(d))=\alpha'(M'_d)$, say $t-\alpha(m)=\alpha(m')$ with $m'\in M'_d\subseteq M_d$, and $t=\alpha(m+m')$. Thus $P(M)$ holds.

\emph{General case.} For $M$ finitely generated over the noetherian graded ring $S$, there is a filtration $0=M^0\subseteq\dots\subseteq M^r=M$ by graded submodules with $M^i/M^{i-1}\cong(S/\fp_i)(l_i)$, $\fp_i$ homogeneous primes (I, 7.4). By the domain case and induction on $r$ (using the extension step), $P(M)$ holds.

\textbf{(c)} Interpretation: the category $\mathcal C$ has as objects the quasi-finitely generated graded $S$-modules, and $\operatorname{Hom}_{\mathcal C}(M,N)=\varinjlim_d\operatorname{Hom}_{S\text{-gr}}(M_{\ge d},N_{\ge d})$ (graded homomorphisms of degree $0$, transition maps by restriction). Two objects are isomorphic in $\mathcal C$ iff $M\approx N$.

\emph{Step 1.} For any graded $M$ and $d$, $\widetilde{M_{\ge d}}\to\widetilde M$ is an isomorphism: for $\fp\in X$ choose $x\in S_1\setminus\fp$ (exists as $S_+$ is generated by $S_1$); a degree zero element $m/f$ of $(M/M_{\ge d})_{(\fp)}$ equals $x^dm/(x^df)=0$ since $x^dm\in M_{\ge d}$. Hence if $M\approx N$ with $N$ finitely generated, then $\widetilde M\cong\widetilde N$ is coherent (5.11(c), $S$ noetherian), and by naturality of $\alpha$ and (b), $\alpha_d(M)$ is an isomorphism for all $d\ge d_1(M)$.

\emph{Step 2: the functor.} A morphism in $\mathcal C$ represented by $\varphi:M_{\ge d}\to N_{\ge d}$ gives $\widetilde\varphi:\widetilde M=\widetilde{M_{\ge d}}\to\widetilde N$, independent of the representative. So $M\mapsto\widetilde M$ is a functor $\mathcal C\to\operatorname{Coh}(X)$.

\emph{Step 3: fully faithful.} Let $\Theta:\operatorname{Hom}_{\mathcal C}(M,N)\to\operatorname{Hom}(\widetilde M,\widetilde N)$ be the map. Let $d_1$ be such that $\alpha_n(M)$, $\alpha_n(N)$ are isomorphisms for $n\ge d_1$. \emph{Injective:} if $\widetilde\varphi=0$ for $\varphi:M_{\ge d}\to N_{\ge d}$, then $\alpha_N\varphi=\Gamma_*(\widetilde\varphi)\alpha_M=0$, so $\varphi=0$ in degrees $\ge d_1$. \emph{Surjective:} given $\psi:\widetilde M\to\widetilde N$, put $\varphi=\alpha_N^{-1}\circ\Gamma_*(\psi)\circ\alpha_M:M_{\ge d_1}\to N_{\ge d_1}$, a graded $S$-homomorphism. We check $\widetilde\varphi=\psi$ on $D_+(f)$, $f\in S_1$ (these cover $X$): $\widetilde M(D_+(f))=M_{(f)}\ni m/f^k$, and we may assume $m\in M_k$, $k\ge d_1$. The section $f$ of $\mathcal O(1)$ trivializes $\mathcal O(1)$ on $D_+(f)$, and under the resulting identification $\widetilde M(k)|_{D_+(f)}\cong\widetilde M|_{D_+(f)}$ the section $\alpha(m)$ corresponds to $m/f^k$; $\psi$ commutes with these identifications, so $\psi(m/f^k)$ corresponds to $\Gamma_*(\psi)(\alpha(m))=\alpha(\varphi(m))$, which corresponds to $\varphi(m)/f^k=\widetilde\varphi(m/f^k)$.

\emph{Step 4: essentially surjective.} Let $\mathcal F$ be coherent on $X$. As $S_1$ is a finitely generated $A$-module, $X$ is a closed subscheme of $\mathbf P^r_A$ and $\mathcal O_X(1)$ is the restriction of $\mathcal O_{\mathbf P^r}(1)$ (5.12(c)), so it is very ample; by Serre's theorem (5.17) there are $n$ and finitely many $s_1,\dots,s_t\in\Gamma(X,\mathcal F(n))$ generating $\mathcal F(n)$. Let $M=\Gamma_*(\mathcal F)$ and let $N\subseteq M$ be the graded submodule generated by the $s_j\in M_n$; $N$ is finitely generated. By 5.15, $\widetilde M\cong\mathcal F$. The image of $\widetilde N(n)\to\widetilde M(n)\cong\mathcal F(n)$ is an $\mathcal O_X$-submodule that contains all the $s_j$, hence is all of $\mathcal F(n)$; since $\widetilde N\to\widetilde M$ is injective, $\widetilde N(n)\cong\mathcal F(n)$, and $\widetilde N\cong\mathcal F$. So $\mathcal F\cong\widetilde N$ with $N$ finitely generated, and $N\in\mathcal C$. (Moreover $M=\Gamma_*(\mathcal F)\approx N$: since $\widetilde N=\widetilde M$, the map $\alpha_d(N):N_d\to\Gamma(X,\widetilde N(d))=\Gamma(X,\mathcal F(d))=M_d$ is the inclusion $N_d\subseteq M_d$, and it is bijective for $d\gg0$ by (b) applied to the finitely generated module $N$. So $\Gamma_*(\mathcal F)$ is quasi-finitely generated and, by Step 3, $\Gamma_*$ is a quasi-inverse to $\widetilde{\ }$.)

Hence $\widetilde{\ }:\mathcal C\to\operatorname{Coh}(X)$ is fully faithful and essentially surjective, i.e.\ an equivalence, with quasi-inverse $\Gamma_*$. \hfill$\square$

% ---------------------------------------------------------------- 5.10
\exercise{5.10}{Let $A$ be a ring, let $S=A[x_0,\dots,x_r]$ and let $X=\operatorname{Proj}S$. We have seen that a homogeneous ideal $I$ in $S$ defines a closed subscheme of $X$ (Ex.\ 3.12), and that conversely every closed subscheme of $X$ arises in this way (5.16).
\begin{enumerate}[label=(\alph*),leftmargin=*]
\item For any homogeneous ideal $I\subseteq S$, we define the \emph{saturation} $\bar I$ of $I$ to be $\{s\in S\mid\text{for each }i=0,\dots,r,\text{ there is an }n\text{ such that }x_i^ns\in I\}$. We say that $I$ is \emph{saturated} if $I=\bar I$. Show that $\bar I$ is a homogeneous ideal of $S$.
\item Two homogeneous ideals $I_1$ and $I_2$ of $S$ define the same closed subscheme of $X$ if and only if they have the same saturation.
\item If $Y$ is any closed subscheme of $X$, then the ideal $\Gamma_*(\mathcal I_Y)$ is saturated. Hence it is the largest homogeneous ideal defining the subscheme $Y$.
\item There is a 1-1 correspondence between saturated ideals of $S$ and closed subschemes of $X$.
\end{enumerate}}
\pf
(Throughout $r\ge1$, as in 5.13 and 5.16; this is needed for $\Gamma_*(\mathcal O_X)=S$. For $r=0$ one has $\Gamma_*(\mathcal O_X)=A[x_0,x_0^{-1}]\neq S$, so (c) is not meaningful as stated; we do not treat that case.)

\textbf{(a)} Clearly $I\subseteq\bar I$. If $s,t\in\bar I$ with $x_i^ns\in I$, $x_i^mt\in I$, then $x_i^{\max(n,m)}(s+t)\in I$; if $u\in S$ then $x_i^n(us)\in I$. So $\bar I$ is an ideal. If $s=\sum s_d$ (homogeneous components) lies in $\bar I$ and $x_i^ns\in I$, then all components $x_i^ns_d$ lie in $I$ because $I$ is homogeneous; taking the maximum of the finitely many exponents needed (for each $i$, over the components $s_d$), we get $s_d\in\bar I$. So $\bar I$ is homogeneous.

\textbf{Two facts.} For a homogeneous ideal $I$ let $Y_I\subseteq X$ be the closed subscheme it defines, with ideal sheaf $\widetilde I\subseteq\widetilde S=\mathcal O_X$ (proof of 5.16(a)); a closed subscheme is determined by its ideal sheaf (5.9).

\emph{Fact 1: $\widetilde I=\widetilde{\bar I}$.} The module $\bar I/I$ has $(\bar I/I)_{(\fp)}=0$ for every $\fp\in X$: choose $i$ with $x_i\notin\fp$ ($\fp\not\supseteq S_+$); a degree-$0$ element of $T^{-1}(\bar I/I)$ is $\bar s/f$ with $s\in\bar I$ homogeneous, and $x_i^ns\in I$ for some $n$, so $\bar s/f=\overline{x_i^ns}/(x_i^nf)=0$. As $\widetilde{\ }$ is exact, $\widetilde{\bar I}/\widetilde I=0$.

\emph{Fact 2: $\Gamma_*(\widetilde I)=\bar I$} (as submodules of $\Gamma_*(\mathcal O_X)=S$, 5.13). Since $\widetilde{I(n)}\subseteq\widetilde{S(n)}$, a homogeneous $s\in S_n$ gives a section of $\widetilde I(n)$ iff its restriction to each $D_+(x_i)$ lies in $I(n)_{(x_i)}\subseteq S(n)_{(x_i)}$ (5.11(b)), i.e.\ iff $s\in I_{x_i}$ for all $i$, i.e.\ iff for each $i$ there is $n_i$ with $x_i^{n_i}s\in I$, i.e.\ iff $s\in\bar I$. Both sides are graded, so they agree.

\textbf{(b)} If $I_1,I_2$ have the same saturation, then $\widetilde I_1=\widetilde{\bar I_1}=\widetilde{\bar I_2}=\widetilde I_2$ (Fact 1), so they define the same closed subscheme. Conversely, if $Y_{I_1}=Y_{I_2}$, then $\widetilde I_1=\widetilde I_2$ and by Fact 2, $\bar I_1=\Gamma_*(\widetilde I_1)=\Gamma_*(\widetilde I_2)=\bar I_2$.

\textbf{(c)} Let $Y\subseteq X$ be closed and $I=\Gamma_*(\mathcal I_Y)$, a homogeneous ideal of $S$. By 5.9 $\mathcal I_Y$ is quasi-coherent, and by 5.15 $\widetilde I=\widetilde{\Gamma_*(\mathcal I_Y)}\cong\mathcal I_Y$ (the isomorphism being compatible with the inclusion into $\mathcal O_X$), so $I$ defines $Y$. By Fact 2, $\bar I=\Gamma_*(\widetilde I)=\Gamma_*(\mathcal I_Y)=I$: $I$ is saturated. If $J$ is any homogeneous ideal defining $Y$ (i.e.\ $\widetilde J=\mathcal I_Y$), then $J\subseteq\bar J=\Gamma_*(\widetilde J)=\Gamma_*(\mathcal I_Y)=I$. So $I$ is the largest.

\textbf{(d)} Map $I\mapsto Y_I$ from saturated homogeneous ideals to closed subschemes. Injective by (b): saturated $I_1,I_2$ defining the same subscheme satisfy $I_1=\bar I_1=\bar I_2=I_2$. Surjective by (c): $Y=Y_{I}$ with $I=\Gamma_*(\mathcal I_Y)$ saturated. The inverse is $Y\mapsto\Gamma_*(\mathcal I_Y)$. \hfill$\square$

% ---------------------------------------------------------------- 5.11
\exercise{5.11}{Let $S$ and $T$ be two graded rings with $S_0=T_0=A$. We define the \emph{Cartesian product} $S\times_AT$ to be the graded ring $\bigoplus_{d\ge0}S_d\otimes_AT_d$. If $X=\operatorname{Proj}S$ and $Y=\operatorname{Proj}T$, show that $\operatorname{Proj}(S\times_AT)\cong X\times_AY$, and show that the sheaf $\mathcal O(1)$ on $\operatorname{Proj}(S\times_AT)$ is isomorphic to $p_1^*(\mathcal O_X(1))\otimes p_2^*(\mathcal O_Y(1))$ on $X\times Y$.

The Cartesian product of rings is related to the Segre embedding of projective spaces (I, Ex.\ 2.14) in the following way. If $x_0,\dots,x_r$ is a set of generators for $S_1$ over $A$, corresponding to a projective embedding $X\hookrightarrow\mathbf P^r_A$, and if $y_0,\dots,y_s$ is a set of generators for $T_1$, corresponding to a projective embedding $Y\hookrightarrow\mathbf P^s_A$, then $\{x_i\otimes y_j\}$ is a set of generators for $(S\times_AT)_1$, and hence defines a projective embedding $\operatorname{Proj}(S\times_AT)\hookrightarrow\mathbf P^N_A$, with $N=rs+r+s$. This is just the image of $X\times Y\subseteq\mathbf P^r\times\mathbf P^s$ in $\mathbf P^N$ under its Segre embedding.}
\pf
\emph{Convention.} As in 5.12 (which is needed for $\mathcal O(1)$ to be invertible) and in the discussion of the Segre embedding, we assume that $S$ and $T$ are generated over $A$ by $S_1$ and $T_1$ respectively. (The statement about $\mathcal O(1)$ is meaningless otherwise; we make no claim about the first statement without this hypothesis.) Put $R=S\times_AT$, so $R_d=S_d\otimes_AT_d$ and $R_1=S_1\otimes_AT_1$. $R$ is generated by $R_1$ over $A$ because $R_d$ is spanned by $(x_{i_1}\cdots x_{i_d})\otimes(y_{j_1}\cdots y_{j_d})=\prod_k(x_{i_k}\otimes y_{j_k})$.

\emph{Step 1: affine pieces.} Let $x\in S_a$, $y\in T_a$ ($a\ge1$), so $x\otimes y\in R_a$. Then
\[
R_{(x\otimes y)}\ \cong\ S_{(x)}\otimes_AT_{(y)},\qquad \frac{s\otimes t}{(x\otimes y)^k}\leftrightarrow\frac s{x^k}\otimes\frac t{y^k}\quad(s\in S_{ak},\,t\in T_{ak}).
\]
Indeed, $S_{(x)}=(S_x)_0=\varinjlim_k(S_{ak},\cdot x)$, $T_{(y)}=\varinjlim_k(T_{ak},\cdot y)$, $R_{(x\otimes y)}=\varinjlim_k(R_{ak},\cdot x\otimes y)$, and tensor products commute with filtered colimits, so $S_{(x)}\otimes_AT_{(y)}=\varinjlim_k S_{ak}\otimes_AT_{ak}=\varinjlim_kR_{ak}$; the identification respects the ring structures. Similarly for $M=S(1)$, $N=T(1)$ one gets (for $a=1$) an isomorphism of modules over this ring
\[
R(1)_{(x\otimes y)}=\varinjlim_k(S_{k+1}\otimes T_{k+1})\ \cong\ S(1)_{(x)}\otimes_AT(1)_{(y)}.
\]

\emph{Step 2: the isomorphism of schemes.} Choose generating sets $\{x_\alpha\}$ of $S_1$ and $\{y_\beta\}$ of $T_1$ over $A$. Then $X=\bigcup D_+(x_\alpha)$, $Y=\bigcup D_+(y_\beta)$, $X\times_AY=\bigcup U_{\alpha\beta}$ with $U_{\alpha\beta}=D_+(x_\alpha)\times_AD_+(y_\beta)=\Spec(S_{(x_\alpha)}\otimes T_{(y_\beta)})$ (3.3), and the $x_\alpha\otimes y_\beta$ generate the $A$-module $R_1$, hence the ideal $R_+$, so $Z=\operatorname{Proj}R=\bigcup V_{\alpha\beta}$, $V_{\alpha\beta}=D_+(x_\alpha\otimes y_\beta)=\Spec R_{(x_\alpha\otimes y_\beta)}$. By Step 1 we have isomorphisms $\phi_{\alpha\beta}:V_{\alpha\beta}\to U_{\alpha\beta}$. On overlaps, $V_{\alpha\beta}\cap V_{\alpha'\beta'}=D_+(x_\alpha x_{\alpha'}\otimes y_\beta y_{\beta'})=\Spec R_{(x_\alpha x_{\alpha'}\otimes y_\beta y_{\beta'})}$ and $U_{\alpha\beta}\cap U_{\alpha'\beta'}=D_+(x_\alpha x_{\alpha'})\times D_+(y_\beta y_{\beta'})=\Spec S_{(x_\alpha x_{\alpha'})}\otimes T_{(y_\beta y_{\beta'})}$; the isomorphism of Step 1 (with $a=2$) for this pair restricts from both $\phi_{\alpha\beta}$ and $\phi_{\alpha'\beta'}$, because all these isomorphisms are given by the same formula $\frac{s\otimes t}{(x\otimes y)^k}\mapsto\frac s{x^k}\otimes\frac t{y^k}$ and the localization maps are induced by the identity on $S\times_AT$. (Explicitly, write $x=x_\alpha,x'=x_{\alpha'}$, $y=y_\beta,y'=y_{\beta'}$. The ring map $S_{(x)}\otimes T_{(y)}\to S_{(xx')}\otimes T_{(yy')}\cong R_{(xx'\otimes yy')}$ coming from $U_{\alpha'\beta'}\cap U_{\alpha\beta}\subseteq U_{\alpha\beta}$ sends $\frac s{x^k}\otimes\frac t{y^k}$ to $\frac{sx'^k}{(xx')^k}\otimes\frac{ty'^k}{(yy')^k}\leftrightarrow\frac{(s\otimes t)(x'\otimes y')^k}{((x\otimes y)(x'\otimes y'))^k}=\frac{s\otimes t}{(x\otimes y)^k}$, which is the localization map $R_{(x\otimes y)}\to R_{(xx'\otimes yy')}$ composed with the isomorphism of Step 1 for $(x\otimes y)$. The same holds with the two indices exchanged.) Hence the $\phi_{\alpha\beta}$ glue to an isomorphism $\Phi:Z\to X\times_AY$.

\emph{Step 3: the sheaf $\mathcal O(1)$.} On $V=V_{\alpha\beta}$, $\mathcal O_Z(1)|_V=\widetilde{R(1)_{(x\otimes y)}}$ (5.11(b), 5.12(b)). On $U=U_{\alpha\beta}=\Spec B$, $B=S_{(x)}\otimes_AT_{(y)}$, put $M_1=S(1)_{(x)}$, $N_1=T(1)_{(y)}$; then $\mathcal O_X(1)|_{D_+(x)}=\widetilde{M_1}$, $\mathcal O_Y(1)|_{D_+(y)}=\widetilde{N_1}$ and by 5.2(e), $p_1^*\widetilde{M_1}=(M_1\otimes_{S_{(x)}}B)^\sim=(M_1\otimes_AT_{(y)})^\sim$, $p_2^*\widetilde{N_1}=(S_{(x)}\otimes_AN_1)^\sim$. By 5.2(b), $p_1^*\mathcal O_X(1)\otimes p_2^*\mathcal O_Y(1)|_U=\big((M_1\otimes_AT_{(y)})\otimes_B(S_{(x)}\otimes_AN_1)\big)^\sim=(M_1\otimes_AN_1)^\sim$. By Step 1, $M_1\otimes_AN_1\cong R(1)_{(x\otimes y)}$ as $B\cong R_{(x\otimes y)}$-modules. So $\Phi^*(p_1^*\mathcal O_X(1)\otimes p_2^*\mathcal O_Y(1))|_V\cong\mathcal O_Z(1)|_V$, via the formula $(\frac s{x^k})\otimes(\frac t{y^k})\mapsto\frac{s\otimes t}{(x\otimes y)^k}$ (with $s\in S_{k+1},t\in T_{k+1}$); these local isomorphisms agree on overlaps for the same reason as in Step 2, hence glue to a global isomorphism $\mathcal O_Z(1)\cong\Phi^*\big(p_1^*\mathcal O_X(1)\otimes p_2^*\mathcal O_Y(1)\big)$.

\emph{The Segre embedding.} By the above, $R$ is generated over $A$ by the elements $x_i\otimes y_j$ of degree $1$, so $A[z_{ij}]\to R$, $z_{ij}\mapsto x_i\otimes y_j$, is a surjective graded homomorphism ($N+1=(r+1)(s+1)$ variables), and $\operatorname{Proj}R\hookrightarrow\mathbf P^N_A$ is a closed immersion (Ex.\ 3.12). For $S=A[x_0..x_r]$, $T=A[y_0..y_s]$ this is the Segre embedding $\sigma:\mathbf P^r\times\mathbf P^s\to\mathbf P^N$, and $\sigma^*\mathcal O(1)\cong p_1^*\mathcal O(1)\otimes p_2^*\mathcal O(1)$ by the above. For general $S=A[x]/I$, $T=A[y]/J$, $S\times_AT$ is a quotient of $A[x]\times_AA[y]$, so $\operatorname{Proj}(S\times_AT)$ is a closed subscheme of $\mathbf P^r\times\mathbf P^s\subseteq\mathbf P^N$; by naturality of $\Phi$ with respect to the surjections $A[x]\to S$, $A[y]\to T$ it is $X\times_AY\subseteq\mathbf P^r\times\mathbf P^s$, embedded by $\sigma$. \hfill$\square$

% ---------------------------------------------------------------- 5.12
\exercise{5.12}{
\begin{enumerate}[label=(\alph*),leftmargin=*]
\item Let $X$ be a scheme over a scheme $Y$, and let $\mathcal L,\mathcal M$ be two very ample invertible sheaves on $X$. Show that $\mathcal L\otimes\mathcal M$ is also very ample. [\emph{Hint}: Use a Segre embedding.]
\item Let $f:X\to Y$ and $g:Y\to Z$ be two morphisms of schemes. Let $\mathcal L$ be a very ample invertible sheaf on $X$ relative to $Y$, and let $\mathcal M$ be a very ample invertible sheaf on $Y$ relative to $Z$. Show that $\mathcal L\otimes f^*\mathcal M$ is a very ample invertible sheaf on $X$ relative to $Z$.
\end{enumerate}}
\pf
In (a) ``very ample'' means very ample relative to $Y$: there is an immersion $i:X\to\mathbf P^r_Y$ of $Y$-schemes with $i^*\mathcal O(1)\cong\mathcal L$ (the definition after 5.16).

\emph{The Segre embedding over a base.} Apply Ex.\ 5.11 with $A=\mathbf Z$, $S=\mathbf Z[x_0,\dots,x_r]$, $T=\mathbf Z[y_0,\dots,y_s]$: there is a closed immersion $\sigma_{\mathbf Z}:\mathbf P^r_{\mathbf Z}\times_{\mathbf Z}\mathbf P^s_{\mathbf Z}\to\mathbf P^N_{\mathbf Z}$, $N=rs+r+s$, with $\sigma_{\mathbf Z}^*\mathcal O(1)\cong p_1^*\mathcal O(1)\otimes p_2^*\mathcal O(1)$. For any scheme $W$ put $\mathbf P^n_W=\mathbf P^n_{\mathbf Z}\times W$ and $\mathcal O_{\mathbf P^n_W}(1)=h^*\mathcal O_{\mathbf P^n_{\mathbf Z}}(1)$ with $h$ the projection. Base changing by $W$, $\sigma_W:\mathbf P^r_W\times_W\mathbf P^s_W\to\mathbf P^N_W$ is a closed immersion (base change of a closed immersion), and since $h\circ\sigma_W=\sigma_{\mathbf Z}\circ h'$ for the projection $h'$ and $p_k\circ h'=h\circ q_k$ ($q_k$ the projections of $\mathbf P^r_W\times_W\mathbf P^s_W$),
\[
\sigma_W^*\mathcal O(1)\cong q_1^*\mathcal O_{\mathbf P^r_W}(1)\otimes q_2^*\mathcal O_{\mathbf P^s_W}(1),
\]
using that $u^*$ commutes with $\otimes$ and $(uv)^*=v^*u^*$.

We will use the standard facts: open and closed immersions, hence immersions, are stable under composition and base change; the diagonal $\Delta:X\to X\times_YX$ is an immersion (Sect.~4).

\textbf{(a)} Let $i:X\to\mathbf P^r_Y$, $j:X\to\mathbf P^s_Y$ be immersions with $i^*\mathcal O(1)=\mathcal L$, $j^*\mathcal O(1)=\mathcal M$. Let $\kappa=(i\times_Yj)\circ\Delta:X\to X\times_YX\to\mathbf P^r_Y\times_Y\mathbf P^s_Y$; here $i\times_Yj=(i\times\mathrm{id})\circ(\mathrm{id}\times j)$ is a composition of base changes of immersions, so an immersion, and $\Delta$ is an immersion; hence $\kappa$ is an immersion. Then $k=\sigma_Y\circ\kappa:X\to\mathbf P^N_Y$ is an immersion of $Y$-schemes, and since $q_1\kappa=i$, $q_2\kappa=j$,
\[
k^*\mathcal O(1)=\kappa^*\big(q_1^*\mathcal O(1)\otimes q_2^*\mathcal O(1)\big)=i^*\mathcal O(1)\otimes j^*\mathcal O(1)\cong\mathcal L\otimes\mathcal M.
\]
So $\mathcal L\otimes\mathcal M$ is very ample relative to $Y$ (it is invertible as a tensor product of invertible sheaves).

\textbf{(b)} Let $i:X\to\mathbf P^r_Y$ ($Y$-immersion) with $i^*\mathcal O(1)=\mathcal L$ and $j:Y\to\mathbf P^s_Z$ ($Z$-immersion) with $j^*\mathcal O(1)=\mathcal M$. Write $\mathbf P^r_Y=\mathbf P^r_Z\times_ZY$ with projections $q:\mathbf P^r_Y\to\mathbf P^r_Z$, $\pi:\mathbf P^r_Y\to Y$; by definition $\mathcal O_{\mathbf P^r_Y}(1)=q^*\mathcal O_{\mathbf P^r_Z}(1)$. Let $\tau=\mathrm{id}\times j:\mathbf P^r_Z\times_ZY\to\mathbf P^r_Z\times_Z\mathbf P^s_Z$, an immersion (base change of $j$), and $k=\sigma_Z\circ\tau\circ i:X\to\mathbf P^N_Z$, an immersion of $Z$-schemes ($i$ is an immersion; composites of immersions are immersions). With $p_1\tau=q$, $p_2\tau=j\pi$:
\[
\tau^*\sigma_Z^*\mathcal O(1)=\tau^*(p_1^*\mathcal O(1)\otimes p_2^*\mathcal O(1))=q^*\mathcal O(1)\otimes\pi^*j^*\mathcal O(1)=\mathcal O_{\mathbf P^r_Y}(1)\otimes\pi^*\mathcal M,
\]
and pulling back by $i$, using $\pi\circ i=f$ ($i$ is a morphism over $Y$): $k^*\mathcal O(1)=i^*\mathcal O(1)\otimes i^*\pi^*\mathcal M=\mathcal L\otimes f^*\mathcal M$. \hfill$\square$

% ---------------------------------------------------------------- 5.13
\exercise{5.13}{Let $S$ be a graded ring, generated by $S_1$ as an $S_0$-algebra. For any integer $d>0$, let $S^{(d)}$ be the graded ring $\bigoplus_{n\ge0}S^{(d)}_n$ where $S^{(d)}_n=S_{nd}$. Let $X=\operatorname{Proj}S$. Show that $\operatorname{Proj}S^{(d)}\cong X$, and that the sheaf $\mathcal O(1)$ on $\operatorname{Proj}S^{(d)}$ corresponds via this isomorphism to $\mathcal O_X(d)$.

This construction is related to the $d$-uple embedding (I, Ex.\ 2.12) in the following way. If $x_0,\dots,x_r$ is a set of generators for $S_1$, corresponding to an embedding $X\hookrightarrow\mathbf P^r_A$, then the set of monomials of degree $d$ in the $x_i$ is a set of generators for $S^{(d)}_1=S_d$. These define a projective embedding of $\operatorname{Proj}S^{(d)}$ which is none other than the image of $X$ under the $d$-uple embedding of $\mathbf P^r_A$.}
\pf
Write $A=S_0$. Note that $S^{(d)}$ is generated over $A$ by $S^{(d)}_1=S_d$: $S_{nd}$ is spanned over $A$ by products of $nd$ elements of $S_1$, i.e.\ by products of $n$ elements of $S_d$.

\emph{The schemes.} For $h\in S_d$ let $D_+(h)\subseteq X$ and $D'_+(h)\subseteq X'=\operatorname{Proj}S^{(d)}$ be the corresponding basic open sets ($h$ has degree $d$ in $S$ and degree $1$ in $S^{(d)}$). Then
\[
S^{(d)}_{(h)}=\Big\{\tfrac s{h^k}:s\in S_{dk}\Big\}=S_{(h)}
\]
as subrings of $S_h$: indeed $(S_h)_0$ consists of the elements $s/h^k$ with $s\in S_{dk}$ (the denominators $h^k$ have degree $dk$), and both rings are by definition this set. So $D_+(h)=\Spec S_{(h)}=D'_+(h)$ as schemes. Now $X=\bigcup_{h\in S_d}D_+(h)$: if $\fp\in X$, some $f\in S_1$ is not in $\fp$ ($S_+$ generated by $S_1$), and $f^d\in S_d\setminus\fp$. Likewise $X'=\bigcup_{h\in S_d}D'_+(h)$ because $S^{(d)}$ is generated by $S^{(d)}_1=S_d$. For $h,h'\in S_d$, $D_+(h)\cap D_+(h')=D_+(hh')=\Spec S_{(hh')}$ and $D'_+(h)\cap D'_+(h')=D'_+(hh')=\Spec S^{(d)}_{(hh')}=\Spec S_{(hh')}$ (same computation with $hh'$ of degree $2d$), and the inclusions are given by the same localization maps. Thus $X$ and $X'$ are glued from the same affine schemes $\Spec S_{(h)}$ ($h\in S_d$) along the same open subschemes by the same maps, so there is a canonical isomorphism $X'\cong X$ restricting to the identity of $\Spec S_{(h)}$ for each $h$.

\emph{The sheaf.} On $D'_+(h)$, $\mathcal O_{X'}(1)=\widetilde{S^{(d)}(1)}$ has sections $S^{(d)}(1)_{(h)}=\{s/h^k:s\in S^{(d)}_{k+1}=S_{d(k+1)}\}$, and on $D_+(h)$, $\mathcal O_X(d)=\widetilde{S(d)}$ has sections $S(d)_{(h)}=\{s/h^k: s\in S(d)_{dk}=S_{dk+d}\}$. These are the same $S_{(h)}$-submodules of $S_h$ (the degree $d$ part of $S_h$), compatible with restriction to $D_+(hh')$. So $\mathcal O_{X'}(1)\cong\mathcal O_X(d)$ under $X'\cong X$ (5.11(b), 5.12(b)).

\emph{The $d$-uple embedding.} Let $x_0,\dots,x_r$ generate $S_1$ and $M_0,\dots,M_N$ be the monomials of degree $d$ in them; they generate $S_d$ over $A$. The surjection $A[z_0,\dots,z_N]\to S^{(d)}$, $z_j\mapsto M_j$, gives a closed immersion $X'=\operatorname{Proj}S^{(d)}\hookrightarrow\mathbf P^N_A$ (Ex.\ 3.12). In particular for $S=A[x_0,\dots,x_r]$ this gives $\mathbf P^r_A\cong\operatorname{Proj}S^{(d)}\hookrightarrow\mathbf P^N_A$, which is the $d$-uple embedding $\nu_d$ (the map defined by the monomials of degree $d$, with $\nu_d^*\mathcal O(1)=\mathcal O(d)$ by the above). For general $S=A[x]/I$, the surjection $A[x]\to S$ induces a surjection $A[x]^{(d)}\to S^{(d)}$, so $X'$ is a closed subscheme of $\operatorname{Proj}A[x]^{(d)}\cong\mathbf P^r_A$, and the isomorphisms of the first part are compatible with these maps (they are identities on each $\Spec S_{(h)}$). Hence the embedding of $X'\cong X$ in $\mathbf P^N_A$ is the composite $X\hookrightarrow\mathbf P^r_A\xrightarrow{\nu_d}\mathbf P^N_A$, the image of $X$ under the $d$-uple embedding. \hfill$\square$

% ---------------------------------------------------------------- 5.14
\exercise{5.14}{Let $A$ be a ring, and let $X$ be a closed subscheme of $\mathbf P^r_A$. We define the \emph{homogeneous coordinate ring} $S(X)$ of $X$ for the given embedding to be $A[x_0,\dots,x_r]/I$, where $I$ is the ideal $\Gamma_*(\mathcal I_X)$ constructed in the proof of (5.16). (Of course if $A$ is a field and $X$ a variety, this coincides with the definition given in (I, Sect.~2)!) Recall that a scheme $X$ is \emph{normal} if its local rings are integrally closed domains. A closed subscheme $X\subseteq\mathbf P^r_k$ is \emph{projectively normal} for the given embedding, if its homogeneous coordinate ring $S(X)$ is an integrally closed domain (cf.\ (I, Ex.\ 3.18)). Now assume that $k$ is an algebraically closed field, and that $X$ is a connected, normal closed subscheme of $\mathbf P^r_k$. Show that for some $d>0$, the $d$-uple embedding of $X$ is projectively normal, as follows.
\begin{enumerate}[label=(\alph*),leftmargin=*]
\item Let $S$ be the homogeneous coordinate ring of $X$, and let $S'=\bigoplus_{n\ge0}\Gamma(X,\mathcal O_X(n))$. Show that $S$ is a domain, and that $S'$ is its integral closure. [\emph{Hint}: First show that $X$ is integral. Then regard $S'$ as the global sections of the sheaf of rings $\mathcal S=\bigoplus_{n\ge0}\mathcal O_X(n)$ on $X$, and show that $\mathcal S$ is a sheaf of integrally closed domains.]
\item Use (Ex.\ 5.9) to show that $S_d=S'_d$ for all sufficiently large $d$.
\item Show that $S^{(d)}$ is integrally closed for sufficiently large $d$, and hence conclude that the $d$-uple embedding of $X$ is projectively normal.
\item As a corollary of (a), show that a closed subscheme $X\subseteq\mathbf P^r_k$ is projectively normal if and only if it is normal, and for every $n\ge0$ the natural map $\Gamma(\mathbf P^r,\mathcal O_{\mathbf P^r}(n))\to\Gamma(X,\mathcal O_X(n))$ is surjective.
\end{enumerate}}
\pf
Throughout $X\neq\emptyset$ (connected spaces are nonempty), $P=k[x_0,\dots,x_r]$, $I=\Gamma_*(\mathcal I_X)$ (a saturated homogeneous ideal with $\widetilde I=\mathcal I_X$, by 5.10(c)), $S=P/I$, so $X=\operatorname{Proj}S$ as a closed subscheme (Ex.\ 3.12) and $\mathcal O_X(n)=i^*\mathcal O_{\mathbf P^r}(n)=\widetilde S(n)$ (5.12(c)). $S$ is generated by $S_1$ over $S_0=k$ and $S_1$ is finite dimensional. Also $S_1\neq0$, since $S_1=0$ would give $\operatorname{Proj}S=\emptyset$. We assume $r\ge1$ (needed for 5.10 and 5.13); if $r=0$ then $X$ is the point $\mathbf P^0_k$, $S=k[x_0]$ and $S^{(d)}=k[x_0^d]$ are polynomial rings in one variable, hence integrally closed, and everything is trivial.

\textbf{(a)} \emph{$X$ is integral.} $X$ is noetherian and its local rings are domains. If a point lay on two irreducible components, its local ring would have two minimal primes. So the (finitely many) irreducible components are pairwise disjoint closed sets, hence each is open; as $X$ is connected, $X$ is irreducible. Local rings are domains, so $X$ is reduced. Hence $X$ is integral; with $K$ its function field. (It is a variety.)

\emph{$S$ injects into $S'$.} The exact sequence $0\to\mathcal I_X\to\mathcal O_{\mathbf P}\to i_*\mathcal O_X\to0$ twisted by $\mathcal O(n)$ and the left exactness of $\Gamma$ give $0\to\Gamma(\mathcal I_X(n))=I_n\to P_n\to\Gamma(X,\mathcal O_X(n))$, so $S_n=P_n/I_n\hookrightarrow\Gamma(X,\mathcal O_X(n))$ for $n\ge0$; this is compatible with products, so $S\subseteq S'$ is a subring.

\emph{$\mathcal S$ is a sheaf of integrally closed domains.} Let $\xi$ be the generic point. Each $\mathcal O_X(n)$ is invertible, and $\mathcal S=\bigoplus_{n\ge0}\mathcal O_X(n)$ is a sheaf of graded rings. Let $U=\Spec B$ be an open affine on which $\mathcal O_X(1)|_U=\mathcal O_Ue$ is free; then $\mathcal O(n)|_U=\mathcal O_Ue^n$ and $\mathcal S(U)=B[e]$ is a polynomial ring, with $B$ an integrally closed domain (the local rings of $B$ are integrally closed and $B=\bigcap B_\fm$). A polynomial ring over an integrally closed domain is integrally closed. The stalk $\mathcal S_\xi=K[e]$. For every nonempty open $V$ the restriction map $\mathcal S(V)\to\mathcal S_\xi$ is injective (check on trivializing affine opens, where it is $B[e]\subseteq K[e]$). Since $X$ is quasi-compact, $S'=\Gamma(X,\mathcal S)$, and $S'\hookrightarrow\mathcal S_\xi=K[e]$ is injective, so $S'$ is a domain, and so is $S\subseteq S'$.

\emph{$S'$ is integrally closed in $\operatorname{Frac}S'$.} Let $\theta\in\operatorname{Frac}(S')$ be integral over $S'$. Embed $\operatorname{Frac}S'\subseteq\Omega=K(e)$. For a trivializing open affine $U$ as above, $S'\to\mathcal S(U)\subseteq\mathcal S_\xi$, so $\theta$ is integral over $\mathcal S(U)=B[e]$ and lies in its fraction field $K(e)=\Omega$, hence $\theta\in\mathcal S(U)$ (integrally closed). Call it $\theta_U$. For two such $U,U'$, $\theta_U|_{U\cap U'}$ and $\theta_{U'}|_{U\cap U'}$ have the same germ $\theta$ at $\xi$, so they are equal; the $\theta_U$ glue to $\tilde\theta\in\mathcal S(X)=S'$ with germ $\theta$. So $\theta\in S'$.

\emph{$S'$ is integral over $S$ and $\operatorname{Frac}S'=\operatorname{Frac}S$.} Apply the Lemma in the solution of 5.9 to $T=S$ (a graded domain with $T_0=k$): it shows $S'$ (which equals the $T'$ of the lemma by part (i) of the lemma) is integral over $S$ and every $s'\in S'$ satisfies $x_i^ns'\in S$ with some $x_i\ne0$, so $s'\in\operatorname{Frac}S$. Hence $\operatorname{Frac}S'=\operatorname{Frac}S$. Thus $S'$ is integral over $S$ and integrally closed in $\operatorname{Frac}S$, so $S'$ is the integral closure of $S$.

\textbf{(b)} Apply Ex.\ 5.9(b) to $M=S$ ($S_0=k$, $S_1$ finite-dimensional; $\widetilde S(d)=\mathcal O_X(d)$): $\alpha_d:S_d\to\Gamma(X,\mathcal O_X(d))=S'_d$, which is the inclusion $S\subseteq S'$, is an isomorphism for $d\ge d_0$ (this is also (iv) of the Lemma).

\textbf{(c)} Let $d\ge d_0$. First, $S'_0=\Gamma(X,\mathcal O_X)$ is a domain, and integral over $S_0=k$: an element $a\in S'_0$ is integral over $S$ by the lemma's (iii), and taking degree-$0$ components in an equation of integral dependence $a^q+c_1a^{q-1}+\dots+c_q=0$ ($c_i\in S$) gives one with coefficients in $S_0=k$. So every element of $S'_0$ is algebraic over the algebraically closed field $k$, hence in $k$: $S'_0=k=S_0$. By (b), $S'_{nd}=S_{nd}$ for $n\ge1$. Hence $S^{(d)}=S'^{(d)}$.

\emph{$S'^{(d)}$ is integrally closed.} Let $\theta\in\operatorname{Frac}(S'^{(d)})$ be integral over $S'^{(d)}$. Then $\theta\in\operatorname{Frac}S'$ is integral over $S'$, so $\theta\in S'$ by (a). Write $\theta=\sum\theta_m$ in homogeneous components, and $\theta=a/b$ with $a,b\in S'^{(d)}$, $b\ne0$. We claim all $\theta_m$ with $\theta_m\neq0$ have $d\mid m$. Let $m_0$ be the least $m$ with $\theta_m\ne0$ and $c_0$ the least degree of a nonzero component of $b$ (so $d\mid c_0$). The component of $a=\theta b$ of degree $m_0+c_0$ is $\theta_{m_0}b_{c_0}\neq0$ ($S'$ is a domain), and $a\in S'^{(d)}$, so $d\mid m_0+c_0$, so $d\mid m_0$. Then $\theta_{m_0}\in S'^{(d)}$ and $\theta-\theta_{m_0}=(a-\theta_{m_0}b)/b$ again has the same form with fewer nonzero components; by induction every component has degree divisible by $d$, i.e.\ $\theta\in S'^{(d)}$. So $S^{(d)}=S'^{(d)}$ is an integrally closed domain.

\emph{Conclusion.} The $d$-uple embedding of $X$ is $X\cong\operatorname{Proj}S^{(d)}\hookrightarrow\mathbf P^N_k$ (Ex.\ 5.13), given by the surjection $P_N=k[z_0,\dots,z_N]\to S^{(d)}$, $z_j\mapsto M_j$, with kernel $J$. $J$ is saturated: if $z_i^ns\in J$ for all $i$, then $M_i^n\bar s=0$ in the domain $S^{(d)}\subseteq S$; as $S_d\ne0$ some $M_i\neq0$, so $\bar s=0$, $s\in J$. The closed subscheme $X_d=\operatorname{Proj}(P_N/J)$ of $\mathbf P^N_k$ has ideal sheaf $\widetilde J$, so by Fact 2 in the solution of 5.10, $\Gamma_*(\mathcal I_{X_d})=\Gamma_*(\widetilde J)=\bar J=J$. So the homogeneous coordinate ring of the $d$-uple embedded $X$ is $P_N/J\cong S^{(d)}$, an integrally closed domain: the $d$-uple embedding is projectively normal.

\textbf{(d)} ($\Rightarrow$) Let $S=S(X)$ be an integrally closed domain (so $S\ne0$ and $X\neq\emptyset$). $X=\operatorname{Proj}S$ is irreducible (generic point $(0)$), hence connected. For $\fp\in X$ let $\Sigma$ be the set of homogeneous elements not in $\fp$; $\mathcal O_{X,\fp}=S_{(\fp)}=(\Sigma^{-1}S)_0$ is a domain. Let $\theta\in\operatorname{Frac}S_{(\fp)}$ be integral over $S_{(\fp)}$. As $\Sigma^{-1}S$ is an integrally closed domain (localization of $S$), $\theta\in\Sigma^{-1}S$. Write $\theta=a/b$ with $a,b\in S_{(\fp)}$, $b\ne0$ homogeneous of degree $0$; then $\theta b=a$ has degree $0$, so $\theta_mb=0$ for the components $\theta_m$, $m\neq0$, hence $\theta_m=0$ and $\theta\in S_{(\fp)}$. So $X$ is normal. Now $X$ is a connected normal closed subscheme, so by (a) $S'$ is the integral closure of $S$, hence $S'=S$, i.e.\ $P_n\to S_n=\Gamma(X,\mathcal O_X(n))$ is surjective for all $n\ge0$.

($\Leftarrow$) Suppose $X$ is normal and $P_n=\Gamma(\mathbf P^r,\mathcal O(n))\to\Gamma(X,\mathcal O_X(n))$ is surjective for all $n\ge0$ (and $X\ne\emptyset$). For $n=0$ this says $\Gamma(X,\mathcal O_X)=k$; if $X$ were a disjoint union of two nonempty open-closed subsets, $\Gamma(X,\mathcal O_X)$ would contain $k\times k$, of dimension $\ge2$. So $X$ is connected, and (a) applies: $S$ is a domain and $S'$ its integral closure. Surjectivity means $S_n=S'_n$ for all $n\ge0$, i.e.\ $S=S'$ is integrally closed. So $X$ is projectively normal. \hfill$\square$

% ---------------------------------------------------------------- 5.15
\exercise{5.15}{\emph{Extension of Coherent Sheaves.} We will prove the following theorem in several steps: Let $X$ be a noetherian scheme, let $U$ be an open subset, and let $\mathcal F$ be a coherent sheaf on $U$. Then there is a coherent sheaf $\mathcal F'$ on $X$ such that $\mathcal F'|_U\cong\mathcal F$.
\begin{enumerate}[label=(\alph*),leftmargin=*]
\item On a noetherian affine scheme, every quasi-coherent sheaf is the union of its coherent subsheaves. We say a sheaf $\mathcal F$ is the \emph{union} of its subsheaves $\mathcal F_\alpha$ if for every open set $U$, the group $\mathcal F(U)$ is the union of the subgroups $\mathcal F_\alpha(U)$.
\item Let $X$ be an affine noetherian scheme, $U$ an open subset, and $\mathcal F$ coherent on $U$. Then there exists a coherent sheaf $\mathcal F'$ on $X$ with $\mathcal F'|_U\cong\mathcal F$. [\emph{Hint}: Let $i:U\to X$ be the inclusion map. Show that $i_*\mathcal F$ is quasi-coherent, then use (a).]
\item With $X,U,\mathcal F$ as in (b), suppose furthermore we are given a quasi-coherent sheaf $\mathcal G$ on $X$ such that $\mathcal F\subseteq\mathcal G|_U$. Show that we can find $\mathcal F'$ a coherent subsheaf of $\mathcal G$, with $\mathcal F'|_U\cong\mathcal F$. [\emph{Hint}: Use the same method, but replace $i_*\mathcal F$ by $\rho^{-1}(i_*\mathcal F)$, where $\rho$ is the natural map $\mathcal G\to i_*(\mathcal G|_U)$.]
\item Now let $X$ be a noetherian scheme, $U$ an open subset, $\mathcal F$ a coherent sheaf on $U$, and $\mathcal G$ a quasi-coherent sheaf on $X$ such that $\mathcal F\subseteq\mathcal G|_U$. Show that there is a coherent subsheaf $\mathcal F'\subseteq\mathcal G$ on $X$ with $\mathcal F'|_U\cong\mathcal F$. Taking $\mathcal G=i_*\mathcal F$ proves the result announced at the beginning. [\emph{Hint}: Cover $X$ with open affines, and extend over one of them at a time.]
\item As an extra corollary, show that on a noetherian scheme $X$, every quasi-coherent sheaf $\mathcal F$ is the union of its coherent subsheaves. [\emph{Hint}: If $s$ is a section of $\mathcal F$ over an open set $U$, apply (d) to the subsheaf of $\mathcal F|_U$ generated by $s$.]
\end{enumerate}}
\pf
\textbf{(a)} Let $X=\Spec A$, $A$ noetherian, $\mathcal F=\widetilde M$ (5.4/5.5). $M$ is the union of its finitely generated submodules $N$; $\widetilde N\subseteq\widetilde M$ is a subsheaf (exactness of $\widetilde{\ }$) and is coherent. Let $U\subseteq X$ be open and $s\in\mathcal F(U)$. $U$ is quasi-compact ($X$ noetherian), so $U=D(f_1)\cup\dots\cup D(f_n)$. Then $s|_{D(f_i)}\in M_{f_i}$ has the form $m_i/f_i^{k_i}$ with $m_i\in M$. Let $N=\sum Am_i$, finitely generated. Then $s|_{D(f_i)}\in(N_{f_i})=\widetilde N(D(f_i))$ for all $i$; as $\widetilde N$ is a subsheaf of $\mathcal F$ these sections glue (the restrictions to overlaps agree in $\mathcal F$, hence in $\widetilde N$) to a section of $\widetilde N(U)$ mapping to $s$. So $s\in\widetilde N(U)$.

We note: if $\mathcal G_1,\mathcal G_2\subseteq\mathcal F$ are coherent subsheaves on a noetherian scheme, so is $\mathcal G_1+\mathcal G_2$ (image of $\mathcal G_1\oplus\mathcal G_2\to\mathcal F$, coherent by 5.7), so coherent subsheaves form a directed family.

\textbf{(b)} Let $i:U\to X$ be the inclusion. Since $U$ is noetherian and $\mathcal F$ is quasi-coherent, $i_*\mathcal F$ is quasi-coherent by 5.8(c). Moreover $(i_*\mathcal F)|_U=\mathcal F$. Part (b) is therefore the special case $\mathcal G=i_*\mathcal F$ of (c): (c) gives a coherent $\mathcal F'\subseteq i_*\mathcal F$ with $\mathcal F'|_U=\mathcal F$. (The proof of (c) below is exactly the method of the hint, applied to $\mathcal H=\rho^{-1}(i_*\mathcal F)$, which for $\mathcal G=i_*\mathcal F$ is $i_*\mathcal F$ itself.)

\textbf{(c)} Let $\mathcal H=\rho^{-1}(i_*\mathcal F)$, i.e.\ $\mathcal H(W)=\{g\in\mathcal G(W):g|_{W\cap U}\in\mathcal F(W\cap U)\}$; this is the kernel of the composite $\mathcal G\xrightarrow{\rho}i_*(\mathcal G|_U)\to i_*(\mathcal G|_U)/i_*\mathcal F$. Now $i_*(\mathcal G|_U)$ is quasi-coherent by 5.8(c) ($U$ noetherian), and so is $i_*\mathcal F$ (as $\mathcal F$ is quasi-coherent); $i_*\mathcal F\to i_*(\mathcal G|_U)$ is a morphism of quasi-coherent sheaves, so the quotient is quasi-coherent (5.7), and $\mathcal H$, the kernel of a morphism of quasi-coherent sheaves, is quasi-coherent (5.7). Moreover, for $W\subseteq U$, $\mathcal H(W)=\mathcal F(W)$, i.e.\ $\mathcal H|_U=\mathcal F$ (equal as subsheaves of $\mathcal G|_U$).

By (a), $\mathcal H$ is the union of its coherent subsheaves. Since $U$ is noetherian, cover it by finitely many open affines $V_j=\Spec B_j$ with $\mathcal F|_{V_j}=\widetilde{N_j}$, $N_j$ finitely generated by $n_{j1},\dots,n_{jt_j}\in\mathcal F(V_j)=\mathcal H(V_j)$. Each of these finitely many sections lies in $\mathcal G_\alpha(V_j)$ for some coherent subsheaf $\mathcal G_\alpha\subseteq\mathcal H$; as the family of coherent subsheaves is directed, a single coherent $\mathcal F'\subseteq\mathcal H\subseteq\mathcal G$ contains all of them. Then $\mathcal F'|_U\subseteq\mathcal H|_U=\mathcal F$, and $\mathcal F'|_{V_j}$ contains the sections $n_{jk}$, which generate $\mathcal F|_{V_j}$ (their germs generate the stalks), so $\mathcal F'|_{V_j}=\mathcal F|_{V_j}$. Hence $\mathcal F'|_U=\mathcal F$ as subsheaves.

\textbf{(d)} Let $X=V_1\cup\dots\cup V_n$ with $V_k$ open affine, and $U_k=U\cup V_1\cup\dots\cup V_k$ ($U_0=U$). We construct, by induction on $k$, coherent subsheaves $\mathcal F_k\subseteq\mathcal G|_{U_k}$ with $\mathcal F_k|_{U_{k-1}}=\mathcal F_{k-1}$ and $\mathcal F_0=\mathcal F$. Given $\mathcal F_{k-1}$, let $V=V_k$ (affine noetherian) and $V'=V\cap U_{k-1}$ open in $V$. Then $\mathcal F_{k-1}|_{V'}$ is a coherent subsheaf of $(\mathcal G|_V)|_{V'}$, with $\mathcal G|_V$ quasi-coherent. By (c) there is a coherent $\mathcal F''\subseteq\mathcal G|_V$ with $\mathcal F''|_{V'}=\mathcal F_{k-1}|_{V'}$ (equality of subsheaves, as shown in (c)). Define $\mathcal F_k\subseteq\mathcal G|_{U_k}$ by $\mathcal F_k(W)=\{g\in\mathcal G(W):g|_{W\cap U_{k-1}}\in\mathcal F_{k-1},\ g|_{W\cap V}\in\mathcal F''\}$; this is a subsheaf with $\mathcal F_k|_{U_{k-1}}=\mathcal F_{k-1}$ and $\mathcal F_k|_V=\mathcal F''$ (the two descriptions agree on $V'$), and it is coherent since coherence is local. Then $\mathcal F'=\mathcal F_n\subseteq\mathcal G$ is coherent on $X=U_n$ with $\mathcal F'|_U=\mathcal F$.

For the announced theorem, let $\mathcal F$ be coherent on $U$ and $\mathcal G=i_*\mathcal F$, quasi-coherent by 5.8(c), so that $(i_*\mathcal F)|_U=\mathcal F$ contains $\mathcal F$. This yields a coherent $\mathcal F'$ on $X$ with $\mathcal F'|_U=\mathcal F$.

\textbf{(e)} Let $\mathcal F$ be quasi-coherent on $X$, $U\subseteq X$ open and $s\in\mathcal F(U)$. Let $\mathcal F_s\subseteq\mathcal F|_U$ be the subsheaf generated by $s$, the image of $\mathcal O_U\to\mathcal F|_U$, $1\mapsto s$. On an open affine $V=\Spec A\subseteq U$ with $\mathcal F|_V=\widetilde M$ we get $\mathcal F_s|_V=\widetilde{As|_V}$, which is coherent ($A$ noetherian, $As|_V$ cyclic). So $\mathcal F_s$ is a coherent subsheaf of $\mathcal F|_U$. By (d) (with $\mathcal G=\mathcal F$) there is a coherent $\mathcal F'\subseteq\mathcal F$ with $\mathcal F'|_U=\mathcal F_s$. Then $s\in\mathcal F_s(U)=\mathcal F'(U)$. \hfill$\square$

% ---------------------------------------------------------------- 5.16
\exercise{5.16}{\emph{Tensor Operations on Sheaves.} First we recall the definitions of various tensor operations on a module. Let $A$ be a ring, and let $M$ be an $A$-module. Let $T^n(M)$ be the tensor product $M\otimes\dots\otimes M$ of $M$ with itself $n$ times, for $n\ge1$. For $n=0$ we put $T^0(M)=A$. Then $T(M)=\bigoplus_{n\ge0}T^n(M)$ is a (noncommutative) $A$-algebra, which we call the \emph{tensor algebra} of $M$. We define the \emph{symmetric algebra} $S(M)=\bigoplus_{n\ge0}S^n(M)$ of $M$ to be the quotient of $T(M)$ by the two-sided ideal generated by all expressions $x\otimes y-y\otimes x$, for all $x,y\in M$. Then $S(M)$ is a commutative $A$-algebra. Its component in degree $n$ is called the $n$th \emph{symmetric product} of $M$. We denote the image of $x\otimes y$ in $S(M)$ by $xy$, $x,y\in M$. As an example, note that if $M$ is a free $A$-module of rank $r$, then $S(M)\cong A[x_1,\dots,x_r]$.

We define the \emph{exterior algebra} $\bigwedge(M)$ of $M$ to be the quotient of $T(M)$ by the two-sided ideal generated by all expressions $x\otimes x$ for $x\in M$. Note that this expression contains all expressions of the form $x\otimes y+y\otimes x$, so that $\bigwedge(M)$ is a skew commutative graded $A$-algebra. This means that if $u\in\bigwedge^r(M)$ and $v\in\bigwedge^s(M)$, then $u\wedge v=(-1)^{rs}v\wedge u$ (here we denote the multiplication in this algebra by $\wedge$, i.e.\ the image of $x\otimes y$ in $\bigwedge^2(M)$ is denoted by $x\wedge y$). The $n$th component $\bigwedge^n(M)$ is called the $n$th \emph{exterior power} of $M$.

Now let $(X,\mathcal O_X)$ be a ringed space, and let $\mathcal F$ be a sheaf of $\mathcal O_X$-modules. We define the \emph{tensor algebra, symmetric algebra, and exterior algebra} of $\mathcal F$ by taking the sheaves associated to the presheaf which to an open set $U$ assigns the corresponding tensor operation applied to $\mathcal F(U)$ as an $\mathcal O_X(U)$-module. The results are $\mathcal O_X$-algebras, and their components are $\mathcal O_X$-modules.
\begin{enumerate}[label=(\alph*),leftmargin=*]
\item Suppose $\mathcal F$ is locally free of rank $n$. Then $T^r(\mathcal F)$, $S^r(\mathcal F)$, and $\bigwedge^r(\mathcal F)$ are also locally free, of ranks $n^r$, $\binom{n+r-1}{n-1}$, and $\binom nr$ respectively.
\item Again let $\mathcal F$ be locally free of rank $n$. Then the multiplication map $\bigwedge^r\mathcal F\otimes\bigwedge^{n-r}\mathcal F\to\bigwedge^n\mathcal F$ is a perfect pairing for any $r$, i.e., it induces an isomorphism of $\bigwedge^r\mathcal F$ with $(\bigwedge^{n-r}\mathcal F)^\vee\otimes\bigwedge^n\mathcal F$. As a special case, note that if $\mathcal F$ has rank 2, then $\mathcal F\cong\mathcal F^\vee\otimes\bigwedge^2\mathcal F$.
\item Let $0\to\mathcal F'\to\mathcal F\to\mathcal F''\to0$ be an exact sequence of locally free sheaves. Then for any $r$ there is a finite filtration of $S^r(\mathcal F)$,
\[
S^r(\mathcal F)=F^0\supseteq F^1\supseteq\dots\supseteq F^r\supseteq F^{r+1}=0
\]
with quotients $F^p/F^{p+1}\cong S^p(\mathcal F')\otimes S^{r-p}(\mathcal F'')$ for each $p$.
\item Same statement as (c), with exterior powers instead of symmetric powers. In particular, if $\mathcal F',\mathcal F,\mathcal F''$ have ranks $n',n,n''$ respectively, there is an isomorphism $\bigwedge^n\mathcal F\cong\bigwedge^{n'}\mathcal F'\otimes\bigwedge^{n''}\mathcal F''$.
\item Let $f:X\to Y$ be a morphism of ringed spaces, and let $\mathcal F$ be an $\mathcal O_Y$-module. Then $f^*$ commutes with all the tensor operations on $\mathcal F$: $f^*(S^n(\mathcal F))\cong S^n(f^*\mathcal F)$, etc.
\end{enumerate}}
\pf
(The binomial coefficient in (a) is printed as $\binom{n+r-1}{n-1}$; this equals $\binom{n+r-1}{r}$, the rank of $S^r$ of a free module of rank $n$.) We use that the construction commutes with restriction to open subsets, and it is natural for isomorphisms of $\mathcal F$.

\textbf{(a)} Let $U$ be open with $\mathcal F|_U\cong\mathcal O_U^n$, basis $e_1,\dots,e_n$. For every open $U'\subseteq U$, $\mathcal F(U')=\mathcal O(U')^n$ is free over $B=\mathcal O(U')$ with basis $e_i|_{U'}$, so $T^r(\mathcal F(U'))$ is free with basis $e_{i_1}\otimes\dots\otimes e_{i_r}$ ($n^r$ elements), $S^r(\mathcal F(U'))$ is free with basis the monomials $e_1^{a_1}\cdots e_n^{a_n}$, $\sum a_i=r$ ($\binom{n+r-1}{r}$ elements, as $S(B^n)=B[e_1,\dots,e_n]$), and $\bigwedge^r(\mathcal F(U'))$ is free with basis $e_{i_1}\wedge\dots\wedge e_{i_r}$, $i_1<\dots<i_r$ ($\binom nr$ elements). The resulting presheaf, e.g.\ $U'\mapsto\mathcal O(U')^{N}$, is already a sheaf ($\mathcal O_U^N$) compatible with the restriction maps (which send basis elements to basis elements), so its sheafification is $\mathcal O_U^N$ with $N=n^r,\binom{n+r-1}r,\binom nr$ respectively.

\textbf{(b)} The wedge product on presheaves induces a morphism $\mathcal F_r\otimes\mathcal F_{n-r}\to\mathcal F_n$ where $\mathcal F_j=\bigwedge^j\mathcal F$; by 5.1(c) (applied with $\mathcal E=\mathcal F_{n-r}$, which is locally free of finite rank by (a)) it corresponds to $\phi:\mathcal F_r\to\mathcal{H}om(\mathcal F_{n-r},\mathcal F_n)$, $\phi(s)(t)=s\wedge t$, and by 5.1(b) $\mathcal{H}om(\mathcal F_{n-r},\mathcal F_n)\cong(\bigwedge^{n-r}\mathcal F)^\vee\otimes\bigwedge^n\mathcal F$. To show $\phi$ is an isomorphism we may work on $U$ with $\mathcal F|_U$ free with basis $e_1,\dots,e_n$. Then $\mathcal F_r$, $\mathcal F_{n-r}$, $\mathcal F_n$ have bases $e_I$ ($|I|=r$), $e_J$ ($|J|=n-r$), $e_{[n]}=e_1\wedge\dots\wedge e_n$, and $e_I\wedge e_J=\pm e_{[n]}$ if $J=I^c$ and $0$ otherwise. Hence $\phi(e_I)=\pm\,e_{I^c}^*\otimes e_{[n]}$, i.e.\ $\phi$ maps a basis to $\pm$ a basis: an isomorphism. For $n=2,r=1$: $\mathcal F\cong\mathcal F^\vee\otimes\bigwedge^2\mathcal F$.

\textbf{(c)} Put $\mathcal F^p=$ the image of $S^p(\mathcal F')\otimes S^{r-p}(\mathcal F)\to S^r(\mathcal F)$ (induced by the multiplication $S^p(\mathcal F)\otimes S^{r-p}(\mathcal F)\to S^r(\mathcal F)$ and $\mathcal F'\to\mathcal F$), for $0\le p\le r$, and $F^{r+1}=0$. Clearly $F^0=S^r(\mathcal F)$ and $F^p\supseteq F^{p+1}$.

\emph{Local structure.} Let $x\in X$. Since $\mathcal F',\mathcal F''$ are locally free of finite rank we can find an open neighborhood $U$ of $x$ with $\mathcal F'|_U$, $\mathcal F''|_U$ free, with bases $e'_1,\dots,e'_{n'}$, $g''_1,\dots,g''_{n''}$. As $\mathcal F\to\mathcal F''$ is surjective, each $g''_j$ lifts to a section of $\mathcal F$ locally; shrinking $U$ we may assume lifts $g_j\in\mathcal F(U)$ exist. Then $\mathcal F'|_U\oplus\mathcal O_U^{n''}\to\mathcal F|_U$, $(e'_i,\text{basis})\mapsto(e_i,g_j)$ is a morphism of short exact sequences (identity on the ends), hence an isomorphism (5-lemma on stalks). So $\mathcal F|_U$ is free with basis $e_1,\dots,e_{n'},g_1,\dots,g_{n''}$ ($e_i$ the image of $e'_i$), and $S^r(\mathcal F)|_U$ is free with basis the monomials $e^ag^b$, $|a|+|b|=r$. The image of $S^p(\mathcal F')\otimes S^{r-p}(\mathcal F)$ is spanned by the products of $p$ elements $e_i$ with a monomial of degree $r-p$, i.e.\ by the monomials with $|a|\ge p$. So $F^p|_U$ is free on $\{e^ag^b:|a|\ge p\}$, $F^p/F^{p+1}$ is free on $\{e^ag^b:|a|=p\}$.

\emph{Global isomorphism.} The composite $\Psi:S^p(\mathcal F')\otimes S^{r-p}(\mathcal F)\to F^p\to F^p/F^{p+1}$ vanishes on $S^p(\mathcal F')\otimes K$ where $K=\mathcal F'\cdot S^{r-p-1}(\mathcal F)=$ the kernel of $S^{r-p}(\mathcal F)\to S^{r-p}(\mathcal F'')$ (locally: the span of the monomials with $|a|\ge1$, which is the kernel of the map killing $e_i$): indeed $x\cdot(f'y)=(xf')y\in F^{p+1}$. By right exactness of $\otimes$, $\Psi$ induces $\bar\Psi:S^p(\mathcal F')\otimes S^{r-p}(\mathcal F'')\to F^p/F^{p+1}$; locally it sends the basis element $e^a\otimes g^b$ ($|a|=p$, $|b|=r-p$) to the class of $e^ag^b$, which is a basis of $F^p/F^{p+1}$. Hence $\bar\Psi$ is an isomorphism.

\textbf{(d)} Same proof: $F^p=$ image of $\bigwedge^p\mathcal F'\otimes\bigwedge^{r-p}\mathcal F\to\bigwedge^r\mathcal F$, $F^{r+1}=0$. Locally $\bigwedge^r\mathcal F$ has basis $e_I\wedge g_J$ ($|I|+|J|=r$), $F^p$ is spanned by the $e_I\wedge g_J$ with $|I|\ge p$, so $F^p/F^{p+1}$ is free on those with $|I|=p$. The map $\bigwedge^p\mathcal F'\otimes\bigwedge^{r-p}\mathcal F\to F^p/F^{p+1}$ kills $\bigwedge^p\mathcal F'\otimes(\mathcal F'\wedge\bigwedge^{r-p-1}\mathcal F)$ (as $x\wedge(f'\wedge y)=\pm(x\wedge f')\wedge y\in F^{p+1}$), which is the kernel of $\bigwedge^{r-p}\mathcal F\to\bigwedge^{r-p}\mathcal F''$ (locally: spanned by the $e_I\wedge g_J$ with $I\ne\emptyset$), so we get $\bar\Psi:\bigwedge^p\mathcal F'\otimes\bigwedge^{r-p}\mathcal F''\to F^p/F^{p+1}$ sending $e_I\otimes g_J$ to $e_I\wedge g_J$, an isomorphism.

If the ranks are $n',n,n''$ ($n=n'+n''$) take $r=n$: $F^p/F^{p+1}\cong\bigwedge^p\mathcal F'\otimes\bigwedge^{n-p}\mathcal F''$ is $0$ unless $p\le n'$ and $n-p\le n''$, i.e.\ unless $p=n'$. So $F^0=\dots=F^{n'}$ and $F^{n'+1}=\dots=F^{n+1}=0$, giving $\bigwedge^n\mathcal F=F^{n'}=F^{n'}/F^{n'+1}\cong\bigwedge^{n'}\mathcal F'\otimes\bigwedge^{n''}\mathcal F''$.

\textbf{(e)} Let $\Lambda^n$ denote one of $T^n,S^n,\bigwedge^n$. For $\mathcal O_Y$-modules $\mathcal A,\mathcal B$ let $\operatorname{Mult}^n_\Lambda(\mathcal A,\mathcal B)$ be the set of compatible families of $n$-linear maps $\mathcal A(V)^n\to\mathcal B(V)$ over $\mathcal O(V)$ (all $V$), which are, for $\Lambda=S$, symmetric, and for $\Lambda=\bigwedge$, alternating (the value is $0$ whenever two adjacent arguments coincide; this is exactly the condition for the map $T^n\to\mathcal B$ to kill the two-sided ideal generated by the $x\otimes x$). By the universal property of the tensor power of modules and of sheafification, $\operatorname{Hom}(\Lambda^n\mathcal A,\mathcal B)=\operatorname{Mult}^n_\Lambda(\mathcal A,\mathcal B)$, naturally in $\mathcal B$. Now fix an $\mathcal O_X$-module $\mathcal H$. We claim there is a natural bijection
\[
\operatorname{Mult}^n_\Lambda(\mathcal F,f_*\mathcal H)\ \cong\ \operatorname{Mult}^n_\Lambda(f^*\mathcal F,\mathcal H),\qquad g\mapsto g^\flat,\quad g^\flat(f^*x_1,\dots,f^*x_n)=g(x_1,\dots,x_n),
\]
where $x_i\in\mathcal F(V)$ and $f^*x$ denotes the image of $x$ in $(f^*\mathcal F)(f^{-1}V)$. Granting this, by the adjunction $\operatorname{Hom}_{\mathcal O_X}(f^*\mathcal A,\mathcal B)=\operatorname{Hom}_{\mathcal O_Y}(\mathcal A,f_*\mathcal B)$ we get natural bijections $\operatorname{Hom}(f^*\Lambda^n\mathcal F,\mathcal H)=\operatorname{Hom}(\Lambda^n\mathcal F,f_*\mathcal H)=\operatorname{Mult}_\Lambda^n(\mathcal F,f_*\mathcal H)=\operatorname{Mult}_\Lambda^n(f^*\mathcal F,\mathcal H)=\operatorname{Hom}(\Lambda^n(f^*\mathcal F),\mathcal H)$, natural in $\mathcal H$, so $f^*\Lambda^n\mathcal F\cong\Lambda^n(f^*\mathcal F)$ by Yoneda.

\emph{Proof of the claim.} For ordinary $n$-linear maps, a family $g$ corresponds to a morphism $\mathcal F\to\mathcal{H}om(\mathcal F,\mathcal{H}om(\dots,\mathcal{H}om(\mathcal F,f_*\mathcal H)))$ (5.1(c), $n-1$ internal Homs; it holds for all $\mathcal E$), and for $\mathcal O_Y$-modules $\mathcal A$ there is the natural isomorphism $\mathcal{H}om(\mathcal A,f_*\mathcal H)\cong f_*\mathcal{H}om(f^*\mathcal A,\mathcal H)$ (the adjunction on each open set $V$, applied to $\mathcal A|_V$, $f^{-1}V\to V$). Applying these and the adjunction repeatedly gives a bijection $g\mapsto g^\flat$ between $n$-linear maps $\mathcal F^n\to f_*\mathcal H$ and $n$-linear maps $(f^*\mathcal F)^n\to\mathcal H$, with $g^\flat(f^*x_1,\dots,f^*x_n)=g(x_1,\dots,x_n)$; since every local section of $f^*\mathcal F$ is locally a finite sum $\sum c_kf^*x_k$ with $c_k\in\mathcal O_X$, $g^\flat$ is determined by this formula. Symmetry of $g$ clearly gives symmetry of $g^\flat$ (expand multilinearly) and conversely. For alternating: if $g^\flat$ is alternating then so is $g$ (restrict to the $f^*x_i$). Conversely if $g$ is alternating, then $g$ is skew-symmetric ($g(\dots,x,z,\dots)+g(\dots,z,x,\dots)=g(\dots,x+z,x+z,\dots)-g(\dots,x,x,\dots)-g(\dots,z,z,\dots)=0$), and for a local section $y=\sum_kc_kf^*x_k$ in two adjacent slots,
\[
g^\flat(\dots,y,y,\dots)=\sum_kc_k^2g(\dots,x_k,x_k,\dots)+\sum_{k<l}c_kc_l\big(g(\dots,x_k,x_l,\dots)+g(\dots,x_l,x_k,\dots)\big)=0.
\]
\hfill$\square$

% ---------------------------------------------------------------- 5.17
\exercise{5.17}{\emph{Affine Morphisms.} A morphism $f:X\to Y$ of schemes is \emph{affine} if there is an open affine cover $\{V_i\}$ of $Y$ such that $f^{-1}(V_i)$ is affine for each $i$.
\begin{enumerate}[label=(\alph*),leftmargin=*]
\item Show that $f:X\to Y$ is an affine morphism if and only if for \emph{every} open affine $V\subseteq Y$, $f^{-1}(V)$ is affine. [\emph{Hint}: Reduce to the case $Y$ affine and use (Ex.\ 2.17).]
\item An affine morphism is quasi-compact and separated. Any finite morphism is affine.
\item Let $Y$ be a scheme, and let $\mathcal A$ be a quasi-coherent sheaf of $\mathcal O_Y$-algebras (i.e., a sheaf of rings which is at the same time a quasi-coherent $\mathcal O_Y$-module). Show that there is a unique scheme $X$, and a morphism $f:X\to Y$, such that for every open affine $V\subseteq Y$, $f^{-1}(V)\cong\Spec\mathcal A(V)$, and for every inclusion $U\hookrightarrow V$ of open affines of $Y$, the morphism $f^{-1}(U)\hookrightarrow f^{-1}(V)$ corresponds to the restriction homomorphism $\mathcal A(V)\to\mathcal A(U)$. The scheme $X$ is called $\mathbf{Spec}\,\mathcal A$. [\emph{Hint}: Construct $X$ by glueing the schemes $\Spec\mathcal A(V)$, for $V$ open affine in $Y$.]
\item If $\mathcal A$ is a quasi-coherent $\mathcal O_Y$-algebra, then $f:X=\mathbf{Spec}\,\mathcal A\to Y$ is an affine morphism, and $\mathcal A\cong f_*\mathcal O_X$. Conversely, if $f:X\to Y$ is an affine morphism, then $\mathcal A=f_*\mathcal O_X$ is a quasi-coherent sheaf of $\mathcal O_Y$-algebras, and $X\cong\mathbf{Spec}\,\mathcal A$ (over $Y$).
\item Let $f:X\to Y$ be an affine morphism, and let $\mathcal A=f_*\mathcal O_X$. Show that $f_*$ induces an equivalence of categories from the category of quasi-coherent $\mathcal O_X$-modules to the category of quasi-coherent $\mathcal A$-modules (i.e., quasi-coherent $\mathcal O_Y$-modules having a structure of $\mathcal A$-module). [\emph{Hint}: For any quasi-coherent $\mathcal A$-module $\mathcal M$, construct a quasi-coherent $\mathcal O_X$-module $\widetilde{\mathcal M}$, and show that the functors $f_*$ and $\widetilde{\ }$ are inverse to each other.]
\end{enumerate}}
\pf
\textbf{(a)} ($\Leftarrow$) is trivial. ($\Rightarrow$) Let $\{V_i\}$ be an affine cover of $Y$ with $f^{-1}V_i$ affine, and let $V=\Spec R$ be any open affine of $Y$. Let $y\in V\cap V_i$. Choose $h\in\mathcal O(V_i)$ with $y\in D_{V_i}(h)\subseteq V\cap V_i$; then $f^{-1}(D_{V_i}(h))=D(f^\#h)$ is a basic open subset of the affine scheme $f^{-1}V_i$, hence affine. Now choose $g\in R=\mathcal O(V)$ with $y\in D_V(g)\subseteq D_{V_i}(h)$; $D_V(g)=D(g|_{D_{V_i}(h)})$ is a basic open subset of the affine scheme $D_{V_i}(h)$, so $f^{-1}(D_V(g))=D(f^\#g)$ is a basic open of the affine scheme $f^{-1}(D_{V_i}(h))$, hence affine. Thus $V$ is covered by basic opens $D(g)$ with $f^{-1}(D(g))$ affine. As $V$ is quasi-compact, $V=D(g_1)\cup\dots\cup D(g_m)$, so $(g_1,\dots,g_m)=R$. Put $W=f^{-1}V$ and $a_k=f^\#(g_k)\in\Gamma(W,\mathcal O_W)$. Then $a_1,\dots,a_m$ generate the unit ideal of $\Gamma(W,\mathcal O_W)$ (the same relation $\sum c_kg_k=1$ maps to $\sum f^\#(c_k)a_k=1$), and $W_{a_k}=f^{-1}(D(g_k))$ is affine. By Ex.\ 2.17(b), $W=f^{-1}(V)$ is affine.

\textbf{(b)} \emph{Quasi-compact:} $f$ affine; the affine cover $\{V_i\}$ of $Y$ has $f^{-1}V_i$ affine, hence quasi-compact (Ex.\ 3.2). \emph{Separated:} the preimage of $V_i$ in $X\times_YX$ is $f^{-1}V_i\times_{V_i}f^{-1}V_i=\Spec(B_i\otimes_{A_i}B_i)$ ($V_i=\Spec A_i$, $f^{-1}V_i=\Spec B_i$), these cover $X\times_YX$, and the restriction of $\Delta$ over each is the morphism $\Spec B_i\to\Spec(B_i\otimes_{A_i}B_i)$ induced by the multiplication $B_i\otimes_{A_i}B_i\to B_i$, which is surjective, hence a closed immersion. Being a closed immersion is local on the target, so $\Delta$ is a closed immersion. \emph{Finite $\Rightarrow$ affine:} by definition, $Y$ has an open affine cover $V_i=\Spec A_i$ with $f^{-1}(V_i)=\Spec B_i$ affine ($B_i$ a finite $A_i$-module); so $f$ is affine.

\textbf{(c)} For an open affine $V\subseteq Y$ let $X_V=\Spec\mathcal A(V)$ with $f_V:X_V\to V=\Spec\mathcal O_Y(V)$ induced by $\mathcal O_Y(V)\to\mathcal A(V)$. If $U\subseteq V$ are open affines, then since $\mathcal A|_V=\widetilde{\mathcal A(V)}$ (quasi-coherent on the affine $V$, 5.4/5.5) and $\widetilde M|_U=(M\otimes_{\mathcal O(V)}\mathcal O(U))^\sim$ (5.2(e) applied to the open immersion $U\to V$; see also 5.1(c)), we get $\mathcal A(U)=\mathcal A(V)\otimes_{\mathcal O(V)}\mathcal O(U)$, i.e.\ $X_U=X_V\times_VU=f_V^{-1}(U)$, an open subscheme of $X_V$, and the corresponding morphism $X_U\to X_V$ is the open immersion given by the restriction map $\mathcal A(V)\to\mathcal A(U)$.

\emph{Existence.} Glue the schemes $X_V$ ($V$ ranging over open affines of $Y$): for $V,V'$ let $X_{V,V'}=f_V^{-1}(V\cap V')\subseteq X_V$, open. For every open affine $W\subseteq V\cap V'$ the open subscheme $f_V^{-1}(W)$ of $X_V$ and $f_{V'}^{-1}(W)\subseteq X_{V'}$ are both canonically identified with $X_W$ (by the previous paragraph). These identifications are compatible for $W_1\subseteq W_2$ (they come from the restriction maps of the sheaf $\mathcal A$, which compose), so they glue to an isomorphism $\varphi_{V'V}:f_V^{-1}(V\cap V')\to f_{V'}^{-1}(V\cap V')$ of $Y$-schemes. Again by compatibility with restriction, $\varphi_{V''V}=\varphi_{V''V'}\circ\varphi_{V'V}$ on triple overlaps (check on $f_V^{-1}(W)$ for affine $W\subseteq V\cap V'\cap V''$, where all three are the identity of $X_W$). So we can glue the $X_V$ to a scheme $X$ with a morphism $f:X\to Y$ (the $f_V$ agree on overlaps), in which each $X_V$ is an open subscheme. Moreover $f^{-1}(V)=X_V$: a point of $X$ over $V$ lies in some $X_{V'}$, hence in $f_{V'}^{-1}(V\cap V')$, which was identified with $f_V^{-1}(V\cap V')\subseteq X_V$. So $f^{-1}(V)=X_V\cong\Spec\mathcal A(V)$ for every open affine $V$, and for open affines $U\subseteq V$ the inclusion $f^{-1}U\subseteq f^{-1}V$ is $X_U\to X_V$, which corresponds to the restriction $\mathcal A(V)\to\mathcal A(U)$ by construction.

\emph{Uniqueness.} If $(X',f')$ is another such pair, the isomorphisms $f'^{-1}(V)\cong\Spec\mathcal A(V)=f^{-1}(V)$ for open affine $V$ (over $V$) agree on $f'^{-1}(U)$ for affine $U\subseteq V$, by the compatibility with the inclusions $U\subseteq V$; so they glue to an isomorphism $X'\cong X$ over $Y$.

\textbf{(d)} For $X=\mathbf{Spec}\mathcal A$, $f^{-1}(V)=\Spec\mathcal A(V)$ is affine for each open affine $V$, so $f$ is affine. Also $(f_*\mathcal O_X)(V)=\mathcal O_X(f^{-1}V)=\mathcal A(V)$, compatible with the restriction maps (by the compatibility in (c)); as open affines form a basis, this gives an isomorphism of sheaves of $\mathcal O_Y$-algebras $f_*\mathcal O_X\cong\mathcal A$.

Conversely, let $f$ be affine. By (b) it is quasi-compact and separated, and $\mathcal O_X$ is quasi-coherent, so $\mathcal A=f_*\mathcal O_X$ is a quasi-coherent $\mathcal O_Y$-module (5.8(c)) and a sheaf of $\mathcal O_Y$-algebras. For $V\subseteq Y$ open affine, $f^{-1}(V)$ is affine by (a), with ring $\mathcal O_X(f^{-1}V)=\mathcal A(V)$, so $f^{-1}(V)\cong\Spec\mathcal A(V)$ over $V$; for $U\subseteq V$ open affine, $f^{-1}U\subseteq f^{-1}V$ corresponds to the restriction $\mathcal A(V)\to\mathcal A(U)$ (a morphism of affine schemes is determined by the ring homomorphism). By the uniqueness in (c), $X\cong\mathbf{Spec}\mathcal A$ over $Y$.

\textbf{(e)} \emph{The functor $f_*$.} For $\mathcal F$ quasi-coherent on $X$, $f_*\mathcal F$ is quasi-coherent (5.8(c) with (b)) and a module over $f_*\mathcal O_X=\mathcal A$.

\emph{The functor $\widetilde{\ }$.} Let $\mathcal M$ be a quasi-coherent $\mathcal A$-module. For $V\subseteq Y$ open affine, $f^{-1}(V)=\Spec\mathcal A(V)$ and $\mathcal M(V)$ is an $\mathcal A(V)$-module; let $\widetilde{\mathcal M}_V=\widetilde{\mathcal M(V)}$ be the associated quasi-coherent sheaf on $f^{-1}(V)$ (relative to the ring $\mathcal A(V)$). For affine $U\subseteq V$: $\mathcal M(U)=\mathcal M(V)\otimes_{\mathcal O(V)}\mathcal O(U)$ (quasi-coherence of $\mathcal M$) and $\mathcal A(U)=\mathcal A(V)\otimes_{\mathcal O(V)}\mathcal O(U)$, so $\mathcal M(V)\otimes_{\mathcal A(V)}\mathcal A(U)=\mathcal M(V)\otimes_{\mathcal O(V)}\mathcal O(U)=\mathcal M(U)$. By 5.2(e) (for $\Spec\mathcal A(U)=f^{-1}U\hookrightarrow f^{-1}V$ an open immersion) $\widetilde{\mathcal M}_V|_{f^{-1}U}\cong\widetilde{\mathcal M}_U$, canonically and compatibly with further restrictions. Hence the $\widetilde{\mathcal M}_V$ glue to a quasi-coherent $\mathcal O_X$-module $\widetilde{\mathcal M}$ with $\widetilde{\mathcal M}|_{f^{-1}V}=\widetilde{\mathcal M(V)}$; the construction is functorial in $\mathcal M$.

\emph{Inverse functors.} $f_*(\widetilde{\mathcal M})(V)=\widetilde{\mathcal M}(f^{-1}V)=\widetilde{\mathcal M(V)}(\Spec\mathcal A(V))=\mathcal M(V)$ (5.1(d)), compatibly with restrictions and with the $\mathcal A$-module structures, for open affine $V$; as the open affines form a basis, $f_*\widetilde{\mathcal M}\cong\mathcal M$ functorially. Conversely for $\mathcal F$ quasi-coherent on $X$ and $V$ open affine in $Y$, $\mathcal F|_{f^{-1}V}=\widetilde P$ with $P=\mathcal F(f^{-1}V)=(f_*\mathcal F)(V)$ an $\mathcal A(V)$-module (5.4/5.5), so $\widetilde{f_*\mathcal F}|_{f^{-1}V}=\widetilde{(f_*\mathcal F)(V)}=\mathcal F|_{f^{-1}V}$; these identifications are compatible on overlaps and natural, so $\widetilde{f_*\mathcal F}\cong\mathcal F$ functorially. Therefore $f_*$ is an equivalence with quasi-inverse $\widetilde{\ }$. \hfill$\square$

% ---------------------------------------------------------------- 5.18
\exercise{5.18}{\emph{Vector Bundles.} Let $Y$ be a scheme. A (geometric) \emph{vector bundle} of rank $n$ over $Y$ is a scheme $X$ and a morphism $f:X\to Y$, together with additional data consisting of an open covering $\{U_i\}$ of $Y$, and isomorphisms $\psi_i:f^{-1}(U_i)\to\mathbf A^n_{U_i}$, such that for any $i,j$, and for any open affine subset $V=\Spec A\subseteq U_i\cap U_j$, the automorphism $\psi=\psi_j\circ\psi_i^{-1}$ of $\mathbf A^n_V=\Spec A[x_1,\dots,x_n]$ is given by a \emph{linear} automorphism $\theta$ of $A[x_1,\dots,x_n]$, i.e., $\theta(a)=a$ for any $a\in A$, and $\theta(x_i)=\sum a_{ij}x_j$ for suitable $a_{ij}\in A$.

An \emph{isomorphism} $g:(X,f,\{U_i\},\{\psi_i\})\to(X',f',\{U'_i\},\{\psi'_i\})$ of one vector bundle of rank $n$ to another is an isomorphism $g:X\to X'$ of the underlying schemes, such that $f=f'\circ g$, and such that $X,f$, together with the covering of $Y$ consisting of all the $U_i$ and $U'_i$, and the isomorphisms $\psi_i$ and $\psi'_i\circ g$, is also a vector bundle structure on $X$.
\begin{enumerate}[label=(\alph*),leftmargin=*]
\item Let $\mathcal E$ be a locally free sheaf of rank $n$ on a scheme $Y$. Let $S(\mathcal E)$ be the symmetric algebra on $\mathcal E$, and let $X=\mathbf{Spec}\,S(\mathcal E)$, with the projection morphism $f:X\to Y$. For each open affine subset $U\subseteq Y$ for which $\mathcal E|_U$ is free, choose a basis of $\mathcal E|_U$, and let $\psi:f^{-1}(U)\to\mathbf A^n_U$ be the isomorphism resulting from the identification of $S(\mathcal E(U))$ with $\mathcal O(U)[x_1,\dots,x_n]$. Then $(X,f,\{U\},\{\psi\})$ is a vector bundle of rank $n$ over $Y$, which (up to isomorphism) does not depend on the bases of the $\mathcal E|_U$ chosen. We call it the \emph{geometric vector bundle associated to $\mathcal E$}, and denote it by $\mathbf V(\mathcal E)$.
\item For any morphism $f:X\to Y$, a \emph{section of $f$ over an open set $U\subseteq Y$} is a morphism $s:U\to X$ such that $f\circ s=\mathrm{id}_U$. It is clear how to restrict sections to smaller open sets, or how to glue them together, so we have a sheaf of sets on $Y$, which we denote $\mathcal S(X/Y)$. If $f:X\to Y$ is a vector bundle of rank $n$, show that the sheaf of sections $\mathcal S(X/Y)$ has a natural structure of an $\mathcal O_Y$-module, which is locally free of rank $n$. [\emph{Hint}: It is enough to define the module structure locally, so we can assume $Y=\Spec A$ affine and $X=\mathbf A^n_Y$. Then a section $s:Y\to X$ comes from an $A$-algebra homomorphism $\theta:A[x_1,\dots,x_n]\to A$, which in turn determines an ordered $n$-tuple $\langle\theta(x_1),\dots,\theta(x_n)\rangle$ of elements of $A$. Use this correspondence between sections $s$ and $n$-tuples of elements of $A$ to define the module structure.]
\item Again let $\mathcal E$ be a locally free sheaf of rank $n$ on $Y$, let $X=\mathbf V(\mathcal E)$, and let $\mathcal S=\mathcal S(X/Y)$ be the sheaf of sections of $X$ over $Y$. Show that $\mathcal S\cong\mathcal E^\vee$, as follows. Given a section $s\in\Gamma(V,\mathcal E^\vee)$ over an open set $V$, we think of $s$ as an element of $\operatorname{Hom}(\mathcal E|_V,\mathcal O_V)$. So $s$ determines an $\mathcal O_V$-algebra homomorphism $S(\mathcal E|_V)\to\mathcal O_V$. This determines a morphism of spectra $V=\mathbf{Spec}\,\mathcal O_V\to\mathbf{Spec}\,S(\mathcal E|_V)=f^{-1}(V)$, which is a section of $X/Y$. Show that this construction gives an isomorphism of $\mathcal E^\vee$ to $\mathcal S$.
\item Summing up, show that we have established a one-to-one correspondence between isomorphism classes of locally free sheaves of rank $n$ on $Y$, and isomorphism classes of vector bundles of rank $n$ over $Y$. Because of this, we sometimes use the terms ``locally free sheaf'' and ``vector bundle'' interchangeably, if no confusion seems likely to result.
\end{enumerate}}
\pf
We use freely: morphisms $Z\to\mathbf A^n_W=W\times\Spec\mathbf Z[x_1,\dots,x_n]$ over $W$ correspond bijectively to $n$-tuples of elements of $\Gamma(Z,\mathcal O_Z)$ (Ex.\ 2.4 applied to $\Spec\mathbf Z[x]$).

\textbf{(a)} $S(\mathcal E)$ is a quasi-coherent sheaf of $\mathcal O_Y$-algebras: for $U=\Spec R$ affine and $\mathcal E|_U=\widetilde N$, we have for each $g\in R$ that $S(\mathcal E(D(g)))=S_{R_g}(N_g)=S_R(N)_g$, and the stalks are $S(\mathcal E)_\fp=S(N_\fp)=S_R(N)_\fp$; hence $S(\mathcal E)|_U=\widetilde{S_R(N)}$ and $S(\mathcal E)(U)=S(\mathcal E(U))$. So $X=\mathbf{Spec}\,S(\mathcal E)$ exists (5.17(c)) and for $U$ affine with $\mathcal E|_U$ free on $e_1,\dots,e_n$, $f^{-1}(U)=\Spec S(\mathcal E(U))=\Spec\mathcal O(U)[e_1,\dots,e_n]=\mathbf A^n_U$; call this $\psi_U$. These $U$ form an open cover of $Y$. Let $V=\Spec A\subseteq U\cap U'$ be affine. The restriction of $\mathcal E$ to $V$ is free with the two bases $\{e_i|_V\}$ and $\{e'_i|_V\}$ coming from $U,U'$; so $e'_j=\sum_ia_{ji}e_i$ with $(a_{ji})\in GL_n(A)$, and $S(\mathcal E(V))=A[e_1,\dots,e_n]=A[e'_1,\dots,e'_n]$. The automorphism $\psi_{U'}\psi_U^{-1}$ of $\mathbf A^n_V$ corresponds to the $A$-algebra automorphism $\theta$ of $A[x_1..x_n]$ with $\theta(x_j)=\sum a_{ji}x_i$ (identifying $x_i\leftrightarrow e_i$, $x_j\leftrightarrow e_j'$), which is linear. Thus we have a vector bundle structure.

\emph{Independence.} If $\{U'\}$ with bases $e'$ is another choice, then $g=\mathrm{id}_X$ with the union of the two coverings and the maps $\psi_U$, $\psi'_{U'}$ is again a vector bundle structure by the same computation (for $V\subseteq U\cap U'$ the transition is given by the change of bases of $\mathcal E(V)$, a linear automorphism). So $\mathrm{id}_X$ is an isomorphism between the two bundles. (An isomorphism $\mathcal E\cong\mathcal E'$ likewise induces $S(\mathcal E)\cong S(\mathcal E')$ and an isomorphism of the bundles, by the same argument; so $\mathbf V(\mathcal E)$ depends up to isomorphism only on the isomorphism class of $\mathcal E$.)

\textbf{(b)} Let $(X,f,\{U_i\},\{\psi_i\})$ be a vector bundle. For $W\subseteq U_i$ open, a section $s:W\to f^{-1}(U_i)\cong\mathbf A^n_{U_i}$ with $fs=\mathrm{incl}$ is the same as a morphism $W\to\mathbf A^n_{U_i}$ over $U_i$, i.e.\ an $n$-tuple $(s^\#(x_1),\dots,s^\#(x_n))\in\mathcal O(W)^n$. Thus $\mathcal S|_{U_i}\cong\mathcal O_{U_i}^n$ as sheaves of sets and we transport the $\mathcal O_{U_i}$-module structure (componentwise addition and scalar multiplication). We must check that the structures from $U_i$ and $U_j$ agree on $U_i\cap U_j$. Both module structures are determined (compatibility with restriction) by their values on affine $V=\Spec A\subseteq U_i\cap U_j$ and sections over open subsets $W\subseteq V$. On such $W$ the coordinates for the two charts are related by the linear automorphism $\theta$ of $A[x_1,\dots,x_n]$ with $\theta(x_j)=\sum_ia_{ji}x_i$: $s^\#(x'_j)=s^\#(\theta(x_j))=\sum_ia_{ji}\,s^\#(x_i)$ with $a_{ji}\in A$. So the identification $\mathcal O(W)^n\to\mathcal O(W)^n$ between the two coordinatizations is given by a matrix with entries in $A$, hence is $\mathcal O(W)$-linear, and the two structures coincide. So $\mathcal S(X/Y)$ is an $\mathcal O_Y$-module, locally isomorphic to $\mathcal O^n$: locally free of rank $n$.

\textbf{(c)} Let $\mathcal E$ be locally free of rank $n$, $X=\mathbf V(\mathcal E)$. For $s\in\mathcal E^\vee(V)=\operatorname{Hom}(\mathcal E|_V,\mathcal O_V)$ let $\varphi_s:S(\mathcal E|_V)\to\mathcal O_V$ be the $\mathcal O_V$-algebra homomorphism extending $s$. For every open affine $W\subseteq V$, $\varphi_s$ gives $S(\mathcal E)(W)=S(\mathcal E(W))\to\mathcal O(W)$, compatible with restrictions, hence a morphism $W\to f^{-1}(W)$ over $W$ (a section) and these glue to a section $\sigma_s:V\to f^{-1}(V)\subseteq X$ of $f$. The map $\Phi:\mathcal E^\vee\to\mathcal S$, $s\mapsto\sigma_s$, is compatible with restriction. It is an isomorphism of $\mathcal O_Y$-modules, which can be checked locally. Let $U=\Spec R$ be an open affine with $\mathcal E|_U$ free on $e_1,\dots,e_n$, $X|_U=\Spec R[e_1,\dots,e_n]$ with coordinates $x_i=e_i$, and let $W\subseteq U$ be open. Then $\mathcal E^\vee|_U$ is free on the dual basis $e_i^*$, and for $s=\sum c_ie_i^*\in\mathcal E^\vee(W)=\mathcal O(W)^n$ with $c_i=s(e_i)\in\mathcal O(W)$, $\varphi_s:\mathcal O(W)[e_1..e_n]\to\mathcal O(W)$ is $e_i\mapsto c_i$, so $\sigma_s^\#(x_i)=c_i$. That is, under the chart $\psi_U$ the section $\sigma_s$ is the $n$-tuple $(c_1,\dots,c_n)$, which is exactly the identification $\mathcal S|_U\cong\mathcal O_U^n$ used to define the module structure in (b). Hence $\Phi|_U:\mathcal E^\vee|_U=\mathcal O_U^n\to\mathcal S|_U\cong\mathcal O_U^n$ is the identity, an $\mathcal O$-linear isomorphism. So $\Phi:\mathcal E^\vee\cong\mathcal S(\mathbf V(\mathcal E)/Y)$.

\textbf{(d)} Let $\mathrm{Loc}_n(Y)$ be the set of isomorphism classes of locally free sheaves of rank $n$ and $\mathrm{Vec}_n(Y)$ the set of isomorphism classes of vector bundles of rank $n$. Define $\alpha:\mathrm{Loc}_n\to\mathrm{Vec}_n$, $[\mathcal E]\mapsto[\mathbf V(\mathcal E)]$ (well defined by (a)) and $\beta:\mathrm{Vec}_n\to\mathrm{Loc}_n$, $[X]\mapsto[\mathcal S(X/Y)^\vee]$. Here $\beta$ is well defined: an isomorphism $g:X\to X'$ of vector bundles induces the bijections $s\mapsto g\circ s$ of sections, which are $\mathcal O$-linear because in the charts $\psi_i$, $\psi'_i\circ g$ (both bundle structures on $X$) the transition maps are linear, as in (b); the dual of a locally free sheaf of rank $n$ is locally free of rank $n$ (5.1).

\emph{$\beta\alpha=\mathrm{id}$:} by (c) and 5.1(a), $\mathcal S(\mathbf V(\mathcal E)/Y)^\vee\cong\mathcal E^{\vee\vee}\cong\mathcal E$. In particular $\alpha$ is injective and $\beta$ surjective.

\emph{$\alpha\beta=\mathrm{id}$:} let $X$ be a vector bundle with charts $\psi_i:f^{-1}(U_i)\to\mathbf A^n_{U_i}$, let $\mathcal S=\mathcal S(X/Y)$ and $\mathcal E=\mathcal S^\vee$ (locally free of rank $n$). We construct an isomorphism $X\cong\mathbf V(\mathcal E)$ of vector bundles.

\emph{Linear functions.} Over $U_i$ let $\varepsilon_1,\dots,\varepsilon_n\in\mathcal S(U_i)$ be the sections with coordinates $(1,0,\dots,0),\dots,(0,\dots,0,1)$ in the chart $\psi_i$; they form a basis of $\mathcal S|_{U_i}\cong\mathcal O^n$, and $x_1,\dots,x_n\in\mathcal O_X(f^{-1}U_i)$ denote the coordinate functions of $\psi_i$. For $V\subseteq U_i$ open and $\lambda\in\mathcal E(V)=\operatorname{Hom}(\mathcal S|_V,\mathcal O_V)$ put $c_i(\lambda)=\sum_k\lambda(\varepsilon_k)\,x_k\in\mathcal O_X(f^{-1}V)$; it is $\mathcal O(V)$-linear in $\lambda$ and compatible with restrictions. It is independent of the chart: let $V=\Spec A\subseteq U_i\cap U_j$ be affine and $\psi_j\psi_i^{-1}$ given by $\theta$, $\theta(x'_j)=\sum_ia_{ji}x_i$. The matrix $(a_{ji})$ lies in $GL_n(A)$: $\theta$ is an automorphism preserving degrees, so it is surjective on the degree-$1$ part, i.e.\ $(a_{ji})$ is a surjective endomorphism of $A^n$, hence invertible. A section $s$ with $i$-coordinates $u_i=s^\#(x_i)$ has $j$-coordinates $u'_j=s^\#(x'_j)=\sum_ia_{ji}u_i$. So the unit section $\varepsilon'_l$ of the chart $j$ satisfies $\sum_ia_{ji}u_i=\delta_{jl}$, i.e.\ $u_i=(a^{-1})_{il}$ and $\varepsilon'_l=\sum_i(a^{-1})_{il}\varepsilon_i$. Therefore
\[
\sum_l\lambda(\varepsilon'_l)x'_l=\sum_{l,i,j}(a^{-1})_{il}\,\lambda(\varepsilon_i)\,a_{lj}\,x_j=\sum_i\lambda(\varepsilon_i)x_i .
\]
So the $c_i$ glue (they agree on affine opens, hence on all opens) to an $\mathcal O_Y$-linear morphism $c:\mathcal E\to f_*\mathcal O_X$. By the universal property of the symmetric algebra, $c$ extends to a homomorphism of $\mathcal O_Y$-algebras $\gamma:S(\mathcal E)\to f_*\mathcal O_X$.

\emph{The isomorphism.} The morphism $f$ is affine (for every open affine $V\subseteq U_i$, $f^{-1}(V)\cong\mathbf A^n_V$ is affine, and these $V$ cover $Y$), so $f_*\mathcal O_X$ is a quasi-coherent $\mathcal O_Y$-algebra and $X\cong\mathbf{Spec}f_*\mathcal O_X$ (5.17(d)). Hence $\gamma$ induces a morphism $g:X\to\mathbf{Spec}\,S(\mathcal E)=\mathbf V(\mathcal E)$ over $Y$. Over an open affine $U\subseteq U_i$, $\mathcal E|_U$ is free with the basis $\varepsilon_k^*$ dual to $\varepsilon_k$; use it to define the chart $\psi^{\mathcal E}_U:\mathbf V(\mathcal E)|_U\cong\mathbf A^n_U$ of (a), so that the coordinate functions are $e_k=\varepsilon_k^*$. Then $\gamma:\mathcal O(U)[e_1,\dots,e_n]\to\mathcal O_X(f^{-1}U)=\mathcal O(U)[x_1,\dots,x_n]$ sends $e_k\mapsto c(\varepsilon_k^*)=\sum_l\varepsilon_k^*(\varepsilon_l)x_l=x_k$, an isomorphism; i.e.\ over $U$, $g$ is an isomorphism and $\psi^{\mathcal E}_U\circ g=\psi_i|_{f^{-1}U}$. Thus $g$ is an isomorphism of schemes over $Y$, and the vector bundle structure of $X$ together with the charts $\psi^{\mathcal E}_{U'}\circ g$ ($U'$ affine with $\mathcal E|_{U'}$ free, with arbitrary bases) is again a vector bundle structure: for $\psi_i$ vs.\ $\psi^{\mathcal E}_{U'}\circ g$ over an affine $V\subseteq U_i\cap U'$ the transition is determined by the change of basis of the free module $\mathcal E(V)$ from $\varepsilon^*$ to the chosen basis of $\mathcal E|_{U'}$, hence linear, as in (a). So $g$ is an isomorphism of vector bundles $X\cong\mathbf V(\mathcal E)$, and $\alpha\beta[X]=[X]$.

Thus $\alpha$ and $\beta$ are inverse bijections. (Note that the sheaf of sections of $\mathbf V(\mathcal E)$ is $\mathcal E^\vee$, not $\mathcal E$: in Hartshorne's convention $\mathbf V(\mathcal E)$ corresponds to the \emph{dual} of its sheaf of sections.) \hfill$\square$

\section*{6. Divisors}

Throughout, results of Hartshorne are cited by number (``(6.5)'', ``Ex.~6.2'', ``Thm.~I.7.7'', \dots). Condition $(*)$ is the standing hypothesis of the section: $X$ is a noetherian integral separated scheme which is regular in codimension one. For a prime divisor $Y$ of such an $X$ we write $\eta_Y$ for its generic point, $v_Y$ for the valuation of the discrete valuation ring $\mathcal O_{X,\eta_Y}$, and $K=K(X)$ for the function field. In the exercises that deal with varieties over a field $k$ (6.2, 6.3, 6.5--6.7, 6.9, 6.11, 6.12) $k$ is algebraically closed, as in the text; where an extra hypothesis (characteristic, $\sqrt{-1}\in k$) is needed it is stated. For a linear form or homogeneous polynomial $h$ I write $\bar h$ for its image in a homogeneous coordinate ring.

\exercise{6.1}{Let $X$ be a scheme satisfying $(*)$. Then $X\times\mathbf P^n$ also satisfies $(*)$, and $\operatorname{Cl}(X\times\mathbf P^n)\cong(\operatorname{Cl}X)\times\mathbf Z$.}
\pf
Write $P:=X\times_{\operatorname{Spec}\mathbf Z}\mathbf P^n_{\mathbf Z}=\mathbf P^n_X$, with projection $p:P\to X$ and homogeneous coordinates $x_0,\dots,x_n$. The open sets $U_i=D_+(x_i)$ cover $P$ and $U_i\cong X\times\mathbf A^n=\mathbf A^n_X$.

\emph{Step 1: $P$ satisfies $(*)$.} By (6.6), applied $n$ times, $X\times\mathbf A^n$ satisfies $(*)$, in particular it is integral, noetherian and regular in codimension one. Any two $U_i,U_j$ meet (their intersection contains $X\times\mathbf G_m^{\,n}\neq\emptyset$ over any open of $X$), so $P$ is covered by integral open subschemes that pairwise meet; hence $P$ is irreducible and reduced, i.e.\ integral. It is noetherian (finite type over $X$) and separated ($P\to X$ is separated and $X$ is separated). Regularity in codimension one is a statement about local rings, hence can be checked on the cover by the $U_i$. So $(*)$ holds.

\emph{Step 2: an exact sequence.} Let $H=V_+(x_0)\cong\mathbf P^{n-1}_X$. By Step 1 (for $n-1$; for $n=1$ it is $\cong X$) $H$ is integral, and over an affine open $\operatorname{Spec}A\subseteq X$, in the chart $\operatorname{Spec}A[t_0,t_2,\dots,t_n]\subseteq U_1$ ($t_j=x_j/x_1$), it is cut out by the prime element $t_0$ (the quotient is a polynomial ring over the domain $A$). By Krull's Hauptidealsatz it has codimension $1$; so $H$ is a prime divisor, and $P-H=U_0\cong X\times\mathbf A^n$. By (6.5)(c) there is an exact sequence
\[
\mathbf Z\xrightarrow{\;1\mapsto[H]\;}\operatorname{Cl}P\longrightarrow\operatorname{Cl}(X\times\mathbf A^n)\to0,
\]
and by (6.6), applied $n$ times, $\operatorname{Cl}(X\times\mathbf A^n)\cong\operatorname{Cl}X$.

\emph{Step 3: a retraction $\rho:\operatorname{Cl}P\to\mathbf Z$ with $\rho[H]=1$.} Let $K=K(X)=\mathcal O_{X,\eta}$, $\eta$ the generic point, and let $P_K=\mathbf P^n_K=P\times_X\operatorname{Spec}K$ be the generic fibre. The points of $P_K$ are the points $\xi\in P$ with $p(\xi)=\eta$, and for such $\xi$, $\mathcal O_{P_K,\xi}=\mathcal O_{P,\xi}$ (localization). Hence the prime divisors of $P_K$ correspond bijectively to the prime divisors $W$ of $P$ which \emph{dominate} $X$ (generic point over $\eta$), $W\mapsto W_K$, with the same valuation: $v_{W_K}=v_W$ on $K(P)=K(P_K)$. Define $\rho_0:\operatorname{Div}P\to\operatorname{Div}P_K$ by $W\mapsto W_K$ if $W$ dominates $X$ and $W\mapsto 0$ otherwise. For $f\in K(P)^*$ we get $\rho_0((f))=(f)_{P_K}$, so $\rho_0$ induces $\operatorname{Cl}P\to\operatorname{Cl}P_K\cong\mathbf Z$, the last isomorphism being the degree of (6.4) (valid over the field $K$). The hyperplane $H$ dominates $X$ and $H_K$ is a hyperplane in $\mathbf P^n_K$, of degree $1$; so $\rho[H]=1$.

\emph{Conclusion.} Since $\rho\circ(1\mapsto[H])=\operatorname{id}_{\mathbf Z}$, the map $\mathbf Z\to\operatorname{Cl}P$ of Step 2 is injective and the short exact sequence $0\to\mathbf Z\to\operatorname{Cl}P\to\operatorname{Cl}X\to 0$ splits: $\operatorname{Cl}P=\mathbf Z[H]\oplus\ker\rho\cong\mathbf Z\oplus\operatorname{Cl}X$.\hfill$\square$

\exercise{6.2}{\emph{Varieties in projective space.} Let $k$ be an algebraically closed field and let $X$ be a closed subvariety of $\mathbf P^n_k$ which is nonsingular in codimension one (hence satisfies $(*)$). For any divisor $D=\sum n_iY_i$ on $X$, we define the \emph{degree} of $D$ to be $\sum n_i\deg Y_i$, where $\deg Y_i$ is the degree of $Y_i$ as a projective variety itself (I, Section~7). (Starred in the book.)
\begin{enumerate}[label=(\alph*),leftmargin=*]
\item Let $V$ be an irreducible hypersurface in $\mathbf P^n$ which does not contain $X$, and let $Y_i$ be the irreducible components of $V\cap X$. They all have codimension $1$ in $X$ (I, Ex.~1.8). For each $i$, let $f_i$ be a local equation for $V$ on some open set $U_i$ on which $Y_i\cap U_i\ne\emptyset$, and let $n_i=v_{Y_i}(\bar f_i)$, where $\bar f_i$ is the restriction of $f_i$ to $U_i\cap X$. Then we define the \emph{divisor} $V.X$ to be $\sum n_iY_i$. Extend by linearity, and show that this gives a well-defined homomorphism from the subgroup of $\operatorname{Div}\mathbf P^n$ consisting of divisors none of whose components contain $X$, to $\operatorname{Div}X$.
\item If $D$ is a principal divisor on $\mathbf P^n$, for which $D.X$ is defined as in (a), show that $D.X$ is principal on $X$. Thus we get a homomorphism $\operatorname{Cl}\mathbf P^n\to\operatorname{Cl}X$.
\item Show that the integer $n_i$ defined in (a) is the same as the intersection multiplicity $i(X,V;Y_i)$ defined in (I, Section~7). Then use the generalized Bezout theorem (I, 7.7) to show that for any divisor $D$ on $\mathbf P^n$, none of whose components contain $X$, $\deg(D.X)=(\deg D)\cdot(\deg X)$.
\item If $D$ is a principal divisor on $X$, show that there is a rational function $f$ on $\mathbf P^n$ such that $D=(f).X$. Conclude that $\deg D=0$. Thus the degree function defines a homomorphism $\deg:\operatorname{Cl}X\to\mathbf Z$. (This gives another proof of (6.10), since any complete nonsingular curve is projective.) Finally, there is a commutative diagram with top row $\operatorname{Cl}\mathbf P^n\to\operatorname{Cl}X$, left vertical map $\deg:\operatorname{Cl}\mathbf P^n\xrightarrow{\ \cong\ }\mathbf Z$, right vertical map $\deg:\operatorname{Cl}X\to\mathbf Z$, and bottom row $\cdot(\deg X):\mathbf Z\to\mathbf Z$, and in particular we see that the map $\operatorname{Cl}\mathbf P^n\to\operatorname{Cl}X$ is injective.
\end{enumerate}}
\pf
Notation: $S=k[x_0,\dots,x_n]$, $I_X\subset S$ the (prime, homogeneous) ideal of $X$, $S_X=S/I_X$ the homogeneous coordinate ring; we assume $\dim X\ge1$ (otherwise $\operatorname{Div}X=0$ and there is nothing to prove). For a prime divisor $Y$ of $X$ let $\mathfrak p_Y\subset S_X$ be its homogeneous prime. Then $Y\cap D_+(x_i)\neq\emptyset$ iff $\bar x_i\notin\mathfrak p_Y$ iff $\eta_Y\in D_+(x_i)$. Let $\mathcal G\subseteq\operatorname{Div}\mathbf P^n$ be the subgroup of divisors none of whose components contain $X$; it is free on the irreducible hypersurfaces $V\not\supseteq X$.

\textbf{(a)} An irreducible hypersurface $V=V_+(h)$, $h$ irreducible of degree $d$. On $D_+(x_i)$ the ideal sheaf of $V$ is generated by $h/x_i^{d}$, a nonzerodivisor, so locally $V$ is given by one equation (this is also clear since $\mathbf P^n$ is locally factorial, (6.11)). Since $V\not\supseteq X$, $h\notin I_X$, so $\bar h\ne0$ in the domain $S_X$ and every $\bar f_i\neq0$.

\emph{Components.} $V\cap X\subsetneq X$ is a proper closed subset of the variety $X$, so each of its components has dimension $\le\dim X-1$; by I.Ex.1.8 it is $\ge\dim X-1$. So the $Y_i$ are prime divisors of $X$.

\emph{Independence of choices.} Let $f$ and $f'$ be two local equations of $V$ defined on open sets $U,U'$ both containing $\eta=\eta_{Y_i}$. Then $f,f'$ both generate the stalk of the ideal sheaf $\mathcal I_V$ at $\eta$, so $f=uf'$ with $u\in\mathcal O_{\mathbf P^n,\eta}^*$. Restricting to $X$ gives $\bar f=\bar u\bar f'$ with $\bar u\in\mathcal O_{X,\eta}^*$, so $v_{Y_i}(\bar f)=v_{Y_i}(\bar f')$. Moreover $\eta\in V$ gives $f\in\mathfrak m_{\mathbf P^n,\eta}$, so $n_i=v_{Y_i}(\bar f_i)\ge1$. Thus $n_i$ depends only on $V$ and $Y_i$, and $V.X:=\sum n_iY_i$ is a well-defined effective divisor. Defining $D.X=\sum m_j(V_j.X)$ for $D=\sum m_jV_j\in\mathcal G$ is then a homomorphism $\mathcal G\to\operatorname{Div}X$, because $\mathcal G$ is free on the $V_j$. Note that if $Y$ is a prime divisor of $X$ and $\eta_Y\notin V$, then $\bar f$ is a unit at $\eta_Y$, so the \emph{coefficient of $Y$ in $V.X$ is $v_Y(\overline{h/x_i^d})$ for every $i$ with $\eta_Y\in D_+(x_i)$}, whether or not $Y\subseteq V$ (it is $0$ in the latter case). I record this:
\[
\operatorname{coeff}_Y(V.X)=v_Y\bigl(\bar h/\bar x_i^{\,d}\bigr)\quad\text{for all }i\text{ with }\bar x_i\notin\mathfrak p_Y.\qquad(1)
\]

\textbf{(b)} \emph{Lemma.} Let $A,B\in S$ be homogeneous of the same degree $d$ with $A,B\notin I_X$, and $f=A/B\in K(\mathbf P^n)^*$. Factor $A=c\prod h_j^{a_j}$, $B=c'\prod h_j^{b_j}$ into irreducibles and set $(A)=\sum a_jV_+(h_j)$, $(B)=\sum b_jV_+(h_j)$. Then $(f)=(A)-(B)\in\mathcal G$ (no $h_j$ lies in $I_X$, since $I_X$ is prime and $A,B\notin I_X$), and
\[
(f).X=\bigl(\bar A/\bar B\bigr)\quad\text{as divisors on }X,
\]
where $\bar A/\bar B\in K(X)^*$. Indeed, let $Y$ be a prime divisor of $X$ and choose $i$ with $\bar x_i\notin\mathfrak p_Y$. By (1),
\[
\operatorname{coeff}_Y((A).X)=\sum_ja_jv_Y(\bar h_j/\bar x_i^{d_j})=v_Y\Bigl(\prod(\bar h_j/\bar x_i^{d_j})^{a_j}\Bigr)=v_Y(\bar A/\bar x_i^d),
\]
(using $\sum a_jd_j=\deg A=d$ and that $c$ is a unit), and similarly for $B$. Subtracting, $\operatorname{coeff}_Y((f).X)=v_Y(\bar A/\bar B)$.

Now let $D\in\mathcal G$ be principal, $D=(f)$. Write $f=A/B$ with $A,B$ coprime forms of the same degree. Then $(A)$ and $(B)$ are effective without common components and $D=(A)-(B)$, so the components of $(A),(B)$ are those of $D$; none contains $X$, so $A,B\notin I_X$, and the Lemma gives $D.X=(\bar A/\bar B)$, principal. So $\mathcal G\cap\{\text{principal}\}\to\{\text{principal divisors of }X\}$, and we get $\mathcal G/(\mathcal G\cap\operatorname{Prin})\to\operatorname{Cl}X$. By (6.4)(a), every class of $\operatorname{Cl}\mathbf P^n$ contains $d\cdot H'$ with $H'$ any hyperplane, and we can choose $H'=V_+(x_i)$ with $x_i\notin I_X$ (possible since $X\ne\emptyset$); so $\mathcal G\to\operatorname{Cl}\mathbf P^n$ is surjective with kernel $\mathcal G\cap\operatorname{Prin}$, i.e.\ $\operatorname{Cl}\mathbf P^n=\mathcal G/(\mathcal G\cap\operatorname{Prin})$. We obtain the homomorphism $\varphi:\operatorname{Cl}\mathbf P^n\to\operatorname{Cl}X$, $[D]\mapsto[D.X]$ for $D\in\mathcal G$.

\textbf{(c)} \emph{Local computation.} Let $Y$ be a prime divisor of $X$, $\mathfrak p=\mathfrak p_Y$ (homogeneous prime of $S_X$, of height one), $x\in S_1$ with $\bar x\notin\mathfrak p$, $A=(S_X)_{(x)}$ (degree-$0$ part of $(S_X)_x$, the coordinate ring of $X\cap D_+(x)$), $R=\mathcal O_{X,\eta_Y}=A_{\mathfrak q}$ where $\mathfrak q=\mathfrak p(S_X)_x\cap A$. Since $\bar x$ is homogeneous of degree one and a unit, $(S_X)_x=A[x,x^{-1}]$ is a Laurent polynomial ring (every homogeneous element of degree $m$ is $(s/x^m)x^m$ with $s/x^m\in A$, and $x$ is algebraically independent over $A$ by degree reasons). As $\mathfrak p$ is homogeneous, so is the ideal $\mathfrak p(S_X)_x$, whence $\mathfrak p(S_X)_x=\bigoplus_m(\mathfrak p(S_X)_x)_m=\bigoplus_m\mathfrak q x^m=\mathfrak qA[x,x^{-1}]$. Let $\pi$ be a uniformizer of the DVR $R$ and $\kappa=R/\pi R$. Localizing, $T:=(S_X)_{\mathfrak p}=R[x,x^{-1}]_{\pi R[x,x^{-1}]}$. Here $\pi R[x,x^{-1}]$ is prime (quotient $\kappa[x,x^{-1}]$), so $T$ is a noetherian local domain whose maximal ideal $\pi T$ is principal and nonzero, hence \emph{$T$ is a DVR with uniformizer $\pi$}, and $v_T|_R=v_R$. Consequently, for $h\in S_d$ with $\bar h\ne0$ in $S_X$,
\[
\operatorname{length}_{S_{\mathfrak p}}\bigl(S/(I_X+(h))\bigr)_{\mathfrak p}=\operatorname{length}_T\bigl(T/\bar hT\bigr)=v_T(\bar h)=v_R(\bar h/\bar x^d),\qquad(2)
\]
since $(S/(I_X+(h)))_{\mathfrak p}=T/\bar hT$ is a module over the quotient $S_{\mathfrak p}\to T$ of $S_{\mathfrak p}$ (so lengths over $S_{\mathfrak p}$ and $T$ agree), $T/\pi^nT$ has length $n$, and $\bar x$ is a unit of $T$.

Now take $V=V_+(h)$, $h$ irreducible of degree $d$, $V\not\supseteq X$, and $Y=Y_i$ a component of $V\cap X$. The intersection multiplicity of (I, Section~7) is $i(X,V;Y)=\mu_{\mathfrak p}(S/(I_X+I_V))=\operatorname{length}_{S_{\mathfrak p}}(S/(I_X+(h)))_{\mathfrak p}$ ($I_V=(h)$ as $h$ is irreducible). By (2) and (1) this equals $v_Y(\bar h/\bar x^d)=n_i$. So $n_i=i(X,V;Y_i)$.

\emph{Bezout.} Let $V$ be an irreducible hypersurface not containing $X$; by Thm.~I.7.7, $\sum_ii(X,V;Y_i)\deg Y_i=(\deg X)(\deg V)$, where $\deg V=\deg h$ by (I, 7.6(d)). The left side is $\deg(V.X)$ by the first part. For $D=\sum m_jV_j\in\mathcal G$, linearity gives $\deg(D.X)=\sum m_j\deg(V_j.X)=\deg X\sum m_j\deg V_j=(\deg X)(\deg D)$, where $\deg D$ is as in (6.4).

\textbf{(d)} Let $D=(g)$ with $g\in K(X)^*$. Write $g=\bar A/\bar B$ with $A,B$ forms of the same degree, $A,B\notin I_X$ (this is the description of $K(X)$ for a projective variety). Put $f=A/B\in K(\mathbf P^n)^*$. By the Lemma in (b), $(f)\in\mathcal G$ and $(f).X=(g)=D$. By (c) and (6.4)(b), $\deg D=\deg((f).X)=\deg(f)\cdot\deg X=0$. So $\deg$ vanishes on principal divisors and defines $\deg:\operatorname{Cl}X\to\mathbf Z$.

\emph{Diagram.} For $D\in\mathcal G$, (c) says $\deg(\varphi[D])=\deg(D.X)=(\deg X)\deg D$, i.e.\ $\deg_X\circ\varphi=(\deg X)\cdot\deg_{\mathbf P^n}$, which is the commutativity. Since $\deg_{\mathbf P^n}$ is an isomorphism (6.4)(c) and $\deg X>0$ (I, 7.6(a)), multiplication by $\deg X$ is injective, hence $\varphi$ is injective.\hfill$\square$

\exercise{6.3}{\emph{Cones.} In this exercise we compare the class group of a projective variety $V$ and the class group of its cone (I, Ex.~2.10). So let $V$ be a projective variety in $\mathbf P^n$, which is of dimension $\ge1$ and nonsingular in codimension $1$. Let $X=C(V)$ be the affine cone over $V$ in $\mathbf A^{n+1}$, and let $\bar X$ be its projective closure in $\mathbf P^{n+1}$. Let $P\in X$ be the vertex of the cone. (Starred in the book.)
\begin{enumerate}[label=(\alph*),leftmargin=*]
\item Let $\pi:\bar X-P\to V$ be the projection map. Show that $V$ can be covered by open subsets $U_i$ such that $\pi^{-1}(U_i)\cong U_i\times\mathbf A^1$ for each $i$, and then show as in (6.6) that $\pi^*:\operatorname{Cl}V\to\operatorname{Cl}(\bar X-P)$ is an isomorphism. Since $\operatorname{Cl}\bar X\to\operatorname{Cl}(\bar X-P)$ is an isomorphism, we have also $\operatorname{Cl}V\cong\operatorname{Cl}\bar X$.
\item We have $V\subseteq\bar X$ as the hyperplane section at infinity. Show that the class of the divisor $V$ in $\operatorname{Cl}\bar X$ is equal to $\pi^*(\text{class of }V.H)$, where $H$ is any hyperplane of $\mathbf P^n$ not containing $V$. Thus conclude using (6.5) that there is an exact sequence
\[0\to\mathbf Z\to\operatorname{Cl}V\to\operatorname{Cl}X\to0,\]
where the first map sends $1\mapsto V.H$, and the second is $\pi^*$ followed by the restriction to $X-P$ and inclusion in $X$. (The injectivity of the first arrow follows from the previous exercise.)
\item Let $S(V)$ be the homogeneous coordinate ring of $V$ (which is also the affine coordinate ring of $X$). Show that $S(V)$ is a unique factorization domain if and only if (1) $V$ is projectively normal (Ex.~5.14), and (2) $\operatorname{Cl}V\cong\mathbf Z$ and is generated by the class of $V.H$.
\item Let $\mathcal O_P$ be the local ring of $P$ on $X$. Show that the natural restriction map induces an isomorphism $\operatorname{Cl}X\to\operatorname{Cl}(\operatorname{Spec}\mathcal O_P)$.
\end{enumerate}}
\pf
\emph{Setup.} Let $S=S(V)=k[x_0,\dots,x_n]/I(V)$, a graded domain generated by $S_1$ with $S_0=k$, $\mathfrak m=S_+$, $\dim S=\dim V+1\ge2$. Then $X=\operatorname{Spec}S$ and $P$ is the closed point $\mathfrak m$. Let $\mathbf P^{n+1}$ have coordinates $x_0,\dots,x_n,w$, so that $\mathbf A^{n+1}=\{w\ne0\}$ and the hyperplane at infinity $\{w=0\}$ is a copy of $\mathbf P^n\supseteq V$. The homogeneous coordinate ring $S[w]$ of $\bar X=V_+(I(V))$ is a domain, $\bar X\cap\{w\ne0\}=X$ and $\bar X\cap\{w=0\}=V$; $P=(0:\dots:0:1)$ and $\bar X-P=\bigcup_iD_+(x_i)\cap\bar X$. Since $P$ is the only point at which all $x_i$ vanish, $\pi(x:w)=(x)$ is a morphism $\bar X-P\to V$.

\textbf{(a)} Let $U_i=V\cap D_+(x_i)=\operatorname{Spec}S_{(x_i)}$. The ring of $\bar X\cap D_+(x_i)$ is $(S[w]_{x_i})_0=S_{(x_i)}[t_i]$, $t_i=w/x_i$, a polynomial ring ($w$ is algebraically independent over $S$). Hence $\pi^{-1}(U_i)\cong U_i\times\mathbf A^1$, $\pi$ being the projection. Put $L=\bar X-P$; it is a noetherian integral separated scheme, regular in codimension one (locally $U_i\times\mathbf A^1$, and $V$ is regular in codimension one, so (6.6) applies to the local rings), i.e.\ $L$ satisfies $(*)$. Also $\bar X$ satisfies $(*)$: all codimension one points lie in $L$ since $\operatorname{codim}(P,\bar X)=\dim V+1\ge2$. So by (6.5)(b), $\operatorname{Cl}\bar X\to\operatorname{Cl}L$ is an isomorphism.

\emph{The map $\pi^*$.} For a prime divisor $Y\subset V$, $\pi^{-1}(Y)\cap\pi^{-1}(U_i)=(Y\cap U_i)\times\mathbf A^1$ is integral of codimension one in $U_i\times\mathbf A^1$ whenever $Y\cap U_i\ne\emptyset$; any two such open pieces meet, so $\pi^{-1}(Y)$ is a prime divisor of $L$ (``type~1''), with $v_{\pi^{-1}Y}|_{K(V)}=v_Y$ (as in (6.6): $\mathcal O_{L,\pi^{-1}Y}=\mathcal O_{V,Y}[t]_{(\pi_Y)}$, a DVR with the same uniformizer $\pi_Y$). Define $\pi^*:\operatorname{Div}V\to\operatorname{Div}L$ by $Y\mapsto\pi^{-1}(Y)$. The prime divisors of $L$ are of two kinds: type~1, $\pi^{-1}(Y)$, and type~2, those whose generic point maps to the generic point of $V$ (if $Z$ is a prime divisor, $\overline{\pi(Z)}$ is $V$ or a prime divisor of $V$ by a dimension count, $\dim Z=\dim V$ and $Z\subseteq\pi^{-1}(\overline{\pi(Z)})$; in the latter case $Z=\pi^{-1}(\overline{\pi (Z)})$). For $f\in K(V)^*$, $\pi^*f\in K(L)^*$ has $(\pi^*f)=\pi^*((f))$: at type~1 primes the valuation is $v_Y(f)$; at type~2 primes the local ring contains $K(V)$, so $\pi^*f$ is a unit. Thus $\pi^*:\operatorname{Cl}V\to\operatorname{Cl}L$ is a well-defined homomorphism.

\emph{Injective.} Suppose $\pi^*D=(g)$, $g\in K(L)^*=K(V)(t_i)$. On $\pi^{-1}(U_i)=U_i\times\mathbf A^1$ the type~2 primes correspond to the irreducible polynomials of $K(V)[t_i]$ (the closed points of the generic fibre $\operatorname{Spec}K(V)[t_i]$), and for such a prime $Z_q=(q)$, $v_{Z_q}(g)$ is the exponent of $q$ in $g=a/b$ ($a,b\in K(V)[t_i]$ coprime). Since $(g)=\pi^*D$ has no type~2 component, $a,b$ have no irreducible factor of positive degree, so $g\in K(V)^*$. Then $\pi^*(D-(g))=0$, so $D=(g)$ because $\pi^*$ is injective on $\operatorname{Div}$ (distinct $Y$ give distinct $\pi^{-1}Y$).

\emph{Surjective.} Type~1 divisors are in the image. If $Z$ is of type~2, in the chart $U_i\times\mathbf A^1$ it is $Z_q$ for an irreducible $q\in K(V)[t_i]$; viewing $q\in K(L)^*$, $v_{Z_q}(q)=1$, $v_{Z'}(q)=0$ for every other type~2 prime $Z'$, so $Z-(q)$ is a combination of type~1 primes. Hence $\pi^*$ is onto. Therefore $\pi^*:\operatorname{Cl}V\xrightarrow{\ \cong\ }\operatorname{Cl}L\xleftarrow{\ \cong\ }\operatorname{Cl}\bar X$.

\textbf{(b)} Choose a linear form $\ell\in k[x_0,\dots,x_n]$ with $\ell\notin I(V)$, $H=V_+(\ell)$. In $K(\bar X)$ consider $f=\ell/w$. On $\pi^{-1}(U_i)=U_i\times\mathbf A^1$ we have $f=\pi^*(\ell/x_i)\cdot t_i^{-1}$. The divisor of $t_i$ is $\{t_i=0\}=V\cap\pi^{-1}(U_i)$ with multiplicity one ($t_i$ is a coordinate), and $(\pi^*(\ell/x_i))=\pi^*((\ell/x_i))$ with $(\ell/x_i)|_{U_i}=(V.H)|_{U_i}$ by the definition of $V.H$ in Ex.~6.2(a). So $(f)=\pi^*(V.H)-V$ on every $\pi^{-1}(U_i)$, hence on $L$. Therefore $[V]=\pi^*[V.H]$ in $\operatorname{Cl}L\cong\operatorname{Cl}\bar X$.

Now $V$ is an irreducible closed subset of codimension one of $\bar X$ with $\bar X-V=X$. By (6.5)(c),
\[
\mathbf Z\xrightarrow{1\mapsto[V]}\operatorname{Cl}\bar X\longrightarrow\operatorname{Cl}X\to0
\]
is exact. Composing with the isomorphism $\pi^*:\operatorname{Cl}V\cong\operatorname{Cl}\bar X$ of (a) the first map becomes $1\mapsto[V.H]$, and the second map $\operatorname{Cl}\bar X\to\operatorname{Cl}X$ (restriction) is, after restricting further to $X-P$, the map described in the statement (note $\operatorname{Cl}X\cong\operatorname{Cl}(X-P)$ by (6.5)(b) since $\operatorname{codim}(P,X)\ge2$). Finally $1\mapsto[V.H]$ is injective: by Ex.~6.2(c),(d) applied to $V\subseteq\mathbf P^n$ and the hyperplane $H\not\supseteq V$ (of degree one), $\deg(V.H)=\deg V>0$, so $m[V.H]=0$ forces $m=0$. This gives the exact sequence $0\to\mathbf Z\to\operatorname{Cl}V\to\operatorname{Cl}X\to0$.

\textbf{(c)} By (6.2), the noetherian domain $S$ is a UFD iff $X=\operatorname{Spec}S$ is normal and $\operatorname{Cl}X=0$. $X$ is normal iff $S$ is integrally closed, i.e.\ iff $V$ is projectively normal (Ex.~5.14: $V$ is projectively normal iff $S(V)$ is an integrally closed domain). By (b), $\operatorname{Cl}X=0$ iff $\operatorname{Cl}V$ is generated by the image of $\mathbf Z$, i.e.\ by $[V.H]$; since $\mathbf Z\to\operatorname{Cl}V$ is injective, this holds iff $\operatorname{Cl}V\cong\mathbf Z$ with generator $[V.H]$. This is the claim.

\textbf{(d)} Let $\mathcal O_P=S_{\mathfrak m}$. The scheme $\operatorname{Spec}\mathcal O_P$ is noetherian, integral, separated and regular in codimension one (its local rings are local rings of $X$), so satisfies $(*)$; its prime divisors are $Y\cap\operatorname{Spec}\mathcal O_P$ for prime divisors $Y\ni P$ of $X$, i.e.\ correspond to height-one primes $\mathfrak p\subseteq\mathfrak m$ of $S$, and for these $\mathcal O_{X,Y}$ is also the local ring of $\operatorname{Spec}\mathcal O_P$ at the corresponding point. Define $\operatorname{res}:\operatorname{Div}X\to\operatorname{Div}\operatorname{Spec}\mathcal O_P$ by $Y\mapsto Y\cap\operatorname{Spec}\mathcal O_P$ (which is $0$ if $P\notin Y$). It is surjective (take closures), and $\operatorname{res}((f))=(f)$ for $f\in K^*$ since valuations are unchanged. So res induces a surjection $r:\operatorname{Cl}X\to\operatorname{Cl}\operatorname{Spec}\mathcal O_P$. The kernel of $r$: if $\operatorname{res}D=(f)$ then $D-(f)$ lies in $\ker(\operatorname{res})$, the group generated by prime divisors \emph{not} through $P$. So it suffices to prove:

\emph{Claim.} Every prime divisor $Y\not\ni P$ of $X$ is principal.

Granting the Claim, $D-(f)$ is principal, so $[D]=0$, and $r$ is injective. 

\emph{Proof of the Claim.} Let $\mathfrak p\subset S$ be the height-one prime of $Y$; $P\notin Y$ means $\mathfrak p\not\subseteq\mathfrak m$. Let $\mathcal H$ be the set of nonzero homogeneous elements of $S$. We use two standard facts: a minimal prime over a homogeneous ideal of a graded ring is homogeneous, and a homogeneous proper ideal of $S$ is contained in $\mathfrak m$ (its degree-$0$ part is a proper ideal of $S_0=k$).

(i) \emph{$\mathfrak p\cap\mathcal H=\emptyset$.} If $h\in\mathfrak p\cap\mathcal H$, then $\mathfrak p$ (of height one in the domain $S$) is minimal over $(h)$, hence homogeneous, hence $\subseteq\mathfrak m$; contradiction. Likewise, $\mathfrak p\not\subseteq\mathfrak m$ gives an element $g=g_0+g_1+\dots+g_D\in\mathfrak p$ ($g_j\in S_j$) with $g_0\ne0$; we normalize $g_0=1$.

(ii) \emph{The element $G$.} Let $S_h=\mathcal H^{-1}S$ and fix $0\ne x\in S_1$. As in Ex.~6.2(c), $S_h=K_0[x,x^{-1}]$ where $K_0=(S_h)_0$ is a field (and $x$ is transcendental over it). The prime $\mathfrak pS_h$ is proper (by (i)) and nonzero, hence a maximal ideal of the PID $K_0[x,x^{-1}]$. Now $g=\sum a_jx^j\in K_0[x]$ with $a_j=g_j/x^j\in A:=S_{(x)}\subseteq K_0$, $a_0=1$, and $g\in\mathfrak pS_h$. Let $G\in K_0[x]$ be a generator of $\mathfrak pS_h$ which is a polynomial with $G(0)=1$: take an irreducible factor of $g$ in $K_0[x]$ lying in $\mathfrak pS_h$ (it is not associate to $x$ as $g(0)\ne0$), multiplied by a constant. It is unique: two generators differ by a unit $\lambda x^m$ of $K_0[x,x^{-1}]$, and both being polynomials with constant term $1$ forces $m=0,\lambda=1$. By this uniqueness, $G$ does not depend on the choice of $x$ (replacing $x$ by $x'=ux$, $u\in K_0^*$, the polynomial $G(x'/u)$ is again a polynomial in $x'$ with constant term $1$ generating $\mathfrak pS_h$). Moreover $G\mid g$ in $K_0[x]$.

(iii) \emph{$(G)=Y$ as divisors on $X$.} Let $\mathfrak q\subset S$ be any height-one prime ($Z=V(\mathfrak q)$). 

\emph{Case $\mathfrak q\cap\mathcal H=\emptyset$.} Then $S_h\subseteq S_{\mathfrak q}$ and $S_{\mathfrak q}=(S_h)_{\mathfrak qS_h}$, with $\mathfrak qS_h$ a maximal ideal of $K_0[x,x^{-1}]$. If $\mathfrak q=\mathfrak p$, $S_{\mathfrak p}$ is the localization of the PID at the maximal ideal $(G)$, so $v_Y(G)=1$. If $\mathfrak q\ne\mathfrak p$, then $\mathfrak qS_h\neq\mathfrak pS_h$ (they contract to $\mathfrak q,\mathfrak p$), so $G\notin\mathfrak qS_h$ and $v_Z(G)=0$.

\emph{Case $\mathfrak q\cap\mathcal H\neq\emptyset$.} Then $\mathfrak q$ is homogeneous (minimal over a homogeneous element) and $\mathfrak q\ne\mathfrak m$ (as $\operatorname{ht}\mathfrak m=\dim S\ge2$), so some $x'\in S_1$ is not in $\mathfrak q$ (else $\mathfrak m\subseteq\mathfrak q$). Use this $x'$ in place of $x$: $G=\sum_{j=0}^ec_jx'^{\,j}$ with $c_j\in K_0$, $c_0=1$, and $G\mid g$ in $K_0[x']$, $g=\sum a_jx'^j$, $a_j=g_j/x'^j\in A'=S_{(x')}$. With $\tau$ an indeterminate, $g^*(\tau)=\tau^D+a_1\tau^{D-1}+\dots+a_D$ is monic over $A'$ and divisible by $G^*(\tau)=\tau^e+c_1\tau^{e-1}+\dots+c_e\in K_0[\tau]$ (reverse the polynomials). Hence every root of $G^*$ in an algebraic closure of $K_0$ is integral over $A'$, and so are the coefficients $c_j$ (symmetric functions of the roots). As in Ex.~6.2(c) we have $\mathfrak q(S_{x'})=\mathfrak q_0S_{x'}$ with $\mathfrak q_0=\mathfrak qS_{x'}\cap A'$ prime (of height one, since $S_{x'}=A'[x',x'^{-1}]$ and $\mathfrak qS_{x'}=\mathfrak q_0A'[x',x'^{-1}]$), $R=A'_{\mathfrak q_0}$ is the local ring of $V$ along the prime divisor $V_+(\mathfrak q)$ of $V$, hence a DVR (as $V$ is nonsingular in codimension one), with uniformizer $\pi$, and $T:=S_{\mathfrak q}=\mathcal O_{X,Z}=R[x',x'^{-1}]_{(\pi)}$ is a DVR with uniformizer $\pi$ and residue field $\operatorname{Frac}(\kappa[x',x'^{-1}])$, $\kappa=R/\pi R$. The $c_j\in K_0$ are integral over $A'\subseteq T$ and $T$ is integrally closed with fraction field $K\supseteq K_0$, so $c_j\in T\cap K_0$. Now $R\to T$ is a flat local homomorphism ($R[x',x'^{-1}]$ is free over $R$ and $T$ is a localization of it), hence faithfully flat, so $aT\cap R=aR$ for $a\in R$; therefore $T\cap K_0=R$ (write $c\in T\cap K_0$ as $a/b$, $a,b\in R$; then $a\in bT\cap R=bR$). So $c_j\in R$ for all $j$, and $G\in T$ has image $\sum\bar c_jx'^j$ in $T/\pi T$ with $\bar c_0=1$, i.e.\ $G\notin\pi T$. Hence $G$ is a unit of $T$ and $v_Z(G)=0$.

Thus $(G)=1\cdot Y$ and $Y$ is principal. This proves the Claim and (d).\hfill$\square$

\exercise{6.4}{Let $k$ be a field of characteristic $\ne2$. Let $f\in k[x_1,\dots,x_n]$ be a \emph{square-free} nonconstant polynomial, i.e.\ in the unique factorization of $f$ into irreducible polynomials, there are no repeated factors. Let $A=k[x_1,\dots,x_n,z]/(z^2-f)$. Show that $A$ is an integrally closed ring. [\emph{Hint}: The quotient field $K$ of $A$ is just $k(x_1,\dots,x_n)[z]/(z^2-f)$. It is a Galois extension of $k(x_1,\dots,x_n)$ with Galois group $\mathbf Z/2\mathbf Z$ generated by $z\mapsto-z$. If $\alpha=g+hz\in K$, where $g,h\in k(x_1,\dots,x_n)$, then the minimal polynomial of $\alpha$ is $X^2-2gX+(g^2-h^2f)$. Now show that $\alpha$ is integral over $k[x_1,\dots,x_n]$ if and only if $g,h\in k[x_1,\dots,x_n]$. Conclude that $A$ is the integral closure of $k[x_1,\dots,x_n]$ in $K$.]}
\pf
Put $R=k[x_1,\dots,x_n]$ (a UFD, hence integrally closed) and $F=\operatorname{Frac}R=k(x_1,\dots,x_n)$.

\emph{$A$ is a domain, $A=R\oplus Rz$.} Since $f$ is nonconstant and square-free, it is not a square in $F$ (some irreducible factor occurs to the first power, and valuations of squares are even); so $z^2-f$ is irreducible over $F$, hence (Gauss) irreducible in $R[z]$ (it is monic). Thus $A=R[z]/(z^2-f)$ is a domain, free over $R$ with basis $1,z$, so $A\subset K:=F[z]/(z^2-f)$ and $K=\operatorname{Frac}A$, $[K:F]=2$. As $\operatorname{char}k\ne2$, $\sigma:z\mapsto-z$ is a nontrivial $F$-automorphism of $K$, so $K/F$ is Galois with group $\{1,\sigma\}$.

\emph{Integrality criterion.} Let $\alpha=g+hz\in K$ ($g,h\in F$). Then $\sigma\alpha=g-hz$, $\alpha+\sigma\alpha=2g$, $\alpha\sigma\alpha=g^2-h^2f$, and $\alpha$ is a root of $X^2-2gX+(g^2-h^2f)\in F[X]$. Since $A$ is integral over $R$ ($z^2=f$), an element of $K$ is integral over $A$ iff it is integral over $R$. If $\alpha$ is integral over $R$ then so is $\sigma\alpha$ (apply $\sigma$ to a monic equation with coefficients in $R\subset F$), hence so are $2g=\alpha+\sigma\alpha$ and $g^2-h^2f=\alpha\sigma\alpha$; these lie in $F$, and $R$ is integrally closed, so $2g\in R$ and $g^2-h^2f\in R$. Conversely if $2g,\,g^2-h^2f\in R$ then $\alpha$ is a root of a monic polynomial over $R$.

Let $\alpha$ be integral over $A$. Then $2g\in R$, so $g\in R$ (2 is a unit), hence $h^2f=g^2-(g^2-h^2f)\in R$. Write $h=u/v$ with $u,v\in R$ coprime. Then $v^2\mid u^2f$, hence $v^2\mid f$. If $v$ had an irreducible factor $q$, then $q^2\mid f$, contradicting square-freeness. So $v$ is a unit and $h\in R$, i.e.\ $\alpha=g+hz\in A$. Hence $A$ is integrally closed in $K=\operatorname{Frac}A$, and, being integral over $R$ with the same fraction field extension, it is the integral closure of $R$ in $K$ (an element of $K$ integral over $R$ is integral over $A$, hence in $A$).\hfill$\square$

\exercise{6.5}{\emph{Quadric Hypersurfaces.} Let $\operatorname{char}k\neq2$, and let $X$ be the affine quadric hypersurface $\operatorname{Spec}k[x_0,\dots,x_n]/(x_0^2+x_1^2+\dots+x_r^2)$ (cf.\ I, Ex.~5.12). (Starred in the book.)
\begin{enumerate}[label=(\alph*),leftmargin=*]
\item Show that $X$ is normal if $r\ge2$ (use Ex.~6.4).
\item Show by a suitable linear change of coordinates that the equation of $X$ could be written as $x_0x_1=x_2^2+\dots+x_r^2$. Now imitate the method of (6.5.2) to show that: (1) If $r=2$, then $\operatorname{Cl}X\cong\mathbf Z/2\mathbf Z$; (2) If $r=3$, then $\operatorname{Cl}X\cong\mathbf Z$ (use (6.6.1) and Ex.~6.3 above); (3) If $r\ge4$ then $\operatorname{Cl}X=0$.
\item Now let $Q$ be the projective quadric hypersurface in $\mathbf P^n$ defined by the same equation. Show that: (1) If $r=2$, $\operatorname{Cl}Q\cong\mathbf Z$, and the class of a hyperplane section $Q.H$ is twice the generator; (2) If $r=3$, $\operatorname{Cl}Q\cong\mathbf Z\oplus\mathbf Z$; (3) If $r\ge4$, $\operatorname{Cl}Q\cong\mathbf Z$, generated by $Q.H$.
\item Prove Klein's theorem, which says that if $r\ge4$, and if $Y$ is an irreducible subvariety of codimension $1$ on $Q$, then there is an irreducible hypersurface $V\subseteq\mathbf P^n$ such that $V\cap Q=Y$, with multiplicity one. In other words, $Y$ is a complete intersection. (First show that for $r\ge4$, the homogeneous coordinate ring $S(Q)=k[x_0,\dots,x_n]/(x_0^2+\dots+x_r^2)$ is a UFD.)
\end{enumerate}}
\pf
\emph{Reading of the statement.} Here $0\le r\le n$. Part (a) holds over any field of characteristic $\ne2$. Parts (b)--(d) need $\sqrt{-1}\in k$ (for the coordinate change) and, for $r=3$, the fact $Q_3\cong\mathbf P^1\times\mathbf P^1$ with its classes over an algebraically closed field (6.6.1), and the geometric Ex.~6.2, 6.3 are for $k$ algebraically closed; so from (b) on I assume, as in the neighbouring exercises, that $k$ is algebraically closed of characteristic $\ne2$. (For $k=\mathbf R$ and $r=2$, for instance, $x_0^2+x_1^2+x_2^2$ cannot be brought to the form $x_0x_1-x_2^2$.)

\textbf{(a)} Let $r\ge2$ and $q=x_1^2+\dots+x_r^2\in k[x_1,\dots,x_n]$. Then $x_0^2+q=0$ defines $A=k[x_1,\dots,x_n,z]/(z^2-f)$ with $z=x_0$ and $f=-q$. By Ex.~6.4 it suffices to see that $f$ (equivalently $q$) is nonconstant and square-free. It is nonconstant. Suppose an irreducible $a$ with $a^2\mid q$, $q=a^2b$. Then $\partial q/\partial x_i=2x_i=a(2\partial_ia\cdot b+a\partial_ib)$ for $i=1,\dots,r$, so $a\mid x_1$ and $a\mid x_2$ (here $r\ge2$ and $2$ is a unit); the irreducible $a$ would then be associate to both $x_1$ and $x_2$, absurd. So $q$ is square-free, $A$ is an integrally closed domain, and $X$ is normal. In particular $X$ satisfies $(*)$.

\textbf{(b)} Since $\sqrt{-1}=:i\in k$, $x_0^2+x_1^2=(x_0+ix_1)(x_0-ix_1)$ and $x_j^2=-(ix_j)^2$. The linear change $y_0=x_0+ix_1$, $y_1=x_0-ix_1$, $y_j=ix_j$ ($2\le j\le r$), $y_j=x_j$ ($j>r$) is invertible (determinant $-2i\cdot i^{\,r-1}\ne0$, as $\operatorname{char}k\ne2$), and $x_0^2+\dots+x_r^2=y_0y_1-(y_2^2+\dots+y_r^2)$ because $x_j^2=-y_j^2$ for $2\le j\le r$. So after renaming, $X=\operatorname{Spec}A$ with $A=k[x_0,\dots,x_n]/(x_0x_1-q)$, $q=x_2^2+\dots+x_r^2$. Let $m=(x_0,\dots,x_n)$. Since $x_0x_1-q\in m^2$, $m/m^2$ has basis $\bar x_0,\dots,\bar x_n$. Observe also
\[
A_{x_1}=k[x_0,\dots,x_n,x_1^{-1}]/(x_0x_1-q)\cong k[x_1,x_1^{-1},x_2,\dots,x_n]\qquad(x_0=q/x_1).\qquad(1)
\]
This is a Laurent polynomial ring over a polynomial ring, hence a UFD, so $\operatorname{Cl}\operatorname{Spec}A_{x_1}=0$ by (6.2). Note $X$ is normal by (a) (the change of coordinates is an isomorphism), so (6.2), (6.5), (6.6) apply.

(1) $r=2$. $X:x_0x_1=x_2^2$ in $\mathbf A^{n+1}$. Let $Y=V(x_1,x_2)$, $\mathfrak p=(x_1,x_2)$; $A/\mathfrak p=k[x_0,x_3,\dots,x_n]$ is a domain of dimension $n-1=\dim A-1$, so $Y$ is a prime divisor. Set-theoretically $V(x_1)=Y$ ($x_1=0\Rightarrow x_2^2=0$). Since $x_0\notin\mathfrak p$, $x_0$ is a unit in $A_{\mathfrak p}$ and $x_1=x_2^2/x_0$, so $\mathfrak pA_{\mathfrak p}=(x_2)$ and $v_Y(x_1)=2$. So $(x_1)=2Y$, i.e.\ $2[Y]=0$. By (6.5)(c) applied to the irreducible $Y$ and $X-Y=D(x_1)=\operatorname{Spec}A_{x_1}$ we have an exact sequence $\mathbf Z\to\operatorname{Cl}X\to\operatorname{Cl}(X-Y)=0$, so $\operatorname{Cl}X$ is generated by $[Y]$, of order $\le2$. It remains to see $[Y]\ne0$. If $Y=(g)$, $g\in K$, then $g\in A_{\mathfrak p'}$ for every height-one prime $\mathfrak p'$, so $g\in A$ by (6.3A), and every $a\in\mathfrak p$ has $a/g\in A_{\mathfrak p'}$ for all $\mathfrak p'$, so $a\in gA$; also $g\in\mathfrak pA_{\mathfrak p}\cap A=\mathfrak p$. Hence $\mathfrak p=gA$ is principal (this is the argument of (6.2)). But then the image of $\mathfrak p$ in $m/m^2$ would be spanned by the class of $g$, of dimension $\le1$, whereas it contains the linearly independent classes $\bar x_1,\bar x_2$. Contradiction. Hence $\operatorname{Cl}X=\mathbf Z/2\mathbf Z$, generated by $[Y]$.

(2) $r=3$. Put $u=x_2+ix_3$, $v=x_2-ix_3$; then $q=uv$ and $X:x_0x_1=uv$, i.e.\ $X\cong X_3\times\mathbf A^{n-3}$ where $X_3=\operatorname{Spec}k[x_0,x_1,u,v]/(x_0x_1-uv)$. By (6.6) (iterated; $X_3$ satisfies $(*)$ being normal by (a)), $\operatorname{Cl}X\cong\operatorname{Cl}X_3$. Now $X_3=C(Q_3)$ is the affine cone over the quadric surface $Q_3=V_+(x_0x_1-uv)\subset\mathbf P^3$, which is nonsingular, of dimension $2$. By (6.6.1) $Q_3\cong\mathbf P^1\times\mathbf P^1$ and $\operatorname{Cl}Q_3=\mathbf Z\oplus\mathbf Z$, and by (6.6.2) a hyperplane section $Q_3.H$ is of type $(1,1)$. Ex.~6.3(b) gives $\operatorname{Cl}X_3\cong(\mathbf Z\oplus\mathbf Z)/\mathbf Z(1,1)\cong\mathbf Z$ (via $(a,b)\mapsto a-b$). So $\operatorname{Cl}X\cong\mathbf Z$.

(3) $r\ge4$. Now $\mathfrak p=(x_1)\subset A$ has quotient $A/(x_1)=k[x_0,x_2,\dots,x_n]/(q)$ with $q=x_2^2+\dots+x_r^2$ a quadratic form of rank $r-1\ge3$, hence irreducible (a reducible quadratic form is a product of two linear forms, so has rank $\le2$). So $\mathfrak p=(x_1)$ is a principal prime, $Y=V(x_1)$ is a prime divisor, and $(x_1)=Y$ (as $\mathfrak pA_{\mathfrak p}=x_1A_{\mathfrak p}$, $v_Y(x_1)=1$, and $x_1$ is a unit elsewhere), so $[Y]=0$. By (6.5)(c), $\mathbf Z\xrightarrow{1\mapsto[Y]}\operatorname{Cl}X\to\operatorname{Cl}(X-Y)\to0$ and $\operatorname{Cl}(X-Y)=\operatorname{Cl}\operatorname{Spec}A_{x_1}=0$ by (1); hence $\operatorname{Cl}X=0$.

\textbf{(c)} For $r\ge2$, $Q\subset\mathbf P^n$ is a projective variety of dimension $n-1\ge1$ (irreducible by (a)), normal (its homogeneous coordinate ring is integrally closed by (a), so the local rings of $Q$, which are localizations of degree-$0$ parts, are integrally closed), hence nonsingular in codimension one; and $X=\operatorname{Spec}S(Q)$ is the affine cone $C(Q)$. So Ex.~6.3(b) gives an exact sequence
\[
0\to\mathbf Z\xrightarrow{1\mapsto Q.H}\operatorname{Cl}Q\to\operatorname{Cl}X\to0.\qquad(2)
\]
(3) $r\ge4$: $\operatorname{Cl}X=0$ by (b)(3), so $\mathbf Z\cong\operatorname{Cl}Q$, generated by $Q.H$.

(2) $r=3$: $\operatorname{Cl}X\cong\mathbf Z$ by (b)(2), which is free; so (2) splits and $\operatorname{Cl}Q\cong\mathbf Z\oplus\mathbf Z$.

(1) $r=2$: $\bar Y:=V_+(x_1,x_2)\subset Q$ is a linear space $\mathbf P^{n-2}$, an irreducible subvariety of codimension $1$ in $Q$ and of degree $1$; the cone over it is $Y$, and the map $\operatorname{Cl}Q\to\operatorname{Cl}X$ of Ex.~6.3(b) sends $[\bar Y]\mapsto[Y]$, a generator of $\operatorname{Cl}X$ by (b)(1). By exactness of (2), $\operatorname{Cl}Q$ is generated by $[\bar Y]$ and $[Q.H]$. Take $H=V_+(x_1)\not\supseteq Q$. In the chart $x_0=1$, $Q$ is $\operatorname{Spec}k[x_2',x_3',\dots,x_n']\cong\mathbf A^{n-1}$ ($x_1'=x_2'^2$), $\bar Y$ is $\{x_2'=0\}$ and $x_1'$ has valuation $2$ along $\bar Y$; moreover $V(x_1)\cap Q=\bar Y$ as a set. So $Q.H=2\bar Y$ and $\operatorname{Cl}Q$ is cyclic, generated by $[\bar Y]$. By Ex.~6.2(d) there is a degree homomorphism $\operatorname{Cl}Q\to\mathbf Z$ with $\deg\bar Y=1$; so $\operatorname{Cl}Q\cong\mathbf Z$ generated by $[\bar Y]$, and $[Q.H]=2[\bar Y]$ is twice the generator. (Consistently, $\deg(Q.H)=\deg Q=2$.)

\textbf{(d)} Let $r\ge4$. By (b)(3) $\operatorname{Cl}X=0$ and by (a) $X$ is normal, so $S:=S(Q)=k[x_0,\dots,x_n]/(x_0^2+\dots+x_r^2)$ is a UFD by (6.2). Let $Y\subset Q$ be an irreducible closed subvariety of codimension one. Its ideal $\bar{\mathfrak p}\subset S$ is a homogeneous prime with $\dim S/\bar{\mathfrak p}=\dim Y+1=\dim S-1$, hence of height one. In a UFD a height-one prime is principal: pick a nonzero homogeneous $a\in\bar{\mathfrak p}$ (possible, $\bar{\mathfrak p}$ being homogeneous and nonzero); some prime factor $\varpi$ of $a$ lies in $\bar{\mathfrak p}$, and $(\varpi)\subseteq\bar{\mathfrak p}$ forces $\bar{\mathfrak p}=(\varpi)$ by height. The factors of a homogeneous element of a graded domain are homogeneous (compare lowest and highest degree components of a factorization), so $\varpi=:\bar g$ is homogeneous, of degree $d\ge1$.

Lift $\bar g$ to a homogeneous $g\in k[x_0,\dots,x_n]$ of degree $d$. \emph{$g$ is irreducible:} if $g=g_1g_2$ with homogeneous $g_i$ of positive degrees $d_i<d$, then $\bar g=\bar g_1\bar g_2$ and, $\bar g$ being prime, $\bar g\mid\bar g_1$ say; $\bar g_1=\bar g\bar w$ with $\bar w$ homogeneous of degree $d_1-d<0$ forces $\bar g_1=0$, hence $\bar g=0$, absurd. So $V=V_+(g)$ is an irreducible hypersurface; $Q\not\subseteq V$ since $\bar g\ne0$. The ideal $(f_Q,g)$, $f_Q=x_0^2+\dots+x_r^2$, has quotient $S/(\bar g)=S/\bar{\mathfrak p}$, a domain; so it is prime and $V\cap Q=Y$ as sets, with the single component $Y$. Finally the multiplicity: by Ex.~6.2(c), the coefficient of $Y$ in $Q.V$ is $i(Q,V;Y)=\operatorname{length}_{S'_{\mathfrak P}}(S'/(f_Q,g))_{\mathfrak P}$ with $S'=k[x_0,\dots,x_n]$, $\mathfrak P$ the preimage of $\bar{\mathfrak p}$; this length is that of $S_{\bar{\mathfrak p}}/\bar gS_{\bar{\mathfrak p}}=S_{\bar{\mathfrak p}}/\bar{\mathfrak p}S_{\bar{\mathfrak p}}$, a field, so it is $1$. Therefore $Q.V=Y$: $V\cap Q=Y$ with multiplicity one.\hfill$\square$

\exercise{6.6}{Let $X$ be the nonsingular plane cubic curve $y^2z=x^3-xz^2$ of (6.10.2).
\begin{enumerate}[label=(\alph*),leftmargin=*]
\item Show that three points $P,Q,R$ of $X$ are collinear if and only if $P+Q+R=0$ in the group law on $X$. (Note that the point $P_0=(0,1,0)$ is the zero element in the group structure on $X$.)
\item A point $P\in X$ has order $2$ in the group law on $X$ if and only if the tangent line at $P$ passes through $P_0$.
\item A point $P\in X$ has order $3$ in the group law on $X$ if and only if $P$ is an inflection point. (An \emph{inflection point} of a plane curve is a nonsingular point $P$ of the curve, whose tangent line (I, Ex.~7.3) has intersection multiplicity $\ge3$ with the curve at $P$.)
\item Let $k=\mathbf C$. Show that the points of $X$ with coordinates in $\mathbf Q$ form a subgroup of the group of all points. Can you determine the structure of this subgroup explicitly?
\end{enumerate}}
\pf
\emph{Setting.} $k$ is algebraically closed with $\operatorname{char}k\ne2$, as in (6.10.2) (for (d), $k=\mathbf C$). $X\subset\mathbf P^2$ is a nonsingular cubic with $P_0=(0,1,0)$ and group law defined as in (6.10.2): $P\mapsto[P-P_0]$ is a bijection from the closed points to $\ker(\deg:\operatorname{Cl}X\to\mathbf Z)$, and the sum $S=P+Q$ in the group is the point with $[S-P_0]=[P-P_0]+[Q-P_0]$ in $\operatorname{Cl}X$, i.e.\ $S\sim P+Q-P_0$. The injectivity of $P\mapsto[P-P_0]$ was shown in (6.10.2): $X$ is not rational, so two distinct points are not linearly equivalent (6.10.1). For a line $L\subset\mathbf P^2$ ($L\not\subseteq X$ since $X$ is irreducible), $L.X$ (Ex.~6.2) is an effective divisor of degree $\deg L\cdot\deg X=3$, whose coefficient at $P$ is the intersection multiplicity $i(X,L;P)$ (Ex.~6.2(c)). ``$P,Q,R$ are collinear'' means: there is a line $L$ with $L.X=P+Q+R$ (so points are counted with multiplicity; e.g.\ for $P=Q$ the line is tangent). For a nonsingular point $P$ of $X$ and a line $L\ni P$, $i(X,L;P)\ge2$ iff $L$ is the tangent line $T_P$ (the order of vanishing at $P$ of the equation of $X$ restricted to $L$ is $\ge2$ iff the gradient of the equation kills the direction of $L$). All lines are linearly equivalent: $L.X\sim L_0.X$ with $L_0=\{z=0\}$, since $(\ell/z).X=L.X-L_0.X$ is principal by Ex.~6.2(b); and $L_0.X=3P_0$ (put $z=0$ in the equation: $x^3=0$). Hence for every line $L$,
\[
L.X\sim3P_0.\qquad(3)
\]
(This is the relation $P+Q+R\sim3P_0$ of (6.10.2).) Finally, given $P,Q\in X$, let $L$ be the line $PQ$ if $P\ne Q$ and $T_P$ if $P=Q$; then $L.X-P-Q\ge0$, say $L.X=P+Q+R'$ (for $P=Q$, $i(X,T_P;P)\ge2$), which defines the ``third point'' $R'$.

\textbf{(a)} If $P,Q,R$ are collinear, $P+Q+R=L.X\sim3P_0$ by (3), i.e.\ $(P-P_0)+(Q-P_0)+(R-P_0)\sim0$, which says $P+Q+R=0$ in the group. Conversely let $P+Q+R=0$ in the group, i.e.\ $(P-P_0)+(Q-P_0)+(R-P_0)\sim0$, so $P+Q+R\sim3P_0$. Let $L$ be the line $PQ$ (tangent line if $P=Q$) and $L.X=P+Q+R'$. By (3), $P+Q+R'\sim3P_0$, so $R\sim R'$, hence $R=R'$ by injectivity of $P\mapsto[P-P_0]$ (two points linearly equivalent are equal). So $L.X=P+Q+R$.

\textbf{(b)} Reading: the statement is about $P\ne P_0$ ($P_0$ is the zero element, of order $1$, though its tangent $z=0$ passes through $P_0$). $P$ has order $2$ iff $P\neq P_0$ and $2P=0$, iff (since $P_0$ is the zero) $P+P+P_0=0$, iff by (a) there is a line $L$ with $L.X=2P+P_0$. If such $L$ exists, $L\ni P,P_0$ and $i(X,L;P)\ge2$, so $L=T_P$ and $T_P\ni P_0$. Conversely if $P\ne P_0$ and $P_0\in T_P$ then $T_P.X\ge2P+P_0$ (the multiplicities at $P$ and at $P_0\neq P$), and has degree $3$, so $T_P.X=2P+P_0$.

\textbf{(c)} Again $P_0$ itself satisfies the equivalence in the form ``$3P=0$ iff $P$ is an inflection point'' (so ``order $3$'' means $P\ne P_0$). Now $3P=0$ iff $P+P+P=0$ iff (a) there is a line $L$ with $L.X=3P$. If so, $i(X,L;P)=3\ge2$ so $L=T_P$ and $i(X,T_P;P)=3$: $P$ is an inflection point. Conversely if $i(X,T_P;P)\ge3$ then $T_P.X=3P$ (degree $3$).

\textbf{(d)} \emph{Subgroup.} Let $X(\mathbf Q)$ be the set of points with coordinates in $\mathbf Q$ (the equation has rational coefficients). $P_0\in X(\mathbf Q)$. Let $P,Q\in X(\mathbf Q)$. The line $L=PQ$ has rational coefficients; if $P=Q$, $T_P$ is given by the gradient of the defining polynomial at $P$, again rational. Choose a rational point $Q'\in L$ with $Q'\ne P$ (take $Q'=Q$ if $Q\ne P$; if $P=Q$, $L=T_P$ is a line with rational coefficients and has other rational points). Parametrizing $L$ by $(s,t)\mapsto sP+tQ'$, the cubic form $F|_L$ in $(s,t)$ has rational coefficients and vanishes at the rational root $(1:0)$ corresponding to $P$, to order $\ge2$ if $P=Q$ and $L=T_P$, and also at $(0:1)$ if $Q'=Q$; dividing out the rational linear factors corresponding to the known roots $P,Q$ (counted with multiplicity), the remaining linear factor is rational, so the third point $R'$ of (a) is rational. Hence the third point of the line through two rational points is rational. By (a), $-P$ is the third point of the line $PP_0$ (since $P+P_0+(-P)=0$ means $P,P_0,-P$ collinear), so $-P\in X(\mathbf Q)$; and $P+Q=-R'$ with $R'$ the third point of $PQ$, so $P+Q\in X(\mathbf Q)$. Hence $X(\mathbf Q)$ is a subgroup.

\emph{Structure.} Claim: $X(\mathbf Q)=\{P_0,\ (0,0,1),\ (1,0,1),\ (-1,0,1)\}\cong\mathbf Z/2\times\mathbf Z/2$.

The four points are on $X$ ($y=0$ and $x\in\{0,\pm1\}$ gives $0=x^3-x$), rational. Each $(e,0,1)$ has tangent line $x=ez$ (the equation is $F=y^2z-x^3+xz^2$, $F_y=2yz=0$ at such a point, so the tangent $F_xx+F_yy+F_zz=0$ is a line through $(0,1,0)=P_0$), so by (b) each of the three has order $2$. A group of order $4$ in which three elements have order $2$ is $\mathbf Z/2\times\mathbf Z/2$. It remains to show there are no other rational points, i.e.\ that $y^2=x^3-x$ has no rational solution with $y\ne0$ (the only point at infinity is $P_0$: $z=0$ forces $x=0$).

Suppose $(x,y)\in\mathbf Q^2$, $y\ne0$, $y^2=x^3-x=x(x-1)(x+1)>0$. Then $x\in(-1,0)\cup(1,\infty)$. If $x\in(-1,0)$, then $(x',y')=(-1/x,\ y/x^2)$ is a rational point with $y'\ne0$ and $x'>1$ (check: $y'^2=(x^3-x)/x^4=(x^2-1)/x^3=x'^3-x'$). So it suffices to exclude $x>1$. Write $x=a/d^2$ with $a,d\in\mathbf Z$, $d\ge1$, $\gcd(a,d)=1$ (possible: if $x=a/b_0$ in lowest terms, $b_0\ge1$, then $x^3-x=a(a^2-b_0^2)/b_0^3$ is again in lowest terms, because $\gcd(a,b_0)=1$; it equals the square $y^2$ of a rational number, whose reduced denominator is a perfect square; so $b_0^3$, hence $b_0$, is a perfect square $d^2$). Then $(yd^3)^2=a(a^2-d^4)=a(a-d^2)(a+d^2)\in\mathbf Z$, so $b:=yd^3\in\mathbf Z$ and $b^2=a(a-d^2)(a+d^2)$, with $a>d^2>0$ since $x>1$. Put $B=a-d^2$, $C=a+d^2$. Then $\gcd(a,B)=\gcd(a,C)=\gcd(a,d^2)=1$ and $\gcd(B,C)\mid\gcd(2a,2d^2)=2$.
\begin{itemize}
\item If $a$ is even, or $d$ is even ($a$ odd): $B,C$ are odd, so $a,B,C$ are pairwise coprime positive integers with product a square; hence $a=s^2$, $B=t^2$, $C=u^2$. If $a$ even then $d$ is odd (coprime); in both cases $s^2-d^2=t^2$, $s^2+d^2=u^2$, whence $s^4-d^4=(tu)^2$ with $s,d,tu$ positive integers.
\item If $a$ and $d$ are both odd: $B,C$ are even and $\gcd(B,C)=2$. As $a$ is coprime to $BC$ and $a\cdot BC=b^2$, $a=s^2$ and $BC$ is a square. Put $B=2B'$, $C=2C'$, $\gcd(B',C')=1$, $B'C'=BC/4$ a square, so $B'=p^2$, $C'=q^2$ with $p,q$ positive integers. Then $s^2=a=B'+C'=p^2+q^2$ and $d^2=C'-B'=q^2-p^2$, so $q^4-p^4=(sd)^2$ with $q,p,sd$ positive integers.
\end{itemize}
Both cases contradict Fermat's theorem (by infinite descent) that $X^4-Y^4=Z^2$ has no solution in positive integers. Hence $X(\mathbf Q)$ consists of the four points listed. So $X(\mathbf Q)\cong\mathbf Z/2\times\mathbf Z/2$: the group of rational points is finite.\hfill$\square$

\exercise{6.7}{Let $X$ be the nodal cubic curve $y^2z=x^3+x^2z$ in $\mathbf P^2$. Imitate (6.11.4) and show that the group of Cartier divisors of degree $0$, $\operatorname{CaCl}^0X$, is naturally isomorphic to the multiplicative group $\mathbf G_m$. (Starred in the book.)}
\pf
\emph{Setting.} $k$ is algebraically closed, $\operatorname{char}k\ne2$ (so that $y^2=x^2(1+x)$ is really a node). $\operatorname{CaCl}X$ is the group of Cartier divisors modulo principal ones, $\operatorname{CaCl}^0X$ the kernel of the degree $\operatorname{CaCl}X\to\mathbf Z$ defined as in (6.11.4); I recall and verify that definition below. Let $Z=(0,0,1)$ and $P_0=(0,1,0)$.

\emph{Step 1: the normalization.} Define $\varphi:\mathbf P^1\to X$, $(T:S)\mapsto(S(T^2-S^2):T(T^2-S^2):S^3)$. In the coordinate $t=T/S$ it is $t\mapsto(t^2-1,\ t(t^2-1))$ on the affine part $z=1$ (the line $y=tx$ through the node meets the curve again at $x=t^2-1$), and $\varphi(\infty)=P_0$. It lands in $X$: $y^2z=T^2(T^2-S^2)^2S^3=x^3+x^2z$ (both equal $S^3(T^2-S^2)^2T^2$ after the substitution, using $(T^2-S^2)+S^2=T^2$). $\varphi^{-1}(Z)=\{t=\pm1\}$. The only singular point of $X$ is $Z$: with $F=y^2z-x^3-x^2z$, $F_x=-3x^2-2xz$, $F_y=2yz$, $F_z=y^2-x^2$; $F_y=0$ means $y=0$ or $z=0$; $z=0$ forces $x=0$ on $X$, and then $F_z=y^2\ne0$ at $P_0$; $y=0$ forces $x^2(x+z)=0$ and $F_z=-x^2=0$, so $x=0$, the point $Z$. The map $\varphi$ is a nonconstant morphism of complete curves, hence finite (6.8), and birational ($t=y/x$). Therefore $\varphi^{-1}(X-Z)\to X-Z$ is a finite birational morphism onto a nonsingular, hence normal, curve, hence an isomorphism. So $\varphi$ is the normalization, $X-Z\cong\mathbf P^1-\{1,-1\}$, and with the coordinate
\[u=\frac{t+1}{t-1}=\frac{y+x}{y-x}\]
(which sends $t=-1,\,1,\,\infty$ to $0,\,\infty,\,1$) we get an isomorphism $u:X-Z\to\mathbf G_m=\mathbf P^1-\{0,\infty\}$ with $u(P_0)=1$.

\emph{Step 2: the local ring at $Z$.} Let $A=k[x,y]/(y^2-x^3-x^2)$ (the affine part $z=1$). Via $\varphi$, $A\to k[t]$ is $x\mapsto t^2-1$, $y\mapsto t(t^2-1)$. Its image is $B=\{g\in k[t]:g(1)=g(-1)\}$: indeed with $w=t^2-1$ one has $k[t]=k[w]\oplus t\,k[w]$ and $B=k+(t^2-1)k[t]=k[w]\oplus tw\,k[w]$, which is generated as an algebra by $w$ and $v=tw$ (note $v^2=w^2(w+1)$), i.e.\ equals the image of $A$; and $A\to B$ is injective: the kernel of $k[x,y]\to k[t]$ is a prime of height one (the image has dimension one) containing the prime $(y^2-x^3-x^2)$ of height one (irreducible, as $x^2(x+1)$ is not a square), hence equal to it. The local ring is $\mathcal O_{X,Z}=B_{\mathfrak m}$, $\mathfrak m=\{g\in B:g(1)=g(-1)=0\}$. Claim:
\[
\mathcal O_{X,Z}=\{f\in k(t):f\text{ regular at }t=\pm1,\ f(1)=f(-1)\}.
\]
``$\subseteq$'' is clear. Conversely, take such $f=a/b$ with $a,b\in k[t]$, $b(1)b(-1)\neq0$ (the semilocal ring of $\mathbf P^1$ at $\{\pm1\}$). Choose a polynomial $\tilde b$ of degree $\le1$ with $\tilde b(1)=b(-1)$, $\tilde b(-1)=b(1)$ (both nonzero). Then $h=b\tilde b\in B$, $h(1)=h(-1)\ne0$, and $g=f\,h=a\tilde b\in k[t]$ has $g(1)=f(1)h(1)=f(-1)h(-1)=g(-1)$, so $g\in B$ and $f=g/h\in B_{\mathfrak m}$. Consequently $\mathcal O_{X,Z}^*=\{f\ \text{regular at }t=\pm1,\ f(1)=f(-1)\ne0\}$, i.e.\ in the coordinate $u$: $f$ regular at $u=0,\infty$ with $f(0)=f(\infty)\ne0$. (This is just the statement that the node is obtained from $\mathbf P^1$ by gluing $t=1$ and $t=-1$.)

\emph{Step 3: Cartier divisors and degree.} Let $\mathcal C$ be the group of Cartier divisors of $X$ which are trivial in a neighbourhood of $Z$. Every Cartier divisor $D$ is linearly equivalent to one in $\mathcal C$: if $f_Z\in K^*$ is a local equation of $D$ near $Z$, then $D-(f_Z)$ is given by $1$ near $Z$. A Cartier divisor in $\mathcal C$ is the same as a Cartier divisor on the nonsingular curve $X-Z$ with finite support (its support is a finite set $T\subset X-Z$ and it is trivial on $X-T\ni Z$); by (6.11) $\mathcal C\cong\operatorname{Div}(X-Z)=\bigoplus_{P\in X-Z}\mathbf Z$. A principal Cartier divisor $(f)$, $f\in K^*$, lies in $\mathcal C$ iff $f$ is a unit in a neighbourhood of $Z$, iff $f\in\mathcal O_{X,Z}^*$. Hence
\[
\operatorname{CaCl}X=\operatorname{Div}(X-Z)\big/\{(f)|_{X-Z}:f\in\mathcal O_{X,Z}^*\}.
\]
The degree of $D=\sum n_iP_i\in\operatorname{Div}(X-Z)$ is $\sum n_i$; it vanishes on the principal divisors in this quotient (as in (6.11.4)): for $f\in\mathcal O_{X,Z}^*$ viewed as a rational function on $\mathbf P^1$, the divisor $(f)_{\mathbf P^1}$ has degree $0$ (6.10 for $\mathbf P^1$) and no support at $t=\pm1$, so equals $(f)|_{X-Z}$. This is the degree of (6.11.4). Thus $\operatorname{CaCl}^0X=\operatorname{Div}^0(X-Z)/\{(f)|_{X-Z}:f\in\mathcal O_{X,Z}^*\}$.

\emph{Step 4: the isomorphism.} Transport by $u$: $\operatorname{Div}(X-Z)=\operatorname{Div}(\mathbf G_m)$, points $a\in k^*$ (note $P_0\leftrightarrow a=1$). Define
\[
\Psi:\operatorname{Div}^0(\mathbf G_m)\to\mathbf G_m,\qquad\sum n_i[a_i]\mapsto\prod a_i^{n_i}\quad(\textstyle\sum n_i=0).
\]
It is a surjective homomorphism ($[a]-[1]\mapsto a$). We show $\ker\Psi=\{(f)|_{\mathbf G_m}:f\in\mathcal O_{X,Z}^*\}$. 

Let $f\in\mathcal O_{X,Z}^*$. Write $f=c\,u^m\prod(u-a_i)^{n_i}$ with $c\in k^*$, $a_i\in k^*$ distinct. Then $v_0(f)=m$ and $v_\infty(f)=-(m+\sum n_i)$; $f$ being a unit at $u=0$ and $u=\infty$ forces $m=0$ and $\sum n_i=0$. So $(f)|_{\mathbf G_m}=\sum n_i[a_i]\in\operatorname{Div}^0$, $f(\infty)=c$ and $f(0)=c\prod(-a_i)^{n_i}=c\prod a_i^{n_i}$ (as $\sum n_i=0$). The condition $f(0)=f(\infty)$ is $\prod a_i^{n_i}=1$, i.e.\ $\Psi((f))=1$. Conversely if $\Psi(D)=1$, $D=\sum n_i[a_i]$, then $f=\prod(u-a_i)^{n_i}$ satisfies $f(\infty)=1=f(0)$, so $f\in\mathcal O_{X,Z}^*$ with $(f)|_{\mathbf G_m}=D$. Hence $\Psi$ induces an isomorphism
\[
\operatorname{CaCl}^0X\xrightarrow{\ \cong\ }\mathbf G_m,\qquad[P-P_0]\mapsto u(P)=\frac{y+x}{y-x}(P)\quad(P\ne Z).
\]
It is natural: it is the identification of the group of nonsingular points (the group variety $X-Z\cong\mathbf G_m$, with identity $P_0$) with $\operatorname{CaCl}^0X$, exactly as in (6.10.2).\hfill$\square$

\exercise{6.8}{
\begin{enumerate}[label=(\alph*),leftmargin=*]
\item Let $f:X\to Y$ be a morphism of schemes. Show that $\mathcal L\mapsto f^*\mathcal L$ induces a homomorphism of Picard groups, $f^*:\operatorname{Pic}Y\to\operatorname{Pic}X$.
\item If $f$ is a finite morphism of nonsingular curves, show that this homomorphism corresponds to the homomorphism $f^*:\operatorname{Cl}Y\to\operatorname{Cl}X$ defined in the text, via the isomorphisms of (6.16).
\item If $X$ is a locally factorial integral closed subscheme of $\mathbf P^n_k$, and if $f:X\to\mathbf P^n$ is the inclusion map, then $f^*$ on $\operatorname{Pic}$ agrees with the homomorphism on divisor class groups defined in Ex.~6.2, via the isomorphisms of (6.16).
\end{enumerate}}
\pf
\textbf{(a)} If $\mathcal L$ is invertible on $Y$, choose an open cover $\{U_i\}$ with $\mathcal L|_{U_i}\cong\mathcal O_{U_i}$. Then $(f^*\mathcal L)|_{f^{-1}U_i}=(f|_{f^{-1}U_i})^*(\mathcal L|_{U_i})\cong f^*\mathcal O_{U_i}=\mathcal O_{f^{-1}U_i}$, so $f^*\mathcal L$ is invertible. Isomorphic sheaves have isomorphic pull-backs, $f^*\mathcal O_Y=\mathcal O_X$, and $f^*(\mathcal L\otimes\mathcal M)\cong f^*\mathcal L\otimes f^*\mathcal M$ (Prop.~II.5.2(e)). Hence $[\mathcal L]\mapsto[f^*\mathcal L]$ is a well-defined homomorphism $\operatorname{Pic}Y\to\operatorname{Pic}X$ (unit and inverse are preserved).

\textbf{(b)} Let $f:X\to Y$ be a finite morphism of nonsingular curves; $f$ is finite, hence has finite fibres, so it is not constant; thus $f(X)$ is dense in $Y$, $f$ is dominant and $K(Y)\subseteq K(X)$ via $f$. Both $X,Y$ are integral, separated, noetherian and locally factorial (regular local rings are UFD's, (6.11A)); by (6.11), (6.15), (6.16) we have isomorphisms $\operatorname{Cl}\cong\operatorname{CaCl}\cong\operatorname{Pic}$, $D\mapsto\mathcal L(D)$. So we must show $f^*\mathcal L(D)\cong\mathcal L(f^*D)$ for every divisor $D$ on $Y$, where $f^*D$ is the divisor of the text.

Let $D$ be the Cartier divisor $\{(U_i,g_i)\}$ on $Y$, $g_i\in K(Y)^*$. Define the Cartier divisor $f^\#D=\{(f^{-1}U_i,g_i)\}$ on $X$ ($g_i\in K(Y)^*\subseteq K(X)^*$; $g_i/g_j\in\mathcal O^*(U_i\cap U_j)$ pulls back to a unit). The sheaf $\mathcal L(D)\subseteq\mathcal K_Y$ is generated by $g_i^{-1}$ on $U_i$. The natural map $f^*\mathcal L(D)=f^{-1}\mathcal L(D)\otimes_{f^{-1}\mathcal O_Y}\mathcal O_X\to\mathcal K_X$, $s\otimes a\mapsto sa$ is, on $f^{-1}U_i$, a map from the free module generated by $g_i^{-1}\otimes1$ to $\mathcal K_X$, sending the generator to the nonzero element $g_i^{-1}$; since $\mathcal O_X$ is a sheaf of domains this is injective, with image $\mathcal O_{f^{-1}U_i}g_i^{-1}=\mathcal L(f^\#D)|_{f^{-1}U_i}$. Hence $f^*\mathcal L(D)\cong\mathcal L(f^\#D)$.

It remains to see that the Weil divisor of $f^\#D$ is the $f^*D$ of the text. Both sides are additive in $D$, so it suffices to take $D=Q$ a closed point. Choose an open neighbourhood $U$ of $Q$ in $Y$ and $t\in K(Y)$ a local parameter at $Q$ with $t$ a unit on $U-\{Q\}$ (shrink $U$), and use the cover $\{U,\ Y-\{Q\}\}$ with local equations $t$ and $1$. For $P\in X$ with $f(P)=Q$, the coefficient of $P$ in the Weil divisor of $f^\#Q$ is $v_P(t)$, since $P\in f^{-1}U$ where the local equation is $t$ (coefficient $v_P$ of the local equation, by the definition in the proof of (6.11)); for $f(P)\ne Q$ it is $v_P(1)=0$. This is $f^*Q=\sum_{f(P)=Q}v_P(t)P$ by definition. Hence the diagram $\operatorname{Cl}Y\to\operatorname{Pic}Y$, $\operatorname{Cl}X\to\operatorname{Pic}X$ with vertical maps $f^*$ commutes.

\textbf{(c)} Here $k$ is algebraically closed, $X$ is a closed integral subscheme of $\mathbf P^n$ and locally factorial; hence normal, so regular in codimension one, and $X$ is a closed subvariety of $\mathbf P^n$ nonsingular in codimension one: the homomorphism $\varphi:\operatorname{Cl}\mathbf P^n\to\operatorname{Cl}X$ of Ex.~6.2 is defined. The scheme $\mathbf P^n$ is locally factorial, so $\operatorname{Cl}\mathbf P^n\cong\operatorname{CaCl}\mathbf P^n\cong\operatorname{Pic}\mathbf P^n$ by (6.16) (and (6.17)). Let $D=\sum m_jV_j\in\mathcal G$ (Ex.~6.2): no $V_j=V_+(h_j)$ contains $X$; every class of $\operatorname{Cl}\mathbf P^n$ has such a representative (Ex.~6.2(b)). With $U_i=D_+(x_i)$, the Cartier divisor of $D$ is $\{(U_i,g_i)\}$, $g_i=\prod_j(h_j/x_i^{d_j})^{m_j}$, $d_j=\deg h_j$. Since no $h_j$ vanishes identically on $X$, restricting gives the Cartier divisor $f^\#D=\{(X\cap U_i,\bar g_i)\}$ on $X$ (with $\bar g_i\in K(X)^*$; when $X\cap U_i=\emptyset$ the $i$-th term is omitted). The sheaf $\mathcal L(D)$ is invertible with transition functions $u_{ij}=g_j/g_i\in\mathcal O^*(U_i\cap U_j)$ with respect to the generators $g_i^{-1}$, and $f^*\mathcal L(D)$ is the invertible sheaf with the restricted transition functions $\bar u_{ij}=\bar g_j/\bar g_i$ with respect to the generators $g_i^{-1}\otimes1$; these are also the transition functions of $\mathcal L(f^\#D)$ with respect to $\bar g_i^{-1}$. Hence $f^*\mathcal L(D)\cong\mathcal L(f^\#D)$ (isomorphic invertible sheaves, as they are trivialized on the same cover with the same cocycle). Finally, the Weil divisor of $f^\#D$ has, for a prime divisor $Y$ of $X$, coefficient $v_Y(\bar g_i)$ ($i$ with $\bar x_i\notin\mathfrak p_Y$), which by (1) of Ex.~6.2(a) and linearity is $\sum_jm_j\operatorname{coeff}_Y(V_j.X)=\operatorname{coeff}_Y(D.X)$. So the Weil divisor of $f^\#D$ is $D.X$, and $f^*\mathcal L(D)\cong\mathcal L(D.X)$: under $\operatorname{Cl}\cong\operatorname{Pic}$, $f^*$ corresponds to $\varphi$.\hfill$\square$

\exercise{6.9}{\emph{Singular Curves.} Here we give another method of calculating the Picard group of a singular curve. Let $X$ be a projective curve over $k$, let $\tilde X$ be its normalization, let $\pi:\tilde X\to X$ be the projection map (Ex.~3.8). For each point $P\in X$, let $\mathcal O_P$ be its local ring, and let $\tilde{\mathcal O}_P$ be the integral closure of $\mathcal O_P$. We use $*$ to denote the group of units in a ring. (Starred in the book.)
\begin{enumerate}[label=(\alph*),leftmargin=*]
\item Show there is an exact sequence
\[0\to\bigoplus_{P\in X}\tilde{\mathcal O}_P^{*}/\mathcal O_P^{*}\to\operatorname{Pic}X\xrightarrow{\ \pi^*\ }\operatorname{Pic}\tilde X\to0.\]
[\emph{Hint}: Represent $\operatorname{Pic}X$ and $\operatorname{Pic}\tilde X$ as the groups of Cartier divisors modulo principal divisors, and use the exact sequence of sheaves $0\to\pi_*\tilde{\mathcal O}_{\tilde X}^*/\mathcal O_X^*\to\mathcal K^*/\mathcal O_X^*\to\mathcal K^*/\pi_*\tilde{\mathcal O}_{\tilde X}^*\to0$.]
\item Use (a) to give another proof of the fact that if $X$ is a plane cuspidal cubic curve, then there is an exact sequence $0\to\mathbf G_a\to\operatorname{Pic}X\to\mathbf Z\to0$, and if $X$ is a plane nodal cubic curve, there is an exact sequence $0\to\mathbf G_m\to\operatorname{Pic}X\to\mathbf Z\to0$.
\end{enumerate}}
\pf
\emph{Setting.} $k$ is algebraically closed, a curve is integral, $K=K(X)=K(\tilde X)$, $\tilde X$ is a nonsingular projective curve and $\pi$ is finite and birational (Ex.~3.8). For $P\in X$, $\tilde{\mathcal O}_P=(\pi_*\mathcal O_{\tilde X})_P=\bigcap_{Q\in\pi^{-1}(P)}\mathcal O_{\tilde X,Q}$, a semilocal Dedekind domain, hence a PID. Since $\pi$ is an isomorphism over the (dense, open) nonsingular locus, the set $\Sigma$ of $P$ with $\mathcal O_P\ne\tilde{\mathcal O}_P$ is finite; so the sum in (a) is finite. As $X$ and $\tilde X$ are integral, (6.13)--(6.15) identify $\operatorname{Pic}$ with $\operatorname{CaCl}=\operatorname{Cart}/\operatorname{Prin}$, where $\operatorname{Cart}X=\Gamma(X,\mathcal K^*/\mathcal O_X^*)$ and $\operatorname{Prin}$ is the image of $K^*$.

\textbf{(a)} Define $\pi^\#:\operatorname{Cart}X\to\operatorname{Cart}\tilde X$ by $\{(U_i,f_i)\}\mapsto\{(\pi^{-1}U_i,f_i)\}$ (well defined since units pull back to units; $\pi^\#$ is induced by $\mathcal K^*/\mathcal O_X^*\to\pi_*(\mathcal K^*/\mathcal O_{\tilde X}^*)$). It maps principal to principal, $\pi^\#(g)=(g)$, and, as in Ex.~6.8(b) ($\pi$ is dominant), $\mathcal L(\pi^\#D)\cong\pi^*\mathcal L(D)$, so on $\operatorname{Pic}=\operatorname{Cart}/\operatorname{Prin}$ the map induced by $\pi^\#$ is $\pi^*$. I show:
\[0\to N\to\operatorname{Cart}X\xrightarrow{\pi^\#}\operatorname{Cart}\tilde X\to0\qquad\text{is exact, with }N\cong\bigoplus_P\tilde{\mathcal O}_P^*/\mathcal O_P^*.\qquad(4)\]
(Recall that this is just the long exact sequence of the hint, with $\mathcal N=\pi_*\tilde{\mathcal O}^*_{\tilde X}/\mathcal O_X^*$ a skyscraper sheaf supported on $\Sigma$, so flasque; I give the elementary argument.)

\emph{Kernel.} A Cartier divisor $D$ is determined by its germs $f_P\in K^*/\mathcal O_P^*$, and $\pi^\#D=0$ iff $f_P\in\mathcal O_{\tilde X,Q}^*$ for all $Q\mapsto P$, i.e.\ $f_P\in\bigcap_Q\mathcal O_{\tilde X,Q}^*=\tilde{\mathcal O}_P^*$ (if $f,f^{-1}\in\bigcap\mathcal O_Q=\tilde{\mathcal O}_P$). So $D\mapsto(f_P\bmod\mathcal O_P^*)_P$ is an injection of $N:=\ker\pi^\#$ into $\bigoplus_{P\in\Sigma}\tilde{\mathcal O}_P^*/\mathcal O_P^*$ (the germs vanish at $P\notin\Sigma$). It is onto: given $u_P\in\tilde{\mathcal O}_P^*$ ($P\in\Sigma$), $u_P^{\pm1}$ are regular on a neighbourhood of $\pi^{-1}(P)$, hence $u_P$ is a unit of $\tilde{\mathcal O}_{P'}$ for $P'$ in an open neighbourhood $W_P$ of $P$ which we may take with $W_P\cap\Sigma=\{P\}$. The cover $\{W_P\}_{P\in\Sigma}\cup\{W_0=X-\Sigma\}$ with local equations $u_P$ and $1$ is a Cartier divisor: on overlaps we have only points $P'\notin\Sigma$, where $\mathcal O_{P'}=\tilde{\mathcal O}_{P'}$ and $u_P$ is a unit. It has germs $u_P$ at $P$.

\emph{Surjectivity of $\pi^\#$.} By (6.11) ($\tilde X$ nonsingular), $\operatorname{Cart}\tilde X=\operatorname{Div}\tilde X$. Let $\tilde D=\sum n_QQ\in\operatorname{Div}\tilde X$. The finite set $T_\Sigma=\pi^{-1}(\Sigma)$ lies in an affine open $\operatorname{Spec}C$ of $\tilde X$ (complement of a hyperplane section of the projective curve $\tilde X$ missing $T_\Sigma$). Let $B=\bigcap_{Q\in T_\Sigma}\mathcal O_{\tilde X,Q}$, the localization of $C$ at the complement of $\bigcup_Q\mathfrak m_Q\cap C$; it is a semilocal Dedekind domain, hence a PID, with maximal ideals $\mathfrak m_QB$ and $B_{\mathfrak m_QB}=\mathcal O_{\tilde X,Q}$. The fractional ideal $\prod_Q(\mathfrak m_QB)^{n_Q}$ is principal, $=gB$ with $g\in K^*$, and $v_Q(g)=n_Q$ for all $Q\in T_\Sigma$ (compare in each localization $B_{\mathfrak m_QB}$). Then $\tilde D'=\tilde D-(g)$ has no component over $\Sigma$; its (finite) support $T$ lies in $\pi^{-1}(X-\Sigma)\cong X-\Sigma$. Let $D'$ be the Cartier divisor on $X$ which on the open set $X-\Sigma\cong\pi^{-1}(X-\Sigma)$ is $\tilde D'$, and which is $0$ on the open set $X-\pi(T)\supseteq\Sigma$; they agree on the overlap. Then $\pi^\#D'=\tilde D'$, so $\tilde D=\pi^\#D'+(g)=\pi^\#(D'+(g))$.

\emph{Passing to $\operatorname{Pic}$.} We now pass to the quotient by principal divisors. The map $\operatorname{Pic}X\to\operatorname{Pic}\tilde X$ is surjective by (4). If $\pi^\#D=(g)$ is principal, $g\in K^*$, then $D-(g)\in N$, so $\ker(\operatorname{Pic}X\to\operatorname{Pic}\tilde X)$ is the image of $N$. Finally $N\to\operatorname{Pic}X$ is injective: if $n\in N$ is principal, $n=(g)$ in $\operatorname{Cart}X$, then $0=\pi^\#n=(g)$ in $\operatorname{Cart}\tilde X$, so $g\in\Gamma(\tilde X,\mathcal O^*)=k^*$ ($\tilde X$ is a complete integral variety over algebraically closed $k$), and then $(g)=0$ in $\operatorname{Cart}X$, i.e.\ $n=0$. This proves the exact sequence of (a).

\textbf{(b)} Both curves have exactly one singular point $Z$, and $\tilde X=\mathbf P^1$, so $\operatorname{Pic}\tilde X=\mathbf Z$ (6.17) and the sum in (a) reduces to $\tilde{\mathcal O}_Z^*/\mathcal O_Z^*$. So it suffices to show that this group is $\mathbf G_a$ (cusp), resp.\ $\mathbf G_m$ (node). We use the normal forms $y^2z=x^3$ and $y^2z=x^3+x^2z$ of (6.11.4) and Ex.~6.7 (the cuspidal, resp.\ nodal, cubic of the statement).

\emph{Node} ($y^2z=x^3+x^2z$, $\operatorname{char}k\ne2$): by Ex.~6.7, $\tilde{\mathcal O}_Z=k[t]_{S}$ ($S$ = polynomials not vanishing at $\pm1$), and $\mathcal O_Z=\{f\in\tilde{\mathcal O}_Z:f(1)=f(-1)\}$. The map $\tilde{\mathcal O}_Z^*\to k^*$, $f\mapsto f(1)/f(-1)$, is a homomorphism, surjective (for $c\in k^*$ the polynomial $f=\frac{c+1}2+\frac{c-1}2t$ has $f(1)=c$, $f(-1)=1$), with kernel $\{f:f(1)=f(-1)\}\cap\tilde{\mathcal O}_Z^*=\mathcal O_Z^*$ (if $f\in\mathcal O_Z$ is a unit of $\tilde{\mathcal O}_Z$, then $f(\pm1)\ne0$ and $1/f\in\mathcal O_Z$). So $\tilde{\mathcal O}_Z^*/\mathcal O_Z^*\cong\mathbf G_m$.

\emph{Cusp} ($y^2z=x^3$): the parametrization $t\mapsto(t^2,t^3)$, i.e.\ $\varphi(T:S)=(T^2S:T^3:S^3)$ ($\varphi(\infty)=(0:1:0)=P_0$), is the normalization $\mathbf P^1\to X$ (finite, birational with $t=y/x$; over $x\ne0$ it is an iso, as $k[t^2,t^3][t^{-2}]=k[t,t^{-1}]$, and near $P_0$, in the chart $y=1$ the curve is $z=x^3$, nonsingular); the singular point is $Z=(0,0,1)$ with the single preimage $t=0$. Here $A=k[x,y]/(y^2-x^3)\cong k[t^2,t^3]=\{g\in k[t]:g'(0)=0\}$ (the coefficient of $t$ vanishes), and $\tilde{\mathcal O}_Z=k[t]_{(t)}$, $\mathcal O_Z=\{f\in\tilde{\mathcal O}_Z:f'(0)=0\}$, where $f'(0)$ is the coefficient of $t$ in the expansion of $f$ at $0$. (``$\subseteq$'' is clear; conversely if $f=a/b$ with $a,b\in k[t]$, $b(0)\ne0$, $f'(0)=0$, put $\tilde b=b(0)-b'(0)t$; then $b\tilde b$ and $a\tilde b$ have no $t$-term: the coefficient is $b(0)b'(0)-b'(0)b(0)=0$, resp.\ $a'(0)b(0)-a(0)b'(0)=b(0)^2f'(0)=0$. So $f=(a\tilde b)/(b\tilde b)$ with numerator and denominator in $A$, $b\tilde b(0)=b(0)^2\ne0$.) The map $\tilde{\mathcal O}_Z^*\to(k,+)$, $f\mapsto f'(0)/f(0)$ is a homomorphism (logarithmic derivative), surjective ($f=1+ct\mapsto c$), with kernel $\{f:f'(0)=0\}\cap\tilde{\mathcal O}_Z^*=\mathcal O_Z^*$ (if $f\in\mathcal O_Z$ is a unit of $\tilde{\mathcal O}_Z$ then $(1/f)'(0)=-f'(0)/f(0)^2=0$). So $\tilde{\mathcal O}_Z^*/\mathcal O_Z^*\cong\mathbf G_a$.

In both cases (a) gives the stated exact sequences $0\to\mathbf G_a\to\operatorname{Pic}X\to\mathbf Z\to0$, $0\to\mathbf G_m\to\operatorname{Pic}X\to\mathbf Z\to0$.\hfill$\square$

\exercise{6.10}{\emph{The Grothendieck Group $K(X)$.} Let $X$ be a noetherian scheme. We define $K(X)$ to be the quotient of the free abelian group generated by all the coherent sheaves on $X$, by the subgroup generated by all expressions $\mathcal F-\mathcal F'-\mathcal F''$, whenever there is an exact sequence $0\to\mathcal F'\to\mathcal F\to\mathcal F''\to0$ of coherent sheaves on $X$. If $\mathcal F$ is a coherent sheaf, we denote by $\gamma(\mathcal F)$ its image in $K(X)$.
\begin{enumerate}[label=(\alph*),leftmargin=*]
\item If $X=\mathbf A^1_k$, then $K(X)\cong\mathbf Z$.
\item If $X$ is any integral scheme, define the \emph{rank} of a coherent sheaf $\mathcal F$ to be $\dim_K\mathcal F_\xi$, where $\xi$ is the generic point of $X$, and $K=\mathcal O_\xi$ is the function field of $X$. Show that the rank function defines a surjective homomorphism $\operatorname{rank}:K(X)\to\mathbf Z$.
\item If $Y$ is a closed subscheme of $X$, there is an exact sequence $K(Y)\to K(X)\to K(X-Y)\to0$, where the first map is extension by zero, and the second is restriction. [\emph{Hint}: For exactness in the middle, show that if $\mathcal F$ is a coherent sheaf on $X$, whose support is contained in $Y$, then there is a finite filtration $\mathcal F=\mathcal F_0\supseteq\mathcal F_1\supseteq\dots\supseteq\mathcal F_n=0$ such that each $\mathcal F_i/\mathcal F_{i+1}$ is an $\mathcal O_Y$-module. To show surjectivity on the right use (Ex.~5.15).]
\end{enumerate}}
\pf
\textbf{(b)} The stalk functor $\mathcal F\mapsto\mathcal F_\xi$ is exact, and $\mathcal F_\xi$ is a finite-dimensional $K$-vector space ($\mathcal F$ coherent), so $\dim_K$ is additive on short exact sequences. Hence $\operatorname{rank}$ defined on generators $\mathcal F$ kills the relations $\mathcal F-\mathcal F'-\mathcal F''$ and descends to a homomorphism $K(X)\to\mathbf Z$. It is surjective since $\operatorname{rank}\mathcal O_X=1$.

\textbf{(a)} Let $X=\operatorname{Spec}A$, $A=k[t]$, a PID. Coherent sheaves on $X$ are the $\widetilde M$ with $M$ a finitely generated $A$-module, and exact sequences of such sheaves correspond to exact sequences of modules (II.5.4, 5.5). By the structure theorem, $M\cong A^r\oplus\bigoplus_iA/(f_i)$ with $f_i\ne0$. For $f\ne0$, the exact sequence $0\to A\xrightarrow{f}A\to A/(f)\to0$ gives $\gamma(\widetilde{A/(f)})=0$ in $K(X)$, and $\gamma(\widetilde M)=\sum\gamma(\text{summands})$ by the exact sequences $0\to M'\to M'\oplus M''\to M''\to0$. So $\gamma(\widetilde M)=r\gamma(\mathcal O_X)$, $K(X)$ is generated by $\gamma(\mathcal O_X)$, and by (b) $\operatorname{rank}:K(X)\to\mathbf Z$ is onto with $\operatorname{rank}\gamma(\mathcal O_X)=1$. Thus $\mathbf Z\to K(X)$, $n\mapsto n\gamma(\mathcal O_X)$ is surjective and splits by $\operatorname{rank}$; hence it is an isomorphism and $K(X)\cong\mathbf Z$.

\textbf{(c)} Let $i:Y\to X$ be the closed immersion, $\mathcal J$ its ideal sheaf, $U=X-Y$, $j:U\to X$. 

\emph{The maps.} For $\mathcal G$ coherent on $Y$, $i_*\mathcal G$ is coherent on $X$ (on an affine open $\operatorname{Spec}B$ of $X$ with $Y\cap\operatorname{Spec}B=\operatorname{Spec}B/I$, $i_*\widetilde N=\widetilde N$ with $N$ a finitely generated $B/I$-module, hence a finitely generated $B$-module), and $i_*$ is exact, so $i_*:K(Y)\to K(X)$ is a homomorphism; restriction to $U$ is exact on coherent sheaves, so $\rho:K(X)\to K(U)$ is a homomorphism; $\rho\circ i_*=0$ since $(i_*\mathcal G)|_U=0$.

\emph{Surjectivity of $\rho$.} By Ex.~5.15, every coherent $\mathcal G$ on $U$ extends to a coherent $\bar{\mathcal G}$ on $X$ with $\bar{\mathcal G}|_U\cong\mathcal G$; so $\gamma(\mathcal G)=\rho(\gamma(\bar{\mathcal G}))$, and these generate $K(U)$.

\emph{Supports in $Y$.} Let $\mathcal F$ be coherent with $\operatorname{Supp}\mathcal F\subseteq Y$. Then $\mathcal J^N\mathcal F=0$ for some $N$: on an affine $\operatorname{Spec}B$, $\mathcal F=\widetilde M$ with $M$ finitely generated and $V(\operatorname{Ann}M)=\operatorname{Supp}M\subseteq V(I)$, $I=\mathcal J(B)$ finitely generated, so $I\subseteq\sqrt{\operatorname{Ann}M}$ and $I^N\subseteq\operatorname{Ann}M$; take the maximum over a finite affine cover. The filtration $\mathcal F\supseteq\mathcal J\mathcal F\supseteq\dots\supseteq\mathcal J^N\mathcal F=0$ has coherent quotients $\mathcal F_m/\mathcal F_{m+1}$ annihilated by $\mathcal J$, i.e.\ of the form $i_*\mathcal G_m$ with $\mathcal G_m$ coherent on $Y$ (a coherent $\mathcal O_X$-module $\mathcal M$ with $\mathcal J\mathcal M=0$ is locally $\widetilde M$ with $M$ a finitely generated module over $B/I$, and by II.5.2 this is the direct image of $\widetilde M$ on $\operatorname{Spec}B/I$; so $\mathcal M=i_*\mathcal G$ with $\mathcal G=i^{-1}\mathcal M$ coherent on $Y$). Hence $\gamma(\mathcal F)=\sum_m i_*\gamma(\mathcal G_m)\in\operatorname{im}(i_*)$.

\emph{Exactness in the middle.} Let $Q=K(X)/\operatorname{im}(i_*)$. We construct $\theta:K(U)\to Q$ with $\theta\circ\rho=$ the quotient map $K(X)\to Q$; then $\rho(x)=0\Rightarrow\bar x=\theta(\rho x)=0$, i.e.\ $\ker\rho\subseteq\operatorname{im}(i_*)$, as required.

(i) \emph{Two extensions agree in $Q$.} Let $\mathcal F_1,\mathcal F_2$ be coherent on $X$ and $\alpha:\mathcal F_1|_U\xrightarrow{\ \sim\ }\mathcal F_2|_U$. The graph $\Gamma\subseteq(\mathcal F_1\oplus\mathcal F_2)|_U$ of $\alpha$ is a coherent subsheaf, so by Ex.~5.15 there is a coherent subsheaf $\bar\Gamma\subseteq\mathcal F_1\oplus\mathcal F_2$ with $\bar\Gamma|_U=\Gamma$. The projections $p_\nu:\bar\Gamma\to\mathcal F_\nu$ restrict to isomorphisms over $U$ (they restrict to $\Gamma\to\mathcal F_\nu|_U$, which are isomorphisms). So $\ker p_\nu$ and $\operatorname{coker}p_\nu$ are coherent with support in $Y$, hence in $\operatorname{im}(i_*)$ by the above. From $0\to\ker p_\nu\to\bar\Gamma\to\operatorname{im}p_\nu\to0$ and $0\to\operatorname{im}p_\nu\to\mathcal F_\nu\to\operatorname{coker}p_\nu\to0$ we get $\gamma(\bar\Gamma)-\gamma(\mathcal F_\nu)=\gamma(\ker p_\nu)-\gamma(\operatorname{coker}p_\nu)\in\operatorname{im}(i_*)$ for $\nu=1,2$. So $\gamma(\mathcal F_1)\equiv\gamma(\mathcal F_2)$ in $Q$.

(ii) For $\mathcal G$ coherent on $U$ put $\theta(\gamma(\mathcal G))=\overline{\gamma(\bar{\mathcal G})}$ for any coherent extension $\bar{\mathcal G}$ (exists by Ex.~5.15; independent of the choice by (i)); this defines $\theta$ on the free abelian group on coherent sheaves of $U$. It kills the relations: for $0\to\mathcal G'\to\mathcal G\to\mathcal G''\to0$ on $U$, extend $\mathcal G$ to $\bar{\mathcal G}$, then (Ex.~5.15(d) applied to $\mathcal G'\subseteq\bar{\mathcal G}|_U$) choose a coherent $\bar{\mathcal G}'\subseteq\bar{\mathcal G}$ with $\bar{\mathcal G}'|_U=\mathcal G'$, and put $\bar{\mathcal G}''=\bar{\mathcal G}/\bar{\mathcal G}'$, a coherent extension of $\mathcal G''$. From $0\to\bar{\mathcal G}'\to\bar{\mathcal G}\to\bar{\mathcal G}''\to0$ we get $\theta\gamma(\mathcal G)=\theta\gamma(\mathcal G')+\theta\gamma(\mathcal G'')$. So $\theta:K(U)\to Q$ is a homomorphism, and $\theta(\rho\gamma(\mathcal F))=\theta(\gamma(\mathcal F|_U))=\overline{\gamma(\mathcal F)}$ (take $\bar{\mathcal G}=\mathcal F$), so $\theta\circ\rho$ is the quotient map. Hence $K(Y)\to K(X)\to K(X-Y)\to0$ is exact.\hfill$\square$

\exercise{6.11}{\emph{The Grothendieck Group of a Nonsingular Curve.} Let $X$ be a nonsingular curve over an algebraically closed field $k$. We will show that $K(X)\cong\operatorname{Pic}X\oplus\mathbf Z$, in several steps. (Starred in the book.)
\begin{enumerate}[label=(\alph*),leftmargin=*]
\item For any divisor $D=\sum n_iP_i$ on $X$, let $\psi(D)=\sum n_i\gamma(k(P_i))\in K(X)$, where $k(P_i)$ is the skyscraper sheaf $k$ at $P_i$ and $0$ elsewhere. If $D$ is an effective divisor, let $\mathcal O_D$ be the structure sheaf of the associated subscheme of codimension $1$, and show that $\psi(D)=\gamma(\mathcal O_D)$. Then use (6.18) to show that for any $D$, $\psi(D)$ depends only on the linear equivalence class of $D$, so $\psi$ defines a homomorphism $\psi:\operatorname{Cl}X\to K(X)$.
\item For any coherent sheaf $\mathcal F$ on $X$, show that there exist locally free sheaves $\mathcal E_0,\mathcal E_1$ and an exact sequence $0\to\mathcal E_1\to\mathcal E_0\to\mathcal F\to0$. Let $r_0=\operatorname{rank}\mathcal E_0$, $r_1=\operatorname{rank}\mathcal E_1$, and define $\det\mathcal F=(\bigwedge^{r_0}\mathcal E_0)\otimes(\bigwedge^{r_1}\mathcal E_1)^{-1}\in\operatorname{Pic}X$. Here $\bigwedge$ denotes the exterior power (Ex.~5.16). Show that $\det\mathcal F$ is independent of the resolution chosen, and that it gives a homomorphism $\det:K(X)\to\operatorname{Pic}X$. Finally show that if $D$ is a divisor, then $\det(\psi(D))=\mathcal L(D)$.
\item If $\mathcal F$ is any coherent sheaf of rank $r$, show that there is a divisor $D$ on $X$ and an exact sequence $0\to\mathcal L(D)^{\oplus r}\to\mathcal F\to\mathcal T\to0$, where $\mathcal T$ is a torsion sheaf. Conclude that if $\mathcal F$ is a coherent sheaf of rank $r$, then $\gamma(\mathcal F)-r\gamma(\mathcal O_X)\in\operatorname{Im}\psi$.
\item Using the maps $\psi$, $\det$, $\operatorname{rank}$, and $1\mapsto\gamma(\mathcal O_X)$ from $\mathbf Z\to K(X)$, show that $K(X)\cong\operatorname{Pic}X\oplus\mathbf Z$.
\end{enumerate}}
\pf
\emph{Preliminaries.} $X$ is integral, noetherian, separated, and all local rings $\mathcal O_P$ ($P$ closed) are DVR's, so $X$ is locally factorial and (6.16) gives $\operatorname{Cl}X\cong\operatorname{Pic}X$, $D\mapsto\mathcal L(D)$. Prime divisors are closed points. $X$ is quasi-projective (it is an open subscheme of a complete nonsingular, hence projective, curve (I.6.7, II.6.7)). A coherent sheaf $\mathcal F$ is a \emph{torsion sheaf} if $\mathcal F_\xi=0$, equivalently its support is a finite set of closed points; then $\mathcal F=\bigoplus_P\mathcal F_P$ (skyscraper) and each $\mathcal F_P$ is a module of finite length $\ell_P$ over $\mathcal O_P$.

\textbf{(a)} Let $D=\sum n_iP_i$ be effective; list its points with multiplicity as $Q_1,\dots,Q_m$ and put $D_j=Q_1+\dots+Q_j$, so $D_0=0$, $D_m=D$, $D_{j+1}=D_j+Q_{j+1}$. By (6.18), $\mathcal I_{D_j}=\mathcal L(-D_j)\subseteq\mathcal O_X$ is the ideal sheaf of the subscheme $D_j$, so $\mathcal O_{D_j}=\mathcal O_X/\mathcal L(-D_j)$. Put $P=Q_{j+1}$. Since $\mathcal L(-D_{j+1})=\mathcal L(-D_j)\otimes\mathcal L(-P)=\mathcal L(-D_j)\cdot\mathcal I_P$ (6.13(b), (6.18)), the kernel of the surjection $\mathcal O_{D_{j+1}}\to\mathcal O_{D_j}$ is $\mathcal L(-D_j)/\mathcal L(-D_j)\mathcal I_P\cong\mathcal L(-D_j)\otimes k(P)\cong k(P)$ (an invertible sheaf is trivial near $P$). So $\gamma(\mathcal O_{D_{j+1}})=\gamma(\mathcal O_{D_j})+\gamma(k(Q_{j+1}))$, and by induction $\gamma(\mathcal O_D)=\sum_j\gamma(k(Q_j))=\sum n_i\gamma(k(P_i))=\psi(D)$.

Let $D,D'$ be effective with $D\sim D'$. By (6.13)(c), $\mathcal L(-D)\cong\mathcal L(-D')$ as abstract sheaves. From $0\to\mathcal L(-D)\to\mathcal O_X\to\mathcal O_D\to0$ we get $\psi(D)=\gamma(\mathcal O_D)=\gamma(\mathcal O_X)-\gamma(\mathcal L(-D))$, which depends only on the isomorphism class of $\mathcal L(-D)$; so $\psi(D)=\psi(D')$. For arbitrary $D\sim D'$, $D-D'=(f)=E_1-E_2$ with $E_1,E_2$ effective (zeros and poles of $f$), and $E_1\sim E_2$, so $\psi(E_1)=\psi(E_2)$; since $\psi$ is additive on $\operatorname{Div}X$ by definition, $\psi(D)-\psi(D')=\psi(E_1)-\psi(E_2)=0$. So $\psi:\operatorname{Cl}X\to K(X)$ is a well-defined homomorphism.

\textbf{(b)} \emph{Existence of a resolution.} Embed $X$ as an open subscheme of a projective curve $\bar X\subseteq\mathbf P^N$. By Ex.~5.15, $\mathcal F$ extends to a coherent sheaf $\bar{\mathcal F}$ on $\bar X$; by (5.18), $\bar{\mathcal F}(m)$ is generated by finitely many global sections for $m\gg0$, i.e.\ there is a surjection $\mathcal O_{\bar X}(-m)^{\oplus s}\to\bar{\mathcal F}$; restricting to $X$ gives a surjection $\mathcal E_0\to\mathcal F$ with $\mathcal E_0$ locally free of finite rank. Let $\mathcal E_1$ be its kernel; it is coherent, and each stalk is a finitely generated submodule of a free module over a DVR (or a field), hence free; so $\mathcal E_1$ is locally free.

We use Ex.~5.16(d): for an exact sequence $0\to\mathcal A\to\mathcal B\to\mathcal C\to0$ of locally free sheaves, $\det\mathcal B:=\bigwedge^{\operatorname{rk}}\mathcal B\cong\det\mathcal A\otimes\det\mathcal C$ (write $\det\mathcal E$ for the top exterior power of a locally free $\mathcal E$). Thus $\det\mathcal F=\det\mathcal E_0\otimes(\det\mathcal E_1)^{-1}$.

\emph{Independence.} Let $0\to\mathcal E_1\to\mathcal E_0\xrightarrow{\ \epsilon\ }\mathcal F\to0$ and $0\to\mathcal E_1'\to\mathcal E_0'\xrightarrow{\ \epsilon'\ }\mathcal F\to0$ be two resolutions. Let $\mathcal E_0''=\ker(\mathcal E_0\oplus\mathcal E_0'\xrightarrow{\epsilon-\epsilon'}\mathcal F)$, locally free (a subsheaf of a locally free sheaf on $X$, by the argument above). The projection $\mathcal E_0''\to\mathcal E_0$ is surjective (given a local section $a$ of $\mathcal E_0$, since $\epsilon'$ is onto locally there is $b$ with $\epsilon'(b)=\epsilon(a)$) with kernel $\{(0,b):\epsilon'(b)=0\}\cong\mathcal E_1'$; symmetrically $\mathcal E_0''\to\mathcal E_0'$ is onto with kernel $\mathcal E_1$. So there are exact sequences of locally free sheaves $0\to\mathcal E_1'\to\mathcal E_0''\to\mathcal E_0\to0$ and $0\to\mathcal E_1\to\mathcal E_0''\to\mathcal E_0'\to0$, whence $\det\mathcal E_0''\cong\det\mathcal E_1'\otimes\det\mathcal E_0\cong\det\mathcal E_1\otimes\det\mathcal E_0'$. Rearranging in the group $\operatorname{Pic}X$: $\det\mathcal E_0\otimes(\det\mathcal E_1)^{-1}\cong\det\mathcal E_0'\otimes(\det\mathcal E_1')^{-1}$. So $\det\mathcal F\in\operatorname{Pic}X$ is well defined.

\emph{Additivity.} Let $0\to\mathcal F'\to\mathcal F\to\mathcal F''\to0$ be exact. Take a resolution $0\to\mathcal E_1\to\mathcal E_0\xrightarrow{\epsilon}\mathcal F\to0$ of $\mathcal F$ and let $\mathcal K''=\ker(\mathcal E_0\to\mathcal F\to\mathcal F'')$, locally free. Then $\epsilon$ maps $\mathcal K''$ onto $\mathcal F'$ with kernel $\mathcal E_1$, so $0\to\mathcal E_1\to\mathcal K''\to\mathcal F'\to0$ is a resolution of $\mathcal F'$, and $0\to\mathcal K''\to\mathcal E_0\to\mathcal F''\to0$ is a resolution of $\mathcal F''$. Hence
\[
\det\mathcal F'\otimes\det\mathcal F''=\bigl(\det\mathcal K''\otimes\det\mathcal E_1^{-1}\bigr)\otimes\bigl(\det\mathcal E_0\otimes\det\mathcal K''^{-1}\bigr)=\det\mathcal E_0\otimes\det\mathcal E_1^{-1}=\det\mathcal F.
\]
By independence of the resolution, $\det$ is additive on exact sequences, so $\det$ defined on generators extends to a homomorphism $\det:K(X)\to\operatorname{Pic}X$.

\emph{$\det\psi(D)=\mathcal L(D)$.} For a point $P$, $0\to\mathcal L(-P)\to\mathcal O_X\to k(P)\to0$ is a resolution of $k(P)$ ((6.18): $\mathcal I_P=\mathcal L(-P)$), so $\det k(P)=\det\mathcal O_X\otimes\mathcal L(-P)^{-1}\cong\mathcal L(P)$ (6.13(b)). Since $\det$ is a homomorphism and $D\mapsto\mathcal L(D)$ is one, $\det\psi(\sum n_iP_i)=\bigotimes\mathcal L(P_i)^{n_i}\cong\mathcal L(D)$.

\textbf{(c)} First a lemma: for every divisor $D$, $\gamma(\mathcal L(D))=\gamma(\mathcal O_X)+\psi(D)$. (\emph{Proof.}) For a point $P$, $0\to\mathcal L(D-P)\to\mathcal L(D)\to\mathcal L(D)\otimes k(P)\to0$ is exact (tensor $0\to\mathcal I_P\to\mathcal O_X\to k(P)\to0$ with the invertible $\mathcal L(D)$, using $\mathcal L(D)\otimes\mathcal I_P=\mathcal L(D-P)$) and $\mathcal L(D)\otimes k(P)\cong k(P)$. So $\gamma(\mathcal L(D-P))=\gamma(\mathcal L(D))-\gamma(k(P))$. The function $\phi(D)=\gamma(\mathcal L(D))-\gamma(\mathcal O_X)-\psi(D)$ therefore satisfies $\phi(D-P)=\phi(D)$ for all $D,P$, hence also $\phi(D+P)=\phi(D)$; as every divisor is reached from $0$ by finitely many such steps, $\phi(D)=\phi(0)=0$.

Now let $\mathcal F$ be coherent of rank $r$. Let $\mathcal T_0\subseteq\mathcal F$ be the torsion subsheaf (sections killed by nonzero elements locally); it is coherent with $(\mathcal T_0)_\xi=0$ (the generic stalk $\mathcal F_\xi$ is a $K$-vector space, torsion-free), so it is a torsion sheaf, supported on a finite set $\Sigma$. The quotient $\mathcal F_0=\mathcal F/\mathcal T_0$ is torsion-free, hence locally free of rank $r$ (stalks are finitely generated torsion-free over a DVR). Choose a basis $e_1,\dots,e_r$ of the $K$-vector space $(\mathcal F_0)_\xi$; since $\mathcal F_0$ is torsion-free, $\mathcal F_0\subseteq$ the constant sheaf $(\mathcal F_0)_\xi=\bigoplus Ke_\nu$. Let $\mathcal G_\nu=\mathcal F_0\cap Ke_\nu$; this is a coherent subsheaf of $\mathcal F_0$ of rank one, torsion-free, hence invertible, so $\mathcal G_\nu\cong\mathcal L(D_\nu)$ for some divisor $D_\nu$ (6.16). Let $S$ be the (finite) union of the supports of $D_1,\dots,D_r$, let $N$ be larger than the absolute values of all coefficients of all $D_\nu$, and put $D=-N\sum_{P\in S}P$, so that $D\le D_\nu$ for all $\nu$ (coefficientwise). Then $\mathcal L(D)\subseteq\mathcal L(D_\nu)$ as sub-$\mathcal O_X$-modules of $\mathcal K$, and composing with the isomorphisms $\mathcal L(D_\nu)\cong\mathcal G_\nu$ we get an injection $\mathcal L(D)^{\oplus r}\hookrightarrow\bigoplus_\nu\mathcal G_\nu\hookrightarrow\mathcal F_0$. At $\xi$ it is an isomorphism ($\mathcal L(D)_\xi=K\to\mathcal G_{\nu,\xi}=Ke_\nu$ is a nonzero map of one-dimensional spaces and $\bigoplus Ke_\nu=(\mathcal F_0)_\xi$), so its cokernel $\mathcal T'$ is a torsion sheaf. (Why $\mathcal G_\nu$ is coherent: it is the kernel of $\mathcal F_0\to\underline{K}^r/Ke_\nu$, a map from a coherent to a quasi-coherent sheaf on a noetherian scheme.)

\emph{Lift to $\mathcal F$.} Let $\iota:\mathcal L(D)^{\oplus r}\to\mathcal F_0$ be this injection. We lift it to $\tilde\iota:\mathcal L(D)^{\oplus r}\to\mathcal F$. On $V=X-\Sigma$, $\mathcal T_0|_V=0$ so $\mathcal F|_V=\mathcal F_0|_V$ and we take $\iota|_V$. For $P\in\Sigma$ choose an open neighbourhood $W_P$ of $P$ with $W_P\cap\Sigma=\{P\}$ and $\mathcal L(D)|_{W_P}$ trivial; then $\iota|_{W_P}$ is given by $r$ sections of $\mathcal F_0(W_P)$, and since $\mathcal F\to\mathcal F_0$ is surjective, after shrinking $W_P$ these lift to sections of $\mathcal F(W_P)$. The liftings on $W_P\cap V=W_P-\{P\}$ agree with $\iota|_V$ because $\mathcal F|_V=\mathcal F_0|_V$. So the local lifts glue to $\tilde\iota$. The composite $\mathcal L(D)^{\oplus r}\to\mathcal F\to\mathcal F_0$ is $\iota$, injective, so $\tilde\iota$ is injective; its cokernel $\mathcal T$ is an extension of $\mathcal T'$ by $\mathcal T_0$ ($0\to\mathcal T_0\to\mathcal T\to\mathcal T'\to0$ from the snake lemma), so a torsion sheaf. This is the exact sequence $0\to\mathcal L(D)^{\oplus r}\to\mathcal F\to\mathcal T\to0$.

\emph{Consequence.} A torsion sheaf $\mathcal T=\bigoplus_P\mathcal T_P$ with $\mathcal T_P$ of finite length $\ell_P$ has a filtration with quotients $k(P)$ (a composition series of $\mathcal T_P$ over $\mathcal O_P$, whose simple module is $k$), so $\gamma(\mathcal T)=\sum\ell_P\gamma(k(P))=\psi(\sum\ell_PP)$. Hence $\gamma(\mathcal F)=r\gamma(\mathcal L(D))+\gamma(\mathcal T)=r\gamma(\mathcal O_X)+\psi(rD)+\psi(\sum\ell_PP)$, i.e.\ $\gamma(\mathcal F)-r\gamma(\mathcal O_X)\in\operatorname{Im}\psi$.

\textbf{(d)} Define $\Theta:\operatorname{Pic}X\oplus\mathbf Z\to K(X)$, $\Theta(\mathcal L(D),m)=\psi([D])+m\gamma(\mathcal O_X)$ (well defined by (6.16) and (a)), and $\Phi=(\det,\operatorname{rank}):K(X)\to\operatorname{Pic}X\oplus\mathbf Z$. Both are homomorphisms. 

$\Phi\Theta=\operatorname{id}$: $\det\psi(D)=\mathcal L(D)$ by (b), $\operatorname{rank}\psi(D)=0$ (torsion sheaves have rank $0$); $\det\gamma(\mathcal O_X)=\mathcal O_X$ (the resolution $0\to0\to\mathcal O_X\to\mathcal O_X\to0$) and $\operatorname{rank}\gamma(\mathcal O_X)=1$. So $\Phi\Theta(\mathcal L(D),m)=(\mathcal L(D),m)$.

$\Theta\Phi=\operatorname{id}$: it suffices to check on the generators $\gamma(\mathcal F)$. Let $r=\operatorname{rank}\mathcal F$. By (c) there is a divisor $D''$ with $\gamma(\mathcal F)-r\gamma(\mathcal O_X)=\psi(D'')$. Applying $\det$ (and using (b), $\det\gamma(\mathcal O_X)=\mathcal O_X$): $\det\mathcal F=\mathcal L(D'')$. So the class $[D'']\in\operatorname{Cl}X$ is the one corresponding to $\det\mathcal F$ under $\operatorname{Cl}\cong\operatorname{Pic}$, and
\[\Theta\Phi(\gamma\mathcal F)=\Theta(\det\mathcal F,r)=\psi(D'')+r\gamma(\mathcal O_X)=\gamma(\mathcal F).\]
Thus $\Theta$ and $\Phi$ are inverse isomorphisms and $K(X)\cong\operatorname{Pic}X\oplus\mathbf Z$.\hfill$\square$

\exercise{6.12}{Let $X$ be a complete nonsingular curve. Show that there is a unique way to define the \emph{degree} of any coherent sheaf on $X$, $\deg\mathcal F\in\mathbf Z$, such that:
\begin{enumerate}[label=(\arabic*),leftmargin=*]
\item If $D$ is a divisor, $\deg\mathcal L(D)=\deg D$;
\item If $\mathcal F$ is a \emph{torsion sheaf} (meaning a sheaf whose stalk at the generic point is zero), then $\deg\mathcal F=\sum_{P\in X}\operatorname{length}(\mathcal F_P)$; and
\item If $0\to\mathcal F'\to\mathcal F\to\mathcal F''\to0$ is an exact sequence of coherent sheaves, then $\deg\mathcal F=\deg\mathcal F'+\deg\mathcal F''$.
\end{enumerate}}
\pf
The curve $X$ is projective (Prop.~6.7) and all the previous results of Ex.~6.10, 6.11 apply; $\deg D=\sum n_i$ for $D=\sum n_iP_i$, and $\deg:\operatorname{Cl}X\to\mathbf Z$ is a well-defined homomorphism (6.10), which with $\operatorname{Cl}X\cong\operatorname{Pic}X$ (6.16) gives $\deg:\operatorname{Pic}X\to\mathbf Z$ with $\deg\mathcal L(D)=\deg D$.

\emph{Existence.} For coherent $\mathcal F$ put $\deg\mathcal F:=\deg(\det\mathcal F)$, with $\det:K(X)\to\operatorname{Pic}X$ of Ex.~6.11(b) (the determinant is defined for all coherent $\mathcal F$). Property (3) holds since $\det\mathcal F\cong\det\mathcal F'\otimes\det\mathcal F''$ and $\deg$ is additive. For (1): $\mathcal L(D)$ is locally free of rank one, with resolution $0\to0\to\mathcal L(D)\to\mathcal L(D)\to0$, so $\det\mathcal L(D)=\mathcal L(D)$ and $\deg\mathcal L(D)=\deg D$. For (2): a torsion sheaf $\mathcal T$ has $\gamma(\mathcal T)=\sum_P\ell_P\gamma(k(P))$, $\ell_P=\operatorname{length}\mathcal T_P$ (a composition series, as in Ex.~6.11(c)), so $\det\mathcal T=\bigotimes\mathcal L(P)^{\ell_P}$ by Ex.~6.11(b), and $\deg\mathcal T=\sum\ell_P$.

\emph{Uniqueness.} Let $\deg'$ be any function satisfying (1)--(3). By (3) it induces a homomorphism $K(X)\to\mathbf Z$. Let $\mathcal F$ be coherent of rank $r$. By Ex.~6.11(c) there is a divisor $D=\sum n_iP_i$ with $\gamma(\mathcal F)=r\gamma(\mathcal O_X)+\sum n_i\gamma(k(P_i))$. Now $\deg'(\mathcal O_X)=\deg'(\mathcal L(0))=\deg0=0$ by (1), and $\deg'(k(P_i))=\operatorname{length}k(P_i)=1$ by (2). Hence $\deg'(\mathcal F)=\sum n_i=\deg D$. In particular this number is the same for every divisor $D$ with $\gamma(\mathcal F)-r\gamma(\mathcal O_X)=\psi(D)$, and equals $\deg'\mathcal F$; since the same computation applies to any other function $\deg''$ with (1)--(3), $\deg''(\mathcal F)=\deg D=\deg'(\mathcal F)$. So the degree is unique.\hfill$\square$

\section*{7. Projective morphisms}

Throughout this section I use the numbering of Hartshorne's Chapter~II: Thm.~7.1 is the correspondence between morphisms to $\mathbf P^n_A$ and generated invertible sheaves, Prop.~7.2/7.3 are the criteria for closed immersions, Thm.~7.6 says $\mathcal L$ is ample iff some power is very ample, Prop.~7.10 and~7.11 concern $\operatorname{Proj}\mathcal S$ and $\mathbf P(\mathcal E)$, Prop.~7.12 is the universal property of $\mathbf P(\mathcal E)$, and Prop.~7.13--7.15 concern blowing up. ``Very ample'' always means very ample relative to the base in the sense of Section~5 of the text. Where a printed statement needs an extra hypothesis (this happens in Ex.~7.3(b), 7.7(b) and 7.9), I say so at that point.

\exercise{7.1}{Let $(X,\mathcal O_X)$ be a locally ringed space, and let $f:\mathcal L\to\mathcal M$ be a surjective map of invertible sheaves on $X$. Show that $f$ is an isomorphism. [\emph{Hint:} Reduce to a question of modules over a local ring by looking at the stalks.]}
\pf
A morphism of sheaves is an isomorphism if and only if it is an isomorphism on every stalk, and surjectivity of $f$ means that every $f_x$ is surjective. Fix $x\in X$. Since $\mathcal L,\mathcal M$ are invertible, $\mathcal L_x\cong\mathcal O_x\cong\mathcal M_x$, so after choosing generators $f_x$ becomes an $\mathcal O_x$-linear surjection $g:\mathcal O_x\to\mathcal O_x$, which is multiplication by $a=g(1)$. Surjectivity gives $b\in\mathcal O_x$ with $g(b)=ab=1$, so $a$ is a unit. Hence $g$, and so $f_x$, is bijective. (Only that $\mathcal O_x$ is a commutative ring is used; locality is not even needed.) Therefore $f$ is an isomorphism.\hfill$\square$

\exercise{7.2}{Let $X$ be a scheme over a field $k$, and let $\mathcal L$ be an invertible sheaf on $X$. Let $\{s_0,\dots,s_n\}$ and $\{t_0,\dots,t_m\}$ be two sets of sections of $\mathcal L$, which generate the same subspace $V\subseteq\Gamma(X,\mathcal L)$, and which generate the sheaf $\mathcal L$ at every point. Suppose $n\le m$. Show that the corresponding morphisms $\varphi:X\to\mathbf P^n_k$ and $\psi:X\to\mathbf P^m_k$ differ by a suitable linear projection $\mathbf P^m-L\to\mathbf P^n$ and an automorphism of $\mathbf P^m$, where $L$ is a linear subspace of $\mathbf P^m$ of dimension $m-n-1$.}
\pf
Let $\sigma:k^{n+1}\to\Gamma(X,\mathcal L)$, $e_i\mapsto s_i$ and $\tau:k^{m+1}\to\Gamma(X,\mathcal L)$, $e'_j\mapsto t_j$; both have image $V$. Put $d=\dim V$, $K=\Ker\sigma$ ($\dim K=n+1-d$) and $K'=\Ker\tau$ ($\dim K'=m+1-d\ge\dim K$).

\textbf{Claim.} There is an injective linear map $\iota:k^{n+1}\to k^{m+1}$ with $\tau\circ\iota=\sigma$.

Choose a complement $U$ of $K$ in $k^{n+1}$, so $\sigma|_U:U\to V$ is an isomorphism, and a linear section $\varsigma:V\to k^{m+1}$ of the surjection $\tau:k^{m+1}\to V$. Choose an injective linear map $j:K\to K'$ (possible since $\dim K\le\dim K'$). Define $\iota|_U=\varsigma\circ\sigma|_U$ and $\iota|_K=j$. Then $\tau\iota=\sigma$ on $U$ (as $\tau\varsigma=\mathrm{id}$) and on $K$ (both sides are $0$). If $\iota(u+\kappa)=0$ with $u\in U,\kappa\in K$, then $\iota(u)=-j(\kappa)\in K'$, so $\sigma(u)=\tau\iota(u)=0$, hence $u=0$, and then $j(\kappa)=0$ gives $\kappa=0$. So $\iota$ is injective.

Extend the linearly independent vectors $\iota(e_0),\dots,\iota(e_n)$ to a basis $b_0=\iota(e_0),\dots,b_n=\iota(e_n),b_{n+1},\dots,b_m$ of $k^{m+1}$, and let $g\in GL_{m+1}(k)$ be defined by $g(e'_j)=b_j=\sum_l g_{lj}e'_l$. Put $t'_j=\tau(b_j)=\sum_l g_{lj}t_l$. Then:
\begin{itemize}
\item $t'_i=\tau(\iota(e_i))=\sigma(e_i)=s_i$ for $0\le i\le n$;
\item $t'_0,\dots,t'_m$ span $V$ (as $g$ is invertible), so they generate $\mathcal L$ at every point and define a morphism $\psi':X\to\mathbf P^m_k$.
\end{itemize}
Let $\alpha$ be the automorphism of $\mathbf P^m_k$ induced by the graded automorphism $x_j\mapsto\sum_l g_{lj}x_l$ of $k[x_0,\dots,x_m]$ (Example~7.1.1); then $\alpha^*\mathcal O(1)\cong\mathcal O(1)$ with $\alpha^*x_j=\sum_lg_{lj}x_l$. Hence $(\alpha\psi)^*\mathcal O(1)\cong\mathcal L$ and $(\alpha\psi)^*x_j=\sum_l g_{lj}t_l=t'_j$, so by the uniqueness in Thm.~7.1(b), $\psi'=\alpha\circ\psi$.

Now let $L=\{x_0=\dots=x_n=0\}\subseteq\mathbf P^m$, a linear subspace of dimension $m-n-1$ (empty if $m=n$), and let $p:\mathbf P^m-L\to\mathbf P^n$, $(x_0:\dots:x_m)\mapsto(x_0:\dots:x_n)$ be the linear projection from $L$; by 7.1 it is the morphism defined by the sections $x_0,\dots,x_n$ of $\mathcal O(1)$, which generate $\mathcal O(1)$ on $\mathbf P^m-L$, so $p^*\mathcal O(1)\cong\mathcal O(1)|_{\mathbf P^m-L}$ and $p^*x_i=x_i$. Since $\psi'^*x_i=s_i$ ($i\le n$) generate $\mathcal L$ at every point $P\in X$, for each $P$ some $x_i$ ($i\le n$) does not vanish at $\psi'(P)$, i.e.\ $\psi'(X)\subseteq\mathbf P^m-L$. Thus $p\circ\psi':X\to\mathbf P^n$ is defined, $(p\psi')^*\mathcal O(1)\cong\psi'^*\mathcal O(1)\cong\mathcal L$ and $(p\psi')^*x_i=\psi'^*x_i=s_i$. By uniqueness in 7.1(b), $p\circ\psi'=\varphi$. Altogether
\[
\varphi=p\circ\alpha\circ\psi ,
\]
a linear projection composed with an automorphism of $\mathbf P^m$ and $\psi$.\hfill$\square$

\exercise{7.3}{Let $\varphi:\mathbf P^n_k\to\mathbf P^m_k$ be a morphism. Then:
\begin{enumerate}[label=(\alph*),leftmargin=*]
\item either $\varphi(\mathbf P^n)=\mathrm{pt}$ or $m\ge n$ and $\dim\varphi(\mathbf P^n)=n$;
\item in the second case, $\varphi$ can be obtained as the composition of (1) a $d$-uple embedding $\mathbf P^n\to\mathbf P^N$ for a uniquely determined $d\ge1$, (2) a linear projection $\mathbf P^N-L\to\mathbf P^m$, and (3) an automorphism of $\mathbf P^m$. Also, $\varphi$ has finite fibres.
\end{enumerate}}
\pf
Here $k$ is any field and $\varphi$ is a $k$-morphism; for $n=0$ everything is trivial, so let $n\ge1$. Let $S=k[x_0,\dots,x_n]$. By Prop.~6.4/(6.17), $\operatorname{Pic}\mathbf P^n_k=\mathbf Z\cdot[\mathcal O(1)]$, so $\varphi^*\mathcal O(1)\cong\mathcal O(d)$ for a unique $d\in\mathbf Z$, and by Thm.~7.1(a) it is generated by the global sections $s_i=\varphi^*x_i$ ($i=0,\dots,m$). Since $\Gamma(\mathbf P^n,\mathcal O(d))=0$ for $d<0$ (Thm.~5.13) and a sheaf generated by zero sections is zero, $d\ge0$; so $s_i=F_i\in S_d$, forms of degree $d$.

\textbf{Case $d=0$.} Then $F_i=c_i\in\Gamma(\mathbf P^n,\mathcal O)=k$, not all zero (the $s_i$ generate $\mathcal O$). The constant morphism $\mathbf P^n\to\operatorname{Spec}k\xrightarrow{(c_0:\dots:c_m)}\mathbf P^m$ has the same $(\mathcal L,s_i)$, so by uniqueness in 7.1(b) it equals $\varphi$. So $\varphi(\mathbf P^n)$ is a point.

\textbf{Case $d>0$.} The $F_i$ have no common zero on the scheme $\mathbf P^n$ (they generate $\mathcal O(d)$ everywhere).

\emph{A lemma.} Let $K$ be a field, $Z\subseteq\mathbf P^n_K$ a nonempty closed subset of dimension $\ge1$, and $F\in K[x_0,\dots,x_n]$ homogeneous of positive degree. Then $Z\cap V_+(F)\neq\emptyset$. Indeed, replace $Z$ by an irreducible component $Z_1$ of dimension $\ge1$; $Z_1=V_+(\mathfrak p)$ for a homogeneous prime $\mathfrak p\not\supseteq S_+$, $B=S_K/\mathfrak p$ is a finitely generated graded $K$-domain with $\dim B=\dim Z_1+1\ge2$, and $\mathfrak m_+=B_+$ is a maximal ideal (with residue field $K$), hence of height $\dim B\ge2$ (finitely generated domains over a field are catenary and equidimensional). If $\bar F=0$ in $B$, then $V_+(F)\supseteq Z_1$. Otherwise suppose $Z_1\cap V_+(F)=\emptyset$; then the only homogeneous prime of $B$ containing $\bar F$ is $\mathfrak m_+$ (a homogeneous prime not containing $B_+$ would be a point of $Z_1\cap V_+(F)$). A minimal prime of $(\bar F)$ is homogeneous, so it equals $\mathfrak m_+$, and Krull's Hauptidealsatz gives $\operatorname{ht}\mathfrak m_+\le1$, a contradiction.

\emph{Finite fibres.} Let $Q\in\mathbf P^m$ be any point, $K=\kappa(Q)$, and choose $j$ with $Q\in D_+(x_j)$. The fibre $X_Q=\mathbf P^n_k\times_{\mathbf P^m_k}\operatorname{Spec}K$ is homeomorphic to $\varphi^{-1}(Q)$ (Ex.~3.10), and, after base change to $K$, is a closed subscheme of $\mathbf P^n_K$ (the $K$-point $Q'$ of $\mathbf P^m_K$ is a closed immersion since $\mathbf P^m_K$ is separated) contained in $\varphi_K^{-1}(D_+(x_j))=D_+(F_j)$. If $X_Q\ne\emptyset$ then $F_j\neq0$, and $X_Q\cap V_+(F_j)=\emptyset$, so by the Lemma $\dim X_Q=0$; a noetherian space of dimension $0$ is finite. Thus all fibres of $\varphi$ are finite.

\emph{The image.} $\mathbf P^n_k$ is proper over $k$ and $\mathbf P^m_k$ is separated, so $\varphi$ is proper (4.8(e)) and closed. Let $\eta$ be the generic point of $\mathbf P^n$, $\xi=\varphi(\eta)$. Then $\varphi(\mathbf P^n)=\overline{\{\xi\}}=:Z$ is an irreducible closed subset with $\dim Z=\operatorname{trdeg}_k\kappa(\xi)$. The fibre over $\xi$ is a finite-type $\kappa(\xi)$-scheme with finitely many points, hence of dimension $0$, and $\eta$ is one of its points, so $\kappa(\eta)=K(\mathbf P^n)$ is a finite extension of $\kappa(\xi)$. Therefore $\dim Z=\operatorname{trdeg}_kK(\mathbf P^n)=n$, and $n=\dim Z\le m$. This proves (a), and the finiteness of the fibres.

\textbf{(b)} We are in the case $d>0$; $d$ is unique since $\varphi^*\mathcal O(1)\cong\mathcal O(d)$ determines it. Let $M_0,\dots,M_N$ be the monomials of degree $d$ ($N+1=\binom{n+d}{d}$) and $\rho_d:\mathbf P^n\to\mathbf P^N$ the $d$-uple embedding, i.e.\ the morphism defined by the sections $M_l$ of $\mathcal O(d)$ (7.1); it is a closed immersion (Ex.~5.13). Write $F_i=\sum_lc_{il}M_l$ and let $\ell_i=\sum_lc_{il}y_l$ ($i=0,\dots,m$) be linear forms on $\mathbf P^N$. Then $\rho_d^*\ell_i=F_i$. Let $L'=V_+(\ell_0,\dots,\ell_m)\subseteq\mathbf P^N$; since the $F_i$ have no common zero, $\rho_d(\mathbf P^n)\cap L'=\emptyset$. The morphism $\lambda=(\ell_0:\dots:\ell_m):\mathbf P^N-L'\to\mathbf P^m$ satisfies $(\lambda\rho_d)^*\mathcal O(1)\cong\mathcal O(d)$, $(\lambda\rho_d)^*x_i=F_i$, so by uniqueness
\[
\varphi=\lambda\circ\rho_d .
\]
Now assume (as is implicit in the statement) that the forms $F_0,\dots,F_m$ are linearly independent, i.e.\ that $\varphi(\mathbf P^n)$ is not contained in a hyperplane. Then $\ell_0,\dots,\ell_m$ are linearly independent, $m\le N$, $L'$ is a linear subspace of dimension $N-m-1$, and $\lambda$ is the linear projection from $L'$ (a projection from $L'$ is defined by a basis of the linear forms vanishing on $L'$; two choices differ by an element of $PGL_{m+1}$, which is the automorphism in (3)). So $\varphi$ is the composite of the $d$-uple embedding, a linear projection and an automorphism of $\mathbf P^m$, as claimed.

\emph{Caveat (the statement needs the independence assumption).} If the $F_i$ are linearly dependent, the printed claim is false: for $\varphi:\mathbf P^1\to\mathbf P^2$, $(x_0:x_1)\mapsto(x_0^2:x_1^2:x_0^2)$, we have $d=2$, $N=m=2$, so $L=\emptyset$ and $(2)$--$(3)$ would make $\varphi$ an automorphism of $\mathbf P^2$ composed with $\rho_2$, whose image is a smooth conic, whereas $\varphi(\mathbf P^1)$ is the line $\{y_0=y_2\}$. In general one proves: let $r+1=\dim\operatorname{span}(F_i)$; then $\varphi=\iota\circ(\text{linear projection }\mathbf P^N-L\to\mathbf P^r)\circ\rho_d$ with $\iota:\mathbf P^r\hookrightarrow\mathbf P^m$ a linear embedding (choose a basis of the span, and re-express the $F_i$ in it); the previous argument applies verbatim to $r$ in place of $m$.\hfill$\square$

\exercise{7.4}{\begin{enumerate}[label=(\alph*),leftmargin=*]
\item Use (7.6) to show that if $X$ is a scheme of finite type over a noetherian ring $A$, and if $X$ admits an ample invertible sheaf, then $X$ is separated.
\item Let $X$ be the affine line over a field $k$ with the origin doubled (4.0.1). Calculate $\operatorname{Pic}X$, determine which invertible sheaves are generated by global sections, and then show directly (without using (a)) that there is no ample invertible sheaf on $X$.
\end{enumerate}}
\pf
(a) Let $\mathcal L$ be ample on $X$. By Thm.~7.6, $\mathcal L^m$ is very ample relative to $\operatorname{Spec}A$ for some $m>0$, i.e.\ there is an immersion $i:X\to\mathbf P^r_A$ over $A$. An immersion is a composite of a closed immersion and an open immersion, hence separated (Cor.~4.6(a),(b)), and $\mathbf P^r_A\to\operatorname{Spec}A$ is separated (it is projective, hence proper, Thm.~4.9). A composite of separated morphisms is separated, so $X\to\operatorname{Spec}A$ is separated.

(b) Write $X=U_1\cup U_2$ with $U_i\cong\operatorname{Spec}k[t]$ glued along $U_{12}=U_1\cap U_2=\operatorname{Spec}k[t,t^{-1}]$ by the identity.

\emph{Pic.} Since $k[t]$ is a PID, every invertible sheaf $\mathcal L$ is trivial on $U_1$ and $U_2$; choose trivializations $\alpha_i:\mathcal L|_{U_i}\cong\mathcal O_{U_i}$. Then $\alpha_1\alpha_2^{-1}$ is multiplication by a unit of $k[t,t^{-1}]$, i.e.\ by $ct^n$, $c\in k^*$, $n\in\mathbf Z$. Replacing $\alpha_2$ by $c\alpha_2$ we may take $c=1$. So $\mathcal L\cong\mathcal L_n$, where $\mathcal L_n$ is obtained by gluing $\mathcal O_{U_1}$ and $\mathcal O_{U_2}$ so that a section over $X$ is a pair $(f_1,f_2)\in k[t]\times k[t]$ with $f_1=t^nf_2$ in $k[t,t^{-1}]$. If $\mathcal L_n\cong\mathcal L_{n'}$, then the two trivializations on each $U_i$ differ by units of $k[t]$, i.e.\ by constants $c_1,c_2\in k^*$, so $t^n=(c_1/c_2)t^{n'}$ and $n=n'$. Clearly $\mathcal L_n\otimes\mathcal L_{n'}\cong\mathcal L_{n+n'}$. Hence $\operatorname{Pic}X\cong\mathbf Z$, generated by $\mathcal L_1$.

\emph{Global sections.} For $n\ge0$ the sections of $\mathcal L_n$ are $(t^nf,f)$, $f\in k[t]$; for $n<0$ they are $(f,t^{|n|}f)$, $f\in k[t]$. In both cases $\Gamma(X,\mathcal L_n)$ is generated by one section $s_n$, namely $s_n=(t^n,1)$ for $n\ge0$, $s_n=(1,t^{|n|})$ for $n<0$. For $n>0$ every global section has first component divisible by $t$, so vanishes at the origin $0_1\in U_1$; for $n<0$ every global section vanishes at the origin $0_2\in U_2$. Thus $\mathcal L_n$ is generated by global sections only for $n=0$: \emph{the only invertible sheaf generated by its global sections is (up to isomorphism) $\mathcal O_X$.}

\emph{No ample sheaf.} Suppose $\mathcal L$ is ample. Applying the definition to the coherent sheaf $\mathcal O_X$, $\mathcal L^n$ is generated by global sections for all $n\ge n_0$, so $\mathcal L^{n_0}\cong\mathcal O_X\cong\mathcal L^{n_0+1}$, whence $\mathcal L\cong\mathcal O_X$. Then for $\mathcal F=\mathcal L_1$ (coherent) the sheaf $\mathcal F\otimes\mathcal L^n\cong\mathcal L_1$ is never generated by global sections, contradicting ampleness. So $X$ has no ample invertible sheaf (consistently with (a), since $X$ is not separated).\hfill$\square$

\exercise{7.5}{Establish the following properties of ample and very ample invertible sheaves on a noetherian scheme $X$. $\mathcal L,\mathcal M$ will denote invertible sheaves, and for (d), (e) we assume furthermore that $X$ is of finite type over a noetherian ring $A$.
\begin{enumerate}[label=(\alph*),leftmargin=*]
\item If $\mathcal L$ is ample and $\mathcal M$ is generated by global sections, then $\mathcal L\otimes\mathcal M$ is ample.
\item If $\mathcal L$ is ample and $\mathcal M$ is arbitrary, then $\mathcal M\otimes\mathcal L^n$ is ample for sufficiently large $n$.
\item If $\mathcal L,\mathcal M$ are both ample, so is $\mathcal L\otimes\mathcal M$.
\item If $\mathcal L$ is very ample and $\mathcal M$ is generated by global sections, then $\mathcal L\otimes\mathcal M$ is very ample.
\item If $\mathcal L$ is ample, then there is an $n_0>0$ such that $\mathcal L^n$ is very ample for all $n\ge n_0$.
\end{enumerate}}
\pf
I use: if $\mathcal F_1,\mathcal F_2$ are generated by global sections, so is $\mathcal F_1\otimes\mathcal F_2$ (if $s_i$ generate $\mathcal F_1$ and $t_j$ generate $\mathcal F_2$, the $s_i\otimes t_j$ generate the stalks of $\mathcal F_1\otimes\mathcal F_2$).

(a) Let $\mathcal F$ be coherent. There is $n_0$ with $\mathcal F\otimes\mathcal L^n$ generated by global sections for $n\ge n_0$. Then $\mathcal F\otimes(\mathcal L\otimes\mathcal M)^n=(\mathcal F\otimes\mathcal L^n)\otimes\mathcal M^n$ is generated by global sections for $n\ge n_0$, as $\mathcal M^n$ is.

(b) Since $\mathcal M$ is coherent and $\mathcal L$ ample, $\mathcal M\otimes\mathcal L^{j}$ is generated by global sections for $j\ge j_0$. For $n\ge j_0+1$ write $\mathcal M\otimes\mathcal L^n=\mathcal L\otimes(\mathcal M\otimes\mathcal L^{n-1})$ and apply (a).

(c) Apply the definition of ample to $\mathcal O_X$: $\mathcal M^{m}$ is generated by global sections for some $m>0$. By Prop.~7.5, $\mathcal L^m$ is ample, so by (a) $\mathcal L^m\otimes\mathcal M^m=(\mathcal L\otimes\mathcal M)^m$ is ample, and by Prop.~7.5 again $\mathcal L\otimes\mathcal M$ is ample.

(d) Let $i:X\to\mathbf P^r_A$ be an immersion with $\mathcal L\cong i^*\mathcal O(1)$, $s_a=i^*x_a$. Since $X$ is quasi-compact, $\mathcal M$ is generated by finitely many global sections $t_0,\dots,t_q$; let $\psi:X\to\mathbf P^q_A$ be the corresponding morphism (7.1), so $\psi^*\mathcal O(1)\cong\mathcal M$, $\psi^*y_b=t_b$. Consider $(i,\psi):X\to\mathbf P^r_A\times_A\mathbf P^q_A$. It is an immersion: it is the composite of the graph $X\to X\times_A\mathbf P^q_A$ of $\psi$, a closed immersion because $\mathbf P^q_A$ is separated over $A$, and $i\times\mathrm{id}:X\times_A\mathbf P^q_A\to\mathbf P^r_A\times_A\mathbf P^q_A$, an immersion as a base change of $i$. Let $\sigma:\mathbf P^r_A\times_A\mathbf P^q_A\to\mathbf P^{(r+1)(q+1)-1}_A$ be the Segre closed immersion (Ex.~5.11), $\sigma^*w_{ab}=x_ay_b$, $\sigma^*\mathcal O(1)=p_1^*\mathcal O(1)\otimes p_2^*\mathcal O(1)$. Then $\sigma\circ(i,\psi)$ is an immersion with $(\sigma(i,\psi))^*\mathcal O(1)\cong\mathcal L\otimes\mathcal M$, so $\mathcal L\otimes\mathcal M$ is very ample.

(e) By Thm.~7.6 there is $m_1>0$ with $\mathcal L^{m_1}$ very ample over $\operatorname{Spec}A$; $X$ is noetherian. Since $\mathcal L$ is ample, there is $j_0$ with $\mathcal L^j$ generated by global sections for all $j\ge j_0$. Let $n\ge n_0:=j_0+2m_1$. Choose $r'$ with $j_0\le r'\le j_0+m_1-1$ and $r'\equiv n\pmod{m_1}$, and write $n=qm_1+r'$; then $q\ge1$. The sheaf $\mathcal L^{m_1}$ is very ample, hence generated by global sections (7.1(a)); by (d) and induction, $(\mathcal L^{m_1})^q$ is very ample for all $q\ge1$. As $\mathcal L^{r'}$ is generated by global sections, (d) shows that $\mathcal L^n=(\mathcal L^{m_1})^q\otimes\mathcal L^{r'}$ is very ample.\hfill$\square$

\exercise{7.6}{\emph{The Riemann--Roch Problem.} Let $X$ be a nonsingular projective variety over an algebraically closed field, and let $D$ be a divisor on $X$. For any $n>0$ we consider the complete linear system $|nD|$. Then the Riemann--Roch problem is to determine $\dim|nD|$ as a function of $n$, and, in particular, its behavior for large $n$. If $\mathcal L$ is the corresponding invertible sheaf, then $\dim|nD|=\dim\Gamma(X,\mathcal L^n)-1$, so an equivalent problem is to determine $\dim\Gamma(X,\mathcal L^n)$ as a function of $n$.
\begin{enumerate}[label=(\alph*),leftmargin=*]
\item Show that if $D$ is very ample, and if $X\hookrightarrow\mathbf P^N$ is the corresponding embedding in projective space, then for all $n$ sufficiently large, $\dim|nD|=P_X(n)-1$, where $P_X$ is the Hilbert polynomial of $X$ (I, Section~7). Thus in this case $\dim|nD|$ is a polynomial function of $n$, for $n\gg0$.
\item If $D$ corresponds to a torsion element of $\operatorname{Pic}X$, of order $r$, then $\dim|nD|=0$ if $r\mid n$, $-1$ otherwise. In this case the function is periodic of period $r$.
\end{enumerate}
It follows from the general Riemann--Roch theorem that $\dim|nD|$ is a polynomial function in $n$ for $n$ large, whenever $D$ is an \emph{ample} divisor. See (IV, 1.3.2), (V, 1.6), and Appendix~A. In the case of algebraic surfaces, Zariski [7] has shown for any effective divisor $D$, that there is a finite set of polynomials $P_1,\dots,P_r$, such that for all $n$ sufficiently large, $\dim|nD|=P_{i(n)}(n)$, where $i(n)\in\{1,2,\dots,r\}$ is a function of $n$.}
\pf
Since $X$ is nonsingular, $\operatorname{Cl}X\cong\operatorname{Pic}X$ (6.16), and $|nD|$ is the projective space of the nonzero sections of $\mathcal L^n=\mathcal L(nD)$ modulo $k^*$ (Prop.~7.7), so $\dim|nD|=\dim_k\Gamma(X,\mathcal L^n)-1$ as in the statement.

(a) $D$ very ample means $\mathcal L\cong i^*\mathcal O(1)$ for an immersion $i:X\to\mathbf P^N$, which is a closed immersion since $X$ is projective (complete); identify $X$ with the closed subvariety $i(X)$ (integral, so with the reduced structure). Then $\mathcal L^n\cong\mathcal O_X(n)$. Let $S(X)=k[x_0,\dots,x_N]/I(X)$ be the homogeneous coordinate ring. It is generated by $S(X)_1$ over $k$, $X=\operatorname{Proj}S(X)$, and $\widetilde{S(X)}=\mathcal O_X$. By Ex.~5.9(b) (with $M=S=S(X)$), $S(X)_n\to\Gamma(X,\mathcal O_X(n))$ is an isomorphism for $n\ge n_0$. By definition of the Hilbert polynomial (I, Section~7), $\dim_kS(X)_n=P_X(n)$ for $n\gg0$. Hence for $n\gg0$,
\[
\dim|nD|=\dim_k\Gamma(X,\mathcal L^n)-1=P_X(n)-1 .
\]

(b) Let $r\ge1$ be minimal with $\mathcal L^r\cong\mathcal O_X$; then $\mathcal L^n\cong\mathcal O_X$ iff $r\mid n$ (the set of such $n$ is the subgroup of $\mathbf Z$ generated by $r$). If $r\mid n$, then $\Gamma(X,\mathcal L^n)\cong\Gamma(X,\mathcal O_X)=k$ (I, Thm.~3.4, $X$ a projective variety), so $\dim|nD|=0$. Let $r\nmid n$ and suppose, for contradiction, that $0\ne s\in\Gamma(X,\mathcal L^n)$. Since $X$ is integral, the germ of $s$ at the generic point is nonzero in the one-dimensional $K(X)$-vector space $(\mathcal L^n)_\eta$, so $s^r\in\Gamma(X,\mathcal L^{nr})\cong\Gamma(X,\mathcal O_X)=k$ is nonzero, i.e.\ a nonzero constant, hence it generates $\mathcal L^{nr}$ at every point. If $s_P\in\mathfrak m_P(\mathcal L^n)_P$ for some $P$, then $s^r_P\in\mathfrak m_P(\mathcal L^{nr})_P$, a contradiction; so $s$ generates $\mathcal L^n$ everywhere, i.e.\ $\mathcal O_X\to\mathcal L^n$, $1\mapsto s$, is surjective, hence an isomorphism by Ex.~7.1. Then $r\mid n$, a contradiction. Hence $\Gamma(X,\mathcal L^n)=0$ and $|nD|=\emptyset$ has dimension $-1$ for $r\nmid n$. The function is periodic of period $r$.\hfill$\square$

\exercise{7.7}{\emph{Some Rational Surfaces.} Let $X=\mathbf P^2_k$, and let $|D|$ be the complete linear system of all divisors of degree $2$ on $X$ (conics). $D$ corresponds to the invertible sheaf $\mathcal O(2)$, whose space of global sections has a basis $x^2,y^2,z^2,xy,xz,yz$, where $x,y,z$ are the homogeneous coordinates of $X$.
\begin{enumerate}[label=(\alph*),leftmargin=*]
\item The complete linear system $|D|$ gives an embedding of $\mathbf P^2$ in $\mathbf P^5$, whose image is the Veronese surface (I, Ex.~2.13).
\item Show that the subsystem defined by $x^2,y^2,z^2,x(y-z),(x-y)z$ gives a closed immersion of $X$ into $\mathbf P^4$. The image is called the Veronese surface in $\mathbf P^4$. Cf.\ (IV, Ex.~3.11).
\item Let $\mathfrak d\subseteq|D|$ be the linear system of all conics passing through a fixed point $P$. Then $\mathfrak d$ gives an immersion of $U=X-P$ into $\mathbf P^4$. Furthermore, if we blow up $P$, to get a surface $\tilde X$, then this map extends to a closed immersion of $\tilde X$ in $\mathbf P^4$. Show that $\tilde X$ is a surface of degree $3$ in $\mathbf P^4$, and that the lines in $X$ through $P$ are transformed into straight lines in $\tilde X$ which do not meet. $\tilde X$ is the union of all these lines, so we say $\tilde X$ is a ruled surface (V, 2.19.1).
\end{enumerate}}
\pf
Throughout $k$ is algebraically closed (as needed for linear systems; the closed-immersion verifications below work over any field except where char${}\ne2$ is flagged). Let $S=k[x,y,z]$. For a form $s$ of positive degree, $X_s=D_+(s)=\operatorname{Spec}S_{(s)}$ is affine, so condition (1) of Prop.~7.2 holds automatically for any family of forms of the same degree; only condition (2) must be checked. For a family $s_0,\dots,s_N\in S_2$ with no common zero, $\varphi$ is the morphism of 7.1, and it is a closed immersion iff for each $i$ the ring $S_{(s_i)}$ is generated as $k$-algebra by the elements $s_j/s_i$.

(a) The six monomials $s_{ab}=x_ax_b$ are the basis of $\Gamma(\mathcal O(2))$, with no common zero. The 2-uple embedding is by definition the morphism they define (7.1). For $s=x_ax_b$, $D_+(x_ax_b)=D_+(x_a)\cap D_+(x_b)$ has ring $k[x_c/x_a\ (c=0,1,2),\ x_a/x_b]$ and each of these generators is a ratio $s'/s$: $x_c/x_a=x_bx_c/(x_ax_b)$ and $x_a/x_b=x_a^2/(x_ax_b)$. So (2) holds, and $\varphi$ is a closed immersion of $\mathbf P^2$ in $\mathbf P^5$, with image the Veronese surface (the image of the $2$-uple embedding, I, Ex.~2.13).

(b) \emph{Hypothesis on the characteristic.} The statement needs $\operatorname{char}k\ne2$ (which I assume): in characteristic~2 the morphism is not an immersion at $(1:1:0)$. Indeed, in the chart $x\ne0$ with $u=y/x,v=z/x$ the five coordinates are $u^2,v^2,u-v,(1-u)v$ and the differentials in characteristic 2 are $0,0,du+dv,(1+u)dv+v\,du$, which at $(u,v)=(1,0)$ span only a line.

So let $\operatorname{char}k\ne2$ and $s_0=x^2,s_1=y^2,s_2=z^2,s_3=x(y-z),s_4=(x-y)z$. These have no common zero ($x^2=y^2=z^2=0$ is impossible). For each $i$ let $R_i\subseteq S_{(s_i)}$ be the $k$-subalgebra generated by the $s_j/s_i$; we show $R_i=S_{(s_i)}$.

\emph{$i=0$:} $S_{(s_0)}=k[u,v]$ with $u=y/x$, $v=z/x$. $R_0$ contains $u^2,v^2,w=u-v,b=(1-u)v$. From $w^2=u^2-2uv+v^2$ we get $uv\in R_0$ ($2$ is invertible), so $v=b+uv\in R_0$ and $u=w+v\in R_0$.

\emph{$i=1$:} $S_{(s_1)}=k[p,q]$, $p=x/y,q=z/y$; $R_1\ni p^2,q^2,p(1-q),(p-1)q$. The sum of the last two is $p-q$; $(p-q)^2$ gives $pq\in R_1$; then $p=p(1-q)+pq$ and $q=pq-(p-1)q$ lie in $R_1$.

\emph{$i=2$:} $S_{(s_2)}=k[a,b]$, $a=x/z$, $b=y/z$; $R_2\ni a^2,b^2,a(b-1),a-b$. From $(a-b)^2$ we get $ab\in R_2$, then $a=ab-a(b-1)$ and $b=a-(a-b)$.

\emph{$i=3$:} Put $X'=x$, $T=y-z$, $Z=z$ (a linear change of coordinates), so $s_3=X'T$ and $S_{(s_3)}=k[\alpha,\beta,\gamma]/(\beta\gamma-1)$ with $\alpha=Z/X'$, $\beta=T/X'$, $\gamma=\beta^{-1}$. Then
\[
\tfrac{s_0}{s_3}=\gamma,\quad \tfrac{s_2}{s_3}=\alpha^2\gamma,\quad \tfrac{s_1}{s_3}=\tfrac{(T+Z)^2}{X'T}=\beta+2\alpha+\alpha^2\gamma,\quad \tfrac{s_4}{s_3}=\tfrac{(X'-T-Z)Z}{X'T}=\alpha\gamma-\alpha-\alpha^2\gamma .
\]
So $R_3\ni\gamma,\alpha^2\gamma$, hence $m=\beta+2\alpha\in R_3$; then $m\gamma=1+2\alpha\gamma$ gives $\alpha\gamma\in R_3$, then $\alpha=\alpha\gamma-\alpha^2\gamma-s_4/s_3\in R_3$ and $\beta=m-2\alpha\in R_3$.

\emph{$i=4$:} Put $r=x-y$, so $s_4=rz$ and $S_{(s_4)}=k[\alpha,\beta,\gamma]/(\beta\gamma-1)$ with $\alpha=y/z$, $\beta=r/z$, $\gamma=\beta^{-1}$, $x/z=\alpha+\beta$. Then
\[
\tfrac{s_2}{s_4}=\gamma,\quad\tfrac{s_1}{s_4}=\alpha^2\gamma,\quad\tfrac{s_0}{s_4}=(\alpha+\beta)^2\gamma=\alpha^2\gamma+2\alpha+\beta,\quad \tfrac{s_3}{s_4}=(\alpha+\beta)(\alpha-1)\gamma=\alpha^2\gamma-\alpha\gamma+\alpha-1 .
\]
So $m=2\alpha+\beta\in R_4$, $m\gamma=2\alpha\gamma+1$ gives $\alpha\gamma\in R_4$, then $\alpha\in R_4$ (from $s_3/s_4$) and $\beta=m-2\alpha\in R_4$.

By Prop.~7.2, $\varphi:X\to\mathbf P^4$ is a closed immersion; its image is the Veronese surface in $\mathbf P^4$ (the projection of the Veronese surface of (a) from a point not on its secant variety).

(c) Choose coordinates with $P=(0:0:1)$. The conics through $P$ are the forms in the $5$-dimensional space $V=\langle x^2,xy,y^2,xz,yz\rangle$ ($=\Gamma(X,\mathcal I_P(2))$), so $\dim\mathfrak d=4$.

\emph{The blow-up.} Let $B\subseteq\mathbf P^2\times\mathbf P^1$ be the closed subscheme $x\eta=y\xi$, with $(\xi:\eta)$ homogeneous coordinates on $\mathbf P^1$, and $\pi:B\to\mathbf P^2$ the projection. I claim $B\to X$ is the blow-up of $P$ (ideal $\mathcal I_P$). Over $D_+(z)=\operatorname{Spec}k[a,b]$ ($a=x/z$, $b=y/z$) with $\mathcal I_P=(a,b)$, the blow-up is the closed subscheme $a\eta=b\xi$ of $\mathbf P^1\times D_+(z)$ (Example~7.12.1), which is $B|_{D_+(z)}$. Over $X-P$ the blow-up is the identity (7.13(b)), and so is $B\to X$ there (a point of $B$ over $(x:y:z)$ with $(x,y)\ne0$ has $(\xi:\eta)=(x:y)$). On the overlap $D_+(z)-P$ both identifications are the unique $X$-isomorphism (both schemes map isomorphically to $D_+(z)-P$). So they glue: $\tilde X\cong B$ over $X$.

\emph{The map.} On $B$ consider the five sections $w_0=x\xi,\ w_1=x\eta(=y\xi),\ w_2=y\eta,\ w_3=z\xi,\ w_4=z\eta$ of $\mathcal O(1,1)|_B$. They have no common zero on $B$: if $z\ne0$ then $\xi=\eta=0$ is impossible; so $z=0$; then if $x\ne0$ again $\xi=\eta=0$; so $x=0$, then $y\ne0$, and $w_2=w_1=0$ force $\eta=0=\xi$, impossible. Let $\Phi:B\to\mathbf P^4$ be the corresponding morphism (7.1). On $B-\pi^{-1}(P)\cong U$ we have $(\xi:\eta)=(x:y)$, so $\mathcal O(1,1)|_U\cong\mathcal O(2)|_U$ with $\xi\leftrightarrow x,\eta\leftrightarrow y$, and the $w_i$ become $x^2,xy,y^2,xz,yz$: thus $\Phi|_U$ is the morphism defined by $\mathfrak d$ on $U$, and $\Phi$ extends it to $B=\tilde X$.

\emph{$\Phi$ is a closed immersion (Prop.~7.2).} The sets $B_{w_i}$ are affine and the charts are as follows (coordinates in each chart; all $B_{w_i}$ are open in $B$):
\begin{itemize}
\item $w_0\ne0$: $x\ne0,\xi\ne0$; set $x=\xi=1$, $u=y,v=z$, then $\eta=u$: $B_{w_0}=\operatorname{Spec}k[u,v]$; the ratios $w_j/w_0$ are $u,u^2,v,uv$, which generate $k[u,v]$.
\item $w_2\ne0$: symmetric: $B_{w_2}=\operatorname{Spec}k[p,q]$ ($p=x/y,q=z/y$), ratios $p^2,p,pq,q$.
\item $w_1\ne0$: $x\ne0,\eta\ne0$; with $x=1$, $u=y,v=z$, $\eta=u\xi$, so $u$ must be invertible: $B_{w_1}=\operatorname{Spec}k[u,u^{-1},v]$; the ratios $w_j/w_1$ are $u^{-1},u,v/u,v$, which generate.
\item $w_3\ne0$: $z\ne0,\xi\ne0$; set $z=\xi=1$, $a=x$, $\eta$ free, and then $y=a\eta$: $B_{w_3}=\operatorname{Spec}k[a,\eta]$; the ratios are $w_0/w_3=a$, $w_1/w_3=a\eta$, $w_2/w_3=a\eta^2$, $w_4/w_3=\eta$, which generate.
\item $w_4\ne0$: $z\ne0,\eta\ne0$; set $z=\eta=1$, $b=y$, $\xi$ free, and then $x=b\xi$: $B_{w_4}=\operatorname{Spec}k[b,\xi]$; the ratios are $w_0/w_4=b\xi^2$, $w_1/w_4=b\xi$, $w_2/w_4=b$, $w_3/w_4=\xi$, which generate.
\end{itemize}
So $\Phi$ is a closed immersion: $\tilde X=B$ is identified with the closed subscheme $\Phi(B)\subseteq\mathbf P^4$; in particular $\Phi|_U$ is an immersion of $U$ (an open subscheme of $\tilde X$).

\emph{Degree $3$.} $\tilde X=\Phi(B)$ is an integral surface. In the chart $B_{w_0}\cong\mathbf A^2$, $\Phi=(1:u:u^2:v:uv)$. Since $k$ is infinite and $B_{w_0}$ is dense in $B$, a form $G(w_0,\dots,w_4)$ of degree $n$ lies in $I(\tilde X)$ iff $G(1,u,u^2,v,uv)=0$ in $k[u,v]$. Hence $S(\tilde X)_n$ is isomorphic to the $k$-span of the monomials $u^av^b$ that are products of $n$ elements of $\{1,u,u^2,v,uv\}$. Such a product has $v$-exponent $b=n_3+n_4\in[0,n]$, and for fixed $b$ the $u$-exponent takes exactly the values $0,1,\dots,2(n-b)+b$ (the $n-b$ factors from $\{1,u,u^2\}$ contribute any value in $[0,2(n-b)]$, the $n_4\in[0,b]$ factors $uv$ contribute $n_4$). So the Hilbert function is
\[
\sum_{b=0}^{n}(2n-b+1)=\frac{(n+1)(3n+2)}2=\frac32n^2+\frac52n+1 ,
\]
the Hilbert polynomial of $\tilde X$. It has degree $2$ and leading coefficient $3/2$, so $\dim\tilde X=2$ and $\deg\tilde X=2!\cdot\frac32=3$. (Check: $\tilde X$ lies on the three $2\times2$ minors $w_0w_2-w_1^2$, $w_0w_4-w_1w_3$, $w_1w_4-w_2w_3$ of $\bigl(\begin{smallmatrix}w_0&w_1&w_3\\ w_1&w_2&w_4\end{smallmatrix}\bigr)$, the cubic scroll.)

\emph{The lines.} The lines of $X$ through $P$ are $\ell_{(\xi_0:\eta_0)}=\{x\eta_0=y\xi_0\}$, $(\xi_0:\eta_0)\in\mathbf P^1$. The strict transform of $\ell_{(\xi_0:\eta_0)}$ is the closure of $\pi^{-1}(\ell-P)$, namely the fibre $F_{(\xi_0:\eta_0)}$ of $B\to\mathbf P^1$ over $(\xi_0:\eta_0)$, i.e.\ $\{((\lambda\xi_0:\lambda\eta_0:z),(\xi_0:\eta_0))\}\cong\mathbf P^1$ with coordinates $(\lambda:z)$. Its image under $\Phi$ is
\[
(\lambda\xi_0^2:\lambda\xi_0\eta_0:\lambda\eta_0^2:z\xi_0:z\eta_0)=\lambda\,(\xi_0^2,\xi_0\eta_0,\eta_0^2,0,0)+z\,(0,0,0,\xi_0,\eta_0),
\]
the straight line through the points $(\xi_0^2:\xi_0\eta_0:\eta_0^2:0:0)$ and $(0:0:0:\xi_0:\eta_0)$. Distinct fibres of $B\to\mathbf P^1$ are disjoint and $\Phi$ is injective, so these lines are pairwise disjoint. Every point of $B$ lies in exactly one fibre, so $\tilde X=\bigcup_{(\xi_0:\eta_0)}\Phi(F_{(\xi_0:\eta_0)})$ is the union of these lines: a ruled surface.\hfill$\square$

\exercise{7.8}{Let $X$ be a noetherian scheme, let $\mathcal E$ be a coherent locally free sheaf on $X$, and let $\pi:\mathbf P(\mathcal E)\to X$ be the corresponding projective space bundle. Show that there is a natural $1$-$1$ correspondence between sections of $\pi$ (i.e., morphisms $\sigma:X\to\mathbf P(\mathcal E)$ such that $\pi\circ\sigma=\mathrm{id}_X$) and quotient invertible sheaves $\mathcal E\to\mathcal L\to0$ of $\mathcal E$.}
\pf
Apply Prop.~7.12 with $Y=X$ and $g=\mathrm{id}_X$: there $\mathbf P(\mathcal E)$ comes with the surjection $\pi^*\mathcal E\to\mathcal O(1)$ (7.11(b)), and a morphism $f:X\to\mathbf P(\mathcal E)$ over $X$ gives the pair $(f^*\mathcal O(1),\ f^*(\pi^*\mathcal E\to\mathcal O(1)))$ with $f^*\pi^*\mathcal E=(\pi f)^*\mathcal E=\mathcal E$. A morphism over $X$ from $X$ is exactly a section $\sigma$ of $\pi$ (the condition $\pi\sigma=\mathrm{id}$ \emph{is} being over $X$).

Define $\Theta(\sigma)=\bigl(\mathcal E\to\sigma^*\mathcal O(1)\bigr)$; it is a surjection onto an invertible sheaf, i.e.\ a quotient invertible sheaf. Two quotients $\mathcal E\to\mathcal L$, $\mathcal E\to\mathcal L'$ are identified if there is an isomorphism $\mathcal L\cong\mathcal L'$ compatible with the maps from $\mathcal E$ (equivalently, if they have the same kernel).

\emph{Surjective.} Given $\mathcal E\to\mathcal L\to0$ with $\mathcal L$ invertible, Prop.~7.12 (proof) gives a morphism $f:X\to\mathbf P(\mathcal E)$ over $X$ with $\mathcal L\cong f^*\mathcal O(1)$ such that the given surjection is $f^*$ of $\pi^*\mathcal E\to\mathcal O(1)$; being over $X$ it is a section, and $\Theta(f)$ is the given quotient.

\emph{Injective.} The uniqueness statement in the proof of 7.12 says that for a given pair $(\mathcal L,\ g^*\mathcal E\to\mathcal L)$ there is a unique $f$ with $\mathcal L\cong f^*\mathcal O(1)$ and with the map obtained from $\pi^*\mathcal E\to\mathcal O(1)$ by $f^*$ (it is deduced from the uniqueness in 7.1 after localizing so that $\mathcal E$ is free). So if $\Theta(\sigma)$ and $\Theta(\sigma')$ are isomorphic quotients, then $\sigma=\sigma'$.

Thus $\sigma\mapsto\Theta(\sigma)$ is a bijection between sections of $\pi$ and isomorphism classes of quotient invertible sheaves of $\mathcal E$; its inverse sends $\mathcal E\to\mathcal L$ to the section determined by $\mathcal L$ and the images of local bases (locally, if $\mathcal E=\mathcal O^{n+1}$, the section $(\mathcal L;\ s_0,\dots,s_n)$ of 7.1). Naturality: it is compatible with restriction to open subsets of $X$ and with isomorphisms of $\mathcal E$.\hfill$\square$

\exercise{7.9}{Let $X$ be a regular noetherian scheme, and $\mathcal E$ a locally free coherent sheaf of rank $\ge2$ on $X$.
\begin{enumerate}[label=(\alph*),leftmargin=*]
\item Show that $\operatorname{Pic}\mathbf P(\mathcal E)\cong\operatorname{Pic}X\times\mathbf Z$.
\item If $\mathcal E'$ is another locally free coherent sheaf on $X$, show that $\mathbf P(\mathcal E)\cong\mathbf P(\mathcal E')$ (over $X$) if and only if there is an invertible sheaf $\mathcal L$ on $X$ such that $\mathcal E'\cong\mathcal E\otimes\mathcal L$.
\end{enumerate}}
\pf
\emph{Conventions.} A regular noetherian scheme is a finite disjoint union of connected, hence integral, regular schemes, so as usual I assume $X$ connected (hence integral); for disconnected $X$ the statement (a) holds componentwise, i.e.\ with $\mathbf Z$ replaced by the group of locally constant functions $X\to\mathbf Z$ (e.g.\ for $X$ two points, $\operatorname{Pic}X=0$ but $\operatorname{Pic}P=\mathbf Z^2$, so connectedness is really needed in (a)). Part (b) needs no such assumption: $\mathbf P(\mathcal E)$ and $\mathbf P(\mathcal E')$ are compatible with the decomposition of $X$ into its finitely many connected components, and invertible sheaves on $X$ are glued from invertible sheaves on the components. Let $\operatorname{rank}\mathcal E=n+1$, $n\ge1$, $P=\mathbf P(\mathcal E)$, $\pi:P\to X$ the projection. Recall $\pi_*\mathcal O_P(l)=S^l(\mathcal E)$ for $l\ge0$ (Prop.~7.11(a)), in particular $\pi_*\mathcal O_P=\mathcal O_X$ and $\pi_*\mathcal O_P(1)=\mathcal E$, and $\pi$ is proper (7.10(a)).

\textbf{Lemma A (local structure of line bundles on $\mathbf P^n$-bundles).} Let $\pi:P\to X$ be a morphism such that $X$ is covered by opens $U$ with $\pi^{-1}(U)\cong\mathbf P^n_U$ ($n\ge1$), $X$ regular noetherian, and let $\mathcal N\in\operatorname{Pic}P$. For $x\in X$ the fibre $P_x\cong\mathbf P^n_{\kappa(x)}$ carries $\mathcal N|_{P_x}\cong\mathcal O(d(x))$ for a unique $d(x)\in\mathbf Z$ (independent of the identification of $P_x$ with $\mathbf P^n_{\kappa(x)}$, since an automorphism of $\mathbf P^n_\kappa$ pulls $\mathcal O(1)$ back to a generator of $\operatorname{Pic}=\mathbf Z$ having sections, i.e.\ to $\mathcal O(1)$). Then every $x\in X$ has an open neighbourhood $U$ with a trivialization $\psi:\pi^{-1}(U)\cong\mathbf P^n_U$ such that $\mathcal N|_{\pi^{-1}U}\cong\psi^*\mathcal O(d(x))$. In particular $d:X\to\mathbf Z$ is locally constant.

\emph{Proof.} Let $R=\mathcal O_{X,x}$, a regular local ring, hence a UFD (6.11A / Auslander--Buchsbaum). The scheme $\mathbf P^n_R$ is noetherian, integral, separated, and regular (it is covered by $\operatorname{Spec}R[y_1,\dots,y_n]$, and polynomial rings over regular rings are regular), hence locally factorial, so $\operatorname{Pic}\mathbf P^n_R=\operatorname{Cl}\mathbf P^n_R$ (6.16). The hyperplane $H=V_+(x_0)$ is a prime divisor with $\mathbf P^n_R-H=\mathbf A^n_R=\operatorname{Spec}R[y_1,\dots,y_n]$, and $R[y_1,\dots,y_n]$ is a UFD, so $\operatorname{Cl}\mathbf A^n_R=0$ (6.2). By (6.5), $\mathbf Z\to\operatorname{Cl}\mathbf P^n_R\to\operatorname{Cl}\mathbf A^n_R=0$ is exact, so $\operatorname{Cl}\mathbf P^n_R$ is generated by $[H]$, i.e.\ by $\mathcal O(1)$. Hence $\mathcal N|_{P_R}\cong\mathcal O(d)$ for some $d$ (here $P_R=P\times_X\operatorname{Spec}R$, trivialized over an affine neighbourhood $U_0=\operatorname{Spec}A$ of $x$), and restricting to the closed fibre gives $d=d(x)$.

The isomorphism is a nowhere vanishing section $\sigma$ of the invertible sheaf $\mathcal G=\mathcal N\otimes\mathcal O(-d)$ over $P_R=\mathbf P^n_R$. Since $\Gamma(\mathbf P^n_R,\mathcal G)=\Gamma(\mathbf P^n_A,\mathcal G)\otimes_AR$ (flat base change; $\mathbf P^n_A$ is quasi-compact and separated), $\sigma$ extends to a section $\sigma_U$ over $\pi^{-1}(U)$ for some smaller affine neighbourhood $U=D(a)$ of $x$. The zero set $Z$ of $\sigma_U$ is closed in $\pi^{-1}(U)$ and misses $\pi^{-1}(x)$; as $\pi$ is proper (locally on $X$ it is a projection $\mathbf P^n_U\to U$, and properness is local on the base, 4.8(f)), $\pi(Z)$ is closed in $U$ and does not contain $x$. Replacing $U$ by $U-\pi(Z)$, $\sigma_U$ is nowhere vanishing, so $\mathcal O\to\mathcal G|_{\pi^{-1}U}$, $1\mapsto\sigma_U$, is a surjection of invertible sheaves, hence an isomorphism (Ex.~7.1), i.e.\ $\mathcal N\cong\mathcal O(d)$ there. Then for $x'\in U$, $d(x')=d$.\hfill$\square$

\textbf{(a)} Define $\Psi:\operatorname{Pic}X\times\mathbf Z\to\operatorname{Pic}P$, $(\mathcal L,d)\mapsto\pi^*\mathcal L\otimes\mathcal O_P(d)$, a homomorphism.

\emph{Injective.} Suppose $\pi^*\mathcal L\otimes\mathcal O_P(d)\cong\mathcal O_P$. Restrict to a fibre $P_x\cong\mathbf P^n_{\kappa(x)}$: $\pi^*\mathcal L$ is trivial there, so $\mathcal O(d)\cong\mathcal O$ on $\mathbf P^n_\kappa$, whence $d=0$ ($n\ge1$). Then $\pi^*\mathcal L\cong\mathcal O_P$ and applying $\pi_*$ and the projection formula, $\mathcal L\cong\mathcal L\otimes\pi_*\mathcal O_P=\pi_*\pi^*\mathcal L\cong\pi_*\mathcal O_P=\mathcal O_X$.

\emph{Surjective.} Let $\mathcal N\in\operatorname{Pic}P$. By Lemma A (applied to the local trivializations coming from free neighbourhoods of $\mathcal E$), $d(x)$ is locally constant, hence equal to a constant $d$ on the connected space $X$, and each point has a neighbourhood $U$ with $\mathcal N(-d)|_{\pi^{-1}U}\cong\mathcal O$. Put $\mathcal L=\pi_*(\mathcal N(-d))$. On such a $U$, $\mathcal L|_U\cong\pi_*\mathcal O=\mathcal O_U$, so $\mathcal L$ is invertible, and the natural map $\pi^*\pi_*\mathcal N(-d)\to\mathcal N(-d)$ is an isomorphism over $U$ (it is the identity of $\mathcal O$ after the trivialization). Hence $\mathcal N\cong\pi^*\mathcal L\otimes\mathcal O_P(d)=\Psi(\mathcal L,d)$. So $\Psi$ is an isomorphism.

\textbf{(b)} ($\Leftarrow$) Let $\mathcal E'\cong\mathcal E\otimes\mathcal L$. Put $\mathcal S=S(\mathcal E)$ and consider $\mathcal S*\mathcal L=\bigoplus_d\mathcal S_d\otimes\mathcal L^d$ as in the Lemma numbered 7.9 of the text. The natural map $\mathcal E\otimes\mathcal L\to(\mathcal S*\mathcal L)_1$ induces an algebra homomorphism $S(\mathcal E\otimes\mathcal L)\to\mathcal S*\mathcal L$, which is an isomorphism (check locally where $\mathcal L$ and $\mathcal E$ are free: $S^d(\mathcal E\otimes\mathcal L)\cong S^d(\mathcal E)\otimes\mathcal L^d$). The Lemma~7.9 gives $\operatorname{Proj}(\mathcal S*\mathcal L)\cong\operatorname{Proj}\mathcal S$ over $X$, i.e.\ $\mathbf P(\mathcal E\otimes\mathcal L)\cong\mathbf P(\mathcal E)$ (and $\mathcal O_{\mathbf P(\mathcal E\otimes\mathcal L)}(1)\cong\mathcal O_{\mathbf P(\mathcal E)}(1)\otimes\pi^*\mathcal L$).

($\Rightarrow$) Let $f:\mathbf P(\mathcal E')\to\mathbf P(\mathcal E)$ be an isomorphism over $X$, with projections $\pi',\pi$ ($\pi f=\pi'$). Put $\mathcal N=f^*\mathcal O_{\mathbf P(\mathcal E)}(1)\in\operatorname{Pic}\mathbf P(\mathcal E')$. By (a) applied to $\mathcal E'$ (rank $\ge2$), $\mathcal N\cong\pi'^*\mathcal L\otimes\mathcal O_{\mathbf P(\mathcal E')}(d)$ for some $\mathcal L,d$. On a fibre over $x$, $f_x$ is an isomorphism of copies of $\mathbf P^n_{\kappa(x)}$, so $f_x^*\mathcal O(1)\cong\mathcal O(1)$ (it is a generator of $\operatorname{Pic}=\mathbf Z$ with nonzero sections), whereas $\mathcal N|_{P'_x}\cong\mathcal O(d)$; thus $d=1$. Apply $\pi'_*$:
\[
\pi'_*\mathcal N=\pi_*f_*f^*\mathcal O(1)=\pi_*\mathcal O_{\mathbf P(\mathcal E)}(1)=\mathcal E,\qquad \pi'_*(\pi'^*\mathcal L\otimes\mathcal O(1))=\mathcal L\otimes\pi'_*\mathcal O(1)=\mathcal L\otimes\mathcal E'
\]
(projection formula, 7.11(a)). So $\mathcal E\cong\mathcal E'\otimes\mathcal L$, i.e.\ $\mathcal E'\cong\mathcal E\otimes\mathcal L^{-1}$.\hfill$\square$

\exercise{7.10}{\emph{$\mathbf P^n$-Bundles Over a Scheme.} Let $X$ be a noetherian scheme.
\begin{enumerate}[label=(\alph*),leftmargin=*]
\item By analogy with the definition of a vector bundle (Ex.~5.18), define the notion of a \emph{projective $n$-space bundle} over $X$, as a scheme $P$ with a morphism $\pi:P\to X$ such that $P$ is locally isomorphic to $U\times\mathbf P^n$, $U\subseteq X$ open, and the transition automorphisms on $\operatorname{Spec}A\times\mathbf P^n$ are given by $A$-linear automorphisms of the homogeneous coordinate ring $A[x_0,\dots,x_n]$ (e.g., $x'_i=\sum a_{ij}x_j$, $a_{ij}\in A$).
\item If $\mathcal E$ is a locally free sheaf of rank $n+1$ on $X$, then $\mathbf P(\mathcal E)$ is a $\mathbf P^n$-bundle over $X$.
\item Assume that $X$ is regular, and show that every $\mathbf P^n$-bundle $P$ over $X$ is isomorphic to $\mathbf P(\mathcal E)$ for some locally free $\mathcal E$ on $X$. [\emph{Hint:} Let $U\subseteq X$ be an open set such that $\pi^{-1}(U)\cong U\times\mathbf P^n$, and let $\mathcal L_0$ be the invertible sheaf $\mathcal O(1)$ on $U\times\mathbf P^n$. Show that $\mathcal L_0$ extends to an invertible sheaf $\mathcal L'$ on $P$. Then show that $\pi_*\mathcal L'=\mathcal E$ is a locally free sheaf on $X$ and that $P\cong\mathbf P(\mathcal E)$.] Can you weaken the hypothesis ``$X$ regular''?
\item Conclude (in the case $X$ regular) that we have a $1$-$1$ correspondence between $\mathbf P^n$-bundles over $X$, and equivalence classes of locally free sheaves of rank $n+1$ under the equivalence relation $\mathcal E\sim\mathcal E'$ if and only if $\mathcal E'\cong\mathcal E\otimes\mathcal M$ for some invertible sheaf $\mathcal M$ on $X$.
\end{enumerate}
(Part (c) is starred in the book.)}
\pf
\textbf{(a) Definition.} A \emph{projective $n$-space bundle} over $X$ is a scheme $P$ with a morphism $\pi:P\to X$, together with an open cover $\{U_i\}$ of $X$ and isomorphisms $\psi_i:\pi^{-1}(U_i)\to\mathbf P^n_{U_i}=U_i\times\mathbf P^n$ over $U_i$ (local trivializations), such that for all $i,j$ the transition automorphism $\psi_i\psi_j^{-1}$ of $\mathbf P^n_{U_i\cap U_j}$ over $U_i\cap U_j$ is \emph{linear}: over every affine open $\operatorname{Spec}A\subseteq U_i\cap U_j$ it is $\operatorname{Proj}$ of an $A$-linear graded automorphism of $A[x_0,\dots,x_n]$, $x_a\mapsto\sum_ba_{ab}x_b$ with $(a_{ab})\in GL_{n+1}(A)$. (This is stable under restriction to smaller affines, so it suffices to ask it on an affine cover.) An isomorphism of $\mathbf P^n$-bundles over $X$ is an isomorphism of $X$-schemes. Exactly as for vector bundles, the transition data are cocycles with values in $PGL_{n+1}$ in place of $GL_n$.

\textbf{(b)} Let $\mathcal E$ be locally free of rank $n+1$. Choose an open cover $\{U_i\}$ of $X$ and isomorphisms $\varepsilon_i:\mathcal E|_{U_i}\cong\mathcal O_{U_i}^{n+1}$. Then $S(\mathcal E)|_{U_i}\cong\mathcal O_{U_i}[T_0,\dots,T_n]$ (the symmetric algebra of a free module of rank $n+1$ is a polynomial ring, see before Ex.~5.16), so $\pi^{-1}(U_i)=\operatorname{Proj}S(\mathcal E|_{U_i})\cong\operatorname{Proj}\mathcal O_{U_i}[T_0,\dots,T_n]=\mathbf P^n_{U_i}$ (Example~7.8.7); call this $\psi_i$. Over an affine $\operatorname{Spec}A\subseteq U_i\cap U_j$, $\varepsilon_i\varepsilon_j^{-1}$ is an automorphism of $A^{n+1}$, i.e.\ a matrix $(a_{ab})\in GL_{n+1}(A)$, and the induced automorphism of $S(A^{n+1})=A[T_0,\dots,T_n]$ is $T_a\mapsto\sum_ba_{ab}T_b$. Hence $\psi_i\psi_j^{-1}=\operatorname{Proj}$ of this $A$-linear graded automorphism. So $\mathbf P(\mathcal E)$ is a $\mathbf P^n$-bundle.

\textbf{(c)} Let $\pi:P\to X$ be a $\mathbf P^n$-bundle over the regular noetherian scheme $X$. If $n=0$ then $\pi$ is an isomorphism ($P$ is locally $U\times\mathbf P^0=U$) and $P=\mathbf P(\mathcal O_X)$; so let $n\ge1$. Treating the connected components separately, assume $X$ connected, hence integral (for disconnected $X$ apply this to each of the finitely many components and take the disjoint union of the resulting sheaves).

\emph{Step 1: $P$ is a regular integral noetherian scheme.} $P$ is locally $\mathbf P^n_U$ with $U$ regular, hence regular and noetherian ($\mathbf P^n_U$ is covered by spectra of polynomial rings over regular rings), so every local ring of $P$ is a UFD (6.11A). It is connected: $\pi$ is open (locally the projection $\mathbf P^n_U\to U$, flat of finite presentation), and a decomposition $P=P_1\sqcup P_2$ into open-closed sets would give $X=\pi(P_1)\sqcup\pi(P_2)$ with both open (each fibre $\mathbf P^n_{\kappa(x)}$ is connected, hence lies in one $P_i$), forcing one of them to be empty. A connected regular noetherian scheme is integral.

\emph{Step 2: extension of $\mathcal L_0$.} Let $U\ne\emptyset$ be one of the $U_i$, $V=\pi^{-1}(U)$, $\psi:V\cong\mathbf P^n_U$, $\mathcal L_0=\psi^*p_2^*\mathcal O(1)\in\operatorname{Pic}V$. (This is the argument of 6.5/6.16, done directly so that no separatedness of $X$ is needed.) Choose a nonzero rational section $s$ of $\mathcal L_0$ (possible since $V$ is integral). For each prime divisor $Y\subset V$ the local ring $\mathcal O_{V,Y}$ is a DVR, and let $n_Y=v_Y(s)$ (valuation with respect to a generator of $(\mathcal L_0)_Y$); $n_Y=0$ for all but finitely many $Y$ since $V$ is noetherian, and $\mathcal L_0\cong\mathcal O_V(D_0)$, $D_0=\sum n_YY$, via $f\mapsto fs$. Let $D=\sum n_Y\overline Y$ on $P$ ($\overline Y$ the closure, a prime divisor of $P$ with $\overline Y\cap V=Y$). Every local ring of $P$ is a UFD, so every prime divisor of $P$ is locally principal (a height-one prime of a UFD is principal), hence $D$ is a Cartier divisor, $\mathcal L'=\mathcal O_P(D)$ is invertible, and $\mathcal L'|_V=\mathcal O_V(D_0)\cong\mathcal L_0$.

\emph{Step 3: $\mathcal L'$ has degree $1$ on all fibres.} Let $\eta$ be the generic point of $X$; $P_\eta\cong\mathbf P^n_{K(X)}$ and $\mathcal L'|_{P_\eta}=\mathcal L_0|_{P_\eta}\cong\mathcal O(1)$, so $d(\eta)=1$ in the notation of Lemma~A (which was proved in Ex.~7.9 and applies to any $\mathbf P^n$-bundle). By Lemma A, $d$ is locally constant, hence $d\equiv1$ on the connected $X$, and every $x\in X$ has an open neighbourhood $U_x$ with a trivialization $\psi_x:\pi^{-1}(U_x)\cong\mathbf P^n_{U_x}$ and $\mathcal L'|_{\pi^{-1}U_x}\cong\psi_x^*\mathcal O(1)$.

\emph{Step 4: $\mathcal E=\pi_*\mathcal L'$ is locally free of rank $n+1$ and $P\cong\mathbf P(\mathcal E)$.} Over $U_x$, $\pi_*\mathcal L'\cong\pi_*\mathcal O_{\mathbf P^n_{U_x}}(1)\cong\mathcal O_{U_x}x_0\oplus\dots\oplus\mathcal O_{U_x}x_n$ (Thm.~5.13 / 7.11(a)), so $\mathcal E$ is locally free of rank $n+1$ (in particular coherent). The natural map $\pi^*\mathcal E=\pi^*\pi_*\mathcal L'\to\mathcal L'$ is, over $U_x$, the map $\mathcal O^{n+1}\to\mathcal O(1)$, $e_a\mapsto x_a$, which is surjective since $x_0,\dots,x_n$ generate $\mathcal O(1)$ (7.11(b)). By Prop.~7.12 (with $Y=P$, $g=\pi$) there is a unique morphism $f:P\to\mathbf P(\mathcal E)$ over $X$ with $f^*\mathcal O(1)\cong\mathcal L'$ compatibly with the surjections from $\pi^*\mathcal E$. Over $U_x$, after identifying $\mathbf P(\mathcal E)|_{U_x}$ with $\mathbf P^n_{U_x}$ via the basis $x_0,\dots,x_n$ of $\mathcal E|_{U_x}$, $f$ is the morphism determined (7.1) by the line bundle $\psi_x^*\mathcal O(1)$ and its sections $\psi_x^*x_a$, which by uniqueness in 7.1 is $\psi_x$ itself. So $f$ is an isomorphism over each $U_x$, hence an isomorphism: $P\cong\mathbf P(\mathcal E)$.

\emph{Weakening the hypothesis.} Regularity was used only to know that (i) the local rings of $X$ are UFD's, so that line bundles on $\mathbf P^n$ over them are $\mathcal O(d)$, and (ii) every local ring of $P$ is a UFD, so that line bundles extend from open sets (Step~2). Both hold if $X$ is merely a noetherian \emph{locally factorial} scheme (all local rings UFD's) (then $\mathbf P^n_R$ and $P$ are again locally factorial). So the statements of (c) and (d), and of Ex.~7.9, hold for locally factorial $X$. (In general the obstruction to $P$ being some $\mathbf P(\mathcal E)$ is a class in $H^2(X,\mathcal O_X^*)$, which vanishes in the locally factorial case; I do not claim a counterexample beyond that case.)

\textbf{(d)} Let $X$ be regular. By (b), $\mathcal E\mapsto\mathbf P(\mathcal E)$ sends rank $n+1$ locally free sheaves to $\mathbf P^n$-bundles; by Ex.~7.9(b)($\Leftarrow$) it is constant on the equivalence classes of $\sim$ ($\mathbf P(\mathcal E\otimes\mathcal M)\cong\mathbf P(\mathcal E)$); by 7.9(b)($\Rightarrow$) (needs $n\ge1$) two sheaves with isomorphic projective bundles are equivalent, so the induced map $\{\text{classes}\}\to\{\mathbf P^n\text{-bundles}\}/\cong$ is injective; and by (c) it is surjective. (For $n=0$ all line bundles are equivalent and $\mathbf P^0$-bundles are $X$ itself.) This gives the $1$-$1$ correspondence.\hfill$\square$

\exercise{7.11}{On a noetherian scheme $X$, different sheaves of ideals can give rise to isomorphic blown up schemes.
\begin{enumerate}[label=(\alph*),leftmargin=*]
\item If $\mathcal I$ is any coherent sheaf of ideals on $X$, show that blowing up $\mathcal I^d$ for any $d\ge1$ gives a scheme isomorphic to the blowing up of $\mathcal I$ (cf.\ Ex.~5.13).
\item If $\mathcal I$ is any coherent sheaf of ideals, and if $\mathcal J$ is an invertible sheaf of ideals, then $\mathcal I$ and $\mathcal I\cdot\mathcal J$ give isomorphic blowings-up.
\item If $X$ is regular, show that (7.17) can be strengthened as follows. Let $U\subseteq X$ be the largest open set such that $f:f^{-1}U\to U$ is an isomorphism. Then $\mathcal I$ can be chosen such that the corresponding closed subscheme $Y$ has support equal to $X-U$.
\end{enumerate}}
\pf
Write $\mathcal S=\bigoplus_{n\ge0}\mathcal I^n$, so $\tilde X=\operatorname{Proj}\mathcal S$.

(a) The sheaf of graded algebras of the blow-up of $\mathcal I^d$ is $\mathcal S'=\bigoplus_n(\mathcal I^d)^n=\bigoplus_n\mathcal I^{dn}=\mathcal S^{(d)}$, the $d$-th Veronese subalgebra. For each open affine $U$, $\mathcal S(U)=\bigoplus I^n$ is generated by its degree-$1$ part over $\mathcal S_0(U)=\mathcal O(U)$, so by Ex.~5.13, $\operatorname{Proj}\mathcal S(U)^{(d)}\cong\operatorname{Proj}\mathcal S(U)$ over $U$ (with $\mathcal O(1)$ corresponding to $\mathcal O(d)$). This isomorphism is canonical: for $f\in S_e$ ($e\ge1$) it identifies $D_+(f)\subseteq\operatorname{Proj}S$ with $D_+(f^d)\subseteq\operatorname{Proj}S^{(d)}$ via the equality of rings $S_{(f)}=(S^{(d)})_{(f^d)}$ (an element $g/f^k$, $g\in S_{ek}$, equals $gf^{(d-1)k}/f^{dk}$ with numerator of degree divisible by $d$). It is therefore compatible with passing to smaller affine opens $U'\subseteq U$, and the isomorphisms glue to an isomorphism of the blow-ups over $X$.

(b) Since $\mathcal J\subseteq\mathcal O_X$ is invertible, locally $\mathcal J=(h)$ with $\mathcal O\to\mathcal J$, $1\mapsto h$ an isomorphism, so $h$ is a nonzerodivisor. Hence multiplication by $h^n$ gives isomorphisms $\mathcal I^n\cong h^n\mathcal I^n=(\mathcal I\mathcal J)^n$ locally, and globally the multiplication map $\mathcal I^n\otimes\mathcal J^{\otimes n}\to\mathcal I^n\mathcal J^n=(\mathcal I\mathcal J)^n$ is an isomorphism (check locally). These are compatible with the products, so $\mathcal S':=\bigoplus_n(\mathcal I\mathcal J)^n\cong\mathcal S*\mathcal J$ (notation of the Lemma numbered 7.9 in the text, with $\mathcal S'_d=\mathcal S_d\otimes\mathcal J^d$). Both satisfy ($\dagger$) (Lemma~7.9). By that Lemma, $\operatorname{Proj}\mathcal S'\cong\operatorname{Proj}\mathcal S$ over $X$: the blow-ups of $\mathcal I\mathcal J$ and of $\mathcal I$ are isomorphic (and $\mathcal O_{\tilde X}(1)$ changes by $\pi^*\mathcal J$).

(c) Let $X$ be a regular noetherian integral scheme (a nonsingular variety), $f:Z\to X$ as in (7.17), and $\mathcal I_0$ a coherent sheaf of ideals with $Z\cong\operatorname{Bl}_{\mathcal I_0}X$ over $X$; then $\mathcal I_0\ne0$. Let $U$ be the largest open subset over which $f$ is an isomorphism (the union of all such opens). I produce $\mathcal I$ with $\operatorname{Bl}_{\mathcal I}X\cong Z$ over $X$ and $\operatorname{Supp}(\mathcal O_X/\mathcal I)=X-U$.

\emph{Step 1: removing the divisorial part.} For each prime divisor $Y\subset X$ (closed integral of codimension~$1$) the ring $\mathcal O_{X,Y}$ is a DVR with valuation $v_Y$; let $m_Y=\min\{v_Y(g):g\in(\mathcal I_0)_Y\}$. Since $\mathcal I_0\ne0$ is coherent, $m_Y=0$ except for finitely many $Y$. Let $D=\sum m_YY$, an effective Cartier divisor ($X$ is locally factorial), and $\mathcal J=\mathcal O_X(-D)\subseteq\mathcal O_X$, an invertible sheaf of ideals, locally generated at $x$ by $h=\prod\pi_Y^{m_Y}$ ($Y\ni x$, $\pi_Y$ a generator of the prime ideal of $Y$ in the UFD $\mathcal O_{X,x}$). If $g\in(\mathcal I_0)_x$ then $v_Y(g)\ge m_Y$ for all $Y\ni x$, so $h\mid g$ in the UFD $\mathcal O_{X,x}$. Thus $\mathcal I_0\subseteq\mathcal J$, and $\mathcal I:=\mathcal J^{-1}\mathcal I_0\subseteq\mathcal O_X$ is a coherent ideal sheaf with $\mathcal I_0=\mathcal J\mathcal I$. By~(b), $\operatorname{Bl}_{\mathcal I}X\cong\operatorname{Bl}_{\mathcal I_0}X\cong Z$ over $X$. Moreover $\mathcal I_Y=\mathcal O_{X,Y}$ for every prime divisor $Y$ (its valuation minimum is $m_Y-m_Y=0$), so \emph{no irreducible component of $V(\mathcal I)$ has codimension $\le1$} (and $\mathcal I\ne0$).

\emph{Step 2: $\operatorname{Supp}(\mathcal O_X/\mathcal I)=X-U$.} Let $\pi:\tilde X=Z\to X$ be the blow-up of $\mathcal I$. By (7.13(b)), $\pi$ is an isomorphism over $X-V(\mathcal I)$, so $X-V(\mathcal I)\subseteq U$, i.e.\ $X-U\subseteq V(\mathcal I)$. Conversely let $x\in V(\mathcal I)$ and suppose $x\in U$; let $W\ni x$ be an open set with $\pi^{-1}W\to W$ an isomorphism. By (7.13(a)) $\mathcal I\cdot\mathcal O_{\tilde X}$ is invertible, and over $W$ it corresponds under the isomorphism to $\mathcal I|_W$; so $\mathcal I_x=(g)$ is a principal ideal of $\mathcal O_{X,x}$, with $g\neq0$ (as $\mathcal I\ne0$) and $g$ a non-unit (as $x\in V(\mathcal I)$). By Krull's Hauptidealsatz a minimal prime over $(g)$ has height~$1$, i.e.\ $V(\mathcal I)$ has an irreducible component of codimension $1$ through $x$, contradicting Step~1. Hence $V(\mathcal I)\subseteq X-U$. So $V(\mathcal I)=X-U$, i.e.\ the closed subscheme $Y$ defined by $\mathcal I$ has support $X-U$.\hfill$\square$

\exercise{7.12}{Let $X$ be a noetherian scheme, and let $Y,Z$ be two closed subschemes, neither one containing the other. Let $\tilde X$ be obtained by blowing up $Y\cap Z$ (defined by the ideal sheaf $\mathcal I_Y+\mathcal I_Z$). Show that the strict transforms $\tilde Y$ and $\tilde Z$ of $Y$ and $Z$ in $\tilde X$ do not meet.}
\pf
The question is local on $X$ (the sets $\pi^{-1}(U)$, $U\subseteq X$ affine open, cover $\tilde X$), so let $X=\operatorname{Spec}A$, $\mathcal I_Y=\mathfrak a$, $\mathcal I_Z=\mathfrak b$, $I=\mathfrak a+\mathfrak b$, $S=\bigoplus_{d\ge0}I^d$ and $\tilde X=\operatorname{Proj}S$.

By Cor.~7.15 and its proof, $\tilde Y$ is the blow-up of $Y$ along the ideal $I\mathcal O_Y=(I+\mathfrak a)/\mathfrak a$ (the ideal of $Y\cap Z$ in $Y$), namely $\operatorname{Proj}S'$ with $S'=\bigoplus_d(I\mathcal O_Y)^d=\bigoplus_d(I^d+\mathfrak a)/\mathfrak a\cong\bigoplus_dI^d/(I^d\cap\mathfrak a)$, embedded in $\operatorname{Proj}S$ via the surjection $S\to S'$ whose kernel is the graded ideal $\mathfrak A=\bigoplus_d(I^d\cap\mathfrak a)$. So $\tilde Y=V_+(\mathfrak A)$ and likewise $\tilde Z=V_+(\mathfrak B)$, $\mathfrak B=\bigoplus_d(I^d\cap\mathfrak b)$. Hence $\tilde Y\cap\tilde Z=V_+(\mathfrak A+\mathfrak B)$.

\textbf{Claim:} $(\mathfrak A+\mathfrak B)_d=I^d$ for every $d\ge1$. We have $\subseteq$ trivially. Conversely $I^d=\sum_{i=0}^d\mathfrak a^i\mathfrak b^{d-i}$. For $i\ge1$, $\mathfrak a^i\mathfrak b^{d-i}\subseteq\mathfrak a\cap I^d$, and $\mathfrak b^d\subseteq\mathfrak b\cap I^d$. So $I^d\subseteq(I^d\cap\mathfrak a)+(I^d\cap\mathfrak b)$.

Hence $\mathfrak A+\mathfrak B\supseteq S_+=\bigoplus_{d\ge1}S_d$, so $V_+(\mathfrak A+\mathfrak B)=\emptyset$ (a homogeneous prime of $S$ containing $\mathfrak A+\mathfrak B$ contains $S_+$, so it is not a point of $\operatorname{Proj}S$). Thus $\tilde Y\cap\tilde Z=\emptyset$. (The hypothesis that neither of $Y,Z$ contains the other is not needed for the disjointness; it only guarantees that the strict transforms are the honest, proper, transforms.)\hfill$\square$


\exercise{7.13}{\emph{A Complete Nonprojective Variety.} Let $k$ be an algebraically closed field of char${}\ne2$. Let $C\subseteq\mathbf P^2_k$ be the nodal cubic curve $y^2z=x^3+x^2z$. If $P_0=(0,0,1)$ is the singular point, then $C-P_0$ is isomorphic to the multiplicative group $\mathbf G_m=\operatorname{Spec}k[t,t^{-1}]$ (Ex.~6.7). For each $a\in k$, $a\ne0$, consider the translation of $\mathbf G_m$ given by $t\mapsto at$. This induces an automorphism of $C$ which we denote by $\varphi_a$.

Now consider $C\times(\mathbf P^1-\{0\})$ and $C\times(\mathbf P^1-\{\infty\})$. We glue their open subsets $C\times(\mathbf P^1-\{0,\infty\})$ by the isomorphism $\varphi:\langle P,u\rangle\mapsto\langle\varphi_u(P),u\rangle$ for $P\in C$, $u\in\mathbf G_m=\mathbf P^1-\{0,\infty\}$. Thus we obtain a scheme $X$, which is our example. The projections to the second factor are compatible with $\varphi$, so there is a natural morphism $\pi:X\to\mathbf P^1$.
\begin{enumerate}[label=(\alph*),leftmargin=*]
\item Show that $\pi$ is a proper morphism, and hence that $X$ is a complete variety over $k$.
\item Use the method of (Ex.~6.9) to show that $\operatorname{Pic}(C\times\mathbf A^1)\cong\mathbf G_m\times\mathbf Z$ and $\operatorname{Pic}(C\times(\mathbf A^1-\{0\}))\cong\mathbf G_m\times\mathbf Z\times\mathbf Z$. [\emph{Hint:} If $A$ is a domain and $*$ denotes the group of units, then $(A[u])^*\cong A^*$ and $(A[u,u^{-1}])^*\cong A^*\times u^{\mathbf Z}$.]
\item Now show that the restriction map $\operatorname{Pic}(C\times\mathbf A^1)\to\operatorname{Pic}(C\times(\mathbf A^1-\{0\}))$ is of the form $\langle t,n\rangle\mapsto\langle t,0,n\rangle$, and that the automorphism $\varphi$ of $C\times(\mathbf A^1-\{0\})$ induces a map of the form $\langle t,d,n\rangle\mapsto\langle t,d+n,n\rangle$ on its Picard group.
\item Conclude that the image of the restriction map $\operatorname{Pic}X\to\operatorname{Pic}(C\times\{0\})$ consists entirely of divisors of degree $0$ on $C$. Hence $X$ is not projective over $k$ and $\pi$ is not a projective morphism.
\end{enumerate}
(Starred in the book.)}
\pf
\textbf{0. Structure of $C$.} Use homogeneous coordinates $(\tau_0:\tau_1)$ on $\mathbf P^1$, $0=(1:0)$, $\infty=(0:1)$, and $t=\tau_1/\tau_0$ (the parameter of $\mathbf G_m=\mathbf P^1-\{0,\infty\}$). The morphism
\[
\nu:\mathbf P^1\to C,\qquad (\tau_0:\tau_1)\mapsto\bigl(4\tau_0\tau_1(\tau_1-\tau_0):4\tau_0\tau_1(\tau_1+\tau_0):(\tau_1-\tau_0)^3\bigr)
\]
lands on $C$ (a direct check: $y^2z=16\tau_0^2\tau_1^2(\tau_1+\tau_0)^2(\tau_1-\tau_0)^3=x^2(x+z)$), and I use the following standard facts about the nodal cubic.
\begin{itemize}
\item[(F1)] $\nu$ is finite and birational, $\nu^{-1}(P_0)=\{0,\infty\}$, and $\nu$ restricts to an isomorphism $\mathbf G_m\cong C-P_0=:C^\circ$ (this is the isomorphism of Ex.~6.7).
\item[(F2)] There is an exact sequence of sheaves on $C$: $0\to\mathcal O_C\to\nu_*\mathcal O_{\mathbf P^1}\to k_{P_0}\to0$, the last map being $F\mapsto F(0)-F(\infty)$ (the difference of the values at the two points over $P_0$).
\end{itemize}
Verification: on the chart $z\ne0$, $C$ is $y^2=x^2(x+1)$; with the slope $s=y/x$ we get $x=s^2-1$, $y=s(s^2-1)$, $s=(t+1)/(t-1)$, the chart is $\nu(\mathbf A^1_s)$ ($z\ne0\iff t\ne1$), and the points $t=\infty,0$ correspond to $s=1,-1$. Since $k[s]=k[x]\oplus s\,k[x]$ ($s^2=x+1$), one finds $k[x,y]/(y^2-x^2(x+1))=k[x]\oplus xs\,k[x]=k+(s^2-1)k[s]=\{f\in k[s]:f(1)=f(-1)\}$ (here $\operatorname{char}k\ne2$ is used: $1\ne-1$). The remaining point $(0:1:0)$ ($t=1$) is a smooth point of $C$ (local coordinate $x/y=1/s$) over which $\nu$ is an isomorphism, and $C-P_0$ is the union of $D(x)$ in the chart and this point.

Let $W$ be $\mathbf A^1=\operatorname{Spec}k[u]$ or $\mathbf G_m=\operatorname{Spec}k[u,u^{-1}]$, $A=\mathcal O(W)$, $Y_W=C\times W$, $\tilde Y_W=\mathbf P^1\times W$, $\nu_W=\nu\times\mathrm{id}$ (finite, birational), $D_W=P_0\times W\cong W\subseteq Y_W$. Tensoring (F2) with $A$ over $k$ (flat base change) gives the exact sequence
\[
0\to\mathcal O_{Y_W}\to\nu_{W*}\mathcal O_{\tilde Y_W}\to i_*\mathcal O_{D_W}\to0,\qquad F\mapsto F|_{\{0\}\times W}-F|_{\{\infty\}\times W}, \tag{$*$}
\]
where $\{0\}\times W$ and $\{\infty\}\times W$ are both identified with $D_W$.

\textbf{The gluing.} The automorphism $\tilde\varphi:\mathbf P^1\times\mathbf G_m\to\mathbf P^1\times\mathbf G_m$, $((\tau_0:\tau_1),u)\mapsto((\tau_0:u\tau_1),u)$ (i.e.\ $t\mapsto ut$) is the identity on $\{0\}\times\mathbf G_m$ and on $\{\infty\}\times\mathbf G_m$, so $\tilde\varphi^*$ preserves the subsheaf $\mathcal O_{Y}\subseteq\nu_*\mathcal O_{\tilde Y}$ in $(*)$ and induces an automorphism $\varphi$ of $Y_{\mathbf G_m}=C\times\mathbf G_m$ over $\mathbf G_m$ with $\nu\tilde\varphi=\varphi\nu$; on $C^\circ\times\mathbf G_m$ it is $(t,u)\mapsto(ut,u)$, and $\varphi_u$ is its fibre over $u$ (the translation of $C$ extending $t\mapsto ut$, fixing $P_0$). This is the isomorphism used in the gluing. Write $X_1=C\times(\mathbf P^1-\{\infty\})=C\times\mathbf A^1_u$ and $X_2=C\times(\mathbf P^1-\{0\})=C\times\mathbf A^1_{u'}$ ($u'=u^{-1}$), glued along $X_{12}=C\times\mathbf G_m$ by identifying $(P,u)\in X_1$ with $\varphi(P,u)\in X_2$.

\textbf{(a)} $\pi^{-1}(\mathbf P^1-\{\infty\})=X_1\to\mathbf P^1-\{\infty\}$ and $\pi^{-1}(\mathbf P^1-\{0\})=X_2\to\mathbf P^1-\{0\}$ are the projections $C\times\mathbf A^1\to\mathbf A^1$, base changes of $C\to\operatorname{Spec}k$, which is proper ($C$ is projective, Thm.~4.9). Properness is local on the base (4.8(f)), so $\pi$ is proper. Since $\mathbf P^1\to\operatorname{Spec}k$ is proper, so is the composite $X\to\operatorname{Spec}k$ (4.8(b)); in particular $X$ is separated and of finite type. $X$ is integral: $X_1,X_2$ are integral ($C$ is an integral variety over an algebraically closed field), and they overlap in the nonempty open set $X_{12}$. Hence $X$ is a complete variety over $k$.

\textbf{(b)} \emph{Claim.} For $W,A$ as above let $\mathcal L_W(\theta,n)$, for $\theta\in A^*$, $n\in\mathbf Z$, be the following subsheaf of $\nu_{W*}\mathcal O_{\tilde Y}(n)$ ($\mathcal O(n)=p_1^*\mathcal O_{\mathbf P^1}(n)$): near $\{0\}\times W$ write a section as $F=f\tau_0^n$, near $\{\infty\}\times W$ as $F=g\tau_1^n$ (the sections $\tau_0^n,\tau_1^n$ generate $\mathcal O(n)$ near $0$, resp.\ $\infty$); then $F\in\mathcal L_W(\theta,n)$ iff $g|_{\{\infty\}\times W}=\theta\cdot f|_{\{0\}\times W}$. Then:
\begin{enumerate}
\item[(i)] $\mathcal L_W(\theta,n)$ is an invertible $\mathcal O_{Y_W}$-module, $\nu_W^*\mathcal L_W(\theta,n)\cong\mathcal O_{\tilde Y_W}(n)$, and $\mathcal L_W(\theta,n)\otimes\mathcal L_W(\theta',n')\cong\mathcal L_W(\theta\theta',n+n')$ (by multiplication of sections);
\item[(ii)] every invertible sheaf on $Y_W$ is isomorphic to $\mathcal L_W(\theta,n)$ for exactly one pair $(\theta,n)\in A^*\times\mathbf Z$; so $\operatorname{Pic}Y_W\cong A^*\times\mathbf Z$.
\end{enumerate}
Away from $D_W$ the sheaf is $\nu_*\mathcal O(n)$, invertible since $\nu_W$ is an isomorphism there. Near $D_W$, on $\nu^{-1}(C_{z\ne0})\times W=\mathbf A^1_s\times W$ the section $\rho=(\tau_1-\tau_0)^n$ trivializes $\mathcal O(n)$, and writing $F=h\rho$ the condition becomes $h(s{=}1)=\theta'h(s{=}-1)$ with $\theta'=(-1)^n\theta$ (recall $t=\infty\leftrightarrow s=1$, $t=0\leftrightarrow s=-1$), while $\mathcal O_Y=\{h:h(1)=h(-1)\}$ by $(*)$. Let $h_0=\frac{1-s}2+\frac{1+s}2\theta'\in A[s]$, so $h_0(-1)=1$, $h_0(1)=\theta'$ are units. The zero set $G$ of $h_0$ misses $\{s=\pm1\}\times W$, so $V=(C_{z\ne0}\times W)-\nu_W(G)$ ($\nu_W$ is closed) is an open neighbourhood of $D_W$, on which $h_0$ is a unit of $\nu_*\mathcal O$ and $\mathcal L_W(\theta,n)|_V=h_0\rho\cdot\mathcal O_V$ (if $h\in\mathcal L$ then $h/h_0$ takes the same value at $s=\pm1$; conversely $h_0\mathcal O_V\subseteq\mathcal L$). So $\mathcal L$ is invertible and $\nu^*\mathcal L=h_0\rho\cdot\mathcal O_{\tilde Y}=\mathcal O(n)$ on $\nu^{-1}V$. The tensor-product statement is checked on these local generators. This proves (i).

For (ii) use the method of Ex.~6.9. By $(*)$ there is an exact sequence of sheaves of abelian groups
\[
1\to\mathcal O_{Y}^*\to\nu_*\mathcal O_{\tilde Y}^*\to i_*\mathcal O_{D}^*\to1,\qquad F\mapsto F|_{\{0\}\times W}\big/F|_{\{\infty\}\times W}
\]
(kernel: units $F$ with equal values at $0,\infty$ lie in $\mathcal O_Y$, together with $F^{-1}$; surjectivity: near a point $(P_0,w)$, given a unit $\vartheta$ of $\mathcal O_{W,w}$, the function $F=\frac{1-s}2+\frac{1+s}2\vartheta$ on $\mathbf A^1_s\times W'$ has $F|_{s=-1}=1$, $F|_{s=1}=\vartheta$, so it is a unit on a neighbourhood of $\{s=\pm1\}\times W'$, which contains $\nu^{-1}$ of a neighbourhood of $(P_0,w)$ ($\nu$ is closed); its ratio of values at $t=0$ and $t=\infty$ is $\vartheta^{-1}$, and $\vartheta$ was arbitrary). Take cohomology on $Y=Y_W$ (Zariski). We have $H^1(Y,\mathcal O_Y^*)=\operatorname{Pic}Y$ (Ex.~III.4.5: line bundles are classified by $H^1(\mathcal O^*)$); $H^0(Y,\nu_*\mathcal O^*_{\tilde Y})=\mathcal O(\tilde Y)^*=A^*$ since $\Gamma(\mathbf P^1\times W,\mathcal O)=A$; $H^0(Y,\mathcal O_Y)=\{F\in A:F|_0=F|_\infty\}=A$ by $(*)$, so $H^0(Y,\mathcal O_Y^*)=A^*$ too; $H^0(Y,i_*\mathcal O^*_D)=A^*$; $H^1(Y,i_*\mathcal O_D^*)=H^1(D,\mathcal O_D^*)=\operatorname{Pic}W=0$ ($A$ is a PID); and $H^1(Y,\nu_*\mathcal O^*_{\tilde Y})=\operatorname{Pic}\tilde Y$, because $R^1\nu_*\mathcal O^*_{\tilde Y}=0$: its stalk at $y$ is the direct limit of $\operatorname{Pic}(\nu^{-1}V)$ over neighbourhoods $V$ of $y$, and a line bundle on the quasi-projective scheme $\tilde Y$ is trivial on a neighbourhood of the finite set $\nu^{-1}(y)$ (semilocal rings have trivial Picard group), which contains some $\nu^{-1}(V')$ as $\nu$ is closed. So we get the exact sequence
\[
1\to A^*\xrightarrow{\ \mathrm{id}\ }A^*\xrightarrow{\ 0\ }A^*\xrightarrow{\ \delta\ }\operatorname{Pic}Y\xrightarrow{\ \nu^*\ }\operatorname{Pic}\tilde Y\to0 ,
\]
where the middle map is zero because $F\in\mathcal O(\tilde Y)^*=A^*$ is constant along $\mathbf P^1$ so $F|_0/F|_\infty=1$. Thus $\delta:A^*\to\operatorname{Pic}Y$ is injective and $\operatorname{Pic}Y/\delta(A^*)\cong\operatorname{Pic}\tilde Y$. Unwinding $\delta$: $\theta\in A^*=H^0(i_*\mathcal O_D^*)$ is lifted locally to units $g_V$ of $\nu_*\mathcal O_{\tilde Y}$ with $g_V|_0/g_V|_\infty=\theta$, and $\delta(\theta)$ is (up to the convention $\delta(\theta)^{\pm1}$, which is immaterial below) the class of the invertible sheaf $\{g_V\mathcal O_V\}=\{F\in\nu_*\mathcal O_{\tilde Y}:F|_\infty=\theta^{-1}F|_0\}=\mathcal L_W(\theta^{-1},0)$. Hence $\delta(A^*)=\{[\mathcal L_W(\theta,0)]:\theta\in A^*\}$, and $\theta\mapsto[\mathcal L_W(\theta,0)]$ is injective. Finally $\operatorname{Pic}\tilde Y=\operatorname{Pic}(\mathbf P^1\times W)=\mathbf Z\cdot[\mathcal O(1)]$: by Ex.~6.1, $\operatorname{Cl}(\mathbf P^1\times W)\cong\operatorname{Cl}W\times\mathbf Z=\mathbf Z$ with generator $V_+(\tau_0)\times W$, i.e.\ $\mathcal O(1)$, and $\operatorname{Pic}=\operatorname{Cl}$ since $\mathbf P^1\times W$ is regular; $\mathcal O(n)$ for $n\ne0$ is nontrivial. By (i), $\nu^*\mathcal L_W(1,n)=\mathcal O(n)$, so $\mathcal L_W(1,n)$ lifts $[\mathcal O(n)]$. Therefore every invertible sheaf is $\mathcal L_W(\theta,0)\otimes\mathcal L_W(1,n)\cong\mathcal L_W(\theta,n)$, and $(\theta,n)$ is unique ($n$ is read off from $\nu^*$, then $\theta$ by injectivity of $\delta$). This proves (ii).

For $W=\mathbf A^1$: $A^*=k[u]^*=k^*=\mathbf G_m$, so $\operatorname{Pic}(C\times\mathbf A^1)\cong\mathbf G_m\times\mathbf Z$, $\langle c,n\rangle\leftrightarrow\mathcal L(c,n)$. For $W=\mathbf G_m$: $A^*=k[u,u^{-1}]^*=k^*\times u^{\mathbf Z}$, so $\operatorname{Pic}(C\times\mathbf G_m)\cong\mathbf G_m\times\mathbf Z\times\mathbf Z$, $\langle c,d,n\rangle\leftrightarrow\mathcal L(cu^d,n)$. (I write $c$ for the first coordinate to avoid a clash with the parameter $t$.)

\textbf{(c)} The sheaf $\mathcal L_{\mathbf A^1}(c,n)$ restricted to $C\times\mathbf G_m$ is, by definition, $\mathcal L_{\mathbf G_m}(c,n)$ (the same description, the constant $c\in k^*\subseteq A^*$). So restriction is $\langle c,n\rangle\mapsto\langle c,0,n\rangle$.

For $\varphi^*$: $\varphi^*\mathcal L_{\mathbf G_m}(\theta,n)$ consists of the pull-backs $\tilde\varphi^*F$ of sections $F$ of $\mathcal L(\theta,n)$, viewed in $\nu_*\mathcal O(n)$ through the identification $\tilde\varphi^*\mathcal O(n)\cong\mathcal O(n)$ for which $\tilde\varphi^*\tau_0^n=\tau_0^n$ and $\tilde\varphi^*\tau_1^n=u^n\tau_1^n$ (the morphism $\tilde\varphi$ is defined by the sections $\tau_0,u\tau_1$ of $\mathcal O(1)$, 7.1). If $F=f\tau_0^n=g\tau_1^n$ near $0$, resp.\ $\infty$, then $\tilde\varphi^*F=f'\tau_0^n=g'\tau_1^n$ with $f'=\tilde\varphi^*f$, $g'=u^n\tilde\varphi^*g$; since $\tilde\varphi$ is the identity on $\{0,\infty\}\times\mathbf G_m$, $f'|_0=f|_0$ and $g'|_\infty=u^ng|_\infty$. So $g|_\infty=\theta f|_0$ becomes $g'|_\infty=u^n\theta f'|_0$:
\[
\varphi^*\mathcal L_{\mathbf G_m}(\theta,n)=\mathcal L_{\mathbf G_m}(u^n\theta,n),\qquad\text{i.e.}\qquad \varphi^*\langle c,d,n\rangle=\langle c,d+n,n\rangle .
\]

\textbf{(d)} Let $\mathcal M\in\operatorname{Pic}X$. By (b), $\mathcal M|_{X_1}\cong\mathcal L_{\mathbf A^1_u}(c,n)=\langle c,n\rangle$ and $\mathcal M|_{X_2}\cong\langle c',n'\rangle$ (with respect to the coordinate $u'$; the constants $k^*\subseteq A^*$ are the same for $u$ and $u'=u^{-1}$). Restricted to $X_{12}$ these are $\langle c,0,n\rangle$ and $\langle c',0,n'\rangle$ by (c). The gluing identifies $X_{12}\subseteq X_1$ with $X_{12}\subseteq X_2$ via $\varphi$, so $\mathcal M|_{X_{12}\subseteq X_1}\cong\varphi^*(\mathcal M|_{X_{12}\subseteq X_2})$, i.e.\ $\langle c,0,n\rangle=\varphi^*\langle c',0,n'\rangle=\langle c',n',n'\rangle$ in $\operatorname{Pic}(C\times\mathbf G_m)$. Comparing coordinates: $c=c'$, $0=n'$, $n=n'$. So $n=n'=0$ and $\mathcal M|_{X_1}\cong\mathcal L_{\mathbf A^1}(c,0)$.

The closed subscheme $C\times\{0\}\subseteq X_1$ is $\cong C$, and $\mathcal L_{\mathbf A^1}(c,0)|_{C\times\{0\}}$ is the sheaf $\{F\in\nu_*\mathcal O_{\mathbf P^1}:F(\infty)=cF(0)\}$ on $C$, whose pull-back to $\mathbf P^1$ is $\mathcal O_{\mathbf P^1}$. So every element of the image of $\operatorname{Pic}X\to\operatorname{Pic}(C\times\{0\})=\operatorname{Pic}C$ has degree $0$, where $\deg\mathcal N:=\deg\nu^*\mathcal N\in\mathbf Z=\operatorname{Pic}\mathbf P^1$ is the degree of the corresponding Cartier divisor class (these form $\operatorname{CaCl}^0C\cong\mathbf G_m$ of Ex.~6.7).

\emph{$X$ is not projective.} Suppose $X$ were projective over $k$. Then there is a very ample invertible sheaf $\mathcal M$ on $X$, i.e.\ a closed immersion $X\hookrightarrow\mathbf P^N$ with $\mathcal M\cong\mathcal O(1)|_X$. Restricting to the closed subscheme $C\times\{0\}\cong C$ gives a closed immersion $j:C\hookrightarrow\mathbf P^N$ with $\mathcal M|_C\cong j^*\mathcal O(1)$. Then $(j\nu)^*\mathcal O(1)\cong\nu^*(\mathcal M|_C)$ has degree $0$ by what we have shown, hence is $\mathcal O_{\mathbf P^1}$, which is generated by the constant sections $(j\nu)^*x_i$; so $j\nu$ is constant. But $\nu$ is nonconstant and $j$ injective, a contradiction. So $X$ is not projective over $k$.

\emph{$\pi$ is not projective.} If $\pi:X\to\mathbf P^1$ were projective, $X$ would be a closed subscheme of $\mathbf P^r_{\mathbf P^1}=\mathbf P^r_k\times_k\mathbf P^1_k$, which is a closed subscheme of a projective space over $k$ by the Segre embedding; then $X$ would be projective over $k$, contradicting the above. Hence $X$ is a complete variety over $k$ that is not projective, and $\pi:X\to\mathbf P^1$ is proper but not projective.\hfill$\square$

\exercise{7.14}{\begin{enumerate}[label=(\alph*),leftmargin=*]
\item Give an example of a noetherian scheme $X$ and a locally free coherent sheaf $\mathcal E$ on $X$, such that the invertible sheaf $\mathcal O(1)$ on $\mathbf P(\mathcal E)$ is not very ample relative to $X$.
\item Let $f:X\to Y$ be a morphism of finite type, let $\mathcal L$ be an ample invertible sheaf on $X$, and let $\mathcal S$ be a sheaf of graded $\mathcal O_X$-algebras satisfying ($\dagger$). Let $P=\operatorname{Proj}\mathcal S$, let $\pi:P\to X$ be the projection, and let $\mathcal O_P(1)$ be the associated invertible sheaf. Show that for all $n\gg0$, the sheaf $\mathcal O_P(1)\otimes\pi^*\mathcal L^n$ is very ample on $P$ relative to $Y$. [\emph{Hint:} Use (7.10) and (Ex.~5.12).]
\end{enumerate}}
\pf
(a) Let $k$ be a field, $X=\mathbf P^1_k$, $\mathcal E=\mathcal O_X(-1)\oplus\mathcal O_X(-1)$ (locally free of rank $2$), $P=\mathbf P(\mathcal E)$. By Prop.~7.11(a) (rank $\ge2$), $\pi_*\mathcal O_P(1)=\mathcal E$, so $\Gamma(P,\mathcal O_P(1))=\Gamma(X,\mathcal E)=\Gamma(\mathbf P^1,\mathcal O(-1))^{2}=0$. Suppose $\mathcal O_P(1)$ were very ample relative to $X$: there would be an immersion $i:P\to\mathbf P^N_X$ over $X$ with $i^*\mathcal O(1)\cong\mathcal O_P(1)$. Then $\mathcal O_P(1)$ would be generated by the global sections $i^*x_0,\dots,i^*x_N$ (7.1(a)), all zero, so $\mathcal O_P(1)=0$, absurd since $P\neq\emptyset$. Hence $\mathcal O_P(1)$ is not very ample relative to $X$. (More generally, if $\mathcal O_P(1)$ is very ample relative to $X$ then $\mathcal E=\pi_*\mathcal O_P(1)$ is generated by finitely many global sections, since the $i^*x_a\in\Gamma(P,\mathcal O_P(1))=\Gamma(X,\mathcal E)$ generate $\mathcal O_P(1)$ and hence, restricted to each fibre $\mathbf P(\mathcal E\otimes\kappa(x))$, span $\mathcal E\otimes\kappa(x)$. Conversely if $\mathcal O^{N+1}\to\mathcal E$ is surjective then $\mathbf P(\mathcal E)\hookrightarrow\mathbf P^N_X$ is a closed immersion with $\mathcal O(1)$ restricting to $\mathcal O_P(1)$. So for affine $X$ the sheaf $\mathcal O_P(1)$ is always very ample, and a non-affine $X$ with $\mathcal E$ not globally generated is the natural place to look for examples.)

(b) \emph{Setup.} All schemes are noetherian ($Y$ noetherian, $X$ of finite type over $Y$). Let $n_0$ be such that $\mathcal S_1\otimes\mathcal L^n$ is generated by global sections for all $n\ge n_0$ ($\mathcal L$ ample, $\mathcal S_1$ coherent). The proof of Prop.~7.10(b) uses only that $\mathcal S_1\otimes\mathcal L^n$ is generated by finitely many global sections; hence, for every $n\ge n_0$, $\mathcal O_P(1)\otimes\pi^*\mathcal L^n$ is very ample on $P$ relative to $X$ (there is an immersion, indeed closed, $P\hookrightarrow\mathbf P^{M}_X$ over $X$ pulling $\mathcal O(1)$ back to it).

\emph{Lemma (ample $\Rightarrow$ some power very ample relative to $Y$).} Let $f:X\to Y$ be of finite type, $Y$ noetherian, and $\mathcal L$ ample on $X$. Then $\mathcal L^{m_1}$ is very ample on $X$ relative to $Y$ for some $m_1>0$. (For $Y=\operatorname{Spec}A$ this is Thm.~7.6; the proof of 7.6 adapts as follows.) Cover $Y$ by finitely many affine opens $U_j=\operatorname{Spec}A_j$. As in the proof of Thm.~7.6 (which uses only that $\mathcal L$ is ample on the noetherian scheme $X$), for every $P\in X$ and every affine open neighbourhood $U$ of $P$ on which $\mathcal L$ is trivial there are $n>0$ and $s\in\Gamma(X,\mathcal L^n)$ with $P\in X_s\subseteq U$ and $X_s$ affine; taking $U\subseteq f^{-1}(U_j)$ for some $j\ni f(P)$ and using quasi-compactness of $X$, we get $s_1,\dots,s_k$, all in $\Gamma(X,\mathcal L^n)$ for one $n$ (replace $s_i$ by powers), with $X_i=X_{s_i}$ affine, $\bigcup X_i=X$ and $X_i\subseteq f^{-1}(U_{j(i)})$. Since $f^{-1}(U_{j(i)})$ is of finite type over $A_{j(i)}$, $B_i=\Gamma(X_i,\mathcal O_X)$ is a finitely generated $A_{j(i)}$-algebra, say by $b_{i1},\dots,b_{il_i}$. By (5.14) there is $N$ such that all $s_i^Nb_{il}$ extend to sections $c_{il}\in\Gamma(X,\mathcal L^{nN})$. The sections $s_i^N$ generate $\mathcal L^{nN}$, so $\{s_i^N,c_{il}\}$ define (7.1) a morphism $g:X\to\mathbf P^M_{\mathbf Z}$ and $h=(g,f):X\to\mathbf P^M_{\mathbf Z}\times Y=\mathbf P^M_Y$ over $Y$, with $h^*\mathcal O(1)=\mathcal L^{nN}$. Let $V_i\subseteq\mathbf P^M_Y$ be the open set where the coordinate of $s_i^N$ is nonzero and the image in $Y$ lies in $U_{j(i)}$; $V_i\cong\mathbf A^M_{A_{j(i)}}=\operatorname{Spec}A_{j(i)}[y_{l}]$, $h^{-1}(V_i)=X_i$, and $h:X_i\to V_i$ corresponds to $y_l\mapsto c_{il}/s_i^N=b_{il}$ (and the $y$'s attached to the other sections go to elements of $B_i$), which is surjective onto $B_i$. So $h$ restricts to closed immersions $X_i\to V_i$; since the $V_i$ cover the open set $W=\bigcup V_i$, $h$ is a closed immersion $X\to W$ and so an immersion $X\to\mathbf P^M_Y$. Thus $\mathcal L^{m_1}$, $m_1=nN$, is very ample relative to $Y$.

\emph{Conclusion.} Apply Ex.~5.12(b) to $P\xrightarrow{\pi}X\xrightarrow{f}Y$: for $a\ge n_0$ the sheaf $\mathcal O_P(1)\otimes\pi^*\mathcal L^{a}$ is very ample on $P$ relative to $X$, and $\mathcal L^{m_1}$ is very ample on $X$ relative to $Y$, so $\bigl(\mathcal O_P(1)\otimes\pi^*\mathcal L^{a}\bigr)\otimes\pi^*\mathcal L^{m_1}=\mathcal O_P(1)\otimes\pi^*\mathcal L^{a+m_1}$ is very ample on $P$ relative to $Y$. Taking $a=n-m_1$, this holds for every $n\ge n_0+m_1$.

(Recall the proof of Ex.~5.12(b) for completeness: if $P\hookrightarrow\mathbf P^M_X$ and $X\hookrightarrow\mathbf P^r_Y$ are immersions, then $P\hookrightarrow\mathbf P^M_X=\mathbf P^M_Y\times_Y X\hookrightarrow\mathbf P^M_Y\times_Y\mathbf P^r_Y\hookrightarrow\mathbf P^{(M+1)(r+1)-1}_Y$, the last map being the Segre embedding, and the pull-back of $\mathcal O(1)$ is the tensor product of the two sheaves.)\hfill$\square$

\section*{8. Differentials}

\noindent\emph{Conventions for this section.} ``Variety'' and ``nonsingular variety'' are in the sense of Chapter~II, so $k$ is algebraically closed whenever the text speaks of varieties (this is the case in Ex.~8.2--8.8; in Ex.~8.1 only a perfect field is needed). Results of the text are cited by number (Prop.~8.7, Thm.~8.17, \dots), $\mathcal N_{Y/X}=\mathcal{H}om(\mathcal I/\mathcal I^2,\mathcal O_Y)$ is the normal sheaf, and $\mathcal O(D)$ is the invertible sheaf of a divisor (written $\mathcal L(D)$ in the book).

\exercise{8.1}{Here we will strengthen the results of the text to include information about the sheaf of differentials at a not necessarily closed point of a scheme $X$.
\begin{enumerate}[label=(\alph*),leftmargin=*]
\item Generalize (8.7) as follows. Let $B$ be a local ring containing a field $k$, and assume that the residue field $k(B)=B/\fm$ of $B$ is a separably generated extension of $k$. Then the exact sequence of (8.4A),
\[0\to \fm/\fm^2\xrightarrow{\ \delta\ }\Omega_{B/k}\otimes_B k(B)\to \Omega_{k(B)/k}\to 0\]
is exact on the left also. [\emph{Hint}: In copying the proof of (8.7), first pass to $B/\fm^2$, which is a complete local ring, and then use (8.25A) to choose a field of representatives for $B/\fm^2$.]
\item Generalize (8.8) as follows. With $B,k$ as above, assume furthermore that $k$ is perfect, and that $B$ is a localization of an algebra of finite type over $k$. Then show that $B$ is a regular local ring if and only if $\Omega_{B/k}$ is free of rank $=\dim B+\operatorname{tr.d.}k(B)/k$.
\item Strengthen (8.15) as follows. Let $X$ be an irreducible scheme of finite type over a perfect field $k$, and let $\dim X=n$. Show that the local ring $\mathcal O_{X,x}$ at any point $x\in X$, not necessarily closed, is a regular local ring if and only if the stalk $(\Omega_{X/k})_x$ of the sheaf of differentials at $x$ is free of rank $n$.
\item Strengthen (8.16) as follows. If $X$ is a variety over an algebraically closed field $k$, then $U=\{x\in X\mid \mathcal O_x \text{ is a regular local ring}\}$ is an open dense subset of $X$.
\end{enumerate}}
\pf
Throughout, $\kappa=k(B)=B/\fm$.

(a) By (8.4A) (applied to the $k$-algebra $B$ and its quotient $B/\fm=\kappa$) the sequence is exact except possibly on the left, so we must show that $\delta$ is injective. Since $\delta$ is a $\kappa$-linear map of $\kappa$-vector spaces, it suffices to show that the transposed map
\[\delta':\Hom_\kappa(\Omega_{B/k}\otimes_B\kappa,\kappa)\to \Hom_\kappa(\fm/\fm^2,\kappa)\]
is surjective: if $\delta(w)=0$ with $w\neq0$ we could choose a linear form $h$ with $h(w)\ne0$ and write $h=\delta'(\varphi)=\varphi\circ\delta$, a contradiction. By the universal property of $\Omega_{B/k}$, $\Hom_\kappa(\Omega_{B/k}\otimes_B\kappa,\kappa)=\Hom_B(\Omega_{B/k},\kappa)=\operatorname{Der}_k(B,\kappa)$, where $\kappa$ is a $B$-module via $B\to\kappa$. If $d$ is such a derivation then $d(ab)=a\,db+b\,da=0$ for $a,b\in\fm$, so $d(\fm^2)=0$; and $\delta'(d)$ is the restriction $d|_\fm:\fm/\fm^2\to\kappa$ (indeed $\delta(\bar x)=dx\otimes1$ and the linear map attached to $d$ sends $dx\otimes 1$ to $d(x)$). Thus we must show: every $\kappa$-linear $h:\fm/\fm^2\to\kappa$ is the restriction of a $k$-derivation $d:B\to\kappa$.

Put $\bar B=B/\fm^2$, a local ring with maximal ideal $\bar\fm=\fm/\fm^2$ satisfying $\bar\fm^2=0$. It is complete (its $\bar\fm$-adic topology is discrete), it contains $k$ (the map $k\to B\to\bar B$ is injective because $k$ is a field and $\bar B\ne0$), and its residue field is $\kappa$, separably generated over $k$. By Cohen's theorem (8.25A) there is a subfield $K\subseteq\bar B$ containing $k$ such that $K\to\bar B/\bar\fm=\kappa$ is an isomorphism. Every element of $\bar B$ is then uniquely $\lambda+c$ with $\lambda\in K$, $c\in\bar\fm$. Define $\bar d:\bar B\to\kappa$ by $\bar d(\lambda+c)=h(c)$. It is additive and vanishes on $K\supseteq k$. For $b=\lambda+c$, $b'=\lambda'+c'$ we have $bb'=\lambda\lambda'+(\lambda c'+\lambda' c)$ because $cc'=0$; here $\lambda c'$ is the $\kappa$-vector space action of $\bar\lambda\in\kappa$ (the residue class of $\lambda$) on $\bar\fm=\fm/\fm^2$. Hence
\[\bar d(bb')=h(\lambda c'+\lambda'c)=\bar\lambda\,h(c')+\bar\lambda'\,h(c)=\bar b\,\bar d(b')+\bar b'\,\bar d(b),\]
where $\bar b,\bar b'\in\kappa$ are residues. So $\bar d$ is a $k$-derivation $\bar B\to\kappa$ with $\bar d|_{\bar\fm}=h$. The composite $d=\bar d\circ(B\to\bar B)$ is a $k$-derivation $B\to\kappa$ whose restriction to $\fm$ is $h$. Hence $\delta'$ is surjective and $\delta$ is injective.

(b) Put $s=\operatorname{tr.d.}\kappa/k$ and $r=\dim B$. Write $B=A_{\fp}$ with $A$ a finitely generated $k$-algebra. Then $\kappa$ is a finitely generated field extension of $k$ (the fraction field of a quotient of $A$), hence separably generated because $k$ is perfect (I, 4.8A). So (a) applies, and by (8.6A) $\dim_\kappa\Omega_{\kappa/k}=s$. The exact sequence of (a) gives
\[\dim_\kappa(\Omega_{B/k}\otimes_B\kappa)=\dim_\kappa\fm/\fm^2+s.\tag{$*$}\]
Also $\Omega_{B/k}$ is a finitely generated $B$-module by (8.5).

\emph{``If''.} Suppose $\Omega_{B/k}$ is free of rank $r+s$. Then the left side of $(*)$ is $r+s$, so $\dim_\kappa\fm/\fm^2=r=\dim B$, i.e.\ $B$ is regular.

\emph{``Only if''.} Suppose $B$ is regular, so $\dim_\kappa\fm/\fm^2=r$ and, by $(*)$, $\dim_\kappa(\Omega_{B/k}\otimes_B\kappa)=r+s$. A regular local ring is a domain; let $L$ be its quotient field. By (8.2A), $\Omega_{B/k}\otimes_BL=\Omega_{L/k}$, and $L$ is a finitely generated, hence (as $k$ is perfect) separably generated, extension of $k$, so $\dim_L\Omega_{L/k}=\operatorname{tr.d.}L/k$ by (8.6A). We claim $\operatorname{tr.d.}L/k=r+s$. Let $\mathfrak q=\ker(A\to B)$, a prime of $A$ because $B$ is a domain; then $B=(A/\mathfrak q)_{\fp/\mathfrak q}$, so we may assume $A$ is a domain with $L=\operatorname{Frac}A$, $\kappa=\operatorname{Frac}(A/\fp)$. Then $r=\operatorname{ht}\fp$ and $s=\dim A/\fp$ by (I, 1.8A(a)), and $\operatorname{ht}\fp+\dim A/\fp=\dim A=\operatorname{tr.d.}L/k$ by (I, 1.8A(b),(a)). So $\dim_L(\Omega_{B/k}\otimes L)=\dim_\kappa(\Omega_{B/k}\otimes\kappa)=r+s$, and Lemma~8.9 (a finitely generated module over a noetherian local domain whose two fibre dimensions agree) shows $\Omega_{B/k}$ is free of rank $r+s$.

(c) Let $x\in X$ and $B=\mathcal O_{X,x}$. Taking an affine neighbourhood $\operatorname{Spec}A$ of $x$, $B$ is a localization of a finitely generated $k$-algebra, and $(\Omega_{X/k})_x=\Omega_{B/k}$ by Remark~8.9.2/(8.2A). Also $k(B)=k(x)$. By (b), $B$ is regular iff $\Omega_{B/k}$ is free of rank $\dim B+\operatorname{tr.d.}k(x)/k$. So it suffices to show $\dim\mathcal O_{X,x}+\operatorname{tr.d.}k(x)/k=n$. Take $x\in\operatorname{Spec}A=U$ affine open, $x=\fp$. $U$ is a nonempty open subset of the irreducible $X$, so $A_{\rm red}$ is a domain; heights and dimensions are topological, so we may replace $A$ by $A_{\rm red}$ (a finitely generated domain). Then $\dim\mathcal O_{X,x}=\operatorname{ht}\fp$, $\operatorname{tr.d.}k(x)/k=\dim A/\fp$, and $\operatorname{ht}\fp+\dim A/\fp=\dim A=\operatorname{tr.d.}K(X_{\rm red})/k$ by (I, 1.8A). Finally $\dim X=\sup_i\dim U_i$ over an affine open cover (a chain of irreducible closed subsets of $X$ meeting $U_i$ restricts to a chain of the same length in $U_i$), and each $\dim U_i=\operatorname{tr.d.}K/k$, so this common value is $n$. Hence $\dim\mathcal O_{X,x}+\operatorname{tr.d.}k(x)/k=n$, and (c) follows.

(d) $X$ is irreducible of finite type over the perfect field $k$, $n=\dim X$. By (c), $x\in U$ iff $(\Omega_{X/k})_x$ is free of rank $n$. The sheaf $\Omega_{X/k}$ is coherent (Remark~8.9.1). If $x\in U$, then by (Ex.~5.7) there is an open neighbourhood $V$ of $x$ on which $\Omega_{X/k}$ is free (of rank $n$); every point of $V$ then lies in $U$ by (c). So $U$ is open. It contains the generic point $\eta$, since $\mathcal O_{X,\eta}$ is a field and fields are regular local rings; so $U$ is a nonempty open subset of the irreducible space $X$, hence dense.\hfill$\square$

\exercise{8.2}{Let $X$ be a variety of dimension $n$ over $k$. Let $\mathcal E$ be a locally free sheaf of rank $>n$ on $X$, and let $V\subseteq\Gamma(X,\mathcal E)$ be a vector space of global sections which generate $\mathcal E$. Then show that there is an element $s\in V$, such that for each $x\in X$, we have $s_x\notin\fm_x\mathcal E_x$. Conclude that there is a morphism $\mathcal O_X\to\mathcal E$ giving rise to an exact sequence
\[0\to\mathcal O_X\to\mathcal E\to\mathcal E'\to0\]
where $\mathcal E'$ is also locally free. [\emph{Hint}: Use a method similar to the proof of Bertini's theorem (8.18).]}
\pf
Let $r=\operatorname{rk}\mathcal E>n$. Write $\mathcal E(x)=\mathcal E_x/\fm_x\mathcal E_x$ for the fibre ($\cong k(x)^r$); for a section $s$ we say $s(x)=0$ if $s_x\in\fm_x\mathcal E_x$.

\emph{Step 0: reductions.} $X$ is noetherian, hence quasi-compact. Since $V$ generates $\mathcal E$, every $x$ has finitely many elements of $V$ generating $\mathcal E_x$, hence (the cokernel of the corresponding map $\mathcal O^m\to\mathcal E$ being coherent with zero stalk at $x$) generating $\mathcal E$ on a neighbourhood of $x$. Finitely many neighbourhoods cover $X$, so a finite-dimensional subspace $V_0\subseteq V$ generates $\mathcal E$; it suffices to find $s\in V_0$. So assume $\dim V=N<\infty$ and fix a basis $s_1,\dots,s_N$. Next, the set $\{x: s(x)=0\}$ is closed (locally, in a trivialization $\mathcal E|_U\cong\mathcal O_U^r$, $s=(f_1,\dots,f_r)$ and the set is $V(f_1,\dots,f_r)$). A nonempty closed subset of the noetherian scheme $X$ contains a closed point (a minimal nonempty closed subset $W$ satisfies $\overline{\{w\}}=W$ for all $w\in W$, so $W$ is one point because schemes are $T_0$). Hence it is enough to find $s\in V$ with $s(x)\ne0$ for all \emph{closed} points $x$.

\emph{Step 1: the incidence scheme.} Let $\mathbf A^N=\operatorname{Spec}k[v_1,\dots,v_N]$; its $k$-points are identified with $V$ via $a\mapsto s_a=\sum a_is_i$. Cover $X$ by finitely many affine opens $U$ with $\mathcal E|_U\cong\mathcal O_U^r$; on $U$ write $s_j=(m_{1j},\dots,m_{rj})$, giving an $r\times N$ matrix $M$ over $\mathcal O(U)$. Since the $s_j$ generate $\mathcal E$, $M(x)$ has rank $r$ at every $x\in U$, so some $r\times r$ minor $\Delta$ (using columns $J$, $|J|=r$) is nonzero at $x$, hence invertible on a neighbourhood. Thus $X$ is covered by finitely many affine opens $W_1,\dots,W_m$ on each of which $\mathcal E$ is trivial and the matrix $M$ has an invertible $r\times r$ minor $M_J$ (columns $J$). Let $Z_a\subseteq W_a\times\mathbf A^N$ be the closed subscheme defined by the $r$ equations $M\cdot v=0$. As $M_J$ is invertible we can solve $v_J=-M_J^{-1}M_{J'}v_{J'}$, so $Z_a\cong W_a\times\mathbf A^{N-r}$ with coordinates $(x,v_{J'})$; $Z_a$ is integral (its ring is a polynomial ring over the domain $\mathcal O(W_a)$) with function field of transcendence degree $n+N-r$ over $k$ ($W_a$ is a nonempty open subset of the variety $X$, so $\operatorname{tr.d.}k(W_a)/k=n$). Let $q:Z_a\to\mathbf A^N$ be the projection and $T_a$ the closure of its image, an integral closed subscheme. The coordinate ring of $T_a$ embeds in $\mathcal O(Z_a)$, so $k(T_a)\subseteq k(Z_a)$ and $\dim T_a=\operatorname{tr.d.}k(T_a)/k\le n+N-r\le N-1$, because $r\ge n+1$. So $T_a$ is a proper closed subset of $\mathbf A^N$, and so is $T_1\cup\dots\cup T_m$ ($\mathbf A^N$ is irreducible). Since $k$ is algebraically closed, the nonempty open complement contains a $k$-point $a$ (Nullstellensatz). Put $s=s_a\in V$.

\emph{Step 2: $s$ works.} Let $x$ be a closed point of $X$, say $x\in W_b$; it is $k$-rational. If $s(x)=0$, then $M(x)a=0$, i.e.\ $(x,a)\in Z_b$, hence $a\in q(Z_b)\subseteq T_b$, contradicting the choice of $a$. So $s(x)\ne0$ for all closed $x$, hence for all $x$ by Step 0.

\emph{The exact sequence.} Define $\sigma:\mathcal O_X\to\mathcal E$ by $1\mapsto s$. Let $x\in X$ and trivialize $\mathcal E$ near $x$, $s=(f_1,\dots,f_r)$. Since $s(x)\ne0$, some $f_i\notin\fm_x$, so $f_i$ is a unit on a neighbourhood $U'$ of $x$. On $U'$ the sections $s,e_j$ ($j\ne i$) form a basis of $\mathcal O_{U'}^r$ (the transition matrix to the standard basis is, up to column order, triangular with the unit $f_i$ on the diagonal). Therefore on $U'$, $\sigma$ is injective and $\mathcal E'=\operatorname{coker}\sigma$ is free of rank $r-1$ (generated by the classes of the $e_j$, $j\ne i$). Thus $0\to\mathcal O_X\xrightarrow{\sigma}\mathcal E\to\mathcal E'\to0$ is exact with $\mathcal E'$ locally free of rank $r-1$.\hfill$\square$

\exercise{8.3}{\emph{Product Schemes.}
\begin{enumerate}[label=(\alph*),leftmargin=*]
\item Let $X$ and $Y$ be schemes over another scheme $S$. Use (8.10) and (8.11) to show that $\Omega_{X\times_SY/S}\cong p_1^*\Omega_{X/S}\oplus p_2^*\Omega_{Y/S}$.
\item If $X$ and $Y$ are nonsingular varieties over a field $k$, show that $\omega_{X\times Y}\cong p_1^*\omega_X\otimes p_2^*\omega_Y$.
\item Let $Y$ be a nonsingular plane cubic curve, and let $X$ be the surface $Y\times Y$. Show that $p_g(X)=1$ but $p_a(X)=-1$ (I, Ex.~7.2). This shows that the arithmetic genus and the geometric genus of a nonsingular projective variety may be different.
\end{enumerate}}
\pf
(a) Let $Z=X\times_SY$ with projections $p_1,p_2$. The morphism $Z\to Y$ is the base extension of $X\to S$ by $Y\to S$, so by (8.10)
\[\Omega_{Z/Y}\cong p_1^*\Omega_{X/S},\qquad\text{and symmetrically}\qquad \Omega_{Z/X}\cong p_2^*\Omega_{Y/S}.\]
By (8.11) applied to $Z\to Y\to S$ and $Z\to X\to S$ we have natural maps $p_2^*\Omega_{Y/S}\to\Omega_{Z/S}\to\Omega_{Z/Y}\cong p_1^*\Omega_{X/S}$ and $p_1^*\Omega_{X/S}\to\Omega_{Z/S}\to\Omega_{Z/X}\cong p_2^*\Omega_{Y/S}$. Let
\[\varphi:p_1^*\Omega_{X/S}\oplus p_2^*\Omega_{Y/S}\to\Omega_{Z/S},\qquad \psi:\Omega_{Z/S}\to\Omega_{Z/Y}\oplus\Omega_{Z/X}\cong p_1^*\Omega_{X/S}\oplus p_2^*\Omega_{Y/S}\]
be the maps built from these. We show $\varphi,\psi$ are inverse isomorphisms; this is local, so let $S=\operatorname{Spec}R$, $X=\operatorname{Spec}A$, $Y=\operatorname{Spec}B$, $C=A\otimes_RB$. Then $\Omega_{C/R}$ is generated as a $C$-module by the elements $d(a\otimes b)=(1\otimes b)\,d(a\otimes1)+(a\otimes1)\,d(1\otimes b)$. Under the identifications of (8.10), $\Omega_{C/B}=\Omega_{A/R}\otimes_AC$ and $\Omega_{C/A}=\Omega_{B/R}\otimes_BC$, the maps are
\[\varphi(da\otimes c)=c\,d(a\otimes1),\quad\varphi(db\otimes c)=c\,d(1\otimes b),\qquad\psi(d\gamma)=(d_{C/B}\gamma,\ d_{C/A}\gamma).\]
Now $\psi\varphi(da\otimes c)=c\,(d_{C/B}(a\otimes1),d_{C/A}(a\otimes 1))=(da\otimes c,\,0)$ since $a\otimes1$ is a constant for $d_{C/A}$; similarly for $db\otimes c$. So $\psi\varphi=\mathrm{id}$. Conversely $\psi(d(a\otimes b))=(da\otimes(1\otimes b),\ db\otimes(a\otimes1))$, so $\varphi\psi(d(a\otimes b))=(1\otimes b)d(a\otimes1)+(a\otimes1)d(1\otimes b)=d(a\otimes b)$; as the $d(a\otimes b)$ generate $\Omega_{C/R}$ and $\varphi\psi$ is $C$-linear, $\varphi\psi=\mathrm{id}$. Since $\varphi$ and $\psi$ are globally defined, they give the isomorphism $\Omega_{X\times_SY/S}\cong p_1^*\Omega_{X/S}\oplus p_2^*\Omega_{Y/S}$.

(b) Let $n=\dim X$, $m=\dim Y$. $Z=X\times_kY$ is integral (II, Ex.~3.15; $k$ algebraically closed), separated and of finite type, so a variety. By (a), $\Omega_{Z/k}\cong p_1^*\Omega_{X/k}\oplus p_2^*\Omega_{Y/k}$ is locally free of rank $n+m$, since $X,Y$ nonsingular (8.15). At the generic point, $\operatorname{rk}\Omega_{Z/k}=\dim_K\Omega_{K/k}=\operatorname{tr.d.}K/k=\dim Z$ ($K$ the function field, separably generated as $k$ is perfect; (8.6A)). Hence $\dim Z=n+m$, and $Z$ is nonsingular by (8.15). Then
\[\omega_Z=\textstyle\bigwedge^{n+m}(p_1^*\Omega_X\oplus p_2^*\Omega_Y)=\bigoplus_{i+j=n+m}\bigwedge^i p_1^*\Omega_X\otimes\bigwedge^jp_2^*\Omega_Y=\bigwedge^np_1^*\Omega_X\otimes\bigwedge^mp_2^*\Omega_Y,\]
because $\bigwedge^i$ of a locally free sheaf of rank $n$ vanishes for $i>n$ (and likewise for $j>m$), so only $(i,j)=(n,m)$ survives (check locally on free modules). As pull-back commutes with exterior powers and tensor products, this is $p_1^*\omega_X\otimes p_2^*\omega_Y$.

(c) $Y\subseteq\mathbf P^2$ is a nonsingular curve of degree $3$, so by (8.20.3) $\omega_Y\cong\mathcal O_Y(3-2-1)=\mathcal O_Y$. By (b), $\omega_X\cong p_1^*\mathcal O_Y\otimes p_2^*\mathcal O_Y=\mathcal O_X$, so $p_g(X)=\dim_k\Gamma(X,\mathcal O_X)$. $X$ is a projective variety (Segre embedding $Y\times Y\subseteq\mathbf P^2\times\mathbf P^2\subseteq\mathbf P^8$), so $\Gamma(X,\mathcal O_X)$ is a finite-dimensional $k$-algebra (II, Thm.~5.19) which is a domain ($X$ integral), hence a finite field extension of $k$, hence $k$. So $p_g(X)=1$.

For $p_a$ we use the Hilbert polynomial of $X\subseteq\mathbf P^8$ (the arithmetic genus is independent of the projective embedding, III, Ex.~5.3; for the definition (I, Ex.~7.2) we simply use this embedding). Let $S_Y=k[x_0,x_1,x_2]/(f)$ with $f$ the cubic; $\dim(S_Y)_m=\binom{m+2}2-\binom{m-1}2=3m$ for $m\ge1$. The homogeneous coordinate ring of $X\subseteq\mathbf P^8$ is the Segre product $R=\bigoplus_m(S_Y)_m\otimes_k(S_Y)_m$: indeed $R$ is a quotient of the Segre ring $\bigoplus_m k[x]_m\otimes k[y]_m=k[z_{ij}]/(\text{Segre relations})$, it is a domain (a subring of $S_Y\otimes_kS_Y$, a domain since $k$ is algebraically closed), so $R=k[z]/J$ with $J$ prime, and the zero set of $J$ is $Y\times Y$ (every element of $J$ vanishes at all points $x\otimes y=(x_iy_j)$ with $x,y$ in the affine cone of $Y$, so $Y\times Y\subseteq V(J)$; conversely $J$ contains the Segre relations, so a point of $V(J)$ has the form $x\otimes y$, and for each $j$ the element $f(z_{0j},z_{1j},z_{2j})$ lies in $J$, because it maps to $y_j^3f(x_0,x_1,x_2)=0$ in $S_Y\otimes S_Y$; hence $f(x)=0$, and symmetrically $f(y)=0$); since $J$ is a homogeneous prime defining $X$, the Nullstellensatz gives $J=I(X)$. Hence $\dim R_m=(3m)^2=9m^2$ for $m\ge1$, so the Hilbert polynomial is $P_X(m)=9m^2$ and $P_X(0)=0$. As $\dim X=2$, $p_a(X)=(-1)^2(P_X(0)-1)=-1$. (Alternatively, by III, Ex.~5.3 and Kunneth, $\chi(\mathcal O_X)=\chi(\mathcal O_Y)^2=0$.) So $p_g(X)=1\ne-1=p_a(X)$.\hfill$\square$

\exercise{8.4}{\emph{Complete Intersections in $\mathbf P^n$.} A closed subscheme $Y$ of $\mathbf P^n_k$ is called a (strict, global) \emph{complete intersection} if the homogeneous ideal $I$ of $Y$ in $S=k[x_0,\dots,x_n]$ can be generated by $r=\operatorname{codim}(Y,\mathbf P^n)$ elements (I, Ex.~2.17).
\begin{enumerate}[label=(\alph*),leftmargin=*]
\item Let $Y$ be a closed subscheme of codimension $r$ in $\mathbf P^n$. Then $Y$ is a complete intersection if and only if there are hypersurfaces (i.e., locally principal subschemes of codimension $1$) $H_1,\dots,H_r$, such that $Y=H_1\cap\dots\cap H_r$ as schemes, i.e., $\mathcal I_Y=\mathcal I_{H_1}+\dots+\mathcal I_{H_r}$. [\emph{Hint}: Use the fact that the unmixedness theorem holds in $S$ (Matsumura [2, p.~107]).]
\item If $Y$ is a complete intersection of dimension $\ge1$ in $\mathbf P^n$, and if $Y$ is normal, then $Y$ is projectively normal (Ex.~5.14). [\emph{Hint}: Apply (8.23) to the affine cone over $Y$.]
\item With the same hypotheses as (b), conclude that for all $l\ge0$, the natural map $\Gamma(\mathbf P^n,\mathcal O_{\mathbf P^n}(l))\to\Gamma(Y,\mathcal O_Y(l))$ is surjective. In particular, taking $l=0$, show that $Y$ is connected.
\item Now suppose given integers $d_1,\dots,d_r\ge1$, with $r<n$. Use Bertini's theorem (8.18) to show that there exist nonsingular hypersurfaces $H_1,\dots,H_r$ in $\mathbf P^n$, with $\deg H_i=d_i$, such that the scheme $Y=H_1\cap\dots\cap H_r$ is irreducible and nonsingular of codimension $r$ in $\mathbf P^n$.
\item If $Y$ is a nonsingular complete intersection as in (d), show that $\omega_Y\cong\mathcal O_Y(\sum d_i-n-1)$.
\item If $Y$ is a nonsingular hypersurface of degree $d$ in $\mathbf P^n$, use (c) and (e) above to show that $p_g(Y)=\binom{d-1}{n}$. Thus $p_g(Y)=p_a(Y)$ (I, Ex.~7.2). In particular, if $Y$ is a nonsingular plane curve of degree $d$, then $p_g(Y)=\frac12(d-1)(d-2)$.
\item If $Y$ is a nonsingular curve in $\mathbf P^3$, which is a complete intersection of nonsingular surfaces of degrees $d,e$, then $p_g(Y)=\frac12de(d+e-4)+1$. Again the geometric genus is the same as the arithmetic genus (I, Ex.~7.2).
\end{enumerate}}
\pf
\emph{Conventions.} $k$ is algebraically closed, $S=k[x_0,\dots,x_n]$, $\fn=(x_0,\dots,x_n)$. The ``homogeneous ideal'' of a closed subscheme $Y$ is $I_Y=\Gamma_*(\mathcal I_Y)=\bigoplus_l\Gamma(\mathbf P^n,\mathcal I_Y(l))$ (II, Ex.~5.9, 5.10): the saturated ideal, with $\widetilde{I_Y}=\mathcal I_Y$ and $Y=\operatorname{Proj}S/I_Y$. We use: \emph{if $J$ is any homogeneous ideal with $\widetilde J=\mathcal I_Y$, then $I_Y=\bigcup_m(J:\fn^m)$; in particular $J=I_Y$ if $\fn\notin\operatorname{Ass}(S/J)$.} (By Ex.~5.9, $J_l=(I_Y)_l$ for $l\gg0$; so for $f\in I_Y$ homogeneous, $\fn^mf\subseteq J$ for some $m$, whence $f\in J$ if $\fn$ is not an associated prime of $S/J$.) A hypersurface $H$ has invertible ideal sheaf, $\mathcal I_H\cong\mathcal O(-d)$ for some $d>0$ (Pic $\mathbf P^n=\mathbf Z$), and the inclusion $\mathcal O(-d)\to\mathcal O$ is multiplication by a nonzero form $F\in S_d=\Gamma(\mathcal O(d))$; thus $H=\operatorname{Proj}S/(F)$, $\mathcal I_H=\widetilde{(F)}$, and conversely every nonconstant form defines a hypersurface (codimension $1$ by Krull). For an irreducible closed $Z\subseteq\mathbf P^n$, the homogeneous prime $\fp_Z$ has $\operatorname{ht}\fp_Z=n-\dim Z=\operatorname{codim}(Z,\mathbf P^n)$, so $\operatorname{ht}I_Y=\operatorname{codim}(Y,\mathbf P^n)$ for $Y\ne\emptyset$.

(a) ($\Rightarrow$) Let $I=I_Y$ be generated by $r$ elements. The $k$-vector space $I/\fn I$ is generated by the images of these, so by the graded Nakayama lemma $I$ is generated by $r'\le r$ homogeneous elements $F_1,\dots,F_{r'}$ (lifts of a basis of $I/\fn I$). By Krull's height theorem $r=\operatorname{ht}I\le r'$, so $r'=r$. No $F_i$ is $0$ (else $r-1$ generators) or a nonzero constant (else $Y=\emptyset$). So $H_i=\operatorname{Proj}S/(F_i)$ are hypersurfaces, and $\mathcal I_Y=\widetilde I=\sum\widetilde{(F_i)}=\sum\mathcal I_{H_i}$.

($\Leftarrow$) Suppose $\mathcal I_Y=\sum\mathcal I_{H_i}$, $H_i=\operatorname{Proj}S/(F_i)$ with $\deg F_i=d_i\geq 1$, and $\operatorname{codim}Y=r$. Put $J=(F_1,\dots,F_r)$; then $\widetilde J=\mathcal I_Y$. As $Y\ne\emptyset$ we have $r\le n$. By Krull every minimal prime of $J$ has height $\le r$. The zero set of $J$ in $\mathbf A^{n+1}$ is the cone over $Y$ together with the origin; every component of $V(J)$ has dimension $\ge n+1-r\ge1$, so the origin is not a component, and the components are the cones over the components of $Y$; the latter have $\dim\le n-r$ (as $\operatorname{codim}Y=r$) and $\ge n-r$ (Krull). Hence all minimal primes of $J$ have height exactly $r$, and $\operatorname{ht}J=r$. Now $S$ is Cohen--Macaulay (regular, (8.21Aa)) and $J$ is generated by $\operatorname{ht}J$ elements, so by the unmixedness theorem every associated prime of $S/J$ has height $r\le n<n+1=\operatorname{ht}\fn$. Hence $\fn\notin\operatorname{Ass}(S/J)$ and $J=I_Y$ is generated by $r$ elements: $Y$ is a complete intersection. (This also shows that, for a complete intersection, $I_Y=(F_1,\dots,F_r)$ with $F_1,\dots,F_r$ forming a regular sequence.)

(b) Let $I=I_Y=(F_1,\dots,F_r)$ (by (a), with homogeneous $F_i$), $S'=S/I$, and let $C=\operatorname{Spec}S'\subseteq\mathbf A^{n+1}$ be the affine cone, $\mathfrak n'=\fn/I$ the vertex. By the proof of (a), $C$ has pure dimension $n+1-r=\dim Y+1\ge2$, i.e.\ $\operatorname{codim}(C,\mathbf A^{n+1})=r$ and $I$ is generated by $r$ elements. So $C$ is a locally complete intersection in the nonsingular variety $\mathbf A^{n+1}$ and (8.23) applies: $C$ is Cohen--Macaulay, and $C$ is normal iff it is regular in codimension~$1$. Moreover $\dim S'_{\mathfrak n'}=n+1-r\ge2$.

\emph{$C\setminus\{\mathfrak n'\}$ is normal.} It is covered by $D(x_j)$, with ring $S'_{x_j}\cong A_j[t,t^{-1}]$, $A_j=(S'_{x_j})_0=\mathcal O(Y\cap D_+(x_j))$, $t=x_j$ (every element of $S'_{x_j}$ is a finite sum of homogeneous parts, and $x_j$ is transcendental over the degree-$0$ part). Now $Y$ is normal, so $A_j$ is a (finite product of) integrally closed domain(s), and $A_j[t,t^{-1}]$ is again a normal ring. Hence $C\setminus\{\mathfrak n'\}$ is normal; in particular it is regular in codimension~$1$ (normal noetherian local rings of dimension $1$ are DVRs). Points of codimension $1$ of $C$ lie in $C\setminus\{\mathfrak n'\}$ because $\dim S'_{\mathfrak n'}\ge2$. Therefore $C$ is regular in codimension~$1$, hence normal by (8.23). So $S'$ is a normal ring; $\operatorname{Spec}S'$ is connected (every minimal prime of the homogeneous ideal $I$ is homogeneous, hence contained in $\mathfrak n'$), and a normal noetherian ring is a finite product of normal domains, so $S'$ is a domain. Thus $S'=S/I_Y$ is an integrally closed domain: $Y$ is projectively normal.

(c) By (b), $S'$ is a normal noetherian domain, and $U=C\setminus\{\mathfrak n'\}=\bigcup D(x_j)$ contains all height-one primes. By (6.3A), $S'=\bigcap_{\operatorname{ht}\fp=1}S'_\fp$ inside $K=\operatorname{Frac}S'$. A section of $\mathcal O_U$ lies in each $S'_\fp$ ($\operatorname{ht}\fp=1$), so $\Gamma(U,\mathcal O_U)=S'$. On the other hand, $\Gamma(U,\mathcal O_U)=\bigoplus_{l\in\mathbf Z}\Gamma(Y,\mathcal O_Y(l))$: a section is a family $(s_j)$, $s_j\in S'_{x_j}$, compatible on overlaps; all localization maps are graded, so the homogeneous components of degree $l$ of the $s_j$ form again a compatible family, i.e.\ a section of $\mathcal O_Y(l)$ (a degree-$l$ element of $S'_{x_j}$ is exactly a section of $\mathcal O_Y(l)$ over $D_+(x_j)\cap Y$; here $Y=\operatorname{Proj}S'$ and $\mathcal O_Y(l)=\mathcal O_{\mathbf P^n}(l)|_Y$), and only finitely many degrees occur since there are finitely many $j$. The isomorphism $S'\to\Gamma(U,\mathcal O_U)$ is graded, so $S'_l\to\Gamma(Y,\mathcal O_Y(l))$ is bijective for all $l\ge0$. As $S_l=\Gamma(\mathbf P^n,\mathcal O(l))\to S'_l$ is surjective, so is $\Gamma(\mathbf P^n,\mathcal O(l))\to\Gamma(Y,\mathcal O_Y(l))$.

For $l=0$: $\Gamma(Y,\mathcal O_Y)=S'_0$ is a quotient of $S_0=k$, so it equals $k$ ($Y\ne\emptyset$). If $Y=Y_1\sqcup Y_2$ with $Y_1,Y_2$ nonempty and closed, then $\Gamma(Y,\mathcal O_Y)=\Gamma(Y_1,\mathcal O)\times\Gamma(Y_2,\mathcal O)$, a product of two nonzero $k$-algebras, of $k$-dimension $\ge2$, contradiction. So $Y$ is connected.

(d) \emph{Setup.} Let $v_d:\mathbf P^n\to\mathbf P^{N}$ be the $d$-uple embedding. Hyperplanes $H'\subset\mathbf P^N$ correspond bijectively (linear isomorphism of projective spaces) to degree-$d$ forms $F$ up to scalars, via $v_d^{-1}(H')=\operatorname{Proj}S/(F)$ (the coordinates of $\mathbf P^N$ are the monomials of degree $d$). If $X\subseteq\mathbf P^n$ is a nonsingular closed subvariety, Bertini (8.18) applied to $v_d(X)\subseteq\mathbf P^N$ gives a nonempty open set $\mathcal U_d(X)$ of forms $F$ of degree $d$ such that $X\not\subseteq V(F)$ and the scheme $X\cap\operatorname{Proj}S/(F)$ is regular at every point ($v_d$ is an isomorphism onto its image, so scheme-theoretic intersections correspond).

\emph{Induction.} We construct $H_1,\dots,H_r$ with $H_i=\operatorname{Proj}S/(F_i)$, $\deg F_i=d_i$, such that, with $Y_i=H_1\cap\dots\cap H_i$ ($Y_0=\mathbf P^n$): $H_i$ is nonsingular, and $Y_i$ is a nonsingular irreducible closed subscheme of codimension $i$ in $\mathbf P^n$. Suppose $Y_{i-1}$ is constructed ($i\le r<n$; $Y_0=\mathbf P^n$). The open sets $\mathcal U_{d_i}(\mathbf P^n)$ and $\mathcal U_{d_i}(Y_{i-1})$ of the irreducible space $\mathbf P(S_{d_i})$ are nonempty, hence meet; pick $F_i$ in the intersection. Then $H_i$ is nonsingular and $Y_i=Y_{i-1}\cap H_i$ is regular. Since $Y_{i-1}$ is an irreducible projective variety of dimension $n-i+1\ge1$ not contained in $H_i$, every component of $Y_i$ has dimension exactly $n-i$ (Krull's Hauptidealsatz and (I, Thm.~7.2)), and $Y_i\ne\emptyset$ (a hypersurface meets every projective variety of positive dimension, (I, Thm.~7.2)). Thus $Y_i=H_1\cap\dots\cap H_i$ has codimension $i$, and by (a) it is a complete intersection. Being regular, it is normal, and $\dim Y_i=n-i\ge n-r\ge1$; so by (c) $Y_i$ is connected. A connected regular scheme is irreducible (its local rings are domains, so distinct irreducible components cannot meet). This completes the induction; $Y=Y_r$ is as required.

(e) Let $Y=H_1\cap\dots\cap H_r$, $H_i=\operatorname{Proj}S/(F_i)$, $\deg F_i=d_i$, be a nonsingular complete intersection of codimension $r$ (nonsingular subvariety of $\mathbf P^n$). The maps $\mathcal O(-d_i)\to\mathcal I_Y$, $1\mapsto F_i$, give a map $\bigoplus\mathcal O_Y(-d_i)\to\mathcal I_Y/\mathcal I_Y^2$, surjective because $\mathcal I_Y=\sum F_i\mathcal O$. By (8.17) the target is locally free of rank $r$, and so is the source; a surjection of locally free sheaves of the same finite rank is an isomorphism (on stalks it is a surjective homomorphism $\mathcal O_y^r\to\mathcal O_y^r$ of finitely generated modules, hence bijective). Hence $\mathcal N_{Y/\mathbf P^n}=\mathcal{H}om(\mathcal I/\mathcal I^2,\mathcal O_Y)\cong\bigoplus\mathcal O_Y(d_i)$ and $\bigwedge^r\mathcal N_{Y/\mathbf P^n}\cong\mathcal O_Y(\sum d_i)$. By (8.20) and (8.20.1),
\[\omega_Y\cong\omega_{\mathbf P^n}\otimes\textstyle\bigwedge^r\mathcal N_{Y/\mathbf P^n}\cong\mathcal O_Y(-n-1)\otimes\mathcal O_Y(\sum d_i)=\mathcal O_Y(\sum d_i-n-1).\]

(f) Let $n\ge2$ and $Y=\operatorname{Proj}S/(F)$ nonsingular of degree $d$ (a complete intersection with $r=1$; $\dim Y=n-1\ge1$; $Y$ regular hence normal). By (e), $\omega_Y\cong\mathcal O_Y(d-n-1)$, so $p_g(Y)=\dim\Gamma(Y,\mathcal O_Y(l))$ with $l=d-n-1$. By the proof of (c), $\Gamma(Y,\mathcal O_Y(l))=S'_l$ for all $l\in\mathbf Z$, where $S'=S/(F)$ (this is $0$ for $l<0$); so $\dim\Gamma(Y,\mathcal O_Y(l))=\dim S_l-\dim S_{l-d}$. Here $l-d=-n-1<0$, so $S_{l-d}=0$, and
\[p_g(Y)=\dim S_{d-n-1}=\tbinom{d-1}{n}\]
(for $d-n-1\ge0$ this is $\binom{d-1}n$; for $d\le n$ both sides are $0$). On the other hand the Hilbert polynomial of $Y$ is $P_Y(m)=\binom{m+n}n-\binom{m-d+n}n$, so $P_Y(0)=1-\binom{n-d}{n}$ (polynomial binomial), and since $\binom{n-d}n=(-1)^n\binom{d-1}n$,
\[p_a(Y)=(-1)^{n-1}(P_Y(0)-1)=(-1)^{n}\tbinom{n-d}n=\tbinom{d-1}n.\]
So $p_g(Y)=p_a(Y)$. For $n=2$ this is $\frac12(d-1)(d-2)$.

(g) Let $Y=H_1\cap H_2$, $H_1=V(F)$, $H_2=V(G)$, $\deg F=d$, $\deg G=e$, a nonsingular curve in $\mathbf P^3$. By (a), $I_Y=(F,G)$. By (e), $\omega_Y\cong\mathcal O_Y(d+e-4)$, and by the proof of (c) (here $Y$ is normal of dimension $1$) $\Gamma(Y,\mathcal O_Y(l))=(S/I_Y)_l$ for all $l\in\mathbf Z$; hence, with $l=d+e-4$, $p_g(Y)=\dim(S/(F,G))_l$. The forms $F,G$ are coprime (a common nonconstant factor $\pi$ would give $V(\pi)\subseteq Y$ with $\dim V(\pi)=2$), so $G$ is a nonzerodivisor on $S/(F)$ ($S$ is a UFD: if $F\mid GH$ and $F\nmid H$ some prime power $\pi^a\,\|\,F$ does not divide $H$, so $\pi\mid G$). Hence $0\to S(-d-e)\xrightarrow{(G,-F)}S(-d)\oplus S(-e)\xrightarrow{(F,G)}S\to S/(F,G)\to0$ is exact, and
\[\dim(S/(F,G))_l=\sigma(l)-\sigma(l-d)-\sigma(l-e)+\sigma(l-d-e),\qquad\sigma(m)=\dim S_m=\tbinom{m+3}3\ (m\ge0),\ 0\ (m<0).\]
At $l=d+e-4$ the last term is $\sigma(-4)=0$, and the sum is $\sigma(d+e-4)-\sigma(e-4)-\sigma(d-4)$ (also correct if $l<0$). Since $d,e\ge1$ we have $\sigma(m)=\frac{(m+3)(m+2)(m+1)}6$ for all $m\ge-3$ (the polynomial vanishes at $-1,-2,-3$), so with $s=d+e$, $p=de$:
\[p_g(Y)=\frac{(s-1)(s-2)(s-3)-(e-1)(e-2)(e-3)-(d-1)(d-2)(d-3)}6=\frac{3p(s-4)+6}6=\tfrac12de(d+e-4)+1,\]
where we used $(s-1)(s-2)(s-3)=s^3-6s^2+11s-6$, $d^3+e^3=s^3-3ps$, $d^2+e^2=s^2-2p$. Finally, for $m\gg0$, $P_Y(m)=\sigma(m)-\sigma(m-d)-\sigma(m-e)+\sigma(m-d-e)$ as polynomials, so $P_Y(0)=1+\frac{(d-1)(d-2)(d-3)+(e-1)(e-2)(e-3)-(s-1)(s-2)(s-3)}6=-\frac12p(s-4)$ (use $\binom{3-d}3=-\frac{(d-1)(d-2)(d-3)}6$ etc.), and $p_a(Y)=(-1)^1(P_Y(0)-1)=1+\frac12de(d+e-4)=p_g(Y)$.\hfill$\square$

\exercise{8.5}{\emph{Blowing up a Nonsingular Subvariety.} As in (8.24), let $X$ be a nonsingular variety. Let $Y$ be a nonsingular subvariety of codimension $r\ge2$, let $\pi:\widetilde X\to X$ be the blow-up of $X$ along $Y$, and let $Y'=\pi^{-1}(Y)$.
\begin{enumerate}[label=(\alph*),leftmargin=*]
\item Show that the maps $\pi^*:\operatorname{Pic}X\to\operatorname{Pic}\widetilde X$, and $\mathbf Z\to\operatorname{Pic}\widetilde X$ defined by $n\mapsto$ class of $nY'$, give rise to an isomorphism $\operatorname{Pic}\widetilde X\cong\operatorname{Pic}X\oplus\mathbf Z$.
\item Show that $\omega_{\widetilde X}\cong f^*\omega_X\otimes\mathcal L((r-1)Y')$. [\emph{Hint}: By (a) we can write in any case $\omega_{\widetilde X}\cong f^*\mathcal M\otimes\mathcal L(qY')$ for some invertible sheaf $\mathcal M$ on $X$, and some integer $q$. By restricting to $\widetilde X-Y'\cong X-Y$, show that $\mathcal M\cong\omega_X$. To determine $q$, proceed as follows. First show that $\omega_{Y'}\cong f^*\omega_X\otimes\mathcal O_{Y'}(-q-1)$. Then take a closed point $y\in Y$ and let $Z$ be the fibre of $Y'$ over $y$. Then show that $\omega_Z\cong\mathcal O_Z(-q-1)$. But since $Z\cong\mathbf P^{r-1}$, we have $\omega_Z\cong\mathcal O_Z(-r)$, so $q=r-1$.]
\end{enumerate}}
\pf
(The book writes $f$ for $\pi$ in (b); $\mathcal L(D)=\mathcal O(D)$.) By (8.24), $\widetilde X$ is a nonsingular variety (of dimension $n=\dim X$), $Y'\cong\mathbf P(\mathcal I/\mathcal I^2)$ is a $\mathbf P^{r-1}$-bundle over $Y$, locally $Y\times\mathbf P^{r-1}$, hence a nonsingular (reduced) subvariety; it is irreducible (it is covered by the opens $Y_i\times\mathbf P^{r-1}$, which are integral, and any two of them meet since $Y$ is irreducible and $Y'\to Y$ is surjective). Its ideal sheaf $\mathcal I_{Y'}=\pi^{-1}\mathcal I\cdot\mathcal O_{\widetilde X}$ is invertible, so $Y'$ is a prime divisor with $\mathcal I_{Y'}=\mathcal O(-Y')$; and by (8.24c) the normal sheaf $\mathcal N_{Y'/\widetilde X}=\mathcal O(Y')|_{Y'}$ corresponds to $\mathcal O_{\mathbf P(\mathcal I/\mathcal I^2)}(-1)$. Hence
\[\mathcal O(nY')|_{Y'}\cong\mathcal O_{Y'}(-n),\tag{1}\]
and its restriction to a fibre $Z\cong\mathbf P^{r-1}$ is $\mathcal O_{\mathbf P^{r-1}}(-n)$. Put $U=\widetilde X\setminus Y'\cong X\setminus Y$.

(a) All of $X,\widetilde X,U$ are nonsingular, so $\operatorname{Pic}=\operatorname{Cl}$ (6.16), compatibly with restriction to open sets. Since $\operatorname{codim}(Y,X)\ge2$, restriction $\operatorname{Pic}X\to\operatorname{Pic}(X\setminus Y)$ is an isomorphism (6.5(b)). Since $Y'$ is an irreducible closed subset of codimension $1$, (6.5(c)) gives an exact sequence
\[\mathbf Z\xrightarrow{\,n\mapsto nY'\,}\operatorname{Pic}\widetilde X\xrightarrow{\ \rho\ }\operatorname{Pic}U\to0.\]
The composite $\rho\circ\pi^*:\operatorname{Pic}X\to\operatorname{Pic}U=\operatorname{Pic}(X\setminus Y)$ is the restriction map, because $\pi|_U:U\to X\setminus Y$ is an isomorphism; so it is an isomorphism. Consider $\Phi:\operatorname{Pic}X\oplus\mathbf Z\to\operatorname{Pic}\widetilde X$, $(L,n)\mapsto\pi^*L\otimes\mathcal O(nY')$, a homomorphism.

\emph{Injective:} if $\pi^*L\otimes\mathcal O(nY')\cong\mathcal O$, restrict to $U$: $\rho\pi^*L=L|_{X\setminus Y}$ is trivial (as $\mathcal O(nY')|_U$ is trivial), so $L$ is trivial, so $\mathcal O(nY')$ is trivial. Restricting to $Y'$ and a fibre $Z\cong\mathbf P^{r-1}$ and using (1), $\mathcal O_{\mathbf P^{r-1}}(-n)$ is trivial; since $r-1\ge1$ this forces $n=0$ (for $n>0$, $\Gamma(\mathcal O(-n))=0\neq\Gamma(\mathcal O)$; for $n<0$, $\dim\Gamma(\mathcal O(-n))>1$).

\emph{Surjective:} let $\mathcal N\in\operatorname{Pic}\widetilde X$. Choose $L\in\operatorname{Pic}X$ with $L|_{X\setminus Y}\cong\mathcal N|_U$ (possible since $\rho\pi^*$ is onto). Then $\mathcal N\otimes\pi^*L^{-1}\in\ker\rho$, which is generated by $\mathcal O(Y')$, so it is $\mathcal O(nY')$ for some $n$, and $\mathcal N=\Phi(L,n)$. Hence $\Phi$ is an isomorphism.

(b) By (a), $\omega_{\widetilde X}\cong\pi^*\mathcal M\otimes\mathcal O(qY')$ for some $\mathcal M\in\operatorname{Pic}X$, $q\in\mathbf Z$. Restrict to $U\cong X\setminus Y$: $\omega_{\widetilde X}|_U=\omega_U$ is identified with $\omega_X|_{X\setminus Y}$ under $\pi|_U$, while $\pi^*\mathcal M|_U=\mathcal M|_{X\setminus Y}$ and $\mathcal O(qY')|_U$ is trivial. So $\mathcal M|_{X\setminus Y}\cong\omega_X|_{X\setminus Y}$ and, since $\operatorname{Pic}X\to\operatorname{Pic}(X\setminus Y)$ is injective (by (a)), $\mathcal M\cong\omega_X$. Therefore $\omega_{\widetilde X}\cong\pi^*\omega_X\otimes\mathcal O(qY')$.

\emph{Determination of $q$.} By (8.20) with codimension $1$ and by (1),
\[\omega_{Y'}\cong\omega_{\widetilde X}|_{Y'}\otimes\mathcal O(Y')|_{Y'}\cong(\pi^*\omega_X)|_{Y'}\otimes\mathcal O_{Y'}(-q)\otimes\mathcal O_{Y'}(-1)=(\pi^*\omega_X)|_{Y'}\otimes\mathcal O_{Y'}(-q-1).\]
Let $y\in Y$ be a closed point, $g:Y'\to Y$ the projection and $Z=g^{-1}(y)\cong\mathbf P^{r-1}$ the fibre. Over a neighbourhood $W$ of $y$ in $Y$ we have $g^{-1}W\cong W\times\mathbf P^{r-1}$ and $Z=\{y\}\times\mathbf P^{r-1}$. If $t_1,\dots,t_s$ ($s=\dim Y$) generate the maximal ideal of $y$ on $W$, then $\mathcal I_Z=(t_1,\dots,t_s)\mathcal O_{g^{-1}W}$; $Z$ is a nonsingular subvariety of $Y'$ of codimension $s$, so $\mathcal I_Z/\mathcal I_Z^2$ is locally free of rank $s$ (8.17), and it is generated by the $s$ global sections $\bar t_i$, hence free on them. So $\mathcal N_{Z/Y'}\cong\mathcal O_Z^s$ and, by (8.20), $\omega_Z\cong\omega_{Y'}|_Z\otimes\bigwedge^s\mathcal N_{Z/Y'}\cong\omega_{Y'}|_Z$. Moreover $Z$ is mapped by $\pi$ to the point $y$, so $(\pi^*\omega_X)|_Z$ is the pull-back of a line bundle on $\operatorname{Spec}k(y)$, i.e.\ trivial. Hence
\[\omega_Z\cong\mathcal O_Z(-q-1).\]
But $Z\cong\mathbf P^{r-1}$, so $\omega_Z\cong\mathcal O_Z(-r)$ (8.20.1). As $\operatorname{Pic}\mathbf P^{r-1}=\mathbf Z$ (generated by $\mathcal O(1)$) and $r-1\ge1$, $q+1=r$, i.e.\ $q=r-1$. Thus $\omega_{\widetilde X}\cong\pi^*\omega_X\otimes\mathcal O((r-1)Y')$.\hfill$\square$

\exercise{8.6}{\emph{The Infinitesimal Lifting Property.} The following result is very important in studying deformations of nonsingular varieties. Let $k$ be an algebraically closed field, let $A$ be a finitely generated $k$-algebra such that $\operatorname{Spec}A$ is a nonsingular variety over $k$. Let $0\to I\to B'\to B\to0$ be an exact sequence, where $B'$ is a $k$-algebra, and $I$ is an ideal with $I^2=0$. Finally suppose given a $k$-algebra homomorphism $f:A\to B$. Then there exists a $k$-algebra homomorphism $g:A\to B'$ making a commutative diagram
\[\begin{array}{ccc}&&0\\&&\downarrow\\&&I\\&&\downarrow\\&&B'\\&\overset{g}{\nearrow}&\downarrow\\A&\xrightarrow{\ f\ }&B\\&&\downarrow\\&&0\end{array}\]
We call this result the \emph{infinitesimal lifting property} for $A$. We prove this result in several steps.
\begin{enumerate}[label=(\alph*),leftmargin=*]
\item First suppose that $g:A\to B'$ is a given homomorphism lifting $f$. If $g':A\to B'$ is another such homomorphism, show that $\theta=g-g'$ is a $k$-derivation of $A$ into $I$, which we can consider as an element of $\Hom_A(\Omega_{A/k},I)$. Note that since $I^2=0$, $I$ has a natural structure of $B$-module, and hence of $A$-module. Conversely, for any $\theta\in\Hom_A(\Omega_{A/k},I)$, $g'=g+\theta$ is another homomorphism lifting $f$. (For this step, you do not need the hypothesis about $\operatorname{Spec}A$ being nonsingular.)
\item Now let $P=k[x_1,\dots,x_n]$ be a polynomial ring over $k$ of which $A$ is a quotient, and let $J$ be the kernel. Show that there does exist a homomorphism $h:P\to B'$ making a commutative diagram
\[\begin{array}{ccccccccc}0&\to&J&\to&P&\to&A&\to&0\\&&\downarrow&&\downarrow{\scriptstyle h}&&\downarrow{\scriptstyle f}&&\\0&\to&I&\to&B'&\to&B&\to&0\end{array}\]
and that $h$ induces an $A$-linear map $\bar h:J/J^2\to I$.
\item Now use the hypothesis $\operatorname{Spec}A$ nonsingular and (8.17) to obtain an exact sequence
\[0\to J/J^2\to\Omega_{P/k}\otimes_PA\to\Omega_{A/k}\to0.\]
Show furthermore that applying the functor $\Hom_A(\,\cdot\,,I)$ gives an exact sequence
\[0\to\Hom_A(\Omega_{A/k},I)\to\Hom_P(\Omega_{P/k},I)\to\Hom_A(J/J^2,I)\to0.\]
Let $\theta\in\Hom_P(\Omega_{P/k},I)$ be an element whose image gives $\bar h\in\Hom_A(J/J^2,I)$. Consider $\theta$ as a derivation of $P$ into $B'$. Let $h'=h-\theta$. Then show that $h'$ is a homomorphism $P\to B'$ such that $h'(J)=0$. Thus $h'$ induces the desired homomorphism $g:A\to B'$.
\end{enumerate}}
\pf
Throughout, $I$ is annihilated by $I$, so multiplication by $b'\in B'$ on $I$ depends only on the class of $b'$ in $B=B'/I$; this makes $I$ a $B$-module, an $A$-module via $f$, and a $P$-module via $P\to A\to B$.

(a) Let $g,g':A\to B'$ both lift $f$. Then $\theta=g-g'$ takes values in $\ker(B'\to B)=I$, is additive and $k$-linear, and for $a,a'\in A$,
\[\theta(aa')=g(a)g(a')-g'(a)g'(a')=g(a)\theta(a')+\theta(a)g'(a')=a\,\theta(a')+a'\,\theta(a),\]
since $g(a)$ and $g'(a')$ act on $I$ as multiplication by $f(a)$ and $f(a')$, i.e.\ by $a$ and $a'$ in the $A$-module structure. So $\theta\in\operatorname{Der}_k(A,I)=\Hom_A(\Omega_{A/k},I)$. Conversely let $\theta\in\operatorname{Der}_k(A,I)$ and $g'=g+\theta$. Then $g'$ is additive, $k$-linear, $g'(1)=1$ (a derivation kills $1$), lifts $f$ (as $\theta(A)\subseteq I$), and
\[g'(a)g'(a')=g(a)g(a')+g(a)\theta(a')+\theta(a)g(a')+\theta(a)\theta(a')=g(aa')+a\theta(a')+a'\theta(a)=g'(aa')\]
because $\theta(a)\theta(a')\in I^2=0$. So $g'$ is a $k$-algebra homomorphism lifting $f$.

(b) Let $\bar x_i$ be the image of $x_i$ in $A$ and choose $b_i'\in B'$ lifting $f(\bar x_i)\in B$. Let $h:P\to B'$ be the $k$-algebra homomorphism with $h(x_i)=b_i'$. The composite $P\to B'\to B$ and $P\to A\xrightarrow fB$ are $k$-algebra homomorphisms agreeing on the $x_i$, hence equal: the right square commutes. Consequently $h(J)\subseteq I$, giving the left vertical map $J\to I$. As $h(J^2)\subseteq I^2=0$, $h$ induces $\bar h:J/J^2\to I$. It is $A$-linear: for $p\in P$, $j\in J$, $h(pj)=h(p)h(j)=\bar p\cdot h(j)$, where $\bar p$ is the image of $p$ in $A$ (acting on $I$ via $f$).

(c) $X=\mathbf A^n_k=\operatorname{Spec}P$ is a nonsingular variety and $Y=\operatorname{Spec}A$ is a nonsingular closed subvariety with ideal $\widetilde J$. By (8.17) (taking global sections on the affine $Y$) we get the exact sequence $0\to J/J^2\xrightarrow{\ \delta\ }\Omega_{P/k}\otimes_PA\to\Omega_{A/k}\to0$ ($\delta(\bar j)=dj\otimes1$), and $\Omega_{A/k}$ is locally free of finite rank (8.15), hence a projective $A$-module. Therefore this sequence splits, and applying $\Hom_A(\,\cdot\,,I)$ gives the exact sequence
\[0\to\Hom_A(\Omega_{A/k},I)\to\Hom_A(\Omega_{P/k}\otimes_PA,I)\to\Hom_A(J/J^2,I)\to0,\]
where $\Hom_A(\Omega_{P/k}\otimes_PA,I)=\Hom_P(\Omega_{P/k},I)$, and the last map is $\theta\mapsto\theta\circ\delta$, i.e.\ $j\mapsto\theta(dj)$ for $j\in J$.

Choose $\theta\in\Hom_P(\Omega_{P/k},I)$ mapping to $\bar h$: $\theta(dj)=h(j)$ for all $j\in J$. Let $D=\theta\circ d:P\to I$, a $k$-derivation of $P$ into $I$ (the $P$-module structure on $I$ is the one above). Apply the converse part of (a) with $P$ in place of $A$, the map $P\to B$ in place of $f$, and the lift $h$: since $-D\in\operatorname{Der}_k(P,I)$, $h'=h-D$ is a $k$-algebra homomorphism $P\to B'$ lifting $P\to B$. For $j\in J$,
\[h'(j)=h(j)-\theta(dj)=h(j)-h(j)=0.\]
So $h'(J)=0$ and $h'$ factors through a $k$-algebra homomorphism $g:A=P/J\to B'$. It lifts $f$: $g(\bar p)=h'(p)\equiv h(p)\equiv f(\bar p)\pmod I$. This proves the infinitesimal lifting property.\hfill$\square$

\exercise{8.7}{As an application of the infinitesimal lifting property, we consider the following general problem. Let $X$ be a scheme of finite type over $k$, and let $\mathcal F$ be a coherent sheaf on $X$. We seek to classify schemes $X'$ over $k$, which have a sheaf of ideals $\mathcal I$ such that $\mathcal I^2=0$ and $(X',\mathcal O_{X'}/\mathcal I)\cong(X,\mathcal O_X)$, and such that $\mathcal I$ with its resulting structure of $\mathcal O_X$-module is isomorphic to the given sheaf $\mathcal F$. We call such an $X'$, together with $\mathcal I$, an \emph{infinitesimal extension} of the scheme $X$ by the sheaf $\mathcal F$. One such extension, the \emph{trivial} one, is obtained as follows. Take $\mathcal O_{X'}=\mathcal O_X\oplus\mathcal F$ as sheaves of abelian groups, and define multiplication by $(a\oplus f)\cdot(a'\oplus f')=aa'\oplus(af'+a'f)$. Then the topological space $X$ with the sheaf of rings $\mathcal O_{X'}$ is an infinitesimal extension of $X$ by $\mathcal F$.

The general problem of classifying extensions of $X$ by $\mathcal F$ can be quite complicated. So for now, just prove the following special case: if $X$ is affine and nonsingular, then any extension of $X$ by a coherent sheaf $\mathcal F$ is isomorphic to the trivial one. (See (III, Ex.~4.10) for another case.)}
\pf
(Here, as in Ex.~8.6, $k$ is algebraically closed and ``nonsingular'' means that $X$ is a nonsingular variety.) Let $X=\operatorname{Spec}A$ with $A$ a finitely generated $k$-algebra, and let $(X',\mathcal I)$ be an extension of $X$ by $\mathcal F$. The underlying space of $X'$ is $|X|$, and we have an exact sequence of sheaves of $k$-algebras
\[0\to\mathcal F\to\mathcal O_{X'}\xrightarrow{\ \pi\ }\mathcal O_X\to0,\]
with $\mathcal F=\mathcal I$, $\mathcal F^2=0$, and the $\mathcal O_X$-module structure on $\mathcal F$ given by $b'\cdot f=\pi(b')f$.

\emph{Step 1: a ring-theoretic section.} Put $B'=\Gamma(X,\mathcal O_{X'})$ and $I=\Gamma(X,\mathcal F)$. Since $X$ is affine and $\mathcal F$ is quasi-coherent, $H^1(X,\mathcal F)=0$ (III, Thm.~3.5; cohomology of a sheaf of abelian groups on the space $X$ is the same whether computed in $\mathfrak{Ab}(X)$ or in $\mathcal O_X$-modules, III, Prop.~2.6). Thus the long exact cohomology sequence gives an exact sequence $0\to I\to B'\to A\to0$ of $k$-algebras, with $I^2=0$. Apply Ex.~8.6 to the finitely generated $k$-algebra $A$ (with $\operatorname{Spec}A$ a nonsingular variety), the quotient $B'\to B:=A$ and $f=\mathrm{id}_A$: there is a $k$-algebra homomorphism $g:A\to B'$ with $\pi\circ g=\mathrm{id}_A$.

\emph{Step 2: a section of sheaves of rings.} For $a\in A$ the element $g(a)|_{D(a)}\in\mathcal O_{X'}(D(a))$ maps to the unit $a\in\mathcal O_X(D(a))=A_a$. An element $x\in\mathcal O_{X'}(U)$ whose image in $\mathcal O_X(U)$ is a unit is itself a unit: at each stalk, the image of $x_z$ is a unit, so (surjectivity of $\mathcal O_{X',z}\to\mathcal O_{X,z}$) there is $y_z$ with $x_zy_z=1+\varepsilon$, $\varepsilon\in\mathcal F_z$, $\varepsilon^2=0$, and then $x_z\cdot y_z(1-\varepsilon)=1$; so $x_z$ is a unit for all $z\in U$, and an element of a sheaf of rings which is a unit in every stalk is a unit (the local inverses are unique, hence glue). Hence by the universal property of localization $g$ extends uniquely to $k$-algebra maps $\sigma_a:A_a\to\mathcal O_{X'}(D(a))$ with $\pi\sigma_a=\mathrm{id}$; by uniqueness they are compatible with restriction between basic open sets, and so define a morphism of sheaves of $k$-algebras $\sigma:\mathcal O_X\to\mathcal O_{X'}$ with $\pi\sigma=\mathrm{id}_{\mathcal O_X}$.

\emph{Step 3: the isomorphism with the trivial extension.} Define $\Phi:\mathcal O_X\oplus\mathcal F\to\mathcal O_{X'}$, $\Phi(a\oplus f)=\sigma(a)+f$. It is additive, and
\[\Phi(a\oplus f)\Phi(a'\oplus f')=\sigma(aa')+\sigma(a)f'+\sigma(a')f+ff'=\sigma(aa')+af'+a'f=\Phi\big((a\oplus f)(a'\oplus f')\big),\]
because $ff'\in\mathcal I^2=0$ and $\sigma(a)f'=\pi(\sigma(a))f'=af'$. It is bijective on every open set: $x=\sigma(\pi x)+(x-\sigma\pi x)$ with $x-\sigma\pi x\in\ker\pi=\mathcal F$, uniquely. So $\Phi$ is an isomorphism of sheaves of rings; it restricts to the identity on $\mathcal F$ and induces the identity modulo $\mathcal F$. Thus $(X',\mathcal I)$ is isomorphic to the trivial extension.\hfill$\square$

\exercise{8.8}{Let $X$ be a projective nonsingular variety over $k$. For any $n>0$ we define the $n$th \emph{plurigenus} of $X$ to be $P_n=\dim_k\Gamma(X,\omega_X^{\otimes n})$. Thus in particular $P_1=p_g$. Also, for any $q$, $0\le q\le\dim X$ we define an integer $h^{q,0}=\dim_k\Gamma(X,\Omega^q_{X/k})$ where $\Omega^q_{X/k}=\bigwedge^q\Omega_{X/k}$ is the sheaf of regular $q$-forms on $X$. In particular, for $q=\dim X$, we recover the geometric genus again. The integers $h^{q,0}$ are called \emph{Hodge numbers}.

Using the method of (8.19), show that $P_n$ and $h^{q,0}$ are \emph{birational invariants}, i.e., if $X$ and $X'$ are birationally equivalent nonsingular projective varieties, then $P_n(X)=P_n(X')$ and $h^{q,0}(X)=h^{q,0}(X')$.}
\pf
Let $X,X'$ be birationally equivalent nonsingular projective varieties; they have the same function field, hence the same dimension $d$. Let $\mathcal G$ denote either $\Omega^q_{X/k}$ or $\omega_X^{\otimes n}$, and $\mathcal G'$ the corresponding sheaf on $X'$; all of these are locally free. Let $V\subseteq X$ be the largest open set on which the given birational map is a morphism $f:V\to X'$.

\emph{Step 1: $\operatorname{codim}(X\setminus V,X)\ge2$.} As in the proof of (8.19): if $P\in X$ has codimension $1$, then $\mathcal O_{X,P}$ is a discrete valuation ring ($X$ nonsingular), and $X'$ is proper over $k$, so by the valuative criterion (4.7) the map $\operatorname{Spec}K\to X'$ extends to $\operatorname{Spec}\mathcal O_{X,P}\to X'$, hence the rational map is defined at $P$ and $P\in V$ (a morphism $\operatorname{Spec}\mathcal O_{X,P}\to X'$ of finite type over $k$ extends to a morphism on an open neighbourhood of $P$, which agrees with the given rational map at the generic point, hence on a dense open set, hence by separatedness of $X'$ is the rational map). The generic point is in $V$ as well, so every irreducible component of $X\setminus V$ has codimension $\ge2$.

\emph{Step 2: $\Gamma(X,\mathcal G)\to\Gamma(V,\mathcal G|_V)$ is bijective.} Let $U\subseteq X$ be an affine open with $\mathcal G|_U$ free. Then $\Gamma(U,\mathcal G)\to\Gamma(U\cap V,\mathcal G)$ is, componentwise, the map $\Gamma(U,\mathcal O_U)\to\Gamma(U\cap V,\mathcal O)$, which is bijective: $X$ is nonsingular, hence normal, and $U\setminus U\cap V$ has codimension $\ge2$, so this is a consequence of (6.3A), exactly as in the proof of (8.19). Gluing, the claim follows.

\emph{Step 3: a map $f^*:\Gamma(X',\mathcal G')\to\Gamma(V,\mathcal G|_V)$.} By (8.11) there is a map $f^*\Omega_{X'/k}\to\Omega_{V/k}$ of locally free sheaves of the same rank $d$; taking $q$th exterior powers, resp.\ the top exterior power followed by the $n$th tensor power, we get maps $f^*\mathcal G'\to\mathcal G|_V$, hence $f^*:\Gamma(X',\mathcal G')\to\Gamma(V,f^*\mathcal G')\to\Gamma(V,\mathcal G|_V)$.

\emph{Step 4: $f^*$ is injective.} Since $f$ is birational there is a nonempty open $W\subseteq V$ such that $f(W)$ is open in $X'$ and $f:W\to f(W)$ is an isomorphism (I, 4.5 / its scheme version). An isomorphism induces an isomorphism on differentials, hence $f^*\mathcal G'\to\mathcal G$ is an isomorphism over $W$. Let $s'\in\Gamma(X',\mathcal G')$ with $f^*s'=0$. Then $s'|_{f(W)}=0$ (pull-back by the isomorphism $W\cong f(W)$ is bijective on sections). Now $\mathcal G'$ is locally free on the integral scheme $X'$ and $f(W)$ is a nonempty open set, so restriction $\Gamma(X',\mathcal G')\to\Gamma(f(W),\mathcal G')$ is injective; hence $s'=0$.

Combining Steps 2--4, $\dim\Gamma(X',\mathcal G')\le\dim\Gamma(V,\mathcal G)=\dim\Gamma(X,\mathcal G)$. By symmetry (exchange $X$ and $X'$) the reverse inequality holds. Therefore $P_n(X)=P_n(X')$ for all $n>0$ (taking $\mathcal G=\omega_X^{\otimes n}$) and $h^{q,0}(X)=h^{q,0}(X')$ for all $q$ (taking $\mathcal G=\Omega^q_{X/k}$).\hfill$\square$

\section*{9. Formal schemes}

\noindent\emph{Conventions.} $\widehat X$ denotes the completion of $X$ along $Y$, $\mathfrak X$ a noetherian formal scheme, $\mathfrak J$ an ideal of definition, $Y_n=(\mathfrak X,\mathcal O_{\mathfrak X}/\mathfrak J^n)$. For a finitely generated module $M$ over a noetherian ring $A$ with ideal $I$, $M^\Delta=\widehat{\widetilde M}$ on $\widehat X=\operatorname{Spec}A^{\wedge}$. We use the following facts established in the text: $M\mapsto M^\Delta$ is exact (9.4c) and $M^\Delta\cong\widetilde M\otimes_{\mathcal O_X}\mathcal O_{\widehat X}$ (9.4b); if $A$ is $I$-adically complete then $\Gamma(\widehat X,M^\Delta)=M$ and $M\mapsto M^\Delta$ is an equivalence of finitely generated $A$-modules with coherent $\mathcal O_{\widehat X}$-modules (9.7); for a coherent $\mathfrak F$, $\mathfrak F/\mathfrak J^n\mathfrak F$ is a coherent $\mathcal O_{Y_n}$-module and $\mathfrak F\cong\varprojlim\mathfrak F/\mathfrak J^n\mathfrak F$ (9.6a). As in the proofs in the text, since the questions are local we may (and do) work on an open affine formal subscheme $\mathfrak U=\widehat U$, $U=\operatorname{Spec}A$ with $A$ noetherian and $I$-adically complete (completing $A$ does not change $\widehat U$), with $\mathfrak J|_{\mathfrak U}=I^\Delta$, so that $Y_n\cap\mathfrak U=\operatorname{Spec}A/I^n$. (By (9.5a) any two ideals of definition are cofinal, so nothing depends on the choice of $\mathfrak J$.) Coherence of a sheaf on $\mathfrak X$ is a local property (as in (9.8), (9.9)), and by (9.7) a coherent sheaf restricted to such a $\mathfrak U$ is of the form $M^\Delta$ with $M$ finitely generated.

\exercise{9.1}{Let $X$ be a noetherian scheme, $Y$ a closed subscheme, and $\widehat X$ the completion of $X$ along $Y$. We call the ring $\Gamma(\widehat X,\mathcal O_{\widehat X})$ the ring of \emph{formal-regular functions} on $X$ along $Y$. In this exercise we show that if $Y$ is a connected, nonsingular, positive-dimensional subvariety of $X=\mathbf P^n$ over an algebraically closed field $k$, then $\Gamma(\widehat X,\mathcal O_{\widehat X})=k$.
\begin{enumerate}[label=(\alph*),leftmargin=*]
\item Let $\mathcal I$ be the ideal sheaf of $Y$. Use (8.13) and (8.17) to show that there is an inclusion of sheaves on $Y$, $\mathcal I/\mathcal I^2\hookrightarrow\mathcal O_Y(-1)^{n+1}$.
\item Show that for any $r\ge1$, $\Gamma(Y,\mathcal I^r/\mathcal I^{r+1})=0$.
\item Use the exact sequences
\[0\to\mathcal I^r/\mathcal I^{r+1}\to\mathcal O_X/\mathcal I^{r+1}\to\mathcal O_X/\mathcal I^r\to0\]
and induction on $r$ to show that $\Gamma(Y,\mathcal O_X/\mathcal I^r)=k$ for all $r\ge1$. (Use (8.21Ae).)
\item Conclude that $\Gamma(\widehat X,\mathcal O_{\widehat X})=k$. (Actually, the same result holds without the hypothesis $Y$ nonsingular, but the proof is more difficult---see Hartshorne [3, (7.3)].)
\end{enumerate}}
\pf
$Y$ is connected and nonsingular (all local rings regular, hence domains), so it is irreducible, i.e.\ a nonsingular closed subvariety of the nonsingular variety $\mathbf P^n$; let $q=\operatorname{codim}(Y,\mathbf P^n)$ and $s=\dim Y\ge1$. Then $Y$ is a projective variety, so $\Gamma(Y,\mathcal O_Y)$ is a finite-dimensional $k$-algebra (II, Thm.~5.19) and a domain, hence a finite field extension of $k$, hence equal to $k$.

(a) By (8.17) there is an exact sequence $0\to\mathcal I/\mathcal I^2\to\Omega_{\mathbf P^n/k}\otimes\mathcal O_Y\to\Omega_{Y/k}\to0$, with $\Omega_{Y/k}$ locally free. Tensoring the Euler sequence (8.13) $0\to\Omega_{\mathbf P^n}\to\mathcal O(-1)^{n+1}\to\mathcal O\to0$ with $\mathcal O_Y$ stays exact (the third term is locally free), giving $0\to\Omega_{\mathbf P^n}|_Y\to\mathcal O_Y(-1)^{n+1}\to\mathcal O_Y\to0$. The composite
\[\mathcal I/\mathcal I^2\hookrightarrow\Omega_{\mathbf P^n}|_Y\hookrightarrow\mathcal O_Y(-1)^{n+1}\]
is injective. Both inclusions have locally free cokernels ($\Omega_{Y/k}$ and $\mathcal O_Y$), hence are locally split; so the composite is locally split, with cokernel locally free (extension of locally free sheaves by locally free).

(b) For $y\in Y$, $\mathcal O_{X,y}$ is regular (so Cohen--Macaulay, (8.21Aa)) and $\mathcal O_{X,y}/\mathcal I_y=\mathcal O_{Y,y}$ is regular of dimension $\dim\mathcal O_{X,y}-q$, so $\mathcal I_y$ is generated by $q$ elements forming part of a regular system of parameters, which is a regular sequence (8.21Ac), exactly as in the proof of (8.24); by (8.21Ae) the natural map $S^r(\mathcal I/\mathcal I^2)\to\mathcal I^r/\mathcal I^{r+1}$ is an isomorphism on every stalk, hence an isomorphism of sheaves on $Y$ (as in the proof of (8.24)). By (a), $\mathcal I/\mathcal I^2$ is locally a direct summand of $\mathcal O_Y(-1)^{n+1}$, so $S^r$ of the inclusion is also an inclusion (locally a direct summand):
\[\mathcal I^r/\mathcal I^{r+1}\cong S^r(\mathcal I/\mathcal I^2)\hookrightarrow S^r(\mathcal O_Y(-1)^{n+1})\cong\mathcal O_Y(-r)^{N},\qquad N=\tbinom{n+r}r.\]
So $\Gamma(Y,\mathcal I^r/\mathcal I^{r+1})\hookrightarrow\Gamma(Y,\mathcal O_Y(-r))^N$, and it suffices to prove $\Gamma(Y,\mathcal O_Y(-r))=0$ for $r\ge1$. Suppose $0\ne t\in\Gamma(Y,\mathcal O_Y(-r))$. For each monomial $u_\alpha=x^\alpha$ of degree $r$ (a section of $\mathcal O_Y(r)$) we get $tu_\alpha\in\Gamma(Y,\mathcal O_Y)=k$, say $tu_\alpha=c_\alpha$. The $u_\alpha$ generate $\mathcal O_Y(r)$ and $t\ne0$ at the generic point of the integral scheme $Y$, so some $c_\alpha\ne0$. Then $t$ and $u_\alpha$ are nowhere vanishing, and for every $\beta$, $u_\beta=(c_\beta/c_\alpha)u_\alpha$ (multiply $tu_\beta=c_\beta$ and $tu_\alpha=c_\alpha$ and cancel the unit $t$). So all monomials of degree $r$ are proportional on $Y$, i.e.\ the $r$-uple Veronese embedding maps $Y$ to a single point; as it is injective, $Y$ is a point, contradicting $\dim Y\ge1$. So $\Gamma(Y,\mathcal O_Y(-r))=0$.

(c) For $r=1$: $\mathcal O_X/\mathcal I=\mathcal O_Y$ and $\Gamma(Y,\mathcal O_Y)=k$. Suppose $\Gamma(Y,\mathcal O_X/\mathcal I^r)=k$ (the constants, i.e.\ the image of $k$; the map $k\to\Gamma$ is injective as $Y\ne\emptyset$). Left exactness of $\Gamma$ and (b) give an exact sequence $0=\Gamma(\mathcal I^r/\mathcal I^{r+1})\to\Gamma(Y,\mathcal O_X/\mathcal I^{r+1})\to\Gamma(Y,\mathcal O_X/\mathcal I^r)=k$. So $\Gamma(Y,\mathcal O_X/\mathcal I^{r+1})$ injects into $k$ and contains $k$; hence it equals $k$, and the map to $\Gamma(Y,\mathcal O_X/\mathcal I^r)=k$ is the identity of $k$.

(d) By definition $\mathcal O_{\widehat X}=\varprojlim\mathcal O_X/\mathcal I^r$, so by (9.2) $\Gamma(\widehat X,\mathcal O_{\widehat X})=\varprojlim_r\Gamma(Y,\mathcal O_X/\mathcal I^r)=\varprojlim(k\xleftarrow{\ \mathrm{id}\ }k\leftarrow\cdots)=k$.\hfill$\square$

\exercise{9.2}{Use the result of (Ex.~9.1) to prove the following geometric result. Let $Y\subseteq X=\mathbf P^n_k$ be as above, and let $f:X\to Z$ be a morphism of $k$-varieties. Suppose that $f(Y)$ is a single closed point $P\in Z$. Then $f(X)=P$ also.}
\pf
Here $Y$ is connected, nonsingular, of positive dimension (so an irreducible nonsingular subvariety), as in Ex.~9.1, and $\mathcal I$ is its ideal sheaf. Let $W=\operatorname{Spec}C$ be an affine open neighbourhood of $P$ in $Z$, $\fm\subseteq C$ the maximal ideal of $P$, and $X_0=f^{-1}(W)$, an open subset of $X$ containing $Y$. The morphism $f|_{X_0}$ corresponds to a $k$-algebra homomorphism $f^\#:C\to\Gamma(X_0,\mathcal O_X)$.

\emph{Claim 1: $f^\#(\fm^m)\subseteq\Gamma(X_0,\mathcal I^m)$ for all $m\ge1$.} Let $c\in\fm$ and $g=f^\#(c)$. For every point $y\in Y$ we have $f(y)=P$ and $c\in\fm_P$, so $g(y)=c(P)=0$ in $k(y)$ (the residue field map $k(P)\to k(y)$ is a ring map). So $g|_Y\in\Gamma(Y,\mathcal O_Y)$ vanishes at every point of the reduced scheme $Y$, hence is $0$, i.e.\ $g\in\Gamma(X_0,\mathcal I)$. As $f^\#$ is a ring homomorphism and $\mathcal I^m$ is the $m$th power, $f^\#(\fm^m)\subseteq\Gamma(X_0,\mathcal I^m)$.

\emph{Claim 2: for every $c\in C$ there is $\lambda(c)\in k$ with $f^\#(c)-\lambda(c)\cdot1=0$.} Fix $c$. For each $m$, the image of $f^\#(c)$ in $\Gamma(X_0,\mathcal O_X/\mathcal I^m)=\Gamma(Y,\mathcal O_X/\mathcal I^m)$ (the sheaf $\mathcal O_X/\mathcal I^m$ is supported on $Y\subseteq X_0$) lies in $k$ by Ex.~9.1(c): it equals $\lambda_m\cdot 1$ with $\lambda_m\in k$; since the transition maps $\Gamma(\mathcal O_X/\mathcal I^{m+1})\to\Gamma(\mathcal O_X/\mathcal I^m)$ are the identity of $k$ (proof of 9.1(c)), $\lambda_m=\lambda$ is independent of $m$. Let $e=f^\#(c)-\lambda\cdot1\in\Gamma(X_0,\mathcal O_X)$. Its image in $\Gamma(X_0,\mathcal O_X/\mathcal I^m)$ is $0$ for all $m$, so $e_y\in\mathcal I_y^m$ for every $y\in Y$ and all $m$. By Krull's intersection theorem in the noetherian local ring $\mathcal O_{X,y}$ (with $\mathcal I_y\subseteq\fm_y$), $\bigcap_m\mathcal I_y^m=0$, so $e_y=0$. Thus $e$ vanishes near each point of $Y$; in particular at the generic point $\eta$ of $X$ (any open neighbourhood of a point of $Y$ contains $\eta$): $e_\eta=0$. As $X_0$ is a nonempty open subset of the integral scheme $X$, $\Gamma(X_0,\mathcal O_X)\to\mathcal O_{X,\eta}$ is injective, so $e=0$.

Hence $f^\#(C)\subseteq k\cdot1$. This means that $f|_{X_0}:X_0\to\operatorname{Spec}C$ factors through $\operatorname{Spec}k\to\operatorname{Spec}C$, so $f(X_0)$ is a single point; it contains $f(Y)=P$ (as $Y\ne\emptyset$), so $f(X_0)=\{P\}$. Since $X$ is irreducible, $X_0$ is dense, and $f(X)\subseteq\overline{f(X_0)}=\overline{\{P\}}=\{P\}$ ($P$ is closed in $Z$). Thus $f(X)=P$.\hfill$\square$

\exercise{9.3}{Prove the analogue of (5.6) for formal schemes, which says, if $\mathfrak X$ is an affine formal scheme, and if
\[0\to\mathfrak F'\to\mathfrak F\to\mathfrak F''\to0\]
is an exact sequence of $\mathcal O_{\mathfrak X}$-modules, and if $\mathfrak F'$ is coherent, then the sequence of global sections
\[0\to\Gamma(\mathfrak X,\mathfrak F')\to\Gamma(\mathfrak X,\mathfrak F)\to\Gamma(\mathfrak X,\mathfrak F'')\to0\]
is exact. For the proof, proceed in the following steps.
\begin{enumerate}[label=(\alph*),leftmargin=*]
\item Let $\mathfrak J$ be an ideal of definition for $\mathfrak X$, and for each $n>0$ consider the exact sequence
\[0\to\mathfrak F'/\mathfrak J^n\mathfrak F'\to\mathfrak F/\mathfrak J^n\mathfrak F'\to\mathfrak F''\to0.\]
Use (5.6), slightly modified, to show that for every open affine subset $\mathfrak U\subseteq\mathfrak X$, the sequence
\[0\to\Gamma(\mathfrak U,\mathfrak F'/\mathfrak J^n\mathfrak F')\to\Gamma(\mathfrak U,\mathfrak F/\mathfrak J^n\mathfrak F')\to\Gamma(\mathfrak U,\mathfrak F'')\to0\]
is exact.
\item Now pass to the limit, using (9.1), (9.2), and (9.6). Conclude that $\mathfrak F\cong\varprojlim\mathfrak F/\mathfrak J^n\mathfrak F'$ and that the sequence of global sections above is exact.
\end{enumerate}}
\pf
(The quotient in the middle is $\mathfrak F/\mathfrak J^n\mathfrak F'$, where $\mathfrak J^n\mathfrak F'\subseteq\mathfrak F'\subseteq\mathfrak F$; this is what makes the first sequence exact, and it is the reading used.) Write $\mathcal G'_n=\mathfrak F'/\mathfrak J^n\mathfrak F'$ and $\mathcal G_n=\mathfrak F/\mathfrak J^n\mathfrak F'$; the transition maps $\mathcal G_{n+1}\to\mathcal G_n$, $\mathcal G'_{n+1}\to\mathcal G'_n$ are the obvious surjections, and the sequences
\[0\to\mathcal G'_n\to\mathcal G_n\to\mathfrak F''\to0\tag{1}\]
are exact (the kernel of $\mathfrak F/\mathfrak J^n\mathfrak F'\to\mathfrak F''$ is $\mathfrak F'/\mathfrak J^n\mathfrak F'$) and compatible with the transition maps (the maps on $\mathfrak F''$ being identities).

(a) Let $\mathfrak U=\widehat U$ be an open affine formal subscheme, $U=\operatorname{Spec}B$; then $(\mathfrak U,\mathcal O_{Y_n}|_{\mathfrak U})$ is an affine noetherian scheme: by (9.5a) $\mathfrak J^n\supseteq(I^\Delta)^N$ for some $N$ on $\mathfrak U$, so it is a closed subscheme of $(\mathfrak U,\mathcal O/(I^\Delta)^N)=\operatorname{Spec}B/I^N$ (9.4a), hence affine. (The same holds for any open $\mathfrak U$ for which $(\mathfrak U,\mathcal O_{Y_1}|_{\mathfrak U})$ is affine, since a scheme whose reduction is affine is affine, III, Ex.~3.1.) By (9.6a), $\mathcal G'_n=\mathfrak F'/\mathfrak J^n\mathfrak F'$ is a coherent, in particular quasi-coherent, sheaf on this affine scheme, so $H^1(\mathfrak U,\mathcal G'_n)=0$ (III, Thm.~3.5). Cohomology of a sheaf of abelian groups is the same computed in abelian groups or in modules over a sheaf of rings (III, Prop.~2.6), so the long exact sequence of cohomology of (1) on $\mathfrak U$ (in the category of sheaves of abelian groups) gives
\[0\to\Gamma(\mathfrak U,\mathcal G'_n)\to\Gamma(\mathfrak U,\mathcal G_n)\to\Gamma(\mathfrak U,\mathfrak F'')\to H^1(\mathfrak U,\mathcal G'_n)=0.\]
(This is exactly (5.6): the proof of (5.6) only needs that the subsheaf is quasi-coherent on an affine scheme with quasi-compact base of principal opens $D(f)$, and here the sheaves $\mathcal G_n,\mathfrak F''$ are merely sheaves of abelian groups on the underlying space of the affine scheme $\operatorname{Spec}B/I^nB$; that is the ``slight modification''.)

(b) Fix $\mathfrak U$ as above. The transition map $\Gamma(\mathfrak U,\mathcal G'_{n+1})\to\Gamma(\mathfrak U,\mathcal G'_n)$ is surjective: its kernel sheaf is $\mathfrak J^n\mathfrak F'/\mathfrak J^{n+1}\mathfrak F'$, a coherent (kernel of a surjection of coherent $\mathcal O_{Y_{n+1}}$-modules, (9.6a)), in particular quasi-coherent, sheaf on the affine noetherian scheme $(\mathfrak U,\mathcal O_{Y_{n+1}})$, so its $H^1$ vanishes (III, Thm.~3.5, 2.6). Hence the inverse system $(\Gamma(\mathfrak U,\mathcal G'_n))_n$ has surjective transition maps and satisfies (ML) (Example 9.1.1). Take the inverse limit of the exact sequences of (a); by (9.1b) we obtain an exact sequence
\[0\to\varprojlim\Gamma(\mathfrak U,\mathcal G'_n)\to\varprojlim\Gamma(\mathfrak U,\mathcal G_n)\to\varprojlim\Gamma(\mathfrak U,\mathfrak F'')\to0.\]
The last system is constant, with limit $\Gamma(\mathfrak U,\mathfrak F'')$; by (9.2), $\varprojlim\Gamma(\mathfrak U,\mathcal G_n)=\Gamma(\mathfrak U,\varprojlim\mathcal G_n)$ and $\varprojlim\Gamma(\mathfrak U,\mathcal G'_n)=\Gamma(\mathfrak U,\varprojlim\mathcal G'_n)$. So for all such $\mathfrak U$ (which form a base of the topology),
\[0\to\Gamma(\mathfrak U,\varprojlim\mathcal G'_n)\to\Gamma(\mathfrak U,\varprojlim\mathcal G_n)\to\Gamma(\mathfrak U,\mathfrak F'')\to0\]
is exact; exactness on sections over a base of opens implies exactness of the sequence of sheaves (stalks are direct limits over the base). Now $\varprojlim\mathcal G'_n=\varprojlim\mathfrak F'/\mathfrak J^n\mathfrak F'\cong\mathfrak F'$ by (9.6a) ($\mathfrak F'$ is coherent), and the natural maps $\mathfrak F\to\mathcal G_n$ are compatible, giving a morphism $\mathfrak F\to\varprojlim\mathcal G_n$ (9.2) and a commutative diagram with exact rows
\[\begin{array}{ccccccccc}0&\to&\mathfrak F'&\to&\mathfrak F&\to&\mathfrak F''&\to&0\\&&\downarrow\cong&&\downarrow&&\|&&\\0&\to&\varprojlim\mathcal G'_n&\to&\varprojlim\mathcal G_n&\to&\mathfrak F''&\to&0\end{array}\]
By the five lemma, $\mathfrak F\cong\varprojlim\mathfrak F/\mathfrak J^n\mathfrak F'$. Finally, take $\mathfrak U=\mathfrak X$ (which is affine): $\Gamma(\mathfrak X,\mathfrak F)=\Gamma(\mathfrak X,\varprojlim\mathcal G_n)$ and $\Gamma(\mathfrak X,\mathfrak F')=\Gamma(\mathfrak X,\varprojlim\mathcal G'_n)$, so the middle row above on global sections is $0\to\Gamma(\mathfrak X,\mathfrak F')\to\Gamma(\mathfrak X,\mathfrak F)\to\Gamma(\mathfrak X,\mathfrak F'')\to0$, which is therefore exact.\hfill$\square$

\exercise{9.4}{Use (Ex.~9.3) to prove that if
\[0\to\mathfrak F'\to\mathfrak F\to\mathfrak F''\to0\]
is an exact sequence of $\mathcal O_{\mathfrak X}$-modules on a noetherian formal scheme $\mathfrak X$, and if $\mathfrak F',\mathfrak F''$ are coherent, then $\mathfrak F$ is coherent also.}
\pf
Coherence is local. Cover $\mathfrak X$ by finitely many open affine formal subschemes $\mathfrak U=\widehat U$, $U=\operatorname{Spec}A$, $A$ complete, on which $\mathfrak F'$ and $\mathfrak F''$ are both coherent (refine the covers on which each is coherent by smaller affine formal opens), i.e.\ by (9.7) $\mathfrak F'|_{\mathfrak U}\cong M'^\Delta$, $\mathfrak F''|_{\mathfrak U}\cong M''^\Delta$ with $M',M''$ finitely generated $A$-modules. It suffices to show $\mathfrak F|_{\mathfrak U}\cong M^\Delta$ for a finitely generated $A$-module $M$.

Let $M=\Gamma(\mathfrak U,\mathfrak F)$, an $A$-module ($A=\Gamma(\mathfrak U,\mathcal O_{\mathfrak U})$ by (9.3.1)). By (9.7), $M'=\Gamma(\mathfrak U,\mathfrak F')$ and $M''=\Gamma(\mathfrak U,\mathfrak F'')$; by Ex.~9.3 applied to $\mathfrak U$ (affine) the sequence
\[0\to M'\to M\to M''\to0\]
is exact. Hence $M$ is finitely generated over the noetherian ring $A$. For any $\mathcal O_{\mathfrak U}$-module $\mathfrak G$ with $N=\Gamma(\mathfrak U,\mathfrak G)$ there is a natural map $\psi_{\mathfrak G}:N\otimes_A\mathcal O_{\mathfrak U}\to\mathfrak G$, $m\otimes g\mapsto g\cdot m|_V$ on a small open set $V$, natural in $\mathfrak G$. If $N$ is finitely generated, $N\otimes_A\mathcal O_{\mathfrak U}=N^\Delta$ (both are right exact in $N$ and agree for $N=A$, by (9.4b)). For $\mathfrak G=\mathfrak F',\mathfrak F''$, $\psi_{\mathfrak G}$ is the isomorphism of (9.7) (the counit of the equivalence $M\leftrightarrow M^\Delta$). We obtain a commutative diagram
\[\begin{array}{ccccccccc}0&\to&M'^\Delta&\to&M^\Delta&\to&M''^\Delta&\to&0\\&&\downarrow\cong&&\downarrow\psi_{\mathfrak F}&&\downarrow\cong&&\\0&\to&\mathfrak F'&\to&\mathfrak F&\to&\mathfrak F''&\to&0\end{array}\]
whose top row is exact by (9.4c) (exactness of $\Delta$) and whose bottom row is exact by hypothesis. By the five lemma $\psi_{\mathfrak F}$ is an isomorphism, so $\mathfrak F|_{\mathfrak U}\cong M^\Delta$ is coherent. Since $\mathfrak X$ is covered by finitely many such $\mathfrak U$, $\mathfrak F$ is coherent.\hfill$\square$

\exercise{9.5}{If $\mathfrak F$ is a coherent sheaf on a noetherian formal scheme $\mathfrak X$, which can be generated by global sections, show in fact that it can be generated by a finite number of its global sections.}
\pf
Cover $\mathfrak X$ by finitely many open affine formal subschemes $\mathfrak U_i=\widehat{U_i}$, $U_i=\operatorname{Spec}A_i$ with $A_i$ complete, on which $\mathfrak F|_{\mathfrak U_i}=M_i^\Delta$, $M_i=\Gamma(\mathfrak U_i,\mathfrak F)$ finitely generated over $A_i$ (9.7). Fix $i$ and let $M_i'\subseteq M_i$ be the $A_i$-submodule generated by the images of $\Gamma(\mathfrak X,\mathfrak F)\to\Gamma(\mathfrak U_i,\mathfrak F)=M_i$. Since $A_i$ is noetherian, $M_i'$ is finitely generated, so it is generated by the images of finitely many global sections $s_{i1},\dots,s_{ir_i}\in\Gamma(\mathfrak X,\mathfrak F)$ (each of finitely many generators is a finite $A_i$-combination of images of global sections).

\emph{Claim: $M_i'=M_i$.} The natural map $M_i'^\Delta\to M_i^\Delta=\mathfrak F|_{\mathfrak U_i}$ is injective (exactness of $\Delta$) and its image contains $(s|_{\mathfrak U_i})_x$ for every global section $s$ and every $x\in\mathfrak U_i$; since the global sections generate $\mathfrak F$, these generate every stalk, so the map is surjective. Hence $(M_i/M_i')^\Delta=\operatorname{coker}=0$, and since $A_i$ is complete, $M_i/M_i'=\Gamma(\mathfrak U_i,(M_i/M_i')^\Delta)=0$ by (9.7). So $M_i'=M_i$.

Therefore the sections $s_{ij}$ ($j\le r_i$) restricted to $\mathfrak U_i$ generate $M_i$, hence generate $M_i\otimes_{A_i}\mathcal O_{\mathfrak U_i}=\mathfrak F|_{\mathfrak U_i}$ (a surjection $\mathcal O^{r_i}\to M_i^\Delta$ is obtained from generators of $M_i$). The finitely many global sections $\{s_{ij}\}_{i,j}$ then generate $\mathfrak F$ on every $\mathfrak U_i$, i.e.\ on $\mathfrak X$.\hfill$\square$

\exercise{9.6}{Let $\mathfrak X$ be a noetherian formal scheme, let $\mathfrak J$ be an ideal of definition, and for each $n$, let $Y_n$ be the scheme $(\mathfrak X,\mathcal O_{\mathfrak X}/\mathfrak J^n)$. Assume that the inverse system of groups $(\Gamma(Y_n,\mathcal O_{Y_n}))$ satisfies the Mittag-Leffler condition. Then prove that $\operatorname{Pic}\mathfrak X=\varprojlim\operatorname{Pic}Y_n$. As in the case of a scheme, we define $\operatorname{Pic}\mathfrak X$ to be the group of locally free $\mathcal O_{\mathfrak X}$-modules of rank $1$ under the operation $\otimes$. Proceed in the following steps.
\begin{enumerate}[label=(\alph*),leftmargin=*]
\item Use the fact that $\ker(\Gamma(Y_{n+1},\mathcal O_{Y_{n+1}})\to\Gamma(Y_n,\mathcal O_{Y_n}))$ is a nilpotent ideal to show that the inverse system $(\Gamma(Y_n,\mathcal O^*_{Y_n}))$ of units also satisfies (ML).
\item Let $\mathfrak F$ be a coherent sheaf of $\mathcal O_{\mathfrak X}$-modules, and assume that for each $n$, there is some isomorphism $\varphi_n:\mathfrak F/\mathfrak J^n\mathfrak F\cong\mathcal O_{Y_n}$. Then show that there is an isomorphism $\mathfrak F\cong\mathcal O_{\mathfrak X}$. Be careful, because the $\varphi_n$ may not be compatible with the maps of the two inverse systems $(\mathfrak F/\mathfrak J^n\mathfrak F)$ and $(\mathcal O_{Y_n})$! Conclude that the natural map $\operatorname{Pic}\mathfrak X\to\varprojlim\operatorname{Pic}Y_n$ is injective.
\item Given an invertible sheaf $\mathcal L_n$ on $Y_n$ for each $n$, and given isomorphisms $\mathcal L_{n+1}\otimes\mathcal O_{Y_n}\cong\mathcal L_n$, construct maps $\mathcal L_{n'}\to\mathcal L_n$ for each $n'\ge n$ so as to make an inverse system, and show that $\mathfrak L=\varprojlim\mathcal L_n$ is a coherent sheaf on $\mathfrak X$. Then show that $\mathfrak L$ is locally free of rank $1$, and thus conclude that the map $\operatorname{Pic}\mathfrak X\to\varprojlim\operatorname{Pic}Y_n$ is surjective. Again, be careful, because even though each $\mathcal L_n$ is locally free of rank $1$, the open sets needed to make them free might get smaller and smaller with $n$.
\item Show that the hypothesis ``$(\Gamma(Y_n,\mathcal O_{Y_n}))$ satisfies (ML)'' is satisfied if either $\mathfrak X$ is affine, or each $Y_n$ is projective over a field $k$.
\end{enumerate}
\emph{Note}: See (III, Ex.~11.5--11.7) for further examples and applications.}
\pf
Put $R_n=\Gamma(Y_n,\mathcal O_{Y_n})$. The map $\operatorname{Pic}\mathfrak X\to\varprojlim\operatorname{Pic}Y_n$ is $\mathfrak L\mapsto(\mathfrak L/\mathfrak J^n\mathfrak L)_n$; it is a homomorphism since $(\mathfrak L\otimes\mathfrak L')/\mathfrak J^n=\mathfrak L/\mathfrak J^n\otimes\mathfrak L'/\mathfrak J^n$, and an invertible $\mathfrak L$ is coherent ($\mathcal O_{\mathfrak X}$ is) with $\mathfrak L/\mathfrak J^n\mathfrak L$ invertible on $Y_n$. Note that $\mathfrak F_{n+1}\otimes\mathcal O_{Y_n}=\mathfrak F/\mathfrak J^n\mathfrak F=\mathfrak F_n$ for $\mathfrak F_n:=\mathfrak F/\mathfrak J^n\mathfrak F$.

(a) The kernel of $\mathcal O_{Y_{n'}}\to\mathcal O_{Y_n}$ is $\mathfrak J^n/\mathfrak J^{n'}$, and $(\mathfrak J^n/\mathfrak J^{n'})^m=0$ for $nm\ge n'$; by left exactness of $\Gamma$, $K_{n'n}:=\ker(R_{n'}\to R_n)=\Gamma(Y_{n'},\mathfrak J^n/\mathfrak J^{n'})$ is a nilpotent ideal of $R_{n'}$. Let $S_n(n')=\operatorname{im}(R_{n'}\to R_n)\cong R_{n'}/K_{n'n}$, a subring of $R_n$. By (ML) for $(R_n)$, for each $n$ there is $n_0=n_0(n)\ge n$ such that $S_n(n')=S_n$ is the same subring for all $n'\ge n_0$. Since units lift modulo a nilpotent ideal ($y\in R_{n'}$ is a unit iff its image in $R_{n'}/K_{n'n}$ is a unit), we get
\[\operatorname{im}(R_{n'}^*\to R_n^*)=\{\,\text{units of the ring }S_n(n')\,\}=S_n^*\quad(n'\ge n_0).\]
Here a unit of $S_n(n')=R_{n'}/K_{n'n}$ means an element with an inverse \emph{in $S_n(n')$}, and is the image of a unit of $R_{n'}$. In particular the images $\operatorname{im}(R_{n'}^*\to R_n^*)$ are constant for $n'\ge n_0(n)$, so $(R_n^*)$ satisfies (ML). Call this stable subgroup $U_n=S_n^*\subseteq R_n^*$.

(b) Let $\mathfrak F$ be coherent with $\mathfrak F_n:=\mathfrak F/\mathfrak J^n\mathfrak F\cong\mathcal O_{Y_n}$ for all $n$. Let $E_n=\operatorname{Isom}_{\mathcal O_{Y_n}}(\mathfrak F_n,\mathcal O_{Y_n})\ne\emptyset$. The group $R_n^*=\operatorname{Aut}(\mathcal O_{Y_n})$ acts simply transitively on $E_n$ (two isomorphisms differ by an automorphism of $\mathcal O_{Y_n}$, i.e.\ multiplication by a unit of $R_n$). Reduction $\psi\mapsto\psi\otimes\mathcal O_{Y_n}$ gives maps $E_{n+1}\to E_n$ (as $\mathfrak F_{n+1}\otimes\mathcal O_{Y_n}=\mathfrak F_n$), compatible with the actions via $R^*_{n+1}\to R^*_n$. So $(E_n)$ is an inverse system of nonempty sets.

\emph{$(E_n)$ satisfies (ML).} Fix $n$ and $n'\ge n_0(n)$, and choose any $e\in E_{n'}$ with image $\bar e\in E_n$. Then $\operatorname{im}(E_{n'}\to E_n)=\{\bar u\bar e:u\in R_{n'}^*\}=U_n\cdot\bar e$ by (a), a coset of the fixed subgroup $U_n$ of $R_n^*$ acting on $E_n$. For $n''\ge n'\ge n_0$ we have $\operatorname{im}(E_{n''}\to E_n)\subseteq\operatorname{im}(E_{n'}\to E_n)$, and both are cosets of the same subgroup $U_n$; a coset contained in another coset of the same subgroup equals it. So the images stabilize.

\emph{$\varprojlim E_n\neq\emptyset$.} Let $E_n'$ be the stable image of $(E_{n'})_{n'\ge n}$ in $E_n$; these are nonempty, form an inverse system with surjective transition maps, so we choose successively $e_1\in E_1'$, $e_2\in E_2'$ mapping to $e_1$, etc. (the argument of (9.1b)). A compatible family $(\psi_n)\in\varprojlim E_n$ is an isomorphism of inverse systems $(\mathfrak F_n)\to(\mathcal O_{Y_n})$, hence induces an isomorphism on the limits. By (9.6a) $\varprojlim\mathfrak F_n\cong\mathfrak F$ ($\mathfrak F$ coherent) and $\varprojlim\mathcal O_{Y_n}=\mathcal O_{\mathfrak X}$. Thus $\mathfrak F\cong\mathcal O_{\mathfrak X}$.

\emph{Injectivity.} If $\mathfrak L\in\operatorname{Pic}\mathfrak X$ maps to the identity of $\varprojlim\operatorname{Pic}Y_n$, then $\mathfrak L/\mathfrak J^n\mathfrak L\cong\mathcal O_{Y_n}$ for each $n$, so $\mathfrak L\cong\mathcal O_{\mathfrak X}$ by the above. So the kernel of the homomorphism is trivial.

(c) An element of $\varprojlim\operatorname{Pic}Y_n$ is a sequence of classes $[\mathcal L_n]$ with $[\mathcal L_{n+1}\otimes\mathcal O_{Y_n}]=[\mathcal L_n]$; choose invertible sheaves $\mathcal L_n$ and isomorphisms $\theta_n:\mathcal L_{n+1}\otimes\mathcal O_{Y_n}\to\mathcal L_n$. Let $\rho_{n+1,n}:\mathcal L_{n+1}\to\mathcal L_{n+1}\otimes\mathcal O_{Y_n}=\mathcal L_{n+1}/\mathfrak J^n\mathcal L_{n+1}\xrightarrow{\theta_n}\mathcal L_n$, a surjection with kernel $\mathfrak J^n\mathcal L_{n+1}$, and let $\rho_{n'n}=\rho_{n+1,n}\circ\cdots\circ\rho_{n',n'-1}$ for $n'>n$ (and $\rho_{nn}=\mathrm{id}$). This makes $(\mathcal L_n)$ an inverse system with surjective maps; and $\ker\rho_{n'n}=\mathfrak J^n\mathcal L_{n'}$ (induction on $n'-n$: $\ker\rho_{n'n}$ is the preimage under $\rho_{n',n+1}$ of $\ker\rho_{n+1,n}=\mathfrak J^n\mathcal L_{n+1}$; since $\rho_{n',n+1}$ is onto with kernel $\mathfrak J^{n+1}\mathcal L_{n'}\subseteq\mathfrak J^n\mathcal L_{n'}$, this preimage is $\mathfrak J^n\mathcal L_{n'}$). Each $\mathcal L_n$ is a coherent $\mathcal O_{Y_n}$-module. By (9.6b), $\mathfrak L=\varprojlim\mathcal L_n$ is a coherent $\mathcal O_{\mathfrak X}$-module and $\mathfrak L/\mathfrak J^n\mathfrak L\cong\mathcal L_n$ for all $n$, compatibly with the maps.

\emph{$\mathfrak L$ is locally free of rank $1$.} Let $x\in\mathfrak X$; take an open affine formal subscheme $\mathfrak U=\widehat U\ni x$, $U=\operatorname{Spec}A$, $\mathfrak J=I^\Delta$ on $\mathfrak U$ and $Y_n\cap\mathfrak U=\operatorname{Spec}A/I^n$ as in the conventions. Since $\mathcal L_1$ is invertible on $Y_1$, there is $g\in A$ with $x\in D=D(g)\cap Y_1$ and $\mathcal L_1|_D\cong\mathcal O_{Y_1}|_D$, via a generator $m_1\in\Gamma(D,\mathcal L_1)$. The key point: $D$ is the same open set for all $n$; it is an affine open $D_n=\operatorname{Spec}(A/I^n)_g$ of $Y_n$. Because $D_{n+1}$ is affine and $\mathcal L_{n+1}\to\mathcal L_n$ is a surjection of quasi-coherent sheaves on $D_{n+1}$ ($\mathcal L_n$ viewed on $Y_{n+1}$), sections lift: choose successively $m_{n+1}\in\Gamma(D,\mathcal L_{n+1})$ with $\rho_{n+1,n}(m_{n+1})=m_n$. Consider $\mu_n:\mathcal O_{Y_n}|_D\to\mathcal L_n|_D$, $1\mapsto m_n$. At each $y\in D$ the stalk $\mathcal L_{n,y}$ modulo $\mathfrak J_y$ is $\mathcal L_{1,y}$, generated by $m_{1,y}$, and $\mathfrak J_y$ is nilpotent in $\mathcal O_{Y_n,y}$, so by Nakayama $m_{n,y}$ generates $\mathcal L_{n,y}$. Thus $\mu_n$ is a surjection from $\mathcal O_{Y_n}|_D$ onto an invertible sheaf, hence an isomorphism (stalkwise a surjective map $\mathcal O_{Y_n,y}\to\mathcal O_{Y_n,y}$ of finitely generated modules is bijective). The $\mu_n$ are compatible with the transition maps, because $\rho_{n+1,n}(m_{n+1})=m_n$. Passing to the inverse limit, $\mathcal O_{\mathfrak X}|_D=\varprojlim\mathcal O_{Y_n}|_D\xrightarrow{\ \cong\ }\varprojlim\mathcal L_n|_D=\mathfrak L|_D$. Thus $\mathfrak L$ is free of rank $1$ on the neighbourhood $D$ of $x$ (open in $\mathfrak X$, equal to the formal scheme $\widehat{D(g)}$). So $\mathfrak L\in\operatorname{Pic}\mathfrak X$ maps to $([\mathfrak L/\mathfrak J^n\mathfrak L])=([\mathcal L_n])$. Hence $\operatorname{Pic}\mathfrak X\to\varprojlim\operatorname{Pic}Y_n$ is surjective, and by (b) it is an isomorphism: $\operatorname{Pic}\mathfrak X=\varprojlim\operatorname{Pic}Y_n$ (under (ML) for $(R_n)$, which was used in (a) and (b)).

(d) \emph{$\mathfrak X$ affine.} Then $\mathfrak X=\widehat X$, $X=\operatorname{Spec}A$, and with $\mathfrak J=I^\Delta$ we have $Y_n=\operatorname{Spec}A/I^n$ and $R_n=A/I^n$ (this is (9.4a) and its proof). The transition maps $A/I^{n+1}\to A/I^n$ are surjective, so (ML) holds (Example 9.1.1). (For another ideal of definition $\mathfrak J'$ the inverse systems $(Y_n)$ and $(Y'_n)$ are cofinal by (9.5a), so the systems $(R_n)$ and $(R_n')$ are cofinal and (ML) for one gives (ML) for the other; alternatively $\Gamma(Y'_n,\mathcal O)=\Gamma(\mathfrak X,\mathcal O)/\Gamma(\mathfrak X,\mathfrak J'^n)$ is a quotient by Ex.~9.3, again with surjective transition maps.)

\emph{Each $Y_n$ projective over a field $k$} (compatibly, i.e.\ the closed immersions $Y_n\hookrightarrow Y_{n+1}$ are $k$-morphisms, so the maps $R_{n+1}\to R_n$ are $k$-algebra homomorphisms). By (II, Thm.~5.19), $R_n=\Gamma(Y_n,\mathcal O_{Y_n})$ is a finite-dimensional $k$-vector space. By Example 9.1.2, an inverse system of finite-dimensional vector spaces satisfies (ML): for fixed $n$ the decreasing chain of subspaces $\operatorname{im}(R_{n'}\to R_n)$ of the finite-dimensional space $R_n$ is stationary.\hfill$\square$

\end{document}
